\documentclass[12pt,twoside,spanish]{book}

\usepackage[spanish]{babel}
\usepackage{csquotes}

\usepackage[backend=biber]{biblatex}
%\usepackage[backend=biber,style=apa]{biblatex}
\addbibresource{bib_file.bib}
\bibliography{bib_file}

\usepackage{fontspec}
\input{esConfigDir/esConfigFonts}
\usepackage{authblk}
\DefineBibliographyStrings{spanish}{andothers={\itshape{et~al}\adddot}}
\makeatletter
%\renewcommand*{\apablx@ifrevnameappcomma}[2]{#2}
\makeatother
\DefineBibliographyExtras{spanish}{\let\finalandcomma=\empty}
\DeclareDelimFormat[bib,textcite]{finalnamedelim}{
 \ifnumgreater{\value{liststop}}{2}{}{}\addspace\bibstring{and}\space}
\DeclareDelimFormat[parencite]{finalnamedelim}{\addspace{y}\space}
\usepackage[letterspace=5]{microtype}
\usepackage{graphicx}
\def\texttitulocorto{Problemas resueltos de Integración Múltiple}
\def\loauthors{J. N. Zambrano, E. Baquero}

%% Estilo LaTeX Espacios Editorial UD
%% Autores: Editores UDFJC
%% Versión:  v0.1 30MAR2024
%% Compilar con LuaLaTeX, i.e: lualatex main.tex
\usepackage[spanish]{babel}
\usepackage{csquotes}
\usepackage{caption}
\usepackage{titlesec}
\usepackage{titletoc}
\usepackage[titles]{tocloft}

%\setlength{\parindent}{1.25cm}
%\setlength\parindent{0pt}

\usepackage{xcolor}
\definecolor{MSBlack}{rgb}{0.0, 0.0, 0.0}
%\definecolor{MSBLack}{rgb}{.16,.16,.16}
\definecolor{MSLightBlack}{rgb}{.10,.10,.10}
\definecolor{MSCustomGray}{RGB}{144, 147, 150}
\definecolor{MSCustomBlackT}{RGB}{20, 20, 20}
\definecolor{MSCustomBlackST}{RGB}{58, 59, 66}

\newcommand\fontstit{\medseries}
\newcommand\fontschap{\regseries}
\newcommand\fontssec{\regseries}
\newcommand\fontssubsec{\regseries}
\newcommand\fontssubsubsec{\regseries}
\newcommand\fontsparagraph{\regseries\itshape}

% Small caps
\newcommand{\smallcaps}[1]{\textls[25]{\textsc{\MakeLowercase{#1}}}}

%configuracion para titulos de partes, capitulos, secciones, subsecciones, subsubsecciones
\newcommand{\chapterlabel}{\regseries Capítulo~\thechapter.}

% \titleformat{<command>}[<shape>]{<format>}{<label>}{<sep>}{<before-code>}[<after-code>]

\titleformat{\part}[display]{\color{MSBlack}\fontsize{25pt}{2.0mm}\selectfont\fontstit}{\thepart}{12pt}{\raggedright}
\titleformat{\chapter}[block]{\color{MSBlack}\fontsize{25pt}{2.5mm}\selectfont\fontschap}{\chapterlabel}{0.3em}{\raggedright}
\titleformat{\section}[block]{\color{MSBlack}\fontsize{16pt}{2.0mm}\selectfont\fontssec}{\thesection}{0.5em}{\raggedright}
\titleformat{\subsection}[block]{\color{MSBlack}\fontsize{13pt}{2.0mm}\selectfont\fontssubsec}{\thesubsection}{0.5em}{\raggedright}
\titleformat{\subsubsection}[block]{\color{black}\fontsize{12pt}{2.0mm}\selectfont\fontssubsubsec}{\thesubsubsection}{0.5em}{\raggedright}
\titleformat{\paragraph}[block]{\color{black}\fontsize{12pt}{2.0mm}\selectfont\fontsparagraph}{\theparagraph}{0.5em}{\raggedright}

%Espacios para titulos
\titlespacing{\part}{0cm}{10cm}{2.8cm}
\titlespacing{\chapter}{0cm}{1cm}{2.8cm}
\titlespacing{\section}{0cm}{7mm}{2mm}
\titlespacing{\subsection}{0cm}{5mm}{2mm} 
\titlespacing{\subsubsection}{0cm}{5mm}{0mm}
\titlespacing{\paragraph}{0cm}{5mm}{0mm}

% Configuración para apéndices
\makeatletter
\g@addto@macro{\appendix}{%
\titleformat{\chapter}[block]{\color{MSBlack}\fontsize{25pt}{2.5mm}\selectfont\fontschap}{Apéndice~\thechapter.}{0.3em}{\raggedright}
  \titlecontents{chapter}[0em]
    {\vspace{1.5mm}\medseries\normalsize}
    {Apéndice~\thecontentslabel.\enspace}
    {}
    {\hfill\small\lhseries\contentspage\vspace{2.5mm}}[]
}
\makeatother

%configuracion tabla de contenido
\titlecontents{part}[0em]
{\bf \sbseries\normalsize}%
{}% numbered sections formatting
{}% unnumbered sections formatting
{\titlerule*[1pc]{}\small\lhseries\contentspage \vspace{1.0mm}}[]%

\titlecontents{chapter}[0em]
{\vspace{1.5mm}\medseries\normalsize}%
{Capítulo~\thecontentslabel.\enspace}% numbered sections formatting
{}% unnumbered sections formatting
{\titlerule*[1pc]{}\footnotesize\lhseries\contentspage \vspace{1mm}}[]%

\titlecontents{section}[1.5em]
{\small\lhseries}%
{}% numbered sections formatting
{}% unnumbered sections formatting
{\titlerule*[1pc]{}\footnotesize\lhseries\contentspage \vspace{0.5mm}}[]%

\titlecontents{subsection}[3em]
{\lhseries}%
{}% numbered sections formatting
{}% unnumbered sections formatting
{\titlerule*[1pc]{}\small\lhseries\contentspage \vspace{0.5mm}}[]%

\usepackage{setspace}
\linespread{1.1}
\usepackage[skip=3.5pt, indent=20pt]{parskip}

\usepackage{fancyvrb}
\usepackage{fancyhdr}

\fancypagestyle{esstyle}{%
\fancyhead{}
  \fancyfoot{}
  \renewcommand{\headrulewidth}{0pt}
%   \newcommand{\runningtitle}{\fontsparagraph \scriptsize \texttitulocorto}
%  \fancyhf[HLE]{\lhseries\scriptsize Author's Name}
  \fancyhf[HLE]{\lhseries\scriptsize \loauthors}
%  \fancyhf[HRO]{\runningtitle}
  \fancyhf[HRO]{\fontsparagraph \scriptsize \texttitulocorto}
  \fancyfoot[LE]{{\fontsize{10.5}{1}\usefont{T1}{Roboto-TLF}{thin}{}\selectfont{E}}\hspace{-0.5mm}{\footnotesize\usefont{T1}{Inter-Sup}{thin}{}\selectfont{\reflectbox{S}}}{\usefont{T1}{Roboto-TLF}{thin}{}\selectfont{\thinspace|\thinspace}}{\lhseries\scriptsize\thepage}}

  \fancyfoot[RO]{{\lhseries\scriptsize\thepage}{\usefont{T1}{Roboto-TLF}{thin}{}\selectfont{\thinspace|\thinspace}}{\fontsize{10.5}{1}\usefont{T1}{Roboto-TLF}{thin}{}\selectfont{E}}\hspace{-0.5mm}{\footnotesize\usefont{T1}{Inter-Sup}{thin}{}\selectfont{\reflectbox{S}}}}
  \fancyfoot[C]{}
}

%https://www.latex-project.org/help/documentation/fntguide.pdf

\fancypagestyle{plain}{%
  \fancyhead{}
  \fancyfoot{}
  \renewcommand{\headrulewidth}{0pt}
  \fancyfoot[LE]{{\fontsize{10.5}{1}\usefont{T1}{Roboto-TLF}{thin}{}\selectfont{E}}\hspace{-0.5mm}{\footnotesize\usefont{T1}{Inter-Sup}{thin}{}\selectfont{\reflectbox{S}}}{\usefont{T1}{Roboto-TLF}{thin}{}\selectfont{\thinspace|\thinspace}}{\lhseries\scriptsize\thepage}}

  \fancyfoot[RO]{{\lhseries\scriptsize\thepage}{\usefont{T1}{Roboto-TLF}{thin}{}\selectfont{\thinspace|\thinspace}}{\fontsize{10.5}{1}\usefont{T1}{Roboto-TLF}{thin}{}\selectfont{E}}\hspace{-0.5mm}{\footnotesize\usefont{T1}{Inter-Sup}{thin}{}\selectfont{\reflectbox{S}}}}
  \fancyfoot[C]{}
}

\assignpagestyle{\chapter}{plain}

%\setcounter{table}{0}
\counterwithout{table}{chapter}
\counterwithout{figure}{chapter}
\counterwithout{equation}{chapter}
\captionsetup[figure]{labelfont=bf, labelsep=period, font=footnotesize, justification=centering, width = 0.8\textwidth}
\captionsetup[table]{labelfont=bf, labelsep=period, font=footnotesize, justification=centering, width = 0.8\textwidth}

%footnote format
\makeatletter
\renewcommand{\@makefntext}[1]{%
  \begin{list}{}{%
    \setlength{\labelwidth}{1.5em}%
    \setlength{\leftmargin}{\labelwidth}%
    \setlength{\labelsep}{3pt}%
    \setlength{\itemsep}{0pt}%
    \setlength{\partopsep}{0pt}%
    \setlength{\topsep}{0pt}%
    \footnotesize}%
  \item[\@makefnmark\hfil]#1%
  \end{list}%
}
\makeatother

\newcommand{\nota}[1]{\parbox{10cm}{\centering\scriptsize\color{MSBlack} {\bf Nota: }{#1}}}
\newcommand{\fuente}[1]{\parbox{10cm}{\centering\scriptsize\color{MSBlack} {\bf Fuente: }{#1}}}

\usepackage{wrapfig}

\newcommand{\espart}[6]{\part[Parte #1]{\protect\thispagestyle{empty}
%{
%\hspace{-2.75cm}\centering\includegraphics[width=17cm, height=8cm]{parte1_img}
%\begin{figure}[!h] %this figure will be at the right
\vspace{-14.0cm}
   \begin{figure}[ht!]\hspace{-2.9cm}\includegraphics[width=1.045\paperwidth, height=0.5\paperheight]{#2}\end{figure}
%\end{figure}
\flushleft
  \medseries{\Large\noindent Parte {#3}.}
  ~\\
  ~\\
  ~\\
  ~\\
  \lhseries{#4}
  ~\\
  ~\\
  ~\\
  ~\\
  \lhseries{\Large #5}
  ~\\
  ~\\
  ~\\
  ~\\
  \lhseries{\normalsize #6}
}}

%%%%%%%%%%%%%%%%%%


%%%%%%%%%% PORTADILLA 1 %%%%%%%%%%%%%%%
\newcommand{\portadillauno}[2]{
\vspace*{6.5cm}
\begin{flushleft}
  \linespread{3}
%  \color{MSCustomBlackT}\lsstyle
  \color{MSBlack}\lsstyle
  \fontsize{22pt}{3.0mm} \lhseries

  \leftskip=2mm
  #1   
  \vspace{2mm}
  \leftskip=-5mm

  %{\color{white}\colorbox{MSCustomGray}
    %{\parbox[][1.5cm][c]{15.2cm}{
        %\parbox{12.5cm}{\leftskip=7mm \fontsize{17pt}{2.5mm} \lhseries \vspace{1.0mm}
%          Consed ut escimos num sequid ut aut mos {\break} excerro et et, tempos quam, optat fuga.\vspace{0.5mm}
          %#2 \vspace{0.5mm}
        %}
      %}
    %}
  %}

\end{flushleft}}
%%%%%%%%%% FIN PORTADILLA 1 %%%%%%%%%%%%%%%

%%%%%%%%%% PORTADILLA 2 %%%%%%%%%%%%%%%
\newcommand{\portadillados}[2]{
\vspace*{6.5cm}
\begin{flushleft}
  \linespread{3}
%  \color{MSCustomBlackT}\lsstyle
  \color{MSBlack}\lsstyle
  \fontsize{22pt}{3.0mm} \lhseries

  \leftskip=2mm
  #1
   
  \vspace{2mm}
  \leftskip=-5mm

  {\parbox[][1.5cm][c]{15.2cm}{
      \parbox{12.5cm}{\leftskip=7mm \fontsize{17pt}{2.5mm} \lhseries \vspace{1.0mm}
%          Consed ut escimos num sequid ut aut mos {\break} excerro et et, tempos quam, optat fuga.\vspace{0.5mm}
          #2 \vspace{0.5mm}
      }
    }
  }
\end{flushleft}}
%%%%%%%%%% PORTADILLA 2 %%%%%%%%%%%%%%%

\makeatletter
\newcommand*\reftoanex[1]{\def\@currentlabelname{#1}}
\makeatother

%% COLOFON ESTILOS
\definecolor{MSGreyCOLOFON}{rgb}{.35,.35,.35}
\definecolor{MSGreyLATEX}{rgb}{.55,.55,.55}

%\usepackage{libertine}
\DeclareRobustCommand\myLaTeX{%
    L\kern-.3em%
    \raisebox{.4ex}{\scshape a}\kern-.13em%
    T\kern-.2em%
    \raisebox{-.5ex}{E}\kern-.13em%
    X%
}

\newcommand{\escolofon}[2]{
\thispagestyle{empty}
\vspace*{10.0cm}
\begin{tikzpicture}[remember picture,overlay]
  \node[text=gray,rotate=0,anchor=east, yshift=5mm] at (15.5,11.5){
    \includegraphics[width=5.0cm, height=4.0cm]{esIcons/logoESColofon_corner.png}
  };
\end{tikzpicture}

\vspace{-4.5cm}
\begin{center}
{\fontsize{7pt}{7pt}\selectfont
\color{MSGreyCOLOFON}
Este libro fue editado y diagramado en} {\fontsize{13pt}{13pt}\selectfont\color{MSGreyLATEX}{\myLaTeX}}
\vspace{3.5cm}

\includegraphics[scale=0.05]{esIcons/logoESColofon.png}
  
  \vspace{0.5cm}
  {\scriptsize
    \color{MSGreyCOLOFON}Este libro se término de imprimir en #1 de #2\\ en la Editorial UD. Bogotá, Colombia.}
\end{center}
}
%%%%%%%%%%%%%%%%%%

\usepackage[paperheight=24.0cm,paperwidth=17.0cm,textwidth=12.5cm,tmargin=2.5cm,bmargin=2.5cm,left=2.5cm,right=2.0cm]{geometry}

\decimalpoint
\usepackage{lipsum}  
\usepackage[pangram]{blindtext}

\usepackage{amsmath}
\usepackage{longtable}
\usepackage{subcaption}
\usepackage{array}
\newcolumntype{P}[1]{>{\centering\arraybackslash}p{#1}}

%%% inicio paquetes
%\usepackage[affil-it]{authblk}
\usepackage{amsmath,amsthm,array,graphicx,comment,ifthen,amsfonts}
\usepackage{graphicx,amsmath,multicol,pst-coil,bigints,caption}
\usepackage{pstricks-add,pst-plot,pst-node,pst-3dplot,hyperref,appendix}
\usepackage{enumitem,bm,calculus,calculator}
\renewcommand{\proofname}{Demostración}

\usepackage{makeidx}
\makeindex

\usepackage{tcolorbox} % hacer cajas de color gris
\usepackage{makeidx}
buselibrary{breakable,skins}

\newcommand{}{bset{colback=gray!10!white,colframe=black!85!blue}
\begin{tcolorbox}[breakable,enhanced,title=\bf Parametrizando el sólido]}
\newcommand{\ptr}{bset{colback=gray!10!white,colframe=black!85!blue}
\begin{tcolorbox}[breakable,enhanced,title=\bf Parametrizando la región de Integración]}
\newcommand{}{\begin{tcolorbox}[breakable,enhanced,title=\bf Instrucción en {\tt Mathematica}]}
\newcommand{}{\end{tcolorbox}}

\author{\Large \textit{\textbf{Juan Nepomuceno Zambrano Caviedes}} \\ \textit{\textbf{Efrén Ricardo Baquero Torres}}}
\date{}

\usepackage{fancyhdr}
\fancyhf{}
\fancyhead{} 
\makeatletter

%\renewcommand{\sectionmark}[1]{\markright{\thesection~~#1}}
%\renewcommand{\chaptermark}[1]{\markboth{\if@mainmatter\chaptername\ \thechapter~~ - \fi#1}{}}
\makeatother
\fancyhead[LE]{\thepage}
\fancyhead[RE]{\itshape\nouppercase  \leftmark}
\fancyhead[LO]{\itshape\nouppercase  \rightmark}
\fancyhead[RO]{\thepage}
\fancyfoot[]{} % \fancyfoot[R]{\thepage}
\renewcommand{\headrulewidth}{0.4pt} % Default \headrulewidth is 0.4pt
\renewcommand{\footrulewidth}{0.4pt}

\addto{\captionsspanish}{\renewcommand{\contentsname}{Contenido}}
\addto{\captionsspanish}{\renewcommand{\bibname}{Referencias}}
\psset{showorigin=false}

\usepackage[hyphenbreaks]{breakurl}
%\usepackage[hyphens]{url}
%\renewcommand{\footrulewidth}{1pt}
%\fancyfoot[CE,CO]{Material en preparación - Efrén Ricardo Baquero}
%\fancyfoot[LE,RO]{}

\newtheorem{teorema}{Teorema}[chapter]
\newtheorem{prop}{Proposición}[chapter]
\newtheorem{defi}{Definición}[chapter]
\newtheorem{ejer}{Ejercicio}[chapter]
\newtheorem{ejemplo}{Ejemplo}[chapter]
\newcommand{\sech}{{\rm sech\,}}
\newcommand{\csch}{{\rm csch\,}}
\newcommand{\arcsinh}{{\rm arcsinh\,}}
\newcommand{\arccosh}{{\rm arccosh\,}}
\newcommand{\arctanh}{{\rm arctanh\,}}
\newcommand{\arccoth}{{\rm arccoth\,}}
\newcommand{\arcsech}{{\rm arcsech\,}}
\newcommand{\arccsch}{{\rm arccsch\,}}
\newcommand{\cn}[3]{\pscircle(#1,#2){5pt}\rput(#1,#2){$#3$}}
\newcommand{\fig}{\begin{figure}}
\newcommand{\ef}[1]{\caption{#1} \end{figure}}
\newcommand{\rot}[1]{{\rm rot}\,{\bf #1}}
\newcommand{\la}{\left \langle}
\newcommand{\ra}{\right \rangle}
\newcommand{\vi}{\boldsymbol{\hat{\imath}}}
\newcommand{\vj}{\boldsymbol{\hat{\jmath}}}
\newcommand{\vk}{\boldsymbol{\hat{k}}}
\newcommand{\ei}[1]{{\bf e}_{#1}}
\newcommand{\ve}[1]{\overrightarrow{#1}}
\newcommand{\bi}{\bigintsss}
\newcommand{\ds}{\displaystyle}
\newcommand{\mc}[1]{\begin{multicols}{#1}}
\newcommand{\emc}{\end{multicols}}
% \usepackage[Lenny]{fncychap}
% \usepackage[Glenn]{fncychap}
%\usepackage[Sonny]{fncychap}
% \ChNameVar{\Large\rm} \ChNumVar{\Huge} \ChTitleVar{\Large\rm}
% \ChRuleWidth{0.3pt} \ChNameUpperCase

\setcounter{tocdepth}{1}
\pagestyle{empty}
\usepackage{fancyvrb}
\newcommand{\letraverbatim}{\fontsize{9}{11.5}\selectfont}
\newenvironment{Figure}
{\par\medskip\noindent\minipage{\linewidth}}
{\endminipage\par\medskip}

\hyphenation{car-dioi-di-co instrucciones}
%%% fin paquetes

\usepackage{shapepar}
%\usepackage{tikz}
%\usepackage{arydshln}
\usepackage{colortbl}

%\usepackage[usenames,dvipsnames]{color}
%\usepackage[dvipsnames]{xcolor}
\usepackage{xcolor}
\definecolor{orcidlogocol}{HTML}{A6CE39}

\usepackage{tabularray}
%\usepackage{orcidlink}
%\usepackage[toc,page]{appendix}
\usepackage{nameref}
\usepackage{fontawesome}
\usepackage{academicons}

\usepackage[axes,cam,width=190truemm,height=260truemm,pdftex,center]{crop}

\setlength\parindent{0pt}

\pagestyle{empty}

% \hyphenation{este sencillo cualquiera segmento segmentos superficie siguientes mientras volumen formando obtiene derivando requiere integrando expresiones para menor considerar integrales iniciando sustituye donde paraboloide porciones simplifica simplificando parametrizada parametriza parametrizan como mismo cilindro cilindros planteamiento circunferencia circunferencias forma determinan recorre interior corresponde correspondientes presentan representan coordenada  coordenadas coordenados sobre determina puntos tanto sentencias ayuda anteriormente estos convexa funciones reescribe establecer naturaleza permite primer ejecutando completo corresponden elemento debe mediante proporciona adecuadamente respectivas punto desde positiva evaluando oeden parte radio resultado recorrer centrada advertir suponiendo vector hiperboloide desigualdades previamente dentro puede hacen diferencial sintaxis vectorial denominaremos denominada diferenciable primeros}
	
\begin{document}

%%%%%%%% EDITAR INICIO PORTADILLA 1 %%%%%%%%%%%%%%%
\portadillauno{Problemas resueltos de Integración Múltiple}{con aplicaciones en {\tt Wolfram Mathematica}}
%%%%%%%% FIN PORTADILLA 1 %%%%%%%%%%%%%%%%%%%%%%%%%
\newpage
\newpage

%%%%%%%% EDITAR INICIO PORTADILLA 2 %%%%%%%%%%%%%%%
\portadillados{Problemas resueltos de Integración Múltiple}{con aplicaciones en {\tt Wolfram Mathematica}}
\vspace*{5.5mm}
\begin{flushleft}
  \leftskip=2mm
      {\color{MSBlack}\lsstyle \fontsize{12pt}{6.0mm} \exlseries Juan Nepomuceno Zambrano Caviedes \linebreak 
      Efrén Ricardo Baquero Torres
      }
\end{flushleft}

\begin{center}
  \vspace{3.5cm}
  \includegraphics[scale=0.06]{esIcons/logoES.png}
\end{center}
%%%%%%%% FIN PORTADILLA 2 %%%%%%%%%%%%%%%%%%%%%%%%%

\newpage
\newpage

%%%%%%%% EDITAR INICIO RESUMENES EN ESPAÑOL E INGLÉS %%%%%%%%%%%%%%%
\chapter*{{\medseries \large Resumen}}
\thispagestyle{empty}
\vspace*{-1.4cm}
\noindent \thispagestyle{empty}
\section*{\fontschap Resumen \footnote{¿Cómo citar este libro? / How to cite this book?
		Zambrano Caviedes, J. y Baquero Torres, E. (2025). Problemas resueltos de integración múltiple con aplicaciones en Wolfram Mathematica. Universidad Distrital Francisco José de Caldas-Editorial UD. DOI: XXXX}}

Presentamos en estas notas de clase, algunos problemas 
rutinarios sobre integración para un curso de Cálculo Multivariado. En ellas se trabaja, básicamente, la integración de una función sobre un objeto que puede ser una curva, superficie o sólido, con el fin de parametrizar la región de integración.
Esta idea es diferenciadora con respecto a la mayoría de textos de cálculo multivariado. Cada problema ofrece un gráfico asociado en diversas vistas y el cálculo que se requiere. Además se incluyen las instrucciones utilizadas en el programa  \emph{Wolfram Mathematica}. El cambio de sistema de coordenadas para evaluar integrales triples, se hace tradicionalmente utilizando coordenadas esféricas y cilindricas; no obstante, se incluyen otros sistemas de coordenadas en los anexos y una breve referencia en los ejercicios propuestos para motivar al lector al estudio de estos tópicos.

\fontschap Palabras clave:\normalfont integración múltiple, parametrización, campos escalares, campos vectoriales, Wolfram Mathematica.
\vspace{-0.7em} % Espacio antes de la línea

\noindent
\hspace*{-0.1cm} % Ajusta este valor para mover la línea hacia la derecha
\rule{2cm}{0.5pt} % Línea horizontal

\vspace{-0.5em} % Espacio después de la línea



\section*{\fontschap Abstract}

We present in these class notes some problems 
routines on integration for a multivariate calculus course. In them, we work, basically, integrating a function on an object that can be a curve, surface or solid, parameterizing the region of integration.
This idea is different from most multivariate calculus texts. Each problem offers an associated graph in various views and the calculation that is required, the instructions used in the {\tt Wolfram Mathematica} program are also included. Changing the coordinate system to evaluate triple integrals is traditionally done using spherical and cylindrical coordinates. However, other coordinate systems are included in the annexes and a brief reference in the proposed exercises to motivate the reader to study these topics.

\fontschap Keywords:\normalfont multiple integration, parameterization, scalar fields, vector fields, Wolfram Mathematica.

\begin{spacing}{1.4}
	\begin{tabular}{p{0.26\textwidth}!{\vrule width 0.9pt}p{0.69\textwidth}}
		& \hspace{2mm}{\medseries \fontsize{11pt}{9pt}\selectfont{¿Como citar este libro?~/ How to cite this book?}} \\[3.5pt]
		& \hangindent=2mm \hangafter=0
		\fontsize{9pt}{9pt}\selectfont{Salamanca Bernal, J. A., Rodríguez Patarroyo, D. J. y Pantoja Benavides, J. F. (2025). \emph{Columna total de ozono estratosférico y manchas solares. Una relación físico-estadística mundial desde la minería de datos geoespaciales NASA}. Universidad Distrital Francisco José de Caldas-Editorial UD. https://doi.org/10.69740/UD.9789587877724.9789587877748}\\
	\end{tabular}
\end{spacing}
\vskip1cm
{\medseries Palabras clave:} Integración múltiple, Parametrización, campos escalares, campos vectoriales, Wolfram Mathematica.

\vfill
\begin{spacing}{1.4}
  \begin{tabular}{p{0.26\textwidth}!{\vrule width 0.9pt}p{0.69\textwidth}}
    & \hspace{2mm}{\medseries \large ¿Como citar este libro?} \\[3.5pt]
    & \hangindent=2mm \hangafter=0
    \fontsize{9pt}{9pt}\selectfont{Zambrano Caviedes, J.N. \& Baquero Torres, E. R. (2025).\linebreak \emph{Problemas resueltos de Integración Múltiple con aplicaciones en {\tt Wolfram Mathematica}}. Universidad Distrital Francisco José de Caldas-Editorial UD. \url{}}\\
  \end{tabular}
\end{spacing}

\newpage

\selectlanguage{english}
\chapter*{{\medseries \large Abstract}}
\thispagestyle{empty}
\vspace*{-1.4cm}
\noindent \chapter*{Abstract}

We present in these class notes some problems 
routines on integration for a multivariate calculus course. In them, we work, basically, integrating a function on an object that can be a curve, surface or solid, parameterizing the region of integration.
This idea is different from most multivariate calculus texts. Each problem offers an associated graph in various views and the calculation that is required, the instructions used in the {\tt Wolfram Mathematica} program are also included. 
\vskip0.2cm
Changing the coordinate system to evaluate triple integrals is traditionally done using spherical and cylindrical coordinates; However, other coordinate systems are included in the annexes and a brief reference in the proposed exercises to motivate the reader to study these topics.

\textbf{Keywords:} multiple integration, parameterization, scalar fields, vector fields, Wolfram Mathematica.

\selectlanguage{spanish}

\vskip1cm
{\medseries Keywords:} Multiple integration, Parameterization, scalar fields, vector fields, Wolfram Mathematica.
\selectlanguage{spanish}

\vfill
\begin{spacing}{1.4}
  \begin{tabular}{p{0.26\textwidth}!{\vrule width 0.9pt}p{0.69\textwidth}}
    & \hspace{2mm}{\medseries \large How to cite this book?} \\[3.5pt]
    & \hangindent=2mm \hangafter=0
    \fontsize{9pt}{9pt}\selectfont{Zambrano Caviedes, J.N. \& Baquero Torres, E. R. (2025). \linebreak \emph{Problemas resueltos de Integración Múltiple con aplicaciones en {\tt Wolfram Mathematica}}. Universidad Distrital Francisco José de Caldas-Editorial UD. \url{}}\\
  \end{tabular}
\end{spacing}
%%%%%%%% FIN RESUMENES EN ESPAÑOL E INGLÉS %%%%%%%%%%%%%%%

\newpage
%%%%%%% EDITAR DEDICATORIA O FRASE CELEBRE %%%%%%%%%%%%%%%%%%%%%%
\vspace*{0.0cm}
\emph{
\begin{flushright}
Las matemáticas son la música de la razón.\vskip0.2cm 
James Joseph Sylvester
\end{flushright}
}
%%%%%%% FIN DEDICATORIA O FRASE CELEBRE %%%%%%%%%%%%%%%%%%

\newpage

\renewcommand*\contentsname{\lhseries Contenido}
\renewcommand{\cfttoctitlefont}{\hfill\Huge}
\renewcommand{\cftdot}{}
\renewcommand*{\tablename}{\bf Tabla} %"Cuadro" por "Tabla"
\renewcommand*{\figurename}{\bf Figura} %"Cuadro" por "Tabla"
\renewcommand*{\bibname}{\fontschap Referencias}

%% Editar re-definiciones de comandos para el usuario%%
%% panga aquí sus definiciones
%%
%%
%%%%%%%%%%%%%%

\tableofcontents
\addtocontents{toc}{\protect\thispagestyle{empty}}
\newpage

\pagestyle{esstyle}

%%%%%%% EDITAR AGRADECIMIENTOS %%%%%%%%%%%%%%%%%%%%%%%%%%%%
\chapter*{\lhseries Agradecimientos}
\addcontentsline{toc}{chapter}{Agradecimientos}
\vspace*{-1.4cm}
\chapteralt{Agradecimientos}
\addcontentsline{toc}{chapter}{Agradecimientos}

Extendemos un agradecimiento especial al grupo de profesores y estudiantes
de los cursos Cálculo Multivariado y Cálculo Vectorial de la Universidad Distrital Francisco José de Caldas y la Universidad Escuela Colombiana de Ingeniería Julio Garavito, por sus valiosos aportes y observaciones.


\begin{flushright}
  {\bfseries  Juan N. Zambrano Caviedes \linebreak 
  	Efrén Baquero T.}
\end{flushright}
%%%%%%% FIN AGRADECIMIENTOS %%%%%%%%%%%%%%%%%%%%%%%%%%%%%%%
\newpage

\chapter*{Introducción}
\addcontentsline{toc}{chapter}{Introducción}
\chapter*{Introducción}
\addcontentsline{toc}{chapter}{Introducción}

La Universidad Distrital Francisco José de Caldas ha procurado que sus estudiantes movilicen, entre otras competencias,
el uso de la tecnología. Además, en el proyecto curricular de Ingeniería Eléctrica de la Facultad Tecnológica, algunos documentos institucionales señalan que se propone la implementación de un laboratorio de cómputo especializado en matemáticas, el cual contará con paquetes tales
como Matlab, Mathematica y \emph{software} libre.
También se busca que el estudiante, en su proceso de aprendizaje autónomo, refuerce los conocimientos vistos en clase.
Pensando en ello, se presentan estas notas de clase que contienen ejercicios
rutinarios sobre integración para un curso de Cálculo Multivariado, en ellas se pretende trabajar básicamente tres elementos:
\begin{itemize}
  \item Integrar una función sobre un objeto: curva, superficie o sólido, parametrizando la región de integración.
    Esta idea es diferenciadora con respecto a la mayoría de textos de cálculo multivariado, pues la evaluación de integrales triples parametrizando el sólido es una estrategia muy poco utilizada, una ventaja es la obtención de límites constantes de integración.
  \item Se usan herramientas tecnológicas para abreviar el tiempo de ejecución, en particular, se ofrecen las instrucciones utilizadas con el asistente computacional Wolfram Mathematica para visualizar la región de integración o para la evaluación de integrales en diversas formas; dicha evaluación se hace incluso de manera numérica.
  \item Se ofrecen varias formas de resolver un mismo problema, en particular, se presenta la evaluación de integrales múltiples en diferentes órdenes.
\end{itemize}
Estos apuntes no pretenden, desde ningún punto de vista convertirse en un texto guía para un curso de Cálculo Multivariado.
Se desea hacer una compilación de problemas que ilustren su solución y que puedan servir como un material complementario de estudio.
La parte inicial de este material incluye los elementos teóricos básicos sobre los cuales se desarrollan los problemas.
Se agrega un capítulo con los comandos elementales para el uso del programa Wolfram Mathematica
y en los problemas presentados posteriormente se muestra las sintaxis requerida para utilizarlo.
En el tercer capítulo se introducen algunas parametrizaciones sencillas que se aplicarán posteriormente en los problemas que se
presentan en los capítulos 4, 5 y 6.
La integración de campos vectoriales y los teoremas relacionados se presenta en el capítulo 7.

Cada problema ofrece la figura asociada en diversas vistas y el cálculo que se requiere, se muestra de varias maneras.
Algunas de estas alternativas se ven muy engorrosas, pero es el resultado de la experiencia adquirida a traves de los años
en que los autores han orientado este curso.

El cambio de sistema de coordenadas para evaluar integrales triples, se hace tradicionalmente utilizando coordenadas esféricas y cilíndricas; no obstante, se incluyen otros sistemas de coordenadas en los anexos y una breve referencia en los ejercicios propuestos para motivar al lector al estudio de estos tópicos.

Este documento se editó en \LaTeX, sus más de 700 figuras se elaboraron en \verb"Mathematica" o con ayuda del paquete \verb"PSTricks" y su versión evolucionada \verb"psplotThreeD". La documentación puede consultarse en \textcite{vignault_pst-solides3d_2023}, \textcite{van_zandt_pst-plot_2016} y \textcite{vos_pst-func_2024}.

\begin{flushright}
  \smallcaps{Los autores}
\end{flushright}

\newpage
\chapter{Preliminares} \chapter{Elementos básicos de cálculo multivariado}

Presentamos en este capítulo algunos de los elementos teóricos básicos que se requieren para abordar 
los problemas planteados en los capítulos posteriores.
\vskip 0.1cm
Usaremos el término \emph{conjunto} para referirnos a una colección bien definida de objetos. 
Una \emph{función} $f$ de un conjunto $A$ 
en un conjunto $B$, es una regla que asigna a cada elemento de $A$ un único elemento de $B.$
Seguramente el lector está familiarizado con el conjunto de números reales $\mathbb{R}.$
En estos apuntes consideraremos funciones cuyo dominio es el espacio el espacio $n-$dimensional $\mathbb{R}^n,$
o un subconjunto de este.

\section{Vectores}
En \cite{marsden_calculo_2013} se define "vector como un segmento de recta dirigido que comienza en el origen,
esto es, un segmento de recta con magnitud y dirección específicos con punto inicial en el origen".\\
Los puntos en $\mathbb{R}^n$ se presentan en la forma $P=(a_1,a_2,\cdots, a_n)$ 
y el vector que va desde el origen hasta el punto $P$ se denotará por $\langle a_1,a_2,\cdots, a_n\rangle.$ 
Particularmente a los vectores $\vi=\langle 1,0,0 \rangle,$ $\vj=\langle 0,1,0 \rangle,$ $\vk=\langle 0,0,1 \rangle$ los denominaremos \emph{vectores coordenados} en $\mathbb{R}^3.$ 
\vskip0.2cm
Para dos vectores $\ve{u}=\langle a_1,a_2,\cdots, a_n\ra$ y $\ve{v}=\langle b_1,b_2,\cdots, b_n\ra$
el \emph{producto punto} se denota por $\ve{u}\cdot \ve{v} = a_1b_1+a_2b_2+\cdots+a_nb_n$; 
la \emph{norma} $\ve{u}$ se denota por $\Vert \ve{u} \Vert$ y se define
por  $\Vert \ve{u}\Vert=\sqrt{a_1^2+a_2^2+\cdots+a_n^2};$ nótese que
$\ve{u}\cdot \ve{u}=\Vert \ve{u}\Vert^2.$
El \'{a}ngulo $\theta $ entre los vectores $\ve{u}$ y $\ve{v}$ satisface la relación
$\ve{u} \cdot \ve{v}= \left\|\ve{u}\right\| \left\| \ve{v}\right\| \cos \theta .$
\vskip0.2cm
El \emph{producto cruz} entre dos vectores  $\ve{u}=\langle a_{1},a_{2},a_{3}\ra$ y 
$\ve{v}=\langle b_{1},b_{2},b_{3}\ra$ en $\mathbb{R}^3$ es el vector 
$$\ve{u}\times \ve{v}=\langle a_{2}b_{3}-a_{3}b_{2},a_{3}b_{1}-a_{1}b_{3},a_{1}b_{2}-a_{2}b_{1}\ra,$$
que es ortogonal tanto a $\ve{u}$ como a $\ve{v}.$ 
Adem\'{a}s $\left\Vert \ve{u}\times \ve{v}\right\Vert =\left\Vert \ve{u}\right\Vert\left\Vert \ve{v}\right\Vert\sin\theta;$ estos resultados se pueden consultar en detalle en \textcite{apostol_calculus_1999}.

La siguiente es una conocida forma de evaluar este producto,
$$\ve{u}\times \ve{v}=\left| \begin{array}{ccc} \vi & \vj & \vk \\ a_1 &a_2&a_3 \\b_1 &b_2&b_3 \\ \end{array} \right|  $$
El producto cruz es útil en el desarrollo de algunos de los siguientes resultados.
\subsection*{Proyección ortogonal}
%
La proyecci\'{o}n ortogonal del vector $\ve{v}$ sobre el vector 
$\ve{u},$ se denota por ${\rm Proy}_{\ve{u}}\ve{v}$
y es un múltiplo escalar del vector $\ve{u}.$ Es decir 
${\rm Proy}_{\ve{u}}\ve{v}=k\ve{u}.$
De la figura se observa que 
$$\cos \theta =\frac{\left\| k\ve{u}\right\|}{\left\|
\ve{v}\right\|}=\frac{\left| k\right| \left\| \ve{u} \right\|}{\left\|\ve{v}\right\|}.$$


El valor $\left| k\right| $ ser\'{a} positivo si 
${\rm Proy}_{\ve{u}}\ve{v}$ posee el mismo sentido de $\ve{u},$ de lo
contrario, ser\'{a} negativo. Al igualar las expresiones para $\cos \theta $ y considerando el caso $k>0,$
se concluye que $k=\frac{\ve{u}\cdot \ve{v}}{\left\| \ve{u}\right\| ^{2}},$ entonces 
la proyección del vector $\ve{v}$ sobre $\ve{u}$ está dada por
$${\rm Proy}_{\ve{u}}\ve{v}=\frac{\ve{u}\cdot\ve{v}}{\left\| \ve{u}\right\|^{2}}\ve{u}$$
En el gr\'{a}fico la línea discontinua corresponde al vector 
$\ve{v}-{\rm Proy}_{\ve{u}}\ve{v}$. Aunque nos apoyamos en un gráfico bidimensional, este resultado se puede presentar en espacios 
de dimensión superior.

\vskip0.2cm
\noindent
\begin{minipage}{\textwidth}
	\centering
	\captionof{figure}{Proyección ortogonal}
	\label{fig:proy-ortogonal}
	\begin{minipage}{0.30\linewidth}
		\centering
		\includegraphics[width=\linewidth]{Fig/1.pdf}
	\end{minipage}
	
	\vspace{0.3em}
	\parbox{0.95\linewidth}{\centering\fontsize{8pt}{10pt}\selectfont\textbf{Fuente:} elaboración propia.}
\end{minipage}
\vskip0.2cm


\begin{prop}
Para todo $\ve{v}\in \mathbb{R}^n$, el vector $\ve{v}-{\rm Proy}_{\ve{u}}\ve{v}$ es ortogonal
tanto a $\ve{u}$ como a ${\rm Proy}_{\ve{u}}\ve{v}.$
\end{prop}
\begin{proof}
Veamos que $\ve{v}-{\rm Proy}_{\ve{u}}\ve{v}$ es ortogonal a $\ve{u}$
\begin{eqnarray*}
\ve{u} \cdot \left(\ve{v}-{\rm Proy}_{\ve{u}}\ve{v}\right) &=& \ve{u} \cdot \left(\ve{v} - \frac{\ve{u}\cdot \ve{v}}{\Vert\ve{u}\Vert^2}\ve{u}\right) \\ 
&=& \ve{u}\cdot\ve{v}-\left (\frac{\ve{u}\cdot\ve{v}}{\Vert\ve{u}\Vert^2}\right)\ve{u}\cdot\ve{u} \\
&=& \ve{u}\cdot\ve{v}-\left (\frac{\ve{u}\cdot\ve{v}}{\Vert\ve{u}\Vert^2}\right) \Vert\ve{u}\Vert^2 \\
&=& \ve{u}\cdot\ve{v}-\ve{u}\cdot\ve{v}=0.
\end{eqnarray*}

Como ${\rm Proy}_{\ve{u}}\ve{v}$ es múltiplo escalar de $\ve{u},$ con el resultado anterior se sigue que
$\ve{v}-{\rm Proy}_{\ve{u}}\ve{v}$ es ortogonal a ${\rm Proy}_{\ve{u}}\ve{v}.$
\end{proof}
\subsection*{Área de paralelogramo}
El área $A$, del paralelogramo cuyas aristas son los vectores $\ve{u}$ y $\ve{v}$
está dada por $\Vert \ve{u} \Vert h = \Vert \ve{u} \Vert \Vert \ve{v} \Vert \sin \theta . $
De la relación $\Vert \ve{u} \times \ve{v} \Vert =\Vert \ve{u} \Vert \Vert \ve{v} \Vert \sin \theta$
se sigue que el área del paralelogramo es
$A=\Vert \ve{u} \times \ve{v} \Vert.$


\begin{figure}
    \centering
    \caption{Área de un paralelogramo y volumen de un paralelepípedo}
    \label{paralelogramo_paralelipedo}
    \includegraphics[width=0.38\linewidth]{Fig/2.pdf}
    \hspace{1em}
    \includegraphics[width=0.38\linewidth]{Fig/3.pdf}  
    \parbox{0.95\linewidth}{\centering\fontsize{9pt}{10pt}\selectfont\textbf{Fuente:} elaboración propia.}
\end{figure}

\subsection*{Volumen de un paralelepípedo}
El volumen del paralelepípedo cuyas aristas son los vectores $\ve{u},$ $\ve{v}$ y $\ve{w},$
se determina con el siguiente procedimiento: 
el área de la base del paralelepípedo está dada por $\Vert \ve{u} \times \ve{v}  \Vert$,
la altura $h$ del paralelepípedo corresponde a la norma de la proyección del vector $\ve{w}$
sobre el vector $\ve{u} \times \ve{v}.$
Es decir,
$ h  = \Vert {\rm Proy}_{\ve{u} \times \ve{v}} \ve{w} \Vert =\left\Vert \frac{\ve{w} \cdot (\ve{u} \times \ve{v})}{\Vert\ve{u} \times \ve{v} \Vert^2 }\,(\ve{u} \times \ve{v} )\right \Vert = \frac{|\ve{w} \cdot (\ve{u} \times \ve{v})|}{\Vert\ve{u} \times \ve{v} \Vert }$. 
Por lo anterior, el volumen $V$ del paralelepípedo es $h \Vert\ve{u} \times \ve{v} \Vert$
entonces, 
\begin{equation} \label{preliminares4}
V= |\ve{w} \cdot (\ve{u} \times \ve{v}) |
\end{equation}

\section{Funciones en el espacio} \label{funcionesenelespacio}
Para $a\in \mathbb{R}^{n},$ al conjunto $B_{r}\left(a\right) =\{x\in \mathbb{R}^{n}|\left\Vert x-a\right\Vert <r\}$ 
lo denominaremos la bola con centro en $a$ y radio $r>0.$ \index{Bola} 
Diremos que un conjunto $U\subseteq \mathbb{R}^{n}$ es \emph{abierto} si para todo $u\in U$ 
existe $r>0,$ tal que, $B_{r}(u) \subset U.$  \index{Conjunto!abierto}
Un sinónimo de conjunto abierto es región abierta. \index{Región abierta} 
Un conjunto $B$ se dice \emph{cerrado} si su complemento es un conjunto abierto. \index{Conjunto!cerrado}
\begin{teorema}
$B_{r}(a)$ es un conjunto abierto.
\end{teorema}
La demostración de este resultado se encuentra en \textcite{marsden_calculo_2013}.
\begin{ejemplo}
El conjunto $A=\{(x,y) \in \mathbb{R}^{2}|x>0\}$ es abierto pues
para todo $a \in A$ se satisface la definición dada.
\vskip0.2cm
El conjunto $B=\{(x,y) \in \mathbb{R}^{2}|xy\geq 0\}$ no es abierto pues
para $\left( 0,0\right) \in B$ no es posible determinar
$r>0$ tal que $B_{r}\left( \left( 0,0\right) \right) \subset B.$
\end{ejemplo}
De manera general abordaremos funciones de la forma
\begin{equation}
f: A\subseteq \mathbb{R}^n \to B\subseteq \mathbb{R}^m
\end{equation}
El conjunto de puntos para los cuales la función $f$ está definida se denomina \emph{dominio} de $f$ 
y el conjunto de puntos de $\mathbb{R}^m$ que son imagen de algún punto del dominio se denomina \emph{rango}
de $f$. De acuerdo con \textcite{apostol_calculus_1999}, a las funciones definidas cuando $n=m=1,$ se denominan \emph{funciones reales,} las cuales se estudian en los cursos de \emph{cálculo diferencial} y \emph{cálculo integral.}
Si $n>1$ y $m=1$ obtendremos \emph{campos escalares};  \index{Campos escalares} \index{Campos vectoriales}
llamaremos \emph{funciones vectoriales} a las obtenidas cuando $n=1$ y $m>1$; mientras que los \emph{campos vectoriales} son funciones obtenidas cuando $n>1$ y $m>1.$ \index{Funciones vectoriales}
\vskip0.2cm 
Para presentar funciones vectoriales, campos vectoriales y en general para vectores, usaremos la notación $\ve{r}(t),$ ${\bf r}(t)$ o ${\bf F}(x,y,z)$ indistintamente y según el caso. 
Por ejemplo:

\begin{itemize}
\item $\overrightarrow{r}(t) =\left\langle \ln t,\frac{1}{t},t^{2}\right\rangle ,$ es 
una funci\'{o}n vectorial en $\mathbb{R}^{3}$ con dominio $(0,\infty ).$ 
% El rango es el conjunto $B={\rm Rango}f=\{(x,y,z)\in\mathbb{R}^{3}/ y\neq 0, z \geq 0\}.$
\item El campo escalar definido en $\mathbb{R}^{2}$ por $f(x,y) =\sqrt{9-x^{2}}+\sqrt{y^{2}-4}$ posee
dominio, ${\rm dom}f=\{(x,y) \in \mathbb{R}^{2}/\left\vert x\right\vert \leq 3,\left\vert y\right\vert \geq 2\}$ y el rango de $f$ es el conjunto $[0,\infty).$
\item $\overrightarrow{F}(x,y,z) =\left\langle y,\frac{x}{x-y},e^{-xz}\right\rangle$ 
es un campo vectorial definido en $\mathbb{R}^{3},$ su dominio es el conjunto 
${\rm dom}f=\{(x,y,z) \in \mathbb{R}^{3}/x\neq y\}.$
\end{itemize}
Supondremos que el lector ha estudiado los conceptos de límite, continuidad, diferenciabilidad e integrabilidad
en funciones reles, de tal manera que extenderemos estos elementos a los otros tipos de funciones. 
\vskip0.2cm 
En cuanto al dominio de la funci\'{o}n $f(x,y) =\frac{x^{4}-y^{4}}{x^{2}+y^{2}},$
${\rm Dom}f=\{(x,y) \in \mathbb{R}^{2}/(x,y) \neq \left( 0,0\right)\}$
a pesar que $\left( 0,0\right) $ no es un punto del dominio de $f,$
es natural preguntarnos: ¿qué sucede si nos acercamos al origen ? En la siguiente sección podremos 
responder a este interrogante.

\subsection*{Límites}

\begin{defi} \label{preliminares5}
Si {\small $\ve{F}: A \subseteq\mathbb{R}^n \to  \mathbb{R}^m$} es una función con dominio $A,$
se dice que $\ve{L}$ es el límite de $\ve{F}(x_{1},x_{2},\cdots ,x_{n})$ cuando 
${\bf x}=\left(x_{1},x_{2},\cdots ,x_{n}\right)$ tiende a ${\bf a}= \left( a_{1},a_{2},\cdots,a_{n}\right) ,$ 
lo cual se denota como {\small $ \lim_{{\bf x}\to {\bf a}}\ve{F}({\bf x})=\ve{L}=\langle l_{1},l_{2},\cdots ,l_{m}\ra,$}
si para todo {\small $\varepsilon >0$} existe $\delta >0$ tal que {\small $\left\Vert \ve{F}\left( x_{1},x_{2},\cdots ,x_{n}\right) -\ve{L}\right\Vert <\varepsilon,$}
siempre que $\left\Vert {\bf x}- {\bf a} \right\Vert <\delta .$
\end{defi} 

\begin{itemize}[leftmargin=0cm,itemindent=0.4cm,labelwidth=\itemindent,labelsep=0cm,align=left]
\item Usando la definición, demostrar que $\lim\limits_{(x,y) \rightarrow \left( 1,1\right) }\left\langle
x^{2},x+y\right\rangle =\left\langle 1,2\right\rangle .$
\vskip0.1cm

Observemos que  para todo $\varepsilon >0$ existe $\delta >0$ tal que si 
$\left\Vert(x,y) -\left( 1,1\right) \right\Vert <\delta $
entonces $\left\Vert \left( x^{2},x+y\right) -\left( 1,2\right) \right\Vert<\varepsilon ;$ esto se evidencia verificando las siguientes desigualdades:
\begin{eqnarray}\label{preliminares}
\left\vert x-1\right\vert &=&\sqrt{(x-1) ^{2}}\leq \sqrt{\left(
x-1\right) ^{2}+(y-1) ^{2}} \nonumber \\ &=&\left\Vert \left(x,y\right) -\left( 1,1\right) \right\Vert <\delta
\end{eqnarray}
De esta relaci\'{o}n se obtiene
\begin{multline}\label{preliminares1}
|x-1|<\delta \rightarrow -\delta <x-1<\delta
\rightarrow 2-\delta <x+1<\delta  \\
\Rightarrow \left\vert x+1\right\vert<2-\delta 
\end{multline}
De manera an\'{a}loga,
\begin{eqnarray}\label{preliminares2}
\left\vert y-1\right\vert &=& \sqrt{(y-1) ^{2}}\leq \sqrt{\left(x-1\right) ^{2}+(y-1)  ^{2}} \nonumber \\ &=& \left\Vert \left(x,y\right) -\left( 1,1\right) \right\Vert <\delta 
\end{eqnarray}
Utilizando los resultados (\ref{preliminares}), (\ref{preliminares1}) y (\ref{preliminares2}), se sigue que
\begin{eqnarray*}
\left\Vert \langle x^{2},x+y\rangle -\langle 1,2\rangle \right\Vert &=& \sqrt{\left( x^{2}-1\right) ^{2}+\left( x+y-2\right) ^{2}} \\
&\leq & \sqrt{\left( x^{2}-1\right) ^{2}} +\sqrt{\left(x+y-2\right) ^{2}} \\
&=& \left\vert x^{2}-1\right\vert  +\left\vert x+y-2\right\vert \\
&=& \left\vert x-1\right\vert  \left\vert x+1\right\vert +\left\vert (x-1) +(y-1) \right\vert \\
&<& \delta (2-\delta) +2\delta = \delta (4-\delta). 
\end{eqnarray*}
Tomando $\varepsilon =\delta \left( 4-\delta \right) $, se obtiene que 
$\delta =2-\sqrt{4-\varepsilon };$ con $0<\varepsilon < 4$ la definici\'{o}n se satisface.

\vskip0.3cm

\item Mostrar que $\lim\limits_{t\rightarrow 1}\left\langle \frac{t^{2}-1}{t-1},\frac{t+1}{t},t^{2}\right\rangle =\left\langle 2,2,1\right\rangle .$ Veamos que para todo $\varepsilon >0$ existe 
$\delta >0$ tal que si $\left\vert t-1\right\vert <\delta, $ 
entonces $\left\Vert \left\langle \frac{t^{2}-1}{t-1},\frac{t+1}{t},t^{2}\right\rangle -\left\langle 2,2,1\right\rangle \right\Vert<\varepsilon .$
De la relación
\begin{eqnarray*}
\left\Vert \left\langle \frac{t^{2}-1}{t-1},\frac{t+1}{t}%
,t^{2}\right\rangle -\left\langle 2,2,1\right\rangle \right\Vert  &=&
\sqrt{\left( t-1\right) ^{2}\left( t^{2}+2t+3\right) } \\ 
&=&\left\vert t-1\right\vert \sqrt{\left( t+1\right) ^{2}+2} \\
&\leq& \left\vert t-1\right\vert \left( \left\vert t-1\right\vert +\sqrt{2}
\right)  \\ &<& \delta \left( \delta +\sqrt{2}\right) 
\end{eqnarray*}

si $\delta \left( \delta +\sqrt{2}\right) =\varepsilon $. Entonces 
$\delta =\frac{\sqrt{2}}{2}\left( \sqrt{2\varepsilon +1}-1\right) ;$ lo cual hace que
la definici\'{o}n se satisfaga.
\end{itemize}
Una consecuencia de la definición (\ref{preliminares5}) es la siguiente.
\begin{teorema}\label{preliminares3}
Para  $\ve{F}: A \subseteq\mathbb{R}^n \to  \mathbb{R}^m,$ 
$$ \lim_{{\bf x}\to {\bf a}}\ve{F}({\bf x})=\ve{L}=\langle l_{1},l_{2},\cdots ,l_{m}\rangle $$ existe si y solo si 
cada uno de los límites de la funciones coordenadas existe.
\end{teorema}


\begin{proof} 
La función $\ve{F}$ tiene la forma
$$\ve{F}\left( x_{1},x_{2},\cdots ,x_{n}\right) =\left\langle f_{1}\left( x_{1},x_{2},\cdots ,x_{n}\right) ,\cdots ,f_{m}\left( x_{1},x_{2},\cdots ,x_{n}\right)\right\rangle $$
\end{proof}
\begin{itemize}[leftmargin=0cm,itemindent=.2cm,labelwidth=\itemindent,labelsep=0cm,align=left]
	
\item $(\to)$ Supongamos que $\lim_{{\bf x}\to {\bf a}}\ve{F}({\bf x})= \ve{L}$, 
esto significa que para todo $\varepsilon >0$ existe $\delta >0$  tal que si 
$\left\Vert\left( x_{1},x_{2},\cdots ,x_{n}\right) -\left( a_{1},a_{2},\cdots ,a_{n}\right) \right\Vert <\delta $
entonces $\left\Vert \ve{F}\left( x_{1},x_{2},\cdots ,x_{n}\right) -\ve{L}\right\Vert <\varepsilon .$ Por lo tanto:


\begin{eqnarray*}
\left\vert f_{i}\left( x_{1},x_{2},\cdots ,x_{n}\right) -l_{i}\right\vert &=&
\sqrt{\left( f_{i}\left( x_{1},x_{2},\cdots ,x_{n}\right) -l_{i}\right) ^{2}} \\
&\leq& \sqrt{\sum_{i=1}^{m}\left( f_{i}\left(x_{1},x_{2},\cdots ,x_{n}\right) -l_{i}\right) ^{2}} \\
&=& \left\Vert \ve{F}\left( x_{1},x_{2},\cdots ,x_{n}\right) -\ve{L}\right\Vert <\varepsilon
\end{eqnarray*}
lo cual implica que $\lim_{{\bf x}\to {\bf a}} f_{i}(x_{1},x_{2},\cdots ,x_{n})=l_i, $
para $i=1,2, \cdots, m. $
\vskip0.2cm

\item $(\leftarrow)$ Ahora supongamos que $\lim_{{\bf x}\to {\bf a}} f_{i}(x_{1},x_{2},\cdots ,x_{n})=l_i,$ 
esto significa que para todo $\frac{\varepsilon}{m}>0$ existe $\delta>0$ tal que
$|f_i( x_{1},x_{2},\cdots ,x_{n}) -l_i| <\varepsilon/m ,$ siempre que
$\left\Vert\left( x_{1},x_{2},\cdots ,x_{n}\right) -\left( a_{1},a_{2},\cdots ,a_{n}\right) \right\Vert <\delta .$
Por lo tanto:
{\small
\begin{eqnarray*}
\left\Vert \ve{F}\left( x_{1},x_{2},\cdots ,x_{n}\right) - \ve{L}\right\Vert &=& 
\sqrt{\sum_{i=1}^{m}\left( f_{i}\left(x_{1},x_{2},\cdots ,x_{n}\right) -l_{i}\right) ^{2}} \\ 
&\leq & \sum_{i=1}^{m}\sqrt{\left( f_{i}\left( x_{1},x_{2},\cdots ,x_{n}\right)-l_{i}\right) ^{2}}\\
&\leq& \sum_{i=1}^{m}\left\vert f_{i}\left( x_{1},x_{2},\cdots ,x_{n}\right)-l_{i} \right\vert\\ & <& \sum_{i=1}^{m}\frac{\varepsilon }{m}=\varepsilon.
\end{eqnarray*}
}

es decir $\left\Vert \ve{F}\left( x_{1},x_{2},\cdots ,x_{n}\right) -\ve{L}\right\Vert <\varepsilon ,$
lo cual implica que $\lim_{{\bf x}\to {\bf a}}\ve{F}({\bf x})$ existe y es igual a $\ve{L}.$
\end{itemize}

\vskip0.2cm
Este resultado permite hacer cálculos directos como se ilustra a continuación.
\vskip0.2cm

\begin{itemize}[leftmargin=*]
\item {\small $\ds\lim_{t\rightarrow 0}\left\langle \frac{\sin t}{t},\sqrt{t},\frac{e^{t}-1}{t}\right\rangle= \langle \lim_{t\rightarrow 0}\frac{\sin t}{t} , \lim_{t\rightarrow 0}\sqrt{1+t} ,
\lim_{t\rightarrow 0}\frac{e^{t}-1}{t}\rangle=\langle 1 ,1,1 \ra$}
\item Un límite con un campo vectorial.
{\small
\begin{multline*} \hskip-0.2cm 
\lim\limits_{(x,y) \rightarrow \left( 1,1\right)
}\left\langle \frac{x-y}{\sqrt{x}-\sqrt{y}}, \frac{x^{3}-y^{3}}{x-y}\right\rangle 
=\lim\limits_{(x,y) \rightarrow (1,1) } \left\langle \sqrt{x}+\sqrt{y},x^{2}+xy+y^{2}\right\rangle \\
=\left\langle \lim\limits_{(x,y)
\rightarrow \left( 1,1\right) }\left( \sqrt{x}+\sqrt{y}\right)
,\lim\limits_{(x,y) \rightarrow \left( 1,1\right) }\left(
x^{2}+xy+y^{2}\right) \right\rangle = \left\langle 2,3\right\rangle 
\end{multline*}
}
{\small $$
\lim_{(x,y) \rightarrow \left( 1,1\right) }\frac{x-y}{x^{2}-y^{2}}=\lim_{(x,y) \rightarrow \left( 1,1\right) }\frac{x-y}{\left( x-y\right) \left( x+y\right) }=\lim_{(x,y)\rightarrow\left(1,1\right) }\frac{1}{\left( x+y\right) }=\frac{1}{2}
$$}
\end{itemize}

La imagen del conjunto $A$ por medio de la funci\'{o}n $\ve{F}$, se denota por $\ve{F}(A)$, es decir 
$\ve{F}(A)=\{\ve{F}(\mathbf{x})/\mathbf{x}\in A\}.$
Para funciones $\ve{F},\ve{G}:A\subseteq \mathbb{R}^n\to \mathbb{R}^m$
y $h:\mathbb{R}^n \to \mathbb{R}$ usaremos las siguientes abreviaturas: 
\begin{eqnarray*}
\left( \ve{F}\pm \ve{G}\right) \left( \mathbf{x} \right) &=& \ve{F}(\mathbf{x}) \pm \ve{G}(\mathbf{x})\\
 &=& \left\langle f_{1}(\mathbf{x}) \pm g_{1}(\mathbf{x}) ,f_{2}(\mathbf{x}) \pm g_{2}\left( 
\mathbf{x}\right) ,\cdots ,f_{m}(\mathbf{x}) \pm g_{m}\left(  \mathbf{x}\right) \right\rangle 
\end{eqnarray*}

\begin{eqnarray*}
\left( \ve{F}\cdot \ve{G}\right) \left( \mathbf{x} \right) &=&\ve{F}(\mathbf{x}) \cdot \ve{G}%
(\mathbf{x}) \\ &=& f_{1}(\mathbf{x}) g_{1}\left( \mathbf{x}\right) +f_{2}(\mathbf{x}) g_{2}(\mathbf{x})
+\cdots +f_{m}(\mathbf{x}) g_{m}(\mathbf{x}) \\
&=& \sum_{k=1}^{m}f_{k}(\mathbf{x}) g_{k}(\mathbf{x}) \\[0.1cm]
\left\Vert \ve{F}(\mathbf{x}) \right\Vert &=& \left( f_{1}^{2}(\mathbf{x}) +f_{2}^{2}(\mathbf{x}) +\cdots
+f_{m}^{2}(\mathbf{x}) \right) ^{1/2}=\sqrt{\sum_{k=1}^{m}f_{k}^{2}(\mathbf{x}) }\\[0.1cm]
h\ve{F}(\mathbf{x}) &=& h(\mathbf{x}) \ve{F}(\mathbf{x}) =\left\langle h\left( \mathbf{x} \right) f_{1}(\mathbf{x}) ,h(\mathbf{x}) f_{2}\left( \mathbf{x}\right) ,\cdots ,h(\mathbf{x}) f_{m}\left( \mathbf{x}\right) \right\rangle 
\end{eqnarray*}

Al considerar esta notación se presenta el siguiente resultado.

\begin{teorema}\label{teo1}
Sean $\ve{F},\ve{G}$ y $h$ funciones que satisfacen que 
$\ve{F},\ve{G}:A\subseteq \mathbb{R}^n\to \mathbb{R}^m$ y $h:A\subseteq \mathbb{R}^n\to \mathbb{R}$ tal que
$\lim_{\mathbf{x\to x}_{0}}\ve{F}\left( \mathbf{x}%
\right) =\mathbf{a,}$ $\lim_{\mathbf{x\to x}_{0}}\ve{G}%
(\mathbf{x}) =\mathbf{b,}$ y $\lim_{\mathbf{x\to x}%
_{0}}h(\mathbf{x}) =k;$ entonces 

\begin{enumerate}[leftmargin=0cm,itemindent=0.4cm,labelwidth=\itemindent,labelsep=0cm,align=left]

\item $\ds \lim_{\mathbf{x\to x}_{0}}\left( \ve{F}\left( \mathbf{x}%
\right) \pm \ve{G}(\mathbf{x}) \right) =\mathbf{a\pm b}$

\item $\ds\lim_{\mathbf{x\to x}_{0}}\left( h(\mathbf{x})
\ve{F}(\mathbf{x}) \right) =k\mathbf{a}$

\item $\ds\lim_{\mathbf{x\to x}_{0}}\left( \ve{F}(\mathbf{x}) \cdot \ve{G}(\mathbf{x}) \right) =\mathbf{a\cdot b}$

\item $\ds\lim_{\mathbf{x\to x}_{0}}\left\Vert \ve{F}(\mathbf{x}) \right\Vert =\left\Vert \mathbf{a}\right\Vert $
\end{enumerate}
\end{teorema}

\begin{proof}
Con la definición obtenemos las siguientes afirmaciones
\begin{enumerate}[leftmargin=0cm,itemindent=0.4cm,labelwidth=\itemindent,labelsep=0cm,align=left]
\item $\lim_{\mathbf{x\to x}_{0}}\ve{F}(\mathbf{x}) =\mathbf{a,}$
significa que para todo $\varepsilon >0$ existe $\delta >0,$ 
es decir, si $\left\Vert \mathbf{x-x}_{0}\right\Vert <\delta $ entonces 
$\left\Vert \ve{F}(\mathbf{x}) \mathbf{-a}\right\Vert <\frac{\varepsilon }{2}.$ 
Por la misma raz\'{o}n 
$\left\Vert \ve{G}(\mathbf{x}) \mathbf{-b}\right\Vert <\frac{\varepsilon }{2},$
entonces
\begin{multline*}
\left\Vert \left( \ve{F}(\mathbf{x}) \pm \ve{G}(\mathbf{x}) \right) -\left( \mathbf{a\pm b}
\right) \right\Vert = \left\Vert \left( \ve{F}\left( \mathbf{x} \right) -\mathbf{a}\right) \pm \left( \ve{G}\left( \mathbf{x} \right) -\mathbf{b}\right) \right\Vert \\
\leq \left\Vert \ve{F} (\mathbf{x}) -\mathbf{a}\right\Vert +\left\Vert \ve{G}(\mathbf{x}) -\mathbf{b}\right\Vert <\frac{\varepsilon }{2}+\frac{\varepsilon }{2}= \varepsilon 
\end{multline*}

\item Es consecuencia de la siguiente igualdad
$$h(\mathbf{x}) \ve{F}(\mathbf{x}) -k \mathbf{a=}\left( h(\mathbf{x}) -k\right) \left( \ve{%
F}(\mathbf{x}) -\mathbf{a}\right) +\left( h\left( \mathbf{x} \right) -k\right) \mathbf{a}+k\left( \ve{F}\left( \mathbf{x} \right) -\mathbf{a}\right),$$
por lo tanto,
\begin{multline*}
\left\Vert h(\mathbf{x}) \ve{F}\left( \mathbf{x}\right) -k\mathbf{a}\right\Vert \mathbf{\leq }\left\vert h\left( \mathbf{x} \right) -k\right\vert \left\Vert \ve{F}(\mathbf{x}) - \mathbf{a}\right\Vert \\ +\left\vert h(\mathbf{x}) -k\right\vert \left\Vert \mathbf{a}\right\Vert +\left\vert k\right\vert \left\Vert \ve{F}(\mathbf{x}) -\mathbf{a}\right\Vert 
\end{multline*}
Los t\'{e}rminos de la derecha tienden a cero cuando $\mathbf{x}\to \mathbf{x}_{0}.$

\item Utilizando la desigualdad de Cauchy-Schwarz, 
$0\leq \left\vert \ve{u}\cdot \ve{v}\right\vert \leq \left\Vert 
\ve{u}\right\Vert \cdot \left\Vert \ve{v}\right\Vert $
y la siguiente identidad
\begin{multline*}
\ve{F}(\mathbf{x}) \cdot \ve{G}\left( \mathbf{x}\right) -\mathbf{a\cdot b=}\left( \ve{F}\left( \mathbf{x}\right) -\mathbf{a}\right) \cdot \left( \ve{G}\left( \mathbf{x} \right) -\mathbf{b}\right) \\ 
+\mathbf{a\cdot }\left( \ve{G}\left( \mathbf{x}\right) -\mathbf{b}\right) +\mathbf{b\cdot }\left( \ve{F}(\mathbf{x}) -\mathbf{a}\right), 
\end{multline*}
se sigue que 
\begin{multline*}
\left\vert \ve{F}(\mathbf{x}) \cdot \ve{G}(\mathbf{x}) -\mathbf{a\cdot b}\right\vert \mathbf{\leq }%
\left\Vert \ve{F}(\mathbf{x}) -\mathbf{a}\right\Vert \left\Vert \ve{G}(\mathbf{x}) -\mathbf{b}\right\Vert \\
+\left\Vert \mathbf{a}\right\Vert \left\Vert \ve{G}\left( \mathbf{x}\right) -\mathbf{b}\right\Vert +\left\Vert \mathbf{b}\right\Vert \left\Vert \ve{F}(\mathbf{x}) -\mathbf{a}\right\Vert.
\end{multline*}
De las hipótesis se puede establecer que cuando $\mathbf{x\to x }_{0},$
$\left\Vert \ve{F}(\mathbf{x}) -\mathbf{a} \right\Vert \to 0,$
y $\left\Vert \ve{G}\left( \mathbf{x} \right) -\mathbf{b}\right\Vert \to 0;$ 
por lo tanto
$\left\vert \ve{F}(\mathbf{x}) \cdot \ve{G}(\mathbf{x}) -\mathbf{a\cdot b}\right\vert \to 0,$ 
cuando $\mathbf{x\to x}_{0}.$
\item El resultado es consecuencia de la siguiente desigualdad
$$\left\vert \left\Vert \ve{F}(\mathbf{x})\right\Vert -\left\Vert \mathbf{a}\right\Vert \right\vert \leq \left\Vert 
\ve{F}(\mathbf{x}) -\mathbf{a}\right\Vert $$
\end{enumerate}
\end{proof}

La siguiente es una característica especial que deseamos estudiar en la mayoría de las funciones que se consideran
en estas notas.

\subsection*{Funciones continuas}
\begin{defi} \index{Función!continua}
Se dice que una funci\'{o}n $\ve{F}:A\subseteq \mathbb{R}^n\to \mathbb{R}^m$ es continua en un punto 
$\mathbf{a}\in A$, si para todo $\varepsilon >0$ existe $\delta >0$ tal que si 
$\left\Vert \mathbf{x}-\mathbf{a}\right\Vert <\delta, $ entonces 
$\left\Vert \ve{F}(\mathbf{x}) -\ve{F}\left( \mathbf{a}\right) \right\Vert <\varepsilon .$
\vskip0.1cm
Diremos que $\ve{F}$ es continua en $A$, cuando $\ve{F}$ es continua en todo punto de $A.$
\end{defi}
Los siguientes teoremas son consecuencias de los teoremas de límites y de la definición anterior.

\begin{teorema}\label{preliminares6} Para $\ve{F}:A \subseteq\mathbb{R}^{n}\rightarrow \mathbb{R}^{m},$
$\ve{F}$ es continua si y solo si cada una sus funciones componentes es continua.
\end{teorema}

\begin{proof}
Las funciones componentes de $\ve{F}$ son campos escalares, es decir
{\small $\ve{F}\left( x_{1},x_{2},\cdots ,x_{n}\right) =\left\langle f_{1}\left(x_{1},x_{2},\cdots ,x_{n}\right) ,\cdots ,f_{m}\left( x_{1},x_{2},\cdots ,x_{n}\right)\right\rangle.$} \\
El resultado es consecuencia de aplicar el teorema (\ref{preliminares3}).
\end{proof} 

En el siguiente resultado $\ve{G}\circ \ve{F}$ denota la composición de las funciones $\ve{G}$ y $\ve{F}.$

\begin{teorema}
\normalfont Sea $\ve{F}:A\subseteq \mathbb{R}^n\to \mathbb{R}^m$ y 
$\ve{G}:B\subseteq \mathbb{R}^m\to \mathbb{R}^p$, con 
$\ve{F}\left( \mathbf{a}\right) \in B$. Si $\ve{F}$ es continua en 
$\mathbf{a}\in A$ y $\ve{G}$ es continua en $\ve{F}\left( \mathbf{a}\right) $ 
entonces $\ve{G}\circ \ve{F}:A\subseteq \mathbb{R}^n\to \mathbb{R}^p$ es continua en $\mathbf{a}$.
\end{teorema}

\begin{proof}
Por ser $\ve{G}$ continua en $\ve{F}\left( \mathbf{a}\right) \in B,$ dado $\varepsilon >0$ 
existe $\lambda >0,$ para $\mathbf{y} \in B,$ si $\left\Vert \mathbf{y-}\ve{F}\left( \mathbf{a}\right)
\right\Vert <\lambda $ entonces $\left\Vert \ve{G}(\mathbf{y})-\ve{G}\left( \ve{F}\left( \mathbf{a}\right) \right) \right\Vert <\varepsilon .$
\vskip0.1cm 
Como $\ve{F}$ continua en $\mathbf{a}\in A,$ esto implica que dado 
$\delta >0,$ existe $\lambda >0,$  si $\left\Vert \mathbf{x-a}\right\Vert <\delta ,$ para $\mathbf{x}\in A,$
entonces $\left\Vert \ve{F}(\mathbf{x}) -\ve{F}\left( \mathbf{a}\right) \right\Vert <\lambda ,$ 
lo anterior implica que 
$\left\Vert \ve{G}\left( \ve{F}\left( \mathbf{x} \right) \right) -\ve{G}\left( \ve{F}\left( 
\mathbf{a}\right) \right) \right\Vert <\varepsilon .$
En síntesis, dado $\varepsilon >0,$ existe $\delta >0$ tal que si 
$\left\Vert \mathbf{x-a}\right\Vert <\delta $ entonces 
$\left\Vert \ve{G} \left( \ve{F}(\mathbf{x}) \right) -\ve{G}\left( \ve{F}\left( \mathbf{a}\right) \right)
\right\Vert <\varepsilon ;$ es decir, $\ve{G}\circ \ve{F}$ es continua en $\mathbf{a.}$
\end{proof} 

\begin{teorema}
Sea $\ve{F},\ve{G}:A\subseteq \mathbb{R}^n\to \mathbb{R}^m$ y $h:A\subseteq \mathbb{R}^n\to \mathbb{R}$, 
funciones continuas en $\mathbf{a}\in A,$ entonces las siguientes funciones son continuas en $\mathbf{a}$ 
%
$\ve{F}\pm \ve{G}:A\subseteq \mathbb{R}^n\to \mathbb{R}^m,$ \\
$ \ve{F}\cdot \ve{G}:A\subseteq \mathbb{R}^n\to \mathbb{R},$\\
$\left\Vert  \ve{F}\right\Vert :A\subseteq \mathbb{R}^n\to \mathbb{R},$\\
$h\ve{F}:A\subseteq \mathbb{R}^n\to \mathbb{R}^m.$ 
%
\end{teorema}
La prueba es consecuencia de aplicar los teoremas (\ref{preliminares5}) y (\ref{teo1}).
En particular, el lector puede verificar la continuidad de la función ${\bf u}(t) \times {\bf v}(t)$ para dos funciones vectoriales continuas 
${\bf u}, {\bf v}: \mathbb{R}\to \mathbb{R}^3.$

\section{Diferenciabilidad} \label{diferencibilidad}
Para funciones reales $f:\mathbb{R}\rightarrow \mathbb{R}$, el lector seguramente
estar\'{a} familiarizado con la definici\'{o}n de la derivada en un punto $a$,
$$f^{\prime}(a) =\lim_{h\rightarrow 0}\frac{f(a+h)-f(a) }{h}.$$
Al considerar la transformaci\'{o}n lineal $A(h)=f^{\prime}(a)\cdot h.$
Entonces, la expresi\'{o}n anterior se reescribe como 
$$\lim_{h\rightarrow 0}\frac{f\left( a+h\right) -f\left( a\right) -A\left( h\right) }{h}=0.$$

En el caso general una funci\'{o}n ${\bf F}:\mathbb{R}^{n}\rightarrow \mathbb{R}^{m}$ se dice
diferenciable en $a\in \mathbb{R}^{n}$
si existe una transformaci\'{o}n lineal $\lambda :\mathbb{R}^{n}\rightarrow \mathbb{R}^{m},$ tal que 
$$\lim_{h\rightarrow 0}\frac{\left\Vert {\bf F}(a+h)-{\bf F}(a)-\lambda(h)\right\Vert }{\left\Vert h\right\Vert }=0.$$
Es de anotar que $\left\Vert {\bf F}\left( a+h\right) -{\bf F}(a) -\lambda(h) \right\Vert$ representan la norma en 
$\mathbb{R}^{m}$ y $\left\Vert h\right\Vert $ es la norma en $\mathbb{R}^{n}.$ \index{Función!diferenciable}
Se dice que ${\bf F}$ es diferenciable en $A$, si es diferenciable para todo $a\in A.$
La mencionada transformación lineal es la denominada \emph{matriz jacobiana} de ${\bf F}$ en $a$,
que se denota por $D({\bf F})(a)$. Para ${\bf F}:\mathbb{R}^n\to {\mathbb R}^m$ y $a=(a_1,a_2,\cdots, a_n)$, es claro que
${\bf F}(a)=\langle f_1(a),f_2(a), \cdots, f_m(a) \rangle. $
La entrada $ij$ de la matriz jacobiana está dada por
$$\frac{\partial f_i}{\partial x_j}(a)=\lim_{h\rightarrow 0}\frac{f_i(a_1,a_2,\cdots,a_j+h,\cdots,a_n)-f_i(a_1,a_2,\cdots,a_j,\cdots, a_n)}{h}.$$
que la denominaremos la derivada parcial de $f_i$ con respecto a $x_j$. \index{Derivada parcial}
La matriz jacobiana de ${\bf F}$, $D{\bf F}(a)$ es 
$$D{\bf F}(a)=\left(\begin{array}{cccc}
\frac{\partial f_1}{\partial x_1}(a)&\frac{\partial f_1}{\partial x_2}(a)&\cdots & \frac{\partial f_1}{\partial x_n}(a)\\[0.2cm]
\frac{\partial f_2}{\partial x_1}(a)&\frac{\partial f_2}{\partial x_2}(a)&\cdots & \frac{\partial f_2}{\partial x_n}(a)\\[0.2cm]
\vdots \\[0.2cm]
\frac{\partial f_m}{\partial x_1}(a)&\frac{\partial f_m}{\partial x_2}(a)&\cdots & \frac{\partial f_m}{\partial x_n}(a)
\end{array} \right)$$
\begin{teorema}
Si ${\bf F}:\mathbb{R}^n \to \mathbb{R}^m$ es diferenciable en $a$ y 
${\bf G}:\mathbb{R}^m \to \mathbb{R}^p$ es diferenciable en ${\bf F}(a)$ entonces la función compuesta
${\bf G}\circ {\bf F}:\mathbb{R}^n \to \mathbb{R}^p$ es diferenciable en $a$. Además
$D({\bf G}\circ {\bf F})(a)=D({\bf G}( {\bf F}(a))\cdot D{\bf F}(a). $
\end{teorema}
La prueba de este resultado puede consultarse en \textcite{spivak_calculus_2018}. 
En el caso de campos escalares $f:\mathbb{R}^n \to \mathbb{R}$ a $Df(a)$ se denota también por 
$\nabla f(a)$ y se denomina el \emph{vector gradiente} de $f$. \index{Gradiente}
\vskip0.2cm
Un campo vectorial ${\bf F}(x,y,z)$ con dominio $D \subseteq \mathbb{R}^3$ 
se denomina \emph{conservativo} si existe un campo escalar
$\phi(x,y,z)$ que se denomina \emph{función potencial} con dominio $D$ y satisface que $\nabla \phi ={\bf F}.$ \index{Campo vectorial!conservativo} \index{Función!potencial}
En \cite{thomas_calculo_2015} se presenta el siguiente criterio, el cual es útil para determinar si un campo vectorial cumple esta condición.
\begin{prop}
Sea ${\bf F} = P(x, y,z)\vi +Q(x,y,z)\vj +P(x,y,z)\vk$ un campo en un dominio abierto simplemente
conexo, cuyas funciones componentes tienen sus primeras derivadas parciales continuas. 
De esta forma, ${\bf F}$ es conservativo si y sólo se satisfacen la siguientes condiciones:
$$\frac{\partial P}{\partial y}=\frac{\partial Q}{\partial x},\hspace{1.5cm}
\frac{\partial P}{\partial z}=\frac{\partial R}{\partial x},\hspace{1.5cm}
\frac{\partial Q}{\partial z}=\frac{\partial R}{\partial y}.$$
\end{prop}
\section{Funciones vectoriales} \index{Funciones vectoriales}
De manera general, una aplicación que asigna a cada número real un vector en $\mathbb{R}^n$
se denomina función vectorial, esto es ${\bf r}(t)= \langle r_1(t),r_2(t),\cdots, r_n(t)\ra$.
\vskip0.2cm
Por ejemplo, la función ${\bf r}(t)= \langle 3\cos(2t),1-\cos(t),2\sin(t)\ra$ está definida para $t\in [0,2\pi]$ su 
gráfico en $\mathbb{R}^3$ se representa a continuación.


\vskip0.2cm
\noindent
\begin{minipage}{\textwidth}
	\centering
	\captionof{figure}{Curva determinada por ${\bf r}(t)= \langle 3\cos(2t),1-\cos(t),2\sin(t)\ra$}
	\begin{minipage}{0.38\linewidth}
		\centering
		\includegraphics[width=\linewidth]{Fig/4.pdf}
	\end{minipage}
	\begin{minipage}{0.38\linewidth}
		\centering
		\includegraphics[width=\linewidth]{Fig/5.pdf}
	\end{minipage}
	
	\vspace{0.3em}
	\parbox{0.95\linewidth}{\centering\fontsize{8pt}{10pt}\selectfont\textbf{Fuente:} elaboración propia.}
\end{minipage}
\vskip0.2cm


Para una función vectorial ${\bf r}(t)=\langle r_1(t),r_2(t),\cdots, r_m(t) \rangle ,$
los dos primeros de los siguientes resultados se obtienen 
como caso particular de los teoremas (\ref{preliminares3}) y (\ref{preliminares6}), con $n=1.$ 

\begin{itemize}
\item $\lim_{t\to a}{\bf r}(t)$ existe si y solo si $\lim_{t\to a}r_i(t)$ para toda $i=1,2,\cdots ,m$.
\item ${\bf r}(t)$ es continua en $a$ si y solo si cada una de las funciones componentes $r_i(t)$ es continua en $a$, para toda $i=1,2,\cdots ,m$.
\item ${\bf r}(t)$ es diferenciable en $a$ si y solo si cada una de las funciones componentes $r_i(t)$ es diferenciable en $a$, para toda $i=1,2,\cdots ,m$; además ${\bf r}^\prime(t)= \langle r^\prime_1(t),r^\prime_2(t),\cdots, r^\prime_m(t)\ra.$
\end{itemize}

Presentamos la demostración del último de estos resultados.
\begin{eqnarray*}
\mathbf{r}^{\prime}(t) &=& \lim_{h\rightarrow 0}\frac{1}{h}\left( \mathbf{r}\left( t+h\right) -\mathbf{r}(t) \right) \\
&=& \lim_{h\rightarrow 0}\frac{1}{h}\left\langle r_{1}\left(t+h\right) -r_{1}(t) ,\cdots ,r_{m}\left( t+h\right) -r_{m}(t) \right\rangle \\
&=&\left\langle \lim_{h\rightarrow 0}\frac{r_{1}\left( t+h\right) -r_{1}(t) }{h},\cdots ,\lim_{h\rightarrow 0}\frac{r_{m}\left( t+h\right) -r_{m}(t) }{h}\right\rangle \\
&=& \left\langle r_{1}^{\prime}(t),r_{2}^{\prime}(t) ,\cdots ,r_{m}^{\prime}\left(t\right) \right\rangle 
\end{eqnarray*}

Si $\mathbf{u}(t) =\left\langle u_{1}(t),u_{2}(t) ,\cdots ,u_{n}(t) \right\rangle $ y 
$\mathbf{v}(t) =\left\langle v_{1}(t) ,v_{2}\left(t\right) ,\cdots ,v_{n}(t) \right\rangle $
son dos funciones vectoriales entonces el producto punto entre funciones vectoriales está dado por
$\mathbf{u}(t) \cdot \mathbf{v}(t) =u_{1}\left(t\right) v_{1}(t) +u_{2}(t) v_{2}(t)+\cdots +u_{n}(t) v_{n}(t) .$
Derivando se obtiene las siguientes igualdades
\begin{eqnarray*}
\left( \mathbf{u}(t) \cdot \mathbf{v}(t) \right)^{\prime} 
%&=& u_{1}^{\prime(t) v_{1}(t) +u_{1}(t) v_{1}^{\prime }(t) +u_{2}^{\prime }(t)v_{2}(t) +u_{2}(t) v_{2}^{\prime }(t)+ \cdots + u_{n}^{\prime }(t) v_{n}(t) +u_{n}\left(t\right) v_{n}^{\prime }(t) \\
&=& \big( u_{1}^{\prime }(t) v_{1}(t) +u_{2}^{\prime}(t) v_{2}(t) +\cdots +u_{n}^{\prime }(t) v_{n}(t) \big) + \\ 
&& \big( u_{1}(t)v_{1}^{\prime }(t) +  u_{2}(t) v_{2}^{\prime }(t) +\cdots +u_{n}(t) v_{n}^{\prime }(t)\big) \\
&=& \mathbf{u}^{\prime }(t) \cdot \mathbf{v}(t) +\mathbf{u}(t) \cdot \mathbf{v}^{\prime }(t) 
\end{eqnarray*}
Lo cual se resume en la siguiente fórmula
\begin{equation}
\left( \mathbf{u}(t) \cdot \mathbf{v}(t) \right)^{\prime }=\mathbf{u}^{\prime }(t) \cdot \mathbf{v}\left(
t\right) +\mathbf{u}(t) \cdot \mathbf{v}^{\prime }\left(t\right).
\end{equation}
Otras fórmulas útiles, son las siguientes
\begin{eqnarray*}
(\left\Vert \mathbf{u}(t) \right\Vert)^{\prime } &=& \frac{d}{dt}\left( \left( \mathbf{u}(t) \cdot \mathbf{u}(t)
\right) ^{1/2}\right) \\
&=& \frac{1}{2}\left( \mathbf{u}(t) \cdot \mathbf{u}(t) \right) ^{-1/2}\left( \mathbf{u}^{\prime }(t) \cdot 
\mathbf{u}(t) +\mathbf{u}(t) \cdot \mathbf{u}^{\prime }(t) \right) = \frac{\mathbf{u}^{\prime }(t)\cdot \mathbf{u}(t) }{\left\Vert \mathbf{u}(t)\right\Vert }
\end{eqnarray*}

La derivada del producto cruz entre dos funciones vectoriales se evalua como sigue
\begin{prop}
Si $\mathbf{u}(t)$ y $\mathbf{v}(t)$ son funciones vectoriales definidas en $\mathbb{R}^3$ entonces
$\left(\mathbf{u}(t)\times\mathbf{v}(t)\right)^{\prime}= \mathbf{u}^{\prime}(t)\times\mathbf{v}(t)+\mathbf{u}(t)\times \mathbf{v}^{\prime}(t)$.
\end{prop}

\begin{proof}
Si $\mathbf{u}(t)=\left\langle u_{1}(t),u_{2}(t),u_{3}(t) \right\rangle $ y 
$\mathbf{v}(t)=\left\langle v_{1}(t),v_{2}(t),v_{3}(t)\right\rangle ,$
obviaremos el uso de la variable $t$ para hacer menos extensa la escritura, es decir
\begin{eqnarray*}
\frac{d}{dt}\left( \mathbf{u}(t)\times \mathbf{v}(t)\right) &=& \frac{d}{dt} \langle u_{2} v_{3} -u_{3} v_{2} ,u_{3} v_{1} -u_{1} v_{3} ,u_{1} v_{2} -u_{2} v_{1} \rangle \\
&=& \big\langle u_{2}^{\prime} v_{3}+u_{2} v_{3}^{\prime} -u_{3}^{\prime} v_{2} -u_{3} v_{2}^{\prime},u_{3}^{\prime} v_{1} +u_{3}v_{1}^{\prime} \\
&& -u_{1}^{\prime}v_{3} -u_{1} v_{3}^{\prime},   u_{1}^{\prime}v_{2}+u_{1}v_{2}^{\prime}  -u_{2}^{\prime} v_{1} -u_{2} v_{1}^{\prime} \big\rangle \\
&=& \left\langle u_{2}^{\prime} v_{3}-u_{3}^{\prime} v_{2} ,u_{3}^{\prime}(t) v_{1} -u_{1}^{\prime} v_{3},u_{1}^{\prime} v_{2} -u_{2}^{\prime} v_{1} \right\rangle  \\
&& +\left\langle u_{2} v_{3}^{\prime} -u_{3}v_{2}^{\prime} ,u_{3} v_{1}^{\prime}(t) -u_{1} v_{3}^{\prime} ,u_{1} v_{2}^{\prime} -u_{2} v_{1}^{\prime} \right\rangle \\
&=& \mathbf{u}^{\prime}(t)\times \mathbf{v}(t)+\mathbf{u}(t)\times \mathbf{v}^{\prime}(t) 
\end{eqnarray*}
\end{proof}
La siguiente expresión es producto de aplicar los resultados anteriores
\begin{multline}\left( \left( \mathbf{u}(t)\times \mathbf{v}(t)\right)\cdot \mathbf{w}(t)\right)^{\prime} \\ = \left( \left( \mathbf{u}^{\prime}(t)\times \mathbf{v}(t)+\mathbf{u}(t)\times \mathbf{v}^{\prime}(t)\right) \cdot \mathbf{w}(t)\right) +\left( \left( \mathbf{u}(t)\times \mathbf{v}(t)\right)\cdot \mathbf{w}^{\prime}(t)\right) 
\end{multline}

\section{Regiones parametrizadas} \label{regionesparametrizadas}
En estas notas trabajaremos básicamente sobre tres tipos de objetos en el espacio: curvas, superficies y sólidos.
En esta sección presentaremos una manera de representar dichos objetos, la idea consiste en \emph{parametrizar} la región en mención. Para este propósito se utiliza, según el caso, uno, dos o tres parámetros, respectivamente.  
\index{Parametrizar}

\vspace{-1mm}

\begin{tcolorbox}[breakable,enhanced,title=\bf Curva parametrizada]
La funci\'on $\ve{r}(t)$ se denomina una \emph{parametrizaci\'on} de una curva $C$ si 
al variar el parametro $t\in [a,b]$ la imagen $\ve{r}(t)$ describe los puntos de la curva $C$. 
Si $\ve{r}(a)=\ve{r}(b)$ la curva se denomina cerrada y la curva se denomina simple si $\ve{r}(t)\neq \ve{r}(s),$
para $a\neq t<s\neq b$.
\end{tcolorbox}

Las funciones vectoriales son un ejemplo de parametrización de curvas.
El concepto de curva simple hace referencia a que la curva no se autointerseca, posiblemente solo en los puntos extremos.
La curva se denomina \emph{suave} si $\ve{r}^\prime(t) \neq \ve{0},$ para todo $t.$ \index{Curva!suave}

\begin{tcolorbox}[breakable,enhanced,title=\bf Superficie parametrizada]
Una \emph{parametrización} de una superficie $\sigma $
es una aplicación $$\ve{r}:[a,b]\times \lbrack c,d]\rightarrow \mathbb{R}^{3}$$
tal que al variar los parámetros $u\in \lbrack a,b]$ y $v\in \lbrack c,d]$
su imagen $\ve{r}(u,v)$ va describiendo los puntos de la superficie $\sigma $.
\end{tcolorbox}
La superficie se denomina \emph{suave} si $\frac{\partial \ve{r}}{\partial u}$ y 
$\frac{\partial \ve{r}}{\partial v} $ son continuas y 
$\frac{\partial \ve{r}}{\partial u} \times \frac{\partial \ve{r}}{\partial u} \neq \ve{0}.$ 
\index{Superficie!suave}
%\vskip0.2cm
Con los elementos anteriormente descritos es posible construir funciones
que describan un sólido en el espacio tridimensional. 
\begin{tcolorbox}[breakable,enhanced,title=\bf Parametrización de un sólido]
Se dice que la función $${\bf r}(u,v,w)= \langle x(u,v,w), y(u,v,w),z(u,v,w) \rangle $$  
parametriza el sólido $E$ si al variar los parámetros se recorre completamente a $E$,  
con $a \leq u \leq b$, $c \leq v \leq d$ y $p \leq u \leq q$.  
\end{tcolorbox}


Una ventaja de usar este tipo de representación es el hecho de que al usar límites constantes el planteamiento de integrales
resulta, en algunos casos, bastante sencillo.

\section{Elementos básicos de integración}

Presentamos en esta sección algunos de los elementos teóricos propios de un curso de 
Cálculo Vectorial que serán usados en los capítulos posteriores. La exposición teórica completa puede ser consultada 
por el lector con más detalle en \cite{apostol_alisis_1981}, \cite{marsden_calculo_2013}, \cite{stewart_calculo_2018} o \cite{thomas_calculo_2015}.
\vskip0.2cm
Posteriormente se utilizan estos elementos en algunas aplicaciones de la integración principalmente en 
el c\'{a}lculo de longitud de curvas, \'{a}reas, volumenes y centros de masa. 
Se estudiar\'{a} tambi\'{e}n el teorema de cambio de variables,
particularmente cuando se utilizan coordenadas polares, cilindricas o esf\'{e}ricas, pero se har\'{a} enfasis en 
parametrizar la regi\'{o}n de integraci\'{o}n.
\vskip0.2cm
Una función $g:\mathbb{R}\to\mathbb{R}$ se dice \emph{integrable} en $[a,b]$ en el sentido señalado en \textcite{apostol_calculus_1999};
utilizaremos esta idea en los siguientes resultados. 

\subsection*{Integrales de línea en campos escalares}

Consideremos una curva $\alpha$ parametrizada por una función 
${\bf r}:\mathbb{R}\to\mathbb{R}^3$ continua, inyectiva y derivable en $(a,b)$, con ${\bf r}(t)=\langle x(t),y(t),z(t) \rangle .$ 
% \index{Partición}
\vskip0.2cm

\noindent
\begin{minipage}{\textwidth}
	\centering
	\captionof{figure}{Curva en $\mathbb{R}^3$}
	\label{Fig/6.png}
	\begin{minipage}{0.32\linewidth}
		\centering
		\includegraphics[width=\linewidth]{Fig/6.png}
	\end{minipage}
	
	\vspace{0.3em}
	\parbox{0.95\linewidth}{\centering\fontsize{8pt}{10pt}\selectfont\textbf{Fuente:} elaboración propia.}
\end{minipage}

\vskip0.2cm
Una \emph{partición} $P$ del intervalo $[a,b]$ es un conjunto finito de puntos 
$P=\{ t_0, t_1, t_2, \cdots , t_{n-1},t_n \}$ que satisface 
$$a = t_0< t_1<  t_2< \cdots < t_{n-1}< t_n=b;$$
la cual determina una partición de la curva $\alpha$ en $n$ puntos de la forma
$Q_i(x(t_{i-1}), x(y_{i-1}),z(t_{i-1})),$ para $i=1,2, \cdots , n.$
%

Denotaremos por $$\Delta\alpha_i=\int_{t_i}^{t_{i+1}}\Vert {\bf r}^\prime(t)\Vert dt $$ 
a la longitud de la curva $\alpha$ desde el punto $Q_i$ hasta $Q_{i+1}$.
Para una función continua $f: \alpha \to \mathbb{R}$ consideremos la suma 
$$\sum_{i=1}^n f({\bf r}(t_i^*))\Delta \alpha_i=\sum_{i=1}^n f(x(t_i^*),y(t_i^*),z(t_i^*))\Delta \alpha_i$$
en donde $t_i^* \in [t_{i-1},t_i],$ con los elementos anteriores se define la integral de una campo escalar sobre una curva.
\begin{defi}
Consideremos $f:\mathbb{R}^3\to \mathbb{R}$ continua. 
La integral de $f$ sobre la curva suave $\alpha$ con respecto a la parametrización 
${\bf r}(t)$ está dada por
\begin{equation} \label{intdelinea}
\int_a^b f({\bf r}(t))\Vert{\bf r}^\prime(t)\Vert dt=\lim_{n\to\infty}\left(\sum_{i=1}^n f({\bf r}(t_i^*))\Delta \alpha_i\right) \end{equation} \index{Integrales!de línea}
\end{defi}
En particular, si la función $f$ de la integral (\ref{intdelinea}) es $1,$ obtendremos la fórmula para la longitud de la 
curva (35) y si $f$ es la densidad $\delta(x,y,z)$ de un alambre o varilla delgada cuya forma está determinada por la curva $\alpha$ en el punto $(x,y,z),$ entonces la masa $M$ de la varilla está dada por
\begin{equation} 
M=\int_a^b \delta({\bf r}(t))\Vert{\bf r}^\prime(t)\Vert dt \end{equation}
Los momentos con respecto a los planos coordenados se definen como
\begin{eqnarray*} 
\mu_{xy}&=&\int_a^b z(t)\,\delta({\bf r}(t))\Vert{\bf r}^\prime(t)\Vert dt,\\
\mu_{yz}&=&\int_a^b x(t)\,\delta({\bf r}(t))\Vert{\bf r}^\prime(t)\Vert dt,\\
\mu_{xz}&=&\int_a^b y(t)\,\delta({\bf r}(t))\Vert{\bf r}^\prime(t)\Vert dt.
\end{eqnarray*}
Las coordenadas del centro de masa de la varilla son 
\index{Centro de masa} 
$$(\overline{x},\overline{y},\overline{z})= (\frac{\mu_{yz}}{M},\frac{\mu_{xz}}{M},\frac{\mu_{xy}}{M}).$$

\subsection*{Integración de funciones vectoriales} \index{Función!integrable}
Para funciones vectoriales ${\bf f}(t)=\langle f_1(t),f_2(t),\cdots,f_n(t)\ra,$ se dice que $f$ es integrable en $[a,b]$ 
si y solo si cada una de las funciones componentes $f_i(t)$ es integrable en $[a,b]$, para toda $i=1,2,\cdots ,n$.  
Además $$\int_a^b{\bf f}(t)dt= \langle \int_a^b f_1(t)dt,\int_a^b f_2(t)dt,\cdots, \int_a^b f_n(t)dt\ra.$$

\subsection*{Integrales dobles} \label{integralesdobles}  \index{Integrales!dobles}
Presentaremos ahora la integral doble de una funci\'{o}n de
dos variables $f(x,y)$ definida sobre una regi\'{o}n del plano como el 
l\'{\i}mite de las aproximaciones a trav\'{e}s de sumas de Riemann, posteriormente
se presentar\'{a}n de manera similar las integrales triples. 
La evaluaci\'{o}n de las integrales dobles o triples se puede obtener
utilizando el teorema fundamental de c\'{a}lculo; véase \cite{apostol_calculus_1999}.
\vskip0.2cm
En el plano $xy$ consideraremos un rectángulo $R=[a,b]\times \lbrack c,d],$ es decir
$$R=\{(x,y) \in \mathbb{R}^{2}\, | \,a\leq x\leq b,c\leq y\leq d\}.$$
Si $f(x,y)$ es una funci\'{o}n continua definida sobre el rect\'{a}ngulo $R,$ 
subdividiremos a $R$ en $nm$ rect\'{a}ngulos m\'{a}s peque\~{n}os a partir
de rectas paralelas a los ejes coordenados. Esto es, para el intervalo $[a,b],$ consideremos
$$a=x_{0}<x_{1}<x_{2}<\cdots <x_{i}<\cdots <x_{n-1}<x_{n}=b,$$ mientras que
para el intervalo $[c,d],$ consideremos
$$c=y_{0}<y_{1}<y_{2}<\cdots <y_{k}<\cdots <y_{m-1}<y_{m}=d.$$ A esta división de $R$, lo denominaremos una partición de $R$.
\index{Partición}
Si $n$ y $m$ crece entonces el ancho y el largo de estos subintervalos decrecen.
\vskip0.2cm
Denotaremos $\Delta x_{i}$ al ancho del $i-$\'{e}simo subintervalo
sobre el eje $x,$ y $\Delta y_{k}$ es el alto $k-$\'{e}simo subintervalo sobre el eje $y$, 
es decir $\Delta x_{i}=x_{i+1}-x_{i},$ de la misma forma $\Delta
y_{k}=y_{k+1}-y_{k}.$ El \'{a}rea de uno de estos subintervalos de ancho $\Delta x_{i}$ y altura $\Delta y_{k}$ es
$\Delta A_{ik}=\Delta x_{i}\Delta y_{k}$.
La norma de la partici\'{o}n se denota por $\left\Vert P\right\Vert$ y es el m\'{a}ximo sobre todos los valores 
$i$ y $k$, de todas las dimensiones de los rect\'{a}ngulos en que se divide la regi\'{o}n $R.$
\vskip0.2cm
Cuando todos estos subintervalos son de igual longitud, los denotaremos por $\Delta x=\frac{b-a}{n}$ y $\Delta y=\frac{d-c}{m}.$
\vskip0.2cm \index{Suma doble de Riemann}
Una suma de Riemann de $f(x,y) $ sobre el rect\'{a}ngulo $R$ se obtiene al considerar un punto cualquiera 
$\left( x_{i}^{\ast },y_{k}^{\ast}\right) ,$ en $ik$-\'{e}simo subrect\'{a}ngulo, esto 
$\left( x_{i}^{\ast},y_{k}^{\ast }\right) \in \lbrack x_{i-1},x_{i}]\times \lbrack y_{k-1},y_{k}],$ 
o de manera similar $x_{i-1}<x_{i}^{\ast }<x_{i},$ y $y_{k-1}<y_{k}^{\ast }<y_{k};$  
El valor de la funci\'{o}n es este punto es $f\left( x_{i}^{\ast},y_{k}^{\ast }\right) $ y este valor se 
multiplica por el \'{a}rea del $ik$-\'{e}simo rect\'{a}ngulos; se suman todos estos productos para obtener una
expresi\'{o}n de la forma
$$S_{nm}=\sum_{i=1}^{n}\sum_{k=1}^{m}f\left( x_{i}^{\ast },y_{k}^{\ast}\right) \Delta x_{i}\Delta y_{k}.$$

Dependiendo de la selecci\'{o}n $\left( x_{i}^{\ast },y_{k}^{\ast }\right) ,$
es posible obtener diversos valores de $S_{nm}.$
\vskip0.2cm
Es de especial inter\'{e}s el conocer el comportamiento de las sumas de
Riemann cuando la norma de la partici\'{o}n tiende a cero, que equivale a
que el ancho y la altura de cada subrect\'{a}ngulo de $R$ tienden a cero.
Si el l\'{\i}mite de las sumas $S_{nm}$ existe sin importar la selecci\'{o}n
que se haga de $\left( x_{i}^{\ast },y_{k}^{\ast }\right) ,$ entonces
diremos que la funci\'{o}n $f(x,y) $ es integrable sobre el rect\'{a}ngulo $R$ 
y este valor se denota por 
$\ds \iint\limits_{R}f\left(x,y\right) dA$ o tambi\'{e}n $\ds \iint\limits_{R}f(x,y) dydx.$
Lo anterior nos motiva a precisar este concepto.

\begin{defi}
La integral doble de una función continua $f(x,y)$ sobre el rectángulo $R$ se define como
\begin{eqnarray} \label{integralesdobles1}
\lim _{\left\Vert P\right\Vert \to 0} S_{nm} &=& \lim _{n,m \to \infty}\left(\sum_{i=1}^{n}\sum_{k=1}^{m}f\left( x_{i}^{\ast },y_{k}^{\ast}\right) \Delta x_{i}\Delta y_{k}\right) \nonumber \\ &=& \ds \iint\limits_{R}f(x,y) dA
\end{eqnarray}
cuando este límite existe.
\end{defi}

\subsection*{Aplicaciones de las integrales dobles}

Si la funci\'{o}n $f(x,y)$ es positiva sobre el rect\'{a}ngulo $R$ del plano $xy$ la integral doble de $f$ sobre $R$ puede
interpretarse como el volumen del s\'{o}lido $E$ acotado en la base por $R,$
por la parte superior la superficie $z=f(x,y),$ y en los lados regiones contenidas en
los planos $x=a,$ $x=b$, $y=c$, $y=d.$ N\'{o}tese que el t\'{e}rmino
$f(x_{i}^{\ast },y_{k}^{\ast })\Delta x_{i}\Delta y_{k}$ corresponde al
volumen del paralelep\'{\i}pedo con \'{a}rea de la base $\Delta x_{i}\Delta y_{k}$ y altura $f(x_{i}^{\ast },y_{k}^{\ast });$
la suma $S_{nm}$ es una aproximaci\'{o}n al volumen del s\'{o}lido $E.$
Es por lo anterior que el volumen de $E$ se define como 
$$\text{volumen:}\lim _{n,m \to \infty}\left(\sum_{i=1}^{n}\sum_{k=1}^{m}f\left( x_{i}^{\ast },y_{k}^{\ast}\right) \Delta x_{i}\Delta y_{k}\right)= \ds \iint\limits_{R}f(x,y) dA.$$ \index{Masa de una región}

Si la función $f(x,y)$ determina la densidad en el punto $(x,y)$ entonces la integral de la fórmula (\ref{integralesdobles1})
representa la masa ($M$) de la región $R$. Los momentos con respecto al eje $x$ y al eje $y$ están dados  por
{\small $\mu_{x}=\ds \iint_R y f(x,y)dA$} y {\small $\mu_{y}=\ds \iint_R x f(x,y)dA$,} los cuales determinan las coordenadas del 
centro de masa de la región $R$, esto es $$(\overline{x},\overline{y})= \left(\frac{\mu_y}{M},\frac{\mu_x}{M}\right).$$
\index{Centro de masa}

Si la región de integración en la fórmula (\ref{integralesdobles1}) es un rectángulo está se evalúa
$$\ds \iint\limits_{R}f(x,y) dA= \int_c^d \left (\int_a^b f(x,y) dx \right )dy.$$
\vskip0.2cm
El teorema de Fubini, cuya demostración puede consultarse en \cite{marsden_calculo_2013} o \cite{apostol_alisis_1981},
es el resultado que permite intercambiar el orden de integración incluso en regiones más generales. 
\vskip0.2cm
Si la región de integración no es un rectángulo, 
la integral doble de una función $f(x,y)$ sobre una región $R$ puede evaluarse mediante la aplicación del Teorema de 
Fubini por alguna de las siguientes expresiones, dependiendo de la forma de la región $R$.
\vskip0.2cm
Si la región del plano en mención, se describe en la forma $\phi_{1}(x) \leq y \leq \phi_{2}(x)$ con $a \leq x \leq b$, 
diremos que la región es de tipo I; mientras que la región del tipo II es aquella que puede ser descrita como
$\varphi_{1}(y) \leq x \leq \varphi_{2}(y)$ con $c \leq y \leq d.$ 

\vskip0.2cm

De acuerdo con lo anterior, el planteamiento de la integral doble es, según el caso, el siguiente:

\vskip0.2cm
\noindent
\begin{minipage}{\textwidth}
	\centering
	\captionof{figure}{Regiones de integración en el plano}
	\label{Regiones_integracion_plano}
	\begin{minipage}{0.35\linewidth}
		\centering
		\includegraphics[width=\linewidth]{Fig/7.pdf}
	\end{minipage}
	\begin{minipage}{0.35\linewidth}
		\centering
		\includegraphics[width=\linewidth]{Fig/8.pdf}
	\end{minipage}
	
	\vspace{0.3em}
	\parbox{0.95\linewidth}{\centering\fontsize{8pt}{10pt}\selectfont\textbf{Fuente:} elaboración propia.}
\end{minipage}
\vskip0.2cm



$$ \iint_{R}f(x,y) dA=\int_{a}^{b}\int_{\phi _{1}(y) }^{\phi_{2}(y)}f(x,y)dydx,$$

o también
\vskip0.2cm
$$ \iint_{R}f(x,y)dA=\int_{c}^{d}\int_{\varphi _{1}(y) }^{\varphi _{2}(y)}f(x,y)dxdy.$$

%\vskip0.2cm

\begin{ejemplo}
Consideremos la función $f(x,y)=8+2x-x^3-4y$ y algunas aproximaciones para diferentes valores de $n$ y $m$;
con $x\in [0,2]$, $y \in [0,1].$

%\captionsetup{hypcap=false}
%\captionof{figure}{\small Aproximaciones al volumen bajo $f(x,y)$}
\vspace{4mm}
\noindent % Evita sangrías no deseadas
\begin{minipage}{\textwidth}
	\centering
	% Usamos \captionof porque estamos fuera de un entorno figure
	\captionof{figure}{Aproximaciones para el volumen bajo $f(x,y) = 8+2x-x^3-4y$.}
	\label{fig:aproximaciones_volumen}
	
	\vspace{0.3cm} % Espacio entre título e imágenes
	
	\begin{tabular}{ccc}
		% PRIMERA FILA: Las imágenes con escalas homogeneizadas
		\includegraphics[scale=0.31]{Img/T_51a.pdf} & 
		\includegraphics[scale=0.28]{Img/T_51b.pdf} &
		\includegraphics[scale=0.28]{Img/T_51c.pdf} \\
		\small $n=5, m=9$ & \small $n=6, m=11$ & \small $n=17, m=29$
	\end{tabular}
	
	\vskip0.2cm
	\fuente{elaboración propia.} 
\end{minipage}
\vspace{3em} % Espacio para separar la minipage del siguiente párrafo

\end{ejemplo}

%Los gráficos anteriores se obtienen con las siguientes instrucciones. 
%\begin{Verbatim}[fontsize=\letraverbatim]
%DiscretePlot3D[8+2*x-x^3-4*y,{x,0,2,0.05},{y,0,1,0.05},
%ExtentSize->Left,BoxRatios->{1,1,1}] \end{Verbatim} 
\vspace{-8mm}
El volumen bajo la figura de $f(x,y)$ sobre el rectángulo $R$ puede verse también como la integral de $A(x),$ 
la cual representa el área de las secciones transversales del sólido;
por tanto el volumen está dado por {\small $\ds \int_a ^b A(x) dx= \int_a ^b \int_c ^d f(x,y) dy dx $.} O de manera similar
{\small $\ds \int_c ^d A(y) dy= \int_c ^d \int_a ^b  f(x,y) dx dy.$}

\vskip0.4cm

Ejemplo de lo anterior es la siguiente integral
que representa el volumen del sólido en mención $$\ds \int_0^2\int_0^1 (8+2x-x^3-4y) dydx = \int_0^1 \int_0^2(8+2x-x^3-4y) dxdy.$$

\begin{figure}[H] 
	\centering
	\caption{Secciones transversales bajo $z = 8+2x - x^3 - 4y$, con $x \in [0,2]$, $y \in [0,1]$.}
	\label{fig:secciones_transversales_bajo}
	\begin{minipage}[b]{0.22\linewidth}
		\centering
		\includegraphics[width=\linewidth]{Img/T_52a.pdf}
	\end{minipage}%
	\begin{minipage}[b]{0.22\linewidth}
		\centering
		\includegraphics[width=\linewidth]{Img/T_52b.pdf}
	\end{minipage}%
	\begin{minipage}[b]{0.22\linewidth}
		\centering
		\includegraphics[width=\linewidth]{Img/T_52c.pdf}
	\end{minipage}%
	\begin{minipage}[b]{0.22\linewidth}
		\centering
		\includegraphics[width=\linewidth]{Img/T_52d.pdf}
	\end{minipage}
	\fuente{elaboración propia.} 
\end{figure}

\subsection*{Integrales triples} \index{Integrales!triples}

Consideraremos una caja rectangular cerrada y acotada $E$ en el espacio con
caras paralelas a los planos coordenados $E=[a,b]\times \lbrack c,d]\times
\lbrack r,s];$ es decir, 
\begin{equation*}
E=\{(x,y,z) \in \mathbb{R}^{3}\,|\,a\leq x\leq b,c\leq y\leq
d,r\leq z\leq s\}.
\end{equation*}
Si $f(x,y,z)$ es una funci\'{o}n continua definida sobre la caja rectangular $E,$
subdividiremos a $E$ en $lnm$ cajas rectangulares m\'{a}s peque\~{n}as a
partir de planos paralelos a los planos coordenados. Esto es, para el
intervalo $[a,b],$ consideremos 
\begin{equation*}
x_{0}=a<x_{1}<x_{2}<\cdots <x_{i}<\cdots <x_{n-1}<x_{n}=b,
\end{equation*}
para el intervalo $[c,d],$ 
\begin{equation*}
y_{0}=c<y_{1}<y_{2}<\cdots <y_{j}<\cdots <y_{m-1}<y_{m}=d,
\end{equation*}
y para el intervalo $[r,s],$ consideremos 
\begin{equation*}
z_{0}=r<z_{1}<z_{2}<\cdots <z_{k}<\cdots <y_{l-1}<y_{l}=s.
\end{equation*}

Si $n,$ $l$ y $m$ crecen entonces el ancho, el largo y el alto de estas
cajas rectangulares decrecen.
Denotaremos $\Delta x_{i}$ al ancho del $i$-\'{e}simo subintervalo sobre el
eje $x,$ $\Delta y_{j}$ es el largo $j$-\'{e}simo subintervalo sobre el eje $y$, y $\Delta z_{k}$ es el alto $k$-\'{e}simo subintervalo sobre el eje $z$, es decir $\Delta x_{i}=x_{i+1}-x_{i},\Delta y_{j}=x_{j+1}-x_{j},$ $\Delta
z_{k}=z_{k+1}-z_{k}.$ 
El volumen de cada una de estas cajas de ancho $\Delta x_{i}$, largo $\Delta y_{j}$ y alto $\Delta z_{k}$ es 
$\Delta V_{ijk}=\Delta x_{i}\Delta y_{j}\Delta z_{k}$. 
La norma de la partici\'{o}n se denota por $\left\Vert P\right\Vert $ y es
el m\'{a}ximo sobre todos los valores $i,$ $j$ y $k$ de todas las
dimensiones de las cajas en que se divide la regi\'{o}n $E.$ 
\vskip0.2cm
Una suma de Riemann de $f(x,y,z) $ sobre la caja rectangular $E$ se consigue al considerar un punto cualquiera 
$\left( x_{i}^{\ast},y_{j}^{\ast },z_{k}^{\ast }\right) ,$ en la $ijk$-\'{e}sima caja
rectangular, esto es $\left( x_{i}^{\ast },y_{j}^{\ast },z_{k}^{\ast}\right) \in \lbrack x_{i-1},x_{i}]\times \lbrack y_{j-1},y_{j}]\times \lbrack z_{k-1},z_{k}],$ lo cual equivale a que $x_{i-1}<x_{i}^{\ast }<x_{i},$ 
$y_{j-1}<y_{j}^{\ast }<y_{j}$ y $z_{k-1}<z_{k}^{\ast }<z_{k}.$ El valor de la
funci\'{o}n es este punto es $f\left( x_{i}^{\ast },y_{j}^{\ast},z_{k}^{\ast }\right) $ 
y este valor se multiplica por el volumen de la $ijk$-\'{e}sima caja rectangular; se suman todos estos productos para obtener
una expresi\'{o}n de la forma 
\begin{equation*}
S_{nml}=\sum_{i=1}^{n}\sum_{j=1}^{m}\sum_{k=1}^{l}f\left( x_{i}^{\ast},y_{j}^{\ast },z_{k}^{\ast }\right) \Delta x_{i}\Delta y_{j}\Delta z_{k}.
\end{equation*}
Dependiendo de la elecci\'{o}n de $\left( x_{i}^{\ast },y_{j}^{\ast },z_{k}^{\ast }\right)$
es posible obtener diversos valores de $S_{nml}.$ 
De manera similar a como se defini\'{o} la integral doble, si el l\'{\i}mite de
las sumas $S_{nml}$ existe para toda elecci\'{o}n que se haga de 
$\left(x_{i}^{\ast },y_{j}^{\ast },z_{k}^{\ast }\right),$ entonces se dice que la
funci\'{o}n $f(x,y,z) $ es integrable sobre la caja rectangular $E,$ y este valor se denota por 
$\ds\iiint\limits_{E}f(x,y,z) dV$ o tambi\'{e}n $\ds \iiint\limits_{E}f(x,y,z) dzdydx.$ Lo anterior nos
motiva a precisar este concepto.

\begin{defi}
La integral triple de una función continua $f(x,y,z)$ 
sobre la caja rectangular $E$ se define como 
\begin{eqnarray}
\lim_{\left\Vert P\right\Vert \rightarrow 0}S_{nml} &=& \lim_{l,m,n \to\infty}\left(\sum_{i=1}^{n}\sum_{j=1}^{m}\sum_{k=1}^{l}f\left( x_{i}^{\ast },y_{j}^{\ast },z_{k}^{\ast}\right) \Delta x_{i}\Delta y_{j}\Delta z_{k} \right)\nonumber \\ &=&\iiint\limits_{E}f\left(x,y,z\right) dV 
\end{eqnarray}
cuando este l\'{\i}mite existe.
\end{defi}
Ilustremos lo anterior con dos integrales triples.
\begin{ejemplo}
Evaluar {\small $\ds \int_{-1}^1\int_0^1\int_0^1  yz^2 e^{x} dzdydx$} mediante sumas de Riemann.
\vskip0.2cm
Al hacer una partición de la región de integración en $nml$ cajas rectangulares se sigue que 
$\Delta x= \frac{2}{n}$ $\Delta y= \frac{1}{m}$ y $\Delta z= \frac{1}{l}$. Además
$x_i= -1+\frac{2i}{n},$ $y_j=\frac{j}{m}$ y $z_k=\frac{k}{l}.$
Por lo tanto 
$$f(x_i,y_j,z_k)\Delta x \Delta y\Delta z=\frac{2 j k^2 e^{\frac{2 i}{n}-1}}{l^3 m^2 n}$$
Con lo cual 
$$\sum_{i=1}^n\sum_{j=1}^m\sum_{k=1}^l f(x_i,y_j,z_k)\Delta x \Delta y\Delta z=\frac{\left(e^2-1\right) (l+1)(2l+1)(m+1)}{6nml^2 \left(1-e^{-2/n}\right)e}$$ entonces 
\begin{multline*} \int_{-1}^1\int_0^1\int_0^1  yz^2 e^{x} dzdydx= \lim_{n,m,l \to \infty}\sum_{i=1}^n\sum_{j=1}^m\sum_{k=1}^l f(x_i,y_j,z_k)\Delta x \Delta y\Delta z \\ =\frac{1}{6} \left(e-\frac{1}{e}\right)\end{multline*}
\end{ejemplo}

\begin{comment}
Con las siguientes instrucciones se evalúan los cálculos anteriores.
 \begin{Verbatim}[fontsize=\letraverbatim]
f[x_,y_,z_]:=z^2*y*Exp[x]
{xi,yj,zk}={-1+2*i/n,j/m,k/l};
{dx,dy,dz}={2/n,1/m,1/l};
suma=Sum[f[xi,yj,zk]*dx*dy*dz,{i,1,n},{j,1,m},{k,1,l}]
Limit[suma,{n,m,l}->{Infinity,Infinity,Infinity}]
Integrate[f[x,y,z],{x,-1,1},{y,0,1},{z,0,1}]
\end{Verbatim}

\end{comment}

\begin{ejemplo}
\begin{sloppypar}
Para $f(x,y,z)= \cos (x) \sin (y)+z^2 \cos^2 (2 y)$ y $E$ el \mbox{paralelepípedo} determinado por 
$[-\frac{\pi}{2},\frac{\pi}{2}]\times [0,\frac{\pi}{4}]\times[0,1]$; evaluar 
{\footnotesize $\ds \iiint_Ef(x,y,z) dzdydx.$}
\vskip0.2cm
Al hacer una partición de $E$ en $nml$ cajas rectangulares es posible establecer que
$\Delta x= \frac{\frac{\pi}{2}-\left(\frac{-\pi}{2}\right)}{n}=\frac{\pi}{2n}$
$\Delta y= \frac{\frac{\pi}{4}-0}{m}=\frac{\pi}{4m}$ y $\Delta z= \frac{1}{l}$. Además
$x_i= -\frac{\pi}{2}+\frac{\pi i}{2n},$ $y_j=\frac{\pi j}{4m}$ y $z_k=\frac{k}{l}.$
Por lo tanto: 
$$f(x_i,y_j,z_k)=\sin\left(\frac{\pi i}{n}\right)\sin\left(\frac{\pi j}{4m}\right)+\frac{k^2}{l^2}\cos^2\left(\frac{\pi j}{2m}\right).$$
Con lo anterior, se establece que $S_{nml}=\sum\limits_{i=1}^n\sum\limits_{j=1}^m\sum\limits_{k=1}^l f(x_i,y_j,z_k)\Delta x \Delta y\Delta z,$ es
{\small 
$$\frac{\pi ^2}{48 m}\left(\frac{(l+1) (2 l+1) (m-1)}{l^2}+\frac{12 \sin \left(\frac{\pi}{8}\right) \sin \left(\frac{\pi  (m+1)}{8 m}\right) \csc \left(\frac{\pi }{8 m}\right)\cot \left(\frac{\pi }{2 n}\right)}{n}\right).$$
}
De modo que $\lim_{n,m,l \to \infty}S_{nml}$ es
\begin{multline*}
\int _{-\frac{\pi }{2}}^{\frac{\pi }{2}}\int _0^{\frac{\pi }{4}}\int _0^1\left(\cos(x)\sin (y)+z^2 \cos ^2(2 y)\right)dzdydx  =2-\sqrt{2}+\frac{\pi ^2}{24}
\end{multline*}
\end{sloppypar}
\end{ejemplo}

\begin{comment}
 \begin{Verbatim}[fontsize=\letraverbatim]
f[x_,y_,z_]:=Cos[x]*Sin[y]+z^2*Cos[2*y]^2
{xi,yj,zk}={-Pi/2+Pi*i/n,Pi/4*j/m,k/l};
{dx,dy,dz}={Pi/n,Pi/(4*m),1/l};
suma=Simplify[
Sum[f[xi,yj,zk]*dx*dy*dz,{i,1,n},{j,1,m},{k,1,l}]];
Limit[suma,{n,m,l}->{Infinity,Infinity,Infinity}]
Integrate[f[x,y,z],{x,-Pi/2,Pi/2},{y,0,Pi/4},{z,0,1}]
\end{Verbatim} 
\end{comment}
Utilizando \verb"Mathematica" es posible simplificar y evaluar las integrales anteriores. 
Su uso se explicará en detalle en los próximos capítulos.
\vskip0.1cm 
%

Para una región acotada sólida $E$ en el espacio, se sigue un proceso análogo al utilizado en las 
integrales dobles. Nótese que $z$ recorre los puntos entre las superficies $z=u_1(x,y)$ y
$z=u_2(x,y)$. La proyección de $E$ en el plano $xy$ muestra que $y$ recorre los puntos que están incluidos entre las curvas $y=h_1(x)$ y $y=h_2(x),$ con $x\in [a,b].$ 


\vskip0.2cm
\noindent
\begin{minipage}{\textwidth}
	\centering
	\captionof{figure}{Integración en el espacio}
	\begin{minipage}{0.39\linewidth}
		\centering
		\includegraphics[width=\linewidth]{Fig/9.pdf}
	\end{minipage}
	
	\vspace{0.3em}
	\parbox{0.95\linewidth}{\centering\fontsize{8pt}{10pt}\selectfont\textbf{Fuente:} elaboración propia.}
\end{minipage}
\vskip0.2cm


Es por esto que la integral respectiva se escribe en la forma
$$\int_a^b\int_{h_1(x)}^{h_2(x)} \int_{u_1(x,y)} ^{u_2(x,y)} f(x,y,z) dz dydx.$$ 

Se debe observar que en algunos casos la proyección del sólido en el plano $xy$ puede corresponder a regiones del tipo I,
tipo II, las cuales fueron descritas anteriormente, o también combinaciones de estas.
\vskip0.2cm
En estos apuntes se consideran otras posibilidades e intercambios en el orden de integración, que dependiendo de la naturaleza de cada problema harán más o menos complicado el tratamiento a estas integrales.

\subsection*{Integrales múltiples sobre regiones parametrizadas}

En \cite{acosta_gempeler_doce_2020}, se considera una función continua $f(x,y,z)$ y una superficie $\sigma$
parametrizada en la forma ${\bf r}(u,v)= \langle x(u,v), y(u,v), z(u,v) \ra,$
con $u \in [a,b]$ y $v \in [c,d];$ cuyas funciones componentes poseen derivadas parciales continuas en 
$[a,b] \times [c,d].$ Al realizar una partición de $\sigma$ en la forma
$$\sigma_{ij}= \{ (x(u_i^\ast,v_j^\ast),y(u_i^\ast,v_j^\ast),z(u_i^\ast,v_j^\ast) )\,|\, u_{i-1} \leq u_i^\ast \leq u_{i}, v_{j-1} \leq v_{j}^\ast \leq v_{j}\} ;$$
considerando el punto $P_{ij}(x(u_i^\ast,v_j^\ast),y(u_i^\ast,v_j^\ast),z(u_i^\ast,v_j^\ast) )$ en cada trozo $\sigma_{ij}$
de la superficie, se obtiene la suma de Riemann de $f$ sobre la superficie $\sigma$ parametrizada por ${\bf r}(u,v)$ y con la partición de $\sigma$,
$$ S_{nm}=\sum_{i=1}^n \sum_{j=1}^m f(x(u_i^\ast,v_j^\ast),y(u_i^\ast,v_j^\ast),z(u_i^\ast,v_j^\ast) ) \Vert {\bf r}_u \times {\bf r}_v \Vert \Delta u_i \Delta v_j .$$
El término $\Vert {\bf r}_u \times {\bf r}_v \Vert $ es el área del paralelogramo determinado por
los vectores ${\bf r}_u$ y ${\bf r}_v.$
Con un procedimiento como el utilizado cuando presentamos integrales dobles, se sigue que en este caso la integral de $f$ 
sobre la superficie $\sigma$ con la parametrización ${\bf r}(u,v)$ está dada por
\begin{multline}
\iint_\sigma f dA = \lim_{n,m \to \infty} S_{nm}  = \int_a^b\int_c^d f({\bf r}(u,v)) \Vert {\bf r}_u \times {\bf r}_v \Vert dudv  \label{integralesdobles2}
\end{multline}
Si la parametrización ${\bf r}(u,v,w)= \langle x(u,v,w), y(u,v,w), z(u,v,w) \ra,$
de un sólido $E$, con $u \in [a,b],$ $v \in [c,d]$ y $w \in [p,q]$;
con funciones componentes que poseen derivadas parciales continuas en 
$[a,b] \times [c,d] \times [p,q].$

\begin{figure}[H]
	\centering
	\caption{Transformación de una región en el espacio}
	\label{fig:transform-region}
	\begin{minipage}[b]{0.30\linewidth}
		\centering
		\includegraphics[width=\linewidth]{Fig/10.pdf}
	\end{minipage}
	\begin{minipage}[b]{0.30\linewidth}
		\centering
		\includegraphics[width=\linewidth]{Fig/11.pdf}
	\end{minipage}
	\begin{minipage}[b]{0.30\linewidth}
		\centering
		\includegraphics[width=\linewidth]{Fig/12.pdf}
	\end{minipage}
	\parbox{0.95\linewidth}{\centering\fontsize{9pt}{10pt}\selectfont\textbf{Fuente:} elaboración propia.}
\end{figure}
\vskip0.2cm

Es posible realizar una partición de $E$ en la forma
\[E_{ijk}= \{ P_{ijk}^\ast\,|\, u_{i-1} \leq u_i^\ast \leq u_{i}, v_{j-1} \leq v_{j}^\ast \leq v_{j}, w_{k-1} \leq w_{k}^\ast \leq w_{k}\};\]
En donde $P_{ijk}^\ast=(x(u_i^\ast,v_j^\ast,w_k^\ast),y(u_i^\ast,v_j^\ast,w_k^\ast),z(u_i^\ast,v_j^\ast,w_k^\ast) );$
al considerar el punto $P_{ijk}(x(u_i^\ast,v_j^\ast,w_k^\ast),y(u_i^\ast,v_j^\ast,w_k^\ast),z(u_i^\ast,v_j^\ast,w_k^\ast) )$
en cada trozo $E_{ijk}$ del sólido, se obtiene la suma de Riemann de $f$ sobre el sólido $E$ parametrizado por 
${\bf r}(u,v,w)$ y con la partición de $E$,
$$ S_{nml}=\sum_{i=1}^l \sum_{i=1}^n \sum_{j=1}^m f(P_{ijk}^\ast) | {\bf r}_u \cdot ({\bf r}_v \times {\bf r}_w) | \Delta u_i \Delta v_j \Delta w_k.$$
El término $| {\bf r}_u \cdot ({\bf r}_v \times {\bf r}_w) |$ es el volumen del paralelepípedo determinado por
los vectores ${\bf r}_u,$ ${\bf r}_v$ y ${\bf r}_w$. 
De manera análoga al procedimiento utilizado cuando presentamos integrales triples, se sigue que la integral de $f$ 
sobre $E$ con la parametrización ${\bf r}(u,v,w)$ está dada por
\begin{multline} \label{math1}
\ds \iiint_E f dV = \lim_{l,m,n \to \infty} S_{nml}  \\ 
=\int_a^b\int_c^d \int_p^q f({\bf r}(u,v,w)) | {\bf r}_u \cdot ({\bf r}_v \times {\bf r}_w) | dudvdw 
\end{multline}
Al término $|{\bf r}_u \cdot ({\bf r}_v \times {\bf r}_w) |dudvdw$ lo denominaremos elemento diferencial de volumen.
\index{Elemento diferencial!de volumen}
Esta última expresión puede verse como la integral de la función $f$ si el sólido $E^*$ es
parametrizado por la función ${\bf r}(u,v,w).$ Si en la expresión (\ref{math1}) $f(x,y,z)=1$ esta integral determina 
el volumen $V$ del sólido $E$, es decir
\begin{eqnarray} \label{math2}
V=\iiint_{E}|{\bf r}_u \cdot ({\bf r}_v \times {\bf r}_w) | dudvdw
\end{eqnarray}

\subsection*{Aplicaciones de las integrales triples}

De manera similar al caso bidimensional la masa $M,$ de un sólido $E$ con densidad $\rho (x,y,z)$ se determina por 
$$M=\ds \iiint_E \rho (x,y,z) dV$$ y los momentos con respecto a los planos coordenados se determinan por 
$$\ds \mu_{xy}=\iiint_E z\rho(x,y,z) dV,\,\,\, \mu_{xz}=\iiint_E y\rho(x,y,z) dV,$$
$$\ds  \mu_{yz}=\iiint_E x\rho(x,y,z)dV.$$
De tal forma que el centro de masa de $E$ es
$(\overline{x},\overline{y},\overline{z})= (\frac{\mu_{yz}}{M}, \frac{\mu_{xz}}{M},\frac{\mu_{xy}}{M}) .$
Una exposición más detalladas de estos elementos y otros usos de las integrales múltiples puede consultarse en \cite{thomas_calculo_2015}.
\index{Centro de masa}
\subsection*{Integrales de superficie}
Una parametrizaci\'{o}n de una superficie es una funci\'{o}n 
$\phi:D\subseteq \mathbb{R}^{2}\rightarrow \mathbb{R}^{3}.$
La superficie $S$ correspondiente a la
funci\'{o}n $\phi $ es su imagen $\phi \left( D\right) =S$ se escribe 
$\phi(u,v)=\left\langle x\left( u,v\right) ,y\left( u,v\right),z\left( u,v\right) \right\rangle .$
La integral de un campo escalar $f$ sobre una superficie $S,$  está dada por

\begin{equation} \label{areasup1}
\iint_{S}f(x,y,z) dS=\iint_{D}f\left( \phi \left( u,v\right)
\right) \left\| {\bf T}_{u}\times {\bf T}_{v}\right\| dudv
\end{equation}

siendo
$${\bf T}_{u}=\frac{\partial \phi }{\partial u}=\frac{\partial x}{\partial u}\vi+\frac{\partial x}{\partial u}\vj+\frac{\partial x}{\partial u}\vk$$
y $${\bf T}_{v}=\frac{\partial \phi }{\partial v}=\frac{\partial x}{\partial v}\vi+\frac{\partial x}{\partial v}\vj+\frac{\partial x}{\partial v}\vk.$$

La integraci\'{o}n de una campo vectorial ${\bf f}$ sobre $S,$ está dada por
\begin{equation}
\iint_{S} {\bf f}(x,y,z) \cdot dS=\iint_{D} {\bf f} \left( \phi \left(
u,v\right) \right) \cdot \left( {\bf T}_{u}\times {\bf T}_{v}\right) dudv
\end{equation}
Una exposición más detallada de esta formulación puede encontrarse en \cite{marsden_calculo_2013}.
\vskip0.2cm
Si en la expresión (\ref{areasup1}) $f(x,y,z)=1,$ esta integral se reduce al área $A$ de la superficie $D,$ es decir
\begin{equation} \label{areasup}
A=\iint_{D} \left\| {\bf T}_{u}\times {\bf T}_{v}\right\| dudv
\end{equation}
Al término $\left\| {\bf T}_{u}\times {\bf T}_{v}\right\|dudv$ lo denominaremos elemento diferencial de área.
\index{Elemento diferencial!de área}

\subsection*{Cambio de variables en integrales dobles} \label{cambioidobles}
\vskip0.2cm
Con el propósito de evaluar la integral $\ds \iint\limits_{R}f(x,y) dA$
se hace el cambio $(x,y) ={\bf r}(u,v)$ para transformar la región $R$ en el sistema 
$(x,y) $ en una región $R^{\ast }$ en el sistema coordenado  $(u,v).$
Un rectángulo en el sistema $uv$ con lados de longitud $\Delta u$ de largo y $\Delta v$ de ancho se transforma mediante la función ${\bf r}.$
\vskip0.2cm
Asumiendo que ${\bf r}$ es diferenciable, si $(u,v) $ son las coordenadas del rectángulo, los v\'{e}rtices
adyacentes son $( u+\Delta u,v) $ y $(u,v+\Delta v)$; mediante la aplicación ${\bf r},$ el rectángulo
se transforma en una región cuya área se puede aproximar mediante el área de un paralelogramo. Además, 
aplicando el teorema del valor medio, se sigue que

\begin{eqnarray*}
{\bf r}(u+\Delta u,v) -{\bf r}(u,v) &=& \frac{{\bf r}( u+\Delta u,v)-{\bf r}(u,v)}{\Delta u}\Delta u 
\approx \frac{\partial {\bf r}}{\partial u}\left( u^{\ast},v\right)\Delta u \\
{\bf r}(u,v+\Delta v) -{\bf r}(u,v) &=& \frac{{\bf r}(u,v+\Delta v)-{\bf r}(u,v)}{\Delta v}\Delta v 
\approx \frac{\partial{\bf r}}{\partial u}( u,v^{\ast})\Delta v.
\end{eqnarray*}

El área de esta pequeña región est\'{a} dada por 
\begin{equation}
\Delta A=\left\Vert \frac{\partial {\bf r}}{\partial u}\Delta u\times 
\frac{\partial {\bf r}}{\partial v}\Delta v\right\Vert =\left\Vert \frac{\partial {\bf r}}{\partial u}\times \frac{\partial {\bf r}}{\partial v}\right\Vert \Delta u\Delta v
\end{equation}
que se puede abreviar en la forma
$\Delta A=\left\vert \frac{\partial (x,y)}{\partial (u,v)}\right \vert \Delta u\Delta v.$
Por simplicidad $x=x(u,v)$ y $y=y(u,v)$, entonces
\begin{equation}
\left\vert \frac{\partial (x,y) }{\partial(u,v)}\right\vert =\left\Vert\frac{\partial{\bf r}}{\partial u}\times \frac{\partial{\bf r}}{\partial v}\right\Vert =
\left\vert \begin{array}{cc}\frac{\partial x}{\partial u}&\frac{\partial x}{\partial v}\\[.1cm]  
\frac{\partial y}{\partial u} & \frac{\partial y}{\partial v}
\end{array} \right\vert
\end{equation}

En algunos libros como \cite{stewart_calculo_2018}, el t\'{e}rmino $\left\vert \frac{\partial\left(x,y\right)}{\partial(u,v)}\right\vert$ 
se denomina el jacobiano de la transformación ${\bf r}$. \index{Jacobiano}
Con los elementos anteriores se establece la siguiente igualdad.
\begin{equation}
\iint\limits_{R}f(x,y) dA=\iint\limits_{R^{\ast }}f\left({\bf r}(u,v) \right) \left\vert \frac{\partial (x,y) }{\partial (u,v)}\right\vert dudv \end{equation}

\subsection*{Cambio de variables en integrales triples}
Con el propósito de evaluar la integral 
$$\iiint\limits_{R}f(x,y,z) dV$$
se puede hacer el cambio $(x,y,z) ={\bf r}(u,v,w)$ para transformar la región $R$ en el sistema 
de coordenadas $(x,y,z) $ en una region $R^{\ast }$ en el sistema coordenado  $(u,v,w) .$
Consideremos un cubo en el sistema $uvw,$ con aristas de longitud $\Delta u$ de largo, $\Delta v$ de ancho y $\Delta w$ de altura, que se transforma mediante la función ${\bf r}$. \vskip0.2cm 
Asumiendo que ${\bf r}$ es diferenciable,
si $(u,v,w) $ son las coordenadas de un vértice del cubo, entonces los v\'{e}rtices
adyacentes son $\left( u+\Delta u,v,w\right) ,$ $\left( u,v+\Delta v,w\right)$ y $\left( u,v,w+\Delta w\right) ;$ mediante la aplicación ${\bf r},$ el cubo se transforma en una región cuyo volumen se puede aproximar mediante un \mbox{paralelep\'{\i}pedo}.
Aplicando el teorema del valor medio, se sigue que

\begin{eqnarray*}
{\bf r}\left( u+\Delta u,v,w\right) -{\bf r}(u,v,w) &=& \frac{{\bf r}\left( u+\Delta u,v,w\right) -{\bf r}(u,v,w) }{\Delta u}\Delta u \\[0.1cm] 
&\approx &   \frac{\partial {\bf r}}{\partial u}\left( u^{\ast},v,w\right) \Delta u \\
{\bf r}\left( u,v+\Delta v,w\right) -{\bf r}(u,v,w) &=& \frac{{\bf r}(u,v+\Delta v,w)-{\bf r}(u,v,w)}{\Delta v}\Delta v \\[0.1cm]
&\approx& \frac{\partial{\bf r}}{\partial v}\left( u,v^{\ast},w\right)\Delta v\\
{\bf r}\left( u,v,w+\Delta w\right) -{\bf r}(u,v,w) &=& \frac{{\bf r}\left(
u,v,w+\Delta w\right) -{\bf r}(u,v,w)}{\Delta w}\Delta w \\[0.1cm]
&\approx & \frac{\partial {\bf r}}{\partial w}\left( u,v,w^{\ast}\right) \Delta w,
\end{eqnarray*}

El volumen de esta región, $\Delta V$, est\'{a} dada por 
\begin{equation}
\left\vert \left( \frac{\partial {\bf r}}{\partial u}\Delta u\times 
\frac{\partial {\bf r}}{\partial v}\Delta v\right) \cdot \frac{\partial {\bf r}}{\partial w}\Delta w\right\vert =\left\vert \left( \frac{\partial {\bf r}}{\partial u}\times \frac{\partial {\bf r}}{\partial v}\right) \cdot \frac{\partial {\bf r}}{\partial w}\right\vert \Delta u\Delta v\Delta w
\end{equation}
que se puede abreviar en la forma
$\Delta V=\left\vert \frac{\partial (x,y,z) }{\partial \left(
u,v,w\right) }\right\vert \Delta u\Delta v\Delta w.$
Por simplicidad, si $x=x(u,v,w) ,$ $y=y(u,v,w) $ y $z=z(u,v,w) $ entonces
\begin{equation}
\left\vert \frac{\partial (x,y,z) }{\partial(u,v,w)}\right\vert =\left\vert\left(\frac{\partial{\bf r}}{\partial u}\times \frac{\partial{\bf r}}{\partial v}\right) \cdot \frac{\partial{\bf r}}{\partial v}\right\vert =
\left\vert 
\begin{array}{ccc}\frac{\partial x}{\partial u}&\frac{\partial x}{\partial v}&\frac{\partial x}{\partial w}\\[.1cm]  
\frac{\partial y}{\partial u}&\frac{\partial y}{\partial v}&\frac{\partial y}{\partial w} \\[0.1cm] 
\frac{\partial z}{\partial u}&\frac{\partial z}{\partial v}&\frac{\partial z}{\partial w}\end{array}\right\vert 
\end{equation}

En algunos libros el t\'{e}rmino $\left\vert \frac{\partial\left(x,y,z\right)}{\partial(u,v,w)}\right\vert$ 
se denomina el jacobiano de la transformación ${\bf r}$. \index{Jacobiano}
Esto muestra la siguiente igualdad.
\begin{equation}
\iiint\limits_{R}f(x,y,z) dV=\iiint\limits_{R^{\ast}}f\left({\bf r}(u,v,w) \right)\left\vert\frac{\partial (x,y,z)}{\partial (u,v,w) }\right\vert dudvdw
\end{equation}

Los siguientes casos particulares de cambio de variables en integrales dobles y triples son utilizados 
de manera frecuente.
\subsection*{Coordenadas polares, cilíndricas y esféricas} \label{esfericas}
\index{Coordenadas!polares} \index{Coordenadas!cilíndricas} \index{Coordenadas!esféricas}
Para una circunferencia de radio $r,$ la longitud $s$, de un arco de $\theta$ radianes es $r\theta$.
El área de un sector circular de arco $\Delta \theta$ y radio $r$ está dado por 
$\frac{1}{2}r^2 \Delta \theta$ y el área $\Delta A,$ de la región acotada por dos sectores 
circulares de radio $r_i$ y $r_{i+1}$ (con $r_i<r_{i+1}$) está dada por
\begin{eqnarray*}
\Delta A &=& \frac{1}{2}r_{i+1}^2 \Delta \theta-\frac{1}{2}r_{i}^2 \Delta \theta 
= \frac{r_{i+1}+r_i}{2}(r_{i+1}-r_i)\Delta \theta  = r_i^*\Delta r \Delta \theta 
\end{eqnarray*}

%\begin{figure}
    %\centering
    %\includegraphics[width=0.5\linewidth]{Fig/12.pdf}
    %\caption{Transformación a coordenadas polares}
    %\label{fig:polar-transform}
%\end{figure}
\noindent
\begin{minipage}{\textwidth}
	\centering
	\captionof{figure}{Transformación a coordenadas polares}
	\label{fig:falta_título}
	\begin{minipage}{0.28\linewidth}
		\centering
		\includegraphics[width=\linewidth]{Fig/13.pdf}
	\end{minipage}
	\hspace{1em}
	\begin{minipage}{0.28\linewidth}
		\centering
		\includegraphics[width=\linewidth]{Fig/14.pdf}
	\end{minipage}
	
	\vspace{0.3em}
	\parbox{0.95\linewidth}{\centering\fontsize{9pt}{10pt}\selectfont\textbf{Fuente:} elaboración propia.}
\end{minipage}




Teniendo en cuenta las relaciones anteriores y el teorema de cambio de variables se puede establecer que en el sistema de coordenadas polares, $x=r \cos \theta$, $y=r \sin \theta$,  la transformación de una integral doble es
{\small
\begin{equation*}\label{integraldoblepolares}
\iint_R f(x,y) dydx = \iint_S f(r \cos \theta,r \sin \theta) r drd\theta.
\end{equation*}
}
\vspace{-7.9mm}


Utilizando el sistema de coordenadas cilíndricas $x=r\cos \theta ,$ 
$y=r\sin \theta ,$ $z=z;$ o el sistema de coordenadas esf\'{e}ricas,
$x=\rho \cos \theta \sin \phi ,$ $y=\rho \sin \theta \sin \phi ,$ $z=\rho\cos \phi ,$ 
con los siguientes gráficos se ilustra la obtención del elemento diferencial de volumen en estos casos.


\vskip0.2cm
\noindent
\begin{minipage}{\textwidth}
	\centering
	\captionof{figure}{Vectores ${\bf F}(x,y,z),$ ${\bf T}$ y ${\bf n}$}
	\begin{minipage}{0.31\linewidth}
		\centering
		\includegraphics[width=\linewidth]{Fig/15.pdf}
	\end{minipage}
	
	\vspace{0.3em}
	\parbox{0.95\linewidth}{\centering\fontsize{8pt}{10pt}\selectfont\textbf{Fuente:} elaboración propia.}
\end{minipage}
\vskip0.2cm


El jacobiano de la transformación a coordenadas cilíndricas se obtiene a partir de la relación
$$\left\vert \begin{array}{ccc}
\cos \theta  & \sin \theta  & 0 \\ 
-r\sin \theta  & r\cos \theta  & 0 \\ 
0 & 0 & 1\end{array}\right\vert = r\cos ^{2}\theta +r\sin ^{2}\theta = r,$$
y el elemento diferencial de Volumen es $dV= r \, dr d\theta dz,$ 
por lo tanto, la fórmula \ref{math1} se escribe como
\index{Integrales triples!en coordenadas cilíndricas} 
$$\iiint_{E}f(x,y,z)dV=\iiint_{E^{\ast }}f(r\cos \theta ,r\sin \theta,z)r\, dzdrd\theta .$$


\vspace{-6mm}
\vskip0.2cm
\noindent
\begin{minipage}{\textwidth}
	\centering
	\vskip0.2cm
	\captionof{figure}{Elemento diferencial de volumen en coordenadas esféricas}
	\label{fig:volumen-esfericas}
	\begin{minipage}{0.31\linewidth}
		\centering
		\includegraphics[width=\linewidth]{Fig/16.pdf}
	\end{minipage}%
	\begin{minipage}{0.31\linewidth}
		\centering
		\includegraphics[width=\linewidth]{Fig/17.pdf}
		\label{fig:esfericas}
	\end{minipage}
	
	\vspace{0.3em}
	\parbox{0.95\linewidth}{\centering\fontsize{8pt}{10pt}\selectfont\textbf{Fuente:} elaboración propia.}
\end{minipage}
\vskip0.2cm


El jacobiano en el caso de las coordenadas esf\'{e}ricas corresponde a 
$$\left\vert \begin{array}{ccc}
\cos \theta \sin \phi  & \sin \theta \sin \phi  & \cos \phi  \\ 
-\rho \sin \theta \sin \phi  & \rho \cos \theta \sin \phi  & 0 \\ 
\rho \cos \theta \cos \phi  & \rho \sin \theta \cos \phi  & -\rho \sin \phi 
\end{array}\right\vert = -\rho ^{2}\sin \phi,$$
que determina el elemento diferencial de volumen $dV= \rho ^2 \sin \phi d\rho d\theta d\phi$
y permite obtener la siguiente igualdad:
\begin{multline}
\iiint_{E}f(x,y,z)dV \\ =\iiint_{E^{\ast }}f(\rho \cos \theta \sin \phi ,\rho
\sin \theta \sin \phi ,\rho \cos \phi )\rho ^{2}\sin \phi\, d\rho d\theta d\phi 
\end{multline}
En el apéndice {B} el lector puede consultar las generalidades de otros sistemas de coordenadas.
\index{Integrales triples!en coordenadas Esféricas}

\begin{ejemplo}
Si $E$ es el sólido acotado por los paraboloides $z=2-x^2-y^2$ y $z=x^2+y^2.$
Evaluar la integral $\displaystyle \iiint_E y^2 dV.$ La intersección de los paraboloides está dada por $x^2+y^2=1,$ que corresponde a un círculo unitario 
en el plano $xy;$ con estos elementos se plantea y evalúa esta integral:

\vspace{-4mm}

\begin{multline*}
\int_{-1}^{1}\int_{-\sqrt{1-x^{2}}}^{\sqrt{1-x^{2}}}\int_{x^{2}+y^{2}}^{2-x^{2}-y^{2}}y^{2}dzdydx = \int_{-1}^{1}\int_{-\sqrt{1-x^{2}}}^{\sqrt{1-x^{2}}}\bigg[ y^{3}\bigg]_{z=x^{2}+y^{2}}^{z=2-x^{2}-y^{2}} dydx \\
=\int_{-1}^{1}\int_{-\sqrt{1-x^{2}}}^{\sqrt{1-x^{2}}}\left( 2y^{2}\left(1-x^{2}-y^{2}\right) \right) dydx
= \int_{-1}^{1}\frac{8}{15}\left( 1-x^{2}\right) ^{\frac{5}{2}}dx= \frac{\pi}{6}.
\end{multline*}

Aplicando un cambio a coordenadas cilíndricas se obtiene
\begin{multline*}
\int_{0}^{2\pi }\int_{0}^{1}\int_{r^{2}}^{2-r^{2}}r\left( r\sin \theta\right) ^{2}dzdrd\theta \\
= \int_{0}^{2\pi}\int_{0}^{1}2r^{3}\left( 1-r^{2}\right) \sin ^{2}\theta drd\theta 
= \int_{0}^{2\pi }\frac{1}{6}\sin ^{2}\theta d\theta = \frac{\pi}{6} 
\end{multline*}
cuyos valores son iguales, como se había establecido teóricamente.
\end{ejemplo}
\section{Divergencia y rotacional} \label{divergenciayrotacional}
Para un campo vectorial 
${\bf F}(x,y,z)= \langle P(x,y,z),Q(x,y,z),R(x,y,z) \ra,$ definido en $\mathbb{R}^3,$ cuyas funciones componentes poseen derivadas parciales continuas,
se define el rotacional y la divergencia de ${\bf F}$ que se denotan por $\rot{F}$ y ${\rm div\,}{\bf F},$ y 
están dados, respectivamente, por 
\index{Divergencia} \index{Rotacional} 
$$\rot{F}(x,y,z)= \langle  \frac{\partial R}{\partial y}-\frac{\partial Q}{\partial z}, \frac{\partial P}{\partial z}-\frac{\partial R}{\partial x}, \frac{\partial Q}{\partial x}-\frac{\partial P}{\partial y} \rangle $$
y $${\rm div\,}{\bf F}(x,y,z)= \frac{\partial P}{\partial x}+\frac{\partial Q}{\partial y}+\frac{\partial R}{\partial z}.$$
Nótese que $\rot{F}(x,y,z)$ es un campo vectorial y ${\rm div\,}{\bf F}(x,y,z)$ es un campo escalar. 
Existe una manera de recordar brevemente dichos cálculos, en términos del operador $\nabla$, que se presentó en la sección 3.3.
\begin{eqnarray*}
{\rm rot}{\bf F} &=& \nabla \times{\bf F}  = \left\vert \begin{array}{ccc}
\vi & \vj &  \vk \\  \frac{\partial }{\partial x}  & \frac{\partial }{\partial y}
  & \frac{\partial }{\partial y}  \\ P & Q & R \end{array}\right\vert\\
{\rm div}{\bf F} &=& \nabla \cdot{\bf F}.
\end{eqnarray*}

\begin{ejemplo}
Para el campo vectorial ${\bf F}(x,y,z) =\left\langle xye^{z},z\cos (xy) ,\sin (xy) \right\rangle,$
es claro que 
$${\rm div}{\bf F}=\frac{\partial }{\partial x} \left( xye^{z}\right) +\frac{\partial }{\partial y} \left( z\cos (xy) \right) +\frac{\partial }{\partial z}(\sin(xy))= ye^{z}-xz\sin xy,$$ y también
$${\rm rot}{\bf F} = \left\langle (x-1) \cos(xy),y(xe^{z}-\cos(xy)) ,-xe^{z}-yz\sin xy \right\rangle.$$
\end{ejemplo}
La integración de estos tipos especiales de funciones se considerarán en el capítulo \ref{chap:stokes},
así como algunas de sus aplicaciones.
\newpage
\chapter{Elementos básicos en {\tt Mathematica}} \chapter{Elementos básicos en {\tt Mathematica}}

En este capítulo presentamos algunas de las funciones y comandos
básicos que se utilizan en el programa \verb"Wolfram Mathematica", que denominaremos simplemente
como \verb"Mathematica". Es de anotar que toda instrucción hecha en este programa se debe \emph{ejecutar} para obtener el producto señalado. En la versión \verb"13", bajo el sistema operativo \verb"Windows 10" esto se consigue con \verb"shift+enter" o \verb"enter" con el teclado numérico.

En general la sintaxis de las funciones inician con letra mayúscula
y en algunos casos llevan otras letras intermedias también en mayúscula. Los argumentos de cada función se insertan entre
corchetes \verb"[...]", los paréntesis \verb"(...)" se usan para agrupar y las llaves \verb"{...}" se utilizan para insertar \emph{listas}.

\section{Funciones básicas}
Las operaciones básicas: adición, sustracción, producto y cociente, entre dos cantidades $x$ e $y$ se denotan respectivamente como \verb"x+y", \verb"x-y", \verb"x*y" y \verb"x/y". 

\begin{tcolorbox}[breakable,enhanced,title=Advertencia]
Es importante aclarar que en general \verb"xy" no es igual a \verb"x*y".
\end{tcolorbox}
Algunas de las constante más utilizadas son: \verb"Pi", $\pi$; \verb"E", $e \approx 2.71828$;
\verb"I", la unidad imaginaria $i$, $i^{2}= -1;$ \verb"Infinity", $\infty$.

 Las funciones trigonométricas\footnote{Con argumento en radianes} se denotan como
{\small
\verb"Sin[x]," \verb"Cos[x]," \verb"Tan[x]," \verb"Cot[x]," 
\verb"Sec[x]," \verb"Csc[x]"
}
y sus respectivas funciones inversas
{\small
\verb"ArcSin[x]," \verb"ArcCos[x]," \verb"ArcTan[x]", \verb"ArcCot[x]," \verb"ArcSec[x]," \verb"ArcCsc[x]".
}

Para un punto $(x,y) \in \mathbb{R}^2$ se puede usar  \verb"ArcTan[x,y]", esta función posee como dominio a $[-\pi,\pi].$ 
Análogamente se escriben las funciones hiperbólicas y sus inversas \verb"Sinh[y]",  \verb"Cosh[y]", \verb"Tanh[y]", \verb"ArcSinh[y]", \verb"ArcCosh[y]," \verb"ArcTanh[y]", etc.

Existen tres tipos de igualdad, para asignación: \verb"x=2";  
prueba lógica: \verb"2+2==4", usada en ecuaciones y como una manera de probar dos expresiones son \emph{lo mismo}, 
\verb"x*x === x^2".
Otras funciones elementales son las siguientes.
\begin{tcolorbox}[breakable,enhanced,title=Algunas funciones básicas en Wolfram]
\begin{tabular}{rl}
\verb"x^y"& Potencia $y$ de $x$, $x^y$\\
\verb"Sqrt[x]"& Raíz cuadrada de $x$\\
\verb"Exp[x]"& Exponencial de $x$, $e^{x}$\\
\verb"Log[x]"& Logaritmo natural de $x$\\
\verb"Log[b,x]"& Logaritmo en base $b$ de $x$\\
\verb"n!"& Factorial de $n$\\
\verb"Abs[x]"& Valor absoluto de $x$\\
\verb"Max[x,y,...]", \verb"Min[x,y,...]"& Máximo o mínimo de $x, y,\cdots$\\
\verb"FactorInteger[n]"& Factores primos de $n$
\end{tabular}
\end{tcolorbox}

Para factorizar, simplificar, expandir, descomponer en fracciones parciales, operar fracciones o borrar una \verb"EXPRESIÓN"
se usan respectivamente, las siguientes funciones. \index{Together} \index{Simplify} \index{Apart}\index{Expand}
\begin{tcolorbox}[breakable, enhanced, title=Funciones básicas]
\begin{multicols}{2} % Activa las dos columnas
\begin{Verbatim}[formatcom=\letraverbatim]
Factor[EXPRESION]
Simplify[EXPRESION]
Expand[EXPRESION]
Apart[FRACCION]
Together[FRACCIONES]
Clear[VARIABLES]
\end{Verbatim}
\end{multicols}
\end{tcolorbox}

Por ejemplo, el usuario puede ejecutar las siguientes instrucciones:
\begin{tcolorbox}[breakable, enhanced]
\begin{multicols}{2}
\begin{Verbatim}[formatcom=\letraverbatim] 
Factor[x^16-y^16]
Apart[4*y^3/(y^4-16)]
Together[1/x+y/(y+1)]
Simplify[(x^9-1)/(x^3-1)]
Expand[(x^9-y^9)*(x^3-y^3)]
\end{Verbatim}
\end{multicols}
\end{tcolorbox}

Para extraer el numerador o el denominador de una expresión racional se pueden utilizar
las funciones \verb"Numerator" o \verb"Denominator". 
Utilizando las funciones \verb"ExpandDenominator" o \verb"ExpandNumerator"
se puede hacer la expansión en una expresión de los términos respectivos.
Las siguientes sentencias ilustran el uso de estos operadores.
\index{Numerator} \index{Denominator} \index{ExpandDenominator} \index{ExpandNumerator}
\begin{tcolorbox}[breakable,enhanced,title=Usos de \verb"ExpandDenominator" y \verb"ExpandNumerator"]
\begin{Verbatim}[formatcom=\letraverbatim]
Numerator[Together[x/y+y/x]]
Denominator[Together[x/y+y/x]]
ExpandDenominator[Together[y/(x+y)+x/(x-y)]]
ExpandNumerator[Together[y/(x+y)+x/(x-y)]]
\end{Verbatim}
\end{tcolorbox}

\index{TrigExpand} \index{TrigFactor} \index{TrigReduce}
Si una \verb"EXPRESIÓN" es una combinación de funciones trigonométricas, la instrucción
\verb"TrigExpand[EXPRESION]" la reescribe en términos de las funciones $\sin(\cdot)$ y $\cos(\cdot)$. 
La sentencia \verb"TrigFactor[EXPRESION]" factoriza usando funciones trigonométricas y
\verb"TrigReduce[EXPRESION]" convierte producto de funciones trigonométricas en suma. 
Por ejemplo:

\begin{tcolorbox}[breakable,enhanced]
\begin{multicols}{2}
\begin{Verbatim}[formatcom=\letraverbatim]
TrigExpand[Sin[5*x]-Cos[6*x]]
TrigReduce[Sin[2*x]*Cos[3*x]]
TrigFactor[Sin[5*x]-Cos[5*x]]
TrigFactor[Sin[x-y]-Cos[x+y]]
TrigExpand[Sin[x]^2*Cos[2*x]^3]
TrigReduce[Sin[3x]^3+Cos[2*x]^3]
\end{Verbatim}
\end{multicols}
\end{tcolorbox}

\subsection*{Tablas}
Para generar una lista con $n$ términos de la expresión $f(k)$ se usa la expresión \verb"Table[f[k],n]", además 
con \verb"Table[f[k],{k,a,b}]" se consigue que la lista satisfaga que $a \leq k \leq b.$ \index{Table}
Ilustremos esto con dos ejemplos sencillos.

\begin{tcolorbox}[breakable,enhanced]
\verb"Table[k^2,{k,2,20}"\\
\verb"Table[Prime[k],{k,2,20}"
\end{tcolorbox}

\subsection*{Vectores y funciones vectoriales}
El vector ${\bf u}= \langle a,b,c \rangle$ se define en \verb"Mathematica" como una lista \verb"u={a,b,c}".
El producto punto, el producto cruz y la norma de los vectores ${\bf u}$ y ${\bf v}$ se determina como sigue
\begin{tcolorbox}[breakable,enhanced,title=Vectores]
\verb"Dot[u,v]" \hskip1.75cm \verb"Cross[u,v]" \hskip1.75cm \verb"Norm[u]"
\end{tcolorbox}

\begin{tcolorbox}[breakable,enhanced,title=Funciones vectoriales]
Una función vectorial ${\bf r}(t)= \la x(t), y(t), z(t)\ra$ se define como un vector
\verb"r={x[t],y[t],z[t]}". 
\end{tcolorbox}
La instrucción  \verb"r={Cos[t],Sin[t],Cos[2*t]}" define una función vectorial que se denomina \verb"r".
Para extraer cada una de las tres componentes, 
se debe usar, según el caso, las siguientes instrucciones

%\verb"r[[1]]" \\ \verb"r[[2]]" \\ \verb"r[[3]]"
\begin{center}
\verb"r[[1]]" \hspace{3cm} \verb"r[[2]]" \hspace{3cm} \verb"r[[3]]"
\end{center}


Se pueden hacer algunos procedimientos asumiendo ciertas condiciones. Las funciones
\verb"Assuming" y \verb"Assumptions", permiten este tratamiento.
\index{Assumptions} \index{Assuming}
En los siguientes ejemplos se evidencia la diferencia de usar estas posibilidades.

\begin{tcolorbox}[breakable,enhanced,title=Usos de \verb"Assuming" y \verb"Assumptions"]
\begin{Verbatim}[formatcom=\letraverbatim]
Norm[{Cos[t],Sin[t],t}]
Norm[{Cosh[t],Sinh[t],t}]
Assuming[t>0,Simplify[Norm[{Cos[t],Sin[t],t}]]]
Assuming[t>0,Simplify[Norm[{Cosh[t],Sinh[t],t}]]]
Simplify[Norm[{Cosh[t],Sinh[t],t}],Assumptions->t>0]
\end{Verbatim}
\end{tcolorbox}

El símbolo $\%$ recupera el cálculo inmediatamente anterior, por ejemplo: 

\begin{tcolorbox}[breakable,enhanced]
\begin{multicols}{2}
\begin{Verbatim}[formatcom=\letraverbatim]
u={Exp[t],Exp[t],t}
v={-Exp[t],-Exp[t],t}
Norm[u+v]
Simplify[%,Assumptions->t>0]
\end{Verbatim}
\end{multicols}
\end{tcolorbox}

\subsection*{Sumas}
El comando \verb"Sum[Expresión,{k,a,b}]" obtiene la suma de una \verb"Expresión"
con variable de sumación $k$ la cual inicia en $a$ y termina en $b.$ 
Ejemplo de lo anterior son las siguientes sumas.
$$\sum_{k=1}^{10} k^2 \hskip1.5cm
\sum_{k=1}^n k^3 \hskip1.5cm
\sum_{k=1}^\infty \frac{1}{k^2} \hskip1.5cm
\sum_{r=0}^\infty \frac{2^r}{r!}$$
Para obtener su valor exacto se usan las siguientes instrucciones.

\begin{tcolorbox}[breakable, enhanced, title=Funciones \verb"Mathematica"]
\begin{multicols}{2} % Activa las dos columnas
\begin{Verbatim}[formatcom=\letraverbatim]
Sum[k^3,{k,1,n}]
Sum[k^2,{k,1,10}]
Sum[1/k^2,{k,1,Infinity}]
Sum[2^r/r!,{r,0,Infinity}]
\end{Verbatim}
\end{multicols}
\end{tcolorbox}
 
 
\subsection*{Límites}
Utilizando la sentencia \verb"Limit[F[x],x->a]" se evalúa en \verb"Mathematica" el $\lim_{x\to a} F(x)$.
También es posible calcular límites en varias variables como se ilustra enseguida.

{\small $$\lim_{x\to0}\frac{\sin(3x)}{\sin(2x)}\hskip1.5cm
\lim_{(x,y)\to (0,0)}\frac{x^2y}{x^2+y^2} \hskip1.5cm
\lim_{(x,y)\to(0,0)}\frac{\sin(x^2+y^2)}{x^2+y^2}$$}

Los límites anteriores se evalúan con las siguientes instrucciones.
\begin{tcolorbox}[breakable,enhanced]
\begin{Verbatim}[formatcom=\letraverbatim]
Limit[Sin[3*x]/Sin[2x],x->0]
Limit[x^2*y/(x^2+y^2),{x,y}->{0,0}]
Limit[Sin[x^2+y^2]/(x^2+y^2),{x,y}->{0,0}]
\end{Verbatim}
\end{tcolorbox}

\subsection*{Precaución}
Aunque es posible evaluar límites de manera iterada, es conveniente advertir que el uso 
incorrecto de esta sintaxis puede llevar a conclusiones arroneas. 
Por eso es importante usar la indicación adecuada. Ilustraremos esta situación.
\begin{itemize}%[leftmargin=0in]
\item Correcta \verb"Limit[n/(n+m),{n,m}->{Infinity,Infinity}]"
\item Incorrecta \verb"Limit[Limit[n/(n+m),m->Infinity],n->Infinity]"
\item Incorrecta \verb"Limit[Limit[n/(n+m),n->Infinity],m->Infinity]"
\end{itemize}

\subsection*{Operadores básicos}
%\begin{tcolorbox}[breakable,enhanced,title=Operadores Básicos en \verb"Mathematica"]
\index{Solve} \index{NSolve}
\begin{itemize}
\item Para resolver la ecuación $f(x)=0,$ en donde $x$ es la incognita, se usa la
instrucción: \verb"Solve[f[x]==0,x]". La solución numérica de ecuaciones se resuelve utilizando \verb"NSolve".
{\bf No olvidar $"=="$}.
\item La derivada de $f(x)$ con respecto a $x$, se obtiene con \verb"D[f[x],x]", 
la segunda derivada de $f(x)$ se consigue con \verb"D[f[x],x,x]".
La $n-$ésima derivada de $f(x)$ con respecto a $x$, $\frac{\partial^n f}{\partial x^n}$, se obtiene con la instrucción
\verb"D[f[x],{x,n}]". 
\item La derivada parcial de $f(x,y,z)$ con respecto a $x$, se consigue ejecutando \verb"D[f[x,y,z],x]", mientras que 
la diferencial total de $x^y,$ (suponiendo que $y$ depende de $x$) se obtiene mediante la expresión 
\verb"Dt[x^y,x]". \index{Diferencial total}
\item La integral indefinida de la función $f(t)$ con respecto a $t$, 
se evalúa mediante la instrucción: \hskip0.1cm \verb"Integrate[f(t),t]".
Con la instrucción \\\verb"Integrate[f(t),{t,a,b}]" se evalúa la integral definida $\ds \int_a^bf(t)dt$.
  \index{Integrate}
\end{itemize} 

%\end{tcolorbox} 
Ejemplo de lo anterior son las siguientes instrucciones:
\begin{tcolorbox}[breakable,enhanced]
\begin{multicols}{2}
\begin{Verbatim}[formatcom=\letraverbatim]
Solve[x^4==1,x]
Solve[x^4y^2==1,y,Reals]
D[ArcSin[x],x]
D[Sqrt[1/(1+r)], r]
D[Sqrt[x^3+x-2], x,x]
Integrate[1/(1+r),{r,0,1}]
\end{Verbatim}
%\end{itemize}
\end{multicols}
\end{tcolorbox}

\subsection*{Animación de objetos}
Con los siguientes comandos se puede dar \emph{animación} a un \verb"TERMINO", la variable es $k$, con $n \leq k \leq m $
e incrementos de $r$ unidades.
\index{Manipulate}\index{Animate}
\begin{tcolorbox}[breakable,enhanced,title=Las funciones {\tt Manipulate} y {\tt Animate} ]

\begin{Verbatim}
Manipulate[TERMINO,{k,n,m,r}]
Animate[TERMINO,{k,n,m,r}]
\end{Verbatim}

\end{tcolorbox}
La $r$ se puede omitir. % en el caso de la función \verb"Animate". 
El lector puede ejecutar las siguientes instrucciones
para evidenciar un poco mejor, el uso de las funciones mencionadas
\begin{tcolorbox}[breakable,enhanced,title=Ejemplos con {\tt Manipulate} y {\tt Animate}]
\begin{Verbatim}
Manipulate[Expand[(a+b)^k],{k,1,20,1}]
Animate[Expand[(a+b)^k],{k,1,20,1}]
Manipulate[Integrate[x^k,x],{k,-5,5,1/2}]
Manipulate[Integrate[(Sin[x])^k,x],{k,-2,10,1}]
\end{Verbatim}
\end{tcolorbox}

\subsection*{Transformación a coordenadas polares}
La función \verb"ToPolarCoordinates" transforma el punto $(a,b)$ escrito en coordenadas cartesianas 
al punto $(r, t)$ escrito en coordenadas polares. 
Con la función \verb"FromPolarCoordinates" se realiza la transformación inversa.
\index{ToPolarCoordinates}  \index{FromPolarCoordinates}

\begin{tcolorbox}[breakable,enhanced]
\verb"ToPolarCoordinates[{a,b}]"\\
\verb"FromPolarCoordinates[{r,t}]"
\end{tcolorbox}
 
El argumento $t$ satisface que $-\pi \leq t \leq \pi$. Por ejemplo:
\begin{tcolorbox}[breakable,enhanced]
\verb"ToPolarCoordinates[{-1,-1}]"\\
\verb"FromPolarCoordinates[{2,-3*Pi/4}]"
\end{tcolorbox}

\subsection*{Definición de funciones}
Es posible definir una función y luego usarla para cálculos posteriores,
esta puede depender de una o varias variables, la sintaxis básica es:

\begin{tcolorbox}[breakable,enhanced,title=Definición de funciones]
\verb"f[x_] := Expresión" para funciones de una variable $x.$ \\
\verb"f[x_,y_] := Expresión" para funciones de dos variables $x,y.$
\end{tcolorbox}
% Al no definir la función en un conjunto el valor por defecto es $0.$ 
Con la siguiente sentencia se define la función vectorial ${\bf r}(t,u)$, \\
\verb"r[t_,u_]:={(1-u)*Cos[2t]+u*Sin[3t],u*Cos[3t],t(1-u)}" \\
para evalúarla en $t=\pi/3$ y $u=1/2$, se utiliza \verb"r[Pi/3,1/2]".

Se pueden definir funciones a trozos, con $k$ condiciones y $x \in [a,b]$, 
\begin{tcolorbox}[breakable,enhanced,title=Definición de funciones a trozos]
\verb"Piecewise[{{Expresión 1,condición 1},{Expresión 2,"
\verb"condición 2},...,{Expresión k,condición k}},{x,a,b}]"
\end{tcolorbox}  \index{Piecewise}

\enlargethispage{\baselineskip}

Ahora se muestran algunos ejemplos:
\begin{itemize}[leftmargin=0cm,itemindent=.5cm,labelwidth=\itemindent,labelsep=0cm,align=left]
\item Para definir la función $f(x,y)=x^2-3y^4$ y evaluarla en el punto $(2,3)$ se ejecutan 
las instrucciones \verb"f[x_,y_]:=x^2-3*y^4", \verb"f[2,3]".
\item La función 
$$ f(x,y) = \left \{ \begin{array}{cl} \frac{xy}{x^2-y^2} & \text{ si } x>y \\[0.2cm] \sin(x+y) & \text{en caso contrario} \end{array} \right.$$
se define ejecutando la siguiente instrucción:\\
\verb"f[x_,y_]:=Piecewise[{{x*y/(x^2-y^2),x>y}},Sin[x+y]]"
\item Con la instrucción \verb"f[x_]:=Piecewise[{{4-x^2,x<=1}," \\
\verb"{2-x^3,1<x<2},{2*Sin[x],x>=2}}]", se define la función 
$$ f(x,y) = \left \{
 \begin{array}{cl} 4-x^2 & \text{ si } x\leq 1 \\[0.2cm]
2-x^3 & \text{ si } x\in(1,2)\\[0.2cm]
2 \sin(x) & \text{ si } x\geq 2
\end{array} \right.$$
\end{itemize}

\section{Gráficos}
Presentamos las principales posibilidades que ofrece \verb"Mathematica" 
para obtener la representación gráfica de objetos matemáticos.

\subsection*{Gráficas de funciones en el plano}
Se ofrece una manera de obtener la representación gráfica en el plano de objetos o funciones 
escritas en forma explícita, implícita o paramétrica.

\begin{itemize}[leftmargin=0cm,itemindent=.5cm,labelwidth=\itemindent,labelsep=0cm,align=left]
\item \verb"Plot[f[x],{x,a,b}]" dibuja la función $y=f(x)$ con $x \in[a,b].$ \index{Plot}
\item \verb"PolarPlot[r[t],{t,a,b}]" grafica la función $r(t)$ escrita en coordenadas polares, con $t \in[a,b].$
\item \verb"ContourPlot[f[x,y]==0,{x,a,b},{y,c,d}]",\index{ContourPlot}
dibuja la ecuación de dos variables $f(x,y)=0$, con $x \in[a,b]$ y $y \in[c,d],$
se debe incluir el símbolo de igualdad y no olvidar $"=="$.
\item \verb"ParametricPlot[{x[t],y[t]},{t,a,b}]", dibuja la curva parametri- \\ zada $\ve{r}(t)=\left\langle x(t),y(t)\right\rangle ,\,\, t\in \lbrack a,b].$
\index{ParametricPlot}
\end{itemize}
% \end{tcolorbox}
Si se desean incluir varias expresiones en las funciones anteriores,
se puede hacer una lista con estas y encerrarlas en llaves (\verb"{...}").
Las siguientes instrucciones son ejemplos de las anteriores funciones:

\begin{tcolorbox}[breakable,enhanced,title=Algunos gráficos en el plano]
\begin{Verbatim}
Plot[Exp[-10*t]*(2+20t)+1,{t,0,1}]
PolarPlot[1-Cos[t],{t,0,2*Pi}] 
PolarPlot[3/(6+Cos[10*u]),{u,0,2*Pi}]
ContourPlot[{y^2+(1-x^2)^2==4,y^2+x^2==4},{x,-2,2},
{y,-2,2}] 
\end{Verbatim}
\begin{Verbatim}
ParametricPlot[{Cos[t],Sin[2*t]},{t,0,2Pi}]
Manipulate[PolarPlot[t*Sin[k*t],{t,0,2Pi}],
{k,1,5,0.5}]
ParametricPlot[{{2*Cos[t],2*Sin[t]},{Cos[t],Sin[t]},
{Cos[t]*(1-Cos[t]),Sin[t]*(1-Cos[t])}},{t,0,2*Pi}]
\end{Verbatim}
\end{tcolorbox}

\subsection*{Gráficos de funciones en el espacio}
De manera análoga a las funciones descritas anteriormente para
conseguir gráficos en el plano, se puede agregar la extensión \verb"3D" y unas opciones adicionales
con el fin de obtener gráficos en el espacio.
Presentamos unas de las funciones más utilizadas en  \verb"Mathematica".
\\

\begin{itemize}[leftmargin=*]
\item \verb"Plot3D[f[x,y],{x,a,b},{y,c,d}]", permite obtener el gráfico de la función 
$z=f(x,y)$, para $(x,y) \in [a,b] \times [c,d]$. 
\\


\item \texttt{Plot3D[\allowbreak\{f[x,y],\allowbreak g[x,y],\allowbreak ..., h[x,y]\},\allowbreak \{x,a,b\},\allowbreak \{y,c,d\}]} \\
permite obtener el gráfico de varias funciones.
\\
%Mientras que con \verb"Plot3D[{f[x,y],g[x,y],...,h[x,y]},{x,a,b},{y,c,d}]", 
%se obtiene el gráfico de varias funciones. \index{Plot3D}
\item \Verb!ParametricPlot3D[{x[t],y[t],z[t]},{t,a,b}]!,
dibuja la curva parametrizada $\ve{r}(t)=\langle x(t),y(t),z(t) \rangle $, 
con $t\in[a,b].$  \index{ParametricPlot3D}
\\

\item \verb"ParametricPlot3D[{x[u,v],y[u,v],z[u,v]},{u,a,b},{v,c,d}]",
se utiliza para graficar superficies parametrizadas por medio de una función
$\ve{r}(u,v)=\langle x(u,v),y(u,v),z(u,v) \rangle $, con
$u\in[a,b],$ $v\in[c,d].$
\\
\item \verb"ContourPlot3D[f[x,y,z]==0,{x,a,b},{y,c,d},{z,p,q}]", se usa
para graficar conjuntos de nivel (contornos).
Este comando presenta el gráfico de una ecuación de tres variables, con $x\in[a,b],$
$y\in[c,d]$ y $z\in[p,q].$  \index{ContourPlot3D}
\\
\end{itemize}

Los gráficos se pueden personalizar en sus atributos con el uso de algunas opciones, 
presentamos las más utilizadas:

\vspace{-2pt}

\begin{tcolorbox}[breakable,enhanced,title=Opciones de los gráficos]
\begin{tabular}{rl}
\verb"Axes" & Muestra los ejes, \verb"True/False"\\
\verb"AxesLabel->{u,v,w}" & Etiqueta los ejes como $u,$ $v$ y $w$ \\
\verb"AxesOrigin->{a,b,c}" & Muestra los ejes con centro en $(a,b,c)$ \\
\verb"Blend[{color1,...}]"& Permite combinar varios colores \\
\verb"Boxed" & Muestra gráfico en caja, \verb"True/False"\\
\verb"BoxRatios->{a,b,c}" & Proporciones de la caja en un gráfico\\
\verb"ColorFunction" & Determina el color de la función\\
\verb"ContourStyle" & Opciones en \verb"ContourPlot3D" \\
\verb"Dashed" & Presenta líneas discontinuas\\
\verb"Mesh->{n,m}" & Crea la malla, con $n$ y $m$ líneas \\
\verb"Mesh" & Activa la malla, \verb"True/False" \\ 
\verb"MeshStyle" & Estilo de la malla (color)\\
\verb"Opacity[a]" &  Opacidad, se usa en plotStyle, $0 \leq a \leq 1$\\
\verb"PlotStyle" &  Opciones del gráfico, color, líneas, etc.\\
\verb"PlotRange" &  Rango de valores en el gráfico\\
\verb"PlotPoints" & Número de puntos sobre el dibujo\\
\verb"RGBColor[a,b,c]" & Crea un color con proporciones $a$, $b$, $c$\\
\verb"Thick" & Crea una línea más gruesa \\
\verb"Thickness" & Grosor de la línea \\
\verb"Ticks" &  Indica las marcas deseadas en la caja\\
\verb"Tube[a]" & Configura un tubo de radio $a$. \\
\end{tabular}
\end{tcolorbox}

El lector puede además de ejecutar las instrucciones, modificarlas.
{\small
\begin{itemize}
\item \verb"Plot3D[{5-y^2,x^2+y^2},{x,-4,4},{y,-4,4},PlotPoints->60,"\\
\verb"PlotStyle->Opacity[0.6]]"
\item \verb"ContourPlot3D[{z==6-(x-1)^2-y^2,4*x^2+y^2==36,"\\
\verb"x^2+4*y^2==4},{x,-6,10},{y,-7,7},{z,-35,6},"\\
\verb"ContourStyle->{Red,Yellow,Green}]"
\item \verb"ParametricPlot3D[{Cos[u]*Cosh[v],Sin[u]*Cosh[v],Sinh[v]},"\\
\verb"{u,0,2*Pi},{v,-2,2},PlotStyle->Orange,MeshStyle->Gray]"
\item \verb"ContourPlot3D[z==3*x-x^3-2*y^2+y^4,{x,-3,3},{y,-2,3},"\\
\verb"{z,-3,3},MeshStyle->Gray]"
\item \begin{Verbatim} 
ContourPlot3D[z==3*x-x^3-2*y^2+y^4,{x,-3,3},{y,-2,3},
{z,-3,3},ColorFunction->Function[{x,y,z},
ColorData["Rainbow"][z]],MeshStyle->Gray]
\end{Verbatim}
\end{itemize}
}

El siguiente gráfico se consigue utilizando las instrucciones señaladas.


\begin{figure}[H]
\centering
\caption{Integración en el espacio}
\includegraphics[width=0.4\textwidth]{Img/math_01a.pdf} \hspace{5mm}
\includegraphics[width=0.4\textwidth]{Img/math_01b.pdf}
\fuentepropia
\end{figure}

\begin{tcolorbox}[breakable,enhanced, title=Instrucción en \verb"Mathematica"]
\begin{Verbatim}[fontsize=\small]
ContourPlot3D[z==(y^2-x)*Exp[1-x^2-y^2],{x,-2,2},
{y,-2,2},{z,-1,2},Mesh->{6,8},Boxed->False,
MeshStyle->RGBColor[0.5,0.5,0.5],AxesOrigin->{0,0,0},
ContourStyle->Orange,Ticks->{{-1,1,2},{-1,1,2},
{-1,1}},AxesStyle->Black,AxesLabel->{u,v,w}]
\end{Verbatim}
\end{tcolorbox}


La composición de varias funciones puede hacer complicado
el tratamiento, la función \verb"Show" permite ver, en un solo gráfico, varios objetos obtenidos con 
diferentes funciones, por ejemplo: \index{Show}
\begin{tcolorbox}[breakable,enhanced,title=Uso de la función \verb"Show"]
\begin{Verbatim} 
Show[ParametricPlot3D[{Cos[t],Sin[t],1},{t,0,2Pi},
PlotStyle->Red],Graphics3D[{Blue,Arrowheads[0.05],
Arrow[Tube[{{1,0,1},{1,1,1}},0.02]]}],
AxesOrigin->{0,0,0},
PlotRange->{{-2,2},{-2,2},{-1,2}}]
\end{Verbatim}
\end{tcolorbox}
Otra posibilidad es ejecutar las instrucciones de manera separada y luego mostrarlos en un solo 
gráfico con la función \verb"Show". En el siguiente ejemplo, las primeras dos instrucciones se ejecutan, 
se nombran como \verb"A" y \verb"B", y luego con la función \verb"Show" se muestran ambas.%\\
\begin{tcolorbox}[breakable,enhanced,title=Ejecución de gráficos por separado]
\begin{Verbatim} 
A=ContourPlot3D[{x==4-y-4z+2z^2,(1-z)^2+(2-y)^2==1}
{x,-3,3},{y,-1,4},{z,-1/2,5/2},Mesh->False];
B=ParametricPlot3D[{Sin[t]+2(Cos[t])^2,2-Sin[t],
1-Cos[t]},{t,0,2Pi},
PlotStyle->{Red,Thickness->0.02}];
Show[A,B]
\end{Verbatim}
\end{tcolorbox}
Otras posibilidades muy útiles para obtener gráficos son las siguientes.

\subsection*{Gráfico de puntos, segmentos y vectores}

\begin{tcolorbox}[breakable,enhanced,title=Gráfico de puntos]
La instrucción \verb"Graphics[Point[{a,b}]]" dibuja el punto de
coordenadas $\left(a,b\right).$ Con \verb"Graphics3D[Point[{a,b,c}]]" se dibuja el 
punto de coordenadas $\left(a,b,c\right) .$
\end{tcolorbox}
El tamaño de un punto se modifica con el comando \verb"PointSize". \index{PointSize}
Para dibujar de color rojo los puntos $\left( 1,1,1\right) $ y $\left( 2,0,1\right)$
se ejecuta la instrucción \verb"Graphics3D[{Red,PointSize[0.1],Point[{{1,1,1},{2,0,1}}]}]".

El segmento que une el punto $(a,b)$ con el punto $(p,q)$ se representa utilizando la siguiente instrucción:
\begin{tcolorbox}[breakable,enhanced,title=Gráfico de segmentos]
\begin{Verbatim}
Graphics[Line[{{a,b},{p,q}}]]
\end{Verbatim}
\end{tcolorbox} \index{Line}

Con la instrucción \texttt{Graphics3D[\allowbreak Line[\allowbreak\{\{a,b,c\},\allowbreak\{p,q,r\}\}]]}, se obtiene el gráfico de un segmento tridimensional. Por ejemplo, con la instrucción \texttt{Graphics3D[\allowbreak\{\allowbreak Dashed,\allowbreak Blue,\allowbreak Line[\allowbreak\{\{1,-1,1\},\allowbreak\{2,3,-1\}\}]\}]}, se obtiene el segmento de color azul en línea discontinua que une los puntos $(1,-1,1)$ y $(2,3,-1).$

También es posible dibujar el vector que va del punto $(a,b,c)$ al punto $(p,q,r).$
\begin{tcolorbox}[breakable,enhanced,title=Gráfico de vectores en el espacio]
\begin{Verbatim}
Graphics3D[Arrow[{{a,b,c},{p,q,r}}]]
\end{Verbatim}
\end{tcolorbox}
Si se desea agregar determinado grosor $u$ al vector, la instrucción que se debe ejecutar es 
\verb"Graphics3D[Arrow[{{a,b,c},{p,q,r}},u]]". \index{Graphics3D}
El vector que une el punto $(1,1,-1)$ con el punto $(2,2,0),$ y algunos atributos personalizados, 
se dibuja en  utilizando las siguientes instrucciones:
\begin{tcolorbox}[breakable,enhanced]
\begin{Verbatim}[formatcom=\letraverbatim]
Graphics3D[{Black,Arrowheads[0.1],
Arrow[Tube[{{1,1,-1},{2,2,0}},0.01]]}]
\end{Verbatim}
\end{tcolorbox}

\begin{figure}[H]
	\begin{center}
		\caption{El vector $\langle 1,1,1\rangle$ se dibuja en el punto $(1,1,-1)$.}
		\includegraphics[height=4.5cm]{Img/T_67a.pdf} \hskip0.2cm
		\includegraphics[height=4.5cm]{Img/T_67b.pdf}
\fuentepropia
	\end{center}  
\end{figure}

\subsection*{Regiones en el plano y el espacio}
La función \verb"RegionPlot" muestra el gráfico de una región dimensional, la sintaxis puede incluir desigualdades 
y los operadores lógicos: \emph{y} (\verb"&&"), \emph{o} (\verb"||"). En los siguientes ejemplos se incorpora el uso 
de las opciones \verb"BoundaryStyle" y \verb"Thickness" que controla el grosor y color de la frontera.
\index{Operadores lógicos} \index{BoundaryStyle} \index{Thickness} 


\begin{figure}[H]
\centering
\caption{\small Uso de la función {\tt RegionPlot}}
\includegraphics[width=0.4\textwidth]{Img/T_49a.pdf} \hspace{5mm}
\includegraphics[width=0.4\textwidth]{Img/T_49b.pdf}

\includegraphics[width=0.4\textwidth]{Img/T_49c.pdf} \hspace{5mm}
\includegraphics[width=0.4\textwidth]{Img/T_49d.pdf}
\fuentepropia
\end{figure}

La versión tridimensional es \verb"RegionPlot3D".



\begin{tcolorbox}[breakable,enhanced]
\begin{Verbatim}[formatcom=\letraverbatim] 
RegionPlot[1<=x^2+2*y^2<=9&&y>=-x^2,{x,-3,3},
{y,-2,5/2},PlotPoints->50,BoundaryStyle->{Black,
Thickness->0.01},PlotStyle->Gray,Frame->False]
RegionPlot[(x+1)^2+2*y^2<3||(x-1)^2+y^2<3,{x,-3,3},
{y,-2,2},BoundaryStyle->{Black,Thickness->0.01},
PlotStyle->Gray,Frame->False]
RegionPlot[{(x+1)^2+2*y^2<3,(x-1)^2+y^2<3},{x,-3,3},
{y,2,2},PlotStyle->Gray,BoundaryStyle->{Black,
Thickness->0.01},Frame->False]
RegionPlot[x^2+y^2<9&&Abs[x]+Abs[y]>2,{x,-3,3},
{y,-3,3},Frame->False,PlotStyle->Gray,
BoundaryStyle->{Black,Thickness->0.01}]
\end{Verbatim}
\end{tcolorbox}



Para representar un sólido se utiliza la función \verb"RegionPlot3D". \index{RegionPlot3D}
Este comando dibuja la región a partir de ciertas \verb"DESIGUALADES", con $x\in[a,b]$, $y\in[c,d]$, $z\in[u,v].$ 
Las inecuaciones se concatenan con los símbolos \verb"&&".
La sintaxis básica es:



\begin{tcolorbox}[breakable,enhanced,title=Gráfico de una región sólida]
\begin{Verbatim}[formatcom=\letraverbatim]
RegionPlot3D[DESIGUALDADES,{x,a,b},{y,c,d},{z,u,v}]
\end{Verbatim}
\end{tcolorbox}


Presentamos ahora, algunas situaciones que ilustran el uso de esta función.

(1) El sólido acotado por las figuras de $z=0,$ $z=1-y,$ $x^2=y,$ $y=1.$


\begin{figure}[H]
\centering
\caption{Cono truncado oblicuo determinado por ${\bf r}(u,t)$}
\includegraphics[width=0.4\textwidth]{Img/math_06c.pdf} \hspace{5mm}
\includegraphics[width=0.4\textwidth]{Img/math_06d.pdf}
\fuentepropia
\end{figure}

Esta región se obtiene con las siguientes instrucciones:

\begin{tcolorbox}
\begin{Verbatim}[formatcom=\letraverbatim]
RegionPlot3D[0<z<1-y&&x^2<y<1,{x,-1,1},{y,0,1},
{z,0,1},PlotPoints->90,PlotStyle->Orange,
Mesh->False,Ticks->{{-1,0,1},{-1,1},{0,1}}]
\end{Verbatim}
\end{tcolorbox}

%\clearpage
(2) El sólido acotado por las figuras de $(x-1)^2+y^2= 1,$ $(x-3/2)^2+y^2=4,$ $z=0,$ $2z+y=4.$



\begin{figure}[H]
\centering
\caption{\small Sólido acotado por $x^2+2x+y^2=0,$ $x^2+y^2-3x=\frac{7}{4},$ $2z+y=4.$}
\includegraphics[width=0.4\textwidth]{Img/math_06.pdf} \hspace{5mm}
\includegraphics[width=0.4\textwidth]{Img/math_06a.pdf}

\includegraphics[width=0.4\textwidth]{Img/math_06b.pdf}
\fuentepropia
\end{figure}

Para obtenerlo se ejecutó la siguiente instrucción:
%
\begin{tcolorbox}
\begin{Verbatim}[formatcom=\letraverbatim] 
RegionPlot3D[z>0&&z<4-y/2&&(x-1)^2+y^2>1&&(x-1.5)^2
+y^2<4,{x,-1,4},{y,-2,2},{z,0,5},Mesh->False,
PlotPoints->90]
\end{Verbatim}
\end{tcolorbox}
%
(3) Las expresiones $(x^2+y^2)^2=z+1, \,\,\, z= 4(1-x^2-y^2), $ 
acotan el sólido que se muestra enseguida. 

\begin{figure}[H]
	\begin{center}
		\caption{Superficie determinada por $(x^2+y^2)^2=z+1$ y $z= 4(1-x^2-y^2)$.}
		\includegraphics[height=4.5cm]{Img/math_06e.pdf} \hspace{0.3cm}
		\includegraphics[height=4.5cm]{Img/math_06f.pdf}
\fuentepropia

\end{center}
\end{figure}

Las siguientes instrucciones permiten obtenerlo:
\begin{tcolorbox}
\begin{Verbatim}[formatcom=\letraverbatim]
RegionPlot3D[(x^2+y^2)^2-1<z<4*(1-x^2-y^2)^2&&
-Sqrt[1-x^2]<y<Sqrt[1-x^2],{x,-1,1},{y,-1,1},
{z,-1,4},PlotPoints->90,Mesh->False]
\end{Verbatim}
\end{tcolorbox}

(4) La región del espacio interior a la esfera $x^2+y^2+z^2=4$ y al hiperboloide de dos mantos $x^2+y^2=2+z^2$, exterior al cilindro $x^2+y^2=\frac{2}{5}$ se muestra en el siguiente gráfico.


\begin{figure}[H]
\centering
\caption{\small Sólido acotado por $x^2+y^2+z^2=4,$ $x^2+y^2=2+z^2,$ $x^2+y^2=\frac{2}{5}.$}
\includegraphics[width=0.4\textwidth]{Img/math_06n.pdf} \hspace{5mm}
\includegraphics[width=0.4\textwidth]{Img/math_06p.pdf}
\fuentepropia
\end{figure}

Este se consigue con las siguientes instrucciones:

\begin{tcolorbox}
\begin{Verbatim}[formatcom=\letraverbatim]
RegionPlot3D[x^2+y^2+z^2<4&&x^2+y^2<2+z^2&&
x^2+y^2>0.4,{x,-2,2},{y,-2,2},{z,-2,2},PlotPoints->90]
\end{Verbatim}
\end{tcolorbox}


(5) En coordenadas cilíndricas, las desigualdades 
$1 \leq r \leq 2,\,\, r \leq z \leq 2, $
determinan la región acotada por un cilindro, un semicono y un plano.


\begin{figure}[H]
\centering
\caption{\small Sólido acotado por $1 \leq r \leq 2,$ $r \leq z \leq 2.$}
\includegraphics[width=0.4\textwidth]{Img/math_06g.pdf} \hspace{5mm}
\includegraphics[width=0.4\textwidth]{Img/math_06h.pdf}
\end{figure}

Para obtener estas figuras se definen $r$ y $\theta$, 
luego se escriben las desigualdades y se agregan otras opciones:
\begin{tcolorbox} 
\begin{Verbatim}[formatcom=\letraverbatim]
{r,\[Theta]}={Sqrt[x^2+y^2],ArcTan[y/Abs[x]]};
RegionPlot3D[1<=r<=2&&r<=z<=2,{x,-2,2},{y,-2,2},
{z,1,2},PlotPoints->90,Mesh->True,AxesLabel->
Automatic,PlotRange->{{-2,2},{-2,2},{1,2.2}}]
\end{Verbatim}
\end{tcolorbox}


\subsection*{Gráficos en coordenadas esféricas}
En coordenadas esféricas se puede utilizar la 
función \verb"SphericalPlot3D", con una sintaxis genérica para graficar una función de la forma 
$\rho = f(\theta,\phi)$, con  $\theta \in [a,b]$ y $\phi \in [c,d];$
esto es: \index{SphericalPlot3D}

\begin{tcolorbox}[breakable,enhanced,title=Gráficos en coordenadas esféricas]
\begin{lstlisting}[language=Mathematica, basicstyle=\ttfamily]
SphericalPlot3D$[f[\[Theta],\[Phi]], 
{\[Theta],a,b}, {\[Phi],c,d}]
\end{lstlisting}
\end{tcolorbox}

Así, se ejecutan las siguientes instrucciones:

\begin{tcolorbox}[breakable,enhanced]
\begin{Verbatim}[formatcom=\letraverbatim]
SphericalPlot3D[1,{\[Theta],Pi/6,Pi},{\[Phi],
Pi/6,5*Pi/3}]
SphericalPlot3D[Cos[\[Phi]],{\[Theta],0,Pi},{\[Phi],
0,Pi}]
SphericalPlot3D[\[Theta],{\[Theta],0,Pi},{\[Phi],
0,2*Pi}]
\end{Verbatim}
\end{tcolorbox}



Es así como se obtienen, respectivamente, los resultados de la figura 21:


\begin{figure}[H]
\centering
\caption{\small Regiones en coordenadas esféricas.}
\label{regiones_esfe}
\includegraphics[width=0.4\textwidth]{Img/math_06aa.pdf} \hspace{5mm}
\includegraphics[width=0.4\textwidth]{Img/math_06ab.pdf}

\includegraphics[width=0.4\textwidth]{Img/math_06ac.pdf}
\end{figure}

De manera alternativa, se pueden usar desigualdades definiendo previamente las variables $\rho$, $\theta$ y $\phi.$ 
Por ejemplo, la región limitada por las siguientes expresiones
$1 \leq \rho \leq 2,\,\, 0\leq \theta\leq \frac{3\pi}{4}, \,\,
\frac{\pi}{6} \leq \phi \leq \frac{3\pi}{4}, $
se muestra a continuación. Para conseguirlo, se pueden definir las variables $\rho,\,\, \theta,\,\, \phi$
y luego establecer la región que se ilustra a continuación:


\begin{figure}[H]
\centering
\caption{\small Región determinada por $1\leq \rho \leq 2,$ $0\leq \theta\leq \frac{3\pi}{4}$ y $\frac{\pi}{6} \leq \phi \leq \frac{3\pi}{4}$}
\includegraphics[width=0.4\textwidth]{Img/math_06m.pdf} \hspace{5mm}
\includegraphics[width=0.4\textwidth]{Img/math_06l.pdf}

\includegraphics[width=0.4\textwidth]{Img/math_06k.pdf}
\fuentepropia
\end{figure}

Para conseguirlo, se pueden definir las variables $\rho$, $\theta$, $\phi$ y luego establecer la región que se ilustra a continuación:

\begin{tcolorbox}[breakable,enhanced,title=Regiones en el espacio]
\begin{Verbatim}[formatcom=\letraverbatim]
{\[Rho],\[Theta]}={Sqrt[x^2+y^2+z^2],ArcTan[y/Abs[x]]}
\[Phi]=ArcCos[z/\[Rho]];
RegionPlot3D[1<\[Rho]<2&&0<\[Theta]<3Pi/4&&Pi/6<
\[Phi]<3*Pi/4,{x,-2,2},{y,0,2},{z,-2,2},
PlotPoints->90,Mesh->False]
\end{Verbatim}
\end{tcolorbox}

Utilizando la función \verb"ArcTan[x,y]", se puede obtener la región del espacio determinada por 
las expresiones $1\leq \rho  \leq 2,$ $0\leq \theta \leq 3\pi/2,$  $\pi/4\leq \phi \leq 5\pi/6.$


\begin{figure}[H]
\centering
\caption{\small Región determinada por $1\leq\rho\leq 2,$ $0\leq\theta\leq 3\pi/2,$ $\pi/4\leq\phi\leq 5\pi/6.$}
\includegraphics[width=0.4\textwidth]{Img/T_50a.pdf} \hspace{5mm}
\includegraphics[width=0.4\textwidth]{Img/T_50b.pdf}
\fuentepropia
\end{figure}

La figura se consigue a partir de dos grupos 
de desigualdades que se concatenan con el operador lógico \emph{o}, que se denota por \verb"||".
Es decir, la desigualdad $0\leq \theta \leq 3\pi/2,$ se escribe en \verb"Mathamatica" como 
$0\leq \theta \leq \pi/2 \, || \,-\pi \leq \theta \leq -\pi/2.$  
 
\begin{tcolorbox}[breakable,enhanced,title=Instrucciones en Mathematica]
\begin{Verbatim}[formatcom=\letraverbatim]
{\[Theta],\[Rho]}={ArcTan[x,y],Sqrt[x^2+y^2+z^2]}
\[Phi]=ArcCos[z/\[Rho]]; 
RegionPlot3D[1<\[Rho]<2&&0<\[Theta]<Pi&&Pi/4<
\[Phi]<5*Pi/6||1<\[Rho]<2&&-Pi<\[Theta]<-Pi/2
&&Pi/4<\[Phi]<5*Pi/6,{x,-2,2},{y,-2,2},{z,-2,2},
PlotPoints->90,Mesh->False] 
\end{Verbatim}
\end{tcolorbox}
 
La función \verb"ToSphericalCoordinates[{x,y,z}]" permite transformar coordenadas rectangulares en esféricas, mientras que la transformación inversa se logra con \Verb!FromSphericalCoordinates[{\[Rho],\[Theta],\[Phi]}]!.
Ejemplo de lo anterior son la siguientes sentencias:
\begin{tcolorbox}[breakable,enhanced]
\begin{Verbatim}[formatcom=\letraverbatim]
ToSphericalCoordinates[{1,-1,1}]
FromSphericalCoordinates[{Sqrt[3],ArcTan[1],Pi/4}]
\end{Verbatim}
\end{tcolorbox}
\index{ToSphericalCoordinates} \index{FromSphericalCoordinates} 

\subsection*{Gráficos con restricciones}
Agregando \verb"RegionFunction->Function[{x,y,z},DESIGULADADES]" \\
dentro de las opciones
de un gráfico se pueden agregar restricciones a este. \index{RegionFunction}
Sobre la esfera $x^2+y^2+z^2=9$ se aplica la restricción $-2yz<x^2+z^2<9$
y el paraboloide $z=5-x^2-y^2$ con la condición $2x<x^2+y^2<9$.


\begin{figure}[H]
\centering
\caption{\small Superficies con restricciones en el espacio.}
\includegraphics[width=0.4\textwidth]{Img/T_47a.pdf} \hspace{5mm}
\includegraphics[width=0.4\textwidth]{Img/T_48a.pdf}
\fuentepropia
\end{figure}

Estas figuras se obtienen con las siguientes instrucciones:
%

\begin{tcolorbox}[breakable,enhanced,title=Uso del comando {\tt RegionFunction} ]
\begin{Verbatim}[formatcom=\letraverbatim]
ContourPlot3D[x^2+y^2+z^2==9,{x,-3,3},{y,-3,3},
{z,-3,3},RegionFunction->Function[{x,y,z},
-2*y*z<x^2+z^2<9],Mesh->False]
Plot3D[5-x^2-y^2,{x,-3,3},{y,-3,3},Mesh->False,
RegionFunction->Function[{x,y,z},2*x<x^2+y^2<9]]
\end{Verbatim}
\end{tcolorbox}

\subsection*{Gráficos discretos en el plano}
Con la instrucción \Verb!DiscretePlot[f(x),{x,a,b,r},ExtentSize->Right]! se muestra el gráfico 
discreto de $f(x)$ en $[a,b]$, con subintervalos de longitud $r.$ La disposición (o altura) de cada rectángulo
se controla con \verb"ExtentSize" cuyas opciones son \Verb!Left/Right/None/Full/All/Scaled["$\Delta$\verb"]!; 
también se puede incluir la opción \Verb!ExtentSize->{a,b}!, en donde \verb"a" y \verb"b" son la distancia del punto 
y el ancho del rectángulo.


En la figura 25 algunos puntos sobre la función $y=4-x+3x^2-x^3$, $x\in [-1,3]$ e 
incrementos de $0.2$ unidades. 


\begin{figure}[H]
\centering
\caption{\small Función $f(x)=4-x+3x^2-x^3$.}
\includegraphics[width=0.4\textwidth]{Img/T_43b.pdf} \hspace{5mm}
\includegraphics[width=0.4\textwidth]{Img/T_43a.pdf}
\fuentepropia
\end{figure}

Además de estas, las siguientes instrucciones muestran otras posibilidades.

\begin{tcolorbox}[breakable,enhanced,title=Instrucciones en Mathematica]
\begin{Verbatim}[formatcom=\letraverbatim]
DiscretePlot[4-x+3*x^2-x^3,{x,-1,3,0.2}]
DiscretePlot[4-x+3*x^2-x^3,{x,-1,3,0.2},
ExtentSize->All]
Show[DiscretePlot[4-x+3*x^2-x^3,{x,-1,3,0.2},
ExtentSize->All],Plot[4-x+3*x^2-x^3,{x,-1,3}]]
Show[DiscretePlot[4-x+3*x^2-x^3,{x,-1,3,0.24},
ExtentSize->{0.1,0.1}],Plot[4-x+3*x^2-x^3,
{x,-1,3},PlotStyle->Black]]
\end{Verbatim}
\end{tcolorbox}

\subsection*{Gráficos discretos en el espacio}
La instrucción \verb"ListPoint" permite graficar una lista de puntos en el plano, mientras que \verb"ListPointPlot3D" permite hacerlo en el espacio.
 
\index{ListPoint} \index{ListPointPlot3D}
Las siguientes figuras corresponden a la función $f(x,y)= 2+xy-x^2-y^3.$ Para esto, se genera una tabla con puntos para 
$x \in[-2,2],$ $y \in[-2,2]$ con incrementos en $x$ de $1/5$ e incrementos en $y$ de $1/4.$ 

\begin{figure}[H]
	\begin{center}
		\caption{Función $f(x,y)=2+xy-x^2-y^3$.}
		\includegraphics[height=4.5cm]{Img/T_44a.pdf}
		\includegraphics[height=4.5cm]{Img/T_44b.pdf}
		\includegraphics[height=4.5cm]{Img/T_44c.pdf}
	\end{center}
	\fuentepropia
\end{figure}


Las instrucciones necesarias para obtener la figura 26 son las siguientes: en ellas 
se usó el comando \verb"InterpolationOrder" en dos posibilidades.
\begin{tcolorbox}[breakable,enhanced,title=Instrucciones en Mathematica]
\begin{Verbatim}[formatcom=\letraverbatim]
ListPointPlot3D[Table[2+i*j-i^2-j^3,{i,-2,2,1/5},
{j,-2,2,1/4}]]
ListPlot3D[Table[2+i*j-i^2-j^3,{i,-2,2,1/5},
{j,-2,2,1/4}],
Mesh->None,InterpolationOrder->0]
ListPlot3D[Table[2+i*j-i^2-j^3,{i,-2,2,1/5},
{j,-2,2,1/4}],Mesh->None,InterpolationOrder->2]
\end{Verbatim}
\end{tcolorbox}

A partir de la función $z=f(x,y)$ con $x \in [a,b]$ e incrementos de $r_1$ unidades, $y \in [c,d]$ 
e incrementos de $r_2$ unidades, se puede trazar la figura de la función discreta 
obtenida mediante el uso de la instrucción 

\begin{tcolorbox}[breakable,enhanced,title=Instrucciones en Mathematica]
\begin{Verbatim}[formatcom=\letraverbatim]
DiscretePlot3D[f[x,y],{x,a,b,r1},{y,c,d,r2}]
\end{Verbatim}
\end{tcolorbox}
\index{DiscretePlot3D}

Considerando la función $f(x,y)=3+xy-x^2-y^3$ con $x \in [-1,1],$ $y \in [-1,1],$ e incrementos de $0.1$ y
$0.2$ respectivamente, se obtiene los siguientes gráficos.

 % Asegura alineación con el margen izquierdo


\begin{figure}[H]
\centering
\caption{\small Función $f(x,y)=3+xy-x^2-y^3$.}
\label{fig:funcion_T45}
\includegraphics[width=0.4\textwidth]{Img/T_45a.pdf} \hspace{5mm}
\includegraphics[width=0.4\textwidth]{Img/T_45b.pdf}
\end{figure}

Los gráficos anteriores se obtienen al ejecutar las siguientes sentencias.


\begin{tcolorbox}[breakable,enhanced,title=Instrucciones en Mathematica]
\begin{Verbatim}[formatcom=\letraverbatim]
DiscretePlot3D[3+x*y-x^2-y^3,{x,-1,1,0.2},{y,-1,1,0.1},
ExtentSize->Full]
DiscretePlot3D[3+x*y-x^2-y^3,{x,-1,1,0.2},{y,-1,1,0.2},
ExtentSize->{0.15,0.18},FillingStyle->Opacity[0.4]]
\end{Verbatim}
\end{tcolorbox}

\index{FillingStyle}
Los rectángulos se han separan un poco y con la instrucción \verb"FillingStyle" se hacen un poco traslúcidos.

Se presenta enseguida la grafica discreta de la función $f(x,y)=3+xy-x^2-y^3$ 
con la restricción $\frac{1}{4} \leq x^2+y^2 \leq 2$;
e incrementos en $x$ e $y$ de $0.1$ unidades en ambos casos.


\begin{figure}[H]
\centering
\caption{\small Función $f(x,y)=3+xy-x^2-y^3$.}
\includegraphics[width=0.4\textwidth]{Img/T_46a.pdf} \hspace{5mm}
\includegraphics[width=0.4\textwidth]{Img/T_46b.pdf}
\fuentepropia
\end{figure}

Este gráfico se consigue al ejecutar las siguientes instrucciones. La condición se obtiene al usar la opción \verb"RegionFunction". \index{RegionFunction}

\begin{tcolorbox}[breakable,enhanced,title=Instrucciones en Mathematica]
\begin{Verbatim}[formatcom=\letraverbatim]
DiscretePlot3D[3+x*y-x^2-y^3,{x,-2,2,0.1},{y,-2,2,0.1},
ExtentSize->{0.1,0.1},FillingStyle->Opacity[0.6],
RegionFunction->Function[{x,y,z},1/4<=x^2+y^2<=2]]
\end{Verbatim}
\end{tcolorbox}

\subsection*{Superficie de revolución}
El gráfico de la región que se obtiene de girar la función 
$y=f(x)$ alrededor del eje $y$, con $x\in[a,b]$, se obtiene con ayuda 
de la siguiente instrucción:
\begin{tcolorbox}[breakable,enhanced,title=Sólidos de revolución] 
\verb"RevolutionPlot3D[f(x),{x,a,b}]"
\end{tcolorbox}  
\index{RevolutionPlot3D}
Si la curva está parametrizada en la forma $\langle x(t),y(t),z(t) \rangle$ y esta gira alrededor del vector $\langle p,q,r \rangle,$ con  $t \in[a,b]$, la sintaxis que
se usa es la siguiente:
\begin{tcolorbox}[breakable,enhanced,title=Sólido de revolución con una curva parametrizada]
\verb"RevolutionPlot3D[{x(t),y(t),z(t)},{t,a,b},"\\
\verb"RevolutionAxis->{p,q,r}]"
\end{tcolorbox}
Por ejemplo, la función 
${\bf r}(t)= \langle 1 - \cos(t), \sin(t), t - t^2\rangle $, con $t\in [0, 1]$, gira alrededor del vector 
$\langle  1,1,0\rangle$.


\begin{figure}[H]
\centering
\caption{\small Superficie generada por ${\bf r}(t)= \langle 1 - \cos(t), \sin(t), t - t^2\rangle.$}
\includegraphics[width=0.4\textwidth]{Img/math_06ad.pdf} \hspace{5mm}
\includegraphics[width=0.4\textwidth]{Img/math_06ae.pdf}
\fuentepropia
\end{figure}

Esta región se obtiene al ejecutar la siguiente instrucción:
\begin{tcolorbox}[breakable,enhanced] 
\begin{Verbatim}[formatcom=\letraverbatim]
RevolutionPlot3D[{1-Cos[t],Sin[t],t-t^2},
{t,0,1},Ticks->{{0,0.5},{0,0.5},{-0.2,0.2}},
RevolutionAxis->{1,1,0}]
\end{Verbatim}
\end{tcolorbox}

\section{Regresión e interpolación}
A partir de un conjunto de puntos es posible obtener una función
que los contenga o los aproxime, presentaremos algunas funciones que permiten hacerlo.
 
\textls[-10]{Para obtener una combinación lineal $f(x),$ de las funciones 
$f_{1},f_{2},\cdots ,f_{k},$ que mejor ajusta a los puntos $\left( x_{1},y_{1}\right) ,\left(
x_{2},y_{2}\right) ,\cdots ,\left( x_{n},y_{n}\right)$ se puede usar la función \verb"Fit".  \index{Fit}
La función $f(x)$ minimiza la siguiente suma de los cuadrados}  
\begin{multline*}
\left( a_{1}f_{1}\left( x_{1}\right) +a_{2}f_{2}\left( x_{1}\right) 
 +\cdots +a_{n}f_{n}\left( x_{1}\right) -y_{1}\right) ^{2}+ \\ 
 \cdots +\left(a_{1}f_{1}\left( x_{1}\right) +a_{2}f_{2}\left( x_{1}\right) +
\cdots +a_{n}f_{n}\left( x_{1}\right) -y_{n}\right) ^{2}
\end{multline*}
En este caso se puede utilizar la siguiente instrucción genérica:
\begin{tcolorbox}[breakable,enhanced,title=Ajuste de puntos]
\begin{Verbatim}[formatcom=\letraverbatim]
Fit[{{x1,y1},{x2,y2},...,{xn,yn}},{f1,f2,...,fk},x].
\end{Verbatim}
\end{tcolorbox}
Si el usuario establece la forma de la función que desea, los coeficientes de esta combinación 
se consiguen con la función \verb"FindFit".  \index{FindFit}
 
La combinación lineal de las funciones $1,$ $x^2,$ $\sin x,$
que mejor ajusta los puntos  $(-1, 1),$ $(2,-1)$, $(1,3)$ y $(3,0)$,
se obtiene al ejecutar:
\begin{tcolorbox}[breakable,enhanced]
\begin{Verbatim}[formatcom=\letraverbatim]
Fit[{{-1, 1},{2,-1},{1,3},{3,0}},{1,x^2,Sin[x]},x]
\end{Verbatim}
\end{tcolorbox}
Si se desea obtener una expresión de la forma $f(x)=ax+bx^2+c \sin x$, 
los coeficientes se determinan con la función \verb"FindFit," esto es:
\begin{tcolorbox}[breakable,enhanced]
\begin{Verbatim}[formatcom=\letraverbatim]
FindFit[{{1,1},{2,2},{-1,2},{-2,2}},{a*x+b*x^2+c*Sin[x]},
{a,b,c},x].
\end{Verbatim}
\end{tcolorbox}
En las siguientes tres sentencias se define un conjunto de puntos, se halla la
combinación lineal de las funciones $1,x^{2},x^{3},x^{5}$ que mejor
ajustan estos puntos y se presenta su gráfico.
\begin{tcolorbox}[breakable,enhanced]
\begin{Verbatim}[formatcom=\letraverbatim]
datos={{1,1},{2,2},{-1,2},{-2,2}};
y1=Fit[datos,{1,x^2,x^3,x^5},x]
Show[Graphics[{Red,Point[datos]}],Plot[y1,{x,-2,2}]]
\end{Verbatim}
\end{tcolorbox}
Los siguientes tres ejemplos ilustran el uso de esta función dos y tres
variables independientes:

\begin{tcolorbox}[breakable,enhanced,title=Instrucciones en Mathematica]
\begin{Verbatim}[formatcom=\letraverbatim]
Fit[{{-1,1,2},{2,2,0},{-1,3,3}},{x^2*y,x*y^2},{x,y}]
Fit[{{0,Pi/2,2},{Pi/3,Pi/2,-1},{2*Pi/3,Pi,-2},
{-Pi,2*Pi,1}},{Sin[x],Cos[y],Cos[x*y]},{x,y}]
Fit[{{1,1,2,2},{1,2,2,0},{-1,-2,3,4},{-3,1,2,-1},
{-2,4,0,1}},{x^2*y*z,x*y^2,y*z^2,Sin[x+y],Cos[y+z]},
{x,y,z}]
\end{Verbatim}
\end{tcolorbox}

Para los puntos $(x_1,y_1),$ $(x_2,y_2),\cdots$ $(x_n,y_n),$ existe un polinomio $p(x)$ 
de grado $n-1$, tal que $p(x_i)=y_i$ para $i=1,2,\cdots,n$. 
Este polinomio se denomina \emph{polinomio de interpolación} y se determina 
de la siguiente manera:
\begin{tcolorbox}[breakable,enhanced,title=Polinomio de interpolación]
\begin{Verbatim}
InterpolatingPolynomial[{{x1,y1},{x2,y2},...
,{xn,yn}},x]
\end{Verbatim}
\end{tcolorbox}
\index{InterpolatingPolynomial}

Con las siguientes instrucciones se define un conjunto de puntos, se halla
el polinomio de interpolación y se presenta su gráfico
\begin{tcolorbox}[breakable,enhanced]
\begin{Verbatim}[formatcom=\letraverbatim]
puntos={{-1,1},{0,-1},{2,1},{4,-2}};
y=InterpolatingPolynomial[puntos,x];
Show[Graphics[{Red,Point[puntos]}],Plot[y,{x,-1,4}]]
\end{Verbatim}
\end{tcolorbox}

\section{Integración múltiple}
Para evaluar integrales dobles o triples en \verb"Wolfram Mathematica", 
se puede proceder como sigue.\\[0.2cm]
Para una integral doble de la forma 
$\ds \int_a^b \int_c^d f(x,y)dydx,$ la instrucción es 
\begin{tcolorbox}[breakable,enhanced,title=Integrales dobles]
\begin{Verbatim}[formatcom=\letraverbatim]
Integrate[f[x,y],{x,a,b},{y,c,d}]
\end{Verbatim}
\end{tcolorbox}
De manera alternativa, se puede utilizar la siguiente notación:\\
\verb"Integrate[Integrate[f[x,y,z],{x,a,b}],{y,c,d}]"; nótese en este caso el orden en la escritura. 
El comando \verb"NIntegrate" pemite evaluar numéricamente una integral. \index{NIntegrate}
Si se desea ver la integral sin evaluar se puede 
agregar la sentencia \verb"//HoldForm" y si se desea evaluar numéricamente una expresión se debe agregar
\verb"//N" al final de la instrucción. Presentamos algunas integrales y las instrucciones requeridas para su evaluación:
%\clearpage
\begin{enumerate}[leftmargin=0cm,itemindent=.75cm,labelwidth=\itemindent,labelsep=0cm,align=left,label={{\bf(\arabic*)}}]
\item %La integral 
$$\ds\int_0^2 \int_{-x}^{x^2} xy^2 dxdy,$$ %se evalúa con la siguiente orden
\begin{tcolorbox}[breakable,enhanced] 
\begin{Verbatim}
Integrate[x*y^2,{x,0,2},{y,-x,x^2}]
\end{Verbatim}
\end{tcolorbox}
\item La siguiente integral se evalúa numéricamente 
$$\ds \int_0^2 \int_{\sqrt{y}}^{2-y} 30x^2 dxdy,$$ %se puede usar la siguiente instrucción:
\begin{tcolorbox}[breakable,enhanced] 
\begin{Verbatim}[formatcom=\letraverbatim]
NIntegrate[30x^2,{y,0,1},{x,Sqrt[y],2-y}]
\end{Verbatim}
\end{tcolorbox}
\item %La integral impropia 
$$\int_{-\infty}^{\infty}\int_{0}^{\infty }e^{-x^{2}-y^{2}}dydx= \frac{\pi}{2} $$ % se evalúa como sigue
\begin{tcolorbox}[breakable,enhanced] 
\begin{Verbatim}[formatcom=\letraverbatim]
Integrate[Exp[-x^2-y^2],{x,-Infinity,Infinity},
{y,0,Infinity}]
\end{Verbatim}
\end{tcolorbox}
\end{enumerate}


La integral triple $$\int_a^b \int_c^d \int_e^s f(x,y,z)\, dzdydx,$$
se puede evaluar mediante esta instrucción:
\begin{tcolorbox}[breakable,enhanced,title=Integrales triples]
\begin{Verbatim}[formatcom=\letraverbatim]
Integrate[f[x,y,z],{z,r,s},{y,c,d},{x,a,b}]
\end{Verbatim}
\end{tcolorbox}
Ahora ilustraremos el uso de esta función con unas integrales triples:
\begin{enumerate}[leftmargin=0cm,itemindent=.75cm,labelwidth=\itemindent,labelsep=0cm,align=left,label={{\bf(\arabic*)}}]
\item % La integral, 
$$\ds \int_{-3}^3 \int_0^2 \int_0^1  x^2 y^3 z^4 dzdy dx$$ % se evalúa como sigue 
\begin{tcolorbox}[breakable,enhanced] 
\begin{Verbatim}[formatcom=\letraverbatim]
Integrate[x^2y^3z^4,{x,-3,3},{y,0,2},{z,0,1}]
\end{Verbatim}
\end{tcolorbox}
\item % La evaluación de la integral 
$$\ds \int_{0}^{1}\int_{-\sqrt[2]{x}}^{\sqrt[2]{x}}\int_{0}^{x}x^{2}dzdydx$$ % se hace con la orden
\begin{tcolorbox}[breakable,enhanced] 
\begin{Verbatim}[formatcom=\letraverbatim]
Integrate[x^2,{x,0,1},{y,-Sqrt[x],Sqrt[x]},{z,0,x}]
\end{Verbatim}
\end{tcolorbox}

\item Agregando \verb"//HoldForm" al final de la siguiente sentencia, se puede ver la forma sin evaluar de la integral:
$$\int_0^\pi\int_0^{\pi/6} \int_0^2 \rho^3 \cos(\theta)\sin^2(\phi)\, d\rho d \theta d\phi=4  $$
\begin{tcolorbox}[breakable,enhanced]
\begin{Verbatim}[formatcom=\letraverbatim] 
Integrate[\[Rho]^3*Cos[\[Theta]]*Sin[\[Phi]],
{\[Rho],0,2},{\[Theta],0,Pi/6},{\[Phi],0,Pi}]
\end{Verbatim}
\end{tcolorbox}

%\item La integral doble impropia, $$\int_{-\infty }^{\infty }\int_{-\infty }^{\infty }x^{2}e^{-x^{2}-y^{2}}dydx= \frac{\pi }{2} $$ se evalúa mediante la ejecución de la siguiente instrucción 
%\begin{tcolorbox}[breakable,enhanced]  \begin{Verbatim}[fontsize=\letraverbatim]
%Integrate[x^2*Exp[-x^2-y^2],{x,-Infinity,Infinity}, {y,-Infinity,Infinity}]\end{Verbatim}\end{tcolorbox}
\item % La integral triple impropia
$$\frac{1}{\sqrt{\pi }}\int_{-\infty}^{\infty}\int_{-\infty}^{\infty}\int_{-\infty}^{\infty}e^{-x^{2}-y^{2}-z^{2}}dzdydx= \pi $$
% se evalúa mediante la ejecución de la siguiente instrucción 
\begin{tcolorbox}[breakable,enhanced] 
\begin{Verbatim}[formatcom=\letraverbatim]
1/Sqrt[Pi]*Integrate[Exp[-x^2-y^2-z^2],
{x,-Infinity,Infinity},{y,-Infinity,Infinity},
{z,-Infinity,Infinity}]
\end{Verbatim}
\end{tcolorbox}
\end{enumerate}

Muchos de los cálculos y manipulaciones algebraicas de estas notas se hacen en \verb"Mathematica" 
y se exportan a un documento \LaTeX.
El código de una \verb"EXPRESION" se consigue mediante la instrucción
\verb"TeXForm[EXPRESION]". \index{TeXForm}
\section{Ecuaciones diferenciales con {\tt DSolve}}
Se denomina \emph{ecuación diferencial} a toda ecuación que involucre
una función desconocida y algunas de sus derivadas. La ecuación
diferencial se llama ecuación diferencial ordinaria, si contiene
derivadas ordinarias y se llama ecuación diferencial parcial, si en ella
aparecen algunas derivadas parciales de una función que depende de más de una variable. 
El \emph{orden} de la ecuación diferencial corresponde al orden de la mayor derivada involucrada.
Una \emph{solución} de la ecuación diferencial es una función definida en un conjunto abierto $I$, tal que para todo $x\in I$ sus derivadas existen y al reemplazarse en la ecuación diferencial la expresión es una identidad.
Una introducción a las ecuaciones diferenciales puede consultarse en \textcite{apostol_calculus_1999}.
\index{Ecuaciones diferenciales} \index{DSolve}
 El comando \verb"DSolve" resuelve una o varias ecuaciones diferenciales ordinarias o parciales.
La ecuación diferencial $y'(x)=f(x)$ se resuelve en \verb"Mathematica" mediante la ejecución de la siguiente instrucción
\verb"DSolve[y'[x]==f[x],y[x],x]". \\ Según el caso, se debe usar agregar como una lista de ecuaciones o una ecuación con sus respectivas condiciones.
\begin{tcolorbox}[breakable, enhanced, title={Uso del comando \texttt{DSolve}}, colback=gray!5, colframe=gray!80]
El comando \texttt{DSolve} resuelve una o varias ecuaciones diferenciales ordinarias o parciales.
La ecuación diferencial $y'(x) = f(x)$ se resuelve en \texttt{Mathematica} mediante la instrucción:
\texttt{DSolve[y'[x] == f[x], y[x], x]}.
Según el caso, se debe usar una lista de ecuaciones o una ecuación con sus respectivas condiciones.
\end{tcolorbox}

El comando \verb"NDSolve" permite resolver numéricamente una ecuación diferencial,
una exposición más detallada de esta función puede conseguirse \textcite{sofroniou_advanced_2008}. Ilustraremos lo anterior con unos ejemplos.
\begin{itemize}[leftmargin=0cm,itemindent=.5cm,labelwidth=\itemindent,labelsep=0cm,align=left]
\item La ecuación diferencial ordinaria de segundo orden
$$ y^{\prime \prime }\left( x\right) + y^{\prime }\left(x\right) =-e^{-x},$$ con las condiciones
$y\left( 0\right) =1,$ $y^{\prime }\left( 0\right) =1,$ tiene por solución $y\left( x\right) =xe^{-x}+1.$
\begin{tcolorbox}[breakable,enhanced]
\begin{Verbatim}[formatcom=\letraverbatim] 
DSolve[{y''[t]+y'[t]==-Exp[-t],y[0]==1,y'[0]==1},
y[t],t]
\end{Verbatim}
\end{tcolorbox}

\item La ecuación diferencial parcial de primer orden
$$\frac{\partial u}{\partial x}+\frac{\partial u}{\partial y}=e^{-y}\left(1-x\right) ,$$ con $u\left( 0,y\right) =0,$ $u\left( x,0\right) =x,$ posee a $u\left( x,y\right) =xe^{-y}$ como solución.

\begin{tcolorbox}[breakable,enhanced]
\begin{Verbatim}[formatcom=\letraverbatim] 
DSolve[{D[u[x,y],x]+D[u[x,y],y]==Exp[-y]*(1-x),
u[0,y]==0,u[x,0]==x},u[x,y],{x,y}]
\end{Verbatim}
\end{tcolorbox}

 
\item El sistema de ecuaciones diferenciales ordinarias de primer orden
$$\left. \begin{array}{l}
y'(t)-x'(t)+y(t)=-e^t(t+1)\\
x'(t)+y'(t)-x(t)=e^t-e^{-t}
\end{array}\right \}$$
con las condiciones iniciales $x(0)=0$, $y(0)=1$ posee a $x(t)= te^t$ y $y(t)=e^{-t}$ como solución.
 
\begin{tcolorbox}[breakable,enhanced,title=Instrucciones en Mathematica]
\begin{Verbatim}[formatcom=\letraverbatim]
DSolve[{y'[t]-x'[t]+y[t]==-Exp[t]*(t+1),x'[t]
+y'[t]-x[t]==Exp[t]-Exp[-t],x[0]==0,y[0]==1},
{x[t],y[t]},t]
\end{Verbatim}
\end{tcolorbox}

  
\end{itemize}

\section{Campos vectoriales}
Presentamos algunas funciones del programa \verb"Mathematica" útiles en el tratamiento de campos vectoriales.

\subsection*{Divergencia y rotacional} \index{Curl} \index{Div}
Con las funciones \verb"Curl" y \verb"Div" se obtiene el rotacional y la divergencia de un campo vectorial 
${\bf F}.$ Se debe establecer un orden para las variables.
\index{Divergencia} \index{Rotacional} 

En el siguiente ejemplo se calcula el rotacional y la divergencia de ${\bf F}$, \mbox{después} de definir 
el campo vectorial ${\bf F}(x,y,z)= \la x^2+y^2,y^3+z^2,x e^z \ra.$
\begin{tcolorbox}[breakable,enhanced]
\begin{Verbatim}[formatcom=\letraverbatim]
F[x_,y_,z_]:={x^2+y^2,y^3+z^2,x*Exp[z]}
Curl[F[x,y,z],{x,y,z}]
Div[F[x,y,z],{x,y,z}]
\end{Verbatim}
\end{tcolorbox}

\subsection*{Gráficos de campos vectoriales} \index{VectorPlot}
Con la instrucción \verb"VectorPlot[F(x,y),{x,a,b},{y,c,d}]" se obtiene un gráfico del campo vectorial
${\bf F}(x,y)$, para $x \in [a,b]$, $y\in [c,d]$. 
La sintaxis que se debe usar para obtener una representación del campo vectorial ${\bf F}(x,y,z)$ es
\verb"VectorPlot3D[F[x,y,z],{x,a,b},{y,c,d},{z,p,q}]". \index{VectorPlot3D}
El campo vectorial ${\bf F}(x,y)= \la \frac{-x-y}{\sqrt{x^2+y^2}}, \frac{x-y}{\sqrt{x^2+y^2}} \ra$, 
se representa para $x \in [-1,1],$ $y \in [-1,1].$ También se presentan algunas \emph{líneas de flujo}\footnote{Para un campo vectorial ${\bf F}$, una línea de flujo para ${\bf F}$ es una curva ${\bf r}(t)$ tal que
${\bf r}^\prime(t) = {\bf F}({\bf r}(t)).$ Esto es,
${\bf F}$ corresponde al campo de velocidades de la curva ${\bf r}(t)$}, esto se logra
utilizando la función \verb"StreamPlot".  \index{StreamPlot} \index{StreamStyle}\index{StreamScale}

\begin{figure}[H]
	\begin{center}
		\caption{Campo vectorial $\la \frac{-x+y}{\sqrt{x^2+y^2}}, \frac{x-y}{\sqrt{x^2+y^2}} \ra$.}
		\includegraphics[height=4.5cm]{Img/T_55a.pdf} \hspace{0.5cm}
		\includegraphics[height=4.5cm]{Img/T_55b.pdf}
		\parbox{0.95\linewidth}{\centering\fontsize{10pt}{10pt}\selectfont\textbf{Fuente:} elaboración propia.}
	\end{center}
\end{figure}

Las instrucciones requeridas son, respectivamente las siguientes.

\begin{tcolorbox}[breakable,enhanced,title=Instrucciones en Mathematica]
\begin{Verbatim}[formatcom=\letraverbatim] 
VectorPlot[{-(x+y)/Sqrt[x^2+y^2],(x-y)/
Sqrt[x^2+y^2]},
{x,-1,1},{y,-1,1},VectorStyle->Black]
StreamPlot[{-(x+y)/Sqrt[x^2+y^2],(x-y)/
Sqrt[x^2+y^2]},
{x,-1,1},{y,-1,1},
StreamStyle->Black,StreamScale->0.08]
\end{Verbatim}
\end{tcolorbox}

En estas notas nos ocuparemos de obtener la representación gráfica de un campo vectorial ${\bf F}(x,y,z)$ sobre una superficie determinada por la expresión $g(x,y,z)=0$, para $x\in [a,b],$  $y\in [c,d]$, $z\in [e,f]$. Esto se 
logra en \verb"Mathematica" utilizando la función \verb"SliceVectorPlot3D". La estructura general es la siguiente:
\begin{tcolorbox}[breakable,enhanced,title=Campos vectoriales sobre superficies] 
\begin{Verbatim}[formatcom=\letraverbatim]
SliceVectorPlot3D[F[x,y,x],{g[x,y,z]=0},{x,a,b},
{y,c,d},{z,e,f},Opciones]
\end{Verbatim}
\end{tcolorbox} \index{SliceVectorPlot3D}



La figura del campo vectorial ${\bf F}(x,y,z)= \la y,-x,x+z \ra$, sobre la superficie del sólido 
acotado por las figuras de las expresiones $z=2(1-x^2-y^2)$ y $z+1=(x^2+y^2)^2$ se obtiene ejecutando 
las siguientes instrucciones. La intersección está determinada por la ecuación $x^2+y^2=1$,
por este motivo se agregó la condición $x^2+y^2\leq 1$.



\begin{tcolorbox}[breakable,enhanced,title=Instrucciones en Mathematica]
\begin{Verbatim}[formatcom=\letraverbatim]
SliceVectorPlot3D[{y,-x,x+z},{(x^2+y^2)^2==z+1,
z==2*(1-x^2-y^2)},{x,-1,1},{y,-1,1},{z,-1,2.2},
RegionFunction->Function[{x,y,z},x^2+y^2<=1],
VectorScale->{0.18,Scaled[0.3],Automatic},
VectorStyle->Black,
VectorPoints->16,PlotStyle->Orange]
\end{Verbatim}
\end{tcolorbox}


\begin{figure}[H]
\centering
\caption{\small Campo vectorial ${\bf F}(x,y,z)= \la y,-x,x+z \ra$.}
\includegraphics[width=0.4\textwidth]{Img/T_28a.pdf} \hspace{5mm}
\includegraphics[width=0.4\textwidth]{Img/T_28b.pdf}
\end{figure}

La figura del campo vectorial ${\bf F}(x,y,z)= \la -y,x,z/2 \ra$ sobre la superficie del paraboloide $4z=4-x^2-y^2$ con la condición $x^2+y^2\leq 4$. 


\begin{figure}[H]
\centering
\caption{\small Campo vectorial ${\bf F}(x,y,z)= \la -y,x,z/2 \ra$.}
\includegraphics[width=0.4\textwidth]{Img/T_27a.pdf} \hspace{5mm}
\includegraphics[width=0.4\textwidth]{Img/T_27b.pdf}
\end{figure}

Esta figura se obtiene compilando la siguiente instrucción. En esta se usó el comando \verb"RegionFunction" que restringe el dominio en la figura de la superficie. \index{RegionFunction}



\begin{tcolorbox}[breakable,enhanced,title=Instrucciones en Mathematica]
\begin{Verbatim}[formatcom=\letraverbatim]
SliceVectorPlot3D[{-y,x,z/2},{4*z==4-x^2-y^2},
{x,-2,2},
{y,-2,2},{z,0,1.1},RegionFunction->Function[{x,y,z},
x^2+y^2<=4],
VectorScale->{0.15,Scaled[0.25],Automatic},
VectorStyle->Black,VectorPoints->16,
PlotStyle->Orange,
Ticks->{{-2,-1,0,1,2},{-2,-1,0,1,2},{0,1}}]
\end{Verbatim}
\end{tcolorbox}

También se puede obtener la representación gráfica de un campo vectorial con
la opción \verb"ListSliceVectorPlot3D." Si se considera el campo vectorial ${\bf F}(x,y,z)$
para $x\in[a,b],$ $y\in[c, d],$ $z\in[e,f]$, se crea una tabla que denominaremos datos,
posteriormente se dibuja el campo vectorial sobre la superficie determinada por
la expresión $g(x,y,z)=0$. En este caso la sintaxis básica es la siguiente:

\enlargethispage{\baselineskip}

\begin{tcolorbox}[breakable,enhanced,title=Instrucciones en Mathematica]
\begin{Verbatim}[formatcom=\letraverbatim] 
datos=Table[F[x,y,z],{x,a,b},{y,c,d},{z,e,f}];
ListSliceVectorPlot3D[datos,
ImplicitRegion[g[x,y,z]==0,{x,y,z}],
DataRange->{{a,b},{c,d},{e,f}}]
\end{Verbatim}
\end{tcolorbox}

El campo vectorial ${\bf F}(x,y,z) =\la x+z,y,2x-z \ra$ se representa 
sobre la superficie $z= (y^2-x)e^{1-x^2-y^2}$. 


\begin{figure}[H]
\centering
\caption{\small Campo vectorial ${\bf F}(x,y,z) =\la x+z,y,2x-z \ra$.}
\includegraphics[width=0.4\textwidth]{Img/T_33a.pdf} \hspace{5mm}
\includegraphics[width=0.4\textwidth]{Img/T_33b.pdf}
\end{figure}

Las instrucciones requeridas para obtener esta figura son las siguientes:
%ContourPlot3D[z==(y^2-x)*Exp[1-x^2-y^2],{x,-2,2},{y,-2,2},
%{z,-2,2},MeshStyle->Gray,ContourStyle->Orange]

\begin{tcolorbox}[breakable,enhanced,title=Instrucciones en Mathematica]
\begin{Verbatim}[formatcom=\letraverbatim] 
datos=Table[{x+z,y,2*x-z},{x,-2,2},{y,-2,2},
{z,-1,1}];
ListSliceVectorPlot3D[datos,
ImplicitRegion[z==(y^2-x)*Exp[1-x^2-y^2],{x,y,z}],
VectorPoints->12,DataRange->{{-2,2},{-2,2},{-2,2}},
VectorStyle->{Thin,Black,Thick}]
\end{Verbatim}
\end{tcolorbox}

La representación del campo vectorial ${\bf F}(x,y,z) =\la x-z,y,z+2 \ra$ sobre la superficie
del paraboloide hiperbólico $4z=x^2-y^2$ con la condición $x^2+y^2 \leq 4$, se obtiene con las 
siguientes instrucciones:

\begin{tcolorbox}[breakable,enhanced,title=Instrucciones en Mathematica]
\begin{Verbatim}[formatcom=\letraverbatim] 
datos=Table[{x-z,y,z+2},{x,-2,2},{y,-2,2},{z,-1,1}];
ListSliceVectorPlot3D[datos,
ImplicitRegion[4*z==x^2-y^2&&x^2+y^2<=4,{x,y,z}],
VectorPoints->12,DataRange->{{-2,2},{-2,2},
{-5/4,5/4}},VectorStyle->{Thin,Black,Thick}]
\end{Verbatim}
\end{tcolorbox}


\begin{figure}[H]
\centering
\caption{\small Campo vectorial ${\bf F}(x,y,z) =\la x-z,y,z+2 \ra$.}
\includegraphics[width=0.4\textwidth]{Img/T_54a.pdf} \hspace{5mm}
\includegraphics[width=0.4\textwidth]{Img/T_54b.pdf}
\end{figure}

\subsection*{Campos conservativos y función potencial}
Recordemos que si ${\bf F} = P(x, y,z)\vi +Q(x,y,z)\vj +P(x,y,z)\vk$ 
es un campo vectorial conservativo existe un campo escalar $\phi $ tal que $\nabla \phi ={\bf F}.$ 

%\clearpage
\begin{ejemplo}
El campo vectorial 
{\small ${\bf F}(x,y,z) = \left\langle y^{3}e^{-z},3xy^{2}e^{-z},3z^{2}-xy^{3}e^{-z}\right\rangle,$}
es conservativo. Las funciones $P$, $Q$ y $R$ son respectivamente las funciones componentes de ${\bf F}$. 
Para hallar la función potencial \index{Función!potencial} proceder como sigue, 
se integra $P$ con respecto a $x$, esto es 
$$\phi (x,y,z) = \int \left( y^{3}e^{-z}\right)
dx= xy^{3}e^{-z}+g\left( y,z\right) $$
Ahora se compara $\frac{\partial \phi }{\partial y}$ con $Q\left(x,y,z\right) ,$ esto es
$\frac{\partial \phi }{\partial y}= 3xy^{2}e^{-z}+\frac{\partial g}{\partial y}=3xy^{2}e^{-z},$ entonces 
$\frac{\partial g}{\partial y}=0,$ es decir, $g$ no depende de $y,$ por lo tanto, 
$$\phi (x,y,z) = xy^{3}e^{-z}+g\left(z\right) .$$
Luego se compara $\frac{\partial \phi }{\partial z}$ con 
$R(x,y,z),$ esto es $\frac{\partial \phi }{\partial z}= -xy^{3}e^{-z}+g^{\prime}(z) = 3z^{2}-xy^{3}e^{-z}$
entonces $g^{\prime}\left( z\right) =3z^{2},$ de lo cual se sigue que $g\left( z\right) =z^{3}+c.$ Por lo tanto, 
$$\phi (x,y,z) = xy^{3}e^{-z}+z^{3}+c.$$
\end{ejemplo}

En \verb"Wolfram Matematica" este campo escalar puede hallarse mediante la solución de
un sistema de ecuaciones diferenciales parciales: 
$$\left. \begin{array}{c}
\frac{\partial \phi }{\partial x}(x,y,z)=P(x,y,z) \\[0.15cm] 
\frac{\partial \phi }{\partial y}(x,y,z) =Q(x,y,z) \\[0.15cm] 
\frac{\partial \phi }{\partial z}(x,y,z) =R(x,y,z) 
\end{array}%
\right\} $$


Con las siguientes instrucciones además de determinar si el campo vectorial es conservativo, se halla la función potencial.
Se utiliza la función \verb"If", la cual permite elegir entre dos resultado dependiendo si la prueba es verdadera o falsa. \index{If}

\begin{tcolorbox}[breakable,enhanced,title=Instrucciones en Mathematica]
\begin{Verbatim}[formatcom=\letraverbatim]
{P,Q}={y^3*Exp[-z],3x*y^2*Exp[-z]};
R=3*z^2-x*y^3*Exp[-z];
If[D[P,y]==D[Q,x]&&D[P,z]==D[R,x]&&D[Q,z]==D[R,y],
"El Campo vectorial es conservativo",
"El Campo Vectorial No es conservativo"]
DSolve[{D[\[Phi][x,y,z],x]==P,D[\[Phi][x,y,z],y]==Q,
D[\[Phi][x,y,z],z]==R},\[Phi][x,y,z],x,y,z]
\end{Verbatim}
\end{tcolorbox}
\newpage 
\chapter{Parametrizaciones básicas}  \chapter{Parametrizaciónes básicas}

En este capítulo se presentan algunos ejemplos sencillos sobre la manera de representar curvas,
superficies o sólidos cuyas generalidades se presentaron en la sección \secref{regionesparametrizadas}. \index{Parametrización}
Con las ecuaciones (\ref{param1}) y (\ref{param2}) que corresponden a ideas muy simples, se construirán otras parametrizaciones y luego, con ayuda de estas
funciones, se parametrizarán regiones en el espacio. Se hacen énfasis en la construcción de este 
tipo de funciones, pues en los capítulos posteriores se hará uso de ellas sin profundizar demasiado en su obtención. 

\section{Curvas} \index{Parametrización!de curvas} \index{Curva} 
De acuerdo con \cite{marsden_calculo_2013}, una curva $C$ en $\mathbb{R}^n$ es la imagen de una función
contínua a trozos, $\ve{r}:[a,b]\rightarrow \mathbb{R}^n$. 
A la función $\ve{r}$ la denominaremos \emph{una parametrización} de la curva $C$. 
\vskip0.3cm 
Si $g:[c,d]\rightarrow \lbrack a,b]$ es una función biyectiva y continua en $[c,d]$
tal que $\alpha =r\circ g, \,$ a $\alpha,$ se le denomina una \emph{reparametrización} de $\ve{r}.$
Si la función $g$ es creciente $\left(g^{\prime }(t)>0\right), $
entonces $\alpha $ preserva la orientación, y si $g$ es decreciente
$\left( g^{\prime }(t)<0\right), $ entonces $\alpha $ invierte la orientación de la curva.
\vskip0.1cm
Dada una curva $C$ parametrizada por medio de la función $\ve{r}(t)$, con $t \in [a,b],$
la sustitución 
\begin{equation} \label{param3}
\ve{v}(t)=\ve{r}(a+b-t)
\end{equation}
parametriza la misma curva con sentido contrario para $t \in [a,b].$ \index{Cambio de sentido!Curvas}
%\vskip0.1cm

\vskip0.2cm
\noindent
\begin{minipage}{\textwidth}
	\centering
	\captionof{figure}{${\bf u}(t)=\langle t^3-3t, t \rangle$}
	\begin{minipage}{0.35\linewidth}
		\centering
		\includegraphics[width=\linewidth]{Fig/20.pdf}
	\end{minipage}
	
	\vspace{0.3em}
	\parbox{0.95\linewidth}{\centering\fontsize{10pt}{10pt}\selectfont\textbf{Fuente:} elaboración propia.}
\end{minipage}
\vskip0.2cm
 
 
 
La curva determinada por la ecuación $x=y^3-3y$, con  $y\in[-1,2]$
se puede parametrizar mediante la función 
$\ve{u}(t)=\langle t^3-3t, t \rangle$
con $t \in [-1,2]$. El punto inicial 
es  $(2,-1)$ y el punto final es $(2,2).$

Al cambiar $t$ por $-1+2-t=1-t$ y utilizar la expresión (\ref{param3}), 
se establece que la función $\ve{v}(t)= \la -t^3+3t^2-2,1-t\ra$
con $t \in [-1,2]$, que representa la misma curva, recorrida en sentido contrario.

\vskip0.2cm
\noindent
\begin{minipage}{\textwidth}
	\centering
	\captionof{figure}{${\bf u}(t)=\langle  t,t^4-3t^2 \rangle$}
	\begin{minipage}{0.35\linewidth}
		\centering
		\includegraphics[width=\linewidth]{Fig/21.pdf}
	\end{minipage}
	
	\vspace{0.3em}
	\parbox{0.95\linewidth}{\centering\fontsize{10pt}{10pt}\selectfont\textbf{Fuente:} elaboración propia.}
\end{minipage}
\vskip0.2cm

%\vskip0.1cm 
La función $y= x^4-3x^2,$ con $x\in[-2,2]$
se parametriza mediante la expresión 
$\ve{r}(t)=\langle t,t^4-3t^2 \rangle$ con $t \in [-2,2]$; el punto inicial de la curva descrita 
es  $(-2,4)$ y el punto final es $(2,4).$
Esta representación no es única, pues $\ve{u}(t)=\langle 2t,16t^4-12t^2 \rangle$ 
con $t \in [-1,1]$ representa la misma curva, con el mismo recorrido y los puntos terminales.

\subsection*{Tres ideas elementales}
Las siguientes tres parametrizaciones se usarán en gran parte de los problemas expuestos y servirán
como insumo para construir o representar varias regiones en el espacio.

\begin{tcolorbox}[breakable,enhanced,title=\bf Parametrización de un segmento]
Para dos puntos $a$ y $b$ en $\mathbb{R}^{n},$  la función 
\begin{equation} \label{param1}
\overrightarrow{r}(t)=(1-t) a+tb,\,\, t\in \lbrack 0,1]
\end{equation}
parametriza el segmento rectilíneo que une a $a$ con $b$. 
\end{tcolorbox}
Intercambiando $a$ y $b$ en la fórmula (\ref{param1}) se logra cambiar el punto inicial con el punto final.
Una expresión como esta la denominaremos una 
\emph{combinación convexa} entre los puntos $a$ y $b$. \index{Combinación convexa}

\begin{tcolorbox}[breakable,enhanced,title=\bf Parametrización de la circunferencia unitaria]
La circunferencia unitaria cuya ecuación es $x^2+y^2=1,$
puede parametrizarse en la forma 
\begin{equation} \label{param2}
\ve{r}(t)=\left\langle \cos(t), \sin(t)\right\rangle,\, t\in [0,2\pi].
\end{equation}
La identidad $\cos^2t + \sin^2 t =1$ motiva el uso de esta elección.
\end{tcolorbox}
En la expresión (\ref{param2}) también se puede asignar $x=\sin (t) $ y $y=\cos (t) $, de esta manera se describe la misma curva; la diferencia está en los puntos terminales o en la orientación.


\vskip0.2cm
\noindent
\begin{minipage}{\textwidth}
	\centering
	\captionof{figure}{Gráficos de ${\bf u}(t)=\langle\cos(t),\sin(t)\rangle$ y ${\bf u}(t)=\langle\sin(t),\cos(t)\rangle$}
	\begin{minipage}{0.40\linewidth}
		\centering
		\vspace{-2mm}
		\includegraphics[width=\linewidth]{Fig/22.pdf}
	\end{minipage}
	
	\vspace{0.3em}
	\parbox{0.95\linewidth}{\centering\fontsize{9pt}{10pt}\selectfont\textbf{Fuente:} elaboración propia.}
\end{minipage}
\vskip0.2cm


Para una misma curva, la parametrización presentada podría no ser única.
Por ejemplo, en el caso de la circunferencia unitaria también se puede \mbox{representar} por medio de la función
${\bf u}(t)=\langle \cos(3t), \sin(3t)\rangle$, $t \in [0,2\pi/3]$.

%\vspace{-10mm}

\begin{tcolorbox}[breakable,enhanced,title=\bf Parametrización de la hipérbola unitaria]
La ecuación $x^2-y^2=1$ representa una hipérbola cuyas ramas abren hacia el eje $x$.
Una porción de esta curva se parametriza por medio de la función 
$\ve{r}(t)=\la\cosh(t), \sinh(t)\ra,$ con $t\in [-1,1]$. La definición de esta función es motivada
por la relación $\cosh^2 t - \sinh^2 t =1.$ 
\end{tcolorbox}

%\vspace{-2mm}

La otra rama se representa por la función 
{\small $\ve{r}(t)=\la -\cosh(t), \sinh(t)\ra.$}  
Las siguientes parametrizaciones representan solo una rama de la hipérbola.
\vskip0.2cm
Con los elementos expuestos presentaremos otras curvas con su respectiva parametrización.

\vskip0.2cm
\noindent
\begin{minipage}{\textwidth}
	\centering
	\captionof{figure}{Gráficos de ${\bf r}(t)=\la \cosh(t),\sinh(t)\ra$ y ${\bf r}(t)=\la \sinh(t),\cosh(t)\ra$}
	\begin{minipage}{0.40\linewidth}
		\centering
		\includegraphics[width=\linewidth]{Fig/23.pdf}
	\end{minipage}
	
	\vspace{0.3em}
	\parbox{0.95\linewidth}{\centering\fontsize{9pt}{10pt}\selectfont\textbf{Fuente:} elaboración propia.}
\end{minipage}
\vskip0.2cm


\subsection*{Ejemplos de parametrización de curvas}

\vspace{0pt}
\textcolor{black}{\scriptsize\ding{110}}\quad El segmento de recta que une los puntos $A(1,1)$ y $B(3,2)$, iniciando en $A$ y terminando en $B$, se parametriza al simplificar la combinación convexa
{\small
	\setlength{\abovedisplayskip}{2pt}
	\setlength{\belowdisplayskip}{2pt}
	\[ (1-t)(1,1) + t(3,2), \]
}
la cual define la función
{\small
	\setlength{\abovedisplayskip}{2pt}
	\[ \ve{r}(t) = \left\langle 1+2t,1+t\right\rangle, \quad t \in [0,1]. \]
}
	
%\vskip0.1cm
\noindent
\begin{minipage}{\textwidth}
	\centering
	\captionof{figure}{Segmento en el plano}
	\begin{minipage}{0.24\linewidth}
		\centering
		\includegraphics[width=\linewidth]{Fig/24.pdf}
	\end{minipage}
	
	\vspace{0.3em}
	\parbox{0.95\linewidth}{\centering\fontsize{9pt}{10pt}\selectfont\textbf{Fuente:} elaboración propia.}
\end{minipage}
\vskip0.2cm		


\vspace{0pt}
\textcolor{black}{\scriptsize\ding{110}}\quad El segmento rectilíneo con punto inicial {\small $(2,-1,0)$} y punto final {\small $(-1, 2, 2),$} se parametriza simplificando una combinación convexa de los puntos
\[
(1 - t)(2,-1,0) + t(-1,2,2),
\]
lo cual permite definir la función
\[
\ve{r}(t) = \left\langle 2 - 2t,\,-1 + 3t,\, 2t \right\rangle, \quad t \in [0,1].
\]

\vskip0.2cm
\noindent
\begin{minipage}{\textwidth}
	\centering
	\captionof{figure}{Segmento en el espacio}
	\label{segmento_espacio}
	\begin{minipage}{0.30\linewidth}
		\centering
		\includegraphics[width=\linewidth]{Fig/18.pdf}
	\end{minipage}
	
	\vspace{0.3em}
	\parbox{0.95\linewidth}{\centering\fontsize{8pt}{10pt}\selectfont\textbf{Fuente:} elaboración propia.}
\end{minipage}
\vskip0.2cm




\textcolor{black}{\scriptsize\ding{110}}\quad La elipse centrada en el origen y alargada en el sentido de las $x$  
con ecuación $\frac{x^2}{9} + \frac{y^2}{4}=1$ se parametriza al hacer $x=3 \cos(t)$ y $y=2 \sin(t)$.  
Es decir, esta curva se representa por medio de la función  
$$\ve{r}(t)=\left\langle 3 \cos(t), 2\sin(t)\right\rangle,\ t\in [0,2\pi].$$

\vskip0.2cm
\noindent
\begin{minipage}{\textwidth}
	\centering
	\captionof{figure}{${\bf r}(t)=\langle 3\cos(t),2\sin(t)\rangle$}
	\begin{minipage}{0.28\linewidth}
		\centering
		\includegraphics[width=\linewidth]{Fig/25.pdf}
	\end{minipage}
	
	\vspace{0.3em}
	\parbox{0.95\linewidth}{\centering\fontsize{9pt}{10pt}\selectfont\textbf{Fuente:} elaboración propia.}
\end{minipage}
\vskip0.2cm


\vspace{0pt}
\textcolor{black}{\scriptsize\ding{110}}\quad Con $t \in [0,\pi]$ se describe la mitad de la elipse. Utilizando la fórmula (\ref{param3}), se obtiene $x = -3\cos(t)$ y $y = 2\sin(t)$. De esta manera se describe la misma curva, variando sus puntos terminales y la orientación.

\vskip0.2cm
\noindent
\begin{minipage}{\textwidth}
	\centering
	\captionof{figure}{${\bf r}(t) = \langle -3\cos(t),\, 2\sin(t) \rangle$}
	\label{semielipse_izquierda}
	\begin{minipage}{0.30\linewidth}
		\centering
		\includegraphics[width=\linewidth]{Fig/26.pdf}
	\end{minipage}
	
	\vspace{0.3em}
	\parbox{0.95\linewidth}{\centering\fontsize{8pt}{10pt}\selectfont\textbf{Fuente:} elaboración propia.}
\end{minipage}
\vskip0.2cm


\vspace{0pt}
\textcolor{black}{\scriptsize\ding{110}}\quad En el espacio, la circunferencia de radio $2$ centrada
en el punto $(0, 2, 0)$ y paralela al plano $xz,$
se puede parametrizar como 
$\ve{r}(t) = \la 2\cos t, 2, 2\sin t\ra,$ $t\in[0,2\pi].$ 
Obsérvese que esta curva se halla en el plano $y=2$.

\vskip0.2cm
\noindent
\begin{minipage}{\textwidth}
	\centering
	\captionof{figure}{Circunferencia en el espacio}
	\label{Circunferencia_espacio}
	\begin{minipage}{0.30\linewidth}
		\centering
		\includegraphics[width=\linewidth]{Fig/19.pdf}
	\end{minipage}
	
	\vspace{0.3em}
	\parbox{0.95\linewidth}{\centering\fontsize{8pt}{10pt}\selectfont\textbf{Fuente:} elaboración propia.}
\end{minipage}
\vskip0.2cm


\vspace{0pt}
\textcolor{black}{\scriptsize\ding{110}}\quad Las ecuaciones $x-y+3z=3$ y $x^2+y^2=4$ representan un plano y un cilindro. En el plano $xy$ la ecuación del cilindro corresponde a una circunferencia de radio $2$ que se parametriza con $x=2 \cos t,$ $y=2 \sin t.$ De la ecuación del plano se obtiene $z$.
Luego, la intersección de las superficies está dada por $\ve{r}(t)=\la 2\cos t, 2\sin t,1-\frac{2}{3}(\cos t+\sin t) \ra.$


\vskip0.2cm
\noindent
\begin{minipage}{\textwidth}
	\centering
	\captionof{figure}{Elipse en el espacio}
	\label{Elipse_espacio}
	\begin{minipage}{0.28\linewidth}
		\centering
		\includegraphics[width=\linewidth]{Fig/27.pdf}
	\end{minipage}
	
	\vspace{0.3em}
	\parbox{0.95\linewidth}{\centering\fontsize{8pt}{10pt}\selectfont\textbf{Fuente:} elaboración propia.}
\end{minipage}
\vskip0.2cm


\subsection*{Uniendo dos puntos con una curva}
Para dos puntos $P$ y $Q$ en $\mathbb{R}^n$, si $f$ y $g$ son dos funciones contínuas en $[a,b]$ tales que 
$f(a)=1,$ $g(a)=0$ y $f(b)=0,$ $g(b)=1,$ entonces la función
\begin{equation}
\ve{r}(t)=f(t)P+g(t)Q,\,\, t \in[a,b].
\end{equation}
parametriza la curva que une a $P$ con $Q$. 
%Además $\ve{r}(t)$ satisface que $\ve{r}(a)=P$ y $ \ve{r}(b)=Q;$
La forma de la trayectoria depende de las funciones $f$ y $g.$
Enseguida ilustraremos lo anterior. 


\vspace{0pt}
\textcolor{black}{\scriptsize\ding{110}}\quad Con las funciones $g(t)=\sqrt[4]{t}$ 
y $f(t)=(1-t)^3$, uniremos los puntos $A(1,1)$ y $B(3,2)$ mediante la expresión 
$(1-t)^3 A + \sqrt[4]{t} B.$ 
\vskip0.05cm
Esto permite definir la función 
${\bf r}(t)=(1-t) ^{3}(1,1) +\sqrt[4]{t}(3,2),$
con  $t\in[0,1].$


\vskip0.2cm
\noindent
\begin{minipage}{\textwidth}
	\centering
	\captionof{figure}{${\bf r}(t)=(1-t) ^{3}(1,1) +\sqrt[4]{t}(3,2)$}
	\begin{minipage}{0.26\linewidth}
		\centering
		\includegraphics[width=\linewidth]{Fig/28.pdf}
	\end{minipage}
	
	\vspace{0.3em}
	\parbox{0.95\linewidth}{\centering\fontsize{8pt}{10pt}\selectfont\textbf{Fuente:} elaboración propia.}
\end{minipage}
\vskip0.2cm


\vspace{0pt}
\textcolor{black}{\scriptsize\ding{110}}\quad Con ayuda de las funciones $g(t) = t^6$ y $f(t) = (1 - t^2)^3$, se define la función
\[
\ve{r}(t) = (1 - t^2)^3(1,1) + t^6(3,2),
\]
que une los puntos $A(1,1)$ y $B(3,2)$ mediante la combinación
\[
(1 - t^2)^3 A + t^6 B, \quad t \in [0,1].
\]


\vskip0.2cm
\noindent
\begin{minipage}{\textwidth}
	\centering
	\captionof{figure}{${\bf r}(t) = (1 - t^2)^3(1,1) + t^6(3,2)$}
	\begin{minipage}{0.26\linewidth}
		\centering
		\includegraphics[width=\linewidth]{Fig/29.pdf}
	\end{minipage}
	
	\vspace{0.3em}
	\parbox{0.95\linewidth}{\centering\fontsize{8pt}{10pt}\selectfont\textbf{Fuente:} elaboración propia.}
\end{minipage}
\vskip0.2cm


\vspace{0pt}
\textcolor{black}{\scriptsize\ding{110}}\quad Utilizando una interpolación cuadrática entre los puntos $(1,1),$ $(3,2)$ y un punto intermedio 
entre ellos, se obtienen las siguientes funciones: \\[0.2cm]
$y=\frac{51-61t+16t^2}{6}$, $y=\frac{27-29t+8t^2}{6}$, $y=\frac{22-27t+7t^2}{6}$, 
$y=\frac{-16+33t-7t^2}{10},$ $y=\frac{-5+16t-3t^2}{8}$. \\[0.2cm]
\\Se muestran estas curvas parametrizadas con
$\la t, y(t)\ra $, $t\in[1,3].$

\vskip0.2cm
\noindent
\begin{minipage}{\textwidth}
	\centering
	\captionof{figure}{Cinco polinomios que unen $A$ con $B$}
	\begin{minipage}{0.29\linewidth}
		\centering
		\includegraphics[width=\linewidth]{Fig/30.pdf}
	\end{minipage}
	
	\vspace{0.3em}
	\parbox{0.95\linewidth}{\centering\fontsize{8pt}{10pt}\selectfont\textbf{Fuente:} elaboración propia.}
\end{minipage}
\vskip0.2cm



Con la función \verb"InterpolatingPolynomial" se consiguen dos animaciones: una con el gráfico y la otra con 
las expresiones que lo definen. \index{InterpolatingPolynomial}

\begin{tcolorbox}[breakable,enhanced,title=\bf Instrucción Mathematica] 
\begin{Verbatim}[formatcom=\letraverbatim]
Manipulate[Expand[InterpolatingPolynomial[{{1,1},
{7/5,k*6/5},{3,2}},x]],{k,12/10,2,1/10}]
Manipulate[Plot[Expand[InterpolatingPolynomial[{{1,1},
{7/5,k*6/5},{3,2}},x]],{x,1,3}],{k,-2,2,0.1}]
\end{Verbatim}
\end{tcolorbox}


\vspace{0pt}
\textcolor{black}{\scriptsize\ding{110}}\quad Utilizando las funciones $f(t)=\cos(\frac{\pi}{2}t)$ y $g(t)=\sin(\frac{\pi}{2}t)$, se define
$$\ve{r}(t)=\cos\left( \frac{\pi t}{2}\right)\left(1,1\right) +\sin \left( \frac{\pi t}{2}\right) \left( 3,2\right),$$
que permite unir los puntos $A(1,1)$ y $B(3,2).$
Esta expresión se consigue, al simplificar la expresión
$$\cos\left( \frac{\pi t}{2}\right)\left(1,1\right) +\sin \left( \frac{\pi t}{2}\right) \left( 3,2\right)$$
para $\,\,0\leq t\leq 1$. 
Por tanto, la parametrización escogida es $\ve{r}(t).$


\vskip0.2cm
\noindent
\begin{minipage}{\textwidth}
	\centering
	\captionof{figure}{${\bf r}(t)=\cos\left( \frac{\pi t}{2}\right)A +\sin \left( \frac{\pi t}{2}\right)B$}
	\begin{minipage}{0.28\linewidth}
		\centering
		\includegraphics[width=\linewidth]{Fig/31.pdf}
	\end{minipage}
	
	\vspace{0.3em}
	\parbox{0.95\linewidth}{\centering\fontsize{8pt}{10pt}\selectfont\textbf{Fuente:} elaboración propia.}
\end{minipage}
\vskip0.2cm


\section{Superficies} \index{Parametrización!de superficies} \label{superficies}
En esta sección construiremos algunas parametrizaciones de superficies a partir 
de los elementos usados en el tratamiento de curvas.
\subsection*{Parametrización de regiones planas}

Con las ideas elementales utilizadas para construir curvas, se pueden construir parametrizaciones de regiones en el plano.


\vskip0.2cm
\noindent
\begin{minipage}{\textwidth}
	\centering
	\captionof{figure}{${\bf u}(r,\theta) =\la r \cos \theta, r \sin \theta \ra$}
	\label{fig:32}
	\begin{minipage}{0.30\linewidth}
		\centering
		\includegraphics[width=\linewidth]{Fig/32.pdf}
	\end{minipage}
	
	\vspace{0.3em}
	\parbox{0.95\linewidth}{\centering\fontsize{8pt}{10pt}\selectfont\textbf{Fuente:} elaboración propia.}
\end{minipage}
\vskip0.2cm


Con el propósito de \emph{rellenar} una circunferencia de radio $a$, se puede variar el radio y el arco en las expresiones
$x= r \cos \theta$ y $y= r \sin \theta.$ Este círculo puede parametrizarse mediante la función $\ve{u}(r,\theta)= \la r \cos \theta, r \sin \theta \ra$,
en donde $r \in [0,a],$ $\theta \in [0,2 \pi].$

\vskip0.2cm

La expresión $4x^2+9y^2=36$ representa una elipse, esta curva se puede parametrizar mediante 
las ecuaciones $x=3\cos \theta,$ $y=2\sin \theta$.
La región interior se puede describir
mediante la función
${\bf u}(r,\theta) = \la 3 r \cos \theta,2r \sin\theta \ra,$
con $r\in [0,1],$ $\theta \in [0,2\pi].$


\vskip0.2cm
\noindent
\begin{minipage}{\textwidth}
	\centering
	\captionof{figure}{${\bf u}(r,\theta) =\la 3r \cos \theta, 2r \sin \theta \ra$}
	\label{fig:33}
	\begin{minipage}{0.26\linewidth}
		\centering
		\includegraphics[width=\linewidth]{Fig/33.pdf}
	\end{minipage}
	
	\vspace{0.3em}
	\parbox{0.95\linewidth}{\centering\fontsize{8pt}{10pt}\selectfont\textbf{Fuente:} elaboración propia.}
\end{minipage}
\vskip0.2cm


\vspace{0pt}
El triángulo con vértices $A$, $B$ y $C$, se parametriza a través de segmentos rectilíneos.
Se construye una combinación convexa entre un punto del segmento $\overline{AB}$ con el punto $C$.
Una función que describe esta región, con $u \in [0,1],$ $v \in [0,1],$ se obtiene simplificando la expresión


\vskip0.2cm
\noindent
\begin{minipage}{\textwidth}
	\centering
	\captionof{figure}{Triángul parametrizado $(1-v)C+v\left((1-u)A+uB \right).$}
	\label{fig:161}
	\begin{minipage}{0.38\linewidth}
		\centering
		\includegraphics[width=\linewidth]{Fig/161.pdf}
	\end{minipage}
	
	\vspace{0.3em}
	\parbox{0.95\linewidth}{\centering\fontsize{8pt}{10pt}\selectfont\textbf{Fuente:} elaboración propia.}
\end{minipage}
\vskip0.2cm


$$(1-v)C+v\left((1-u)A+uB \right).$$


El cuadrilátero con vértices $A(-2,-1),$ $B(-1,4)$, $C(3,1),$ $D(4,3),$ 
se parametriza formando dos combinaciones convexas adecuadamente elegidas (evitar crear un \emph{corbatín}) 
con el mismo parámetro; con las expresiones obtenidas se forma otra combinación 
convexa con un parámetro distinto. 


\vskip0.2cm
\noindent
\begin{minipage}{\textwidth}
	\centering
	\captionof{figure}{Cuadrilátero parametrizado}
	\label{fig:Cuadrilátero_parametrizado}
	\begin{minipage}{0.32\linewidth}
		\centering
		\includegraphics[width=\linewidth]{Fig/35.pdf}
	\end{minipage}
	
	\vspace{0.3em}
	\parbox{0.95\linewidth}{\centering\fontsize{8pt}{10pt}\selectfont\textbf{Fuente:} elaboración propia.}
\end{minipage}
\vskip0.2cm

Combinando las expresiones $(1-t)A+tB$ y $(1-t)C+tD$,
se obtiene la función $\ve{r}(t,u)= (1-u)\left ((1-t)A+tB\right )+u \left( (1-t)C+tD\right),$
con $t\in [0,1]$, $u\in [0,1].$ En este caso, %La función obtenida es 
$\ve{r}(t,u)=\la t+5u-2, 5t+2u-3tu -1\ra.$

Con las siguientes instrucciones se consigue la figura \ref{fig:Cuadrilátero_parametrizado}:
\par
\vskip -12pt

\begin{tcolorbox}[breakable,enhanced,title=\bf Instrucción \verb"Mathematica"] 
\begin{Verbatim}[formatcom=\letraverbatim]
ParametricPlot[(1-u)*((1-t)*{-2,-1}+t*{-1,4})
+u*((1-t)*{3,1}+t*{4,3}),{t,0,1},{u,0,1}]
\end{Verbatim}
\end{tcolorbox}


\vspace{0pt}
La región $R$ acotada por las figuras de las funciones $y=f(x)$ y $y=g(x)$,
con $x\in [a,b];$ se parametriza formando una combinación convexa entre dos puntos
$P(x,f(x))$ y $Q(x,g(x)).$ Por tanto, 
una función que describe a $R$ es
$\ve{r}(x,u)= (1-u)(x,f(x))+u(x,g(x)),$
con $u \in [0,1]$, $x \in [a,b].$


\vskip0.2cm
\noindent
\begin{minipage}{\textwidth}
	\centering
	\captionof{figure}{Región parametrizada}
	\label{fig:Región_parametrizada}
	\begin{minipage}{0.32\linewidth}
		\centering
		\includegraphics[width=\linewidth]{Fig/36.pdf}
	\end{minipage}
	
	\vspace{0.3em}
	\parbox{0.95\linewidth}{\centering\fontsize{8pt}{10pt}\selectfont\textbf{Fuente:} elaboración propia.}
\end{minipage}
\vskip0.2cm


De manera análoga al caso anterior, si la región está acotada por las figuras de las expresiones 
$x=u(y)$ y $x=v(y)$, con $y\in [c,d];$ con lo cual, para
$u \in [0,1]$, $y \in [c,d],$ 
una función que describe la dicha región es %\\[0.1cm]
\[
\ve{r}(y,u)= (1-u)\la y,u(y)\ra + u\la y,v(y)\ra
\]

\vskip0.2cm
\noindent
\begin{minipage}{\textwidth}
	\centering
	\captionof{figure}{Región parametrizada}
	\label{fig:Región_parametrizada_3}
	\begin{minipage}{0.32\linewidth}
		\centering
		\includegraphics[width=\linewidth]{Fig/37.pdf}
	\end{minipage}
	
	\vspace{0.3em}
	\parbox{0.95\linewidth}{\centering\fontsize{8pt}{10pt}\selectfont\textbf{Fuente:} elaboración propia.}
\end{minipage}
\vskip0.2cm


\vspace{0pt}
La región $R,$ acotada por las figuras de $y=4-x^2$ y $y=x^3-4x$ para $x\in[-1,2]$
se parametriza simplificando la expresión 
$(1-u)\la x, 4-x^2 \ra +u\la x, x^3-4x \ra $, es decir 
$${\bf r}(x,u)=\la x, (x^2-4)(ux+u-1) \ra,$$ con $x\in[-1,2],$ $u\in[0,1],$
representa a la mencionada región.


\vskip0.2cm
\noindent
\begin{minipage}{\textwidth}
	\centering
	\captionof{figure}{$y=4-x^2$ y $y=x^3-4x$}
	\label{fig:38}
	\begin{minipage}{0.28\linewidth}
		\centering
		\includegraphics[width=\linewidth]{Fig/38.pdf}
	\end{minipage}
	
	\vspace{0.3em}
	\parbox{0.95\linewidth}{\centering\fontsize{8pt}{10pt}\selectfont\textbf{Fuente:} elaboración propia.}
\end{minipage}
\vskip0.2cm


\vspace{0pt}
Con las ecuaciones $y=x+1$ y $y=x^2-1,$ para $x\in[-1,2],$
se forma la combinación convexa $(1-u)\la x, 1+x \ra +u\la x, x^2-1 \ra;$
simplificando se define la función ${\bf r}(x,u)=\la x, (x+1) (1-2u+ux) \ra,$ con $x\in[-1,2],$ $u\in[0,1],$
que determina la región mostrada.
	
\vskip0.2cm
\noindent
\begin{minipage}{\textwidth}
	\centering
	\captionof{figure}{$y=x+1$ y $y=x^2-1$}
	\label{fig:39}
	\begin{minipage}{0.33\linewidth}
		\centering
		\includegraphics[width=\linewidth]{Fig/39.png}
	\end{minipage}
	
	\vspace{0.3em}
	\parbox{0.95\linewidth}{\centering\fontsize{8pt}{10pt}\selectfont\textbf{Fuente:} elaboración propia.}
\end{minipage}
\vskip0.2cm	
	

\vspace{0pt}
La región acotada por las graficas de las expresiones $y=x^4-x^2,$ $y=4-x^2$ y $x=1$, 
se parametriza por medio de la expresión
$(1-u)(x,x^4-x^2)+u(x,4-x^2),$
que al simplificarse, con $x \in [-\sqrt{2},1]$, $u \in [0,1],$ permite definir la función \\
$\ve{r}(x,u)=\la x, u(4-x^4)+x^2(x^2-1) \ra.$


\vskip0.2cm
\noindent
\begin{minipage}{\textwidth}
	\centering
	\captionof{figure}{$\la x, u(4-x^4)+x^2(x^2-1) \ra$}
	\label{fig:region_gris}
	\begin{minipage}{0.35\linewidth}
		\centering
		\includegraphics[width=\linewidth]{Fig/40.pdf}
	\end{minipage}
	
	\vspace{0.3em}
	\parbox{0.95\linewidth}{\centering\fontsize{8pt}{10pt}\selectfont\textbf{Fuente:} elaboración propia.}
\end{minipage}
\vskip0.2cm	


Con las siguientes instrucciones se obtiene la figura \ref{fig:region_gris} en \verb"Mathematica"



\begin{tcolorbox}[breakable,enhanced,title=\bf Instrucción \verb"Mathematica"]  
\begin{Verbatim}[formatcom=\letraverbatim]
B=ParametricPlot[{x,u*(4-x^4)+x^2*(x^2-1)},
{x,-Sqrt[2],1},{u,0,1}]; 
A=ContourPlot[{y==x^4-x^2,y==4-x^2},{x,-2,2},{y,-1,4}];
Show[A,B]
\end{Verbatim}
\end{tcolorbox}


\vspace{0pt}
Las parábolas $x=y^2-5y,$ $x=3y-y^2$ se cortan 
en $(0,0)$ y $(-4,4).$ La región acotada por estas curvas se parametriza 
simplificando la expresión
$(1-v)\la u^2-5u,u\ra +v\la 3u-u^2,u\ra,$
con $u \in [0,4]$, $v \in [0,1]$. Esta región se describe por medio de la función
$\ve{r}(u,v) = \la u (u+8v-2uv-5),u \ra.$


\vskip0.2cm
\noindent
\begin{minipage}{\textwidth}
	\centering
	\captionof{figure}{$ \la (8v-5)u+(1-2v)u^2,u \ra$}
	\label{fig:41}
	\begin{minipage}{0.27\linewidth}
		\centering
		\includegraphics[width=\linewidth]{Fig/41.pdf}
	\end{minipage}
	
	\vspace{0.3em}
	\parbox{0.95\linewidth}{\centering\fontsize{8pt}{10pt}\selectfont\textbf{Fuente:} elaboración propia.}
\end{minipage}
\vskip0.2cm	


\begin{tcolorbox}[breakable,enhanced,title=\bf Instrucción \verb"Mathematica"] 
\begin{Verbatim}[formatcom=\letraverbatim]
A=ContourPlot[{x==y^2-5*y,x==3*y-y^2},{x,-7,4},
{y,-1,5}];
B=ParametricPlot[(1-v){u^2-5*u,u}+v*{3*u-u^2,u},
{v,0,1},{u,0,4}];
Show[A,B]
\end{Verbatim}
\end{tcolorbox}

\subsection*{Superficies parametrizadas} \index{Superficies parametrizadas}
Recordemos un concepto que permite hacer estas construcciones en el espacio.

\begin{tcolorbox}[breakable,enhanced,title=\bf Definición]
Una parametrización de una superficie $\sigma $,
es una aplicación $$\ve{r}:[a,b]\times \lbrack c,d]\rightarrow \mathbb{R}^{3},$$
tal que al variar los parámetros $u\in \lbrack a,b]$ y $v\in \lbrack c,d]$
su imagen $\ve{r}(u,v)$ va describiendo los puntos de la superficie $\sigma $.
\end{tcolorbox}

\textcolor{black}{\scriptsize\ding{110}}\quad Utilizando coordenadas esféricas, la superficie de la esfera unitaria se puede parametrizar 
por medio de la función
{\small
\[{\bf r}(\theta, \phi)=\big \langle \cos \theta \sin \phi, \sin \theta \sin \phi, \cos \phi \big\rangle,\]
}
con $\theta \in [0,2\pi]$ y $ \phi \in [0,\pi]$.

\vskip0.2cm

El elipsoide con ecuación $\frac{x^2}{16}+\frac{y^2}{9}+\frac{z^2}{4}=1,$ 
se parametriza con una sencilla variación de la función anterior
${\bf r}(\theta, \phi)=\la 4\cos \theta \sin \phi,3\sin \theta \sin \phi, 2\cos \phi \ra.$ 

\vskip0.2cm
\noindent
\begin{minipage}{\textwidth}
	\centering
	\captionof{figure}{\small La esfera unitaria y el elipsoide parametrizado por ${\bf r}(\theta, \phi)$}
	\begin{minipage}{0.20\linewidth}
		\centering
		\includegraphics[width=\linewidth]{Img/T_53a.pdf}
	\end{minipage}
	\begin{minipage}{0.30\linewidth}
		\centering
		\includegraphics[width=\linewidth]{Img/T_53b.pdf}
	\end{minipage}
	\begin{minipage}{0.30\linewidth}
		\centering
		\includegraphics[width=\linewidth]{Img/T_53c.pdf}
	\end{minipage}
	
	\vspace{0.3em}
	\parbox{0.95\linewidth}{\centering\fontsize{8pt}{10pt}\selectfont\textbf{Fuente:} elaboración propia.}
\end{minipage}
\vskip0.2cm


Esta figura se obtiene al ejecutar las siguientes instrucciones:

\begin{tcolorbox}[breakable,enhanced,title=\bf Instrucción \verb"Mathematica"]  
\begin{Verbatim}[formatcom=\letraverbatim]
ParametricPlot3D[{Cos[\[Theta]]*Sin[\[Phi]],
Sin[\[Theta]]*Sin[\[Phi]],Cos[\[Phi]]},{\[Theta],0,
2*Pi},{\[Phi],0,Pi},MeshStyle->Gray,PlotStyle->Orange]

ParametricPlot3D[{4Cos[\[Theta]]*Sin[\[Phi]],3*
Sin[\[Theta]]*Sin[\[Phi]],2*Cos[\[Phi]]},{\[Theta],0,
2*Pi},{\[Phi],0,Pi},MeshStyle->Gray,PlotStyle->Orange]
\end{Verbatim}
\end{tcolorbox}


%\vspace{-5mm}
\textcolor{black}{\scriptsize\ding{110}}\quad La superficie $\sigma$ sobre el paraboloide 
$2z=9-x^2-y^2$ y dentro del cilindro $x^2+ \left(y -\frac{1}{3} \right)^2=4,$ se parametriza al asignar $x=u \cos v,$ 
$y= \frac{1}{3}+u\sin v$, con $u \in [0,2]$, $v \in [0,2\pi].$ 
Del paraboloide se obtiene $z = \frac{1}{18} (80-9u^2-6u\sin(t)).$


\vskip0.2cm
\noindent
\begin{minipage}{\textwidth}
	\centering
	\captionof{figure}{Superficie $\sigma$}
	\label{Superficie_sigma}
	\begin{minipage}{0.38\linewidth}
		\centering
		\includegraphics[width=\linewidth]{Fig/42.pdf}
	\end{minipage}
	
	\vspace{0.3em}
	\parbox{0.95\linewidth}{\centering\fontsize{8pt}{10pt}\selectfont\textbf{Fuente:} elaboración propia.}
\end{minipage}
\vskip0.2cm


Entonces, para $u \in [0,2]$, $v \in [0,2\pi],$ la superficie $\sigma$ se parametriza por
\[{\bf r}(u,v)= \langle u\cos v, \frac{1}{3} +u\sin v, \frac{1}{18}(80-9u^2-6u \sin(v)) \rangle.\]
Si $u=2$ la expresión obtenida corresponde a la parametrización de la curva de intersección del paraboloide
y el cilindro, es decir ${\bf r}(v)= \langle 2\cos v, \frac{1}{3} +2\sin v,\frac{1}{9}(22-6\sin(v))\rangle,$ $v \in [0,2\pi].$


\vspace{0pt}
\textcolor{black}{\scriptsize\ding{110}}\quad La superficie del cilindro con ecuación $x^2+y^2=2x$, para $z \in [-1,2],$ se 
parametriza en la forma 
$${\bf r}(u,t)= \la 1+\cos(t),\sin(t), u \ra,$$ con $t \in [0,2\pi],$ $u\in [-1,2].$ 
En el plano $xy$ el cilindro corresponde a la circunferencia $(x-1)^2+y^2=1.$


\vskip0.2cm
\noindent
\begin{minipage}{\textwidth}
	\centering
	\captionof{figure}{${\bf r}(u,t)= \la 1+\cos t,\sin t, u \ra$}
	\label{fig:43}
	\begin{minipage}{0.32\linewidth}
		\centering
		\includegraphics[width=\linewidth]{Fig/43.pdf}
	\end{minipage}
	
	\vspace{0.3em}
	\parbox{0.95\linewidth}{\centering\fontsize{8pt}{10pt}\selectfont\textbf{Fuente:} elaboración propia.}
\end{minipage}
\vskip0.2cm


\textcolor{black}{\scriptsize\ding{110}}\quad Las expresiones $2z=x^{2}-y^{2}$ y $2x^{2}+y^{2}+2y=8$ representan un paraboloide hiperbólico y un cilindro elíptico, respectivamente. La ecuación
del cilindro representa una elipse en el plano $xy$, el interior de esta región 
se parametriza como $x=\frac{3}{\sqrt{2}}u\cos (t),$ $y=-1+3u\sin (t),$ con la ecuación del paraboloide hiperbólico 
se halla $z$, esto es 
$z=\frac{1}{2}\left( x^{2}-y^{2}\right) =\frac{27}{4}u^{2}\cos^{2}(t) +3u\sin (t) -\frac{9}{2}u^{2}-\frac{1}{2}.$ 
Entonces, la función
{\small
$${\bf r}\left( t,u\right) =\left\langle \frac{3}{\sqrt{2}}u\cos (t),3u\sin (t)-1 ,\frac{27}{4}u^{2}\cos ^{2}(t)
+3u\sin (t) -\frac{9u^{2}+1}{2} \right\rangle,$$
}

con $t\in \lbrack 0,2\pi ],$ $u\in \lbrack 0,1],$  parametriza la región sobre el paraboloide 
hiperbólico acotada por el cilindro.
\vskip0.2cm
\noindent
\begin{minipage}{\textwidth}
	\centering
	\captionof{figure}{\small Superficie sobre el paraboloide hiperbólico.}
	\begin{minipage}{0.23\linewidth}
		\centering
		\includegraphics[width=\linewidth]{Img/73a.pdf}
	\end{minipage}
	\hspace{1em}
	\begin{minipage}{0.23\linewidth}
		\centering
		\includegraphics[width=\linewidth]{Img/73b.pdf}
	\end{minipage}
	
	\vspace{0.3em}
	\parbox{0.95\linewidth}{\centering\fontsize{8pt}{10pt}\selectfont\textbf{Nota:} elaboración propia.}
\end{minipage}
\vskip0.2cm

La anterior figura se consigue con las siguientes instrucciones.

\begin{tcolorbox}[breakable,enhanced,title=\bf Instrucción \verb"Mathematica"]  
\begin{Verbatim}[formatcom=\letraverbatim]
a=ContourPlot3D[{2*z==x^2-y^2,2*x^2+y^2+2*y==8},
{x,-5,5},{y,-5,5},{z,-9,5},Mesh->False,ContourStyle
->{Directive[Orange,Opacity[0.6]],Directive[Cyan,
Opacity[0.2]]},PlotPoints->80,Ticks->{{-3,0,3},
{-3,0,3},{-8,-4,-2,2,4}}];
b=ParametricPlot3D[{3*u*Cos[t]/Sqrt[2],-1+3*u*Sin[t],
27*u^2*Cos[t]^2/4+3*u*Sin[t]-9*u^2/2-1/2},{t,0,2*Pi},
{u,0,1},PlotStyle->{Cyan},PlotPoints->80];
Show[a,b]
\end{Verbatim}
\end{tcolorbox}

\textcolor{black}{\scriptsize\ding{110}}\quad Las ecuaciones $x^{2}+\left( y-3\right) ^{2}=4$ y $4z=36-\left(x^{2}+y^{2}\right) $ 
representan un cilindro y un paraboloide invertido con vértice en el punto $(0,0,9)$.
Con el fin de parametrizar la superficie $\sigma$ sobre el paraboloide exterior al cilindro, con $z\geq 0,$
analizaremos la base y luego proyectaremos esta región sobre el paraboloide.
La figura \ref{Superficie_paraboloide_exterior_cilindro} del sólido en mención se obtiene con ayuda de las siguientes instrucciones:

\begin{tcolorbox}[breakable,enhanced,title=\bf Instrucción \verb"Mathematica"] 
\begin{Verbatim}[formatcom=\letraverbatim]
ContourPlot3D[{z==9-(x^2+y^2)/4,x^2+(y-3)^2==4},
{x,-6,6},{y,-6,6},{z,0,9},Mesh->False,ContourStyle
->{Orange,Cyan}]
ParametricPlot3D[{2*(2*u+1)*Sin[t],2*(2*u+1)Cos[t]
-3*u+3,(u-1)*(12*(1+2u)*Cos[t]-23-25u)/4},{u,0,1},
{t,0,2*Pi}]
\end{Verbatim}
\end{tcolorbox}


\noindent
\begin{minipage}{\textwidth}
	\centering
	\captionof{figure}{Superficie sobre el paraboloide exterior al cilindro}
	\label{Superficie_paraboloide_exterior_cilindro}
	\begin{minipage}{0.26\linewidth}
		\centering
		\includegraphics[width=\linewidth]{Img/86a.pdf}
	\end{minipage}
	\hspace{1em}
	\begin{minipage}{0.33\linewidth}
		\centering
		\includegraphics[width=\linewidth]{Fig/44.pdf}
	\end{minipage}
	
	\vspace{0.3em}
	\parbox{0.95\linewidth}{\centering\fontsize{8pt}{10pt}\selectfont\textbf{Fuente:} elaboración propia.}
\end{minipage}
\vskip0.2cm





La base del cilindro se puede parametrizar como
\begin{equation} \label{sup5a}
x=2\sin (t), \,\, y=3+2\cos (t), \, t\in \lbrack 0,2\pi ].
\end{equation}

Mientras que la intersección del paraboloide y el plano $z=0$ se parametriza como $x=6\sin (t) $, $y=6\cos (t) $, con $t\in \lbrack 0,2\pi ].$


Con el fin de parametrizar la proyección de $\sigma$ sobre el plano $xy$, se puede formar 
una combinación convexa entre dos puntos de la forma 
$P = \left( 2\sin(t),\, 3 + 2\cos(t),\, 0 \right)$ y 
$Q = \left( 6\sin(t),\, 6\cos(t),\, 0 \right)$. 
La expresión que permite describir la región acotada por las dos circunferencias está dada por 
\[
(1 - u)\,P + u\,Q.
\]

\vskip0.2cm
\noindent
\begin{minipage}{\textwidth}
	\centering
	\captionof{figure}{Proyección de $\sigma$}
	\label{fig:45}
	\begin{minipage}{0.33\linewidth}
		\centering
		\includegraphics[width=\linewidth]{Fig/45.pdf}
	\end{minipage}
	
	\vspace{0.3em}
	\parbox{0.95\linewidth}{\centering\fontsize{8pt}{10pt}\selectfont\textbf{Fuente:} elaboración propia.}
\end{minipage}
\vskip0.2cm

Por tanto, la función que parametriza la región sobre el plano $xy$ acotada por las dos 
circunferencias presentadas antes es
\begin{equation}
\ve{v}(u,t)=\left( 2\left( 2u+1\right) \sin t,2\left( 2u+1\right) \cos t-3u+3,0\right) \label{sup5}
\end{equation}
La coordenada $z$ sobre la superficie del paraboloide es
\begin{eqnarray*}
z &=&9-\frac{1}{4}\left( \left( 2\left( 2u+1\right) \sin t\right) ^{2}+\left(
2\left( 2u+1\right) \cos t-3u+3\right) ^{2}\right) \\
&=& \frac{1}{4}\left( 1-u\right) \left( 25u-12(2u+1) \cos t+23\right)
\end{eqnarray*}
Por lo anterior, la superficie $\sigma$ se parametriza mediante la función 
\begin{multline*}
{\bf r}(t,u) =\big\langle 2(2u+1) \sin t,2(2u+1)\cos t-3u+3, \\
\frac{1}{4}( 1-u) ( 25u-12(2u+1)\cos t+23)  \big\rangle,
\end{multline*}
donde $t \in [0,2\pi],$ $u \in [0,1].$
Sobre la intersección de las superficies se satisface que 
$z=9-\frac{1}{4}\left( \left( 2\sin(t)\right)^{2}+\left(3+2\cos(t)\right)^{2}\right) = \frac{23}{4}-3\cos(t);$
dicha curva se parametriza como 
$$\ve{v}(t)= \langle 2\sin (t), 3+2\cos (t), \frac{23}{4}-3\cos(t) \rangle, \, t \in [0,2\pi].$$
\vskip0.3cm
\textcolor{black}{\scriptsize\ding{110}}\quad Con algunos de los elementos anteriores construiremos el siguiente ejemplo.
Utilizando la ecuación (\ref{sup5}), se sigue que
la coordenada $z$ de la curva de intersección del cilindro $x^2+(y-3)^2=4$ y el paraboloide hiperbólico 
$5z=x^2-y^2$ está dado por
{\small
\[
z=\frac{x^{2}-y^{2}}{5}=\frac{1}{5}\left((2\sin(t))^{2}-(3+2\cos(t))^{2}\right)=-1-\frac{12}{5}\cos t-\frac{8}{5}\cos ^{2}t.
\]
}

La región determinada por la ecuación (\ref{sup5}), se proyecta sobre el paraboloide hiperbólico con ecuación
$5z=x^2-y^2;$ la coordenada $z$ está dada por
{\small
\begin{eqnarray*}
z &=& \frac{1}{5}\left( x^{2}-y^{2}\right) =\frac{1}{5}\left( \left( 2\left(1+2u\right) \sin t\right) ^{2}-\left( 3-3u+\left( 2+4u\right) \cos t\right)^{2}\right) \\
&=& \frac{1}{5}\left( u+5\right) \left( 7u-1\right) -\frac{8}{5}\left( 1+2u\right) ^{2}\cos ^{2}t+\frac{12}{5}\left( 1+2u\right) \left(u-1\right) \cos t
\end{eqnarray*}
}

Entonces, la función
{\small
\begin{multline*}
{\bf r}(t,u) = \big \langle 2\left( 1+2u\right) \sin t,3-3u+\left( 2+4u\right) \cos t, \frac{1}{5} (u+5) (7u-1)\\
-\frac{8}{5}\left( 1+2u\right) ^{2}\cos ^{2}t+\frac{12}{5}\left( 1+2u\right) \left( u-1\right) \cos t\big \rangle, 
\end{multline*}
}
$u\in \lbrack 0,1],$ $t\in \lbrack 0,2\pi ],$ parametriza la superficie sobre
la silla de montar con ecuación  $5z=x^{2}-y^{2},$ exterior al cilindro con ecuación 
$x^{2}+(y-3)^{2}=4,$ e interior al paraboloide con ecuación 
$z=9-\left( x^{2}-y^{2}\right) /4.$
La figura \ref{Superficie_paraboloide_hiper} ilustra dicha región.

\noindent
\begin{minipage}{\textwidth}
	\centering
	\captionof{figure}{Superficie sobre el paraboloide hiperbólico $5z=x^2-y^2$}
	\label{Superficie_paraboloide_hiper}
	\begin{minipage}{0.47\linewidth}
		\centering
		\includegraphics[width=\linewidth]{Fig/46.pdf}
	\end{minipage}
	
	\vspace{0.3em}
	\parbox{0.95\linewidth}{\centering\fontsize{8pt}{10pt}\selectfont\textbf{Fuente:} elaboración propia.}
\end{minipage}
\vskip0.2cm

La figura de la superficie se obtiene ejecutando las siguientes instrucciones. Este es un ejemplo de parametrización en \verb"Mathematica".

\begin{tcolorbox}[breakable, enhanced, title={Código en \verb"Mathematica"}]
\begin{Verbatim}[formatcom=\letraverbatim]
{x,y}=(1-u)*{2*Sin[t],3+2*Cos[t]}+u*{6*Sin[t],6*Cos[t]};
r={x,y,(x^2-y^2)/5};
ParametricPlot3D[r,{u,0,1},{t,0,2Pi},ViewPoint->{-3,6,2}]
\end{Verbatim}
\end{tcolorbox}

Con los elementos anteriores es posible parametrizar otras superficies.
\begin{tcolorbox}[breakable,enhanced,title=\bf Una manera de parametrizar una superficie]
Para dos funciones vectoriales $\ve{\alpha}$ y $\ve{\beta}$ en $\mathbb{R}^n$, si $f$ y $g$ son
funciones reales contínuas en el intervalo $[a,b]$, que satisfacen que 
$f(a)=1,$ $g(a)=0$ y $f(b)=0,$ $g(b)=1,$ entonces la función vectorial
$$\ve{r}(u,t)=f(u)\ve{\alpha}(t)+g(u)\ve{\beta}$$
es una parametrización de la superficie $\sigma$ acotada por las curvas determinadas por $\ve{\alpha}$ y $\ve{\beta}$.
\end{tcolorbox}

Veamos algunos ejemplos de la expresión anterior.

\textcolor{black}{\scriptsize\ding{110}}\quad La cardioide $r=1-\cos(t),$ en el plano $z=2,$ se define por medio de la función vectorial 
${\bf a}(t)= \la (1-\cos(t))\cos(t),(1-\cos(t))\cos(t),2  \ra,$  con $t \in [0,2\pi],$
mientras que la función 
${\bf b}(t)= \la 3\cos(t), 3\sin(t), 0\ra,$ con $t \in [0,2\pi],$ 
determina una circunferencia en el plano $xy.$ 
Una combinación convexa de estas dos funciones vectoriales determina, mediante segmentos rectilíneos,
la superficie que se muestra a continuación. \index{Cardioide} 
\index{Combinación convexa}

\vskip0.2cm
\noindent
\begin{minipage}{\textwidth}
	\centering
	\captionof{figure}{Superficie que conecta una cardioide y una circunferencia}
	\label{Superficie_cardioide_circunferencia}
	\begin{minipage}{0.21\linewidth}
		\centering
		\includegraphics[width=\linewidth]{Fig/47.pdf}
	\end{minipage}
	\hspace{1em}
	\begin{minipage}{0.23\linewidth}
		\centering
		\includegraphics[width=\linewidth]{Img/42a.pdf}
	\end{minipage}
	\hspace{1em}
	\begin{minipage}{0.25\linewidth}
		\centering
		\includegraphics[width=\linewidth]{Img/42b.pdf}
	\end{minipage}
	
	\vspace{0.3em}
	\parbox{0.95\linewidth}{\centering\fontsize{8pt}{10pt}\selectfont\textbf{Fuente:} elaboración propia.}
\end{minipage}
%\vskip0.2cm
%\clearpage

Simplificando la expresión $(1-u){\bf a}(t)+u{\bf b}(t),$
se obtiene la función que parametriza esta superficie
{\small
\[{\bf r}(u,t)= \langle \cos(t)((u-1)\cos(t)+1+2u), \sin(t)((u-1)\cos(t)+1+2u),2-2u\rangle,\]
}
con $t\in[0,2\pi],\, u\in[0,1],$ la figura anterior se obtiene al ejecutar las siguientes instrucciones:


\begin{tcolorbox}[breakable,enhanced,title=\bf Instrucción \verb"Mathematica"] 
\begin{Verbatim}[formatcom=\letraverbatim]
q={(1-Cos[t])*Cos[t],(1-Cos[t])*Sin[t],2};
r=(1-u)*q+u*{3*Cos[t],3*Sin[t],0};
a=ParametricPlot3D[{q,{3*Cos[t],3*Sin[t],0}},{t,0,
2*Pi},PlotStyle->{Directive[Black,Thickness->0.01],
PlotPoints->90,Directive[Black,Thickness->0.01]}];
b=ParametricPlot3D[r,{t,0,2*Pi},{u,0,1},
PlotStyle->Orange,MeshStyle->Gray];
Show[a,b]
\end{Verbatim}
\end{tcolorbox}

\textcolor{black}{\scriptsize\ding{110}}\quad Con las funciones vectoriales
$\overrightarrow{r}_1(t)=\la 3\cos(t),\sin (t),4+\sin ( t) \ra $ y
$\overrightarrow{r}_2(t)=\la 2\cos(t),3\sin (t),-1+\cos ( t) \ra $ con $ t\in \lbrack 0,2\pi],$
que representan dos elipses en el espacio, se forma una combinación convexa de ellas con parámetro 
$u\in \lbrack 0,1].$
\\
Simplificando la expresión $(1-u)\overrightarrow{r}_1(t)+u\overrightarrow{r}_2(t)$
se obtiene la función:
$${\bf r}(u,t)=\la (3-u)\cos(t),(1+2u)\sin(t),4+(1-u)\sin(t)-5u+u\cos(t)\ra,$$
para $u\in[0,1],\,\, t\in[0,2\pi],$ la cual representa la superficie que se obtiene de unir las dos elipses mediante segmentos rectilíneos. 



\vskip0.2cm
\noindent
\begin{minipage}{\textwidth}
	\centering
	\captionof{figure}{\small Superficie que une dos elipses mediante segmentos}
	\label{conos_67}
	\begin{minipage}{0.17\linewidth}
		\centering
		\includegraphics[width=\linewidth]{Img/33a.pdf}
	\end{minipage}
	\begin{minipage}{0.17\linewidth}
		\centering
		\includegraphics[width=\linewidth]{Img/33b.pdf}
	\end{minipage}
	\begin{minipage}{0.17\linewidth}
		\centering
		\includegraphics[width=\linewidth]{Img/33c.pdf}
	\end{minipage}
	
	\vspace{0.3em}
	\parbox{0.95\linewidth}{\centering\fontsize{8pt}{10pt}\selectfont\textbf{Fuente:} elaboración propia.}
\end{minipage}
\vskip0.2cm


La figura \ref{conos_67} se consigue ejecutando las siguientes instrucciones:

\begin{tcolorbox}[breakable,enhanced,title=\bf Instrucción \verb"Mathematica"]  
\begin{Verbatim}[formatcom=\letraverbatim]
ParametricPlot3D[(1-u)*{3*Cos[t],Sin[t],4+Sin[t]}+
u*{2*Cos[t],3*Sin[t],-1+Cos[t]},{t,0,2*Pi},{u,0,1},
MeshStyle->Blue,PlotStyle->{Opacity[0.9]}]
\end{Verbatim}
\end{tcolorbox}

También se pueden generar dos figuras por separado llamadas \verb"a" y \verb"b", 
luego, con ayuda de la función \verb"Show", se muestran en una sola.

\begin{tcolorbox}[breakable,enhanced,title=\bf Instrucción \verb"Mathematica"]  
\begin{Verbatim}[formatcom=\letraverbatim]
a=ParametricPlot3D[{{3*Cos[t],Sin[t],4+Sin[t]},
{2*Cos[t],3*Sin[t],-1+Cos[t]}},{t,0,2*Pi},
PlotStyle->{Directive[RGBColor[0.2,0.1,1],
Thickness->0.015],Directive[Black,Thickness->0.015]}];
b=ParametricPlot3D[(1-u)*{3*Cos[t],Sin[t],4+Sin[t]}+
u*{2*Cos[t],3*Sin[t],-1+Cos[t],{t,0,0,2*Pi},{u,0,1},
MeshStyle->Gray,PlotStyle->{Orange,Opacity[0.9]}];
Show[a,b]
\end{Verbatim}
\end{tcolorbox}

%\fontsize{9.5}{12}\selectfont
\textcolor{black}{\scriptsize\ding{110}}\quad Las dos circunferencias en el espacio, representadas por las funciones vectoriales 
$\overrightarrow{\alpha}(t)=\langle 2\cos(t), 2\sin(t),0\rangle $ y $\overrightarrow{\beta}(t)=\langle \cos(t),\sin(t),2\rangle,$
con $t \in [0,2\pi].$ 
Con ellas se forma la siguiente combinación:
$$\frac{2\sqrt[4]{1-u}}{e^u+1}\, \overrightarrow{\alpha}(t) + \sqrt[4]{u}\,\overrightarrow{\beta}(t).$$
Con diferentes valores de $u= \frac{1}{4}, \,\frac{1}{2},\,2$, se logra la figura \ref{superficies_circunferencias68}


\begin{figure}[H]
	\begin{center}
		\captionof{figure}{\small Superficies que conectan dos circunferencias.}
		\label{superficies_circunferencias68}
		\includegraphics[width=2.75cm,height=1.5cm]{Img/8c.pdf}
		\includegraphics[width=2.75cm,height=1.5cm]{Img/8b.pdf}
		\includegraphics[width=2.75cm,height=1.5cm]{Img/8a.pdf}
	\end{center}
	\parbox{0.95\linewidth}{\centering\fontsize{8pt}{10pt}\selectfont\textbf{Fuente:} elaboración propia.}
\end{figure}


Con ayuda de la función \verb"Manipulate" se pueden generar varios elementos como los anteriores. 
Los comandos necesarios son los siguientes:

\begin{tcolorbox}[breakable,enhanced,title=\bf Instrucción \verb"Mathematica"] 
\begin{Verbatim}[formatcom=\letraverbatim]
Manipulate[ParametricPlot3D[2*(1-u)^k/(Exp[u]+1)
*{2*Cos[t],2*Sin[t],0}+u^k*{Cos[t],Sin[t],2},
{t,0,2Pi},{u,0,1}],{k,0.5,2,0.1}]
\end{Verbatim}
\end{tcolorbox}

\textcolor{black}{\scriptsize\ding{110}}\quad Con las funciones $f(t)=te^{t-1}$ y $g(t)=(1-t)e^{t}$, 
es sencillo ver que $f(0)=0=g(1)$ y $f(1)=1=g(0)$. Además, con las funciones vectoriales
\vspace{-2mm}
{\small
\[
\ve{\alpha}(t) = \left\langle 3\cos(t),2\sin(t),2-\frac{\sin(t)}{2}\right\rangle 
\, \text{y} \, 
\ve{\beta}(t) = \left\langle 2\cos(t),\sin(t),\frac{\sin (2t)}{2} \right\rangle
\]
}

se define la función  
${\bf r}(u,t)= f(u)\ve{\alpha}(t)+g(u)\ve{\beta}(t),$ $t\in [0,2\pi],$ $u\in [0,1],$ 
cuya figura es la siguiente:

\noindent
\begin{minipage}{\textwidth}
	\centering
	\captionof{figure}{Superficie que conecta las curvas $\alpha$ y $\beta$}
	\label{superficie_conecta_alfa_beta}
	\begin{minipage}{0.22\linewidth}
		\centering
		\includegraphics[width=\linewidth]{Fig/48.pdf}
	\end{minipage}
	\hspace{1em}
	\begin{minipage}{0.22\linewidth}
		\centering
		\includegraphics[width=\linewidth]{Fig/49.pdf}
	\end{minipage}
	
	\vspace{0.3em}
	\parbox{0.95\linewidth}{\centering\fontsize{8pt}{10pt}\selectfont\textbf{Fuente:} elaboración propia.}
\end{minipage}

\textcolor{black}{\scriptsize\ding{110}}\quad Las siguientes funciones vectoriales
{\small
$\ve{\alpha}(t)= \langle3\cos(t),2\sin(t),2-\frac{1}{2}\sin(t)\rangle$ y 
$\ve{\beta}(t)= \langle 2\cos(t),\sin(t),\frac{1}{2}\sin ( 2t) \rangle,$ $t\in [0,2\pi]$
} determinan dos curvas en el espacio. 
Con ayuda de las funciones $f(t)=1-t^2$ y $g(t)=t^3$, se puede formar la combinación entre ellas y
definir la función ${\bf r}(u,t)= f(u)\ve{\alpha}(t)+g(u)\ve{\beta}(t)$ con $t \in [0,2\pi],$ $u \in [0,1].$ 
Esta une los puntos de las curvas determinadas por $\ve{\alpha}(t)$
y $\ve{\beta}(t),$ a través de una superficie curva.

\vskip0.2cm
\noindent
\begin{minipage}{\textwidth}
	\centering
	\captionof{figure}{Superficie determinada por ${\bf r}(u,t)$}
	\label{Superficie_determinada}
	\begin{minipage}{0.23\linewidth}
		\centering
		\includegraphics[width=\linewidth]{Fig/50.pdf}
	\end{minipage}
	\hspace{1em}
	\begin{minipage}{0.30\linewidth}
		\centering
		\includegraphics[width=\linewidth]{Img/69a.pdf}
	\end{minipage}
	
	\vspace{0.3em}
	\parbox{0.95\linewidth}{\centering\fontsize{8pt}{10pt}\selectfont\textbf{Fuente:} elaboración propia.}
\end{minipage}
\vskip0.2cm

Esta figura se genera con ayuda de las siguientes instrucciones:

\begin{tcolorbox}[breakable,enhanced,title=\bf Instrucción \verb"Mathematica"] 
\begin{Verbatim}[formatcom=\letraverbatim]
a={3*Cos[t],2*Sin[t],2-Sin[t]/2};
b={2*Cos[t],Sin[t],Sin[2*t]/2};
p=ParametricPlot3D[(1-u^2)*a+u^3*b,{t,0,2*Pi},{u,0,1},
PlotStyle->Orange,MeshStyle->Gray,PlotPoints->60];
q=ParametricPlot3D[{a,b},{t,0,2*Pi},PlotPoints->
60,PlotStyle->{Directive[Black,Thickness->0.012],
Directive[Black,Thickness->0.012]}];
Show[p,q]
\end{Verbatim}
\end{tcolorbox}

\vskip0.2cm
\vspace{0pt}
\textcolor{black}{\scriptsize\ding{110}}\quad Las dos circunferencias unitarias ubicadas en el plano $z=2$ y $z=0$,
se unen por medio de curvas y generan la siguiente superficie.
La parametrizacion de esta región se consigue con la ayuda de las funciones 
$f(t)=\frac{2t}{t^2+1}$ y $g(t)=\frac{1-t}{t^2+1}$.


\vskip0.2cm
\noindent
\begin{minipage}{\textwidth}
	\centering
	\captionof{figure}{Circunferencias conectadas}
	\begin{minipage}{0.32\linewidth}
		\centering
		\includegraphics[width=\linewidth]{Fig/51.pdf}
	\end{minipage}
	
	\vspace{0.3em}
	\parbox{0.95\linewidth}{\centering\fontsize{8pt}{10pt}\selectfont\textbf{Fuente:} elaboración propia.}
\end{minipage}
\vskip0.2cm


Estas funciones se utilizan para formar la siguiente combinación.
{\small
\begin{equation*}
f(u)\ve{\alpha}(t)+g(u)\ve{\beta}(t) 
% \frac{2u}{u^2+1}\la \cos(t),\sin(t),2\ra + \frac{1-u}{u^2+1} \la\cos(t),\sin(t),0\ra\\
= \la  \frac{u+1}{u^2+1}\cos t,\frac{u+1}{u^2+1} \sin t,\frac{4u}{u^2+1}\ra,u \in[0,1], t\in[0,2\pi].
\end{equation*}
}

\vspace{0pt}
\textcolor{black}{\scriptsize\ding{110}}\quad Las funciones {\small $\ve{\alpha}(t) = \langle(2-2\cos(t)) \cos(t), (2-2\cos(t))\sin(t),2\rangle$} y {\small $\ve{\beta}(t) = \langle ( 1-\cos(t)) \cos(t),( 1-\cos(t)) \sin(t),0\rangle$} 
representan dos cardioides.

Se forma la combinación {\small $${\bf r}(u,t)= (1-u^{4})\ve{\alpha}(t) +u^{3}\ve{\beta}(t),$$}
con $u \in[0,1],$ $t \in[0,2\pi],$ que parametriza la superficie curva que las une y se muestra en la figura \ref{Dos_cardioides_conectadas}.


\vskip0.2cm
\noindent
\begin{minipage}{\textwidth}
	\centering
	\captionof{figure}{Dos cardioides conectadas}
	\label{Dos_cardioides_conectadas}
	\begin{minipage}{0.30\linewidth}
		\centering
		\includegraphics[width=\linewidth]{Fig/52.pdf}
	\end{minipage}
	
	\vspace{0.3em}
	\parbox{0.95\linewidth}{\centering\fontsize{8pt}{10pt}\selectfont\textbf{Fuente:} elaboración propia.}
\end{minipage}
\vskip0.2cm


\textcolor{black}{\scriptsize\ding{110}}\quad La figura \ref{Dos_cardioides_conectadas} de la función vectorial {\small ${\bf r}(t)= \langle 4 \cos(t)-\cos(4t),4\sin(t)-\sin(4t),0 \rangle,$} con
$t\in [0,2\pi],$ se presenta a continuación.

\vskip0.2cm
\vspace{0pt}
El interior se parametriza con
$x(u,t) = (4 \cos(t)-\cos(4t))u,$
$y(u,t) = (4\sin(t)-\sin(4t))u,$
para $t\in [0,2\pi],$ $u\in [0,1].$
Utilizando las funciones $f(u)= \sqrt[4]{(1-u)^3}$  y $g(u)=\sqrt[4]{u^3}$,
se forma la combinación con el punto $(1,0,3),$
$${\bf r}_2(u,t)=f(u) {\bf r}(t)+ g(u)(1,0,3).$$


\vskip0.2cm
\noindent
\begin{minipage}{\textwidth}
	\centering
	\captionof{figure}{${\bf r}_1(u,t)=\la x(u,t),y(u,t),0 \ra$}
	\label{fig:53}
	\begin{minipage}{0.23\linewidth}
		\centering
		\includegraphics[width=\linewidth]{Fig/53.pdf}
	\end{minipage}
	
	\vspace{0.3em}
	\parbox{0.95\linewidth}{\centering\fontsize{8pt}{10pt}\selectfont\textbf{Fuente:} elaboración propia.}
\end{minipage}
\vskip0.2cm


Lo anterior permite definir la función
{\small
\[
{\bf r}_2(t,u)= \sqrt[4]{(1-u)^3} \left( 4 \cos(t)-\cos(4t),4\sin(t)-\sin(4t),0\right)+ \sqrt[4]{u^3} \left(1,0,3\right).
\]
}

Las figuras de las superficies generadas por las funciones 
${\bf r}_1(t,u)$ y ${\bf r}_2(t,u)$ se muestran con diferentes colores en la siguiente figura. 

\begin{figure}[H]
	\begin{center}
		\captionof{figure}{\small Superficies determinadas por ${\bf r}_1(u,t)$ y ${\bf r}_2(u,t)$.}
		\includegraphics[height=4.5cm]{Img/78c.pdf}
		\includegraphics[height=4.5cm]{Img/78a.pdf}
		\includegraphics[height=4.5cm]{Img/78b.pdf}
	\end{center}
	\parbox{0.95\linewidth}{\centering\fontsize{8pt}{10pt}\selectfont{\bf Fuente:} elaboración propia.}
\end{figure}

Las siguientes instrucciones permiten obtener estos gráficos.

\begin{tcolorbox}[breakable,enhanced,title=\bf Instrucción \verb"Mathematica"]
\begin{Verbatim}[formatcom=\letraverbatim]
p={4*Cos[t]-Cos[4*t],4*Sin[t]-Sin[4*t],0};
r=(1-u)^(3/4)*p+u^(2/4)*{1,1,3};
ParametricPlot3D[{r,p*u},{u,0,1},{t,0,2Pi},
PlotStyle->{Orange,Cyan}]
\end{Verbatim}
\end{tcolorbox}

La función generada es ${\bf r}(u,t)= \langle(2\sqrt{u}+1-u^2)\cos t, (2\sqrt{u}+1-u^2) \sin t,\sqrt{u}\rangle,$
$ t\in [0,2\pi],\, u\in[0,1].$ 

\vskip0.2cm
\vspace{0pt}
\textcolor{black}{\scriptsize\ding{110}}\quad Las circunferencias parametrizadas por
$\ve{\alpha}(t)= \langle \cos(t),\sin(t),0 \rangle $ y $\ve{\beta}(t)= \langle 2\cos(t),2\sin(t),1 \rangle $,
se unen utilizando las funciones $g(u)= \sqrt{u}$ y $f(u)=1-u^2$; mediante la combinación
${\bf r}(u,t)= f(u)\ve{\alpha}(t) +g(u)\ve{\beta}(t)= \langle (2-2u^2+\sqrt{u})\cos t,(2-2u^2+\sqrt{u})\sin t,1- u^2\rangle$
con $t\in [0,2\pi],\,u\in[0,1].$


\vskip0.2cm
\noindent
\begin{minipage}{\textwidth}
	\centering
	\captionof{figure}{Circunferencias conectadas}
	\label{Circunferencias_conectadas}
	\begin{minipage}{0.32\linewidth}
		\centering
		\includegraphics[width=\linewidth]{Fig/54.pdf}
	\end{minipage}
	
	\vspace{0.3em}
	\parbox{0.95\linewidth}{\centering\fontsize{8pt}{10pt}\selectfont\textbf{Fuente:} elaboración propia.}
\end{minipage}
\vskip0.2cm


\textcolor{black}{\scriptsize\ding{110}}\quad La curva parametrizada por las ecuaciones $x(t)= 16\sin ^{3}(t),$ 
$y(t)=13\cos (t) -5\cos \left( 2t\right) -2\cos \left( 3t\right) -\cos \left(4t\right),$
con $t\in \lbrack 0,2\pi ]$, gira alrededor del eje vertical y genera una superficie.
Al hacer coincidir el eje vertical con el eje $z,$ entonces la función que describe la superficie en mención, para $t\in \lbrack 0,2\pi ],$ $u\in \lbrack 0,1],$ es
\[
\resizebox{\textwidth}{!}{$
	\mathbf{r}(t,u) = \langle 16\sin^3(t)\cos u,\ 16\sin^3(t)\sin u,\ 13\cos(t) - 5\cos(2t) - 2\cos(3t) - \cos(4t) \rangle
	$}
\]

\vspace{-4mm}
\begin{figure}[H]
	\centering
	% 1. Reducimos el espacio que el caption deja debajo de sí mismo
	\captionsetup{skip=2pt} 
	\caption{Superficies parametrizadas por $a(t, u)$ y $b(t, u)$}
	\label{fig:superficies}
	
	% 2. Eliminamos los vspace negativos entre imágenes si están en la misma línea
	% y usamos un vspace positivo pequeño para separarlas del caption
	\vspace{2pt} 
	\includegraphics[width=0.21\linewidth]{Img/49a.pdf} \hspace{1em}
	\includegraphics[width=0.21\linewidth]{Img/49b.pdf} \hspace{1em}
	\includegraphics[width=0.21\linewidth]{Fig/55.pdf}
	
	% 3. Controlamos el espacio de la fuente de manera independiente
	\vspace{0.2cm} 
	\parbox{0.95\linewidth}{\centering\fontsize{8pt}{10pt}\selectfont\textbf{Fuente:} elaboración propia.}
\end{figure}


Esta figura se genera ejecutando las siguientes instrucciones:

\begin{tcolorbox}[breakable,enhanced,title=\bf Instrucción \verb"Mathematica"] 
\begin{Verbatim}[formatcom=\letraverbatim]
ParametricPlot[{16*Sin[t]^3,13*Cos[t]-5*Cos[2*t]
-2*Cos[3*t]-Cos[4*t]},{t,0,2Pi}];
ParametricPlot3D[{16*Sin[t]^3*Sin[u],(16*Sin[t]^3)
*Cos[u],13*Cos[t]-5*Cos[2*t]-2*Cos[3*t]-Cos[4*t]},
{t,0,2Pi},{u,0,2*Pi},PlotPoints->80]
\end{Verbatim}
\end{tcolorbox}

\vspace{-0.2em}
Presentamos enseguida otros ejemplos de superficies parametrizadas bajo ciertas condiciones.
\subsection*{Superficie acotada por dos paraboloides} \index{Paraboloide} 
Las ecuaciones $8y=9-5x^{2}-6z^{2}$ y $16y=4z^{2}-x^{2},$
representan gráficamente un paraboloide y un paraboloide hiperbólico, al 
eliminar $y$ de estas expresiones se obtiene la ecuación $9x^{2}+16z^{2}=18,$ que corresponde a una elipse en el plano $xz.$ La región interior a esta curva se parametriza como 
$x=\sqrt{2}u\cos (t) ,$ $z=\frac{3\sqrt{2}}{4}u\sin (t) ,$ para $t\in \lbrack 0,2\pi ],$ $u\in \lbrack 0,1].$ Al reemplazar estas expresiones en las ecuaciones del paraboloide y del paraboloide hiperbólico, se obtiene:
\[
8y=9-5x^{2}-6z^{2}= 9-\frac{13}{4}u^{2}\cos ^{2}(t) -\frac{27}{4}u^{2},
\]
entonces
\[
y= \frac{9}{8}-\frac{13}{32}u^{2}\cos ^{2}(t) -\frac{27}{32}u^{2}.
\]
Por otra parte,
\[
16y=4z^{2}-x^{2}= \frac{9}{2}u^{2}-\frac{13}{2}u^{2}\cos^{2}(t) ,
\]
luego
\[
y= \frac{9}{32}u^{2}-\frac{13}{32}u^{2}\cos ^{2}(t) .
\]
Las superficies señaladas y el sólido  en mención se muestran enseguida.

\begin{figure}[H]
	\centering
	% 1. Centramos la leyenda respecto a todo el ancho de la página
	\caption{Superficies parametrizadas por $a(t,u)$ y $b(t,u)$}
	\label{fig:superficies_parametrizadas}
	\vspace{2pt} % Espacio pequeño para acercar la leyenda a las figuras
	
	% 2. Reducimos el tamaño de las minipages (de 0.32 a 0.28) para disminuir las figuras
	\begin{minipage}[b]{0.28\linewidth}
		\centering
		% Usamos scale o una altura menor (ej. 3.5cm) para homogeneizar
		\includegraphics[width=\linewidth, height=3.0cm, keepaspectratio]{Img/71b.pdf}
	\end{minipage}
	\hfill
	\begin{minipage}[b]{0.28\linewidth}
		\centering
		\includegraphics[width=\linewidth, height=3.0cm, keepaspectratio]{Img/71d.pdf}
	\end{minipage}
	\hfill
	\begin{minipage}[b]{0.28\linewidth}
		\centering
		\includegraphics[width=\linewidth, height=3.5cm, keepaspectratio]{Fig/56.pdf}
	\end{minipage}
	
	% 3. Espacio para la fuente
	\vspace{0.4cm}
	\parbox{0.95\linewidth}{\centering\fontsize{8pt}{10pt}\selectfont\textbf{Fuente:} elaboración propia.}
\end{figure}

\vskip0.3cm
Con lo señalado antes, la superficie sobre el paraboloide acotada por el
paraboloide hiperbólico se parametriza como
{\small
\[
\mathbf{a}(t,u) = \left\langle \sqrt{2}u\cos(t),\,
\frac{9}{8} - \frac{13}{32}u^{2}\cos^{2}(t) - \frac{27}{32}u^{2},\,
\frac{3\sqrt{2}}{4}u\sin(t) \right\rangle,
\]
}
mientras que la superficie sobre el paraboloide hiperbólico acotada por el paraboloide está dada por
{\small
\[
\mathbf{b}(t,u) = \left\langle \sqrt{2}u\cos(t),\,
\frac{9}{32}u^{2} - \frac{13}{32}u^{2}\cos^{2}(t),\,
\frac{3\sqrt{2}}{4}u\sin(t) \right\rangle.
\]
}
En ambos casos, $t \in [0, 2\pi]$ y $u \in [0, 1]$. 
Los gráficos se obtienen al ejecutar las siguientes instrucciones. Al ingresar el código se usó la identidad $2\cos^2(t) = 1 + \cos(2t)$.

\begin{tcolorbox}[breakable,enhanced,title=\bf Instrucción \verb"Mathematica"] 
\begin{Verbatim}[formatcom=\letraverbatim]
ParametricPlot3D[{{Sqrt[2]*u*Cos[t],(72-67*u^2
-13*u^2*Cos[2*t])/64,3*Sqrt[2]/4*u*Sin[t]},
{Sqrt[2]*u*Cos[t],u^2*(5-13*Cos[2*t])/64,
3*Sqrt[2]/4*u*Sin[t]}},{t,0,2*Pi},{u,0,1},
PlotStyle->{Orange,Cyan},PlotPoints->80]
\end{Verbatim}
\end{tcolorbox}

\subsection*{Superficies creadas por una cardioide y una elipse}
La cardioide con ecuación $r=1-\cos(t)$, se puede parametrizar,  para $t\in [0,2\pi]$ como 
$\langle(1-\cos(t))\cos(t),(1-\cos(t))\sin(t)\rangle$ y la elipse parametri- \\ zada en la forma $\langle 3\cos(t),2\sin(t)\rangle,$ $t\in [0,2\pi]$, determinan una región en el plano acotada por sus gráficos.

\vskip0.2cm
\noindent
\begin{minipage}{\textwidth}
	\centering
	\captionof{figure}{Superficie sobre el cilindro parabólico $z=1+(1-y)^2/4$}
	\label{Superficie_cilindro_parabólico}
	\begin{minipage}{0.37\linewidth}
		\centering
		\vspace{-3mm}
		\includegraphics[width=\linewidth]{Fig/57.pdf}
	\end{minipage}
	\hspace{1em}
	\begin{minipage}{0.37\linewidth}
		\centering
		\vspace{-3mm}
		\includegraphics[width=\linewidth]{Fig/58.pdf}
	\end{minipage}
	
	\vspace{0.3em}
	\parbox{0.95\linewidth}{\centering\fontsize{8pt}{10pt}\selectfont\textbf{Fuente:} elaboración propia.}
\end{minipage}
\vskip0.2cm

Una parametrización de la región del plano entre la elipse y el cardioide, se determina
formando una combinación convexa entre dos puntos, uno en cada curva, es decir
$(1-u)\langle(1-\cos(t))\cos(t),(1-\cos(t))\sin(t)\rangle
+u\langle 3\cos(t),2\sin(t)\rangle,$
esta expresión se simplifica en la forma 
{\small $${\bf r}(u,t)=\big \langle((u-1)\cos(t)+1+2u)\cos(t), (1+u+(u-1)\cos(t))\sin(t)\big \rangle,$$}
con $t\in [0,2\pi],\, u\in [0,1]$.
Esta región la proyectaremos sobre diferentes superficies para obtener su respectiva
parametrización. 
\vskip0.3cm
\subsubsection*{Superficie sobre el cilindro con ecuación \texorpdfstring{$z=1+(1-y)^2/4$}{z=1+(1-y)²/4}}

\vskip0.3cm
La ecuación $z=1+(1-y)^2/4$ nos permite llevar esta región del plano al espacio. 
Con lo mostrado anteriormente, solo nos falta determinar la coordenada $z$, esto es
$z=1+\frac{1}{4}(1-y)^2=1+\frac{1}{4}(1-((1+u+(u-1)\cos(t))\sin(t)))^2,$
con lo cual, la función que representa la parte del cilindro sobre la región que se considera es:
{\small
\begin{multline*}
{\bf r}(u,t)=\big \langle((u-1)\cos(t)+1+2u)\cos(t), (1+u+(u-1)\cos(t))\sin(t),\\
1+\frac{1}{4}(1-((1+u+(u-1)\cos(t))\sin(t)))^2 \big \rangle, 
\end{multline*}
}
con $t\in [0,2\pi],\, u\in [0,1].$ Una manera distinta de ver la misma superficie es considerarla como la parte del 
cilindro con ecuación $z=1+(1-y)^2/4$, que se encuentra entre el cilindro cardioidico
\index{Cilindro cardiodico}
con ecuación polar $r=1-\cos(t)$ y el cilindro elíptico con ecuación 
$\frac{x^2}{9}+\frac{y^2}{4}=1.$ La figura de esta situación se muestra a continuación.
\vspace{-1mm}
\vskip0.2cm
\noindent
\begin{minipage}{\textwidth}
	\centering
	\captionof{figure}{\small Superficie sobre el cilindro parabólico $z=1+(1-y)^2/4$}
	\begin{minipage}{0.24\linewidth}
		\centering
		\includegraphics[width=\linewidth]{Img/43b.pdf}
	\end{minipage}
	\hspace{1em}
	\begin{minipage}{0.24\linewidth}
		\centering
		\includegraphics[width=\linewidth]{Img/43a.pdf}
	\end{minipage}
	
	\vspace{0.3em}
	\parbox{0.95\linewidth}{\centering\fontsize{8pt}{10pt}\selectfont\textbf{Fuente:} elaboración propia.}
\end{minipage}

\vskip0.2cm
Estos gráficos se elaboraron en \verb"Mathematica".
En el plano $xy$ se forma una combinación convexa entre la función que determina la elipse y la cardioide.
Esta región se proyecta sobre el paraboloide $z=1+(y-1)^2/4$, la cual se grafica en el espacio.
Los cilindros se grafican por separado, definiendolos como \verb"a" y \verb"b",
que se muestran en un solo gráfico.

\begin{tcolorbox}[breakable,enhanced,title=\bf Instrucción \verb"Mathematica"] 
\begin{Verbatim}[formatcom=\letraverbatim]
{x1,y1}={(1-Cos{t})},*Cos{t},(1-Cos{t})*Sin{t});
{x2,y2}={3*Cos[t],2*Sin[t]};
{x,y}={1-u}*{x1,y1}+u*{x2,y2};
ParametricPlot[{x,y},{t,0,2*Pi},{u,0,1}]
ParametricPlot3D[{x,y,1+(1-y)^2/4},{u,0,1}{t,0,2*Pi},
PlotRange->{{-3, 3}, {-2, 2}},
MeshStyle->Gray]
a=ContourPlot3D[{z==1+(1-y)^2/4},
{x,-4,4},{y,-4,4},{z,-1,5},
ContourStyle->RGBColor[0.1,0.4,0.3], Mesh->False];
b=ParametricPlot3D[{{x1,y1,u},{x2,y2, u}},{u,0,5},
{t,0,2*Pi}, 
Mesh->False];
Show[a, b];
	\end{Verbatim}
\end{tcolorbox}

\subsubsection*{Superficie sobre el cilindro \texorpdfstring{$z=\cos (2x)$}{z=cos(2x)}}
En este caso, la parametrización de la superficie es
{\small
\begin{multline*}
{\bf r}\left( u,t\right) = \big\langle \cos (t) \left( \cos (t)
\left( u-1\right) +1+2u\right) ,\sin (t) ( 1+u+\\ 
(u-1) \cos (t)) , \cos \big( 2\cos (t)\left( \cos (t) \left( u-1\right) +1+2u\right) \big) \big \rangle,
\end{multline*}
}
con $t \in [0,2\pi]$, $u \in [0,1]$, dicha superficie se muestra a continuación.
\vskip0.2cm
\noindent
\begin{minipage}{\textwidth}
	\centering
	\captionof{figure}{\small Superficie sobre el cilindro $z=\cos (2x)$}
	\begin{minipage}{0.24\linewidth}
		\centering
		\includegraphics[width=\linewidth]{Img/46a.pdf}
	\end{minipage}
	\hspace{1em}
	\begin{minipage}{0.24\linewidth}
		\centering
		\includegraphics[width=\linewidth]{Img/46b.pdf}
	\end{minipage}
	
	\vspace{0.3em}
	\parbox{0.95\linewidth}{\centering\fontsize{8pt}{10pt}\selectfont\textbf{Nota:} elaboración propia.}
\end{minipage}

\vspace{0.3cm}
Las instrucciones para conseguirlo son las siguientes:
primero, se define la función ${\bf r}$, que describe la región entre la elipse y la cardioide;
luego, se define la función que parametriza la superficie y 
se muestra su gráfico. 

\begin{tcolorbox}[breakable,enhanced,title=\bf Instrucción \verb"Mathematica"] 
\begin{Verbatim}[formatcom=\letraverbatim]
r=(1-u)*{(1-Cos[t])*Cos[t],(1-Cos[t])*Sin[t]}+
u*{3*Cos[t],2*Sin[t]};
ParametricPlot3D[{r[[1]],r[[2]],Cos[2*r[[1]]]},
{u,0,1},{t,0,2*Pi},PlotRange->{{-3,3},{-3,3},{-1,1}},
PlotStyle->RGBColor[0.8,0.2,0.3],Mesh->False,
PlotPoints->90,Ticks->{{-2,0,2},{-2,0,2},{-1,0,1}}]
\end{Verbatim}
\end{tcolorbox}

\subsubsection*{Superficie sobre el cilindro \texorpdfstring{$z=\sin (2x)$}{z=sin(2x)}}
Sobre este cilindro la parametrización que se utiliza es:
{\small
\begin{multline*}
{\bf r}(u,t) = \big\langle \cos (t) \left( \cos (t)\left( u-1\right) +1+2u\right),
\sin(t)(1+u+\\ (u-1)\cos(t)), \sin \big( 2\cos (t)\left( \cos (t) \left( u-1\right) +1+2u\right)\big) \big\rangle,
\end{multline*}
}
en donde $t \in [0,2\pi]$, $u \in [0,1]$.
Se presenta ahora la figura y las instrucciones
que se requieren para obtenerlo:


\vskip0.2cm
\noindent
\begin{minipage}{\textwidth}
	\centering
	\captionof{figure}{\small Superficie sobre el cilindro $z=\sin(2x)$}
	\begin{minipage}{0.23\linewidth}
		\centering
		\includegraphics[width=\linewidth]{Img/47b.pdf}
	\end{minipage}
	\begin{minipage}{0.23\linewidth}
		\centering
		\includegraphics[width=\linewidth]{Img/47a.pdf}
	\end{minipage}
	
	\vspace{0.3em}
	\parbox{0.95\linewidth}{\centering\fontsize{8pt}{10pt}\selectfont\textbf{Fuente:} elaboración propia.}
\end{minipage}
\vskip0.2cm



\begin{tcolorbox}[breakable,enhanced,title=\bf Instrucción \verb"Mathematica"]  
\begin{Verbatim}[formatcom=\letraverbatim]
r=(1-u)*{(1-Cos[t])*Cos[t],(1-Cos[t])*Sin[t]}+
u*{3*Cos[t],2*Sin[t]};
p={r[[1]],r[[2]],Sin[2*r[[1]]]};
ParametricPlot3D[p,{u,0,1},{t,0,2*Pi},
PlotRange->{{-3,3},{-3,3},{-1,1}},Mesh->False,
PlotStyle->RGBColor[0.8,0.2,0.3],PlotPoints->80]
\end{Verbatim}
\end{tcolorbox}

\subsubsection*{Superficie sobre el cilindro \texorpdfstring{$z=\frac{1}{5}y^{2}-\frac{1}{4}(x-1) ^{2}$}{z=(1/5)y²-(1/4)(x-1)²}}

La superficie mencionada se presenta a continuación.

\vskip0.2cm
\noindent
\begin{minipage}{\textwidth}
	\centering
	\captionof{figure}{\small Superficie sobre el cilindro $z=\frac{1}{5}y^{2}-\frac{1}{4}(x-1) ^{2}$}
	\begin{minipage}{0.24\linewidth}
		\centering
		\includegraphics[width=\linewidth]{Img/45a.pdf}
	\end{minipage}
	\begin{minipage}{0.24\linewidth}
		\centering
		\includegraphics[width=\linewidth]{Img/45b.pdf}
	\end{minipage}
	
	\vspace{0.3em}
	\parbox{0.95\linewidth}{\centering\fontsize{8pt}{10pt}\selectfont\textbf{Fuente:} elaboración propia.}
\end{minipage}
\vskip0.2cm


Como se señaló antes, nos falta determinar la coordenada $z$, esto es 
{\footnotesize
	\begin{align*}
z&=\frac{1}{5}y^{2}-\frac{1}{4}\left( x-1\right) ^{2} \\	&=\frac{1}{5}\big( \sin(t) \left( 1+u+(u-1) \cos(t) \right) \big)^{2}-\frac{1}{4}\big(\cos(t)\left(\cos(t) \left( u-1\right) +1+2u\right)-1\big)^{2},
	\end{align*}
}
entonces la parametrización de la superficie es
{\small
\begin{multline*}
{\bf r}(u,t) = \bigg\langle \cos (t) \left( \cos (t)
\left( u-1\right) +1+2u\right) ,\sin (t)(1+u
+(u-1) \cos (t)) , \\ \frac{1}{5}\big( \sin(t) \left( 1+u+(u-1) \cos(t) \right)\big) ^{2} \\
-\frac{1}{4}\big( \cos (t) \left( \cos (t) \left( u-1\right) +1+2u\right) -1\big) ^{2} \bigg \rangle. 
\end{multline*}
}
Para $t \in [0,2\pi]$, $u \in [0,1]$, las instrucciones 
para conseguir el gráfico se muestran enseguida.

%\vspace{-4mm}
\begin{tcolorbox}[breakable,enhanced,title=\bf Instrucción \verb"Mathematica"] 
\begin{Verbatim}[formatcom=\letraverbatim]
r=(1-u)*{(1-Cos[t])*Cos[t],(1-Cos[t])*Sin[t]}+
u*{3*Cos[t],2*Sin[t]};
p={r[[1]],r[[2]],-1/4*(r[[1]]-1)^2+1/5*(r[[2]])^2};
ParametricPlot3D[p,{u,0,1},{t,0,2*Pi},
PlotRange->{{-3,3},{-2,2},{-4.2,1}},Mesh->False,
PlotStyle->RGBColor[0.2,0.5,0.5],PlotPoints->80]
\end{Verbatim}
\end{tcolorbox}

\subsection*{Superficies acotadas por cilindros elípticos}

Las expresiones $4x^2 + y^2 = 36$ y $x^2 + 4y^2 = 4$ representan dos cilindros elípticos con eje $z$ como eje de simetría. Sin embargo, se puede dar una mejor idea de sus ecuaciones representando dos elipses que pueden parametrizarse respectivamente como $\vec{a}(t) = \langle 3\cos(t),\, 6\sin(t) \rangle$ y $\vec{b}(t) = \langle 2\cos(t),\, \sin(t) \rangle$.




\vskip0.2cm
\noindent
\begin{minipage}{\textwidth}
	\centering
	\captionof{figure}{Región parametrizada por ${\bf r}(u,t)=\left\langle \cos(t)(3-u),\, \sin(t)(6-5u) \right\rangle$}
	\label{fig:59}
	\begin{minipage}{0.23\linewidth}
		\centering
		\includegraphics[width=\linewidth]{Img/1a.pdf}
	\end{minipage}
	\begin{minipage}{0.28\linewidth}
		\centering
		\includegraphics[width=\linewidth]{Fig/59.pdf}
	\end{minipage}
	\begin{minipage}{0.28\linewidth}
		\centering
		\includegraphics[width=\linewidth]{Fig/60.pdf}
	\end{minipage}
	
	\vspace{0.3em}
	\parbox{0.95\linewidth}{\centering\fontsize{8pt}{10pt}\selectfont\textbf{Fuente:} elaboración propia.}
\end{minipage}
\vskip0.2cm


Con dos puntos $P$ y $Q$, simplificando la expresión: 
\begin{align}
	(1 - u)P + uQ &= (1 - u)(3\cos(t),\, 6\sin(t)) + u(2\cos(t),\, \sin(t)) \notag \\
	&= (\cos(t)(3 - u),\, \sin(t)(6 - 5u)). \tag{3.7} \label{1a}
\end{align}
Para $u \in [0,1]$, esta función parametriza la región que se halla entre estas dos curvas. La región en mención, que denominaremos $R$, se proyectará sobre un plano y posteriormente sobre un paraboloide; se presentará su construcción con sus respectivas parametrizaciones.

\subsubsection*{Superficie sobre un plano}
Se puede representar la porción del plano con ecuación $6z=2x+y+12$ sobre la región acotada  por los dos cilindros. 

\begin{figure}[H]
	\centering
	\captionof{figure}{\small Cilindros elípticos intersecados por el plano $6z=2x+y+12$}
	\label{fig:cilindros_intersecados}
	\vspace{2pt} % Acerca la leyenda a las imágenes
	
	% Reducción uniforme: width de 4cm a 3.5cm o 0.28\linewidth
	\includegraphics[width=0.28\linewidth, height=3cm, keepaspectratio]{Img/1c.pdf}
	\hspace{1cm} % Espacio horizontal entre ambas
	\includegraphics[width=0.28\linewidth, height=3cm, keepaspectratio]{Img/1d.pdf}
	
	\vspace{0.4cm} % Separación con la fuente
	\parbox{0.95\linewidth}{\centering\fontsize{8pt}{10pt}\selectfont\textbf{Fuente:} elaboración propia.}
\end{figure}

Para determinar $z$ se puede proceder como sigue
\vspace*{-4mm} 
\[
\begin{aligned}
	z &= \frac{1}{6}\left( 2\cos (t) \left( 3-u\right) +\sin (t) \left( 6-5u\right) +12\right) \\ 
	&= \left( 1-\frac{1}{3}u\right)\cos (t) +\left( 1-\frac{5}{6}u\right) \sin (t) +2
\end{aligned}
\]
Por último, la región $R$ sobre el plano se parametriza mediante la función:
\vspace*{-2mm} 
{\footnotesize
	\[
	\overrightarrow{r}(u,t) =\big\langle \cos (t) \left( 3-u\right) ,\sin (t) \left( 6-5u\right) ,
	\left( 1 - \! \frac{1}{3}u\right) \cos (t) + \left( 1 - \! \frac{5}{6}u\right) \sin (t) +2\big\rangle ,
	\]
}
\vspace*{-1mm} 
con $t\in \lbrack 0,2\pi ],\,\, u\in \lbrack 0,1].$


\vskip0.2cm
\noindent
\begin{minipage}{\textwidth}
	\centering
	\captionof{figure}{Región sobre el plano $6z=2x+y+12$ acotada por los cilindros elípticos}
	\label{Región_cilindros_elípticos}
	\begin{minipage}{0.50\linewidth}
		\centering
		\includegraphics[width=\linewidth]{Fig/61.pdf}
	\end{minipage}
	
	\vspace{0.3em}
	\parbox{0.95\linewidth}{\centering\fontsize{8pt}{10pt}\selectfont\textbf{Fuente:} elaboración propia.}
\end{minipage}
\vskip0.2cm



\subsubsection*{Superficie sobre un paraboloide} \index{Paraboloide} 
Con ayuda de la expresión (\ref{1a}), se puede parametrizar la parte del
parabo- \\ loide con ecuación $z=6-\left( x-1\right) ^2-y^2,$ cuya proyección 
en el plano $xy$ es la región $R,$ para este fin determinaremos $z$, esto es
\begin{eqnarray*} z&=&  6-\left( \left( 3-u\right)\cos(t)-1\right)^2-\left(
\left( 6-5u\right) \sin (t) \right) ^2 \\
&=& -31+\left( 27-54u+24u^{2}\right) \cos ^{2}(t) +\left( 6-2u\right) \cos (t) +60u-25u^{2}
\end{eqnarray*}


\vskip0.2cm
\noindent
\begin{minipage}{\textwidth}
	\centering
	\captionof{figure}{Superficie sobre el paraboloide acotada por los cilindros elípticos}
	\label{Superficie_paraboloide_acotada_cilindros_elípticos}
	\begin{minipage}{0.23\linewidth}
		\centering
		\includegraphics[width=\linewidth]{Img/1e.pdf}
	\end{minipage}
	\begin{minipage}{0.23\linewidth}
		\centering
		\includegraphics[width=\linewidth]{Img/1f.pdf}
	\end{minipage}
	\begin{minipage}{0.28\linewidth}
		\centering
		\includegraphics[width=\linewidth]{Fig/62.pdf}
	\end{minipage}
	
	\vspace{0.3em}
	\parbox{0.95\linewidth}{\centering\fontsize{8pt}{10pt}\selectfont\textbf{Fuente:} elaboración propia.}
\end{minipage}
\vskip0.2cm


Por lo tanto, la función 
\vspace*{-4mm} % Ajuste vertical después del párrafo
\begin{multline*}
	\overrightarrow{r}(u,t) = \big\langle \cos(t)(3-u) ,\sin(t)( 6-5u) , \\
	-31+\left( 27-54u+24u^{2}\right) \cos ^{2}(t) +\left( 6-2u\right) \cos (t) +60u-25u^{2} \big\rangle
\end{multline*}

\noindent con $t\in \lbrack 0,2\pi ],$ $u\in \lbrack 0,1],$
con $ t\in \lbrack 0,2\pi ],$ $u\in \lbrack 0,1],$  permite representar gráficamente la situación mencionada.
% Las curvas de intersección de los cilindros elípticos y el paraboloide se determinan reemplazando en el valor 
% de $z,$ que para los dos casos es:
% $$z=6-\left( 3\cos (t) -1\right) ^{2}-\left( 6\sin (t) \right) ^{2}= -31+27\cos ^{2}(t) +6\cos (t) $$
% $$z=6-\left( 2\cos (t) -1\right) ^{2}-\left( \sin (t) \right) ^{2}= 4-3\cos ^{2}(t) +4\cos (t) $$
Las figuras de esta sección pueden obtenerse ejecutando las siguientes instrucciones:

\begin{tcolorbox}[breakable,enhanced,title=\bf Instrucción \verb"Mathematica"] 
\begin{Verbatim}[formatcom=\letraverbatim]
ContourPlot3D[{4*x^2+y^2==36, x^2+4*y^2==4,6*z==-2*x+y+12},
{x,-6,6},{y,-6,6},{z,0,4},PlotPoints->50]
ContourPlot3D[{6*z==-2*x+y+12, 4*x^2+y^2==36,x^2+4*y^2==4},
{x,-4,4},{y,-7,7},{z,-0.5,4.5},PlotPoints->80,
ContourStyle->{Gray,Yellow,Red}]
ContourPlot3D[{z==6-(x-1)^2-y^2,4*x^2+y^2==36,
x^2+4*y^2==4},
{x,-6,10},{y,-7,7},{z,-35,6},
ContourStyle->{Red,Yellow,Green}]
\end{Verbatim}
\end{tcolorbox}


\begin{tcolorbox}[breakable,enhanced,title=\bf Instrucción \verb"Mathematica"] 
\begin{Verbatim}[formatcom=\letraverbatim]
ParametricPlot3D[{{1 + 3*Cos[t],2*Sin[t],Sqrt[5*(Sin[t])^2 
-6*Cos[t]+6]},{1+3*Cos[t],2*Sin[t],-Sqrt[5*(Sin[t])^2
-6*Cos[t]+6]}},{t,0,2*Pi}, PlotStyle->{Directive[Black,
Thickness[0.015]],Directive[Black,Thickness[0.015]]}]
\end{Verbatim}
\end{tcolorbox}

\subsection*{Dos paraboloides} \index{Paraboloide!hiperbólico}
Las ecuaciones $4y=z^{2}-x^{2}\,$ y $\, y=4-x^{2}-z^{2}$ 
representan un paraboloide y un paraboloide hiperbólico. Estas figuras se muestran enseguida.

\vskip0.2cm
\noindent
\begin{minipage}{\textwidth}
	\centering
	\captionof{figure}{\small $4y=z^{2}-x^{2}$ y $y=4-x^{2}-z^{2}$}
	\begin{minipage}{0.24\linewidth}
		\centering
		\includegraphics[width=\linewidth]{Img/4a.pdf}
	\end{minipage}
	\hspace{1em} 
	\begin{minipage}{0.24\linewidth}
		\centering
		\includegraphics[width=\linewidth]{Img/4b.pdf}
	\end{minipage}
	
	\vspace{0.3em}
	\parbox{0.95\linewidth}{\centering\fontsize{8pt}{10pt}\selectfont\textbf{Fuente:} elaboración propia.}
\end{minipage}
\vskip0.2cm

Estas superficies se intersecan 
en una curva con ecuación cartesiana 
$3x^{2}+5z^{2}=16,$ que equivale a 
$\frac{x^{2}}{16/3}+\frac{z^{2}}{16/5}=1,$ la cual representa una elipse en
el plano $xz.$ 


\vskip0.2cm
\noindent
\begin{minipage}{\textwidth}
	\centering
	\captionof{figure}{$3x^{2}+5z^{2}\leq 16$}
	\label{fig:63}
	\begin{minipage}{0.34\linewidth}
		\centering
		\vspace{-4mm}
		\includegraphics[width=\linewidth]{Fig/63.pdf}
	\end{minipage}
	
	\vspace{4mm}
	\parbox{0.95\linewidth}{\centering\fontsize{8pt}{10pt}\selectfont\textbf{Fuente:} elaboración propia.}
\end{minipage}
\vskip0.2cm




La frontera puede parametrizarse en la forma
$x=\frac{4}{\sqrt{3}}\cos (t) ,\,\, z=\frac{4}{\sqrt{5}}\sin(t),$  $t\in \lbrack 0,2\pi ].$
Mientras que el interior se parametriza en la forma  
$x=\frac{4u}{\sqrt{3}}\cos (t) ,$ $z=\frac{4u}{\sqrt{5}}\sin(t),$ 
$t\in \lbrack 0,2\pi ]$ y $u\in \lbrack 0,1].$ Reemplazando $x$ y $z$ en la ecuación del paraboloide se sigue que
{\footnotesize
	\[
	y = 4-x^{2}-z^{2}= 4-\left( \frac{4}{\sqrt{3}}\cos \left( t\right) \right) ^{2}-\left( \frac{4}{\sqrt{5}}\sin \left( t\right) \right) ^{2} =  \frac{4}{5}-\frac{32}{15}\cos ^{2}\left( t\right)
	\]
}
Con lo anterior, se establece que la función vectorial:
\[
\overrightarrow{r}(t)=\left\langle \frac{4}{\sqrt{3}}\cos (t) ,\frac{4}{5}-\frac{32}{15}\cos ^{2}(t) ,\frac{4}{\sqrt{5}}\sin (t) \right\rangle,
\]
\noindent con $t\in \lbrack 0,2\pi ],$ representa la curva $C$, que es la intersección.



\noindent \textcolor{black}{\scriptsize\ding{110}}\quad La superficie sobre el paraboloide hiperbólico acotado por la curva $C$, se puede 
parametrizar reemplazando $x$ y $z$, para obtener el valor de $y$, 
\[
y=\frac{1}{4}\left(\left(\frac{4u}{\sqrt{5}}\sin (t) \right) ^{2}-\left( \frac{4u}{\sqrt{3}}\cos (t) \right) ^{2}\right) = \frac{4}{5}u^{2}-\frac{32}{15}u^{2}\cos ^{2}(t),
\]
	
por lo tanto, 
\[
\overrightarrow{r}(t,u)=\left\langle \frac{4u}{\sqrt{3}}\cos (t) , \frac{4}{5}u^{2}-\frac{32}{15}u^{2}\cos ^{2}(t) , \frac{4u}{\sqrt{5}}\sin (t)  \right\rangle,
\]
con $t\in [0,2\pi ],$ $u\in [ 0,1 ], $ es una manera de parametrizar la mencionada superficie.
	

%\vskip0.2cm
\noindent
\begin{minipage}{\textwidth}
	\centering
	\captionof{figure}{Superficie determinada por $\mathbf{r}(t,u)$}
	\label{fig:64}
	\begin{minipage}{0.32\linewidth}
		\centering
		\vspace{-3mm}
		\includegraphics[width=\linewidth]{Fig/64.pdf}
	\end{minipage}
	
	\vspace{0.3em}
	\parbox{0.95\linewidth}{\centering\fontsize{8pt}{10pt}\selectfont\textbf{Fuente:} elaboración propia.}
\end{minipage}
\vskip0.2cm

\noindent
\textcolor{black}{\scriptsize\ding{110}}\quad Ahora se desea parametrizar la parte del paraboloide acotado por la curva $C$.
Para este propósito, el valor de $y$ se consigue sustituyendo en la expresión adecuada:

%\vspace*{2mm} 
{\footnotesize
	\[
	y = 4-\left( \frac{4u}{\sqrt{3}}\cos(t)\right) ^{2}-\left(\frac{4u}{\sqrt{5}}\sin (t)\right)^{2} = 4-\frac{32}{15}u^{2}\cos ^{2}(t) -\frac{16}{5}u^{2},
	\]
}
\vspace{2mm} 
por lo tanto,
{\footnotesize
	\[
	\overrightarrow{r}(t,u)=\left\langle \frac{4u}{\sqrt{3}}\cos (t) , 4-\frac{32}{15}u^{2}\cos ^{2}(t) -\frac{16}{5}u^{2} , \frac{4u}{\sqrt{5}}\sin (t) \right\rangle,
	\]
}
con $t\in \lbrack 0,2\pi ],\, u\in \lbrack 0,1 ], $ representa la superficie deseada.

\begin{figure}[H]
	\centering 
	\captionof{figure}{\small Superficie acotada por los paraboloides $4y=z^{2}-x^{2}$ y $y=4-x^{2}-z^{2}$}
	\includegraphics[height=4.5cm]{Img/4c.pdf} \hspace{0.5cm} 
	\includegraphics[height=4.5cm]{Img/4d.pdf} 
	\parbox{0.95\linewidth}{\centering\fontsize{8pt}{10pt}\selectfont\textbf{Fuente:} elaboración propia.}
\end{figure}

\noindent{Esta figura se obtiene ejecutando las siguientes instrucciones:}
\vspace{0.5em}

\begin{tcolorbox}[breakable,enhanced,title=\bf Instrucción \verb"Mathematica"]  
\begin{Verbatim}[formatcom=\letraverbatim]
a=Show[ContourPlot3D[{4*y==z^2-x^2,y==4-x^2-z^2},{x,-3,3},
{y,-3,4},{z,-3,3},Mesh->None,PlotPoints->90,
ContourStyle->{{Gray,Orange},Opacity[0.6]}],
ParametricPlot3D[{4/Sqrt[3]*Cos[t],4-(4/Sqrt[5]*Sin[t])^2
+(4/Sqrt[3]*Cos[t])^2,4/Sqrt[5]*Sin[t]},{t,0,2*Pi},
PlotPoints->80,PlotStyle->{Purple,Thickness->0.015}]];

b=ParametricPlot3D[{4*u/Sqrt[3]*Cos[t],
4-32/15*u^2*Cos[t]^2-
16*u^2/5,4*u/Sqrt[5]*Sin[t]},{t,0,2*Pi},{u,0,1},
PlotPoints->90,
Ticks->{{-2,0,2},{0,2,4},{-1,0,1}},PlotStyle->RGBColor[0.9,
0.1,0.1,1],MeshStyle->RGBColor[0.3,0.1,0.1,1]];

Show[a,b]
\end{Verbatim}
\end{tcolorbox}

\subsection*{Una esfera y un cilindro elíptico}

Consideraremos la esfera 
$x^2 + y^2 + z^2 = 16$ y el cilindro elíptico $\frac{(x - 1)^2}{9} +\frac{y^2}{4} = 1$. 
La región interior al cilindro se parametriza en la forma $x-1=3u\cos(t)$, $y=2u\sin(t)$, 
con $u \in[0,1]$ y $t \in [0,2\pi]$. Mientras que para la frontera usaremos $x=1+3\cos(t)$, $y=2\sin(t)$.
\vskip0.3cm

\begin{figure}[H]
	\begin{center}
		\centering
		\captionof{figure}{\small Esfera y un cilindro elíptico}
		\includegraphics[width=3.5cm,height=3.25cm]{Img/2a.pdf}\hspace{0.5cm}
		\includegraphics[width=3.5cm,height=3.25cm]{Img/2b.pdf}
		\parbox{0.95\linewidth}{\centering\fontsize{8pt}{10pt}\selectfont\textbf{Fuente:} elaboración propia.}
	\end{center}
\end{figure}

Con lo anterior, se obtienen las ecuaciones de la curva de intersección del cilindro y la esfera que se determinan reemplazando el valor de $z$ en la ecuación, es decir
$\left( 1+3\cos (t) \right) ^{2}+\left( 2\sin (t) \right) ^{2}+z^{2}=16$, de donde \linebreak
\[
z=\pm \sqrt{5\sin ^{2}(t) -6\cos (t) +6}.
\]
\vskip0.1cm 
Por lo tanto, una parametrización para las curvas en mención está dada por
{\small
\[
\vec{r}(t) = \left\langle 1 + 3\cos(t),\, 2\sin(t),\, \pm \sqrt{5\sin^2(t) - 6\cos(t) + 6} \right\rangle,
\quad t \in [0, 2\pi].
\]
}

Una parte de la curva determina la raíz positiva, mientras que la raíz negativa representa la otra parte de la curva. La figura de esta situación se muestra a continuación.


\vskip0.2cm
\noindent
\begin{minipage}{\textwidth}
	\centering
	\captionof{figure}{\small Esfera y un cilindro elíptico con su curva de intersección}
	\begin{minipage}{0.24\linewidth}
		\centering
		\includegraphics[width=\linewidth]{Img/2c.pdf}
	\end{minipage}
	\begin{minipage}{0.24\linewidth}
		\centering
		\includegraphics[width=\linewidth]{Img/2d.pdf}
	\end{minipage}
	
	\vspace{0.3em}
	\parbox{0.95\linewidth}{\centering\fontsize{8pt}{10pt}\selectfont\textbf{Fuente:} elaboración propia.}
\end{minipage}
\vskip0.2cm




Para la superficie de la esfera que se halla dentro del cilindro, el valor de $z$ está dado por:
\vspace*{-3mm} % Ajuste vertical para pegar la ecuación al párrafo
\[
\left( 1+3u\cos (t) \right) ^{2}+\left( 2u\sin (t) \right) ^{2}+z^{2}=16,
\]
\vspace{4mm} % Espacio vertical para separar la ecuación del texto siguiente
\noindent de donde, \[z= \pm \sqrt{15-6u\cos (t) -9u^{2}+5u^{2}\sin ^{2}(t)}\] con esto se establece que
$${\bf r}(u,t) =\left\langle 1+3u\cos(t),u\sin(t),\sqrt{15-6u\cos(t)-9u^{2}+5u^{2}\sin^{2}(t)}\right\rangle,$$

con $ t \in [0, 2\pi], \,\, u\in [0,1]  $, es una parametrización de la parte de la esfera dentro del cilindro.

\vskip0.2cm
\noindent
\begin{minipage}{\textwidth} % <--- MINIPAGE: Fija la figura al flujo de texto
	\centering % <--- \centering es más compacto que \begin{center}
		\captionof{figure}{\small Superficie sobre la esfera acotada por el cilindro elíptico}
		\includegraphics[width=2.8cm,height=2.90cm]{Img/2e.pdf}\hspace{1cm}
		\includegraphics[width=2.8cm,height=2.9cm]{Img/2f.pdf}
		\parbox{0.95\linewidth}{\centering\fontsize{8pt}{10pt}\selectfont\textbf{Fuente:} elaboración propia.}
	\end{minipage}

\vspace{0.3cm}
\noindent 
Para representar la parte del cilindro que está dentro de la esfera, podemos tomar dos puntos $P = \langle 1 + 3\cos(t),\, 2\sin(t),\, \sqrt{6 - 6\cos(t) + 5\sin^2(t)} \rangle$ y $Q = \langle 1 + 3\cos(t),\, 2\sin(t),\, -\sqrt{6 - 6\cos(t) + 5\sin^2(t)} \rangle$. Con ellos se parametriza el segmento rectilíneo que los une. Simplificando la expresión $(1 - u)P + uQ$, con $u \in [0, 1]$, se obtiene la función $\vec{r}(u, t)$ dada por:	
	\[
	\vec{r}(u, t) = \langle 1 + 3\cos(t),\, 2\sin(t),\, (1 - 2u)\sqrt{6 - 6\cos(t) + 5\sin^2(t)} \rangle
	\]	
Esta función, con $t \in [0, 2\pi]$ y $u \in [0, 1]$, representa una parametrización de la superficie sobre el cilindro que se encuentra dentro de la esfera.

\begin{figure}[H]
	\centering
	\caption{Superficie sobre el cilindro elíptico acotada por la esfera}
	\label{Superficie_cilindro_elíptico_acotada_esfera}
	\begin{minipage}[b]{0.24\linewidth}
		\centering
		\includegraphics[width=\linewidth]{Img/2g.pdf}
	\end{minipage}
	\hspace{1em}
	\begin{minipage}[b]{0.25\linewidth}
		\centering
		\includegraphics[width=\linewidth]{Img/2h.pdf}
	\end{minipage}
	\hspace{1em}
	\begin{minipage}[b]{0.24\linewidth}
		\centering
		\includegraphics[width=\linewidth]{Fig/65.pdf}
	\end{minipage}
	\parbox{0.95\linewidth}{\centering\fontsize{9pt}{11pt}\selectfont{\bf Fuente:} elaboración propia.}
\end{figure}


Los gráficos se obtienen con ayuda de las siguientes instrucciones:

\begin{tcolorbox}[breakable,enhanced,title=\bf Instrucción \verb"Mathematica"] 
\begin{Verbatim}[formatcom=\letraverbatim]
Show[{ContourPlot3D[{x^2+y^2+z^2==1,(x-1)^2/9+y^2/4==1},
{x,-4,4},{y,-4,4},{z,-4,4.0},ContourStyle->{Directive[Orange,
Opacity[0.15]],Directive[Gray,Opacity[0.15]]},Mesh->False],
ParametricPlot3D[{1+3*u*Cos[t],2*u*Sin[t],Sqrt[15-6*u*Cos[t]
-9*u^2+2+5*u^2*(Sin[t])^2]},{u,0,1}{t,0,2Pi},PlotStyle->Red]}]

Show[{ContourPlot3D[{x^2+y^2+z^2==1,(x-1)^2/9+y^2/4==1},
{x,-4,4},{y,-4,4},{z,-4,4.0},ContourStyle->{Directive[Orange,
Opacity[0.15]],Directive[Gray,Opacity[0.15]]},Mesh->False],
ParametricPlot3D[{1+3*Cos[t], 2*Sin[t],(1-2*u)Sqrt[6
-6*Cos[t]+5*(Sin[t])^2]},{u,0,1},{t,0,2Pi},PlotStyle->Red]}]
\end{Verbatim}
\end{tcolorbox}


\section{Sólidos} \label{p1} \index{Parametrización!de sólidos}
Con los elementos anteriores parametrizamos algunos sólidos.

\subsection*{Sólido acotado por dos paraboloides}

Los paraboloides elípticos con ecuaciones $z=3-2x^2-y^2$ y $z=x^2+2y^2,$ se intersectan en una región determinada por $x^2+y^2=1,$ cuyo interior se representa con $x=u \cos v,$ $y=u \sin v$ con $u\in [0,1]$ y $v\in [0,2\pi].$


%\vspace*{-0.5\baselineskip} 
La coordenada $z$ de dos puntos $P$ y $Q$ sobre los paraboloides, es
\begin{align*}
z_1 &= 3-2(u\cos v)^2-(u\sin v)^2 \\ &= \frac{1}{2}(6-3u^2-u^2 \cos(2v)) \\
z_2 &= (u\cos v)^2+2(u\sin v)^2 \\ &= \frac{u^2}{2}(3-\cos(2v))
\end{align*}


\vskip0.2cm
\noindent
\begin{minipage}{\textwidth}
	\centering
	\captionof{figure}{\small Sólido acotado por paraboloides}
	\label{fig:300}
	\begin{minipage}{0.35\linewidth}
		\centering
		\includegraphics[width=\linewidth]{Fig/300.pdf}
	\end{minipage}
	
	\vspace{0.3em}
	\parbox{0.95\linewidth}{\centering\fontsize{8pt}{10pt}\selectfont\textbf{Fuente:} elaboración propia.}
\end{minipage}
\vskip0.2cm


\noindent
La región se parametriza uniendo los puntos $P(u\cos v,u\sin v,\frac{u^2}{2}(3-\cos(2t)))$ y 
$Q(u\cos v,u\sin v,\frac{1}{2} (6-3u^2-u^2 \cos(2t))),$
mediante segmentos rectilíneos. Es decir, la expresión $(1-w)P+wQ$ se simplifica y permite definir la función
$${\bf r}(u,v,w)= \langle u\cos v,u\sin v, 3(1-w)+u^2(6w-3-\cos(2t))/2 \rangle ,$$
$u\in [0,1],$ $v\in [0,2\pi],$ $w\in [0,1].$  Una representación de estas superficies se obtiene con las siguientes instrucciones:

\begin{tcolorbox}[breakable,enhanced,title=\bf Instrucción \verb"Mathematica"] 
\begin{Verbatim}[formatcom=\letraverbatim]
A=ContourPlot3D[{z==3-2*x^2-y^2,z==x^2+2*y^2},
{x,-2,2},{y,-2,2},{z,0,3},Mesh->False,
ContourStyle->{Directive[Orange,Opacity[0.8]],
Directive[Cyan,Opacity[0.95]]},PlotPoints->80];
B=ParametricPlot3D[{Cos[t],Sin[t],Cos[t]^2+
2*Sin[t]^2},{t,0,2Pi},PlotStyle->Black];
Show[A,B]
RegionPlot3D[x^2+2*y^2<=z<=3-2*x^2-y^2&&x^2+y^2<=1,
{x,-1,1},{y,-1,1},{z,0,3},PlotPoints->80,Mesh->False,
PlotStyle->Orange]
\end{Verbatim}
\end{tcolorbox}

\subsection*{Una cuña circular}
Consideremos el sólido acotado por los gráficos de las ecuaciones
$r=3 \cos\theta,$ $z=-y$ y $z=0.$ Esta región se muestra a continuación

\begin{figure}[H]
	\centering
	\caption{Sólido determinado por $0\leq z \leq -y$, $r \leq 3 \cos \theta.$}
	\label{Sólido_determinado}
	\vspace{-4mm}
	\includegraphics[width=0.27\linewidth]{Fig/66.pdf}\hspace{1mm}
	% Eliminamos los vspace intermedios para evitar saltos de línea accidentales
	\vspace{-4mm}
	\includegraphics[width=0.27\linewidth]{Fig/67.pdf}\hspace{1mm}
	\vspace{-4mm}
	\includegraphics[width=0.27\linewidth]{Fig/68.pdf}
	
	\vspace{0.6cm} % <--- AJUSTA ESTE VALOR para bajar la fuente a tu gusto
	\parbox{0.95\linewidth}{\centering\fontsize{8pt}{10pt}\selectfont\textbf{Fuente:} elaboración propia.}
\end{figure}


La base del sólido es el semicírculo de radio $3$ que se parametriza en la forma $\textstyle\langle 3\cos\theta\cos\theta,\allowbreak\, 3\cos\theta\sin\theta \rangle$, y se simplifica como $\textstyle\langle 3\cos^2\theta,\allowbreak\, \tfrac{3}{2}\sin(2\theta) \rangle$, con $\theta \in [\pi, 2\pi]$. El interior de esta región se representa en el plano $xy$ como $\textstyle\langle 3u\cos^2\theta,\allowbreak\, \tfrac{3}{2}u\sin(2\theta) \rangle$, con $u \in [0,1]$ y $\theta \in [\pi, 2\pi]$. Un punto en la base del sólido tendrá la forma $\textstyle P = (3u\cos^2\theta,\allowbreak\, \tfrac{3}{2}u\sin(2\theta),\, 0)$, mientras que otro punto sobre el plano $z = -y$ tendrá la forma $\textstyle Q = (3u\cos^2\theta,\allowbreak\, \tfrac{3}{2}u\sin(2\theta),\allowbreak\, -\tfrac{3}{2}u\sin(2\theta))$. Formando una combinación convexa $\textstyle(1 - w)P + wQ$ entre estos dos puntos, se obtiene la función:
\vspace{-0.70mm}
{\small
\begin{equation*}
{\bf r}(\theta, u,w) = 
\bigg\langle 3u \cos^2\theta, \frac{3}{2}u \sin(2 \theta) , -\frac{3}{2}uw \sin(2 \theta)\bigg \rangle,
\end{equation*}
}
con $\theta \in  [\pi, 2\pi],$ $u\in [0,1]$ y $w\in [0,1],$ que parametriza el sólido.

\subsection*{Un sólido acotado por un cono y una esfera}
Las ecuaciones $z=\sqrt{x^2+y^2}$ y $x^2-4x+y^2+z^2=6$ 
representan un semicono y una esfera.
El eliminar $z$ de estas expresiones, se obtiene la ecuación
$\left( x-1\right) ^2+y^2=4$ que representa la curva de intersección de las superficies en el plano $xy$.
El interior de esta región se puede parametrizar en la forma
$x-1=2u\cos(t)$, $y=2u\sin(t)$, con $t \in [0,2\pi]$ y $u \in [0,1].$ 

\begin{figure}[H]
	\centering
	\caption{Semiesfera $z=\sqrt{6-x^2-y^2+4x}$ y semicono $z=\sqrt{x^2+y^2}$}
	\label{Semiesfera_y_semicono}
	\includegraphics[width=0.23\linewidth]{Img/74c.pdf}\hspace{3mm}
	\includegraphics[width=0.33\linewidth]{Fig/69.pdf} 
	\parbox{0.95\linewidth}{\centering\fontsize{9pt}{11pt}\selectfont{\bf Fuente:} elaboración propia.}
\end{figure}
\vskip0.2cm
Determinaremos la forma de los puntos sobre el semicono y sobre la esfera, posteriormente se hará una combinación convexa entre estos puntos con el fin de describir completamente el sólido. Sobre el semicono, se satisface que

\vspace{-7mm}
{\small
\[
 z=\sqrt{x^{2}+y^{2}}=\sqrt{\left( 1+2u\cos \left( t\right) \right)^{2}+\left( 2u\sin \left( t\right) \right) ^{2}}= \sqrt{1+4u\cos\left( t\right) +4u^{2}}.
\]
}

Entonces, los puntos sobre el cono (dentro del sólido) se determinan por

{\small
$$P\left( 1+2u\cos (t)  ,2u\sin (t) ,\sqrt{1+4u\cos\left( t\right) +4u^{2}} \right). $$
}

Reemplazando $x$ e $y$ en la ecuación de la esfera {\footnotesize $z=\sqrt{6 - x^{2} + 4x - y^{2}}$,} se obtiene {\footnotesize $z=\sqrt{6 - (1{+}2u\cos t)^2 + 4(1{+}2u\cos t) - (2u\sin t)^2} = \sqrt{9{+}4u\cos t{-}4u^2}$,} entonces los puntos de la forma {\footnotesize $Q(1{+}2u\cos t,\, 2u\sin t,\, \sqrt{9{+}4u\cos t{-}4u^2})$} yacen sobre la esfera. Por lo tanto, el sólido se describe al simplificar la expresión $(1{-}w)P{+}wQ$, para $w\in[0,1]$, de modo que la siguiente función describe el sólido, para $t\in[0,2\pi]$, $u\in[0,1]$ y $w\in[0,1]$. La figura \ref{region_conos} muestra esta situación.
\vspace{-04mm}
\begin{align*}
	\mathbf{r}(t,u,w) &= \left\langle 1{+}2u\cos t,\, 2u\sin t,\right. \\
	&\quad \left.(1{-}w)\sqrt{1{+}4u\cos t{+}4u^2} + w\sqrt{9{+}4u\cos t{-}4u^2} \right\rangle
\end{align*}


\vskip0.2cm
\noindent
\begin{minipage}{\textwidth}
	\centering	
	\captionof{figure}{\small Región determinada por $x^2+y^2 \leq z^2 \leq 6-x^2-y^2+4x$}
	\label{region_conos}
	\begin{minipage}{0.26\linewidth}
		\centering
		\includegraphics[width=\linewidth]{Img/74a.pdf}
	\end{minipage}
	\begin{minipage}{0.26\linewidth}
		\centering
		\includegraphics[width=\linewidth]{Img/74b.pdf}
	\end{minipage}
	
	\vspace{0.3em}
	\parbox{0.95\linewidth}{\centering\fontsize{8pt}{10pt}\selectfont\textbf{Fuente:} elaboración propia.}
\end{minipage}
\vskip0.2cm


La figura se obtiene ejecutando las siguientes instrucciones:



\begin{tcolorbox}[breakable,enhanced,title=\bf Instrucción \verb"Mathematica"] 
\begin{Verbatim}[formatcom=\letraverbatim]
ParametricPlot3D[{{1+2*u*Cos[t],2*u*Sin[t],
Sqrt[9+4*u*Cos[t]-4*u^2]},{1+2*u*Cos[t],
2*u*Sin[t],Sqrt[1+4*u*Cos[t]+4*u^2]}},
{u,0,1},{t,0,2*Pi},
PlotStyle->{RGBColor[0.8,0.3,0.3],Cyan},
MeshStyle->Gray,PlotPoints->60]
\end{Verbatim}
\end{tcolorbox}

\subsection*{Un sólido hueco}
Las expresiones $2x\leq x^{2}+y^{2}$ y $0\leq z\leq 9-x^{2}-y^{2},$
describen el sólido que se muestra a continuación.


\vskip0.2cm
\noindent
\begin{minipage}{\textwidth}
	\centering
	\captionof{figure}{\small Sólido interior al paraboloide y exterior al cilindro}
	\label{conos_naranjas}
	\begin{minipage}{0.22\linewidth}
		\centering
		\includegraphics[width=\linewidth]{Img/79a.pdf}
	\end{minipage}
	\begin{minipage}{0.22\linewidth}
		\centering
		\includegraphics[width=\linewidth]{Img/79c.pdf}
	\end{minipage}
	\begin{minipage}{0.22\linewidth}
		\centering
		\includegraphics[width=\linewidth]{Img/79b.pdf}
	\end{minipage}
	
	\vspace{0.3em}
	\parbox{0.95\linewidth}{\centering\fontsize{8pt}{10pt}\selectfont\textbf{Fuente:} elaboración propia.}
\end{minipage}
\vskip0.2cm


Este objeto consta de los puntos que yacen entre el plano $xy$ y el paraboloide $z=9-x^{2}-y^{2},$
exteriores al cilindro circular recto $2x=x^{2}+y^{2}.$
Esta última ecuación se reescribe como $\left( x-1\right) ^{2}+y^{2}=1$ y representa en el plano $xy$
una circunferencia. La figura \ref{conos_naranjas} se consigue ejecutando las siguientes instrucciones.
\vskip0.2cm

\begin{tcolorbox}[breakable,enhanced,title=\bf Instrucción \verb"Mathematica"]
\begin{Verbatim}[formatcom=\letraverbatim]
RegionPlot3D[x^2+y^2>2*x&&0<z<9-x^2-y^2,{x,-3,3},
{y,-3,3},{z,0,9},PlotPoints->90,Mesh->False,
PlotStyle->Orange]
\end{Verbatim}
\end{tcolorbox}
 

El paraboloide se proyecta en el plano $xy$ como una circunferencia con ecuación $x^{2}+y^{2}=9,$
la proyección del sólido en plano $xy$ contiene a los puntos interiores a una circunferencia y exteriores a la otra.


Para recorrer la circunferencia de radio $3$, usaremos {\small $2t$} como parámetro. La región entre las circunferencias en mención se parametriza
usando una combinación convexa entre dos puntos {\footnotesize $\left( 1+\cos (2t) ,\sin (2t)\right)$} y {\footnotesize $\left(3\cos(2t),3\sin(2t)\right),$} la cual se simplifica como 
\vspace{-3mm}
{\small
\begin{align*}
	{\bf r}(u,t)&=\big\langle u+(3-2 u)\cos (2t),\, (3-2u)\sin(2t)\big\rangle.
\end{align*}
}

\vskip0.2cm
\noindent
\begin{minipage}{\textwidth}
	\centering
	\captionof{figure}{${\bf r}(u,t)$}
	\label{fig:70}
	\begin{minipage}{0.24\linewidth}
		\centering
		\includegraphics[width=\linewidth]{Fig/70.pdf}
	\end{minipage}
	
	\vspace{0.3em}
	\parbox{0.95\linewidth}{\centering\fontsize{8pt}{10pt}\selectfont\textbf{Fuente:} elaboración propia.}
\end{minipage}
\vskip0.2cm


Las siguientes instrucciones muestran una animación el recorrido de las curvas.

\begin{tcolorbox}[breakable,enhanced,title=\bf Instrucción \verb"Mathematica"] 
\begin{Verbatim}[formatcom=\letraverbatim]
Manipulate[ParametricPlot[{{3*Cos[2*t],3*Sin[2*t]},
{1+Cos[2*t],Sin[2*t]}},{t,0,k},PlotRange->{-3,3}],
{k,Pi/20,Pi}]
\end{Verbatim}
\end{tcolorbox}

Con un punto $P(u+(3-2u)\cos(2t),(3-2u)\sin(2t),0)$ en la base del sólido 
se obtiene la coordenada $z$ de un punto $Q$ sobre el paraboloide,
\begin{eqnarray*}
z &=& 9-(u+(3-2u)\cos(2t))^2-((3-2u)\sin(2t))^{2} \\
&=& u (12-5 u+2(2u-3)\cos(2t)).
\end{eqnarray*}
El sólido se parametriza formando segmentos rectilíneos que unen a $P$
con $Q,$ esto se logra mediante la combinación convexa $(1-v)P+vQ,$
en donde $Q$ tiene la forma
$Q(u+(3-2u)\cos(2t),(3-2u)\sin(2t),u (12-5 u+2(2u-3)\cos(2t))).$ Esto nos lleva a definir la función 
\begin{multline} \label{solidohueco}
{\bf r}(t,u,v) = \big \langle u+(3-2u)\cos(2t),(3-2u)\sin(2t), \\ uv(12-5 u+2(2u-3)\cos(2t)) \big \rangle
\end{multline}
En donde $t\in [0,\pi]$, $u\in [0,1]$ y $v \in [0,1]$.

\subsection*{Un sólido hueco II} \label{solidohueco3}

Las ecuaciones $z=4-y^{2},$ $z=4-x,$ $z=0,$ $x=0,$ representan un cilindro parabólico y 
tres planos. Estas superficies acotan un sólido $E,$
que se perfora por el cilindro $\left( z-2\right) ^{2}+y^{2}=1,$ obteniendo la región 
que se muestra en la figura.


\vskip0.2cm
\noindent
\begin{minipage}{\textwidth}
	\centering
	\captionof{figure}{$z=4-y^2$, $z=4-x$}
	\label{fig:71}
	\begin{minipage}{0.30\linewidth}
		\centering
		\includegraphics[width=\linewidth]{Fig/71.pdf}
	\end{minipage}
	
	\vspace{0.3em}
	\parbox{0.95\linewidth}{\centering\fontsize{8pt}{10pt}\selectfont\textbf{Fuente:} elaboración propia.}
\end{minipage}
\vskip0.2cm

 
%\vspace{-4mm}
La figura del sólido se obtiene ejecutando las siguientes instrucciones.

\begin{tcolorbox}[breakable,enhanced,title=\bf Instrucción \verb"Mathematica"] 
\begin{Verbatim}[formatcom=\letraverbatim]
RegionPlot3D[0<x<4-z&&(z-2)^2+y^2>1&&0<z<4-y^2&&
-2<y<2,{x,0,4},{y,-2,2},{z,0,4},PlotPoints->90,
PlotStyle->Orange,Mesh->False]
\end{Verbatim}
\end{tcolorbox}

\noindent
\begin{minipage}{\textwidth}
	\centering
	\captionof{figure}{\small Región determinada por $0\leq z \leq 4-x$, $z\leq 4-y^2,$ $1\leq (z-2)^2+y^2.$}
	\begin{minipage}{0.22\linewidth}
		\centering
		\includegraphics[width=\linewidth]{Img/5a.pdf}
	\end{minipage}
	\hspace{1em}
	\begin{minipage}{0.22\linewidth}
		\centering
		\includegraphics[width=\linewidth]{Img/5c.pdf}
	\end{minipage}
	\hspace{1em}
	\begin{minipage}{0.22\linewidth}
		\centering
		\includegraphics[width=\linewidth]{Img/5b.pdf}
	\end{minipage}
	
	\vspace{0.3em}
	\parbox{0.95\linewidth}{\centering\fontsize{8pt}{10pt}\selectfont\textbf{Fuente:} elaboración propia.}
\end{minipage}


La proyección del sólido en el plano $yz$ corresponde a una circunferencia y una porción de 
parábola, que se representan en la forma $\ve{r}_1(t)=\la 0,\cos (t),2+\sin (t)\ra,$ con 
$t\in [0,2\pi]$ y $\ve{r}_2(t)=\la 0,t,4-t^{2}\ra,$ con $t\in \lbrack -2,2],$ respectivamente.
Inicialmente parametrizaremos la base del sólido en el plano $yz$. Esta región a su vez 
la dividiremos en dos porciones.

\begin{figure}[H] 
	\centering
	\caption{Región determinada por $0\leq z \leq 4-y^2$, $1\leq (z-2)^2+y^2.$}
	\label{Región_determinada}
	\begin{minipage}[t]{0.35\linewidth} 
		\centering
		\includegraphics[width=\linewidth]{Fig/72.pdf}
	\end{minipage}
	\hspace{1em}
	\begin{minipage}[t]{0.35\linewidth} 
		\centering
		\includegraphics[width=\linewidth]{Fig/73.pdf}
	\end{minipage}
	\parbox{0.95\linewidth}{\centering\fontsize{9pt}{11pt}\selectfont{\bf Fuente:} elaboración propia.}
\end{figure}
%\vskip0.1cm
Para describir las regiones anteriores unificaremos el dominio de las parametrizaciones dadas y luego construir segmentos rectilíneos entre las curvas.
\vskip2cm
\vspace{-10mm}

\noindent
\textcolor{black}{\scriptsize\ding{110}}\quad {\bf Región 1.} Para aplicar el intervalo $[-2,2]$ en $[0,\pi ],$ se utiliza la relación
$v=\frac{\pi }{4}\left(u+2\right)$; además, con el intercambio de $t$ por $\pi -t$ se invierte el recorrido. 

\vskip0.1cm
El cambio de $t$ por $2-\frac{4t}{\pi}$ en la parametrización de la parábola,
unifica el dominio de las funciones. Entonces, con una combinación convexa de los puntos 
$P(0,\cos (t),2+\sin(t))$ y $Q(0,2-\frac{4t}{\pi},\frac{16t}{\pi}\left(1-\frac{t}{\pi}\right))$, se obtiene la función que parametriza la región en mención, esto es
\begin{multline*}
{\bf r}_1(t,u)  =\big \langle 0, (1-u)\cos (t) +u\left(2-\frac{4t}{\pi}\right),\\ \frac{16(\pi-t)tu}{\pi^2}-(u-1)(\sin(t)+2) \big\rangle, 
\end{multline*} 
para $t\in [0,\pi ]$ y $u\in [0,1]$.

\vspace{2mm}

\noindent
\textcolor{black}{\scriptsize\ding{110}}\quad {\bf Región 2.} El segmento que une los puntos $\left(0,-2,0 \right) $ con $\left(0,2,0 \right),$
se parametriza como $\langle 0,4(t-\pi)/\pi-2,0,0 \rangle $,
con $t\in [\pi ,2\pi ].$ Trazando segmentos rectilíneos desde el segmento anterior 
hasta la semicircunferencia, se consigue parametrizar el conjunto de interés; esto se logra 
formando una combinación convexa entre los puntos
$(0,\cos (t),2+\sin (t))$ y $(0,4(t-\pi)/\pi-2,0)$ por esta razón, la función 
$${\bf r}_2\left( t,u\right)  = \left\langle 0,\left(
1-u\right) \cos (t) +u\left( \frac{4\left( t-\pi \right) }{\pi }%
-2\right) ,\left( 1-u\right) \left( 2+\sin (t) \right)
\right\rangle ,$$ con $t\in \lbrack \pi ,2\pi ]$ y $u\in \lbrack 0,1],$ parametriza la región considerada.


Con las funciones ${\bf r}_1$ y ${\bf r}_2$ se representará la región sobre el 
plano $z=4-x,$ esto cambia solamente la primera función componente de estas funciones. 
\begin{multline*}
{\bf r}_3(u,t)=\bigg \langle 4-\frac{16 (\pi-t) t u}{\pi ^2}+(u-1) (\sin (t)+2),
(1-u)\cos (t) \\ +u\left(2-\frac{4t}{\pi}\right), \frac{16 (\pi -t) t u}{\pi ^2}-(u-1) (\sin (t)+2) \bigg\rangle, t\in[0,\pi],\, u\in[0,1].
\end{multline*}  \begin{multline*}
{\bf r}_4(u,t)= \big \langle 4-(1-u)(\sin(t)+2),\left(\frac{4(t-\pi)}{\pi}-2\right)u+(1-u)\cos(t),\\
(1-u)(\sin (t)+2)\big \rangle,\,t\in[\pi,2\pi],\, u\in[0,1].
\end{multline*}

Con las siguientes instrucciones se obtiene el gráfico de las superficies que acotan el sólido.


%\fontsize{9.5}{10}\selectfont
\begin{tcolorbox}[breakable,enhanced,title=\bf Instrucción \verb"Mathematica"] 
\begin{Verbatim}[formatcom=\letraverbatim]
y1=(1-u)*Cos[t]+u*(2-4*t/Pi);
z1=16*(Pi-t)*u*t/Pi^2-(u-1)*(2+Sin[t]);
{y2,z2}={(1-u)*Cos[t]+u*(4*(t-Pi)/Pi-2),
(1-u)*(2+Sin[t])};
A1=ParametricPlot3D[{{0,y1,z1},{4-z1,y1,z1}},{t,0,Pi},
{u,0,1},Mesh->False,PlotStyle->Orange];
A2=ParametricPlot3D[{{0,y2,z2},{4-z2,y2,z2}},
{t,Pi,2*Pi},{u,0,1},Mesh->False,PlotStyle->Cyan];
r5=(1-v)*{0,y1,z1}+v*{4-z1,y1,z1}/.u->1;
A3=ParametricPlot3D[r5,{t,0,Pi},{v,0,1},
PlotStyle->Orange,Mesh->False];
Show[A1,A2,A3]
\end{Verbatim}
\end{tcolorbox}

El sólido quedará parametrizado al hacer una combinación convexa entre dos puntos: 
uno en la base del sólido (del plano $yz$) y otro punto sobre el plano $z=4-x$. 
Es por lo expuesto antes, que una parametrización del sólido es:
\begin{equation} \label{solidohueco2}
{\bf r}(u,t,w) = \left \{ \begin{array}{cc}
(1-w){\bf r}_1(u,t)+w {\bf r}_3(u,t) & t \in [0,\pi] \\[0.2cm]
(1-w){\bf r}_2(u,t)+w {\bf r}_4(u,t) & t \in [\pi,2\pi] 
\end{array}\right . 
\end{equation}

\subsection*{Sólidos con una cardioide en su base} \index{Cardioide}
\noindent
\textcolor{black}{\scriptsize\ding{110}}\quad La cardioide con ecuación $r=1+\cos (t) $
se parametriza en el plano $xy$ como  $\left\langle (1+\cos(t)) \cos(t),(1+\cos (t))\sin (t) \right\rangle ,$ 
$t\in [0,2\pi];$ para representar su interior se hace el siguiente cambio sutil
$\big\langle u(1+\cos(t)) \cos(t),u(1+\cos (t))\sin (t) \big\rangle ,$  con $u\in[0,1].$
Estos puntos se unen a $(1,0,1)$ utilizando el par de funciones inyectivas en 
$[0,1],$ $f(v)=\sqrt[5]{(1-v) ^{2}}$ y $g(v) =v^{2},$ lo cual parametriza el sólido que se muestra a continuación.

\begin{figure}[H]
\begin{center}
	\captionof{figure}{\small Gráfico de ${\bf r}\left(t,u,w\right)$}
\includegraphics[height=4.5cm]{Img/76b.pdf} 
\includegraphics[height=4.5cm]{Img/76a.pdf}
\includegraphics[height=4.5cm]{Img/76c.pdf}
\end{center}
\parbox{0.95\linewidth}{\centering\fontsize{9pt}{11pt}\selectfont{\bf Fuente:} elaboración propia.}
\end{figure}

La combinación que genera el sólido es:
$${\bf r}(t,u,w) =\sqrt[5]{(1-v)^2}\, u(1+\cos(t)) (\cos(t),\sin(t),0) +v^{2}(1,0,1),$$
con $t \in [0,2\pi]$, $u \in [0,1]$ y $w \in [0,1].$ 
Las siguientes instrucciones, con $u=1$, permiten obtener el gráfico anterior.

\begin{tcolorbox}[breakable,enhanced,title=\bf Instrucción \verb"Mathematica"] 
\begin{Verbatim}[formatcom=\letraverbatim]
r=(1-v)^(2/5)*{(1+Cos[t])*Cos[t],(1+Cos[t])*Sin[t],0}
+v^2*{1,0,1};
p={r[[1]],r[[2]],0};
ParametricPlot3D[{r,p},{v,0,1},{t,0,2Pi},PlotPoints
->80,PlotStyle->{Orange,Cyan},MeshStyle->Gray]
\end{Verbatim}
\end{tcolorbox}

\vspace{2mm}

\noindent
\textcolor{black}{\scriptsize\ding{110}}\quad De manera similar al ejemplo anterior, se puede utilizar
la función 
$${\bf r}(t,u,v)=\sqrt[4]{(1-v)^3}\, u(1+\cos(t)) (\cos(t),\sin(t),0)+\sqrt[4]{v^3}(1,0,1),$$
con $t \in [0,2\pi]$, $u \in [0,1]$ y $w \in [0,1].$ Para $u=1$ se genera la superficie cuyo gráfico se muestra
a continuación.

\begin{figure}[H]
\begin{center}
	\captionof{figure}{\small Gráfico de ${\bf r}\left(t,u,w\right)$}
\includegraphics[height=4.5cm]{Img/77a.pdf} 
\includegraphics[height=4.5cm]{Img/77b.pdf}
\end{center}
\parbox{0.95\linewidth}{\centering\fontsize{9pt}{10pt}\selectfont{\bf Fuente:} elaboración propia.}
\end{figure}

Con las siguientes instrucciones se puede obtener el gráfico.

\begin{tcolorbox}[breakable,enhanced,title=\bf Instrucción \verb"Mathematica"]
\begin{Verbatim}[formatcom=\letraverbatim]
r=(1-v)^(3/4)*{(1+Cos[t])*Cos[t],(1+Cos[t])*Sin[t],0}
+v^(3/4)*{1,0,1};
p={r[[1]],r[[2]],0};
ParametricPlot3D[{r,p},{v,0,1},{t,0,2Pi},PlotPoints
->80,PlotStyle->{Orange,Cyan},MeshStyle->Gray]
\end{Verbatim}
\end{tcolorbox}
  
\subsection*{Cono cardioidico} \index{Cono cardioidico}
El interior de la cardioide con ecuación polar $r=1+\cos(t),$ \index{Cardioide}
se puede parametrizar en el plano $xy$, como 
$\left\langle u\left(1+\cos (t) \right) \cos (t) ,u\left( 1+\cos \left(
t\right) \right) \sin (t) ,0\right\rangle $. El sólido
que se genera mediante una combinación convexa de los puntos de esta región 
con el punto $\left( 1,0,1\right) $ es un \emph{cono cardioidico}\footnote{Este término fué acuñado por el profesor Juan Zambrano, de la Universidad Distrital Francisco José de Caldas.}.

\begin{figure}[H]
\begin{center}
	\captionof{figure}{\small Cono cardioidico}
\includegraphics[height=4.5cm]{Img/75a.pdf} 
\includegraphics[height=4.5cm]{Img/75b.pdf}
\includegraphics[height=4.5cm]{Img/75c.pdf}
\end{center}
\parbox{0.95\linewidth}{\centering\fontsize{9pt}{10pt}\selectfont\textbf{Fuente:} elaboración propia.}
\end{figure}

Es por esto que al simplificar la combinación 
$$\left( 1-w\right) \left( u\left( 1+\cos (t) \right) \cos \left(
t\right) ,u\left( 1+\cos (t) \right) \sin (t),0\right) +w\left( 1,0,1\right) $$
se obtiene la función 
${\bf r}(t,u,w)=\big\langle u\cos (t) \left(
1+\cos (t) \right) \left( 1-w\right) +w,\linebreak (1-w)u\left( 1+\cos (t) \right) \sin (t),w \big\rangle ,$
que describe el mencionado sólido, para $t\in [0,2\pi]$, $u\in [0,1]$ y $w\in [0,1].$
Los gráficos se obtienen con ayuda de las siguientes instrucciones.

\begin{tcolorbox}[breakable,enhanced,title=\bf Instrucción \verb"Mathematica"]
\begin{Verbatim}[formatcom=\letraverbatim]
p={(1+Cos[t])*Cos[t],(1+Cos[t])*Sin[t],0};
r=(1-u)*p+u*{1,0,1};
ParametricPlot3D[{r,p*u},{u,0,1},{t,0,2Pi},
Ticks->{{0,1,2},{-1,0,1},{0,1}},PlotStyle
->{Orange,Cyan},MeshStyle->Gray,PlotPoints->80]
\end{Verbatim}
\end{tcolorbox}

%\clearpage
\begin{ejer}
Parametrice y grafique en \verb"Mathematica" las siguientes regiones.
\end{ejer}

\begin{enumerate}
\item La región del plano comprendida entre las elipses $\frac{x^2}{16}+\frac{y^2}{9}=1$ y $(x-1)^2+\frac{y^2}{4}=1$.

\vskip0.2cm
\vspace{-6mm}
\begin{figure}[H]
	\centering
	\caption{Región acotada por dos elipses}
	\label{Región_acotada_dos_elipses}
	\includegraphics[width=0.4\linewidth]{Fig/74.pdf}
	\parbox{0.95\linewidth}{\centering\fontsize{8pt}{10pt}\selectfont\textbf{Fuente:} elaboración propia.}
\end{figure}

\item La región del plano acotada por las elipses $\frac{x^2}{16}+\frac{y^2}{9}=1$ y $(x-1)^2+\frac{y^2}{4}=1$,
se proyecta sobre el paraboloide hiperbólico $3z=y^2-x^2-2y$.

\begin{figure}[H]
	\centering
	\caption{Superficie sobre el paraboloide hiperbólico $3z = y^2 - x^2 - 2y$.}
	\label{Superficie_paraboloide_hiperbólico}
	\begin{minipage}[t]{0.30\linewidth}
		\centering
		\includegraphics[width=\linewidth]{Fig/75.pdf}
	\end{minipage}%
	\hspace{1em}
	\begin{minipage}[t]{0.30\linewidth}
		\centering
		\includegraphics[width=\linewidth]{Fig/76.pdf}
	\end{minipage}
	\parbox{0.95\linewidth}{\centering\fontsize{8pt}{10pt}\selectfont\textbf{Fuente:} elaboración propia.}
\end{figure}

\item Para $a_1=\cosh 2$, $b_1=\sinh(2)$, $a_2=\cos(1/2)$ y $b_2=\sin(1/2)$, la región del plano interior a la elipse
$\frac{x^2}{a_1^2}+\frac{y^2}{b_1^2}= 1$ y a la hipérbola $\frac{x^2}{a_2^2}-\frac{y^2}{b_2^2}= 1.$
\vspace{-1mm}
\begin{figure}[H]
	\centering
	\caption{Región acotada por una elipse y una hipérbola}
	\label{Región_acotada_elipse_hipérbola}
	\includegraphics[width=0.3\linewidth]{Fig/77.pdf}
	\parbox{0.95\linewidth}{\centering\fontsize{8pt}{10pt}\selectfont\textbf{Fuente:} elaboración propia.}
\end{figure}
\vspace{-2mm}
{\bf Sugerencia:} Utilice coordenadas curvilíneas
\vskip0.2cm
\index{Coordenadas curvilíneas}
\item Para $v=\pm \ln(5)$ parametrizar, con una sola función la región exterior a las dos circunferencias
con ecuaciones $(x-\coth v)^2+y^2=\frac{24}{144}$.

\begin{figure}[H]
	\begin{center}
		\captionof{figure}{\small Región exterior a dos circunferencias.}
		\includegraphics[height=4.5cm]{Img/T_71a.pdf}
	\end{center}
	\parbox{0.95\linewidth}{\centering\fontsize{8pt}{10pt}\selectfont\textbf{Fuente:} elaboración propia.}
\end{figure}
\vspace{-5mm}
{\bf Sugerencia:} Utilice coordenadas curvilíneas.

\item Para $v=\pm \ln(5)$ parametrizar, con una sola función, la región sobre el paraboloide hiperbólico 
$6z=x^2-y^2$, exterior a los cilindros con ecuaciones $(x-\coth v)^2+y^2=\frac{24}{144}$.
\begin{figure}[H]
\begin{center}
	\captionof{figure}{\small Región sobre el paraboloide hiperbólico.}
\includegraphics[height=4.5cm]{Img/T_71b.pdf} \hskip0.5cm
\includegraphics[height=4.5cm]{Img/T_71c.pdf}
\end{center}
\parbox{0.95\linewidth}{\centering\fontsize{8pt}{10pt}\selectfont\textbf{Fuente:} elaboración propia.}

\end{figure} {\bf Sugerencia:} Utilice coordenadas curvilíneas
\begin{comment}
Clear[x,y,z,u,v];
x=Sinh[v]/(Cosh[v]-Cos[u]);
y=Sin[u]/(Cosh[v]-Cos[u]);
z=(y^2+6x-x^2)/6;
ParametricPlot3D[{x,y,z},{u,0,2*Pi},{v,-Log[3],Log[3]},
PlotStyle->Orange,Mesh->20,PlotPoints->95,
MeshStyle->Gray,PlotRange->{{-4,10},{-4,4},{-4,4.2}},BoxRatios->{1,1,1/2}]
\end{comment}

\vskip0.2cm
\item El sólido acotado por el cilindro elíptico $2x^2+y^2=16$ y los parabo- \\ loides hiperbólicos
$x^2-y^2= 4z+12$ y $x^2-2y^2=4z-20$. 

%\vspace{-4mm}
\vskip0.2cm
\noindent
\begin{minipage}{\textwidth}
	\centering
	\captionof{figure}{\small $2x^2+y^2\leq 16$, $ x^2-y^2-12 \leq 4z\leq 20 + x^2-2y^2.$}
	\begin{minipage}{0.17\linewidth}
		\centering
		\includegraphics[width=\linewidth]{Img/T_56a.pdf}
	\end{minipage}
	\hspace{1em}
	\begin{minipage}{0.17\linewidth}
		\centering
		\includegraphics[width=\linewidth]{Img/T_56b.pdf}
	\end{minipage}
	
	\vspace{0.3em}
	\parbox{0.95\linewidth}{\centering\fontsize{8pt}{10pt}\selectfont\textbf{Fuente:} elaboración propia.}
\end{minipage}
\vskip0.2cm
%RegionPlot3D[2*x^2+y^2<=16&&x^2-y^2-12<=4*z<=20+x^2-2*y^2,{x,-3,3},{y,-4,4},{z,-7,7},PlotPoints->90,Mesh->False]
\item El sólido acotado por las figuras del plano $5z+x+y=20$ y el semicono 
$z=\sqrt{x^2+y^2}.$ 


\vskip0.2cm
\noindent
\begin{minipage}{\textwidth}
	\centering
	\captionof{figure}{\small $\sqrt{x^2+y^2} \leq z \leq (20-x-y)/5$}
	\begin{minipage}{0.22\linewidth}
		\centering
		\includegraphics[width=\linewidth]{Img/T_39a.pdf}
	\end{minipage}
	
	\vspace{0.3em}
	\parbox{0.95\linewidth}{\centering\fontsize{8pt}{10pt}\selectfont\textbf{Fuente:} elaboración propia.}
\end{minipage}
\vskip0.2cm



%\begin{comment}
%RegionPlot3D[z>=Sqrt[x^2+y^2]&&5*z<20-y-x,{x,-6,4},{y,-6,4},{z,0,6},PlotPoints->90,Mesh->False]
%\end{comment}

\item La porción del cilindro elíptico con ecuación $4x^2+9y^2=36$ acotado por los planos 
$3z+x=9$ y $3z-y=-9$.

\vskip0.2cm
\noindent
\begin{minipage}{\textwidth}
	\centering
	\captionof{figure}{\small $4x^2+9y^2=36,$ $-9+y \leq 3z \leq 9-x.$}
	\begin{minipage}{0.14\linewidth}
		\centering
		\includegraphics[width=\linewidth]{Img/T_38a.pdf}
	\end{minipage}
	
	\vspace{0.3em}
	\parbox{0.95\linewidth}{\centering\fontsize{8pt}{10pt}\selectfont\textbf{Fuente:} elaboración propia.}
\end{minipage}
\vskip0.2cm




%\begin{comment}
%\begin{Verbatim}
%a=ParametricPlot3D[{{3*Cos[t],2*Sin[t],3-Cos[t]},{3*Cos[t],2*Sin[t],-3+Sin[t]}},
%{t,0,2*Pi},PlotStyle->{Directive[Black,Thickness->0.015],Directive[Black,Thickness->0.015]}];
%b=ParametricPlot3D[(1-u)*{3*Cos[t],2*Sin[t],3-Cos[t]}+u{3*Cos[t],2*Sin[t],-3+Sin[t]},
%{t,0,2*Pi},{u,0,1},MeshStyle->Gray];
%Show[a,b]
%\end{Verbatim}
%\end{comment}
\end{enumerate}
 
\chapter{Superficies y sus curvas de intersección} \chapter{Superficies y sus curvas de intersección}

En la sección \secref{regionesparametrizadas} se presentaron algunas generalidades
sobre parametrizaciones básicas.
\section{Curvas parametrizadas} \index{Curva!parametrizada}
De acuerdo a \cite{thomas_calculo_2015}, consideremos que $f$ y $g$ son funciones continuas en $[a,b]$.
Si $x=f(t) $ y $y=g(t)$, el conjunto de puntos $(x,y)$ que satisfacen estas ecuaciones se denomina curva paramétrica
(o parametrizada) en $\mathbb{R}^2$ y las ecuaciones que la definen se denominan las ecuaciones paramétricas
de la curva. Si además $f'(t)$ y $g'(t)$ no son simultaneamente cero, la curva se denomina suave.
\index{Curva!suave}
\section{Longitud de curva} \index{Longitud de curva}
Para la función  vectorial $\ve{r}(t)$ usaremos la notación
$\overrightarrow{r}(t)$ o ${\bf r}(t)$, indistintamente para referirnos a
este tipo de funciones.

De acuerdo con \cite{marsden_calculo_2013}, la longitud de arco $L$ de la curva $C$ en $\mathbb{R}^3$ determinada por la función vectorial
$\overrightarrow{r}(t)=\left\langle x(t) ,y(t),z(t) \right\rangle$ para $t\in \lbrack a,b]$, está dada por
\begin{equation} \label{curvas}
L=\int_{a}^{b} \left\| \overrightarrow{r}^{\prime }(t)\right\| \, dt.
\end{equation}
La expresión anterior es un caso particular de la fórmula (\ref{intdelinea}), considerando $f(x,y,z)=1$.
Al término $\left\| \overrightarrow{r}^{\prime }(t)\right\|dt$ lo denominaremos elemento diferencial de longitud.
\index{Elemento diferencial!de longitud}
\\
Recordemos que la integral del campo escalar $f\left( x,y,z\right) $ a lo largo de la curva $C$,
está determinada por
$$\int_{C}fds=\int_{a}^{b}f\left( x\left( t\right) ,y\left( t\right) ,z\left(
t\right) \right) \left\| \overrightarrow{r}^{\prime }\left( t\right)\right\| dt.$$

Presentaremos enseguida algunos ejemplos de superficies, su curva de intersección y sus longitud.

\subsection{Un paraboloide y un plano}
La ecuación $-x+y+z=1$ representa un plano y la expresión $2z=x^{2}+y^{2}$
representa un paraboloide. \label{curvas1}

\begin{figure}[H]
\begin{center}
  \caption{Gráfico de $2z=x^{2}+y^{2}$ y $-x+y+z=1$}
  \includegraphics[height=6cm]{Img/72a.pdf} \hskip0.5cm
  \includegraphics[height=6cm]{Img/72b.pdf}
  \fuentepropia
\end{center}
\end{figure}

Con las siguientes instrucciones, se obtienen las figuras anteriores:

\begin{tcolorbox}[title=\bf Instrucción \verb"Mathematica"]
\begin{Verbatim}[formatcom=\letraverbatim]
a=ContourPlot3D[{2*z==x^2+y^2,-x+y+z==1},{x,-5,4},
{y,-4,5},{z,0,6},ContourStyle->{Orange,Cyan},
PlotPoints->60,Mesh->False,Ticks->{{-3,0,3},
{-3,0,3},{1,2,4,5}}];
b=ParametricPlot3D[{1+2*Cos[t],-1+2*Sin[t],3+2*Cos[t]
-2*Sin[t]},{t,0,2*Pi},
PlotStyle->{Black,Thickness->0.016}];
Show[a,b]
  \end{Verbatim}
\end{tcolorbox}

Al remplazar $z$ de la ecuación del paraboloide en la ecuación del plano se obtiene: $\left(\frac{x^{2}+y^{2}}{2}\right)-x+y=1$, que se reescribe como
$\left( x-1\right) ^{2}+\left( y+1\right) ^{2}=4$. Esta última expresión representa una circunferencia con centro en el punto $\left( 1,-1\right) $ y
radio $2$ en el plano $xy$. Esta curva se puede parametrizar en la forma
$x-1=2\cos (t) $, $y+1=2\sin (t) $.
Al sustituir $z$ en la ecuación del paraboloide, se obtiene:
\[
z = \frac{(1+2\cos(t))^{2}+\left(-1+2\sin(t)\right) ^{2}}{2}= 3+2\cos(t) -2\sin (t) .
\]
Por tanto, la función vectorial
$$\overrightarrow{r}(t)=\left\langle 1+2\cos (t) ,-1+2\sin \left(
t\right) ,3+2\cos (t) -2\sin (t) \right\rangle ,\,\,
t\in \lbrack 0,2\pi ]$$
describe la curva de intersección entre el paraboloide y el plano.
Además $\left\| \overrightarrow{r}^{\prime }(t)\right\| = 2\sqrt{2\sin (t) \cos (t) +2}$,
la longitud de la curva está dada por
$$\int_{0}^{2\pi }2\sqrt{2\sin (t) \cos (t) +2}dt\approx 17.475.$$
El cálculo anterior se puede evaluar numéricamente como sigue:

\begin{tcolorbox}[title=\bf Instrucción \verb"Mathematica"]
\begin{Verbatim}[formatcom=\letraverbatim]
r={1+2*Cos[t],-1+2*Sin[t],3+2*cos[t]-2Sin[t]}
A=Assuming[t>=0,Simplify[Norm[D[r,t]]]]
NIntegrate[A,{t,0,2*Pi}]
\end{Verbatim}
\end{tcolorbox}

\subsection{Un paraboloide y un cilindro}
Las ecuaciones $x-y^{2}+4y-z^{2}+2z=2 $ y $y^{2}-4y+4+z^{2}-4z=0$, se reescriben como
$x+3=(y-2)^2+(z-1)^2$ y $(z-2)^2+(y-2)^2=4$.

Estas expresiones representan un paraboloide
con vértice en el punto $(-3,2,1)$, que abre hacia la parte positiva del eje $x$ y un cilindro de
radio $2$ cuyo eje de simetría es una recta paralela al eje $x$. \label{curvas2}

\begin{figure}[H]
\centering
\caption{Gráfico de $x-y^{2}+4y-z^{2}+2z=2$ y $y^{2}-4y+4+z^{2}=4z$}
\includegraphics[width=0.4\textwidth]{Img/70a.pdf} \hspace{5mm}
\includegraphics[width=0.4\textwidth]{Img/70b.pdf}
\fuentepropia
\end{figure}

Las instrucciones para conseguir el gráfico son las siguientes:

\begin{tcolorbox}[title=\bf Instrucción \verb"Mathematica"]
\begin{Verbatim}[formatcom=\letraverbatim]
a=ContourPlot3D[{x+3==(y-2)^2+(z-1)^2,(z-2)^2+(y-2)^2
==4},{x,-3,7},{y,-2,6},{z,-2.2,4.5},Mesh->False,
ContourStyle->{Yellow,RGBColor[0.8,0.3,0.2]},
PlotPoints->60];
b=ParametricPlot3D[{2+4*Sin[t],2+2*Cos[t],
2+2*Sin[t]},
{t,0,2*Pi},PlotStyle->{Black,Thickness->0.014},
PlotPoints->80];
Show[a,b]
\end{Verbatim}
\end{tcolorbox}

Eliminando de las ecuaciones a $y$, se obtiene que $x=2z-2$, que corresponde a la ecuación
de un plano, la curva de intersección está contenida en este. A partir de la ecuación del cilindro, se puede establecer que $y=2+2\cos(t)$ y $z=2+2\sin(t)$
con lo cual $x=2(2+2\sin(t))-2=2+4\sin(t)$.
Entonces, la función vectorial
\[
{\bf r}(t) = \langle 2+4\sin(t), 2+2\cos(t),2+2\sin(t)\rangle,
\]
con $t \in [0,2\pi], $ parametriza la curva de intersección del paraboloide y el cilindro.
A partir de esta función, se sigue que la magnitud del vector tangente es
\[
\left\| {\bf r}^{\prime }(t)\right\| = 2\sqrt{4\cos ^{2}(t) +1}= 2\sqrt{3+2\cos(2t);}
\]
entonces, la longitud de la curva se determina por 
\[
\int_{0}^{2\pi }2\sqrt{3+2\cos(2t)}\,dt \approx 21.081.
\]
Este cálculo se puede obtener numéricamente con las siguientes sentencias:

\begin{tcolorbox}[title=\bf Instrucción \verb"Mathematica"]
\begin{Verbatim}[formatcom=\letraverbatim]
r={2+4*Cos[t],2+2*Cos[t],2+2*Sin[t]};
A=Assuming[t>=0,Simplify[Norm[D[r,t]]]]
NIntegrate[A,{t,0,2*Pi}]
\end{Verbatim}
\end{tcolorbox}

\subsection{Un cilindro y un paraboloide hiperbólico}
El cilindro con ecuación $x^{2}+y^{2}=4$ y el paraboloide hiperbólico con ecuación
$z=x^{2}-y^{2}$ se intersecan en una curva como lo muestra las siguientes figuras.

\begin{figure}[H]
\centering
\caption{Gráfico de $x^{2}+y^{2}=4$ y $z=x^{2}-y^{2}$}
\label{fig:67}
\includegraphics[width=0.4\textwidth]{Img/67a.pdf} \hspace{5mm}
\includegraphics[width=0.4\textwidth]{Fig/301.pdf}
\fuentepropia
\end{figure}

A partir de la ecuación del cilindro se puede definir $x=2\cos \left(t\right) $ y
$\,y=2\sin (t) $ para reemplazar en la ecuación del
paraboloide hiperbólico
$$z=\left( 2\cos (t) \right)^{2}-\left(2\sin(t)\right) ^{2}= 4\cos(2t) ,$$
es decir, la función vectorial
$\overrightarrow{r}\left(t\right) =\left\langle 2\cos (t) ,2\sin (t) ,4\cos\left( 2t\right) \right\rangle $,
para $t\in \left[ 0,2\pi \right] $, define
la curva de intersección de las superficies se\~{n}aladas.
\\
A partir de $\overrightarrow{r}(t) $, se obtiene
$\overrightarrow{r}^{\prime }(t) =\left\langle -2\sin \left(
t\right) ,2\cos (t) ,-8\sin \left( 2t\right) \right\rangle $  y
{\small $$\left\| \overrightarrow{r}^{\prime }(t) \right\| =
2\sqrt{64\cos ^{2}(t) +1-64\cos ^{4}(t)}= 2\sqrt{9-8\cos \left( 4t\right)}.$$
}
Con lo cual, la longitud de esta curva está dada por
{\small $$\int_{0}^{2\pi }2\sqrt{9-8\cos \left( 4t\right) }dt \approx 35.\, 2571.$$
}
El gráfico y la evaluación numérica de esta integral se obtienen en \linebreak
\verb"Mathematica" mediante el uso se las siguientes instrucciones:

\begin{tcolorbox}[title=\bf Instrucción \verb"Mathematica"]
\begin{Verbatim}[formatcom=\letraverbatim]
r={2*Cos[t],2*Sin[t],4*Cos[2*t];Show[ContourPlot3D[
{x^2+y^2==4,z==x^2-y^2},
{x,-4,4},{y,-4,4},{z,-5,5},
Mesh->False,PlotPoints->50,
ContourStyle->{Orange,RGBColor[0.1,0.6,0.2]}],
ParametricPlot3D[{2*Cos[t],2*Sin[t],4*Cos[2*t]},
{t,0,2*Pi},PlotStyle->{Black,Thickness[0.016]}]]
A=Assuming[t>=0,Simplify[Norm[D[r,t]]]]
NIntegrate[A,{t,0,2*Pi}]
\end{Verbatim}
\end{tcolorbox}
%\vspace*{-\baselineskip}
%

\subsection{Un cilindro parabólico y un paraboloide}
Las ecuaciones $z=x^2+y^2$ y $z=5-y^2$, que representan un paraboloide y un cilindro parabólico, respectivamente.

\begin{figure}[H]
\centering
\caption{Gráfico de $z=x^2+y^2$ y $z=5-y^2$}
\includegraphics[width=0.4\textwidth]{Img/50a.pdf} \hspace{5mm}
\includegraphics[width=0.4\textwidth]{Img/50b.pdf}
\fuentepropia
\end{figure}

Las figuras de estas superficies y su curva de intersección se consiguen con las siguientes instrucciones:

\begin{tcolorbox}[title=\bf Instrucción \verb"Mathematica"]
\begin{Verbatim}[formatcom=\letraverbatim]
A=ContourPlot3D[{z==5-y^2,z==x^2+y^2},{x,-3,3},
{y,-3,3},{z,0,5},
PlotPoints->90,ContourStyle->{Cyan,Orange},
Mesh->False]r={Sqrt[5]*Cos[t],Sqrt[5/2]*Sin[t],
(5/2)*(1+Cos[t]^2)};
B=ParametricPlot3D[r,{t,0,2Pi},
PlotStyle->{Black,
Thickness->0.012}];
Show[A,B]
\end{Verbatim}
\end{tcolorbox}

Con las ecuaciones $z=x^{2}+y^{2}$ y $z=5-y^{2}$, se pueden restar y se obtiene
$x^{2}+2y^{2}=5;$ esta expresión representa una elipse con centro en el origen y alargada en el sentido de las $x$.
La curva de intersección está dada en el plano $xy$ por
$x=\sqrt{5}\cos t$, $y=\sqrt{\frac{5}{2}}\sin t$, donde $t\in \lbrack 0,2\pi ]$.
Para hallar la coordenada $z$ se sustituye en la ecuación del paraboloide
$$z=\left( \sqrt{5}\cos t\right) ^{2}+\left( \sqrt{\frac{5}{2}}
\sin t\right) ^{2}= \frac{5}{2}\cos ^{2}t+\frac{5}{2},$$
o también en la ecuación del cilindro
$z=5-\left( \sqrt{\frac{5}{2}}\sin t\right) ^{2}= \frac{5}{2}\cos^{2}t+\frac{5}{2};$
por lo tanto, la curva de intersección está determinada por
$$\overrightarrow{r}(t)=\left\langle \sqrt{5}\cos t,\sqrt{\frac{5}{2}}\sin t,%
\frac{5}{2}\cos ^{2}t+\frac{5}{2}\right\rangle ,\,\,t\in \lbrack 0,2\pi ].$$
Además
$\left\| \overrightarrow{r}^{\prime}(t)\right\| =\frac{\sqrt{10}}{4}\sqrt{11-2\cos(2t)-5\cos(4t)}$,
entonces la longitud de la curva está dada por
$$L=\int_{0}^{2\pi}\frac{\sqrt{10}}{4}\sqrt{11-2\cos(2t)-5\cos(4t)} \,dt \approx 16.192.$$
La integral anterior se evaluó numéricamente mediante las siguientes tres instrucciones:
%\vspace*{-1.2\baselineskip}
%
\begin{tcolorbox}[title=\bf Instrucción \verb"Mathematica"]
\begin{Verbatim}[formatcom=\letraverbatim]
r={Sqrt[5]*Cos[t],Sqrt[5/2]*Sin[t],(5/2)*(1+Cos[t]^2)};
A=Assuming[t>=0,Simplify[Norm[D[r,t]]]]
NIntegrate[A,{t,0,2Pi}]
\end{Verbatim}
\end{tcolorbox}
%\vspace*{-1.2\baselineskip}
%
\subsection{Dos cilindros parabólicos}

La región del espacio acotado por las figuras de las expresiones
$x=z^{2}$ y $x=4y-y^{2}$ \label{curvas3} que corresponden a dos cilindros parabólicos,
junto con la curva de intersección se muestran a continuación.

\begin{figure}[H]
\centering
\caption{Gráfico de $x=z^{2}$ y $x=4y-y^{2}$}
\includegraphics[width=0.4\textwidth]{Img/53a.pdf} \hspace{5mm}
\includegraphics[width=0.4\textwidth]{Img/53e.pdf}

\includegraphics[width=0.4\textwidth]{Img/53f.pdf}
\fuentepropia
\end{figure}

La figura se obtiene con ayuda de las siguientes instrucciones:

\begin{tcolorbox}[title=\bf Instrucción \verb"Mathematica"]
\begin{Verbatim}[formatcom=\letraverbatim]
A=ContourPlot3D[{x==z^2,x==4*y-y^2},{x,-1,5},{y,-1,5},
{z,-3,3},Mesh->False,PlotPoints->50,ContourStyle->
{Blend[{Green,Yellow},2/3],RGBColor[1,0.3,0.3]}];
B=ParametricPlot3D[{4*Sin[t]^2,2+2*Cos[t],2*Sin[t]},
{t,0,2*Pi},PlotStyle->{Black,Thickness[0.016]}];
Show[A,B]
\end{Verbatim}
\end{tcolorbox}

Con estas expresiones se puede reemplazar el valor de $x$ para determinar su proyección en el plano $yz$.
De este procedimiento se obtiene la ecuación  $(y-2)^2+z^2=4$, que corresponde a una circunferencia de radio $2$ centrada en el punto $(2,0)$.  En este caso puede hacerse $y=2+2\cos(t)$ y $z=2\sin(t)$.
La expresión para $x$ se obtiene sustituyendo $y$ y $z$ en cualquiera de las expresiones que determinan los cilindros.
Es decir, $x=4(2+2\cos t) -( 2+2\cos t) ^{2}= 4-4\cos^{2}(t) =4\sin ^{2}(t) $ o también
$x=( 2\sin t) ^{2}= 4\sin ^{2}(t) $.
Por lo tanto, la curva de intersección está dada por la siguiente función vectorial
$${\bf r}(t)=\langle4\sin^{2}(t),2+2\cos(t),2\sin(t)\rangle, \, t \in [0,2\pi].$$
La longitud de esta curva se calcula integrando la norma de la derivada de la función vectorial, es decir
$${\bf r}'(t)= \langle 8\sin(t)\cos(t),-2\sin (t),2\cos(t)\rangle .$$
Su norma está dada por
$$\|{\bf r}^{\prime}(t)\| = 2\sqrt{16\cos^{2}(t)+1-16\cos^{4}(t)}= 2\sqrt{3-2\cos(4t)},$$
este término se integra numéricamente
$$ \int_{0}^{2\pi}2\sqrt{3-2\cos(4t)}dt \approx 21.081.$$
La evaluación se consigue con ayuda de las siguientes instrucciones:

\begin{tcolorbox}[title=\bf Instrucción \verb"Mathematica"]
\begin{Verbatim}[formatcom=\letraverbatim]
r={4*Sin[t]}^2,2+2*Cos[t],2*Sin[t;]
A=Assuming[t>=0,Simplify[Norm[D[r,t]]]]
NIntegrate[A,{t,0,2Pi}]
\end{Verbatim}
\end{tcolorbox}

\subsection{Dos paraboloides}
Las ecuaciones $8y=9-5x^{2}-6z^{2}$ y $16y=4z^{2}-x^{2}$
representan gráficamente un paraboloide y un paraboloide hiperbólico, respectivamente.

\begin{figure}[H]
\centering
\caption{Gráfico de $8y=9-5x^{2}-6z^{2}$ y $16y=4z^{2}-x^{2}$}
\includegraphics[width=0.4\textwidth]{Img/71c.pdf} \hspace{5mm}
\includegraphics[width=0.4\textwidth]{Img/71a.pdf}
\fuentepropia
\end{figure}

Los gráficos se consiguen ejecutando las siguientes instrucciones:

\begin{tcolorbox}[title=\bf Instrucción \verb"Mathematica"]
\begin{Verbatim}[formatcom=\letraverbatim]
a=ContourPlot3D[{8*y+5*x^2+6*z^2==9,16*y==4*z^2-x^2},
{x,-3,3},{y,-1,2},{z,-2,2},Mesh->False,PlotPoints->60,
ContourStyle->{Yellow,RGBColor[0.6,0.2,0.2]},
Ticks->{{-2,0,2},{-1,0,1},{-2,-1,1,2}}];
b=ParametricPlot3D[{Sqrt[2]*Cos[t],(5-13*Cos[2*t])/64,
3*Sqrt[2]/4*Sin[t]},{t,0,2*Pi},PlotPoints->80,
PlotStyle->{Black,Thickness->0.014}];
Show[a,b]
\end{Verbatim}
\end{tcolorbox}

Al eliminar $y$ de estas expresiones se obtiene la ecuación
$9x^{2}+16z^{2}=18$, que corresponde a una elipse en el plano $xz$, esta curva se puede
parametrizar en la forma $x=\sqrt{2}\cos (t) $, $z=\frac{3\sqrt{2}}{4}\sin \left(
t\right) $.
Al reemplazar estas expresiones en cualquiera de las ecuaciones
del paraboloide o del paraboloide hiperbólico se obtiene:
$$8y=9-5x^{2}-6z^{2}= 9-10\cos ^{2}(t) -\frac{27}{4}\sin ^{2}(t)= \frac{9}{4}-\frac{13}{4}\cos^{2}( t),$$
entonces, $y=\frac{9}{32}-\frac{13}{32}\cos ^{2}(t) $.
Por lo tanto, la curva de intersección se parametriza como
$$\overrightarrow{r}(t) =\left\langle \sqrt{2}\cos \left(
t\right) ,\frac{9}{32}-\frac{13}{32}\cos ^{2}(t) ,\frac{3\sqrt{2}}{4}\sin (t) \right\rangle ,\,\, t\in \lbrack 0,2\pi ].$$
Es sencillo ver que
$\left\| \overrightarrow{r}^{\prime }(t)\right\|= \frac{1}{16}\sqrt{512-55\cos ^{2}(t) -169\cos ^{4}(t)} $,
con lo cual, la longitud de la curva es
{\small $$\int_{0}^{2\pi }\frac{1}{16}\sqrt{512-55\cos ^{2}(t) -169\cos^{4}(t) }dt \approx 8.019.$$}

Las siguientes instrucciones permiten hacer este cálculo numéricamente:

\begin{tcolorbox}[title=\bf Instrucción Mathematica]
\begin{Verbatim}[formatcom=\letraverbatim]
r={Sqrt[2]*Cos[t],9/32-13/32*Cos[t]^2,3*Sqrt[2]/
4*Sin[t]};
A=Assuming[t>=0,Simplify[Norm[D[[r,t]]]]
NIntegrate[A,{t,0,2*Pi}]
\end{Verbatim}
\end{tcolorbox}

\subsection{Dos cilindros elípticos y un paraboloide}

Los cilindros elipticos con
ecuaciones $4x^2+y^2=36$ y $x^2+4y^2=4$, poseen al eje $z$ como eje de simetría.
Estos se pueden proyectar en el plano $xy$, como dos elipses que se parametrizan respectivamente como
$\overrightarrow{u}(t)=\left\langle 3\cos (t) ,6\sin (t) \right\rangle $
y  $\overrightarrow{v}(t)=\left\langle 2\cos (t) ,\sin (t) \right\rangle $,
con  $t\in\lbrack 0,2\pi ]$.

Si estos cilindros los intersecan con el paraboloide determinado por la ecuación
$z=6-\left(x-1\right)^2-y^2$,
se obtienen dos curvas, cuya coordenada $z$ está dada, respectivamente por
{\small
\begin{eqnarray*}
z&=&6-\left( 3\cos (t) -1\right) ^{2}-\left( 6\sin (t) \right) ^{2}= -31+27\cos ^{2}(t) +6\cos (t) \\
z &=& 6-\left( 2\cos (t) -1\right) ^{2}-\left( \sin (t) \right) ^{2}= 4-3\cos ^{2}(t) +4\cos (t).
\end{eqnarray*}
}

\begin{figure}[H]
%
\centering
\caption{Gráfico de $4x^2+y^2=36$ y $x^2+4y^2=4$}
\begin{minipage}[t]{0.34\linewidth}
\centering
\includegraphics[width=0.69\linewidth]{Img/1f.pdf}
\end{minipage}%
\hspace{0.2em}
\begin{minipage}[t]{0.34\linewidth}
\centering
\includegraphics[width=0.69\linewidth]{Fig/78.pdf}
\end{minipage}
\fuentepropia
\end{figure}

Con las siguientes instrucciones se consigue graficar estas superficies:

\begin{tcolorbox}[title=\bf Instrucción \verb"Mathematica"]
\begin{Verbatim}[formatcom=\letraverbatim]
ContourPlot3D[{z==6-(x-1)^2-y^2,4*x^2+y^2==36,
x^2+4*y^2==4},{x,-6,10},{y,-7,7},{z,-35,6},
ContourStyle->{Red,Yellow,Green},PlotPoints->80]
\end{Verbatim}
\end{tcolorbox}

En el espacio estas las dos curvas de intersección se pueden parametrizar como
\begin{eqnarray*}
\overrightarrow{\alpha }(t) &=& \left\langle 3\cos (t) ,6\sin
(t) ,-31+27\cos ^{2}(t) +6\cos (t) \right\rangle \\
\overrightarrow{\beta }(t) &=& \left\langle 2\cos (t) ,\sin \left(
t\right) ,4-3\cos ^{2}(t) +4\cos (t) \right\rangle,
\end{eqnarray*}
con $t \in [0,2\pi]$.
La longitud de dichas curvas se puede determinar con el
siguiente procedimiento:
\begin{eqnarray*}
\overrightarrow{\alpha }^{\prime }(t)&=& \left\langle -3\sin (t),6\cos (t) ,-54\cos (t) \sin (t)
-6\sin (t) \right\rangle \\
\left\| \overrightarrow{\alpha }^{\prime }(t)\right\|
&=&  3\sqrt{5+\cos (t) \left( 323\cos(t) -324\cos ^{3}(t) +72\sin ^{2}(t)\right),}
\end{eqnarray*}
con lo cual, la longitud de la curva mayor está dada por
$$\int_{0}^{2\pi}3\sqrt{5+\cos (t) \left( 323\cos(t) -324\cos ^{3}(t) +72\sin ^{2}(t)\right) }  \,  dt\approx 115. 268;$$
por otro lado,
\begin{eqnarray*}
\overrightarrow{\beta }^{\prime }(t)&=& \left\langle -2\sin \left(
t\right) ,\cos (t) ,6\cos (t) \sin (t)-4\sin (t) \right\rangle \\
\left\| \overrightarrow{\beta }^{\prime }(t)\right\|
&=& \sqrt{20+\cos (t) \left( 17\cos (t)-36\cos ^{3}(t) -48\sin ^{2}(t) \right).}
\end{eqnarray*}
Entonces, la longitud de la segunda curva está dada por
$$\int_{0}^{2\pi}\sqrt{20+\cos (t) \left( 17\cos (t)-36\cos ^{3}(t) -48\sin ^{2}(t) \right)} \, dt\approx 21.161.$$
En los dos casos, estas integrales se evaluaron en \verb"Mathematica"
usando las siguientes instrucciones:

\begin{tcolorbox}[title=\bf Instrucción \verb"Mathematica"]
\begin{Verbatim}[formatcom=\letraverbatim]
r={3*Cos[t],6*Sin[t]};
z=6-(r[[1]]-1)^2-r[[2]]^2;
r={r[[1]],r[[2]],z};
NIntegrate[Norm[D[r,t]],{t,0,2*Pi}]
\end{Verbatim}
\end{tcolorbox}

\begin{tcolorbox}[title=\bf Instrucción \verb"Mathematica"]
\begin{Verbatim}[formatcom=\letraverbatim]
r={2*Cos[t],Sin[t]};
z=6-(r[[1]]-1)^2-r[[2]]^2;
r={r[[1]],r[[2]],z};
NIntegrate[Norm[D[r,t]],{t,0,2*Pi}]
\end{Verbatim}
\end{tcolorbox}

\subsection{Un elipsoide y un hiperboloide} \index{Hiperboloide}

Las ecuaciones $2x^2+y^2=z^2+12$ y $x^2+y^2+8z^2=30$ representan un hiperbólico de un manto y un elipsoide, respectivamente.

\begin{figure}[H]
\begin{center}
\caption{Gráfico de $2x^2+y^2=z^2+12$ y $x^2+y^2+8z^2=30$}
\includegraphics[height=4.5cm]{Img/87a.pdf}
\includegraphics[height=4.5cm]{Img/87b.pdf}
\fuentepropia
\end{center}
\end{figure}

La figura anterior se obtiene con las siguientes instrucciones:

\begin{tcolorbox}[title=\bf Instrucción \verb"Mathematica"]
\begin{Verbatim}[formatcom=\letraverbatim]
A=ContourPlot3D[{z^2+12==2*x^2+y^2,x^2+y^2+8*z^2==30},
{x,-6,6},{y,-6,6},{z,-3,3},Mesh->False,PlotPoints->40,
ContourStyle->{Directive[Orange,Opacity[0.8]],
Directive[Cyan,Opacity[0.76]]},BoxRatios->{6,5,4}];
B=ParametricPlot3D[{{3*Sqrt[2]*Sinh[t],Sqrt[2]*Sqrt[6
+Cosh[t]^2-18*Sinh[t]^2],Sqrt[2]*Cosh[t]},
{3*Sqrt[2]*Sinh[t],-Sqrt[2]*Sqrt[6+Cosh[t]^2-
18*Sinh[t]^2],Sqrt[2]*Cosh[t]}},
{t,-ArcCosh[Sqrt[27/17]],ArcCosh[Sqrt[27/17]]},
PlotStyle->{Directive[Black,Thickness->0.015],
Directive[Black,Thickness->0.015]}];
Show[A,B]
\end{Verbatim}
\end{tcolorbox}

Al restar las ecuaciones del elipsoide y del hiperboloide se elimina a $y$,
esto se reduce a $x^{2}-9z^{2}=-18$, que
equivale a $\frac{z^{2}}{2}-\frac{x^{2}}{18}=1;$ esta ecuación representa
una hipérbola y se parametriza como
$z=\sqrt{2}\cosh \left( t\right) $, $x=3\sqrt{2}\sinh \left( t\right) $.
La coordenada $y$ se obtiene reemplazando $x$ y $z$ en cualquiera de las
ecuaciones iniciales (la del elipsoide o la del hiperboloide). Lo cual implica que
$y=\pm \sqrt{48-34\cosh ^{2}t}$.
Para determinar el dominio, puede verse que si $48-34\cosh ^{2}t=0$, entonces la solución está dada por
$a={\rm arccosh}\left( \frac{2\sqrt{102}}{17}\right) \approx 0.7066$.

Por lo descrito anteriormente, la curva de intersección se parametriza por
$$\overrightarrow{r}\left( t\right) =\left\langle 3\sqrt{2}\sinh \left(
t\right) ,\sqrt{48-34\cosh ^{2}t},\sqrt{2}\cosh \left( t\right)
\right\rangle ,$$
con $t \in[-a,a]$, con lo cual, la longitud de la curva está dada por
$$\bigint_{-a}^{a}\sqrt{\frac{2\left(32\cosh ^{2}t-119\cosh^{4}t+24\right)}{17\cosh ^{2}t-24}}dt\approx 30.5672.$$
Los cálculos que permiten obtener la longitud de la curva
se obtienen ejecutando en \verb"Mathematica" las siguientes instrucciones:

\begin{tcolorbox}[title=\bf Instrucción \verb"Mathematica"]
\begin{Verbatim}[formatcom=\letraverbatim]
z=Sqrt[2]*Cosh[t],Sqrt[2]*Cosh[t]}
x=3*Sqrt[2]*Sinh[t];
y=Sqrt[2]*Sqrt[6+Cosh[t]^2-
18*Sinh[t]^2];
a=ArcCosh[Sqrt[27/17]];
A=Assuming[t>=0,Simplify[Norm[D[{r,t},]]]];
NIntegrate[2*A,{t,-a,a}]
\end{Verbatim}
\end{tcolorbox}

\subsection{Una esfera y dos cilindros perpendiculares}

La esfera centrada en el origen de radio $2$ se interseca con los cilindros cuyas ecuaciones son $x^{2}+y^{2}=2x$ y $x^{2}+z^{2}=2x$.
Hallaremos la curva de intersección y de esta calcularemos su longitud de arco. La siguiente figura muestra dicha situación.

\begin{figure}[H]
\centering
\caption{Gráfico de $x^{2}+y^{2}=2x$ y $x^{2}+z^{2}=2x$}
\begin{minipage}[t]{0.38\linewidth}
\centering
\includegraphics[width=0.75\linewidth]{Img/6a.pdf}
\end{minipage}%
\hspace{0.1em}
\begin{minipage}[t]{0.40\linewidth}
\centering
\includegraphics[width=0.80\linewidth]{Fig/79.pdf}
\end{minipage}
\fuentepropia
\end{figure}

Ejecutando las siguientes instrucciones se obtiene la figura:

\begin{tcolorbox}[title=\bf Instrucción \verb"Mathematica"]
\begin{Verbatim}[formatcom=\letraverbatim]
ContourPlot3D[{x^2+y^2+z^2==4,x^2+y^2==2*x,
x^2+z^2==2*x},{x,-2,2},{y,-2,2},{z,-2,2},
Mesh->False,ContourStyle->{Orange,Yellow,Gray}]
\end{Verbatim}
\end{tcolorbox}

El cilindro $x^{2}+y^{2}=2x$, se reescribe como $\left( x-1\right)^{2}+y^{2}=1$,
la curva de intersección se parametriza al hacer $x=1+\cos (t)$, $y=\sin (t)$, y para hallar $z$
es necesario reemplazar en la ecuación de la esfera, de donde se obtiene que

\[x^{2}+y^{2}+z^{2}= \left( 1+\cos (t) \right)
^{2}+\sin ^{2}(t) +z^{2}= 4,\]
es decir, $2+2\cos(t) +z^{2}=4$,
de donde $z=\pm \sqrt{2-2\cos (t)}$,
entonces la función

$${\bf r}_1(t) =\left\langle 1+\cos (t) ,\sin (t),\sqrt{2-2\cos (t) }\right\rangle ,\,\, t\in \lbrack 0,2\pi ],$$

parametriza la curva de intersección. Es de anotar que son dos porciones que se obtienen por la raiz positiva y la raiz negativa.

Para la expresión  $x^{2}+z^{2}=2x$ (que representa el otro cilindro), el tratamiento es similar, obtiendose la función
$${\bf r}_2(t) =\left\langle 1+\cos (t),\sqrt{2-2\cos (t)},\sin (t)\right\rangle ,$$
donde $t\in \lbrack 0,2\pi ]$.
Por la simetría de la curva y las ecuaciones que la definen
consideraremos solamente ${\bf r}_1(t)$, entonces la longitud ($L$) de la curva es cuatro veces la integral del
término
$$\left\Vert {\bf r}_1^{\prime }(t) \right\Vert =\left\Vert \left\langle -\sin t,\cos t,\frac{1}{2}\frac{\sqrt{2}\sin t}{\sqrt{1-\cos t}}\right\rangle \right\Vert=\frac{1}{2}\sqrt{2}\sqrt{\cos t+3} ,$$
de donde
$$L=4\int_{0}^{2\pi }\frac{1}{2}\sqrt{2}\sqrt{\cos t+3}\,dt \approx 30.562.$$

Este valor se obtuvo en \verb"Mathematica" utilizando las siguientes instrucciones:

\begin{tcolorbox}[title=\bf Instrucción \verb"Mathematica"]\begin{Verbatim}[formatcom=\letraverbatim]
r={1+Cos[t],Sin[t],Sqrt[2-2Cos[t]]};
A=Assuming[t>=0,Simplify[Norm[D[r,t]]]];
NIntegrate[4*A,{t,0,2Pi}]
\end{Verbatim}
\end{tcolorbox}

\subsection{Un cilindro cardióidico y un paraboloide} \label{cilcardiopar}

Consideremos el parabóloide con ecuación $3z=9-(x+2)^{2}-(y+1)^{2}$ y el cilindro cardióidico
$r=1-\cos (t)$. En dos dimensiones, el cardioide se  parametriza como
$\overrightarrow{r}(t)=\left\langle (1-\cos (t))\cos (t),(1-\cos (t))\sin(t)\right\rangle$,
donde $t\in [0,2\pi]$.

\begin{figure}[H]
\centering
\caption{$r=1-\cos t$, $3z=4-4x-x^2-2y-y^2$}
\includegraphics[width=0.4\textwidth]{Img/58a.pdf} \hspace{5mm}
\includegraphics[width=0.4\textwidth]{Img/58b.pdf}
\fuentepropia
\end{figure}

Sobre el paraboloide se cumple que
$z=\frac{9}{2}-2 \sin (t)+\sin (2 t)-2 \cos (t)+\frac{3}{2} \cos (2 t)$.
Por lo tanto, para $t\in [0,2\pi]$, la curva de intersección del paraboloide y el cilindro cardioidico se parametriza por medio de la sigui\-ente función:
\begin{multline*}
\overrightarrow{r}(t)=\big\langle(1-\cos t)\cos t,(1-\cos t)\sin t,\frac{9}{2}-2\sin t+\sin(2t) \\
-2\cos t+\frac{3}{2}\cos(2t)\big\rangle.
\end{multline*}

Puede verse también, que:
\begin{multline*}
\Vert \overrightarrow{r}^\prime(t)\Vert= \frac{\sqrt{2}}{3}\left|\sin\left(\frac{t}{2}\right)\right|\bigg(4(\sin(t)+\sin(2t)+3\sin(3t))+7\cos(t)\\ +6\cos(2t)+5\cos(3t)+32\bigg )^{1/2} ,
\end{multline*}

por tanto:
$$\int_0^{2\pi}\left\Vert \overrightarrow{r}^\prime(t)\right\Vert dt \approx 10.059.$$

Las figuras y los cálculos para determinar la longitud de la curva se obtienen ejecutando las siguientes instrucciones:

\begin{tcolorbox}[title=\bf Instrucción \verb"Mathematica"]\begin{Verbatim}[formatcom=\letraverbatim]
r={(1-Cos[t])*Cos[t],(1-Cos[t])*Sin[t],
1+Cos[t]/3*(3*Cos[t]-2)+2*Sin[t]/3*(Cos[t]-1)};
A=Assuming[t>=0,Simplify[Norm[D[r,t]]]]
NIntegrate[A,{t,0,2*Pi}]
a=ContourPlot3D[{3*z==9-(x+2)^2-(y+1)^2},{x,-5,1},
{y,-4,2},{z,0,3},ContourStyle->{Orange},Mesh->False,
PlotPoints->80];
b=ParametricPlot3D[{r[[1]],r[[2]],u},{t,0,2*Pi},
{u,0,3},Mesh->False,PlotPoints->80,
PlotStyle->{Opacity[0.7],RGBColor[0.1,0.5,0.3]}];
c=ParametricPlot3D[r,{t,0,2*Pi},PlotPoints->80,
PlotStyle->{Black,Thickness->0.015}];
Show[a,b,c]
\end{Verbatim}
\end{tcolorbox}

\subsection{Superficies obtenidas a partir de su curva de intersección}

Para las siguientes funciones vectoriales se construyen dos superficies cuya intersección corresponde
a la curva determinada por la función $\overrightarrow{r}(t)$. Además se halla la longitud ($L$) de la curva para
un intervalo dado.

%

\begin{funciones}[series=ejercicios]
\item \makebox[\linewidth][c]{$\overrightarrow{r}(t)=\la 3t^{3}+t,3t^{2},3t^{3}-t\ra, $ $t \in [-2,2]$.}
\end{funciones}

De las ecuaciones paramétricas se sigue que
$x=3t^{3}+t$, $y=3t^{2}$, $z=3t^{3}-t$.
Entonces, $x-z=2t;$ por tanto, $\left( x-z\right) ^{2}=4t^{2}=\frac{4}{3}\left(
3t^{2}\right) =\frac{4}{3}y$. Además $x+z=6t^{3}=\left( 2t\right) \left( 3t^{2}\right) =\left(
x-z\right) y$.

Por tal razón, utilizaremos las expresiones
$x-z=\frac{4}{3}y$ y $x+z=y\left( x-z\right) $, cuyos gráficos se
presentan a continuación y con éstos también se muestra la curva de intersección.

\begin{figure}[H]
\centering
\caption{Superficies con intersección ${\bf r}(t)=\la 3t^{3}+t,3t^{2},3t^{3}-t\ra$}
\includegraphics[width=0.4\textwidth]{Img/94a.pdf} \hspace{5mm}
\includegraphics[width=0.4\textwidth]{Img/94b.pdf}
\fuentepropia
\end{figure}

Las indicaciones para conseguirlos son las siguientes:

\begin{tcolorbox}[title=\bf Instrucción \verb"Mathematica"]\begin{Verbatim}[formatcom=\letraverbatim]
a=ParametricPlot3D[{3*t^3+t,3*t^2,3*t^3-t},{t,-1,1},
PlotStyle->{Black,Thickness->0.012},PlotPoints->90];
b=ContourPlot3D[{(x-z)^2==4*y/3,x+z==y*(x-z)},
{x,-1,1},{y,-0.5,1},{z,-1,1},Mesh->False,
ContourStyle->{Directive[RGBColor[0.3,0.6,0.3],
Opacity[0.95]],Directive[Orange,Opacity[0.8]]}];
Show[b,a]
\end{Verbatim}
\end{tcolorbox}

A partir de la relación
$\left\| \frac{d}{dt}\la 3t^{3}+t,3t^{2},3t^{3}-t\ra\right\| = \sqrt{2}\left( 9t^{2}+1\right) $, es posible establecer que
la longitud de la arco está dada por
$$L=\int_{-2}^{2}\sqrt{2}\left( 9t^{2}+1\right) dt= 52\sqrt{2}.$$

\begin{funciones}[resume=ejercicios]
\item \makebox[\linewidth][c]{$\overrightarrow{r}(t)=\la \sqrt{t^{3}},\frac{9}{16}t,t^{2}\ra $, $t \in [0,2]$.}
\end{funciones}

A partir de las ecuaciones $x=\sqrt{t^{3}}$, $y=\frac{9}{16}t$, $z=t^{2}$, se sigue que $x^2=t^3$ y $t= \frac{16}{9}y;$
entonces, dos posibles relaciones son $9x^2=16yz$ y $z= \left(\frac{16}{9}y\right)^2$.

Las figuras de estas ecuaciones y la curva de intersección se presentan enseguida:

\begin{figure}[H]
\centering
\caption{Superficies con intersección ${\bf r}(t)=\la \sqrt{t^{3}},\frac{9}{16}t,t^{2}\ra$}
\includegraphics[width=0.4\textwidth]{Img/92a.pdf} \hspace{5mm}
\includegraphics[width=0.4\textwidth]{Img/92b.pdf}
\fuentepropia
\end{figure}

Las instrucciones requeridas para obtener las figuras son las siguientes:

\begin{tcolorbox}[title=\bf Instrucción \verb"Mathematica"]\begin{Verbatim}[formatcom=\letraverbatim]
a=ParametricPlot3D[{t^(3/2),9*t/16, t^2},{t,0,1},
PlotStyle->{Black,Thickness->0.014}];
b=ContourPlot3D[{9*x^2==16*y*z,z==(16*y/9)^2},{x,0,1},
{y,0,1},{z,0,2}, PlotPoints->90,Mesh->False,
ContourStyle->{Directive[Orange,Opacity[0.8]],
Directive[RGBColor[0.5,0.8, 0.5],Opacity[0.85]]}];
Show[b,a]
\end{Verbatim}
\end{tcolorbox}

Además, $\left\| \frac{d}{dt}\la \sqrt{t^{3}},\frac{9}{16}t,t^{2}\ra\right\| = \frac{9}{16}+2t$,
entonces la longitud de la curva está dada por
$$L=\int_0^2 \left ( \frac{9}{16}+2t \right )dt = \frac{41}{8}. $$

\begin{funciones}[resume=ejercicios]
\item \makebox[\linewidth][c] {$\overrightarrow{r}(t)=\la 4t^{3}+2t^{2},\frac{4}{3\sqrt{2}}t,4t^{3}-2t^{2}\ra $, $t \in [-1,1]$.}
\end{funciones}

A partir de las expresiones $x=4t^{3}+2t^{2}$, $y=\frac{4t}{3\sqrt{2}}$, $z=4t^{3}-2t^{2}$, se sigue que
$t=\frac{3\sqrt{2}}{4}y$, $x-z=4t^{2}=4\left( \frac{3y}{4}\sqrt{2}\right)
^{2}= \frac{9}{2}y^{2}$.\\
Además $x+z=8t^{3}=8\left( \frac{3y}{4}\sqrt{2}\right) ^{3}= \frac{27}{4}\sqrt{2}y^{3}$.

Por lo tanto, las figuras de las ecuaciones $x-z= \frac{9}{2}y^{2}$ y $x+z=\frac{27}{4}\sqrt{2}y^{3}$
se muestras a continuación junto con la curva de intersección:

\begin{figure}[H]
\begin{center}
\caption{Superficies con intersección ${\bf r}(t)=\la 4t^{3}+2t^{2},\frac{4}{3\sqrt{2}}t,4t^{3}-2t^{2}\ra$}
\includegraphics[height=4.5cm]{Img/93a.pdf} \hskip0.5cm
\includegraphics[height=4.5cm]{Img/93b.pdf}
\fuentepropia
\end{center}
\end{figure}

Estos figuras se obtienen ejecutando las siguientes instrucciones:

\begin{tcolorbox}[title=\bf Instrucción \verb"Mathematica"]\begin{Verbatim}[formatcom=\letraverbatim]
a=ParametricPlot3D[{4*t^3+2*t^2,4*t/(3*Sqrt[2]),
4*t^3-2*t^2},{t,-2,2},PlotPoints->90,
PlotStyle->{Black,Thickness->0.014}];
b=ContourPlot3D[{x+z==27*Sqrt[2]*y^3/4,x-z==9*y^2/2},
{x,-1,1},{y,-1.5,1},{z,-2,1},PlotPoints->90,
Mesh->False,ContourStyle->{Directive[Orange,
Opacity[0.8]],Directive[RGBColor[0.5,0.8,0.5],
Opacity[0.9]]}];
Show[b,a]
\end{Verbatim}
\end{tcolorbox}

Es sencillo ver que
$ \left\| \frac{d}{dt}\la 4t^{3}+2t^{2},\frac{4t}{3\sqrt{2}},4t^{3}-2t^{2}\ra \right\| = \frac{2}{3}\sqrt{2}\left(
18t^{2}+1\right), $
por lo tanto, la longitud de la curva está dada por
$$L=\int_{-1}^{1}\frac{2}{3}\sqrt{2}\left( 18t^{2}+1\right) dt= \frac{28}{3}\sqrt{2}$$

\begin{funciones}[resume=ejercicios]
\item \makebox[\linewidth][c]
{$\overrightarrow{r}(t) =\la 3t-t^{3},3t^2,3t+t^{3} \ra $, $t \in [-1,1]$.}
\end{funciones}

Para esta función es claro que
$x= 3t-t^{3}$, $y=3t^2$, $z=3t+t^{3}$. De modo que
$x+z=6t$ entonces $(x+z)^2=12y$. Por otro lado, $x^2-z^2=-12t^4$, entonces, $x^2-z^2=-4y^2/3$.
Las figuras de las dos superficies junto con la curva de intersección se muestra a continuación:

\begin{figure}[H]
\centering
\caption{Superficies con intersección ${\bf r}(t)=\la 3t-t^{3},3t^2,3t+t^{3} \ra$}
\includegraphics[width=0.4\textwidth]{Img/91a.pdf} \hspace{5mm}
\includegraphics[width=0.4\textwidth]{Img/91b.pdf}
\fuentepropia
\end{figure}

Las instrucciones para conseguir estas figuras son las siguientes:

\begin{tcolorbox}[title=\bf Instrucción \verb"Mathematica"]\begin{Verbatim}[formatcom=\letraverbatim]
a=ParametricPlot3D[{3*t-t^3,3*t^2,3*t+t^3},
{t,-1.15,1.15},PlotStyle->{Black,Thickness->0.017},
PlotPoints->90];
b=ContourPlot3D[{x^2-z^2+4*y^2/3,(x+z)^2==12*y},
{x,-5,5},{y,-5,5},{z,-5,5},
ContourStyle->{Directive[Orange,Opacity[0.7]],
Directive[RGBColor[0.5,0.8,0.5],Opacity[0.6]]}],
Mesh->False;
Show[b,a]
\end{Verbatim}
\end{tcolorbox}

Para hallar $L$ se determina el elemento diferencial de longitud
$$\left\| \frac{d}{dt}\la 3t-t^{3},3t^2,3t+3t^{3} \ra \right\| = 3 \sqrt{2}(t^2+1),$$

por lo tanto, $$L= \int_{-1}^1 3 \sqrt{2}(t^2+1) dt = 8\sqrt{2}.$$

\begin{funciones}[resume=ejercicios]
\item \makebox[\linewidth][c]
{$\overrightarrow{r}(t) =\la \sec t, \ln (\sec(t)+\tan(t)),\tan(t) \ra $, $t \in [-\frac{\pi}{3},\frac{\pi}{3}]$.}
\end{funciones}

A partir de las funciones componentes y de algunas identidades trigono\-métricas conocidas
se pueden establecer relaciones entre ellas;
dos de éstas funciones son: $x^2=1+z^2$ y $y=\ln(x+z)$.
Los figuras de estas dos superficies junto con la curva de intersección se muestra a continuación:

\begin{figure}[H]
\centering
\caption{Superficies con intersección ${\bf r}(t)=\la \sec t, \ln (\sec(t)+\tan(t)),\tan(t) \ra$}
\includegraphics[width=0.4\textwidth]{Img/95a.pdf} \hspace{5mm}
\includegraphics[width=0.4\textwidth]{Img/95b.pdf}
\fuentepropia
\end{figure}

Las instrucciones para conseguir estas figuras son las siguientes:

\begin{tcolorbox}[title=\bf Instrucción \verb"Mathematica"]\begin{Verbatim}[formatcom=\letraverbatim]
a=ParametricPlot3D[{Sec[t],Log[Sec[t]+Tan[t]],Tan[t]},
{t,-Pi/3,Pi/3},PlotStyle->{Thickness[0.017],Black}];
b=ContourPlot3D[{1+z^2==x^2,y==Log[x+z]},{x,-2,3},
{y,-2,2},{z,-2,2},Mesh->False,
ContourStyle->{Directive[Orange,Opacity[0.95]],
Directive[Cyan,Opacity[0.9]]}];
Show[b,a]
\end{Verbatim}
\end{tcolorbox}

Es claro que
\[
\left\| \overrightarrow{r}^{\prime }(t)\right\| =\left\| \frac{d}{dt}\left\langle \sec (t) ,\ln \left( \sec (t) +\tan (t) \right) ,\tan (t) \right\rangle \right\| =\sqrt{2}\sec ^{2}t\text{,}
\]
entonces, la longitud de la curva está dada por
$$L=\int_{-\pi /3}^{\pi /3}\sqrt{2}\sec ^{2}tdt= 2\sqrt{6}.$$
%\end{enumerate}

\begin{ejer}
\normalfont
Para cada una de las siguientes funciones vectoriales, determine al menos dos
superficies cuya intersección esté dada por $\overrightarrow{r}(t)$.
También halle la longitud de la curva para el intervalo dado.
\begin{enumerate}
\item $ \overrightarrow{r}(t)=\la e^{-t}\sin (t),e^{-t},e^{-t}\cos (t) \ra$, $t \in [0,\pi/2]$.
%$$\left\| \frac{d}{dt}\la e^{-t}\sin (t),e^{-t},e^{-t}\cos (t) \ra \right\| = \sqrt{3}e^{-t}$$
\item $ \overrightarrow{r}(t)=\la e^{-t}\cos \left( 2t\right),e^{-t}\sin \left( 2t\right),e^{-t} \ra $,
$t \in [0,\pi/2]$.
%$$\left\| \frac{d}{dt}\la e^{-t}\sin \left( 2t\right),e^{-t},e^{-t}\cos \left( 2t\right) \ra \right\| = \sqrt{6}e^{-t}$$
\item $\overrightarrow{r}(t)=\la e^{t},e^{t}\sin \left(2t\right) ,e^{t}\cos \left( 2t\right) \ra $,
$t \in [0,\pi/2]$
%$$\left\| \frac{d}{dt}\la e^{t},e^{t}\sin \left( 2t\right),e^{t}\cos \left( 2t\right) \ra \right\| = \sqrt{6}e^{t}$$
\item $ \overrightarrow{r}(t)=\la 3t^{4}+t^{2},\frac{4}{3}\sqrt{6}t^{3},3t^{4}-t^{2}\ra $,
$t \in [-1,1]$.
%$$\left\|\frac{d}{dt}\la 3t^{4}+t^{2},\frac{4}{3}\sqrt{6}t^{3},3t^{4}-t^{2}\ra \right\| =  2t\sqrt{2}\left(6t^{2}+1\right) $$
\item $ \overrightarrow{r}(t)=\la t^{2},\frac{6}{\sqrt{3}}t,\ln t^{3}\ra $,
$t \in [1,e]$.
%$$\left\| \frac{d}{dt}\la t^{2},\frac{6}{\sqrt{3}} t,\ln t^{3}\ra \right\|= \frac{2t^{2}+3}{t}$$
\item $\overrightarrow{r}(t) =\la \frac{1}{3}t^{3}-4t+1,-\frac{4}{t}-2t^{2},2\sqrt{2}\left( t+\ln t^{2}\right) \ra $,
$t \in [1,e]$.
\normalfont
\item $\overrightarrow{r}(t) =\la \cosh(2t), 2t, -\sinh(2t) \ra $, $t\in [-1,1]$.
%$$\left\| \frac{d}{dt}\la \frac{1}{3}t^{3}-4t+1,-\frac{4}{t}-2t^{2},\frac{4}{\sqrt{2}}\left( t+\ln t^{2}\right) \ra \right\| = \frac{\left( t^{2}+2\right) ^{2}}{t^{2}}$$
\item $\overrightarrow{r}(t) =\la t \cos (2t),t \sin (2t),7 t \ra $, $t\in [0,2\pi]$.
\item $\overrightarrow{r}(t) =\la t^2\cos (3t),t^2 \sin (3t),t^2 \ra $, $t\in [0,2\pi]$.
\item $\overrightarrow{r}(t) =\la e^t \cosh(t),e^t \sinh(t),e^t \ra $, $t\in [-1,1]$.
\item $\overrightarrow{r}(t) =\la -\frac{1}{t},\frac{2}{3}t^3,4t+5\ra $, $t\in [1,2]$.
\item $\overrightarrow{r}(t) =\la t,\ln (\cot (t)-\csc (t)),\ln(\sin (t)) \ra $, $t\in [\pi/6,5\pi/6]$.

\end{enumerate}
\end{ejer}

\begin{ejer}\textit {Determine una expresión que defina la curva de intersección de las siguientes superficies y halle su longitud.}

\enlargethispage{\baselineskip}

\begin{enumerate}
\item El hiperboloide $x^2+z^2=4y^2+\frac{1}{2}$ y el paraboloide $x^2+z^2=1-y$.

\begin{figure}[H]
\begin{center}
\caption{Gráfico de $x^2+z^2=4y^2+\frac{1}{2}$ y $x^2+z^2=1-y$}
\includegraphics[height=4.5cm]{Img/T_23a.pdf}\hskip0.5cm
\includegraphics[height=4.5cm]{Img/T_23b.pdf}
\fuentepropia
\end{center}
\end{figure}

\item El hiperboloide $x^2+z^2-2x=y^2-2$ y el cilindro $3x^2+2z^2-4z=34$.

\begin{figure}[H]
\begin{center}
\caption{Gráfico de $x^2+z^2-2x=y^2-2$ y $3x^2+2z^2-4z=34$}
\includegraphics[height=4.5cm]{Img/T_24a.pdf}\hskip0.5cm
\includegraphics[height=4.5cm]{Img/T_24b.pdf}
\fuentepropia
\end{center}
\end{figure}

\item La esfera $x^2+y^2+z^2-2x=31$ y el paraboloide $x^2+y^2+2x=4z-1$.

\begin{figure}[H]
\begin{center}
\caption{Gráfico de $x^2+y^2+z^2-2x=31$ y $x^2+y^2+2x=4z-1$}
\includegraphics[height=4.5cm]{Img/T_25a.pdf}\hskip0.5cm
\includegraphics[height=4.5cm]{Img/T_25b.pdf}
\fuentepropia
\end{center}
\end{figure}

\item El cilindro cardioidico con ecuación {\small $r=1+\cos(t)$ y el plano $x+y+3z=1$.} \index{Cilindro cardiodico}

\vspace*{-4mm}
\begin{figure}[H]
\begin{center}
\caption{Gráfico de $r=1+\cos(t)$ y $x+y+3z=1$}
\includegraphics[height=4.5cm]{Img/T_29a.pdf}\hskip0.25cm
\includegraphics[height=4.5cm]{Img/T_29b.pdf}
\fuentepropia
\end{center}
\end{figure}

\item Las esferas
$x^2 + y^2 + z^2 - 2 x = 4$ y $x^2 + y^2 + z^2 - 2z = 4$.

\begin{figure}[H]
\begin{center}
\caption{Gráfico de $x^2 + y^2 + z^2 - 2 x = 4$ y $x^2 + y^2 + z^2 - 2z = 4$}
\includegraphics[height=4.5cm]{Img/T_72a.pdf}\hskip0.5cm
\includegraphics[height=4.5cm]{Img/T_72b.pdf}
\fuentepropia
\end{center}
\end{figure}

\item El cilindro $r=\sin(2\theta)$ y el plano $x+y+3z=1$.

\begin{figure}[H]
\begin{center}
\caption{Gráfico de $r=\sin(2\theta)$ y $x+y+3z=1$}
\includegraphics[height=4.5cm]{Img/97a.pdf}\hskip0.5cm
\includegraphics[height=4.5cm]{Img/97b.pdf}
\fuentepropia
\end{center}
\end{figure}

\item El cilindro $r=\sin(2\theta)$ y el paraboloide $z=1-x^2-y^2$.

\begin{figure}[H]
\begin{center}
\caption{Gráfico de $r=\sin(2\theta)$ y $z=1-x^2-y^2$}
\includegraphics[height=4.5cm]{Img/98a.pdf}\hskip0.5cm
\includegraphics[height=4.5cm]{Img/98b.pdf}
\fuentepropia
\end{center}
\end{figure}
\end{enumerate}
\end{ejer}
  
\chapter{Integrales dobles y cálculo de áreas} \chapter{Integrales dobles y cálculo de áreas}

Presentamos algunos problemas relativos a las integrales dobles, cuyos principales elementos teóricos se presentaron en la sección \secref{integralesdobles}.
Se ilustrará la manera de plantear sumas de Riemann, el intercambio en el orden de integración y se parametrizará la región. Esto se hará básicamente, aplicando las fórmulas \ref{integralesdobles1} y \ref{integralesdobles2}.
\section{
Sumas de Riemmann
}
\begin{ejemplo}\normalfont
Evaluar mediante una suma de Riemann el volumen del sólido bajo la superficie $f(x,y)=6-\frac{1}{4}x^{3}-y^{2}$ sobre el plano $xy$ con $-1\leq x\leq 2,$ $0\leq y\leq 2.$

El rectángulo $R=\left[ -1,2\right] \times \left[ 0,2\right] $ se dividirá en $nm$ subrectángulos, $n$ para el intervalo $\left[ -1,2\right] $ y  $m$ para el intervalo $\left[0,2\right] .$ Asumiremos que los rectángulos en $x$ e $y,$ son iguales, respectivamente; esto es
$\Delta x=\frac{2-(-1) }{n}= \frac{3}{n},\,\,\Delta y=\frac{2-0}{m}= \frac{2}{m}.$
Por otra parte,


\begin{multline*}
x_{0}=-1,x_{1}=-1+\Delta x,x_{2}=-1+2\Delta x,\,\cdots, \\
x_{i}=-1+i\Delta x,\cdots,  x_{n}=-1+n\Delta x \\[0.1cm]
y_{0}=0,y_{1}=0+\Delta y,y_2=2\Delta y,\,\cdots ,y_{j}=j\Delta y,\,\cdots , \\ y_{m}=m\Delta y=m\left(\frac{2}{m}\right) = 2.
\end{multline*}



\begin{figure}[H]
\centering
\caption{$f(x,y)=6-\frac{1}{4}x^{3}-y^{2}.$}
\includegraphics[width=0.4\textwidth]{Img/T_69c.pdf} \hspace{5mm}
\includegraphics[width=0.4\textwidth]{Img/T_69b.pdf}

\includegraphics[width=0.4\textwidth]{Img/T_69a.pdf}
\end{figure}

En resumen $x_{i}=-1+\frac{3i}{n},$ $y_{j}=\frac{2j}{m}.$
Si $( x_{i}^{\ast },y_{j}^{\ast}) $ es un punto cualquiera en el rectángulo
$\left[ x_{i},x_{i+1}\right] \times \left[ y_{j},y_{j+1}\right],$
en particular consideraremos $\left( x_{i},y_{j}\right) ,$ entonces el volumen aproximado de paralelepípedo está dado por
\begin{eqnarray*}
f\left( x_{i},y_{j}\right) \Delta x\Delta y &=&\left( 6-\frac{1}{4}\left( -1+ \frac{3i}{n}\right) ^{3}-\left( \frac{2j}{m}\right) ^{2}\right) \frac{3}{n} \frac{2}{m} \\
&=& \frac{75}{2mn}-\frac{24j^{2}}{m^{3}n}+\frac{81i^{2}}{2mn^{3}}-\frac{81i^{3}}{2mn^{4}}-\frac{27i}{2mn^{2}}
\end{eqnarray*}
Sumando todos estos términos, obtendremos el volumen aproximado del sólido mencionado.
\begin{multline*}
\sum_{j=1}^m\sum_{i=1}^n f(x_i,y_j) \Delta x\Delta y
% = \sum_{j=1}^{m} \sum_{i=1}^{n}\left(\frac{75}{2mn}-\frac{24j^{2}}{m^{3}n}+\frac{81i^{2}}{2mn^{3}}-\frac{81i^{3}}{2mn^{4}}-\frac{27i}{2mn^{2}}\right)  \\
= \frac{209}{8}-\frac{4}{m^{2}}-\frac{27}{4n}-\frac{27}{8n^{2}}-\frac{12}{m}
\end{multline*}

La aproximación mejora si el número de rectángulos en que se divide $R$ es muy grande; de manera alternativa, si el ancho y el largo de dichos rectángulos tienden a cero; es por este motivo que se considera el siguiente límite:
\[
\lim_{n,m\to \infty }\left(\sum_{j=1}^{m}\sum_{i=1}^{n}f\left( x_{i},y_{j}\right) \Delta x\Delta y\right) = \frac{209}{8}.
\]
Evaluando la integral doble, se obtiene el mismo valor:
\[
\int_{-1}^{2}\int_{0}^{2}\left( 6-\frac{1}{4}x^{3}-y^{2}\right) \,dy\,dx
=\int_{-1}^{2}\left( \frac{28}{3}-\frac{1}{2}x^{3}\right) dx= \frac{209}{8}.
\]
Los gráficos, la suma de Riemann con su correspondiente límite y la integral doble se evalúan en \verb"Mathematica" con ayuda de las siguientes instrucciones:
\end{ejemplo}



\begin{tcolorbox}[breakable,enhanced,title=\bf Instrucción \verb"Mathematica", top=2pt,bottom=2pt,boxsep=0pt,before skip=2pt,after skip=0pt]
\begin{Verbatim}[formatcom=\letraverbatim]
f[x_,y_]:=6-x^3/4-y^2
Plot3D[f[x,y],{x,-1,2},{y,0,2}]
DiscretePlot3D[f[x,y],{x,-1,2,1/6},{y,0,2,1/8},
ExtentSize->Full,PlotStyle->Gray];
{xi,yj,dx,dy}={-1+3*i/n,2*j/m,3/n,2/m};
suma=Expand[Sum[f[xi,yj]*dx*dy,{i,1,n},{j,1,m}]]
Limit[suma,{n,m}->{Infinity,Infinity}]
Integrate[f[x,y],{x,-1,2},{y,0,2}]
\end{Verbatim}
\end{tcolorbox}



\begin{ejemplo}\normalfont
Usando sumas de Riemann, hallar el volumen del sólido bajo la superficie $f(x,y) =3-x^{2}y$ sobre el plano $xy,$ con $-3/2\leq x\leq 3/2,$
$0\leq y\leq 5/4.$

El rectángulo $\left[ -3/2,3/2\right] \times \left[ 0,5/4\right] $ se dividirá en $nm$ subrectángulos.
$n$ para el intervalo $\left[ -3/2,3/2\right] $ y $m$ para el intervalo $\left[ 0,5/4\right] $.




\begin{figure}[H]
\centering
\caption{$f(x,y)=3-x^{2}y$}
\includegraphics[width=0.4\textwidth]{Img/T_70a.pdf} \hspace{5mm}
\includegraphics[width=0.4\textwidth]{Img/T_70b.pdf}

\includegraphics[width=0.4\textwidth]{Img/T_70c.pdf}
\end{figure}

En este caso, la partición de los intervalos está dada por:
\[
\Delta x=\frac{\frac{3}{2}-\left(-\frac{3}{2}\right) }{n}= \frac{3}{n}\qquad \text{y}\qquad \Delta y=\frac{5/4-0}{m}= \frac{5}{4m}\text{.}
\]
Entonces, los puntos de partición son:
\[
x_{i}=-\frac{3}{2}+i\Delta x = -\frac{3}{2}+\frac{3i}{n}\qquad \text{y}\qquad y_{j}=j\Delta y = \frac{5j}{4m}\text{.}
\]
Evaluando la función $f(x, y)$ en estos puntos, se obtiene:
\[
f\left( x_{i},y_{j}\right) =3-x_{i}^{2}y_{j}=3-\left(\frac{3i}{n}-\frac{3}{2}\right)^{2}\left(\frac{5j}{4m}\right) = 3-\frac{45 j i^2}{4 m n^2}+\frac{45 j i}{4 m n}-\frac{45 j}{16 m}.
\]
Aplicando la suma de Riemann, el resultado es:
\begin{equation*}
	\sum_{i=1}^{n}\sum_{j=1}^{m}f\left( x_{i},y_{j}\right) \Delta x\Delta y = \frac{1215}{128}-\frac{225}{64 m n^2}-\frac{225}{128 m}-\frac{225}{64 n^2}.
\end{equation*}
Al tomar el límite de la suma de Riemann, se obtiene:
\[
\lim_{n,m\rightarrow \infty }\left( \sum_{i=1}^{n}\sum_{j=1}^{m}f(x_{i},y_{j}) \Delta x\Delta y\right) = \frac{1215}{128}
\]
Por otro lado, al evaluar la integral doble se obtiene el mismo valor:
\[
\int_{-\frac{3}{2}}^{\frac{3}{2}}\int_{0}^{\frac{5}{4}}\left( 3-x^{2}y\right) dydx=\int_{-\frac{3}{2}}^{\frac{3}{2}}\left[ 3y-\frac{1}{2}x^{2}y^{2}\right] _{y=0}^{y=\frac{5}{4}}dx=\frac{1215}{128}\text{.}
\]
Esta suma de Riemann, su correspondiente límite y la integral doble, se evalúan en \verb"Mathematica" con ayuda de las siguientes instrucciones:
\end{ejemplo}



\begin{tcolorbox}[breakable,enhanced,title=\bf Instrucción \verb"Mathematica", top=2pt,bottom=2pt,boxsep=0pt,before skip=2pt,after skip=0pt]
\begin{Verbatim}[formatcom=\letraverbatim]
f[x_,y_]:=3-x^2*y
Plot3D[f[x,y],{x,-3/2,3/2},{y,0,5/4}]
DiscretePlot3D[f[x,y],{x,-3/2,3/2,1/10,
{y,0,5/4,1/10},ExtentSize->Full,PlotStyle->Gray]
{xk,yj,dx,dy}={-3/2+3*k/n,5/4*j/m,3/n,5/(4*m)};
suma=Sum[f[xk,yj]*dx*dy,{k,1,n},{j,1,m}]
Limit[suma,{n,m}->{Infinity,Infinity}]
Integrate[f[x,y],{x,-3/2,3/2},{y,0,5/4}]
\end{Verbatim}
\end{tcolorbox}




\begin{ejemplo}\normalfont
Utilizando una suma de Riemann, evaluar la integral doble de $f(x,y)=x^2 \sin (2y)+y^2 \cos (x)$ sobre el rectángulo $[0,\pi/2]\times[0,\pi/2].$

\vskip1cm


\begin{figure}[H]
\centering
\caption{$f(x,y)=x^2 \sin (2y)+y^2 \cos (x)$}
\includegraphics[width=0.4\textwidth]{Img/T_68a.pdf} \hspace{5mm}
\includegraphics[width=0.4\textwidth]{Img/T_68b.pdf}

\includegraphics[width=0.4\textwidth]{Img/T_68c.pdf}
\end{figure}

Si los rectángulos en $x$ e $y$ son iguales, entonces
$\Delta x=\frac{\pi}{2n}$ y $\Delta y=\frac{\pi}{2m};$ por lo tanto,
$x_i=\frac{\pi i}{2n}$ y $y_j=\frac{\pi j}{2m}$. Por otra parte,
\[f(x_i,y_j)\Delta x\Delta y=\frac{\pi^4i^2\sin\left(\frac{\pi j}{m}\right)}{4mn^3}+\frac{\pi^4j^2\cos\left(\frac{\pi i}{2n}\right)}{4m^3n}\]
y
\begin{equation*}
	\resizebox{\linewidth}{!}{$
		\displaystyle
		\sum_{i=1}^n\sum_{j=1}^m f(x_i,y_j)\,\Delta x\,\Delta y
		=\frac{\pi^4(2m^2+3m+1)}{192m^2n}\left(\cot\frac{\pi}{4n}-1\right)
		+\frac{\pi^4(2n^2+3n+1)}{96mn^2}\cot\frac{\pi}{2m}.
		$}
\end{equation*}

entonces
\begin{equation*}
	\resizebox{\linewidth}{!}{$
		\displaystyle
		\int_0^{\frac{\pi}{2}}\!\int_0^{\frac{\pi}{2}}\big(x^2\sin(2y)+y^2\cos x\big)\,dy\,dx
		= \lim_{n,m\to\infty}\Big(\sum_{i=1}^n\sum_{j=1}^m f(x_i,y_j)\,\Delta x\,\Delta y\Big)
		= \frac{\pi^3}{12}.
		$}
\end{equation*}

Las figuras y cálculos anteriores se hicieron en \verb"Mathematica" ejecutando las siguientes instrucciones.
\end{ejemplo}



\begin{tcolorbox}[breakable,enhanced,title=\bf Instrucción \verb"Mathematica", top=2pt,bottom=2pt,boxsep=0pt,before skip=2pt,after skip=0pt]
\begin{Verbatim}[formatcom=\letraverbatim]
f[x_,y_]:=y^2*Cos[x]+x^2*Sin[2*y]
DiscretePlot3D[f[x,y],{x,0,Pi/2,Pi/35},
{y,0,Pi/2,Pi/35},ExtentSize->Full,PlotStyle->Gray]
{xi,yj,dx,dy}={Pi/2*i/n,Pi/2*j/m,Pi/(2*n),Pi/(2*m)};
Expand[f[xi,yj]*dx*dy]
suma=Simplify[Sum[f[xi,yj]*dx*dy,{i,1,n},{j,1,m}]]
Limit[suma,{n,m}->{Infinity,Infinity}]
Integrate[f[x,y],{y,0,Pi/2},{x,0,Pi/2}]
\end{Verbatim}
\end{tcolorbox}



Uno de los procedimientos que se utilizan con frecuencia para evaluar integrales dobles se ilustra a continuación.

\section{Cambio en el orden de integración}



\begin{enumerate}[leftmargin=0cm,itemindent=.75cm,labelwidth=\itemindent,labelsep=0cm,align=left,label={{\bf \arabic*.\,}}]
\item Evaluar la integral de la función $f(x,y)=xy$, sobre la región $D$ en el primer cuadrante del plano $xy$ acotada por las figuras de las expresiones $y=x,\,\,\,x^{2}+y^{2}=4$ y $x=2.$La intersección de la recta inclinada y la circunferencia corresponden en el primer cuadrante a $x=\sqrt{2}.$
\end{enumerate}


\begin{figure}[H]
\centering
\caption{$y=x,$ $x^{2}+y^{2}=4$}
\includegraphics[width=0.8\textwidth]{Fig/83.pdf}
\fuentepropia
\end{figure}

Para integrar en el orden $dxdy$, hay que tener en cuenta que hay una parte de la región donde $x$
recorre los puntos (de manera horizontal) desde la circunferencia hasta la recta $x=2$, con
$y \in [0,\sqrt{2}]$ y hay otra región de la región en la cual $x$ recorre los puntos desde la recta $y=x$
hasta la recta $x=2,$ con $y \in [\sqrt{2},2].$ 


En el orden $dydx$ se observa del gráfico que $y$ recorre los puntos (de manera vertical)
desde la circunferencia hasta la recta $y=x$, con $x\in[\sqrt{2},2].$



\begin{figure}[H]
\centering
\caption{$x\leq y \leq \sqrt{4-x^{2}}$}
\includegraphics[width=0.4\textwidth]{Fig/84.pdf} \hspace{5mm}
\includegraphics[width=0.4\textwidth]{Fig/85.pdf}
\fuentepropia
\end{figure}

Apoyándonos en el gráfico, la integral $\displaystyle\iint_{D}xydA$ se puede evaluar como sigue
\[\int_{0}^{\sqrt{2}}\int_{\sqrt{4-y^{2}}}^{2}xydxdy+\int_{\sqrt{2}}^{2}\int_{y}^{2}xydxdy=\int_{\sqrt{2}}^{2}\int_{\sqrt{4-x^{2}}}^{x}xydydx=1.\]



En coordenadas polares, la ecuación $x^2+y^2=4$ corresponde a $r=2$ y la recta $x=2$ se escribe como $r=2\sec \theta$.
Del gráfico se observa que $r$ recorre los puntos desde la circunferencia $r=2$ hasta la recta $r=2\sec \theta$.
La integral, usando la fórmula \ref{integraldoblepolares}, se evalúa de la siguiente manera:
\[\int_{0}^{\pi /4}\int_{2}^{2\sec \theta }r^{3}\sin \theta \cos\theta drd\theta =1.\]


\begin{figure}[H]
\centering
\caption{$2 \leq r \leq 2 \sec \theta.$}
\includegraphics[width=0.5\textwidth]{Fig/86.pdf}
\fuentepropia
\end{figure}

Estos cálculos se hicieron ejecutando las siguientes instrucciones:



\begin{tcolorbox}[breakable,enhanced,title=\bf Instrucción \verb"Mathematica", top=2pt,bottom=2pt,boxsep=0pt,before skip=2pt,after skip=0pt]
\begin{Verbatim}[formatcom=\letraverbatim]
Integrate[x*y,{x,Sqrt[2],2},{y,Sqrt[4-x^2],x}]
Integrate[x*y,{y,0,Sqrt[2]},{x,Sqrt[4-y^2],2}]+
Integrate[x*y,{y,Sqrt[2],2},{x,y,2}]
\end{Verbatim}
\end{tcolorbox}




El mismo valor se obtiene de la siguiente manera.



\begin{tcolorbox}[breakable,enhanced,title=\bf Parametrizando la región de integración, top=2pt,bottom=2pt,boxsep=0pt,before skip=2pt,after skip=0pt]
Un punto sobre la circunferencia tiene la forma {\small $P\left( 2\cos u, 2\sin u\right)$}, con {\small $u\in\left[0,\frac{\pi}{4}\right]$} y un punto sobre la recta $y=x$, fijando la primera coordenada, es de forma {\small $Q\left( 2\cos u, 2\cos u\right)$}. Por lo tanto, una combinación convexa de estos puntos parametriza la región de integración, es decir,

{\footnotesize \[
\mathbf{r}(u,v) = (1-v)P + vQ = \left\langle 2\cos u,2(1-v)\sin u+2v\cos u,0\right\rangle.\]}
Se agregó un $0$ en el tercer componente con el fin de evaluar el producto cruz.
Derivando parcialmente se obtienen los siguientes vectores:

{\footnotesize
\begin{align*}
\mathbf{r}_{u} &=  \left\langle -2\sin u, 2(1-v)\cos u-2v\sin u, 0\right\rangle \\
\mathbf{r}_{v} &= \left\langle 0, 2\cos u-2\sin u, 0\right\rangle \\
\mathbf{r}_{u}\times \mathbf{r}_{v} &=\left\langle 0, 0, -4(\cos u-\sin u) \sin u\right\rangle.
\end{align*}
}
Además, {\small \[xy = 4\left((1-v)\sin u+v\cos u\right)\cos u\]} y
{\small $\left\Vert \mathbf{r}_{u}\times \mathbf{r}_{v}\right\Vert =\left\vert 4( \sin u- \cos u) \sin u\right\vert,$}
de modo que la integral $\iint_D xy dA$ se convierte en
\\[0.1em]
\begin{center}
\small 
$\int_{0}^{1}\int_{0}^{\pi /4}16\left|\cos u((1-v)\sin u+v\cos u)(\sin u-\cos u)\sin u\right| dv du= 1.$
\end{center}
\end{tcolorbox}



\begin{enumerate}[start=2,leftmargin=0cm,itemindent=.75cm,labelwidth=\itemindent,labelsep=0cm,align=left,label={{\bf \arabic*.\,}}]
\item Invertir el orden de integración y evaluar la siguiente integral:
{\footnotesize \[\displaystyle\int_{0}^{1}\int_{\sqrt{y}}^{2-y}xy\, dxdy.\]}

La región de integración corresponde a la figura limitada en el primer cuadrante, por las figuras de las expresiones
$y=0$, $y=x^2$ y $y=2-x$.
\end{enumerate}

\begin{figure}[H]
	\centering
	\caption{$\sqrt{y}\leq x \leq 2-y $}
	
	\includegraphics[width=0.28\linewidth]{Fig/80.pdf}
	\includegraphics[width=0.28\linewidth]{Fig/81.pdf}
	\includegraphics[width=0.25\linewidth]{Fig/82.pdf}
	
	\fuentepropia
\end{figure}
%

Se puede observar que en el orden $dydx$, $y$ recorre (verticalmente) los puntos desde el eje $x$ hasta la parábola, en una parte, y hasta la recta inclinada en otra parte. Para el orden $dxdy$, $x$ recorre (horizontalmente) los puntos desde la parábola hasta la recta. Esto da origen a las siguientes integrales dobles en el orden que se requiere.

{\footnotesize
\[ \int_{0}^{1}\int_{\sqrt{y}}^{2-y}xydxdy =
\int_{0}^{1}\int_{0}^{x^{2}}xydydx+\int_{1}^{2}\int_{0}^{2-x}xydydx= \frac{7}{24}.\]
}
Estas integrales se evalúan con la ayuda de las siguientes instrucciones:



\begin{tcolorbox}[breakable,enhanced,title=\bf Instrucción \verb"Mathematica", top=2pt,bottom=2pt,boxsep=0pt,before skip=2pt,after skip=0pt]
\begin{Verbatim}[formatcom=\letraverbatim]
Integrate[x*y,{y,0,1},{x,Sqrt[y],2-y}]
Integrate[x*y,{x,0,1},{y,0,x^2}]+Integrate[x*y,
{x,1,2},{y,0,2-x}]
\end{Verbatim}
\end{tcolorbox}



Otra manera de evaluar la misma integral es la siguiente:



\begin{tcolorbox}[breakable,enhanced,title=\bf Parametrizando la región de integración, top=2pt,bottom=2pt,boxsep=0pt,before skip=2pt,after skip=0pt]
	
	La porción de parábola se parametriza por $\la \sqrt{t},t\ra,$ mientras que la recta se representa por $\la 2-t,t\ra ,$ para $t\in \lbrack 0,1].$
	Una combinación convexa de estas parametriza la región de interés, esto es
	{\small
	\[{\bf r}(t,u)=(1-u)\la 2-t,t\ra +u\la \sqrt{t},t\ra = \la (u-1)(t-2)+u\sqrt{t},t\ra\]}
	con $u\in \lbrack 0,1].$ De esta función se obtienen los siguientes vectores:
	{\small
	\[{\bf r}_{u} =\la t+\sqrt{t}-2,0\ra,{\bf r}_{t} =\la u+\frac{u}{2\sqrt{t}}-1,1\ra.\]}
	
	Los extendemos a $\mathbb{R}^{3},$ para obtener
	${\bf r}_{u}\times {\bf r}_{t}= \la 0,0,t+\sqrt{t}-2\ra.$
	De lo cual se sigue que $\left\Vert{\bf r}_{u}\times{\bf r}_{t}\right\Vert=\left\vert t+\sqrt{t}-2\right\vert.$
	Con la definición de la función se observa que $xy=(\left(u-1\right) \left( t-2\right) +\sqrt{t}u) t,$ entonces:
	\\[0.6em]
	\resizebox{\linewidth}{!}{$
		\iint_{R}xydA=\int_{0}^{1}\int_{0}^{1}\left((u-1)(t-2)+u\sqrt{t}\right)t\left|t+\sqrt{t}-2\right|\,du\,dt=\frac{7}{24}
		$}
\end{tcolorbox}



\begin{enumerate}[start=3,leftmargin=0cm,itemindent=.75cm,labelwidth=\itemindent,labelsep=0cm,align=left,label={{\bf \arabic*.\,}}]
	\item Determinar el centro de masa de una lámina delgada cuya densidad de área está dada por
$\delta (x,y)=x$ $kg/m^2$ y posee la forma de la región en el primer cuadrante interior a la circunferencia con ecuación $x^{2}+y^{2}=4$ y exterior a la circunferencia $x^{2}+y^{2}=2x.$
\end{enumerate}



Para plantear las integrales dobles utilizando coordenadas rectangulares, se debe tener en cuenta que en el orden $dxdy$, el área considerada se divide en tres porciones. En una de ellas $x$ toma valores desde el eje $y$ hasta la circunferencia de radio $1$.


\begin{figure}[H]
\centering
\caption{$\sqrt{x^2-2x}\leq y \leq \sqrt{4-x^2}$}
\includegraphics[width=0.8\textwidth]{Fig/87.pdf}
\fuentepropia
\end{figure}

En una segunda región $x$ toma valores desde la circunferencia de radio $1$ hasta la circunferencia mayor y en una tercera región $x$ recorre los puntos desde el eje $y$ hasta la circunferencia de radio $2$. En el orden $dydx$, $y$ recorre los puntos desde la circunferencia de radio $1$ hasta la circunferencia de radio $2$ con $x\in [0,2].$

En coordenadas polares la ecuación de la circunferencia $x^2+y^2=4$ se reescribe como $r=2$
y la ecuación $x^2+y^2=2x$ corresponde a $r^2=2r\cos \theta,$ es decir, $r=2 \cos \theta.$ Es por esto que $r$
recorre los puntos desde la circunferencia $r=2 \cos \theta$ hasta $r=2;$
es decir, la región de interés se describe con $2\cos \theta \leq r \leq 2$ y $0\leq  \theta \leq \pi/2$.



\begin{figure}[H]
\centering
\caption{$2\cos \theta \leq r \leq 2$}
\includegraphics[width=0.4\textwidth]{Fig/88.pdf} \hspace{5mm}
\includegraphics[width=0.4\textwidth]{Fig/89.pdf}

\includegraphics[width=0.4\textwidth]{Fig/90.pdf}
\fuentepropia
\end{figure}

Estas figuras ayudan a establecer que la masa $M$ de la lámina en el orden $dxdy$ está dada por
\[
M=\int_{0}^{1}\int_{0}^{1-\sqrt{1-y^{2}}}xdxdy+\int_{0}^{1}\int_{1+\sqrt{1-y^{2}}}^{\sqrt{4-y^{2}}}xdxdy+\int_{1}^{2}\int_{0}^{\sqrt{4-y^{2}}}xdxdy
\]

En el orden $dydx$, la integral correspondiente es
\[
M=\int_{0}^{2}\int_{\sqrt{2x-x^{2}}}^{\sqrt{4-x^{2}}}xdydx=\frac{8}{3}-\frac{\pi}{2}.\]
y en coordenadas polares la masa está dada por
\[M=\ds\int_{0}^{\pi/2}\int_{2\cos\theta}^2\left(r\cos\theta\right)rdrd\theta=\frac{8}{3}-\frac{\pi}{2}.\]
Mientras que los momentos con respecto al eje  $y$ y al eje $x$ están dados por
\begin{eqnarray*}
\mu {y} &=& \int_{0}^{2}\int_{\sqrt{2x-x^{2}}}^{\sqrt{4-x^{2}}}x^2dydx=\int_{0}^{\pi /2}\int_{2\cos \theta }^{2}r^{3}\cos ^{2}\theta drd\theta = \frac{3}{8}\pi \\
\mu {x} &=& \int_{0}^{2}\int_{\sqrt{2x-x^{2}}}^{\sqrt{4-x^{2}}}xydydx=\int_{0}^{\pi /2}\int_{2\cos \theta }^{2}r^{3}\cos \theta \sin \theta drd\theta = \frac{4}{3}.
\end{eqnarray*}
Por lo tanto, el centro de masa de la lámina es
\[\left( \overline{x},\overline{y}\right) = \left( \frac{\mu _{y}}{M},\frac{\mu _{x}}{M}\right) = \left( \frac{\frac{3}{8}\pi }{\frac{8}{3}-\frac{1}{2}\pi },\frac{\frac{4}{3}}{\frac{8}{3}-\frac{1}{2}\pi }\right) = \left( \frac{9\pi }{64-12\pi },\frac{8}{16-3\pi}\right). \]

Los cálculos para obtener los momentos y la masa se hacen usando las siguientes instrucciones:



\begin{tcolorbox}[breakable,enhanced,title=\bf Instrucción en \verb"Mathematica", top=2pt,bottom=2pt,boxsep=0pt,before skip=2pt,after skip=0pt]
\begin{Verbatim}[formatcom=\letraverbatim]
Integrate[x,{x,0,2},{y,Sqrt[2*x-x^2],Sqrt[4-x^2]}]
Integrate[x*y,{x,0,2},{y,Sqrt[2*x-x^2],Sqrt[4-x^2]}]
Integrate[x^2,{x,0,2},{y,Sqrt[2*x-x^2],Sqrt[4-x^2]}]
Integrate[r^2*Cos[t],{t,0,Pi/2},{r,2*Cos[t],2}]
\end{Verbatim}
\end{tcolorbox}



%\clearpage

\begin{enumerate}[start=4,leftmargin=0cm,itemindent=.75cm,labelwidth=\itemindent,labelsep=0cm,align=left,label={{\bf \arabic*.\,}}]
	
\item
Determinar las coordenadas $(\overline{x},\overline{y})$ del centro de masa de una lámina determinada por la región en el primer cuadrante acotada por las figuras de las ecuaciones $x^2+y^2=4$ y $x^2+y^2=2y,$
cuya densidad de masa en cualquier punto $(x,y)$ es $\delta (x,y)=xy^2$. La forma de la lámina se presenta en la figura \ref{fig510}.
\end{enumerate}


\begin{figure}[H]
\centering
\caption{$2\sin \theta \leq r \leq 2$}
\label{fig510}
\includegraphics[width=0.8\textwidth]{Fig/91.pdf}
\fuentepropia
\end{figure}

Con el propósito de plantear las integrales dobles que determinan la masa y los demás elementos requeridos, observemos en  los siguientes gráficos la manera como se recorre la región de integración.


En el orden $dxdy$, $x$ recorre los puntos desde la circunferencia menor hasta la mayor, para $y \in [0,2]$. En el orden $dydx$, la región se divide en tres porciones. En una parte $y$ recorre los puntos desde el eje $x$ hasta la circunferencia menor con $y \in [0,1]$; en una segunda parte $y$ recorre los puntos desde la circunferencia menor hasta la circunferencia mayor, con $y \in [1,2]$, y en la tercera parte $y$ recorre los puntos desde el eje $x$ hasta la circunferencia mayor con $y \in [0,\sqrt{3}]$. También se debe observar que de la ecuación $x^2+y^2=2y$ se desprende $y=1 \pm \sqrt{1-x^2}.$


\begin{figure}[H]
\centering
\caption{$2\sin \theta \leq r \leq 2$}

\includegraphics[width=0.30\linewidth]{Fig/92.pdf}
\includegraphics[width=0.26\linewidth]{Fig/93.pdf}
\includegraphics[width=0.26\linewidth]{Fig/94.pdf}
\fuentepropia
\end{figure}

En coordenadas polares se satisface que $2\sin \theta \leq r \leq 2$ para $0 \leq \theta \leq \pi/2$. Por tanto, la masa $M$ de la lámina se puede hallar evaluando cualquiera de las siguientes integrales.
%
\[
\resizebox{\textwidth}{!}{$
\begin{aligned}
M &= \int_{0}^{1} \int_{0}^{1-\sqrt{1-x^{2}}}xy^{2} \,dy\,dx 
+ \int_{0}^{1} \int_{1+\sqrt{1-x^{2}}}^{\sqrt{4-x^{2}}} xy^{2} \,dy\,dx 
+ \int_{1}^{2} \int_{0}^{\sqrt{4-x^{2}}} xy^{2} \,dy\,dx \\
&= \int_{0}^{2} \int_{\sqrt{2y - y^{2}}}^{\sqrt{4 - y^{2}}} xy^{2} \,dx\,dy 
= \int_{0}^{\frac{\pi}{2}} \int_{2\sin\theta}^{2} r^{4} \sin^{2}\theta \cos\theta \,dr\,d\theta 
= \frac{4}{3}
\end{aligned}
$}
\]

%
Los momentos con respecto al eje $x$, $\mu_{x}$ y, con respecto al eje $y$, $\mu_{y}$,
se pueden determinar de la siguiente manera:

$\mu_{x}=\int_{0}^{2}\int_{\sqrt{2y-y^{2}}}^{\sqrt{4-y^{2}}}xy^{3}dxdy=\int_{0}^{\frac{\pi }{2}}\int_{2\sin \theta }^2r^{5}\cos \theta \sin^{3}\theta drd\theta=\frac{8}{5}$

$\mu_{y}= \int_{0}^{2}\int_{\sqrt{2y-y^{2}}}^{\sqrt{4-y^{2}}}x^{2}y^{2}dxdy=\int_{0}^{\pi /2}\int_{2\sin \theta }^{2} r^5 \cos^2 \theta\sin^2 \theta drd\theta =\frac{25}{48}\pi .$


Con lo anterior, se concluye que el centro de masa está dado por
\[(\overline{x},\overline{y})= \left( \frac{\frac{25}{48}\pi }{\frac{4}{3}},\frac{\frac{8}{5}}{\frac{4}{3}}\right) = \left( \frac{25\pi }{64},\frac{6}{5}\right).\]


El cálculo de la masa se puede evaluar mediante el uso de las siguientes indicaciones.



\begin{tcolorbox}[breakable,enhanced,title=\bf Instrucción en \verb"Mathematica", top=2pt,bottom=2pt,boxsep=0pt,before skip=2pt,after skip=0pt]
\begin{Verbatim}[formatcom=\letraverbatim]
Integrate[x*y^2,{y,0,2},{x,Sqrt[2*y-y^2],
Sqrt[4-y^2]}]
Integrate[r^4*Sin[t]^2*Cos[t],{t,0,Pi/2},
{r,2*Sin[t],2}]
\end{Verbatim}
\end{tcolorbox}



\begin{enumerate}[start=5,leftmargin=0cm,itemindent=.75cm,labelwidth=\itemindent,labelsep=0cm,align=left,label={{\bf \arabic*.\,}}]
\item
Si $f$ es continua en $R$, plantear la integral $\ds\iint_R f(x,y)dA$, siendo $R$ la región en el primer cuadrante acotada por las figuras de las ecuaciones $x^2+y^2=4,$ $y= \sqrt{3}\,x$ y $y=1.$
En coordenadas polares $y=1$ se transforma en $r= \csc \theta$, $y= \sqrt{3}\,x$ equivale a $\theta = \pi /3 $ y la ecuación $x=\sqrt{4-y^2}$ corresponde a un arco de circunferencia de radio $2$ con centro en el origen.
\end{enumerate}


\begin{figure}[H]
\centering
\caption{$\csc \theta \leq r \leq 2,$ $\frac{\pi}{6} \leq\theta \leq \frac{\pi}{3}$}
\includegraphics[width=0.5\textwidth]{Fig/95.pdf}
\fuentepropia
\end{figure}

En los siguientes gráficos se muestran las diversas maneras de recorrer la región $R$. Para integrar en el orden $dxdy$ se debe tener en cuenta que $x$ recorre los puntos desde la recta $y=\sqrt{3}\, x$ hasta la circunferencia. Para el orden $dydx$, en una parte $y$ recorre puntos desde la recta $y=1$ hasta la recta $y=\sqrt{3}\, x,$ y en otra parte $y$ recorre los puntos desde $y=1$ hasta la circunferencia.

% REDUCCIÓN DEL ESPACIO VERTICAL SUPERIOR
\vspace*{1.0em} 
\begin{center} 
	% Usamos un parbox principal para contener y alinear todo el bloque
	\parbox{\linewidth}{
		\centering
		
		% LEYENDA (encima de la figura)
		\footnotesize
		\textbf{Figura 13.} $\csc \theta \leq r \leq 2,\quad \frac{\pi}{6} \leq \theta \leq \frac{\pi}{3}.$
		 % Espacio mínimo entre leyenda e imágenes
		
		% LAS TRES IMÁGENES
		\includegraphics[width=0.28\linewidth]{Fig/96.pdf}
		\includegraphics[width=0.28\linewidth]{Fig/97.pdf}
		\includegraphics[width=0.28\linewidth]{Fig/98.pdf}
		
		% NOTA DE ELABORACIÓN
		\fuentepropia
	}
\end{center}
% REDUCCIÓN DEL ESPACIO VERTICAL INFERIOR (antes del siguiente párrafo)
\vspace*{1.2em} 

En coordenadas polares, $r$ toma valores desde la recta horizontal hasta la circunferencia con $\theta \in [\frac{\pi}{6},\frac{\pi}{3}].$
De acuerdo con lo anterior, se pueden plantear las integrales, como se muestra a continuación:


% REDUCCIÓN DEL ESPACIO DE LA ECUACIÓN
%\setlength{\abovedisplayskip}{3pt} 
%\setlength{\belowdisplayskip}{3pt}
\[
\large
\resizebox{\textwidth}{!}{$
	\begin{array}{rcl}
		\ds\iint_R f(x,y)\,dA &=& \int_{\frac{1}{\sqrt{3}}}^{1}\int_{1}^{\sqrt{3}x}f(x,y)\,dy\,dx 
		+ \int_{1}^{\sqrt{3}}\int_{y}^{\sqrt{4-x^{2}}}f(x,y)\,dy\,dx \\[10pt]
		&=& \int_{1}^{\sqrt{3}}\int_{y/\sqrt{3}}^{\sqrt{4-x^{2}}}f(x,y)\,dx\,dy 
		= \int_{\pi/6}^{\pi/3}\int_{\csc\theta}^{2}f(r,\theta)\,r\,dr\,d\theta
	\end{array}
	$}
\]

%\vspace*{1.2em} 
\begin{enumerate}[start=6,leftmargin=0cm,itemindent=.75cm,labelwidth=\itemindent,labelsep=0cm,align=left,label={{\bf \arabic*.\,}}]
	\item
Escribe en coordenadas rectangulares la siguiente integral: \[\ds\int_{0}^{\pi/4}\int_{0}^{2\sqrt{2}}r^2drd\theta.\]
\end{enumerate}


Los límites de integración son $r=0,$ $r=2\sqrt{2},$ $\theta=0,$ $\theta= \pi/4,$
esto determina una región en el primer cuadrante limitada por las rectas $y=0,$
$y=x,$ y la circunferencia con ecuación \[x^2+y^2=\left(2\sqrt{2}\right)^2=8.\]
Se usó un cambio a coordenadas polares $x= r \cos \theta$  y $y= r \cos \theta$.


\begin{figure}[H]
\centering
\caption{$0 \leq r \leq 2\sqrt{2}.$}
\includegraphics[width=0.8\textwidth]{Fig/99.pdf}
\fuentepropia
\end{figure}

Al evaluar la integral en el orden $dydx,$
la región de integración se divide en dos, en una parte $y$ varía desde el eje $x$
hasta la recta $y=x$ con $x \in[0,2],$ en la otra parte $y$ varía desde el eje $x$
hasta la circunferencia, con $x \in[2,2\sqrt{2}].$
Al integrar en  orden $dxdy$, el $x$ recorre (horizontalmente) los puntos desde la recta hasta la circunferencia, con $y\in[0,2]$.

\begin{figure}[h!]
	\centering
	\caption{$y \leq x \leq \sqrt{8 - x^2}.$}
	\includegraphics[width=0.3\linewidth]{Fig/100.pdf}
	\includegraphics[width=0.3\linewidth]{Fig/101.pdf}
	\includegraphics[width=0.3\linewidth]{Fig/102.pdf}
	\fuentepropia
\end{figure}

Por lo señalado antes, la integral en mención se reescribe como
\begin{multline*}
\displaystyle\int_{0}^{\pi /4}\int_{0}^{2\sqrt{2}}r^2drd\theta = \int_{0}^2\int_{0}^{x}\sqrt{x^2+y^2}dydx + \\ \int_{2}^{2\sqrt{2}}\int_{0}^{\sqrt{8-x^2}}\sqrt{x^2+y^2}dydx = \int_0^2 \int _y ^{\sqrt{8-y^2}} \sqrt{x^2+y^2}\, dxdy = \frac{4\sqrt{2} \pi}{3}
\end{multline*}

Estas integrales se evaluaron por medio de las siguientes instrucciones:

%

\begin{tcolorbox}[breakable,enhanced,title=\bf Instrucción en \verb"Mathematica", top=2pt,bottom=2pt,boxsep=0pt,before skip=2pt,after skip=0pt]
\begin{Verbatim}[formatcom=\letraverbatim]
Integrate[Sqrt[x^2+y^2],{y,0,2},{x,y,Sqrt[8-y^2]}]
Integrate[Sqrt[x^2+y^2],{x,0,2},{y,0,x}]+
Integrate[Sqrt[x^2+y^2],{x,2,2*Sqrt[2]},
{y,0,Sqrt[8-x^2]}]
Integrate[r^2,{\[Theta],0,Pi/4},{r,0,2*Sqrt[2]}]
\end{Verbatim}
\end{tcolorbox}




\begin{enumerate}[start=7,leftmargin=0cm,itemindent=.75cm,labelwidth=\itemindent,labelsep=0cm,align=left,label={{\bf \arabic*.\,}}]
\item 
Evaluar la integral $\ds \iint_{R} |y-x|dA,$ donde $R$ es la región acotada por las figuras de las ecuaciones $y=2-x^{2}$ y $y=x^{2}+2x-2.$


La función $|y-x|$ se expresa como

\[|y-x|= \left \{\begin{array}{cc} y-x & \text{si\,\,} y\geq x \\ x-y & \text{si\,\,} y<x\end{array} \right.\]
\end{enumerate}


\begin{figure}[H]
\centering
\caption{$x^2+2x-2 \leq y \leq 2-x^2$}
\includegraphics[width=0.8\textwidth]{Fig/103.pdf}
\fuentepropia
\end{figure}

El planteamiento de la integral doble en coordenadas rectangulares requiere hacer algunas consideraciones. De la expresión $y=2-x^2$ se obtiene $x= \pm \sqrt{2-y},$ mientras que de la ecuación  $y=x^2+2x-2$ se obtiene que $x=-1\pm \sqrt{y+3}$. Ahora se ilustra gráficamente la región de integración.

\begin{figure}[h!]
	\centering
	\caption{$x^2 + 2x - 2 \leq y \leq 2 - x^2$}
	\includegraphics[width=0.48\linewidth]{Fig/104.pdf}
	\includegraphics[width=0.38\linewidth]{Fig/105.pdf}
	\fuentepropia
\end{figure}

En el orden $dydx$ la región $R$ se divide en dos partes, en una de ellas $y$ recorre los puntos desde la recta $y=x$ hasta la curva $y=2-x^2,$ para $x\in [-2,1].$
En la otra región $y$ recorre los puntos desde la curva
$y=x^2+2x-1$ hasta la recta $y=x$, con $x\in [-2,1].$ En este caso, la integral es
\[\ds \int_{-2}^{1}\int_{x}^{2-x^{2}}\left( y-x\right)
dydx+\int_{-2}^{1}\int_{x^{2}+2x-2}^{x}\left( x-y\right) dyd = \frac{81}{10}\]

Con las siguientes instrucciones se evalúa esta integral.



\begin{tcolorbox}[breakable,enhanced,title=\bf Instrucción en \verb"Mathematica", top=2pt,bottom=2pt,boxsep=0pt,before skip=2pt,after skip=0pt]
\begin{Verbatim}[formatcom=\letraverbatim]
Integrate[y-x,{x,-2,1},{y,x,2-x^2}]+Integrate[x-y,
{x,-2,1},{y,x^2+2*x-2,x}]
\end{Verbatim}
\end{tcolorbox}



Para el orden $dxdy$, $R$ queda dividida en cuatro regiones de acuerdo con el recorrido horizontal que hace $x$.

\begin{itemize}[label=$\bullet$\,\,]
\item  $y\in[1,2]$ se satisface que $-\sqrt{2-y} \leq x \leq \sqrt{2-y}.$
\item Para $y\in[-2,1],$ hay dos regiones: en una parte, $-\sqrt{2-y} \leq x \leq y,$
mientras que en la otra región $y \leq x \leq 1+\sqrt{3+y}.$
\item Para $y\in[-3,-2]$ se cumple que $-1-\sqrt{1+y} \leq x \leq -1+\sqrt{1+y}.$
\end{itemize}
Por lo señalado antes, la integral en este caso es la siguiente:
\begin{multline*}
\int_{1}^{2}\int_{-\sqrt{2-y}}^{\sqrt{2-y}}(y-x)dxdy+\int_{-2}^{1}\int_{-\sqrt{2-y}}^{y}(y-x) dxdy \\
+\int_{-2}^{1}\int_{y}^{-1+\sqrt{3+y}}(x-y)dxdy
+\int_{-3}^{-2}\int_{-1-\sqrt{3+y}}^{-1+\sqrt{3+y}}(x-y)dxdy=\frac{81}{10}.
\end{multline*}
Las integrales anteriores se evalúan con las siguientes instrucciones:



\begin{tcolorbox}[breakable,enhanced,title=\bf Instrucción en \verb"Mathematica", top=2pt,bottom=2pt,boxsep=0pt,before skip=2pt,after skip=0pt]
\begin{Verbatim}[formatcom=\letraverbatim]
Integrate[y-x,{y,1,2},{x,-Sqrt[2-y],Sqrt[2-y]}]+
Integrate[y-x,{y,-2,1},{x,-Sqrt[2-y],y}]+Integrate[x-y,
{y,-2,1},{x,y,-1+Sqrt[3+y]}]+Integrate[x-y,{y,-3,-2},
{x,-1-Sqrt[3+y],-1+Sqrt[3+y]}]
\end{Verbatim}
\end{tcolorbox}



\begin{enumerate}[start=8,leftmargin=0cm,itemindent=.75cm,labelwidth=\itemindent,labelsep=0cm,align=left,label={{\bf \arabic*.\,}}]
	\item 
Evaluar $\ds \iint R x dA,$ donde $R$ es la región acotada por las figuras de las ecuaciones $x^2+y^2=4$
y $y=1-(1+\sqrt{3})x$ con  $x \geq 0.$ \\[0.1cm]

De la ecuación de la recta se sigue que
$x=\frac{1}{2}\left(1-\sqrt{3}\right)(y-1),$ un punto de corte de esta línea con la circunferencia es $(1,-\sqrt{3}).$
\end{enumerate}


\begin{figure}[H]
\centering
\caption{$x^2+y^2=4$}
\includegraphics[width=0.8\textwidth]{Fig/106.pdf}
\fuentepropia
\end{figure}

Para evaluar la integral en el orden $dydx$ se debe tener en cuenta dos regiones: en una de estas $y$ recorre los puntos desde la recta inclinada hasta la circunferencia, con $x \in [0,1],$ en la otra región $y$ recorre puntos desde la parte negativa de la circunferencia hasta la parte positiva de la misma, con $x \in [1,2].$
Al considerar el orden  $dxdy$ tendremos dos regiones, en una de ellas $x$ recorre los puntos desde el eje $y$ hasta la circunferencia,  con $y \in [1,2]$; en la otra región $x$ recorre los puntos desde la recta inclinada hasta la circunferencia, con $y \in [-\sqrt{3},1]$.


\begin{center}
	{\small\textbf{Figura:} $x^2+y^2\leq 4,\quad 1-(1+\sqrt{3})x \geq y$}
	
	
	
	\begin{minipage}{0.29\linewidth}
		\centering
		\includegraphics[width=\linewidth]{Fig/107.pdf}
	\end{minipage}
	\hfill
	\begin{minipage}{0.29\linewidth}
		\centering
		\includegraphics[width=\linewidth]{Fig/108.pdf}
	\end{minipage}
	\hfill
	\begin{minipage}{0.29\linewidth}
		\centering
		\includegraphics[width=\linewidth]{Fig/109.pdf}
	\end{minipage}
	
	
	
	\fuentepropia
\end{center}

Por lo señalado antes, la integral se puede evaluar como sigue.
{\small
\begin{align*}
	\ds \iint R x dA &= \int_{-\sqrt{3}}^{1}\int_{\frac{1}{2}\left(1-\sqrt{3}\right)(y-1)}^{\sqrt{4-y^{2}}}\!\!xdxdy + \int_1^2\int_{0}^{\sqrt{4-y^{2}}}\!\!xdxdy \\
	&= \int_{0}^{1}\int_{1-\left( 1+\sqrt{3}\right)x}^{\sqrt{4-x^{2}}}\!\!xdydx + \int_{1}^{2}\int_{-\sqrt{4-x^{2}}}^{\sqrt{4-x^{2}}}\!\!xdydx = \frac{4}{3}\sqrt{3}+\frac{5}{2}.
\end{align*}
}


Para considerar las coordenadas polares, se halla el ángulo de corte entre la recta inclinada y la circunferencia en el cuarto cuadrante; para ello, usaremos el punto
$(1,-\sqrt{3})$, entonces $\theta=\arctan \left(-\sqrt{3}\right)= -\frac{1}{3}\pi.$ La ecuación de la recta se reescribe como sigue:
$ r\sin \theta =1-\left( 1+\sqrt{3}\right)  r\cos \theta $,
de donde $r=\frac{1}{\sin(\theta)+\left(1+\sqrt{3}\right)\cos(\theta)}.$
Por lo tanto,
\[\ds \iint _R x dA=\int_{-\pi/3}^{\pi /2}\int_{\frac{1}{\sin(\theta)+\left( 1+\sqrt{3}\right) \cos (\theta) }}^{2}\left( r^{2}\cos \theta \right) drd\theta =\frac{5}{2}+\frac{4}{3}\sqrt{3}.\]
Las integrales anteriores se evalúan en \verb"Mathematica" mediante la ejecución de las siguientes instrucciones:




\begin{tcolorbox}[breakable,enhanced,title=\bf Instrucción en \verb"Mathematica", top=2pt,bottom=2pt,boxsep=0pt,before skip=2pt,after skip=0pt]
\begin{Verbatim}[formatcom=\letraverbatim]
Integrate[x,{y,-Sqrt[3],1},{x,(1-Sqrt[3])*(y-1)/2,
Sqrt[4-y^2]}]+Integrate[x,{y,1,2},{x,0,Sqrt[4-y^2]}]
Integrate[x,{x,0,1},{y,1-(1+Sqrt[3])*x,Sqrt[4-x^2]}]+
Integrate[x,{x,1,2},{y,-Sqrt[4-x^2],Sqrt[4-x^2]}]
Integrate[r^2*Cos[t],{t,-Pi/3,Pi/2},
{r,1/((1+Sqrt[3])*Cos[t]+Sin[t]),2}]
\end{Verbatim}
\end{tcolorbox}




%\vspace*{1.2em} 
\begin{enumerate}[start=9,leftmargin=0cm,itemindent=.75cm,labelwidth=\itemindent,labelsep=0cm,align=left,label={{\bf \arabic*.\,}}]
	\item 
Hallar el área de la región encerrada por la curva determinada por las ecuaciones paramétricas $x=2a\cos t-a\sin (2t),$ $y=b\sin t,$ $t\in \lbrack 0,2\pi ].$
\end{enumerate}


De acuerdo al teorema de Green, véase sección \secref{teoremadegreen}, el área $A$, de la región $R$ encerrada por la curva $C,$ puede hallarse por medio de la expresión
\[A=\iint_{R}dA=\int_{C}-\frac{1}{2}ydx+\frac{1}{2}xdy,\]
en este caso $x=2a\cos t-a\sin (2t),$ $y=b\sin (t).$ De estas ecuaciones se obtiene
$dx= -2a\sin t-2a\cos (2t),\,\,\, dy= b\cos t,$
entonces
\begin{multline*}
A=\frac{1}{2}\int_{0}^{2\pi}\bigg((2a\sin(t)+a\cos(2t))(b\sin(t))\\ +(2a\cos(t)-a\sin(2t))(b\cos t) \bigg) dt  = 2\pi ab.
\end{multline*}



\begin{figure}[H]
\centering
\caption{${\bf r}(t)=\langle 2a\cos t-a\sin (2t),b\sin t\rangle,$ $t\in [0,2\pi ]$}
\includegraphics[width=0.8\textwidth]{Fig/110.pdf}
\end{figure}

Otra posibilidad es la parametrización de la región interior a la curva.




\begin{tcolorbox}[breakable,enhanced,title=\bf Instrucción Mathematica, top=2pt,bottom=2pt,boxsep=0pt,
before skip=2pt,after skip=0pt]
Para $u\in [0,1]$, al multiplicar por $u$ las ecuaciones que determinan la curva, se define la siguiente función en $\mathbb{R}^3,$
${\bf r}(t,u)=\left\langle \left( 2a\cos t-\sin (2t) \right) u,ub\sin t,0\right\rangle.$
Con esta función se obtienen los siguientes vectores.
\begin{eqnarray*}
{\bf r}_{t} &=& \left\langle \left( -2a\sin (t)-2\cos(2t)
\right) u,ub\cos(t) ,0\right\rangle \\
{\bf r}_{u} &=& \left\langle 2a\cos (t) -\sin(2t) ,b\sin (t),0\right\rangle,
\end{eqnarray*}
luego, ${\bf r}_{t}\times {\bf r}_{u}= \left\langle 0,0,-ub\left( 2a+2\sin \left(t\right) \cos (2t) -\cos \left( t\right) \sin (2t)\right)  \right\rangle ,$
entonces
$\left\| {\bf r}_{t}\times {\bf r}_{u}\right\| =\left| bu\left( 2a+2\sin \left( t\right)
\cos (2t) -\cos \left( t\right) \sin (2t) \right)\right| ,$
cuya integral nos provee el área de la región en mención
\[\int_{0}^{2\pi }\int_{0}^{1}\left| bu\left( 2a+2\sin \left( t\right) \cos
(2t) -\cos \left( t\right) \sin (2t) \right) \right|
dudt= 2ab\pi.\]
\end{tcolorbox}



%
Para evaluar integrales dobles, es posible plantear un cambio de variable; la fundamentación teórica de este proceso se presentó en la sección \secref{cambioidobles}.
\section{
Cambio de variables en integrales dobles
}

\index{Cambio de variables}
Presentamos algunos ejemplos de aplicación del teorema de cambio de variables en integrales dobles.

\begin{enumerate}[leftmargin=0cm,itemindent=.75cm,labelwidth=\itemindent,labelsep=0cm,align=left,label={{\bf \arabic*.\,}}]
\item Evaluar $\ds \iint_{R}y^{2}dA,$ siendo $R$ la región del plano $xy$ acotada por las figuras de las ecuaciones.
$x=4-y^{2}$ y $x=4y^{2}-1.$


\begin{figure}[H]
	\centering	
	\caption{\small ${\bf T}(u,v)= (4u^2 - v^2, uv)$}
	\label{transformacion_uv}
	\begin{minipage}[b]{0.40\linewidth}
		\centering
		
		\includegraphics[width=\linewidth]{Fig/111.pdf}
	\end{minipage}
	\hspace{0.5cm} 
	\begin{minipage}[b]{0.35\linewidth}
		\centering
		
		\includegraphics[width=\linewidth]{Fig/112.pdf}
	\end{minipage}
	
	
	
	\fuentepropia
\end{figure}

Haremos la sustitución $x=4u^{2}-v^{2},\,\, y=uv.$
En la expresión $x=4-y^{2},$ se hace el cambio y se sigue que
$4u^{2}-v^{2}=4-\left( uv\right) ^{2}$ de donde $u^{2}\left( 4+v^{2}\right)
=4+v^{2},$ es decir,
$4+v^{2}=1$ (lo cual no es soluble en los números reales) o $u^{2}=1,$
entonces $u=\pm 1.$
Sustituyéndose en la ecuación $x=4y^{2}-1,$ se obtiene
$4u^{2}-v^{2}=4\left( uv\right) ^{2}-1,$ es decir, $4u^{2}\left(1-v^{2}\right) =v^{2}-1;$
por lo tanto $u=0$ o $v=\pm 1.$ De lo anterior se obtienen los límites para la región en el plano $uv.$

El jacobiano de la transformación es
\[\left| \begin{array}{cc}
\frac{\partial x}{\partial u} & \frac{\partial x}{\partial v} \\[0.2cm]
\frac{\partial y}{\partial u} & \frac{\partial y}{\partial u}
\end{array}
\right| =\left| \begin{array}{cc}8u & -2v \\[0.2cm]
v & u \end{array} \right| = 8u^{2}+2v^{2}.\]
Entonces,
\begin{eqnarray*}
\int_{-1}^{1}\int_{4y^{2}-1}^{4-y^{2}}y^{2}dxdy &=&
\int_{0}^{1}\int_{-1}^{1}\left( uv\right) ^{2}\left( 8u^{2}+2v^{2}\right) dudv \\
&=& \int_{0}^{1}\left( \frac{16}{5}v^{2}+\frac{4}{3}v^{4}\right), dv = \frac{4}{3}.
\end{eqnarray*}

\item Evaluar $\ds \iint_R xy dA,$ donde $R$ es la región del primer cuadrante acotada por las figuras de las expresiones
$y=1/x$, $y=3x$, $y=x$, $y=3/x$.




\begin{figure}[H]
\centering
\caption{${\bf T}(u,v)= (\sqrt{uv},\sqrt{\frac{u}{v}})$}
\includegraphics[width=0.8\textwidth]{Fig/113.pdf}
\fuentepropia
\end{figure}

Con las expresiones $y=1/x$ y $y=3/x,$ se puede hacer el cambio $xy=u,$ de esta manera $u\in \lbrack 1,3].$ Mientras que con $y=x$ y $y=3x,$ se puede hacer la sustitución $v=y/x,$ de modo que $v\in \lbrack 1,3].$
El jacobiano de la transformación es el recíproco del determinante de la matriz
\[\left| \begin{array}{cc}
\frac{\partial u}{\partial x} & \frac{\partial u}{\partial y} \\[0.1cm]
\frac{\partial v}{\partial x} & \frac{\partial v}{\partial y}
\end{array} \right| =\left| \begin{array}{cc}y & x \\ -y/x^{2} & 1/x\end{array}
\right| = \frac{2y}{x}=2v.\]
Entonces, el jacobiano de la transformación es $\frac{1}{2v};$ ahora bien,
{\footnotesize \[\iint_R xy dA=\int_{1}^{3}\int_{1}^{3}\frac{u}{2v}dudv= 2\ln 3.\]}
Esta integral también puede evaluarse en la forma
{\footnotesize
\begin{eqnarray*}
\iint_R xy dA &=& \int_{\frac{1}{\sqrt{3}}}^{1}\int_{\frac{1}{x}}^{3x}xydydx+\int_{1}^{\sqrt{3}}\int_{x}^{3/x}xydydx \\
&=& \int_{1}^{\sqrt{3}}\int_{\frac{y}{3}}^{y}xydxdy + \int_{\sqrt{3}}^3\int_{\frac{y}{3}}^{\frac{3}{y}}xydxdy =2\ln 3.
\end{eqnarray*}
}

Estas integrales se evalúan ejecutando las siguientes instrucciones:



\begin{tcolorbox}[breakable,enhanced,title=\bf Instrucción en \verb"Mathematica", top=2pt,bottom=2pt,boxsep=0pt,before skip=2pt,after skip=0pt]
\begin{Verbatim}[formatcom=\letraverbatim]
Integrate[x*y,{y,1,Sqrt[3]},{x,1/y,y}]+Integrate[x*y,
{y,Sqrt[3],3},{x,y/3,3/y}]
Integrate[x*y,{x,1/Sqrt[3],1},{y,1/x,3*x}]+
Integrate[x*y,{x,1,Sqrt[3]},{y,x,3/x}]
\end{Verbatim}
\end{tcolorbox}



\item Evaluar $\ds \iint_R xy dA,$ en donde $R$ es la región en el primer cuadrante acotada por las graficas de las expresiones
$x^{2}y=1,$ $x^{2}y=4,$ $y=x^{2}$ y $y=4x^{2}.$



\begin{figure}[H]
\centering
\caption{${\bf T}(u,v)= (\sqrt[4]{\frac{u}{v}},\sqrt{uv})$}
\includegraphics[width=0.8\textwidth]{Fig/114.pdf}
\end{figure}

Los puntos de intersección de estas curvas son $\left( 1,1\right) ,$ $\left( \sqrt{2},2\right) ,$
$\left( \frac{1}{\sqrt{2}},2\right) $ y $\left( 1,4\right)$; con esta información, se puede evaluar esta integral en coordenadas rectangulares, es decir

{\small
\begin{eqnarray*}
\ds \iint_R xy dA &=& \bi_{\frac{1}{\sqrt{2}}}^{1}\bi_{1/x^{2}}^{4x^{2}}xydydx+\bi_{1}^{\sqrt{2}}\bi_{x^{2}}^{4/x^{2}}xydydx \\
&=& \bi_{1}^{2}\bi_{\sqrt{\frac{1}{y}}}^{\sqrt{y}}xydxdy+\bi_{2}^{4}\bi_{\sqrt{\frac{y}{4}}}^{\sqrt{\frac{4}{y}}}xydxdy =\frac{7}{3}
\end{eqnarray*}
}
Estas integrales se evalúan con las siguientes instrucciones:



\begin{tcolorbox}[breakable,enhanced,title=\bf Instrucción en \verb"Mathematica", top=2pt,bottom=2pt,boxsep=0pt,before skip=2pt,after skip=0pt]
\begin{Verbatim}[formatcom=\letraverbatim]
Integrate[x*y,{y,2,4},{x,Sqrt[y/4],Sqrt[4/y]}]+
Integrate[x*y,{y,1,2},{x,Sqrt[1/y],Sqrt[y]}]
		
Integrate[x*y,{x,1/Sqrt[2],1},{y,1/x^2,4*x^2}]+
Integrate[x*y,{x,1,Sqrt[2]},{y,x^2,4/x^2}]
\end{Verbatim}
\end{tcolorbox}




También se puede utilizar el teorema de cambio de variable, definiendo $u=x^{2}y$ y $v=y/x^{2},$ de donde $x=\sqrt[4]{\frac{u}{v}},\,\, y=\sqrt{u}\sqrt{v}$,
por lo tanto, el jacobiano de la transformación es
{\small \[\left| \begin{array}{cc} \frac{\partial }{\partial u}\left( \sqrt[4]{\frac{u}{v}}\right) & \frac{\partial }{\partial v}\left( \sqrt[4]{\frac{u}{v}}\right) \\[0.1cm]
\frac{\partial }{\partial u}\left( \sqrt{uv}\right) & \frac{\partial }{\partial v}\left( \sqrt{uv}\right) \end{array} \right| = \frac{1}{4\sqrt[4]{uv^{3}}}.\]}
A partir de la definición, es posible establecer que $1\leq u \leq 4$ y $1\leq v \leq 4$, respectivamente.
{\small
\begin{multline*}
\iint_R xy dA = \int_{1}^{4}\int_{1}^{4}\left( \sqrt[4]{\frac{u}{v}}\right) \left( \sqrt{u}\sqrt{v}\right) \frac{1}{4\sqrt[4]{uv^{3}}}\, dudv\\
=\int_{1}^{4}\int_{1}^{4}\frac{1}{4}\frac{\sqrt{u}}{\sqrt{v}}\, dudv= \frac{7}{3}.
\end{multline*}
}

\item
Evaluar la integral {\footnotesize $\ds\int_{0}^{1}\int_{2y}^{\sqrt{1+3y^2}}xy^{2}dxdy+\int_{-1}^{0}\int_{-2y}^{\sqrt{1+3y^2}}xy^{2}dxdy.$} La región de integración está limitada por las figuras de las ecuaciones $x= \pm 2y\,$ y
$\, x^2-3y^2=1,$ con $x\geq 0$.



\begin{figure}[H]
\centering
\caption{${\bf T}(u,v)= (u\cosh v, \frac{u}{\sqrt{3}}\sinh v)$}
\includegraphics[width=0.4\textwidth]{Fig/115.pdf} \hspace{5mm}
\includegraphics[width=0.4\textwidth]{Fig/116.pdf}
\end{figure}

Puede probarse que
\[ \int_{0}^{1}\int_{2y}^{\sqrt{1+3y^{2}}}xy^{2}dxdy+\int_{-1}^{0}\int_{-2y}^{\sqrt{1+3y^{2}}}xy^{2}dxdy= \frac{2}{15}.\]
El siguiente proceso permite evaluar las anteriores integrales de una manera distinta.
 
Al hacer $x=u\cosh v,$ $y=\frac{u}{\sqrt{3}}\sinh v,$ se sigue que la expresión $x^{2}-3y^{2}=1$ se transforma en
$\left( u\cosh v\right) ^{2}-3(\frac{u}{\sqrt{3}}\sinh v)^2=1,$
que se reduce a $u^{2}=1,$ es decir, $u=\pm 1.$
Para la recta $y=x/2,$ se sigue que
$2\left( \frac{u}{\sqrt{3}}\right) \sinh v=u\cosh v$, es decir $v=\ln \left(
2+\sqrt{3}\right) ,$ mientras que para la recta $y=-x/2,$ se concluye que
$v=-\ln \left(2+\sqrt{3}\right). $
  
Con el cambio propuesto, el jacobiano de la transformación es
%
\begin{eqnarray*}
\left| \begin{array}{ll}
\frac{\partial }{\partial u}\left( u\cosh v\right) & \frac{\partial }{\partial v}\left( u\cosh v\right)  \\[0.2cm] \frac{\partial }{\partial u}\left( \frac{u}{\sqrt{3}}\sinh v\right)  & \frac{
\partial }{\partial v}\left( \frac{u}{\sqrt{3}}\sinh v\right)
\end{array}\right| &=& \frac{\sqrt{3}}{3}u\cosh ^{2}\left( v\right) -\frac{\sqrt{3}}{3}u\sinh ^{2}\left( v\right) \\&=&  \frac{1}{3}u\sqrt{3}.
\end{eqnarray*}
Entonces, la integral en mención se transforma en
\[
\int_{0}^{1} \int_{-\ln \left( 2+\sqrt{3} \right)}^{\ln \left( 2+\sqrt{3} \right)}
\left( u \cosh v \right)
\left( \frac{u}{\sqrt{3}} \sinh v \right)^2
\left( \frac{\sqrt{3}}{3} u \right)
\, dv \, du = \frac{2}{15}.
\]

%

\begin{tcolorbox}[breakable,enhanced,title=\bf Instrucción en \verb"Mathematica", top=2pt,bottom=2pt,boxsep=0pt,before skip=2pt,after skip=0pt]
\begin{Verbatim}[formatcom=\letraverbatim]
Integrate[x*y^2,{y,0,1},{x,2*y,Sqrt[1+3*y^2]}]+
Integrate[x*y^2,{y,-1,0},{x,-2*y,Sqrt[1+3*y^2]}]

Integrate[u*Cosh[v]*(u/Sqrt[3]/3*u),
{u,0,1},{v,-Log[2+Sqrt[3]],Log[2+Sqrt[3]]}]
\end{Verbatim}
\end{tcolorbox}


 

\item Sea $R$ el trapecio con vértices $\left( 0,2\right),$ $\left( 1,1\right),\left( 2,2\right),$ $\left( 1,3\right).$
En coordenadas rectangulares, la integral $\ds \iint _{R}\frac{x-y}{x+y}dA$ se puede evaluar como se muestra a continuación:
\[I=\int_{0}^{1}\int_{2-x}^{2+x}\frac{x-y}{x+y}dydx+\int_{1}^2\int_{x}^{4-x}\frac{x-y}{x+y}dydx=-\ln \left( 2\right).\]
Al hacer el cambio $u=x+y,$ $v=y-x,$ se sigue que $y=\frac{u+v}{2},$ $x=\frac{u-v}{2}$. 



\begin{figure}[H]
\centering
\caption{\small ${\bf T}(u,v)= \left(\frac{u - v}{2},\, \frac{u + v}{2}\right)$}
\includegraphics[width=0.4\textwidth]{Fig/117.pdf} \hspace{5mm}
\includegraphics[width=0.4\textwidth]{Fig/118.pdf}
\end{figure}

La figura muestra la región transformada; el jacobiano de la transformación es
$J=\det \left( \begin{array}{cc}
\frac{1}{2} & -\frac{1}{2} \\[0.2cm]
\frac{1}{2} & \frac{1}{2} \end{array} \right) = \frac{1}{2}$; luego
$I=\ds \int_{2}^{4}\ds\int_{0}^2\frac{-v}{u}\left( \frac{1}{2}\right) dvdu= -\ln \left( 2\right) .$

\item Evaluar $\ds \iint_R \sqrt{xy}\ln \left( \frac{y}{x}\right) dA,$ en donde $R$ es la región en el plano $xy$ determinada por el cambio de variables
$u^{2}=xy,$ $v=\ln \left( y/x\right) ,$ con $u\in \lbrack 1,2],$ $v\in \lbrack 0,2].$ \\[0.2cm]

Con el cambio de variable se sigue que $e^{v}=y/x,$: es decir, $y=xe^{v},$
que, al reemplazar en la primera ecuación, se obtiene que
$u^{2}=x^{2}e^{v},$ entonces $x=ue^{-v/2},$ y $y=ue^{v/2}.$
Para determinar las ecuaciones de las curvas que acotan la región, lo haremos por casos.
\begin{itemize}[label=$\bullet$\,\,]
\item Si $u=1$ y $0\leq v\leq 2,$
entonces, $x=e^{-v/2}$ y $y=e^{v/2},$ por tanto $xy=1.$
\item Si $u=2$ y $0\leq v\leq 2,$ entonces \
$x=2e^{-v/2}$ y $y=2e^{v/2}$ es decir $xy=4.$
\item Si $v=0$ y $1\leq u\leq 2,$
se sigue que $x=u$ y $y=u$ entonces $x=y.$
\item Si $v=2$ y $1\leq u\leq 2$
Entonces, $x=ue^{-1}$ y $y=ue$; por lo tanto, $y=e^{2}x.$
\end{itemize}

El siguiente gráfico muestra la región de integración. Los puntos de corte de las expresiones que determinan dicho conjunto en el plano $xy$
son $(e^{-1},e),$ $(1,1)$, $(2e^{-1},2e)$ y $(2,2).$ También se muestra la región en el plano $uv,$
que se obtiene al aplicar la transformación.

\begin{figure}[H]
	\centering
	\caption{${\bf T}(u,v)= (ue^{-v/2},ue^{v/2})$}
	\begin{minipage}{0.37\linewidth}
		\centering
		\includegraphics[width=\linewidth]{Fig/119.pdf}
	\end{minipage}
	\hspace{1em}
	\begin{minipage}{0.36\linewidth}
		\centering
		\includegraphics[width=\linewidth]{Fig/120.pdf}
	\end{minipage}
	
	
	\fuentepropia	
\end{figure}

La integral en mención puede evaluarse usando coordenadas rectangulares como sigue
\begin{multline*}
\int_{e^{-1}}^{1}\int_{1/x}^{e^{2}x}\sqrt{xy}\ln \left( \frac{y}{x}\right)
dydx+\int_{1}^{2e^{-1}}\int_{x}^{e^{2}x}\sqrt{xy}\ln \left( \frac{y}{x}\right) dydx\\
+\int_{2e^{-1}}^{2}\int_{x}^{4/x}\sqrt{xy}\ln \left( \frac{y}{x}\right) dydx= \frac{14}{3}
\end{multline*}
Esta integral se evalúa utilizando las siguientes instrucciones.



\begin{tcolorbox}[breakable,enhanced,title=\bf Instrucción en \verb"Mathematica", top=2pt,bottom=2pt,boxsep=0pt,before skip=2pt,after skip=0pt]
\begin{Verbatim}[formatcom=\letraverbatim]
Integrate[Sqrt[x*y]*Log[y/x],{x,Exp[-1],1},
{y,1/x,Exp[2]*x}]+
Integrate[Sqrt[x*y]*Log[y/x],{x,1,2*Exp[-1]},
{y,x,Exp[2]*x}]+
Integrate[Sqrt[x*y]*Log[y/x],{x,2*Exp[-1],2},
{y,x,4/x}]
\end{Verbatim}
\end{tcolorbox}



Utilizando el cambio de variables, el jacobiano de la transformación es
\[
\det \left(
\begin{array}{cc}
\frac{\partial }{\partial u}\left( ue^{-v/2}\right)  & \frac{\partial }{\partial v}\left( ue^{-v/2}\right)  \\[0.2cm]
\frac{\partial }{\partial u}\left( ue^{v/2}\right)  & \frac{\partial }{\partial v}\left( ue^{v/2}\right)
\end{array}
\right)
=
\det \left(
\begin{array}{cc}
e^{-\frac{1}{2}v} & -\frac{1}{2}ue^{-\frac{1}{2}v} \\
e^{\frac{1}{2}v} & \frac{1}{2}ue^{\frac{1}{2}v}
\end{array}
\right)
= u
\]

Con lo cual, la integral se reduce a
\[\iint_R \sqrt{xy}\ln \left( \frac{y}{x}\right) dA= \int_{1}^{2}\int_{0}^{2}u^{2}vdvdu= \frac{14}{3}\]

\item Aplicar el cambio de variables $u=\sqrt{xy},\,\, v=\sqrt{\frac{y}{x}},$ en la integral
\[\int_{1}^{4} \int_{1/y}^{y}\sqrt{\frac{y}{x}}e^{\sqrt{xy}}dxdy.\]
Del cambio se sigue que
\[uv=\sqrt{xy}\sqrt{\frac{y}{x}}= y;\] además, \[x=\frac{u^{2}}{uv}= \frac{u}{v}.\]
El jacobiano en este caso es
\[\left|  \begin{array}{cc}
\frac{1}{v} & v \\[0.1cm]
\frac{-u}{v^{2}} & u
\end{array} \right| = \frac{2u}{v}.\]
Las respectivas regiones en los planos $xy$ y $uv$
se muestran enseguida.

\begin{figure}[h!]
	\centering
	\caption{${\bf T}(u,v)= \left(\frac{u}{v},\, uv\right)$}
	\begin{minipage}{0.35\linewidth}
		\centering
		\includegraphics[width=\linewidth]{Fig/121.pdf}
	\end{minipage}
	\hspace{1em}
	\begin{minipage}{0.40\linewidth}
		\centering
		\includegraphics[width=\linewidth]{Fig/122.pdf}
	\end{minipage}
	
	
	\fuentepropia
\end{figure}

Por ello, la integral se transforma en
\[\int_{1}^{4}\int_{1/y}^{y}\sqrt{\frac{y}{x}}e^{\sqrt{xy}}dxdy
=\int_{1}^{4}\int_{1}^{4/u}2ue^{u}dvdu = 2e^4-8e\]
%\begin{eqnarray*} \int_{1}^{4}\int_{1/y}^{y}\sqrt{\frac{y}{x}}e^{\sqrt{xy}}dxdy &=&
%\int_{1}^{4}\int_{1}^{4/u}ve^{u}\left( \frac{2u}{v}\right)dvdu \\
%&=&\int_{1}^{4}\int_{1}^{4%/u}2ue^{u}dvdu = 2e^4-8e %\end{eqnarray*}

Estas integrales se pueden evaluar con ayuda de las siguientes sentencias.

\vspace{3pt}

\begin{tcolorbox}[breakable,enhanced,title=\bf Instrucción en \verb"Mathematica", top=2pt,bottom=2pt,boxsep=0pt,before skip=2pt,after skip=0pt]
\begin{Verbatim}[formatcom=\letraverbatim]
Integrate[Sqrt[y/x]*Exp[Sqrt[x*y]],{y,1,4},{x,1/y,y}]
Integrate[2*u*Exp[u],{u,1,4},{v,1,4/u}]
\end{Verbatim}
\end{tcolorbox}

\end{enumerate}

\begin{ejer}
Use el teorema de cambio de variable para resolver los siguientes problemas

\begin{enumerate}
\item   Evaluar
$\displaystyle \iint_R (x^2+y^2)dA, $ donde $R$ es la región en el primer cuadrante acotada por las figuras de las expresiones $x^{2}-y^{2}=3,$ $y^{2}-x^{2}=3,$ $xy=2,$ $xy=3\sqrt{2}.$
Muestre también que el área de la región mencionada está dada por $ -6+9\sqrt{2}.$



\begin{figure}[H]
\centering
\caption{${\bf T}(u,v)= (x,y)$}
\includegraphics[width=0.8\textwidth]{Fig/123.pdf}
\end{figure}

\begin{comment}
$ \int_{1}^{\sqrt{3}}\int_{2/x}^{\sqrt{3+x^{2}}}(x^2+y^2)dydx+\int_{\sqrt{3}}^{2}\int_{2/x}^{3\sqrt{2}/x}(x^2+y^2)dydx+\int_{2}^{\sqrt{6}}\int_{\sqrt{x^{2}-3}}^{3\sqrt{2}/x}(x^2+y^2)dydx=-6+9\sqrt{2}$

Es claro que $J= \frac{1}{2(x^2+y^2)}.$
Con lo cual $\int_{-3}^{3}\int_{2}^{3\sqrt{2}} \frac{1}{2} dudv= -6+9\sqrt{2} $
\end{comment}

\item Determine el área de la región $R$ acotada en el primer cuadrante para las curvas
$x^2+y^2=2$, $x^2+y^2=4,$ $x^2-y^2=2$, $x^2+y^2=4$.



\begin{figure}[H]
\centering
\caption{${\bf T}(u,v)= (x,y)$}
\includegraphics[width=0.8\textwidth]{Fig/124.pdf}
\end{figure}

\end{enumerate}
\end{ejer}

\section{Áreas de superficies parametrizadas} \index{Superficies parametrizadas!áreas}
Además de parametrizar algunas superficies, hallaremos su área. Esta tarea se hará básicamente utilizando la expresión
\ref{areasup}.
\subsection*{Un cono truncado oblicuo}
Las figuras de las funciones vectoriales $\overrightarrow{\alpha}(t)=\langle 1+\cos(t) ,\sin(t),2\rangle $ y
$\overrightarrow{\beta}(t)=\langle 3\cos (t),3\sin(t),0\rangle ,$ para $t\in [0,2\pi],$ son dos circunferencias.
Formando una combinación convexa con un punto $P$ sobre la curva $\overrightarrow{\alpha}$ y un punto $Q$ sobre la curva $\overrightarrow{\beta},$ se genera un cono truncado oblicuo.


\begin{figure}[H]
\centering
\caption{Cono truncado oblicuo determinado por ${\bf r}(u,t)$}
\includegraphics[width=0.4\textwidth]{Fig/125.pdf} \hspace{5mm}
\includegraphics[width=0.4\textwidth]{Fig/126.pdf}
\fuentepropia
\end{figure}

Una parametrización para esta región es la siguiente.
\begin{equation*}
\resizebox{0.97\textwidth}{!}{$
\begin{aligned}
{\bf r}(u,t) &= (1-u) \overrightarrow{\alpha}(t) + u\overrightarrow{\beta}(t) \\
&= \langle 1 - u + (1 + 2u)\cos t,\ (1 + 2u)\sin t,\ 2 - 2u \rangle,\!\!\! u \in [0,1],\ t \in [0,2\pi].
\end{aligned}
$}
\end{equation*}
%$u\in[0,1],$ %$t\in[0,2\pi].$
Para calcular el área superficial, se requieren los siguientes vectores
{\small
\begin{eqnarray*}
{\bf r}_{t} &=& \langle -\sin t( 1+2u) ,( 1+2u) \cos t,0\rangle \\
{\bf r}_{u} &=& \langle -1+2\cos t,2\sin t,-2\rangle. \end{eqnarray*}
}
además,
$ {\bf r}_{t}\times {\bf r}_{u}= \langle -2( 1+2u) \cos t,-2( \sin t) (1+2u) ,( \cos t-2) ( 1+2u) \rangle .$
El elemento diferencial de área es
$\| {\bf r}_{t}\times {\bf r}_{u}\| = (1+2u) \sqrt{\cos ^{2}t-4\cos t+8}.$
Por último, el área superficial está dada por
{\small
\[\int_{0}^{1}\int_{0}^{2\pi }( 1+2u) \sqrt{\cos ^{2}t-4\cos t+8}\, dtdu\approx 36.117.\]}
Este valor se consiguió utilizando las siguientes instrucciones.



\begin{tcolorbox}[breakable,enhanced,title=\bf Instrucción en \verb"Mathematica", top=2pt,bottom=2pt,boxsep=0pt,before skip=2pt,after skip=0pt]
\begin{Verbatim}[formatcom=\letraverbatim]
r=(1-u)*{1+Cos[t],Sin[t],2}+u*{3*Cos[t],3*Sin[t],0};
A=Assuming[t>=0&&u>=0,Simplify[Norm[Cross[D[r,u],
D[r,t]]]]]
NIntegrate[A,{t,0,2*Pi},{u,0,1}]
\end{Verbatim}
\end{tcolorbox}



\index{Cilindro cardiodico}
\subsection*{Una esfera y dos cilindros perpendiculares}
La esfera centrada en el origen de radio $2$ y los cilindros con ecuaciones son $x^{2}+y^{2}=2x$ y $x^{2}+z^{2}=2x$
se intersecan. Se determinará la superficie sobre la esfera acotada por los cilindros, y de esta nos interesará su área. % Los siguientes gráficos ilustran la situación mencionada.


\begin{figure}[H]
\centering
\caption{Esfera y dos cilindros}
\includegraphics[width=0.4\textwidth]{Img/6g.pdf} \hspace{5mm}
\includegraphics[width=0.4\textwidth]{Img/6a.pdf}

\includegraphics[width=0.4\textwidth]{Img/6b.pdf}
\fuentepropia
\end{figure}

La porción sobre la esfera interior de los cilindros se parametriza en dos casos:

\newlength{\bulletindent} 
\settowidth{\bulletindent}{\textcolor{black}{\scriptsize\ding{110}}\quad}
\begin{minipage}[t]{\textwidth}
 
\begin{itemize}[
leftmargin=\dimexpr 1.5cm - \bulletindent\relax,itemsep=0pt,parsep=0pt,labelsep=0.4em,label=\textcolor{black}{\scriptsize\ding{110}}]
\item El cilindro $x^2+y^2=2r$ se reescribe como $(x-1)^2+y^2=1$. Para la superficie sobre la esfera se puede hacer $x=1+u\cos(t)$, $y=u\sin(t)$, donde $t\in[0,2\pi]$, lo cual parametriza la región sobre la esfera interior al cilindro proyectada el plano $xy$. Reemplazando en la ecuación de la esfera se obtiene {\small $-3+2u\cos(t)+u^2+z^2=0$}, cuya solución es {\small $z=\pm\sqrt{3-2u\cos(t)-u^2}$}. Por lo tanto la función
\end{itemize}
\end{minipage}

{\footnotesize
\[{\bf r}( u,t) =\left\langle 1+u\cos (t) ,u\sin \left(
t\right) ,\sqrt{3-2u\cos (t) -u^{2}}\right\rangle,\]}
con $t\in \lbrack 0,2\pi ],\,\, u\in \lbrack 0,1],$
representa la parte de la esfera que se halla acotada por el cilindro, la otra parte se determina con la raíz negativa en la tercera componente.

%\clearpage


\begin{figure}[H]
\centering
\caption{Curva de intersección de la esfera y los cilindros}
\includegraphics[width=0.4\textwidth]{Img/6f.pdf} \hspace{5mm}
\includegraphics[width=0.4\textwidth]{Img/6c.pdf}
\fuentepropia
\end{figure}

\settowidth{\bulletindent}{\textcolor{black}{\scriptsize\ding{110}}\quad}
\begin{minipage}[t]{\textwidth}


\begin{itemize}[
leftmargin=\dimexpr 1.5cm - \bulletindent\relax,itemsep=0pt,parsep=0pt,labelsep=0.4em,before=\vspace{-0.4cm},label=\textcolor{black}{\scriptsize\ding{110}}]
\item Para la expresión $x^{2}+z^{2}=2x$ (que representa el otro cilindro), el tratamiento es similar al caso anterior, obteniendo la función

{\footnotesize
\[
{\\bf r}_1(u,t) = \left\langle 1 + u\cos(t),\, \pm \sqrt{3 - 2u\cos(t) - u^{2}},\, u\sin(t) \right\rangle
\]
}

con $t \in [0, 2\pi]$, $u \in [0, 1]$, la cual parametriza la porción de la esfera dentro del cilindro con ecuación $x^{2}+z^{2}=2x$.
\end{itemize}
\end{minipage}

%


\begin{figure}[H]
\centering
\caption{Curva de intersección de la esfera y los cilindros}
\includegraphics[width=0.4\textwidth]{Img/6d.pdf} \hspace{5mm}
\includegraphics[width=0.4\textwidth]{Img/6e.pdf}

\includegraphics[width=0.4\textwidth]{Fig/127.pdf}
\fuentepropia
\end{figure}

Los anteriores gráficos se obtienen ejecutando las siguientes instrucciones:


\begin{enumerate}[leftmargin=0cm,itemindent=.75cm,labelwidth=\itemindent,labelsep=0cm,align=left,label={{\bf \arabic*.\,}}]
\item La esfera y los cilindros

Región sobre la esfera acotada por los cilindros, se definen tres objetos nombrados \verb"a,b,c", luego se muestran en un solo gráfico.



\begin{tcolorbox}[breakable,enhanced,title=\bf Instrucción en \verb"Mathematica", top=2pt,bottom=2pt,boxsep=0pt,before skip=2pt,after skip=0pt]
\begin{Verbatim}[formatcom=\letraverbatim]
ContourPlot3D[{x^2+y^2+z^2==4,x^2+y^2==2*x,x^2+
z^2==2*x},{x,-2,2},{y,-2,2},{z,-2,2},PlotPoints->80,
Mesh->False]
ContourStyle->{Orange,Yellow,LightGreen}
\end{Verbatim}
\end{tcolorbox}




\item Región sobre la esfera acotada por los cilindros, se definen tres objetos nombrados \verb"a,b,c", luego se muestran en una sola figura usando la función \textit{show}.



\begin{tcolorbox}[breakable,enhanced,title=\bf Instrucción en \verb"Mathematica", top=2pt,bottom=2pt,boxsep=0pt,before skip=2pt,after skip=0pt]
\begin{Verbatim}[formatcom=\letraverbatim]
Clear[x,y,z]
a=ContourPlot3D[x^2+y^2+z^2==3.95,{x,-2,2},{y,-2,2},
{z,-2,2},ContourStyle->{Cyan,Opacity[0.4]},Mesh->Fals
PlotPoints->50];
r1={x,y,z}={1+Cos[t],Sqrt[2-2*Cos[t]],Sin[t]};
r5={x1,y1,z1}={1+u*Cos[t],Sqrt[3-2*u*Cos[t]-u^2],
u*Sin[t]};
b=ParametricPlot3D[{r1,{x,-y,z},{x,z,y},{x,z,-y}},
{t,0,2*Pi},PlotStyle->Black,PlotPoints->80];
c=ParametricPlot3D[{r5,{x1,-y1,z1},{x1,z1,y1},
{x1,z1,-y1}},{t,0,2*Pi},{u,0,1},
PlotStyle->Blend[{Yellow,Pink,Green}]];Mesh->False,
Show[a,b,c]
\end{Verbatim}
\end{tcolorbox}



Por la simetría de la región en mención, el área de la porción sobre la esfera interior a los cilindros corresponde a cuatro veces la integral del término
$\left\Vert {\bf r}_{t}\times {\bf r}_{u}\right\Vert,$ que es el elemento diferencial del área.
Para tal fin, se hallan las siguientes funciones vectoriales:

\begin{eqnarray*}
{\bf r}_{t}(t,u) &=& \la -u\sin t,u\cos t,\frac{u\sin t}{\sqrt{3-u^2 -2u\cos t}}\ra \\
{\bf r}_{u}(t,u)&=& \la\cos t,\sin t,-\frac{u+\cos t}{\sqrt{3-u^2-2u\cos t}}\ra \\
{\bf r}_{t}\times{\bf r}_{u}&=&\la-\frac{u^2 \cos t+u}{\sqrt{3-u^2 -2u\cos t}},\frac{-u^2 \sin t}{\sqrt{3-u^2 -2u\cos t}},-u,\ra,
\end{eqnarray*}
\end{enumerate}
entonces
$\left\Vert {\bf r}_{t}\times{\bf r}_{u} \right\Vert =\frac{2u}{\sqrt{3-u^{2}-2u\cos t}}.$
Por lo anterior, el área $A$ de la superficie es cuatro veces la integral del término anterior; es decir,

\[4\int_{0}^{2\pi}\int_{0}^{1}\frac{2u}{\sqrt{3-u^{2}-2u\cos t}}dudt \approx 18. 265.\]

Este cálculo se hace con \verb"Mathematica" ejecutando las siguientes instrucciones:



\begin{tcolorbox}[breakable,enhanced,title=\bf Instrucción en \verb"Mathematica", top=2pt,bottom=2pt,boxsep=0pt,before skip=2pt,after skip=0pt]
\begin{Verbatim}[formatcom=\letraverbatim]
{x,y,z}={1+u*Cos[t],Sqrt}[3-u^2-2*u*Cos[t]]};
r={x,y,z};
A=Norm[Cross[D[r,t],D[r,u]]];
NIntegrate[4*A,{u,0,1},{t,0,2Pi}]
\end{Verbatim}
\end{tcolorbox}



\subsection*{Un cilindro cardiogénico bajo un paraboloide}

Las ecuaciones $3z=9-(x+2)^2-(y+1)^2$ y $r=1-\cos(t)$ representan un paraboloide y un cilindro cardiodónico, respectivamente.
Estas superficies, junto con el plano $z=0$, acotan un sólido.
En el plano $xy$ el cilindro se proyecta como un cardioide que puede parametrizarse como
\[\overrightarrow{s}(t)=\la(1-\cos(t))\cos(t),(1-\cos(t))\sin(t),0\ra, \,\, t \in [0,2\pi].\]


Sobre el paraboloide, la coordenada $z$ tiene la forma

{\small
\begin{align*}
z &= \frac{1}{3}\left(9 - \left((1 - \cos(t))\cos(t) + 2\right)^2 - \left((1 - \cos(t))\sin(t) + 1\right)^2\right) \\
&= \frac{1}{6}(13 - 4\sin(t) + 2\sin(2t) + \cos(2t))
\end{align*}
}




\begin{figure}[H]
\centering
\caption{\small Cilindro cardiogénico y paraboloide}
\includegraphics[width=0.4\textwidth]{Img/40a.pdf} \hspace{5mm}
\includegraphics[width=0.4\textwidth]{Img/40b.pdf}
\fuentepropia
\end{figure}

Por tanto, la curva de intersección se parametriza por medio de la función

{\footnotesize
\begin{multline*} \overrightarrow{r}(t)=\big \langle(1-\cos(t))\cos(t),(1-\cos(t))\sin(t), \\ \frac{1}{3}\left(13-4\sin(t)+2\sin(2t)+\cos(2t)\right)\bigg \rangle,\,\, t \in [0,2\pi].
\end{multline*}
}



\begin{figure}[H]
\centering
\caption{Superficie sobre el cilindro cardiogénico acotada por el paraboloide}
\includegraphics[width=0.4\textwidth]{Fig/128.pdf} \hspace{5mm}
\includegraphics[width=0.4\textwidth]{Img/39c.pdf}
\fuentepropia
\end{figure}

Una combinación convexa de las curvas $\overrightarrow{s}$ y $\overrightarrow{r}$ de la forma $(1-u)\overrightarrow{s}(t)+u\overrightarrow{r}(t),$
parametriza la porción del cilindro que se hallaba bajo el paraboloide y sobre el plano $xy.$
{\small
\begin{multline*} \overrightarrow{p}(t,u)=\big\langle (1-\cos(t))\cos(t),(1-\cos(t))\sin(t), \\ \frac{1}{6}\left(u(13-4\sin(t)+2\sin(2t)+\cos(2t)) \right)\big\rangle,
\end{multline*}
}

Calcularemos el área de esta región. Con esta función es sencillo ver que
\begin{eqnarray*}
\overrightarrow{p}_{t}&=& \bigg \langle-\sin(t)+2\sin(t)\cos(t),\cos(t)-2\cos^{2}(t)+1,\\
&& \hspace{3.8cm} \frac{1}{6} u (-2 \sin (2 t)-4 \cos (t)+4\cos(2t)) \bigg \rangle \\
\overrightarrow{p}_{u}&=& \la 0,0,\frac{1}{6} (-4 \sin (t)+2 \sin (2 t)+\cos (2 t)+13) \ra
\end{eqnarray*}
Con estos vectores obtenemos el siguiente término:
\begin{multline*}
\|\overrightarrow{p}_{t}\times \overrightarrow{p}_{u}\|  =
\frac{1}{6}\bigg |4\cos\left(\frac{t}{2}\right)-6\cos\left(\frac{3t}{2}\right)+2\cos\left(\frac{5t}{2}\right)-22\sin\left(\frac{t}{2}\right)\\
+7\sin\left(\frac{3t}{2}\right)-3\sin\left(\frac{5t}{2}\right)\bigg |
\end{multline*}
Al integrarlo con los límites descritos antes, proporciona el valor del área del cilindro cardiogénico que se encuentra entre el plano $xy$ y el paraboloide, esto es
\[\int_{0}^{2\pi}\int_{0}^{1}\|\overrightarrow{p}_{t}(t,u)\times \overrightarrow{p}_{u}(t,u)\|dudt=\frac{608}{45}.\]
Los cálculos se hacen en \verb"Mathematica", para este fin se define $x$ e $y$, que corresponden a la parametrización de la cardioide en el plano. Sustituyendo esto en la ecuación del paraboloide, se define
$z$, que al unir los puntos permite parametrizar el sólido mediante la expresión $r$.
Luego se obtiene el gráfico de $r$ y se calcula el producto triple que se denomina $A$, este se integra con los límites descritos anteriormente.




\begin{tcolorbox}[breakable,enhanced,title=\bf Instrucción en \verb"Mathematica", top=2pt,bottom=2pt,boxsep=0pt,before skip=2pt,after skip=0pt]
\begin{Verbatim}[formatcom=\letraverbatim]
{x,y}={(1-Cos[t])*Cos[t],(1-Cos[t])*Sin[t]};
z=(9-(x+2)^2-(y+1)^2)/3;
r=Simplify[(1-u)*{x,y,z}+u*{x,y,0}];
ParametricPlot3D[r,{t,0,2*Pi},{u,0,1}];
A=Assuming[t>=0,Simplify[Norm[Cross[D[r,t],D[r,u]]]]]
Integrate[Abs[A],{t,0,2*Pi},{u,0,1}]
\end{Verbatim}
\end{tcolorbox}




\subsection*{Un paraboloide fuera de un cilindro cardiogénico}
Consideremos el paraboloide con ecuación\ $3z=9-x^{2}-y^{2}$ y el cilindro cardiodónico $r=1-\cos(t)$.
La proyección del paraboloide en el plano $xy$ es la circunferencia centrada en el origen del radio $3$,
que puede parametrizarse como $\la 3\cos(t),3\sin(t)\ra $.
Para $t\in[0,2\pi],\, u\in[0,1]$, la región que se encuentra entre la cardioide y la circunferencia se parametriza por:

{\small
\begin{multline*}
(1-u)\la(1-\cos(t))\cos(t),(1-\cos(t))\sin(t)\ra+u\la 3\cos(t),3\sin(t)\ra \\
= \la(1+2u+(u-1)\cos(t))\cos(t),(1+2u+(u-1)\cos(t))\sin(t)\ra.
\end{multline*}
}



\begin{figure}[H]
\centering
\caption{Superficie sobre el paraboloide $3z=9-x^{2}-y^{2}$}
\includegraphics[width=0.4\textwidth]{Fig/129.pdf} \hspace{5mm}
\includegraphics[width=0.4\textwidth]{Fig/130.pdf}
\fuentepropia
\end{figure}

\par 

La coordenada $z$ se obtiene reemplazando $x$ y $z$ en la ecuación del paraboloide, esto es

{\small
\begin{align*}
	z &= \frac{1}{3}(9-x^2-y^2) \\
	&= \frac{1}{3}\left(9 - \left((1+2u + (u-1)\cos(t))\cos(t)\right)^2 \right. \\
	&\quad \left. - \left((1+2u + (u-1)\cos(t))\sin(t)\right)^2\right) \\
	&= \frac{1}{3}\left( (u-1)\cos(t)+2 \right) \left( (1-u)\cos(t) - 2u - 4 \right).
\end{align*}
}
Entonces la superficie sobre el paraboloide que se encuentra fuera del cilindro cardiogénico, y sobre el plano $xy$ se parametriza como:
{\small
\begin{multline*}
{\bf r}(t,u)=\bigg \langle(1+2u+(u-1)\cos(t))\cos(t),(1+2u+(u-1)\cos(t))\sin(t) \\
	\frac{1}{3}(u-1)(\cos(t)+2)((1-u)\cos(t)-2u-4) \bigg \rangle, u\in\lbrack0,1],t\in [0,2\pi].
\end{multline*}
}
%$((\cos(t))(1+2u+(u-1)\cos(t)),(1+2u+(u-1)\cos(t))\sin(t),\frac{1}{3}(u-1)(\cos(t)+2)((1-u)\cos(t)-2u-4))$

\begin{figure}[H]
	
	\centering
	\caption{Superficie sobre el paraboloide \( 3z = 9 - x^2 - y^2 \)}
	\begin{minipage}[t]{0.29\textwidth}
		\centering
		\includegraphics[width=\linewidth]{Img/41a.pdf}
	\end{minipage}
	\hspace{0.02\linewidth}
	\begin{minipage}[t]{0.29\textwidth}
		\centering
		\includegraphics[width=\linewidth]{Img/41b.pdf}
	\end{minipage}

	\fuentepropia
\end{figure}

Con la función ${\bf r}$ hallamos el siguiente vector.
\vspace{-5pt}
{\small
\begin{multline*}
\hspace*{-\multlinegap}
{\bf r}_{t} = \langle \sin(t)(2(1-u)\cos(t)-1-2u),\,
\cos(t)(2(u-1)\cos(t)+1+2u)-u+1, \\
\frac{2}{3}(u-1)\sin(t)(2u+(u-1)\cos(t)+1) \rangle \\[1ex]
{\bf r}_{u} = \langle (\cos(t)+2)\cos(t),\,
(\cos(t)+2)\sin(t),\,
\frac{2}{3}(\cos(t)+2)((1-u)\cos(t)-1-2u) \rangle
\end{multline*}
}

{\small
\begin{multline*}
{\bf r}_{t}\times{\bf r}_{u}=\bigg \langle-\frac{2}{3}\cos(t)(\cos(t)+2)(2u+(u-1)\cos(t)+1)^{2},\\
-\frac{2}{3}\sin(t)(\cos(t)+2)(2u+(u-1)\cos(t)+1)^{2},\\ -(\cos(t)+2)(2u+(u-1)\cos(t)+1)\bigg \rangle
\end{multline*}
}

De modo que el elemento diferencial de área está dado por:


{\small
\[
\lVert \mathbf{r}_t \times \mathbf{r}_u \rVert = \frac{1}{3}(\cos(t) + 2)(2u + (u - 1)\cos(t) + 1)f(u,t),
\]
}

en donde:
\begin{center}
	\resizebox{1.0\textwidth}{!}{%
		$f(u,t) = \sqrt{4\cos(t)(u - 1)((u - 1)\cos(t) + 4u + 2) + 13 + 16u(1 + u)}$%
	}
\end{center}


con lo cual, el área de la superficie es


\[\int_{0}^{1}\int_{0}^{2\pi} \|{\bf r}_{t}\times {\bf r}_{u}\|\,\, dtdu\approx 41.9822.\]

Para hacer los cálculos en \verb"Mathematica" definimos $x$ e $y$, que corresponden a las expresiones para la cardioide. La región entre las curvas se parametriza como \verb"r", mientras que
\verb"p" es la porción sobre el paraboloide cuya base está dada por \verb"r" y \verb"q", que definen la curva de intersección. El gráfico de la curva de intersección y la superficie se logran con las siguientes instrucciones:



\begin{tcolorbox}[breakable,enhanced,title=\bf Instrucción en \verb"Mathematica", top=2pt,bottom=2pt,boxsep=0pt,before skip=2pt,after skip=0pt]
\begin{Verbatim}[formatcom=\letraverbatim]
{x,y}={(1-Cos[t])*Cos[t],(1-Cos[t])*Sin[t]};
z=(9-x^2-y^2)/3;
r=(1-u)*{x,y}+u*{3*Cos[t],3*Sin[t]};
p={r[[1]],r[[2]],(9-r[[1]]^2-r[[2]]^2)/3};
q={x,y,z};


El gráfico de la curva de intersección y la superficie
se logran con las siguientes instrucciones


a=ParametricPlot3D[q,{t,0,2*Pi},PlotStyle->
{Black,Thickness->0.01},PlotPoints->80];
b=ParametricPlot3D[p,{t,0,2*Pi},{u,0,1},Mesh->False,
PlotStyle->{Orange,PlotPoints->120}];


Show[b,a]


Hallamos $\|{\bf r}_{t}\times {\bf r}_{u}\|$ 
que se denomina $A$ y se integra numéricamente,
A=Simplify[Norm[Cross[D[p,t],D[p,u]]]];
NIntegrate[A,{t,0,2*Pi},{u,0,1}]
\end{Verbatim}
\end{tcolorbox}




\subsection*{Un cilindro y dos planos}
Las ecuaciones $4x^{2}+z^{2}=4$, $z=3+2y$ y $z=5-x-3y$ representan un cilindro elíptico y dos planos, respectivamente.
Estas superficies acotan un sólido como se muestra a continuación:




\begin{figure}[H]
\centering
\caption{cilindro $4x^{2}+z^{2}=4$ acotado por dos planos}
\includegraphics[width=0.4\textwidth]{Img/55d.pdf} \hspace{5mm}
\includegraphics[width=0.4\textwidth]{Img/55c.pdf}

\includegraphics[width=0.4\textwidth]{Img/55b.pdf}
\fuentepropia
\end{figure}

Sobre este sólido parametrizado, las tres superficies que lo determinan y calcularemos el área de cada una.

La proyección del cilindro en el plano $xz$ corresponde a la ecuación $4x^{2}+z^{2}=4,$
que puede parametrizarse como $x=\cos(t),$ $z=2\sin (t),$ para $t\in \lbrack 0,2\pi ].$
La coordenada $y$ está determinada por las siguientes expresiones:

En el plano $z=3+2y$ obtenemos $y=\sin (t)-\frac{3}{2}$ y en el plano
$z=5-x-3y$ se sigue que $y=-\frac{2}{3}\sin (t)+\frac{5}{3}-\frac{1}{3}\cos (t).$

Definimos ahora las funciones vectoriales que determinan la intersección del cilindro con cada uno de los dos planos.
\begin{eqnarray*}
\overrightarrow{a}(t)&=&\la \cos(t),\sin(t)-\frac{3}{2},2\sin (t)\ra, \, t\in [0,2\pi] \\
\overrightarrow{b}(t)&=& \la \cos(t) ,-\frac{2}{3}\sin (t)+\frac{5}{3}-\frac{1}{3}\cos (t),2\sin (t)\ra , \, t\in [0,2\pi]
\end{eqnarray*}

\subsubsection*{Área sobre la superficie del cilindro}
Una combinación convexa de las funciones
$\overrightarrow{a}(t)$ y $\overrightarrow{b}(t)$ parametriza la superficie del cilindro que se encuentra entre los dos planos. Es decir,

\begin{multline*}
{\bf r}(u,t) = ( 1-u) \overrightarrow{a}(t)+u\overrightarrow{b}(t) \\
=\la \cos (t),\frac{1}{3}( 3-5u)\sin (t) +\frac{19}{6}u-\frac{1}{3}u\cos (t)-\frac{3}{2},2\sin (t)\ra,
\end{multline*}
con $t\in [0,2\pi],\,u\in [0,1].$
De aquí se obtienen las siguientes funciones.
\begin{eqnarray*}
{\bf r}_{t}&=& \la -\sin(t),\frac{1}{3}(3-5u)\cos(t)+\frac{1}{3}u\sin(t),2\cos(t)\ra \\
{\bf r}_{u}&=& \la 0,-\frac{5}{3}\sin(t)+\frac{19}{6}-\frac{1}{3}\cos(t),0\ra
\end{eqnarray*}

También se calcula
\begin{center}
\resizebox{1.0\textwidth}{!}{%
$\mathbf{r}_{t}\times\mathbf{r}_{u}= \left \langle2\cos(t)\left (\frac{5}{3}\sin(t)-\frac{19}{6}+\frac{1}{3}\cos(t)\right),0,\sin(t)\left (\frac{5}{3}\sin(t)-\frac{19}{6}+\frac{1}{3}\cos(t)\right) \right \rangle,$
	}
\end{center}

cuya magnitud es
\[\|{\bf r}_{t}\times{\bf r}_{u}\|=\frac{\sqrt{5+3\cos(2t)}(8+\cos(t)-4\sin (t))}{3\sqrt{2}};\]
y con esta expresión se obtiene el área, la superficie en mención.
\[\bigintss_{0}^{1}\bigintss_{0}^{2\pi}\frac{\sqrt{5+3\cos(2t)}(8+\cos(t)-4\sin(t))}{3\sqrt{2}} \, dtdu \approx 25.836.\]
Los cálculos anteriores se pueden realizar en \verb"Mathematica"
mediante la ejecución de las siguientes instrucciones:



\begin{tcolorbox}[breakable,enhanced,title=\bf Instrucción en \verb"Mathematica", top=2pt,bottom=2pt,boxsep=0pt,before skip=2pt,after skip=0pt]
\begin{Verbatim}[formatcom=\letraverbatim]
a={Cos[t],(-3+2*Sin[t])/3,2*Sin[t]};
b={Cos[t],(5-Cos[t]-2*Sin[t])/3,2*Sin[t]};
B=Norm[Cross[D[(1-u)*a+u*b,t],D[(1-u)*a+u*b,u]]]
A=Assuming[t>=0&&u>=0,Simplify[B]]
NIntegrate[A,{t,0,2Pi},{u,0,1}]
\end{Verbatim}
\end{tcolorbox}



\subsubsection*{Área sobre los planos}
La región interior a la elipse se parametriza como
$x=u\cos(t),$ $z=2u\sin (t),$ para $t\in \lbrack 0,2\pi ],u\in \lbrack 0,1].$
La coordenada $y$ se puede determinar reemplazando $x$ y $z$ en la ecuación de los planos, de esta manera para  $z=3+2y$ se sigue $y=u\sin (t)-\frac{3}{2},$
mientras que en $z=5-x-3y$ se obtiene que
$y=-\frac{2}{3}u\sin (t)+\frac{5}{3}-\frac{1}{3}u\cos (t).$

Por lo tanto, es natural que se definan las siguientes funciones que parame- \\ trizan la región sobre cada uno de los planos que acotan el cilindro.
\begin{eqnarray*}
\overrightarrow{\alpha }(t,u)&=&\la u\cos (t),u\sin (t)-\frac{3}{2},2u\sin (t)\ra \\
\overrightarrow{\beta }(t,u)&=&\la u\cos (t),-\frac{2}{3}u\sin (t)+\frac{5}{3}-\frac{1}{3}u\cos (t),2u\sin (t)\ra.
\end{eqnarray*}

El área sobre el plano $z=3+2y$ dentro del cilindro puede determinarse con el siguiente procedimiento: primero, se derivan parcialmente las funciones, según el caso.
\begin{eqnarray*}
\overrightarrow{\alpha }{t}&=& \la -u\sin (t),u\cos (t),2u\cos (t)\ra  \\
\overrightarrow{\alpha }{u}&=& \la \cos (t),\sin (t),2\sin (t)\ra.
\end{eqnarray*}
Luego se calcula el elemento diferencial de área 


\begin{align*}
	\|\mathbf{r}_t \times \mathbf{r}_u\| &= \| \langle -u\sin(t), u\cos(t), 2u\cos(t) \rangle \times \langle \cos(t), \sin(t), 2\sin(t) \rangle \| \\
	&= \sqrt{5}u
\end{align*}

Y, por último, este término se integra con los límites establecidos
\[\int_{0}^{2\pi }\int_{0}^{1}\| r_{t}\times r_{u}\|
dudt=\int_{0}^{2\pi }\int_{0}^{1}\sqrt{5}ududt= \sqrt{5}\pi.\]

Para la región sobre el plano con ecuación $z=5-x-3y,$
hallamos los mismos vectores que en el caso anterior
\begin{eqnarray*}
\overrightarrow{\beta}_{t}&=&  \la -u\sin (t),-\frac{1}{3}u( 2\cos (t)-\sin (t))  ,2u\cos (t) \ra  \\
\overrightarrow{\beta}_{u}&=& \la \cos (t),-\frac{2}{3}\sin (t)-\frac{1}{3}\cos (t),2\sin (t)\ra.
\end{eqnarray*}
El elemento diferencial de área en este caso está dado por
\[\| \overrightarrow{\beta}_{t} \times \overrightarrow{\beta}_{u} \|= \frac{2}{3}\sqrt{11}u\]
% $\left \| \la -u\sin (t),-\frac{1}{3}u( 2\cos (t)-\sin (t))  ,2u\cos (t)\ra  \times \la \cos (t),-\frac{2}{3}\sin (t)-\frac{1}{3}\cos (t),2\sin (t)\ra \right \| =  \frac{2}{3}\sqrt{11}u$
y el área de la superficie señalada es
\[\int_{0}^{2\pi }\int_{0}^{1}\| r_{t}\times r_{u}\|
dudt=\int_{0}^{2\pi }\int_{0}^{1}\frac{2}{3}\sqrt{11}ududt= \frac{2}{3}\pi \sqrt{11}.\]

\subsection*{Una porción de paraboloide}
El paraboloide con ecuación $\frac{(z-2)^2}{4}+\frac{(y-3)^2}{9}=1-\frac{x}{4}$
posee vértice en el punto $(4,3,2)$
y abre hacia la parte negativa del eje $x$; consideraremos la porción donde $x \geq 0$ y presentaremos además la proyección en el plano $yz$, que es una elipse.

Para rellenar el interior de la elipse se pueden usar las expresiones \linebreak
$y=3+3u\sin t,z=2+2u\cos t,$ donde $ u\in \lbrack 0,1]$ y $t \in \lbrack 0,2\pi ].$
La coordenada $x$ sobre el paraboloide se obtiene al sustituir $z$ y $y$ en dicha ecuación, es decir,

\begin{figure}[H]
	\centering
	\caption{Paraboloide y su proyección en el plano \( yz \)}
	\begin{minipage}[t]{0.23\textwidth}
		\centering
		
		\includegraphics[width=\linewidth]{Img/32.pdf}
	\end{minipage}
	\hspace{0.02\linewidth}
	\begin{minipage}[t]{0.27\textwidth}
		\centering
		
		\includegraphics[width=\linewidth]{Fig/131.pdf}
	\end{minipage}
	\hspace{0.02\linewidth}
	\begin{minipage}[t]{0.27\textwidth}
		\centering
		
		\includegraphics[width=\linewidth]{Fig/132.pdf}
	\end{minipage}
	\fuentepropia
\end{figure}

\[x=  4\left( 1-\frac{\left( 2u\cos (t) \right) ^{2}}{4}- \frac{\left( 3u\sin (t) \right) ^{2}}{9}\right) =4-4u^{2}.\]

Entonces, la función ${\bf r}(u,t)=\left\langle 4-4u^{2},3+3u\sin t,2+2u\cos
t\right\rangle ,$ $u\in \lbrack 0,1],\,\,t\in \lbrack 0,2\pi ], $
parametriza la superficie del paraboloide, con $x\geq 0.$
Con esta función hallamos los siguientes vectores
\begin{eqnarray*}
{\bf r}_{u} &=& \left\langle -8u,3\sin (t) ,2\cos (t) \right\rangle\\
{\bf r}_{t} &=& \left\langle 0,3u\cos (t) ,-2u\sin (t) \right\rangle.
\end{eqnarray*}
Entonces,
${\bf r}_{t}\times{\bf r}_{u}=\left\langle 6u,16u^{2}\sin (t),24u^{2}\cos (t) \right\rangle,$
por tanto,


\[\left\| {\bf r}_{u}\times {\bf r}_{t}\right\| = 2u\sqrt{9+64u^{2}+80u^{2}\cos ^{2}(t) }.\]
De esta manera, el área de la superficie del paraboloide está dada por
%$\int_{0}^{1}\int_{0}^{2\pi }\left\| {\bf r}_{u}\times {\bf r}_{t}\right\| dtdu  \approx 46.979.$
{\footnotesize
	\begin{equation*}
		\int_{0}^{1}\int_{0}^{2\pi} \|\mathbf{r}_{u} \times \mathbf{r}_{t}\| \, dt\,du =
		\int_{0}^{1}\int_{0}^{2\pi}
		2u\sqrt{9 + 64u^2 + 80u^2 \cos^2(t)} \, dt\,du \approx 46{,}979
	\end{equation*}
}
Este cálculo se hace en \verb"Mathematica" con las siguientes instrucciones



\begin{tcolorbox}[breakable,enhanced,title=\bf Instrucción en \verb"Mathematica", top=2pt,bottom=2pt,boxsep=0pt,before skip=2pt,after skip=0pt]
\begin{Verbatim}[formatcom=\letraverbatim]
y=3+3*u*Sin[t],
z=2+2*u*Cos[t]};
x=Simplify[4*(1-(z-2)^2/4-(y-3)^2/9)]
r={x,y,z};
A=Norm[Cross[D[r,t],D[r,u]]];
NIntegrate[A,{t,0,2Pi},{u,0,1}]
\end{Verbatim}
\end{tcolorbox}



\subsection*{Un cilindro cardióidico y un plano}
La superficie del cilindro cardiogénico con ecuación
$r=2(1+\cos(t))$, acotada por los planos $z=0$ y $z=2-x/2$, se presenta a continuación. También se muestra la porción sobre el plano $z=2-x/2$ acotada por el cilindro cardiogénico. Hallaremos el área de estas dos superficies:


\begin{figure}[H]
\centering
\caption{Superficies sobre el cilindro cardiogénico y sobre el plano \( z = 2 - \frac{x}{2} \)}
\includegraphics[width=0.4\textwidth]{Fig/134.pdf} \hspace{5mm}
\includegraphics[width=0.4\textwidth]{Fig/135.pdf}
\fuentepropia
\end{figure}

Sobre el plano $xy,$ la cardioide $r=2(1+\cos(t))$ se parametriza como
\[\la 2\left( 1+\cos (t) \right) \cos (t) ,2\left(1+\cos (t) \right) \sin (t),0 \ra.\]
En el plano con ecuación $z=2-x/2,$ se satisface que
\[z=2-2\left( 1+\cos (t) \right) \cos(t) /2= 2-\left( 1+\cos t\right) \cos t,\]
por tanto, la curva de intersección entre el plano y el cilindro cardiogénico se parametriza como
\[\la 2\left(1+\cos(t)\right)\cos(t),2\left(1+\cos(t)\right)\sin(t),2-(1+\cos (t))\cos(t) \ra,\]
con lo que forma una combinación convexa entre dos puntos cuyos primer y segundo componente son iguales, uno sobre el plano $z=0$ y el otro sobre el plano
$z=2-x/2.$ Se parametriza la superficie lateral, esto es
\begin{multline*}
{\bf r}(u,t)=\left( 1-u\right) \la 2\left( 1+\cos (t) \right) \cos \left(
t\right) ,2\left( 1+\cos (t) \right) \sin (t) ,0\ra \\
+u\la 2\left( 1+\cos (t) \right) \cos(t) ,2\left( 1+\cos (t) \right) \sin (t)
,2-\left( 1+\cos (t) \right) \cos (t) \\
= \la 2( \cos (t)+1) \cos (t),2\sin (t)+\sin (2t),u(2-\left( 1+\cos (t) \right) \cos (t)) \ra
\end{multline*}
%
$t\in [0,2\pi],$ $u\in [0,1].$ 
\\

Con esta función obtenemos los siguientes vectores:
\begin{eqnarray*}
{\bf r}_{t} &=& \big \langle-2\sin (t) \left( 2\cos (t) +1\right) ,2(\cos(t)+\cos ^{2}(t)-\sin^{2}(t)) , \\
&& \hskip 6cm u\sin (t) \left( 2\cos (t)+1\right) \big \rangle \\
{\bf r}_{u} &=& \la 0,0,\left( \cos (t) +2\right) \left( 1-\cos (t) \right) \ra\\
{\bf r}_{t}\times {\bf r}_{u}&=&\big\langle 2\sin^{2}(t)\left( 2\cos(t) -1\right)\left(\cos(t)+2\right) ,\\
&&\hskip1.5cm 2\sin(t) \left( 2\cos (t) +1\right) \left( \cos (t) +2\right) \left( 1-\cos (t) \right) ,0\big \rangle.
\end{eqnarray*}
Con lo cual,
\[\left\| {\bf r}_{t}\times {\bf r}_{u}\right\| = 2\sin (t) \left( \cos (t) +2\right) \sqrt{2-2\cos(t);}\]
entonces el área lateral del cilindro acotada por los dos planos está dada por

{\small
\[\int_{0}^{1}\int_{0}^{2\pi}\left|2\sin(t)\left(\cos(t)+2\right)\sqrt{2-2\cos(t)}\right|dtdu= \frac{96}{5}.\] 
}
El valor anterior se obtiene en \verb"Mathematica" con ayuda de las siguientes instrucciones:



\begin{tcolorbox}[breakable,enhanced,title=\bf Instrucción en \verb"Mathematica", top=2pt,bottom=2pt,boxsep=0pt,before skip=2pt,after skip=0pt]
\begin{Verbatim}[formatcom=\letraverbatim]
x=2*(1+Cos[t])*Cos[t);
y=2*(1+Cos[t]*Sin[t];
z=2-x/2;
r=(1-u)*{x,y,0}+u*{x,y,z);
A=Norm[Cross[D[r,t],D[r,u]]];
Integrate[A,{u,0,1},{t,0,2Pi}]
\end{Verbatim}
\end{tcolorbox}




Hallaremos ahora el área de la región sobre el plano $z=2-x/2,$ acotada por el cilindro cardiogénico.
El interior de la cardioide en el plano $xy$ se parametriza como
\[\la 2u\left( 1+\cos(t)\right) \cos(t) ,2u\left( 1+\cos(t)\right) \sin(t),0\ra,\]
con $t\in \lbrack 0,2\pi ],$ $u\in \lbrack 0,1].$
Sobre el plano $z=2-x/2,$ se cumple que
\[z=2-2u\left( 1+\cos(t)\right)\cos(t)/2=2-u\left(\cos t+1\right)\cos(t).\]

Entonces, la parte del plano acotada por el cilindro cardiogénico se parame- \\ triza por medio de la siguiente función
\[
\adjustbox{max width=\textwidth}{
	${\bf r}(u,t)=\big\langle 2u\left( 1+\cos(t)\right) \cos(t),2u\left(1+\cos(t)\right)\sin(t),2-u(\cos(t)+1)\cos(t) \big\rangle,$
}
\] con $t\in \lbrack 0,2\pi ],$ $u\in \lbrack 0,1].$ De aquí se obtienen los siguientes vectores:
\begin{align*}
	\mathbf{r}_t &= \left\langle -2u\sin(t)(2\cos(t) + 1), 2u\left(\cos(t) - \sin^2(t) + \cos^2(t)\right), \right. \\
	&\quad \left. u\sin(t)(2\cos(t) + 1) \right\rangle \\
	\mathbf{r}_u &= \left\langle 2(1 + \cos(t))\cos(t), 2(1 + \cos(t))\sin(t), -(1 + \cos(t))\cos(t) \right\rangle \\
	\mathbf{r}_t \times \mathbf{r}_u &= \left\langle -2u(1 + \cos(t))^2, 0, -4u(1 + \cos(t))^2 \right\rangle \\
	&= -2u(1 + \cos(t))^2 \langle 1, 0, 2 \rangle
\end{align*}
Estos vectores, a su vez, nos permiten obtener el elemento diferencial de área:
$\left\|  {\bf r}_{t}\times {\bf r}_{u}\right\|= 2\sqrt{5}\left|u\left( 1+\cos(t)\right)^{2}\right| .$ Por lo tanto, el área de la superficie en mención está dada por
\[\int_{0}^{1}\int_{0}^{2\pi }2\sqrt{5}\left| u\left( 1+\cos (t) \right) ^{2}\right|
dtdu= 3\pi \sqrt{5}.\]

En \verb"Mathematica" este valor se consigue con las siguientes instrucciones.



\begin{tcolorbox}[breakable,enhanced,title=\bf Instrucción en \verb"Mathematica", top=2pt,bottom=2pt,boxsep=0pt,before skip=2pt,after skip=0pt]
\begin{Verbatim}[formatcom=\letraverbatim]
x=2*u*(1+Cos[t])*Cos[t;
y=2*u*(1+Cos[t])*Sin[t];
z=2-x/2;
r=(1-u)*{x,y,0}+u*{x,y,z};
A=Simplify[Norm[Cross[D[r,t],D[r,u]]]];
NIntegrate[A,{u,0,1},{t,0,2Pi}]
\end{Verbatim}
\end{tcolorbox}


%\clearpage

\subsection*{Dos cilindros cardiotónicos entre dos planos}
Las ecuaciones $r = 1-\cos(t)$ y $r = 2- 2\cos(t)$ representan, en el plano, dos cardioides, y en el espacio estas mismas ecuaciones representan dos cilindros cardiotónicos \footnote{Este término fue acuñado por el profesor Juan Zambrano de la Universidad Distrital Francisco José de Caldas.}. Estas cilíndricas se cortan por los planos $z=0$ y $z=2-x/2$.
%
\begin{figure}[H]
	\centering
	\caption{Superficie sobre el plano \( z = 2 - x/2 \) acotada por dos cilindros cardióidicos}
	\begin{minipage}[t]{0.28\linewidth}
		\centering
		\includegraphics[width=\linewidth]{Fig/136.pdf}
	\end{minipage}
	\hspace{0.5em}
	\begin{minipage}[t]{0.28\linewidth}
		\centering
		\includegraphics[width=\linewidth]{Img/36b.pdf}
	\end{minipage}
	\vspace{1ex}
	\fuentepropia
\end{figure}
%
En el plano $xy$ los cilindros cardiotónicos se  parametrizan en la forma:
\begin{eqnarray*}
\overrightarrow{\alpha} (t) &=& \la(1-\cos(t))\cos(t), (1-\cos(t))\sin(t), 0\ra \\
\overrightarrow{\beta} (t)& =& \la(2-2\cos(t))\cos(t), (2-2\cos(t))\sin(t), 0\ra  ,\,\, t \in[0, 2\pi]
\end{eqnarray*}
Ahora formamos una combinación convexa entre estas, que determina la región contenida entre las dos curvas. Esto es, para $u\in[0, 1]$ y $t \in[0, 2\pi]$,
cada una de las curvas $\overrightarrow{\alpha }(t)$ y $\overrightarrow{\beta }(t)$
se proyecta sobre el plano $z=2-x/2$, de la siguiente forma:
\[
\adjustbox{max width=\textwidth}{
$\begin{array}{rcl}
\overrightarrow{p}(t) &=& \big\langle \left( 1-\cos(t) \right)\cos(t),\left( 1-\cos(t) \right)\sin(t), 2-\left( \left( 1-\cos(t) \right)\cos(t) \right)/2 \big\rangle \\[1ex]
\overrightarrow{q}(t) &=& \big\langle \left( 2-2\cos(t) \right)\cos(t), \left( 2-2\cos(t) \right)\sin(t), 2-\left( 2\left( 1-\cos(t) \right)\cos(t) \right)/2 \big\rangle
\end{array}$}
\]

\begin{figure}[H]
	\centering
	\caption{Cilindros cardiotónicos acotados por dos planos}
	\begin{minipage}[t]{0.28\linewidth}
		\centering
		
		\includegraphics[width=\linewidth]{Fig/137.pdf}
	\end{minipage}
	\hspace{1em}
	\begin{minipage}[t]{0.28\linewidth}
		\centering
		
		\includegraphics[width=\linewidth]{Fig/138.pdf}
	\end{minipage}
	\vspace{1ex}
	\fuentepropia
\end{figure}

Con el propósito de parametrizar la parte sobre los cilindros entre los planos $z=0$ y $z=2-x/2$, se hará una combinación convexa entre las expresiones $\overrightarrow{p}(t)$ y $\overrightarrow{\alpha}(t)$ por un lado
y $\overrightarrow{q}(t)$ con $\overrightarrow{\beta}(t)$ por el otro.




\begin{figure}[H]
\centering
\caption{
	Cilindros cardiotónicos
}
\includegraphics[width=0.4\textwidth]{Img/57a.pdf} \hspace{5mm}
\includegraphics[width=0.4\textwidth]{Img/57b.pdf}
\fuentepropia
\end{figure}

Estas figuras se obtienen en \verb"Mathematica" con las siguientes instrucciones

 % espacio entre texto y código



\begin{tcolorbox}[breakable,enhanced,title=\bf Instrucción en \verb"Mathematica", top=2pt,bottom=2pt,boxsep=0pt,before skip=2pt,after skip=0pt]
\begin{Verbatim}[formatcom=\letraverbatim]
ParametricPlot3D[{{2*Cos[t]*(1-Cos[t]),
2*Sin[t]*(1-Cos[t]),u(2-Cos[t]+Cos[t]^2)},
{Cos[t]*(1-Cos[t]),Sin[t]*(1-Cos[t]),
(u/2)*(4-Cos[t]+Cos[t]^2)}},{u,0,1},{t,0,2*Pi},
PlotStyle->{Orange,RGBColor[0.3,0.8,0.4]},
Ticks->{{-2,0,1},{-2,0,1},{-1,0,1}},
Mesh->False,PlotPoints->80]
\end{Verbatim}
\end{tcolorbox}




Ahora calcularemos el área sobre los cilindros.

\subsubsection*{Cilindro menor}
Para $\overrightarrow{p}(t)$ y $\overrightarrow{\alpha }(t)$ se puede simplificar la expresión
$(1-u)\overrightarrow{\alpha }(t)+u\overrightarrow{p}(t),$
la cual proporciona la función que describe la parte del cilindro menor y que se halla entre los planos $z=0$ y $z=2-x/2$.


\begin{multline*}
{\bf r}\left( u,t\right) = \left( 1-u\right) \left\langle \left( 1-\cos (t) \right) \cos
(t) ,\left( 1-\cos (t) \right) \sin (t),0\right\rangle \\
+u\left\langle \left( 1-\cos (t) \right) \cos (t) ,\left( 1-\cos (t) \right) \sin (t).
2-\left( \left( 1-\cos (t) \right) \cos (t) \right)
/2\right\rangle \\
=\left\langle \cos (t) \left( 1-\cos (t) \right)
\sin (t) \left( 1-\cos (t) \right) ,\frac{1}{2}u\left( 4-\cos (t) +\cos ^{2}(t) \right)
\right\rangle,
\end{multline*}
en donde $t\in \lbrack 0,2\pi ],$ $u\in \lbrack 0,1].$
Con el fin de hallar el área superficial, se requieren los siguientes elementos.


\begin{multline*}
{\bf r}_{u}= \la 0,0,2-\frac{1}{2}\cos (t) +\frac{1}{2}\cos ^{2}(t) \ra \\
{\bf r}_{t}= \big\langle \sin (t) \left( 2\cos (t) -1\right) ,\left( 2\cos (t) +1\right) \left( 1-\cos (t) \right) , \hskip1.5cm \\ \frac{1}{2}u\sin (t) \left( 1-2\cos (t) \right) \big\rangle.
\end{multline*}

El elemento diferencial de área es
\begin{equation}
\left\| \mathbf{r}_{t}\times \mathbf{r}_{u}\right\| = \frac{1}{2}\left( \cos ^{2}(t) -\cos (t) +4\right) \sqrt{2-2\cos (t)}.
\end{equation}
Entonces, el área superficial sobre el cilindro menor es
\begin{equation}
\int_{0}^{1}\int_{0}^{2\pi }\frac{1}{2}\left(\cos ^{2}(t)
-\cos (t) +4\right) \sqrt{2-2\cos (t)}dtdu= \frac{96}{5}.
\end{equation}
Estos cálculos se hacen muy simples, ejecutando las siguientes instrucciones:



\begin{tcolorbox}[breakable,enhanced,title=\bf Instrucción en \verb"Mathematica", top=2pt,bottom=2pt,boxsep=0pt,before skip=2pt,after skip=0pt]
\begin{Verbatim}[formatcom=\letraverbatim]
{x,y}={Cos[t]*(1-Cos[t]),Sin[t]*(1-Cos[t])};
a={x,y,0};
p={x,y,2-x/2};
r=(1-u)*a+u+p;
Integrate[Norm[Cross[D[r,u]]],{t,0,2*Pi},{u,0,1}]
\end{Verbatim}
\end{tcolorbox}



\subsubsection*{Cilindro mayor}
Para $\overrightarrow{q}(t)$ y $\overrightarrow{\beta}(t)$, es necesario simplificar la expresión 
\[
(1 - u)\overrightarrow{\beta}(t) + u\overrightarrow{q}(t),
\]
lo que nos provee de la función que parametriza la región entre los planos sobre el cilindro mayor, esto es

\begin{flushleft}
\adjustbox{max width=\textwidth}{
$\begin{array}{l}
\hspace{-0.6em}\mathbf{r}(u,t) = (1 - u)(2(1 - \cos(t))\cos(t),\; 2(1 - \cos(t))\sin(t),\; 0) \\[1ex]
\quad + u((2 - 2\cos(t))\cos(t),\; (2 - 2\cos(t))\sin(t),\; 2 - (2(1 - \cos(t))\cos(t))/2) \\[1ex]
= \langle 2\cos(t)(1 - \cos(t)),\; 2\sin(t)(1 - \cos(t)),\; u(2 - \cos(t) + \cos^2(t)) \rangle
\end{array}$}
\end{flushleft}

Con lo anterior, se halla el elemento diferencial de área:
\[
\left\| {\bf r}_{t} \times {\bf r}_{u} \right\| = \left| 2\left(2 - \cos(t) + \cos^2(t)\right) \sqrt{2 - 2\cos(t)} \right|.
\]

Luego, el área superficial del cilindro mayor está dada por:
\[
\int_{0}^{1} \int_{0}^{2\pi} \left| 2\left(2 - \cos(t) + \cos^2(t)\right) \sqrt{2 - 2\cos(t)} \right| \, dt \, du = \frac{224}{5}.
\]

En \verb"Mathematica", estos cálculos se hacen con las siguientes instrucciones.



\begin{tcolorbox}[breakable,enhanced,title=\bf Instrucción en \verb"Mathematica", top=2pt,bottom=2pt,boxsep=0pt,before skip=2pt,after skip=0pt]
\begin{Verbatim}[formatcom=\letraverbatim]
{x,y}={2*Cos[t]*(1-Cos[t]),2*Sin[t]*(1-Cos[t])};
r=(1-u)*{x,y,0}+u*{x,y,2-x/2};
Integrate[Norm[Cross[D[r,t,[r,u]]],{t,0,2*Pi},{u,0,1}]
\end{Verbatim}
\end{tcolorbox}



\subsection*{Cono oblicuo truncado}
Consideremos las funciones vectoriales $\overrightarrow{\alpha }(t) =\langle 1+\cos \left(
t\right) ,\sin (t) ,0\rangle$ y $\overrightarrow{\beta }(t) =\langle 2+2\cos \left(
t\right) ,2\sin (t) ,3 \rangle,$
que representan dos circunferencias en el espacio. Una combinación convexa de estas define la función r:

\begin{flalign*}
	&\mathbf{r}(t, u) = (1-u)\overrightarrow{\alpha}(t) + u\overrightarrow{\beta}(t) = \small{\left\langle (1+\cos(t))(1+u), \sin(t)(1+u), 3u\right\rangle}, &&
\end{flalign*}

con $t\in \lbrack 0,2\pi ],$ $u\in \lbrack 0,1].$ Esta función parametriza el cono truncado oblicuo que se muestra en la figura.



\begin{figure}[H]
\centering
\caption{Cono oblicuo truncado}
\includegraphics[width=0.5\textwidth]{Img/19a.pdf}
\fuentepropia
\end{figure}

Con la función ${\bf r}\left( t,u\right) $ hallamos${\bf r}_{t}=\left\langle -\sin (t) \left( 1+u\right) ,\cos \left(
t\right) \left( 1+u\right) ,0\right\rangle \,$ \break
y $\,{\bf r}_{u}=\left\langle 1+\cos (t) ,\sin (t),3\right\rangle,$
que determinan el elemento diferencial de área
{\small
\begin{align*}
	\|\mathbf{r}_t \times \mathbf{r}_u\| &= \left\| \left\langle 3\cos(t)(1+u), 3\sin(t)(1+u), -(1+\cos(t))(1+u) \right\rangle \right\| \\
	&= (1+u)\sqrt{\cos^2(t) + 10 + 2\cos(t)},
\end{align*}
}
cuya integral proporciona el área de esta región
\[\int_{0}^{2\pi }\int_{0}^{1}\left( 1+u\right) \sqrt{\cos ^{2}(t) +10+2\cos (t) }dudt \approx 30.4679.\]
La figura y los cálculos se obtienen con las siguientes instrucciones:



\begin{tcolorbox}[breakable,enhanced,title=\bf Instrucción en \verb"Mathematica", top=2pt,bottom=2pt,boxsep=0pt,before skip=2pt,after skip=0pt]
\begin{Verbatim}[formatcom=\letraverbatim]
ParametricPlot3D[{(1+u)*(1+Cos[t]),(1+u)*Sin[t],3*u},
{u,0,1},{t,0,2*Pi},PlotPoints->80,Mesh->True,
PlotStyle->Blend[{Red,Black,Yellow}]]
{a,b}={{1+Cos[t],Sin[t],0},{2+2*Cos[t],2*Sin[t],3}}
A=Norm[Cross[D[(1-u)*a+u*b,t],D[(1-u)*a+u*b,u]]]
NIntegrate[A,{t,0,2*Pi},{u,0,1}]
\end{Verbatim}
\end{tcolorbox}



\subsection*{Superficie obtenida entre dos elipses}

Las funciones vectoriales
$\overrightarrow{a}(t)= \la 2\cos(t),\sin(t),0\ra$ y $\overrightarrow{b}(t)= \big \langle \cos(t), \linebreak 3\sin(t),3\big\rangle$,
con $t \in [0,2\pi]$ representan dos elipses en el espacio.
Con la función
\begin{eqnarray*}
{\bf r}\left( t,u\right) &=&\left( 1-u\right) \left\langle 2\cos (t),\sin (t) ,0\right\rangle +u\left\langle \cos (t),3\sin (t) ,3 \right\rangle \\
&=&  \left\langle \cos (t)(2-u),\sin (t)(1+2u),3u\right\rangle,
\end{eqnarray*}
Con $u \in [0,1]$ y $t \in [0,2\pi]$, se parametriza la superficie que se muestra en la siguiente figura.

%


\begin{figure}[H]
\centering
\caption{Superficie que conecta dos elipses.}
\includegraphics[width=0.4\textwidth]{Fig/139.pdf} \hspace{5mm}
\includegraphics[width=0.4\textwidth]{Img/20a.pdf}

\includegraphics[width=0.4\textwidth]{Img/20b.pdf}
\fuentepropia
\end{figure}

Este gráfico se obtiene con las siguientes instrucciones:



\begin{tcolorbox}[breakable,enhanced,title=\bf Instrucción en \verb"Mathematica", top=2pt,bottom=2pt,boxsep=0pt,before skip=2pt,after skip=0pt]
\begin{Verbatim}[formatcom=\letraverbatim]
ParametricPlot3D[{(2-u)*Cos[t],(1+2*u)*Sin[t],3*u},
{u,0,1},{t,0,2*Pi},PlotPoints->80,Mesh->False]
\end{Verbatim}
\end{tcolorbox}



Con ayuda de la función $\mathbf{r}$ se obtienen los siguientes vectores:\\
$\mathbf{r}_{u} =\left\langle -\cos(t),\, 2\sin(t),\, 3 \right\rangle$ y $\mathbf{r}_{t} = \left\langle -\sin(t)(2 - u),\, \cos(t)(1 + 2u),\, 0 \right\rangle$, los cuales describen las direcciones tangentes de la superficie parametrizada. Además, el producto cruzado\\ $\mathbf{r}_{t} \times \mathbf{r}_{u} = \left\langle 3\cos(t)(1 + 2u),\, 3\sin(t)(2 - u),\, 5\cos^2(t) + 2u - 4 \right\rangle$ nos proporciona el vector normal y define el elemento diferencial de área.
\[
\left\| {\\bf r}_{t} \times {\\bf r}_{u} \right\| = \sqrt{ \cos^2(t)\left(25\cos^2(t) - 67 + 92u + 27u^2\right) + 13(u - 2)^2 }.
\]

La integral que proporciona el área de la superficie en mención es:
\[
\int_{0}^{1} \int_{0}^{2\pi} \sqrt{ \cos^2(t)\left(25\cos^2(t) + P(u)\right) + 13(u - 2)^2 } \, dt \, du \approx 35.72,
\]

Estos cálculos se hacen en \verb"Mathematica" usando las siguientes sentencias.



\begin{tcolorbox}[breakable,enhanced,title=\bf Instrucción en \verb"Mathematica", top=2pt,bottom=2pt,boxsep=0pt,before skip=2pt,after skip=0pt]
\begin{Verbatim}[formatcom=\letraverbatim]
r=Simplify[(1-u)*{2*Cos[t],Sin[t],0}+u*{Cos[t],
3*Sin[t],3}];
A=Simplify[Norm[Cross[D[r,t],D[r,u]]]];
NIntegrate[A,{t,0,2Pi},{u,0,1}]
\end{Verbatim}
\end{tcolorbox}




\index{Cono cardioidico!oblicuo y truncado}
\subsection*{Cono cardiogénico oblicuo y truncado} 
A partir de las cardioides definidas por las funciones vectoriales
\begin{eqnarray*}
\overrightarrow{\alpha }(t) & \left\langle \left( 1+\cos \left(
t\right) \right) \cos (t) ,\left( 1+\cos (t) \right)
\sin (t) ,0\right\rangle \\
\overrightarrow{b}(t) & \left\langle 2\left( 1+\cos \left(
t\right) \right) \cos (t) ,2\left( 1+\cos (t)\right) \sin (t) ,3\right\rangle.
\end{eqnarray*}
Con $t \in [0,2\pi],$ se trazan segmentos rectilíneos entre estas dos curvas, lo que produce una parametrización de un cono truncado oblicuo, cuyos cortes paralelos al plano $xy$ son cardioides.
La función que parametriza la superficie obtenida y su gráfico son las siguientes:
\begin{eqnarray*}
{\bf r}\left( t,u\right) &=& (1-u)\overrightarrow{\alpha }(t)+ u \overrightarrow{\beta}(t) \\
&& \left\langle (1+\cos(t))(1+u) \cos (t),( 1+\cos (t))(1+u) \sin (t) ,3u  \right\rangle.
\end{eqnarray*}




\begin{figure}[H]
\centering
\caption{Cono cardiogénico oblicuo truncado}
\includegraphics[width=0.4\textwidth]{Fig/140.pdf} \hspace{5mm}
\includegraphics[width=0.4\textwidth]{Img/21a.pdf}

\includegraphics[width=0.4\textwidth]{Img/21b.pdf}
\fuentepropia
\end{figure}

La superficie se puede obtener ejecutando las siguientes instrucciones.



\begin{tcolorbox}[breakable,enhanced,title=\bf Instrucción en \verb"Mathematica", top=2pt,bottom=2pt,boxsep=0pt,before skip=2pt,after skip=0pt]
\begin{Verbatim}[formatcom=\letraverbatim]
ParametricPlot3D[{(2-u)*Cos[t]*(1+Cos[t]),
(2-u)*Sin[t]*(1+Cos[t]),3-3*u},{u,0,1},{t,0,2*Pi},
PlotPoints->90,Mesh->False]
\end{Verbatim}
\end{tcolorbox}



Con la función ${\bf r}$, se obtienen los siguientes vectores\\
${\bf r}_{t}= \left\langle -\sin (t) \left( 1+2\cos (t) \right) (1+u) ,\left( 1+\cos(t) \right)
\left( 2\cos(t) -1\right)( 1+u) ,0\right\rangle $\\
$ {\bf r}_{u}=\left\langle \cos (t) \left( 1+\cos \left(
t\right) \right) ,\sin (t) \left( 1+\cos (t) \right),3\right\rangle .$
Además,
\begin{multline*}
{\bf r}_{t}\times {\bf r}_{u} = (1+u) \big\langle 3\left( 1+\cos (t) \right) \left( 2\cos (t) -1\right),3\sin (t) \left( 1+2\cos (t) \right)  , \\
-\left( 1+\cos (t) \right) ^{2} \big \rangle
\end{multline*}

La norma de este último vector es el elemento diferencial de área
\[\left\| {\bf r}_{t}\times {\bf r}_{u}\right\| = (1+u) \sqrt{(1+\cos(t))(\cos ^{3}(t) +3\cos^{2}(t)+3\cos (t) +19)},\]
que, al integrarse, nos proporciona el área de la superficie que se considera
{\small
\[\int_{0}^{2\pi }\int_{0}^{1}\left\| {\bf r}_{t}\times {\bf r}_{u}\right\|  dudt \approx 39.4.\]
}

Este valor se obtiene ejecutando las siguientes instrucciones:



\begin{tcolorbox}[breakable,enhanced,title=\bf Instrucción en \verb"Mathematica", top=2pt,bottom=2pt,boxsep=0pt,before skip=2pt,after skip=0pt]
\begin{Verbatim}[formatcom=\letraverbatim]
r=Simplify[(1-u)*{(1+Cos[t])*Cos[t],
(1+Cos[t])*Sin[t],0}+u*{2*(1+Cos[t])*Cos[t],
2*(1+Cos[t])*Sin[t],3}];
A=Simplify[Norm[Cross[D[r,t],D[r,u]]]]
NIntegrate[A,{t,0,2Pi},{u,0,1}]
\end{Verbatim}
\end{tcolorbox}



\subsection*{Una flor sobre una esfera}
La ecuación $r=\sin(3\theta) $ representa, en coordenadas polares, una rosa de 3 pétalos; la proyectamos sobre la esfera centrada en el origen de radio $\sqrt{2}$.



\begin{figure}[H]
\centering
\caption{La flor $r=3\sin(3\theta)$ sobre una esfera}
\includegraphics[width=0.4\textwidth]{Fig/141.pdf} \hspace{5mm}
\includegraphics[width=0.4\textwidth]{Img/48a.pdf}
\fuentepropia
\end{figure}

En el plano, la flor (con su interior) se parametriza en  forma \\ {\small $x=u\sin (3t) \cos (t) ,$} {\small $y=u\sin (3t) \sin (t)$}
al reemplazar en la ecuación de la esfera,
{\small $x^{2}+y^{2}+z^{2}=2$,} se sigue que
{\small $z=\sqrt{2-u^{2}\sin^{2}(3t)},$}
entonces, la función
{\small ${\bf r}\left( t,u\right) =\left\langle u\sin (3t) \cos (t) ,u\sin (3t) \sin (t) ,\sqrt{2-u^{2}\sin^{2}(3t) }\right\rangle$}
con {\small $u \in [0,1]$, $t \in [0,2\pi]$,}
parametriza la región sobre la esfera cuya proyección en el plano $xy$ es la flor mencionada.
Los gráficos se obtienen ejecutando las siguientes instrucciones:



\begin{tcolorbox}[breakable,enhanced,title=\bf Instrucción en \verb"Mathematica", top=2pt,bottom=2pt,boxsep=0pt,before skip=2pt,after skip=0pt]
\begin{Verbatim}[formatcom=\letraverbatim]
A=ContourPlot3D[x^2+y^2+z^2==2,{x,-1.5,1.5},
{y,-1.5,1.5},{z,-1.5,1.5},ContourStyle->{Orange,
Opacity[0.6]},Mesh->False];
r={u*Sin[3*t]*Cos[t],u*Sin[3*t]*Sin[t]};
p={r[[1]],r[[2]],Sqrt[2-r[[1]]^2-r[[2]]^2]};
b=ParametricPlot3D[p,{u,0,1},{t,0,2*Pi},
PlotRange->{{-1,1},{-1,1},{1,1.5}},Mesh->False,
PlotStyle->RGBColor[0.8,0.3,0.3],PlotPoints->90,
Ticks->{{-1,0,1},{-1,0,1},{1.2,1.4}}];
c=ParametricPlot3D[{Sin[3*t]*Cos[t],Sin[3*t]*Sin[t],
Sqrt[2-Sin[3*t]^2]},{t,0,2*Pi},PlotPoints->80,
PlotStyle->{Black,Thickness->0.012}];
Show[A,c,b]
\end{Verbatim}
\end{tcolorbox}





\begin{figure}[H]
\centering
\caption{La flor $r=3\sin(3\theta)$ sobre una esfera}
\includegraphics[width=0.4\textwidth]{Img/68a.pdf} \hspace{5mm}
\includegraphics[width=0.4\textwidth]{Img/68b.pdf}
\fuentepropia
\end{figure}

Con la función ${\bf r}$ definida antes, hallamos los siguientes vectores:
{\small
\begin{eqnarray*}
{\bf r}_{t} &=& u\bigg\langle \cos(2t)+2\cos(4t),2\sin(4t)-\sin(2t),\frac{-3u\sin(6t)}{2\sqrt{2-u^{2}\sin^{2}(3t)  }} \bigg\rangle \\
{\bf r}_{u} &=& \sin (3t)\left\langle  \cos (t),
\sin (t) ,-\frac{u\sin (3t) }{\sqrt{2-u^{2}\sin^{2}(3t) }}\right\rangle
\end{eqnarray*}
}

{\small
\begin{flalign*}
	&\mathbf{r}_t \times \mathbf{r}_u = \left\langle -\frac{u^2 \cos(t) \sin^3(3t)}{\sqrt{2 - u^2 \sin^2(3t)}}, -\frac{u^2 \sin(t) \sin^3(3t)}{\sqrt{2 - u^2 \sin^2(3t)}}, -u \sin^2(3t) \right\rangle. &
\end{flalign*}
}

Entonces,
{\small
\[\left\| {\bf r}_{t}\times {\bf r}_{u}\right\|=\frac{\sqrt{2}u\sin^{2}(3t)}{\sqrt{2-u^{2}\sin^{2}(3t)}},\]
}
cuya integral nos proporciona claridad sobre la región considerada.
{\small
\[\bigintss_{0}^{2\pi}\bigintss_{0}^{1}\frac{\sqrt{2}u\sin ^{2}(3t) }{\sqrt{2-u^{2}\sin^{2}(3t) } }dudt \approx 1. 761.\]
}
La integral anterior se evalúa numéricamente con las siguientes sentencias.




\begin{tcolorbox}[breakable,enhanced,title=\bf Instrucción en \verb"Mathematica", top=2pt,bottom=2pt,boxsep=0pt,before skip=2pt,after skip=0pt]
\begin{Verbatim}[formatcom=\letraverbatim]
r={u*Sin[3*t]*Cos[t],u*Sin[3*t]*Sin[t]};
p={r[[1]],r[[2]],Sqrt[2-r[[1]]^2-r[[2]]^2]};
A=Norm[Cross[D[p,t],D[p,u]]];
NIntegrate[A,{t,0,2*Pi},{u,0,1}]
\end{Verbatim}
\end{tcolorbox}



\subsection*{Toros cardioidicos} \index{Toro cardioidico}
En esta sección consideraremos la superficie que se obtiene al hacer girar una cardioide alrededor de una recta en el mismo plano que no interseca a la figura.




\begin{figure}[H]
\centering
\caption{Cardioides}
\includegraphics[width=0.4\textwidth]{Fig/143.pdf} \hspace{5mm}
\includegraphics[width=0.4\textwidth]{Fig/144.pdf}

\includegraphics[width=0.4\textwidth]{Fig/145.pdf}
\fuentepropia
\end{figure}

\subsubsection*{Toro cardióidico generado por \boldmath$r=1-\sin(t)$ \unboldmath}
En el plano, la figura de la ecuación $r=1-\sin(t)$ representa una cardioide.
Esta figura tiene área determinada por
\[\int_{0}^{2\pi }\int_{0}^{1-\sin \theta }rdrd\theta = \frac{3}{2}\pi\]
y su longitud de arco se puede hallar como sigue
\begin{eqnarray*}
\int_{0}^{2\pi }\sqrt{r^{2}+\left( \frac{dr}{d\theta }\right) ^{2}}d\theta &=&
\int_{0}^{2\pi }\sqrt{\left( 1-\sin \theta \right) ^{2}+\left( -\cos \left(
\theta \right) \right) ^{2}}d\theta\\
&=&\int_{0}^{2\pi }\sqrt{2-2\sin \left(\theta \right) }d\theta = 8.
\end{eqnarray*}

La figura de la ecuación $r=1-\sin(t)$ se puede parametrizar en el plano en la forma
$\langle (1-\sin(t))\cos(t),(1-\sin(t))\sin(t) \rangle $ con $t\in [0,2\pi].$
Ahora, le sumaremos $3$ a la coordenada $x$, obteniendo la función $\overrightarrow{\alpha}(t)$,
y su gráfico se hará girar dicha figura alrededor de la recta $\theta = \pi/2.$


\begin{figure}[H]
\centering
\caption{Cardioide $\boldsymbol\alpha(t)=\langle 3+(1-\sin(t))\cos(t),(1-\sin(t))\sin(t)\rangle$}
\includegraphics[width=0.8\textwidth]{Fig/146.pdf}
\fuentepropia
\end{figure}

El área de una superficie de revolución está dada por
\begin{equation}\label{toros1}
A=\int_{a}^{b}x(t)\sqrt{\left( \frac{dx}{dt}\right) ^{2}+\left( \frac{dy}{dt}\right) ^{2}}dt
\end{equation}
Si el giro se hace alrededor del eje $y$, en este caso se obtiene

\[2\pi \int_{0}^{2\pi }\left( 3+\left( 1-\sin (t) \right) \cos
(t) \right) \sqrt{2-2\sin (t) }dt= 48\pi.\]

Para hallar el centro de masa de la curva, determinada por la expresión
$r=1-\sin (t) ,$ nótese que el término \[\sqrt{\left( \frac{dx}{dt}\right) ^{2}+\left( \frac{dy}{dt}\right) ^{2}}=\sqrt{2-2\sin (t) }.\]
De esta manera, las coordenadas del centro de masa de la curva están dadas por
\begin{eqnarray*}
\overline{x}&=&\frac{1}{8}\int_{0}^{2\pi }\left( 3+\left( 1-\sin \left(
t\right) \right) \cos (t) \right) \sqrt{2-2\sin (t) }dt= 3 \\
\overline{y} &=& \frac{1}{8}\int_{0}^{2\pi }\left( 1-\sin (t)\right) \sin (t) \sqrt{2-2\sin (t) }dt= -\frac{4}{5}.
\end{eqnarray*}

Usando el teorema de Pappus (\ref{teoremadepappus}), ver página (\pageref{teoremadepappus}), el área de la superficie de revolución de la curva $r=1-\sin(t)$ alrededor de la recta  $\theta = \pi /2$ es.
$2\pi \left( 3\right) \left( 8\right) = 48\pi .$ \index{Teorema de Pappus!área superficial}

Al ubicar la curva de giro en el espacio tridimensional y hacer coincidir el eje de giro con el eje $z$,
se satisface que si un punto sobre esta curva gira alrededor del eje $z$, entonces se obtiene una circunferencia y se genera una superficie que denominaremos \emph{toro cardioidico.}
\index{Toro cardioidico}




\begin{figure}[H]
\centering
\caption{Toro cardiogénico construido a partir de $r=1-\sin(t)$}
\includegraphics[width=0.4\textwidth]{Fig/147.pdf} \hspace{5mm}
\includegraphics[width=0.4\textwidth]{Img/63a.pdf}
\fuentepropia
\end{figure}

Al girar estos puntos sobre el eje $z$, todos los puntos conservan la misma coordenada en $z$,
y la coordenada en $x$ e $y$ se le agregará una $\cos(u)$ y $\sin(u)$ respectivamente, siendo $u$ un parámetro entre $0$ y $2 \pi$. Por lo tanto, esta superficie se parametriza como
\begin{flushleft} % Fuerza todo el bloque a la izquierda.
	\resizebox{1.0\textwidth}{!}{%
		$\mathbf{r}(t,u) = \left\langle (3+(1-\sin(t))\cos(t))\cos(u),(1-\sin(t))\sin(t)\sin(u),(1-\sin(t))\sin(t) \right\rangle$
	}
\end{flushleft}
con $t \in [0,2\pi],$ $u \in [0,2\pi].$




\begin{figure}[H]
\centering
\caption{Toro cardiogénico construido a partir de $r=1-\sin(t)$}
\includegraphics[width=0.4\textwidth]{Img/63b.pdf} \hspace{5mm}
\includegraphics[width=0.4\textwidth]{Img/63c.pdf}
\fuentepropia
\end{figure}

Las figuras anteriores se obtienen en \verb"Mathematica" ejecutando las siguientes instrucciones:



\begin{tcolorbox}[breakable,enhanced,title=\bf Instrucción en \verb"Mathematica", top=2pt,bottom=2pt,boxsep=0pt,before skip=2pt,after skip=0pt]
\begin{Verbatim}[formatcom=\letraverbatim]
{x,y}={3+(1-Sin[t])*Cos[t],(1-Sin[t])*Sin[t]};
ParametricPlot[{x,y},{t,0,2*Pi}]
ParametricPlot3D[{x*Cos[u],x*Sin[u],y},{t,0,2*Pi},
{u,0,5*Pi/3},Ticks->{{-3,-1,1,3},{-2,0,2},{-1,0,1}},
MeshStyle->Gray]
\end{Verbatim}
\end{tcolorbox}



Utilizando la función ${\bf r}(t,u)$  se hallan los siguientes vectores:

{\small
${\bf r}_{t} = \big \langle\cos(u)(\cos(2t)-\sin(t)),\sin(u)(\cos(2t)-\sin(t)),\cos(t)(1-2\sin(t)) \big\rangle.$}

${\bf r}_{u} = \big \langle\sin(u)(\cos(t)\sin(t)-3-\cos(t)),\cos(u)\cos(t)(1-\sin(t)),0\big\rangle$

y

${\bf r}_{u} = \big \langle\sin(u)(\cos(t)\sin(t)-3-\cos(t)),\cos(u)\cos(t)(1-\sin(t)),0\big\rangle.$


\begin{multline*}
{\bf r}_{t}\times {\bf r}_{u}=\bigg\langle\cos^{2}(t)(1-2\sin(t))(\sin(t)-1)\cos(u),\\
\cos(t)(1-2\sin(t))\sin (u)(\cos(t)\sin(t)-3-\cos(t)),\\
(\cos(2t)+\sin(t))(\cos(t)(\sin(t)-1)-3\sin^{2}(u))\bigg\rangle,
\end{multline*}
y con estos se obtiene el elemento diferencial de área
\begin{multline*}
\left\|{\bf r}_{t}\times {\bf r}_{u}\right\|=\frac{1}{2\sqrt{2}}\bigg|-14\cos\left(\frac{t}{2}\right)-3\cos\left(\frac{3t}{2}\right)+\cos\left(\frac{5t}{2}\right)+10\sin\left(\frac{t}{2}\right)\\
+3\sin\left(\frac{3t}{2}\right)+\sin\left(\frac{5t}{2}\right)\bigg|,
\end{multline*}
que se integra y nos proporciona el área superficial de la región en mención.
\begin{multline*}
\bigintss_{0}^{2\pi}\bigintss_{0}^{2\pi}\frac{1}{2\sqrt{2}}\bigg|-14\cos\left(\frac{t}{2}\right)-3\cos\left(\frac{3t}{2}\right)+\cos\left(\frac{5t}{2}\right)+10\sin\left(\frac{t}{2}\right)\\
+3\sin\left(\frac{3t}{2}\right)+\sin\left(\frac{5t}{2}\right)\bigg|dudt=48\pi.
\end{multline*}


Este procedimiento (y la figura de la curva revolucionada) puede evaluarse en \verb"Mathematica"
mediante la ejecución de las siguientes instrucciones:



\begin{tcolorbox}[breakable,enhanced,title=\bf Instrucción en \verb"Mathematica", top=2pt,bottom=2pt,boxsep=0pt,before skip=2pt,after skip=0pt]
\begin{Verbatim}[formatcom=\letraverbatim]
{x,y}={3+(1-Sin[t])*Cos[t],(1-Sin[t])*Sin[t]};
RevolutionPlot3D[{x,0,y},{t,0,2*Pi},MeshStyle->Gray,
RevolutionAxis->{0,0,1},PlotStyle->Orange]
r={x*Cos[u],x*Sin[u],y};
R=Norm[Cross[D[r,t],D[r,u]]];
A=Assuming[t>=0&&u>=0,Simplify[R]]
Integrate[A,{t,0,2*Pi},{u,0,2*Pi}]
\end{Verbatim}
\end{tcolorbox}




\subsubsection*{Toro cardioidico generado por \boldmath$r=1-\cos(t)$ \unboldmath}

El área de la figura de la ecuación $r=1-\cos(t)$ está dada por


{\small
\[\int_{0}^{2\pi }\int_{0}^{1-\cos \theta }rdrd\theta = \frac{3}{2}\pi\]
}
y su longitud de curva es
{\small
\[\int_{0}^{2\pi }\sqrt{\left( 1-\cos (t) \right) ^{2}+\left( \sin (t)\right)^{2}}dt=\int_{0}^{2\pi }\sqrt{2-2\cos (t) }dt= 8.\]
}


De acuerdo con la expresión (\ref{toros1}), el área superficial está determinada por
{\small
\[2\pi \int_{0}^{2\pi }\left( 3+\left( 1-\cos (t) \right) \cos
(t) \right) \sqrt{2-2\cos (t) }dt= \frac{176}{5}\pi.\]
}

El centro de masa de esta curva está dado por

{\small
\begin{eqnarray*}
\overline{x} &=& \frac{1}{8}\int_{0}^{2\pi }\left( 3+\left( 1-\cos (t) \right) \cos \left( t\right) \right) \sqrt{2-2\cos \left( t\right)}dt= \frac{11}{5}\\
\overline{y} &=& \frac{1}{8}\int_{0}^{2\pi }\left( 1-\cos \left( t\right)
\right) \sin \left( t\right) \sqrt{2-2\cos \left( t\right) }dt= 0.
\end{eqnarray*}
}


De acuerdo con el teorema de Pappus, el área de la superficie de revolución es
$2\pi \left( \frac{11}{5}\right) \left( 8\right) = \frac{176}{5}\pi .$
En este caso se usará la parametrización de la superficie:


{\small
\begin{multline*}
{\bf r}\left( t,u\right) =\big\langle( 3+\left( 1-\cos (t) \right)
\cos (t)) \cos (u) ,\left( 1-\cos (t)
\right) \cos (t) \sin \left( u\right) ,\\
\left( 1-\cos (t) \right) \sin (t) \big\rangle, \, t\in [0,2\pi],\, u\in [0,2\pi].
\end{multline*}
}

La figura obtenida es el siguiente:



\begin{figure}[H]
\centering
\caption{Toro cardiogénico construido a partir de $r=1-\cos(t)$}
\includegraphics[width=0.4\textwidth]{Img/59a.pdf} \hspace{5mm}
\includegraphics[width=0.4\textwidth]{Img/59b.pdf}

\includegraphics[width=0.4\textwidth]{Img/59c.pdf}
\fuentepropia
\end{figure}

Las instrucciones para conseguirlo se presentan a continuación.



\begin{tcolorbox}[breakable,enhanced,title=\bf Instrucción en \verb"Mathematica", top=2pt,bottom=2pt,boxsep=0pt,before skip=2pt,after skip=0pt]
\begin{Verbatim}[formatcom=\letraverbatim]
{x,y}={3+(1-Cos[t])*Cos[t],(1-Cos[t])*Sin[t]};
ParametricPlot[{x,y},{t,0,2*Pi}]
ParametricPlot3D[{x*Cos[u],x*Sin[u],y},{t,0,2*Pi},
{u,0,2*Pi}]
\end{Verbatim}
\end{tcolorbox}



Derivando parcialmente la función ${\bf r}$, hallamos los siguientes vectores:


\begin{flushleft}
\resizebox{1.0\textwidth}{!}{%
$\mathbf{r}_{t} = \left\langle (\sin(2t) - \sin t)\cos u,\; (\sin(2t) - \sin t)\sin u,\; (1 - \cos t)(2\cos t + 1) \right\rangle,$
}
\end{flushleft}
\[
{\mathbf r}_{u} = \left\langle (\cos t - 1)\cos t \sin u,\; (1 - \cos t)\cos t \cos u,\; 0 \right\rangle.
\]


Con ellos se halla el elemento diferencial de área:
\[
\left\| {\\bf r}_{t} \times {\\bf r}_{u} \right\| = \frac{1}{2} \left| -8\sin\left( \frac{t}{2} \right) - 3\sin\left( \frac{3t}{2} \right) + \sin\left( \frac{5t}{2} \right) \right|.
\]


La integración nos proporciona el área de la superficie en mención:
\[
\int_{0}^{2\pi} \int_{0}^{2\pi} \frac{1}{2} \left| -8\sin\left( \frac{t}{2} \right) - 3\sin\left( \frac{3t}{2} \right) + \sin\left( \frac{5t}{2} \right) \right| \, du \, dt = \frac{176}{5}\pi.
\]

\subsubsection*{Toro cardioidico generado por \boldmath$r=1+\cos(t)$ \unboldmath}
Al igual que en los casos anteriores, para la cardioide con ecuación
$r=1+\cos (t) ,$ el área está dada por
\[\int_{0}^{2\pi }\int_{0}^{1+\cos t}rdrdt= \frac{3}{2}\pi\]
y la longitud del arco es
\[\int_{0}^{2\pi }\sqrt{\left( 1+\cos t\right) ^{2}+\left( -\sin t\right) ^{2}}dt= 8\]
Si esta curva gira alrededor de la recta $\theta= \pi/2$, usando la expresión (\ref{toros1}),
el área superficial está dada por  
\[2\pi \int_{0}^{2\pi }\left( 3+\left( 1+\cos \left( t\right) \right) \cos
\left( t\right) \right) \sqrt{2+2\cos \left( t\right) }dt= \frac{304}{5}\pi.\]

El centro de masa de la curva está dado por
\begin{eqnarray*}
\overline{x} &=& \frac{1}{8}\int_{0}^{2\pi }\left( 3+\left( 1+\cos \left(
t\right) \right) \cos \left( t\right) \right) \sqrt{2+2\cos \left( t\right) }%
dt= \frac{19}{5} \\
\overline{y} &=& \frac{1}{8}\int_{0}^{2\pi }\left( 1+\cos \left( t\right)
\right) \sin \left( t\right) \sqrt{2+2\cos \left( t\right) }dt= 0
\end{eqnarray*}

El área de la superficie de revolución, de acuerdo con el teorema de Pappus, es  $2\pi \left( \frac{19}{5}\right) \left( 8\right) = \frac{304}{5}\pi. $


Para la superficie que se obtiene al hacer girar la figura de la curva $r=1+\cos(t)$ alrededor de la recta $\theta= \pi/2$, usaremos la parametrización.
\begin{multline*}
{\bf r}\left( t,u\right) =\bigg\langle(3+(1+\cos(t))\cos(t))\cos(u),(1+\cos(t))\cos(t)\sin(u),\\
(1+\cos(t))\sin(t)\bigg\rangle, \, t\in [0,2\pi],\, u\in [0,2\pi].
\end{multline*}

La figura obtenida es la siguiente:




\begin{figure}[H]
\centering
\caption{Toro cardiogénico construido a partir de $r=1+\cos(t)$}
\includegraphics[width=0.4\textwidth]{Img/64a.pdf} \hspace{5mm}
\includegraphics[width=0.4\textwidth]{Img/64b.pdf}

\includegraphics[width=0.4\textwidth]{Img/64c.pdf}
\fuentepropia
\end{figure}

Las siguientes instrucciones permiten obtener las figuras presentadas:



\begin{tcolorbox}[breakable,enhanced,title=\bf Instrucción en \verb"Mathematica", top=2pt,bottom=2pt,boxsep=0pt,before skip=2pt,after skip=0pt]
\begin{Verbatim}[formatcom=\letraverbatim]
x=3+{1+Cos[t]*Cos[t];
y=(1+Cos[t])*Sin[t];
ParametricPlot[{x,y},{t,0,2*Pi}]
ParametricPlot3D[{x*Cos[u],x*Sin[u],y},{t,0,2*Pi},
{u,0,4*Pi/3},MeshStyle->Gray]
\end{Verbatim}
\end{tcolorbox}



Con ayuda de la función ${\bf r}$, hallamos los siguientes vectores \linebreak
\[
\begin{split}
	\mathbf{r}_{t} &= \left\langle -\sin(t)\cos(u)(2\cos(t)+1), -\sin(t)\sin(u)(2\cos(t)+1), \right. \\
	&\quad \left. (1+\cos(t))(2\cos(t)-1) \right\rangle
\end{split}
\]
y 
\[{\bf r}_{u}= \big\langle-(3+(1+\cos (t))\cos (t))\sin (u),\linebreak (1+\cos(t))\cos(t)\cos(u),0 \big\rangle,\]
con ellos se obtiene el elemento diferencial de área
\[\left\| {\bf r}_{t}\times {\bf r}_{u}\right\| =\frac{1}{2}\left| 16\cos \left( \frac{t}{2}\right) +3\cos \left( \frac{3t}{2}\right) +\cos \left( \frac{5t}{2}\right) \right| \]
que, al integrarse, proporciona el área superficial de la región en mención.
\[\bigintss_{0}^{2\pi}\bigintss_{0}^{2\pi }\frac{1}{2}\left| 16\cos \left( \frac{t}{2}\right) +3\cos \left( \frac{3t}{2}\right) +\cos \left( \frac{5t}{2}\right) \right| dudt= \frac{304}{5}\pi.\]
Los cálculos anteriores se pueden evaluar en \verb"Mathematica" mediante la ejecución de las siguientes instrucciones:



\begin{tcolorbox}[breakable,enhanced,title=\bf Instrucción en \verb"Mathematica", top=2pt,bottom=2pt,boxsep=0pt,before skip=2pt,after skip=0pt]
\begin{Verbatim}[formatcom=\letraverbatim]
x=3+(1+Cos[t])*Cos[t];
y=(1+Cos[t])*Sin[t];
ParametricPlot[{x,0,y},{t,0,2*Pi}];RevolutionAxis->
{0,0,1}
MeshStyle->Gray,
Ticks->{{-3,0,3},{-3,0,3},
{-1,0}},
PlotStyle->Orange]
r={x*Cos[u],x*Sin[u],y};
A=Assuming[t>=0&&u>=0,
Simplify[Norm[Cross[D[r,t],D[r,u]]]]]
Integrate[A,{t,0,2*Pi},{u,0,2*Pi}]
\end{Verbatim}
\end{tcolorbox}




\begin{ejer}
Determine una expresión que defina la superficie indicada y calcule su área.
Además, determine una función que represente la curva de intersección de la superficie y halle su longitud.


\begin{enumerate}
\item La superficie sobre el paraboloide $2z=x^2+y^2$ acotada por el cilindro cardiogénico $r=1+\cos(t).$
%



\begin{figure}[H]
\centering
\caption{
	Paraboloide acotada por el cilindro cardiogénico
}
\includegraphics[width=0.4\textwidth]{Img/T_30a.pdf} \hspace{5mm}
\includegraphics[width=0.4\textwidth]{Img/T_30b.pdf}

\includegraphics[width=0.4\textwidth]{Img/T_30c.pdf}
\end{figure}

\begin{comment}
estos gráficos se obtienen ejecutando las siguientes instrucciones

\begin{Verbatim}
a=ParametricPlot3D[{(1+Cos[t])Cos[t],(1+Cos[t])*Sin[t],2*u},{t,0,2Pi},{u,0,2},
Mesh->False,PlotStyle->Orange,
PlotPoints->60];
c=ContourPlot3D[2*z==x^2+y^2,{x,-2,2},{y,-2,2},{z,0,2},
Mesh->False,ContourStyle->Cyan];
x=(1+Cos[t])Cos[t];
y=(1+Cos[t])*Sin[t];
z=(x^2+y^2)/2;
b=ParametricPlot3D[{x,y,z},{t,0,2Pi},
PlotStyle->{Black,Thickness->0.015}];
Show[c,a,b,PlotRange->{{-2,2.05},{-2,2.05},{-0.15,2}}]
\end{Verbatim}

\end{comment}

\item La superficie sobre el paraboloide hiperbólico $2z=y^2-x^2$ acotada por el cilindro cardiogénico $r=1+\cos(t).$



\begin{figure}[H]
\centering
\caption{Un paraboloide hiperbólico y un cilindro cardioidico}
\includegraphics[width=0.4\textwidth]{Img/T_31a.pdf} \hspace{5mm}
\includegraphics[width=0.4\textwidth]{Img/T_31b.pdf}
\end{figure}

\begin{comment}
estos gráficos se obtienen ejecutando las siguientes instrucciones

\begin{Verbatim}
a=ParametricPlot3D[{(1+Cos[t])Cos[t],(1+Cos[t])*Sin[t],
2*u},{t,0,2Pi},{u,-1,1},Mesh->False,
PlotStyle->Orange,PlotPoints->60];
c=ContourPlot3D[
2*z==y^2-x^2,{x,-3,3},{y,-3,3},{z,-2.5,2},
Mesh->False,ContourStyle->Blend[{Red,Green,Blue},3/4]];
x=(1+Cos[t])Cos[t];
y=(1+Cos[t])*Sin[t];
z=(y^2-x^2)/2;
b=ParametricPlot3D[{x,y,z},{t,0,2Pi},
PlotStyle->{Black,Thickness->0.015}];
Show[b,c,a,PlotRange->{{-2.5,2.5},{-2.5,2.5},{-2.1,2.6}}]
\end{Verbatim}

\end{comment}

\item La superficie sobre la esfera con ecuación $x^2+y^2+z^2=4$ acotada por el cilindro cardiogénico $1+\cos(t)$. 


\begin{figure}[H]
\centering
\caption{Esfera y cilindro cardiológico}
\includegraphics[width=0.4\textwidth]{Img/T_32a.pdf} \hspace{5mm}
\includegraphics[width=0.4\textwidth]{Img/T_32b.pdf}

\includegraphics[width=0.4\textwidth]{Img/T_32c.pdf}
\end{figure}

\begin{comment}
estos gráficos se obtienen ejecutando las siguientes instrucciones

\begin{Verbatim}
a=ParametricPlot3D[{(1+Cos[t])Cos[t],(1+Cos[t])*Sin[t],u},{t,0,2Pi},
{u,-2,2},Mesh->False,PlotStyle->Orange,PlotPoints->80];
c=ContourPlot3D[
x^2+y^2+z^2==4,{x,-2.5,2.5},{y,-2.5,2.5},{z,-2.5,2},
Mesh->False,
ContourStyle->{Blend[{Orange,Green},1/4],Opacity[0.6]},
PlotPoints->80];
x=(1+Cos[t])Cos[t];
y=(1+Cos[t])*Sin[t];
z=Sqrt[4-y^2-x^2];
b=ParametricPlot3D[{{x,y,z},{x,y,-z}},{t,0,2Pi},
PlotStyle->{Directive[Black,Thickness->0.01],
Directive[Black,Thickness->0.01]},PlotPoints->90];
Show[b,c,a,PlotRange->{{-2,2},{-2,2},{-2,2}}]
\end{Verbatim}

\end{comment}
%\clearpage

\item Un pétalo de la flor $r=\sin(3\theta)$ gira alrededor de una recta perpendicular al eje polar, que pasa por el polo.
Hallar el área de la superficie obtenida.


\begin{figure}[H]
\centering
\caption{Superficie obtenida de girar $r=\sin(3\theta)$.}
\includegraphics[width=0.4\textwidth]{Img/99a.pdf} \hspace{5mm}
\includegraphics[width=0.4\textwidth]{Img/99b.pdf}

\includegraphics[width=0.4\textwidth]{Img/99c.pdf}
\end{figure}

\begin{comment}

\begin{Verbatim}
x=Sin[3t]*Cos[t]*Cos[u];
y=Sin[3t]*Cos[t]*Sin[u];
z=Sin[3t]*Sin[t];
ParametricPlot3D[{x,z,y},{t,0,Pi/3},{u,0,13Pi/8},
Boxed->False,Axes->False,PlotTheme->"ThickSurface",
PlotStyle->Yellow,MeshStyle->Black,PlotPoints->60]
\end{Verbatim}

\end{comment}

\end{enumerate}
\end{ejer}  
\chapter{Cálculo de volúmenes} \chapter{Cálculo de volúmenes}

En este capítulo se presentan varios ejemplos de sólidos con sus respectivos gráficos y diferentes maneras de hallar su volumen. En primera instancia se pueden usar coordenadas rectangulares y plantear los seis posibles órdenes para evaluar la integral triple; otra posibilidad es utilizar un cambio de coordenadas, ya sean cilíndricas o esféricas; también parametrizar el sólido, esta última idea se apoya en la fórmula \ref{math2}.
En algunos casos se ha incluido la función de densidad del sólido, motivo por el cual se agrega el cálculo de la masa del sólido.
Además, se presenta la manera de obtener los gráficos y evaluar estas integrales con los procedimientos respectivos utilizando el programa \verb"Mathematica".

\section{Sólidos acotados por planos}
Presentamos algunos sólidos acotados principalmente por planos y superficies cuadráticas.

\subsection*{Una pirámide}
A continuación se muestra la figura del sólido acotado por las figuras de los planos $x=2y$, $z=4-x-2y$
en el primer octante.
\vskip0.2cm
Se debe tener en cuenta que en el plano $xy$ este se proyecta en la recta $x=4-2y$. La intersección de los planos proyectada en el plano $yz$ se determina por la ecuación $z=4-4y$ y en el plano $xz$ esta región se determina por $z=4-2x$.

\noindent
\begin{minipage}{\textwidth}
	\centering
	\captionof{figure}{Intersección entre los planos $x + 2y + z = 4$ \quad y \quad $x = 2y$.}
	\label{fig:interseccion_planos}
	\begin{minipage}{0.36\linewidth}
		\centering
		\vspace{-3mm}
		\includegraphics[width=\linewidth]{Fig/148.pdf}
	\end{minipage}
	\hspace{1em}
\begin{minipage}{0.35\linewidth}
	   \centering
	   \vspace{-3mm}
		\includegraphics[width=\linewidth]{Fig/149.pdf}
	\end{minipage}
	
	\vspace{1ex}
	\fuentepropia
\end{minipage}

\vskip0.2cm
Sobre los planos coordenados se determinan las siguientes regiones.

\vskip0.2cm
\noindent
\begin{minipage}{\textwidth}
	\centering
	\captionof{figure}{Proyección de los planos $x + 2y + z = 4$ \quad y \quad $x = 2y$.}
	\label{fig:proyeccion_planos}
	\begin{minipage}{0.28\linewidth}
		\centering
		\includegraphics[width=\linewidth]{Fig/150.pdf}
	\end{minipage}
	\hspace{1em}
	\begin{minipage}{0.28\linewidth}
		\centering
		\includegraphics[width=\linewidth]{Fig/151.pdf}
	\end{minipage}
	\hspace{1em}
	\begin{minipage}{0.28\linewidth}
		\centering
		\includegraphics[width=\linewidth]{Fig/152.pdf}
	\end{minipage}
	
	\vspace{1ex}
	\fuentepropia
\end{minipage}

\vskip0.2cm
Estos gráficos nos ayudan a establecer los límites de las integrales triples que determinan su volumen.
\begin{multline*}
\ds \int_{0}^{1}\int_{0}^{4-4y}\int_{0}^{2y}dxdzdy +\int_{0}^{1}\int_{4-4y}^{4-2y}\int_{0}^{4-z-2y} dxdzdy\\
+ \int_{1}^{2}\int_{0}^{4-2y}\int_{0}^{4-z-2y}dxdzdy=
\ds \int_{0}^{2}\int_{x/2}^{\frac{4-x}{2}}\int_{0}^{4-x-2y}dzdydx\\
= \ds \int_{0}^{2}\int_{0}^{4-2x}\int_{x/2}^{\frac{4-x-z}{2}}dydzdx= \int_{0}^{4}\int_{0}^{\frac{4-z}{2}}\int_{x/2}^{\frac{4-x-z}{2}}dydxdz\\
= \ds \int_{0}^{1}\int_{0}^{2y}\int_{0}^{4-x-2y}dzdxdy+\int_{1}^{2}\int_{0}^{4-2y}\int_{0}^{4-x-2y}dzdxdy \\
= \ds \int_{0}^{4}\int_{0}^{\frac{4-z}{4}}\int_{0}^{2y}dxdydz+\int_{0}^{4}\int_{\frac{4-z}{4}}^{\frac{4-z}{2}}\int_{0}^{4-z-2y}dxdydz =\frac{8}{3}.
\end{multline*}
\vskip0.2cm
La primera de estas integrales se evaluó en \verb"Mathematica" ejecutando las siguientes instrucciones:
\enlargethispage{\baselineskip}
\vspace{0.6em}
\noindent
\begin{tcolorbox}[breakable,enhanced,title=\bf Instrucción en \verb"Mathematica", top=2pt,bottom=2pt,boxsep=0pt,before skip=2pt,after skip=0pt]
\begin{Verbatim}[formatcom=\letraverbatim]
Integrate[1,{y,0,1},{z,0,4-4*y},{x,0,2*y}]+
Integrate[1,{y,0,1},{z,4-4y,4-2*y},{x,0,4-z-2*y}]+
Integrate[1,{y,1,2},{z,0,4-2*y},{x,0,4-z-2*y}]
\end{Verbatim}
\end{tcolorbox}
\vspace{0.7em}
%\noindent
\pagebreak[0]
El volumen también se puede hallar de la siguiente manera.

%\begin{tcolorbox}[title=Parametrizando el sólido]
%\small
	Los vértices de la base son los puntos $(0,0,0)$, $(2,1,0)$ y $(0,2,0)$. Este triángulo se puede parametrizar primero, considerando un segmento entre dos puntos de estos y luego \textbf{formar segmentos rectilíneos con el tercer punto}. Es decir $(1-u)(0,0,0) + u(2,1,0) = (2u, u, 0)$ y luego $(1-v)(2u, u, 0) + v(0,2,0) = (2(1-v)u, (1-v)u + 2v, 0)$. Por tal razón la función $\mathbf{r}(u, v) = \langle 2(1-v)u, (1-v)u + 2v, 0 \rangle$, $u \in [0,1]$, $v \in [0,1]$, representa la base del sólido. El sólido completo se recorre involucrando el vértice $(0,0,4)$, es decir, se debe formar la siguiente combinación convexa: $(1-w)(0,0,4) + w(2(1-v)u, (1-v)u + 2v, 0)$ que al simplificarse, determina la función que parametriza el sólido
	
	\[ \mathbf{r}(u,v,w) = \left\langle 2wu(1-v), wu(1-v)+2wv, 4-4w \right\rangle, \]
	con $u \in [0,1]$, $v \in [0,1]$, $w \in [0,1]$. Es sencillo ver que
	\begin{align*}
		\mathbf{r}_u &= \left\langle 2w(1-v), w(1-v), 0 \right\rangle \\
		\mathbf{r}_v &= \left\langle -2wu, w(-u+2), 0 \right\rangle \\
		\mathbf{r}_w &= \left\langle 2(1-v)u, u-uv+2v, -4 \right\rangle,
	\end{align*}
	con lo que se puede establecer $\mathbf{r}_u \times \mathbf{r}_v = \langle 4w(u-2), -8wu, -4wu \rangle;$ por lo tanto, $|\mathbf{r}_u \cdot (\mathbf{r}_v \times \mathbf{r}_w)| = 16|w^2(v-1)|$, el volumen del sólido está determinado por
	\[ \int_{0}^{1}\int_{0}^{1}\int_{0}^{1} 16|w^2(v-1)|\, dudvdw = \int_{0}^{1}\int_{0}^{1}\int_{0}^{1} 16w^2(1-v)\, dudvdw = \frac{8}{3}. \]
%\end{tcolorbox}
%\clearpage

\noindent
Si la densidad del sólido en un punto $(x,y,z)$ está dada por $\delta(x,y,z) = xyz^{2}$, entonces la masa es:
\begin{align*}
\int_{0}^{2} \int_{x/2}^{\frac{4-x}{2}} \int_{0}^{4 - x - 2y} xyz^{2} \, dz \, dy \, dx 
&= \int_{0}^{2} \int_{0}^{4 - 2x} \int_{x/2}^{\frac{4 - x - z}{2}} xyz^{2} \, dy \, dz \, dx \\
&= \frac{256}{315}.
\end{align*}

\vspace{0.5em}

Con la parametrización escogida anteriormente, la densidad está dada por:
\[
\delta(u,v,w) = 32w^{2}u(1-v)(u - uv + 2v)(1 - w)^{2}.
\]

Y la masa del sólido es:
\begin{center} % Mantiene el centrado si es deseado
	\resizebox{1.0\textwidth}{!}{%
		$\begin{aligned}
			\int_{0}^{1} \int_{0}^{1} \int_{0}^{1} 16 \cdot \left(32w^{2}u(1 - v)^{2}(u - uv + 2v)(1 - w)^{2}\right) w^{2} \, du \, dv \, dw = \frac{256}{315}.
		\end{aligned}$
	}
\end{center}
\subsection*{Una pirámide II}
Las figuras de las expresiones $z=y-x,\,\, z=0,\,\, y=1, \,\, x=0,$
corresponden a cuatro planos que determinan un sólido; hallaremos su volumen.

\begin{figure}[H]
	\centering
	\caption{Intersección entre los planos $x - y + z = 0$ \quad y \quad $y = 1$.}
	\label{fig:interseccion_xy_planos}
	\begin{minipage}[t]{0.48\textwidth}
		\centering
		\includegraphics[width=\linewidth]{Fig/153.pdf}
	\end{minipage}
	\hspace{1em}
	\begin{minipage}[t]{0.35\textwidth}
		\centering
		\includegraphics[width=\linewidth]{Fig/154.pdf}
	\end{minipage}
	
	\vspace{1ex}
	\fuentepropia
\end{figure}



\noindent
El sólido proyectado en cada uno de los planos coordenados proporciona las siguientes figuras:

\noindent
\begin{minipage}{\textwidth}
	\centering
	% Usar \captionof es más robusto cuando no estás en un entorno 'figure' flotante.
	\captionof{figure}{Proyección de los planos $x - y + z = 0$, $y = 1$ y su intersección.} 
	\label{fig:proyeccion_interseccion_planos}
	
	% AJUSTE CRÍTICO DE ANCHOS Y ESPACIOS:
	% (3 * 0.31) + (2 * 0.015) = 0.96 < 1.0 (deja un poco de margen)
	\begin{minipage}{0.31\linewidth} % Ancho reducido
		\centering
		\vspace{-5mm}
		\includegraphics[width=\linewidth]{Fig/155.pdf}
	\end{minipage}
	\hfill % Relleno de espacio flexible
	\begin{minipage}{0.31\linewidth} % Ancho reducido
		\centering
		\vspace{-5mm}
		\includegraphics[width=\linewidth]{Fig/156.pdf}
	\end{minipage}
	\hfill % Relleno de espacio flexible
	\begin{minipage}{0.31\linewidth} % Ancho reducido
		\centering
		\vspace{-5mm}
		\includegraphics[width=\linewidth]{Fig/157.pdf}
	\end{minipage}
	
	\vspace{1ex}
	% Nota al pie con fuente pequeña.
	\fuentepropia
\end{minipage}
\vskip0.2cm


Con ayuda de ellos es posible plantear las integrales triples que determinan el volumen del sólido presentado.
\vspace{-2mm}
\begin{multline*}
\int_{0}^{1}\int_{x}^{1}\int_{0}^{y-x}dzdydx= \int_{0}^{1}\int_{0}^{y}\int_{0}^{y-x}dzdxdy \\ =\int_{0}^{1}\int_{z}^{1}\int_{0}^{y-z}dxdydz
\int_{0}^{1}\int_{0}^{y}\int_{0}^{y-z}dxdzdy \\ =\int_{0}^{1}\int_{0}^{1-z}\int_{x+z}^{1}dydxdz= \int_{0}^{1}\int_{0}^{1-x}\int_{x+z}^{1}dydzdx= \frac{1}{6}
\end{multline*}
Estos cálculos se pueden hacer en \verb"Mathematica" ejecutando las siguientes instrucciones:

\vskip0.2cm
\noindent
\begin{tcolorbox}[breakable,enhanced,title=\bf Instrucción en \verb"Mathematica", top=2pt,bottom=2pt,boxsep=0pt,before skip=2pt,after skip=0pt]
\begin{Verbatim}[formatcom=\letraverbatim]
Integrate[1,{x,0,1},{y,x,1},{z,0,y-x}]
Integrate[1,{y,0,1},{x,0,y},{z,0,y-x}]
Integrate[1,{z,0,1},{y,z,1},{x,0,y-z}]
Integrate[1,{y,0,1},{z,0,y},{x,0,y-z}]
Integrate[1,{z,0,1},{x,0,1-z},{y,x+z,1}]
Integrate[1,{x,0,1},{z,0,1-x},{y,x+z,1}]
\end{Verbatim}
\end{tcolorbox}
\vspace{0.7em}
\noindent

El volumen también se puede hallar con el siguiente procedimiento. En el plano $xy$, el segmento que une los puntos $(0,0)$ con $(0,1)$ se parametrizan como $(0,t)$; y el segmento que une cualquiera de estos puntos con $(1,1)$ es de la forma $(u, t - tu + u)$. En el espacio, el sólido se recorre mediante segmentos que unen los puntos $(u, t - tu + u, 0)$ con $(1, 1, 1)$, esto es $(1-w)(u, t - tu + u, 0) + w(1, 1, 1)$. Por tanto, la función $\mathbf{r}(t, u, w) = \left\langle (u + w(1-u)), t + (t-1)(uw - w - u), w \right\rangle$, con $t \in [0,1]$, $u \in [0,1]$, $w \in [0,1]$, parametriza el sólido. Con esta función hallamos los siguientes vectores.
\vspace{-4mm}
	\begin{align*}
		\mathbf{r}_{t} &= \langle 0, 1 - w + uw - u, 0 \rangle \\
		\mathbf{r}_{u} &= \langle 1 - w, t(1 - w) - (1 - w), 0 \rangle \\
		\mathbf{r}_{w} &= \langle 1 - u, (t - 1)(u - 1), 1 \rangle
	\end{align*}
	
	\noindent
	\textbf{El producto vectorial triple} es $|\mathbf{r}_{t} \cdot (\mathbf{r}_{u} \times \mathbf{r}_{w})| = |(w-1)^{2}(u-1)|$, en este caso es $(1-w)^{2}(1-u)$, al integrarlo nos proporciona el volumen del sólido
	
	\[ \int_{0}^{1}\int_{0}^{1}\int_{0}^{1} (1-w)^{2}(1-u) \, dt\,du\,dw = \frac{1}{6}. \]
	
%\end{tcolorbox}


\vskip0.2cm
El cálculo anterior, incluyendo la parametrización, se consigue con las siguientes instrucciones:
\vskip0.2cm
\noindent
\begin{tcolorbox}[breakable,enhanced,title=\bf Instrucción en \verb"Mathematica", top=2pt,bottom=2pt,boxsep=0pt,before skip=2pt,after skip=0pt]
\begin{Verbatim}[formatcom=\letraverbatim]
r=(1-w)*((1-u){0,t,0}+u*{1,1,0})+w*{1,1,1};
A=Abs[Dot[D[r,t],Cross[D[r,u],D[r,w]]]];
Integrate[A,{u,0,1},{v,0,1},{w,0,1}]
\end{Verbatim}
\end{tcolorbox}
\vspace{0.7em}
\noindent

\subsection*{Tres planos y un cilindro parabólico}
El cilindro parabólico con ecuación $z=4-y^{2}$ y los planos $y=x,$ $x=0$ y $z=0$ limitan en el primer octante del sólido que se muestra a conti\-nuación.

\vskip0.3cm
\noindent
\begin{minipage}{\textwidth}
	\centering
	\captionof{figure}{Sólido determinado por $0 \leq z \leq 4 - y^2$, \quad $0 \leq x \leq y$.}
	\label{fig:solido_parabolico_inclinado_2}
	\begin{minipage}{0.23\linewidth}
		\centering
		\includegraphics[width=\linewidth]{Img/16c.pdf}
	\end{minipage}
	\hspace{1em}
	\begin{minipage}{0.23\linewidth}
		\centering
		\includegraphics[width=\linewidth]{Img/16a.pdf}
	\end{minipage}
	\hspace{1em}
	\begin{minipage}{0.23\linewidth}
		\centering
		\includegraphics[width=\linewidth]{Img/16b.pdf}
	\end{minipage}
	
	\vspace{1ex}
	\fuentepropia
\end{minipage}

\vskip0.3cm

Estas figuras se obtienen ejecutando las siguientes instrucciones:

\vskip0.2cm
\noindent
\begin{tcolorbox}[breakable,enhanced,title=\bf Instrucción en \verb"Mathematica", top=2pt,bottom=2pt,boxsep=0pt,before skip=2pt,after skip=0pt]
\begin{Verbatim}[formatcom=\letraverbatim]
ContourPlot3D[{z==4-y^2,x==y,x==0},{x,-1,2},{y,-2,2},
{z,0,4},PlotPoints->60,Mesh->None,
ContourStyle->{Orange,Blend[{Yellow,Pink,Green}]}]
RegionPlot3D[0<z<4-y^2&&x<y<2&&0<x<2,{x,0,2},{y,0,2},
{z,0,4},PlotPoints->90,Mesh->None]
\end{Verbatim}
\end{tcolorbox}
\vspace{0.7em}
\noindent


\noindent
\begin{minipage}{\textwidth}
	\centering
	\captionof{figure}{Proyección del sólido sobre los planos coordenados.}
	\label{fig:proyeccion_sólido}
	\begin{minipage}{0.29\linewidth}
		\centering
		\vspace{-4mm}
		\includegraphics[width=\linewidth]{Fig/158.pdf}
	\end{minipage}
	\hspace{1em}
	\begin{minipage}{0.29\linewidth}
		\centering
		\vspace{-4mm}
		\includegraphics[width=\linewidth]{Fig/159.pdf}
	\end{minipage}
	\hspace{1em}
	\begin{minipage}{0.29\linewidth}
		\centering
		\vspace{-4mm}
		\includegraphics[width=\linewidth]{Fig/160.pdf}
	\end{minipage}
	
	\vspace{1ex}
	\fuentepropia
\end{minipage}
\vskip0.2cm
Estas figuras nos orientan para plantear las siguientes integrales triples.
\begin{multline*}
\int_{0}^{2}\int_{x}^{2}\int_{0}^{4-y^{2}}dzdydx= \int_{0}^{2}\int_{0}^{y}\int_{0}^{4-y^{2}}dzdxdy \\
= \int_{0}^{4}\int_{0}^{\sqrt{4-z}}\int_{x}^{\sqrt{4-z}}dydxdz
\int_{0}^{2}\int_{0}^{4-x^{2}}\int_{x}^{\sqrt{4-z}}dydzdx \\
\int_{0}^{4}\int_{0}^{\sqrt{4-z}}\int_{0}^{y}dxdydz= \int_{0}^{2}\int_{0}^{4-y^{2}}\int_{0}^{y}dxdzdy= 4.
\end{multline*}
Estas integrales se evalúan en \verb"Mathematica" con las siguientes sentencias:

\vskip0.2cm
\noindent
\begin{tcolorbox}[breakable,enhanced,title=\bf Instrucción en \verb"Mathematica", top=2pt,bottom=2pt,boxsep=0pt,before skip=2pt,after skip=0pt]
\begin{Verbatim}[formatcom=\letraverbatim]
Integrate[1,{x,0,2},{y,x,2},{z,0,4-y^2}]
Integrate[1,{y,0,2},{x,0,y},{z,0,4-y^2}]
Integrate[1,{z,0,4},{x,0,Sqrt[4-z]},{y,x,Sqrt[4-z]}]
Integrate[1,{x,0,2},{z,0,4-x^2},{y,x,Sqrt[4-z]}]
Integrate[1,{z,0,4},{y,0,Sqrt[4-z]},{y,0,y}]
Integrate[1,{y,0,2},{z,0,4-y^2},{x,0,y}]
\end{Verbatim}
\end{tcolorbox}
\vspace{0.7em}
\noindent

Otra posibilidad de obtener este valor es la siguiente.
\vskip0.2cm


%\begin{tcolorbox}[title=\bfseries Parametrizando el sólido]
\noindent	
La base de esta región es un triángulo en el plano $xy$, cuyos vértices son los puntos $(0,0),$ $(0,2)$ y $(2,2).$
Primero, se parametriza el segmento que une los puntos $(0,0)$ con $(2,2),$ con una combinación convexa entre dichos puntos, es decir,
$(1-u)(0,0) +u( 2,2) =( 2u,2u) .$
	
\vspace{1ex} 
	
\noindent
	

Luego se une el punto $(0,2) $ con los puntos de la forma $(2u,2u)$, es decir, se forman los segmentos rectilíneos por medio de la expresión
$(1-v)(0,2)+v ( 2u,2u)=\langle 2uv,2-2v+2uv\rangle.$
Algunos de estos segmentos se muestran a continuación.

	
\vskip0.2cm
\noindent
\begin{minipage}{\textwidth}
	\centering
	\captionof{figure}{$(2uv, 2(1 - v + uv)).$}
	\label{fig:curva_parametrica_uv_2}
	\begin{minipage}{0.32\linewidth}
		\centering
		\includegraphics[width=\linewidth]{Fig/161.pdf}
	\end{minipage}
	
	\vspace{0.3em}
	\fuentepropia
\end{minipage}
\vskip0.2cm
	
	\vspace{1ex} 
	\noindent
	
Sobre la superficie del cilindro parabólico, la coordenada en $z$ se puede obtener con la ecuación que lo define, es decir, \[z= 4-y^{2}=4-(2-2v+2uv)^{2} =4v(1-u)(uv+2-v).\]
%\end{tcolorbox}


%\begin{tcolorbox}
%\small
\begin{figure}[H]
	\centering
	\caption{Sólido determinado por $0 \leq z \leq 4 - y^2$, \quad $x \leq y \leq 2$.}
	\label{fig:solido_parabolico_y}
	\vspace{0.3cm} 
	\begin{minipage}[t]{0.33\textwidth}
		\centering
		\includegraphics[width=\linewidth]{Fig/162.pdf}
	\end{minipage}
	\hspace{1em}
	\begin{minipage}[t]{0.33\textwidth}
		\centering
		\includegraphics[width=\linewidth]{Fig/163.pdf}
	\end{minipage}
	
	\vspace{0.4cm} 
	
	\fuentepropia
\end{figure}

\vspace{-4mm}
Es de anotar que la curva sobre el cilindro en el plano $x=0$, se determina por medio de la expresión
$\langle 0,t,4-t^{2}\rangle$ y sobre el plano $y=x$, se usa $\langle t,t,4-t^{2}\rangle$.\\

Para efectos del gráfico, la parte sobre el cilindro se parametriza con
$(1-u)\langle 0,t,4-t^{2}\rangle +u\langle t,t,4-t^{2}\rangle=\langle ut,t,4-t^{2}\rangle$. \\[0.25cm]
Con lo anterior, se puede establecer que un punto $Q$ sobre la superficie del cilindro tiene la forma
$(2uv,2-2v+2uv,4v(1-u)(uv+2-v))$, mientras que un punto $P$ sobre la base del sólido tiene la forma $(2uv,2-2v+2uv,0).$
Ahora se puede parametrizar el sólido formando una combinación convexa de los puntos $P$ y $Q$,
esto es $(1-w)Q+wP$,
%$(1-w) \langle 2uv,2-2v+2uv,0\rangle+w\langle 2uv,2-2v+2uv,4v(1-u)(uv+2-v)\rangle$
que al simplificarse nos permite definir la función
${\bf r}( u,v,w) = \langle 2uv,2-2v+2uv,4wv(1-u)(uv+2-v)\rangle,$
para $u\in [0,1]$, $v\in [0,1]$, $w\in [0,1].$ De ella se obtienen los siguientes vectores:
\vspace{-3mm}
\begin{eqnarray*}
{\bf r}_{u}&=& \langle 2v,2v,8wv( v-uv-1)\rangle \\
{\bf r}_{v}&=& \langle 2u,-2+2u,-8w(u-1) ( uv-v+1) \rangle \\
{\bf r}_{w}&=& \langle 0,0,4v( 1-u) (uv+2-v) \rangle
\end{eqnarray*}
\vspace{-4pt}
Además, ${\bf r}_{v}\times {\bf r}_{w}= \big\langle -8v( u-1)^{2}(uv+2-v),8uv( u-1)(uv+2-v),0 \big \rangle,$
entonces, $ |{\bf r}_{u}\cdot({\bf r}_{v}\times {\bf r}_{w})|=16|v^{2}(u-1)(uv+2-v)|.$ Por último, el volumen del sólido está dado por
{\small
\[\int_{0}^{1}\int_{0}^{1}\int_{0}^{1}16| v^{2}( u-1)(uv+2-v)|dudvdw= 4.\]
}
%\end{tcolorbox}

%\clearpage
Los cálculos anteriores se pueden efectuar en \verb"Mathematica"
definiendo una parametrización $p$ para la base del sólido, $x$ e $y$ son las coordenadas de $p.$ Se determina una parametrización ${\bf r}$ del sólido y se evalúa el triple producto vectorial, que se integra para encontrar el volumen deseado.
\vskip0.2cm
\noindent
\begin{tcolorbox}[breakable,enhanced,title=\bf Instrucción en \verb"Mathematica",top=2pt,bottom=2pt,boxsep=0pt,before skip=2pt,after skip=0pt]
\begin{Verbatim}[formatcom=\letraverbatim]
p={x,y}=(1-v)*{0,2}+v*{2*u,2*u}
r=(1-w)*{x,y,0}+w*{x,y,4-y^2};
A=Factor[Dot[D[r,u],Cross[D[r,v],D[r,w]]]]
Integrate[Abs[A],{u,0,1},{v,0,1},{w,0,1}]
\end{Verbatim}
\end{tcolorbox}
\vspace{0.7em}
\noindent

\subsection*{Cuatro planos y un cilindro parabólico}

La porción de las figuras de las expresiones $z=2-x$, $z^2=x$, $y=0$, $y=3$ y $z=0$ en el primer octante se muestra a continuación; estas acotan una región en el espacio.

\vskip0.2cm
\noindent
\begin{minipage}{\textwidth}
	\centering
	\captionof{figure}{Superficies delimitantes: $z = 2 - x$, $z^2 = x$, $y = 0$, $y = 3$, $z = 0$.}
	\label{fig:superficies_z2x_z2x_y03_2}
	\begin{minipage}{0.34\linewidth}
		\centering
		\vspace{-3mm}
		\includegraphics[width=\linewidth]{Fig/164.pdf}
	\end{minipage}
	\hspace{1em}
	\begin{minipage}{0.29\linewidth}
		\centering
		\vspace{-3mm}
		\includegraphics[width=\linewidth]{Fig/165.pdf}
	\end{minipage}
	
	\vspace{1ex}
	\fuentepropia
\end{minipage}

\vskip0.2cm
Ahora se muestra el sólido en mención y las instrucciones que lo generan.


\vskip0.2cm
\noindent
\begin{minipage}{\textwidth}
	\centering
	\captionof{figure}{Sólido acotado por $z = 2 - x$, $z^2 = x$, $y = 0$, $y = 3$, $z = 0$}
	\label{fig:solido_z2x_z2x_y03}
	\begin{minipage}{0.19\linewidth}
		\centering
		\includegraphics[width=\linewidth]{Img/15a.pdf}
	\end{minipage}
	\begin{minipage}{0.19\linewidth}
		\centering
		\includegraphics[width=\linewidth]{Img/15b.pdf}
	\end{minipage}
	
	\vspace{0.3em}
	\fuentepropia
\end{minipage}
\vskip0.2cm



\noindent
\begin{tcolorbox}[breakable,enhanced,title=\bf Instrucción en \verb"Mathematica", top=2pt,bottom=2pt,boxsep=0pt,before skip=2pt,after skip=0pt]
\begin{Verbatim}[formatcom=\letraverbatim]
RegionPlot3D[z^2<x<2-z&&0<y<3&&0<x<2,{x,0,2},{y,0,3},
{z,0,1},PlotPoints->90,Mesh->None]
\end{Verbatim}
\end{tcolorbox}
\vspace{0.7em}
\noindent


\noindent
\begin{minipage}{\textwidth}
	\centering
	\captionof{figure}{Proyección del sólido en los planos coordenados.}
	\label{fig:proyeccion_sólido_coordenados}
	\begin{minipage}{0.28\linewidth}
		\centering
		\includegraphics[width=\linewidth]{Fig/166.pdf}
	\end{minipage}
	\hspace{1em}
	\begin{minipage}{0.28\linewidth}
		\centering
		\includegraphics[width=\linewidth]{Fig/167.pdf}
	\end{minipage}
	\hspace{1em}
	\begin{minipage}{0.28\linewidth}
		\centering
		\includegraphics[width=\linewidth]{Fig/168.pdf}
	\end{minipage}
	\vspace{0.2cm}
	\vspace{0.3em}
	\fuentepropia
\end{minipage}
	

El volumen del sólido se puede determinar evaluando, según el caso, alguna de las siguientes integrales.
{\small
\[\int_{0}^{1}\int_{0}^{3}\int_{z^{2}}^{2-z}dxdydz=\int_{1}^{2}\int_{0}^{3}\int_{0}^{2-x}dzdydx+\int_{0}^{1}\int_{0}^{3}\int_{0}^{\sqrt{x}}dzdydx =\frac{7}{2}\]
\[\int_{0}^{3}\int_{1}^{2}\int_{0}^{2-x}dzdxdy+\int_{0}^{3}\int_{0}^{1}\int_{0}^{\sqrt{x}}dzdxdy =\int_{0}^{3}\int_{0}^{1}\int_{z^{2}}^{2-z}dxdzdy=\frac{7}{2}\]
\[\int_{0}^{1}\int_{z^{2}}^{2-z}\int_{0}^{3}dydxdz= \int_{0}^{1}\int_{0}^{\sqrt{x}}\int_{0}^{3}dydzdx
+\int_{1}^{2}\int_{0}^{2-x}\int_{0}^{3}dydzdx= \frac{7}{2}.\]
}

\vskip0.2cm
Dos de estas integrales se pueden evaluar con las siguientes sentencias:

\vskip0.2cm
\noindent
\begin{tcolorbox}[breakable,enhanced,title=\bf Instrucción en \verb"Mathematica", top=2pt,bottom=2pt,boxsep=0pt,before skip=2pt,after skip=0pt]
\begin{Verbatim}[formatcom=\letraverbatim]
Integrate[1,{z,0,1},{y,0,3},{x,z^2,2-z}]
Integrate[1,{x,1,2},{y,0,3},{z,0,2-x}]+
Integrate[1,{x,0,1},{y,0,3},{z,0,Sqrt[x]}]
\end{Verbatim}
\end{tcolorbox}
\vspace{0.7em}
\noindent

Este volumen también se puede obtener de la siguiente manera.
\vskip0.2cm
%\begin{tcolorbox}[title=\bfseries Parametrizando el sólido]
	%\small
	\noindent
	En el plano $xz$, para $t \in [0,1]$, los puntos $A$ y $B$ son de la forma $(t^{2},t)$ y $(2-t,t)$ respectivamente; el segmento que los une se parametriza mediante la combinación convexa $(1-u)A+uB,$ que se simplifica como $(t^{2}(1-u)+u(2-t),t),$ para $u\in [0,1].$
	\vskip0.3cm
	%\noindent
	\begin{minipage}{\textwidth}
		\captionof{figure}{Superficies delimitantes: $z = 2 - x$, $z^2 = x$, $y = 0$, $y = 3$, $z = 0$.}
		\label{fig:superficies_z2x_z2x_y03_3}
		\begin{minipage}{0.31\linewidth}
			\centering
			\vspace{-4mm}
			\includegraphics[width=\linewidth]{Fig/206.pdf}
		\end{minipage}
		\hspace{1em}
		\begin{minipage}{0.37\linewidth}
			\centering
			\vspace{-4mm}
			\includegraphics[width=\linewidth]{Fig/207.pdf}
		\end{minipage}
		%\vspace{0.2cm}
		\vspace{0.3em}
		\fuentepropia
	\end{minipage}
	
	\vspace{1ex} 
	El sólido se parametriza formando segmentos rectilíneos entre dos puntos $P(t^{2}(1-u)+(2-t)u,0,t)$ y $Q(t^{2}(1-u)+(2-t)u,3,t).$
	Simplificando $(1-w)P+wQ$ se define
	 \[\mathbf{r}(t,u,w) = \left\langle t^{2}(1-u)-u(t-2), 3w, t \right\rangle,\]
con $t\in [0,1]$, $u\in [0,1]$ y $w\in [0,1]$. Además,
%\vspace{-10mm}
{\small
\begin{align*}
	\mathbf{r}_t &= \frac{\partial}{\partial t}\left\langle t^{2}(1-u)-u(t-2),3w,t\right\rangle = \left\langle 2t(1-u)-u,0,1\right\rangle \\
	\mathbf{r}_u &= \frac{\partial}{\partial u}\left\langle t^{2}(1-u)-u(t-2),3w,t\right\rangle = \left\langle -t^{2}+2-t,0,0\right\rangle \\
	\mathbf{r}_w &= \frac{\partial}{\partial w}\left\langle t^{2}(1-u)-u(t-2),3w,t\right\rangle = \left\langle 0,3,0\right\rangle.
\end{align*}
}
	
	El elemento diferencial de volumen es $\left| \mathbf{r}_t\cdot (\mathbf{r}_u\times \mathbf{r}_w) \right|
	= \left|6-3t-3t^{2}\right| ,$
	al integrarse determina el volumen del sólido
	
	\vspace{-5mm}
	
	\[ \int_{0}^{1}\int_{0}^{1}\int_{0}^{1} |6-3t-3t^{2}| \, dt\,du\,dw = \int_{0}^{1}\int_{0}^{1}\int_{0}^{1}(6-3t-3t^{2}) \, dt\,du\,dw = \frac{7}{2}. \]

%\end{tcolorbox}


\subsection*{Un cilindro y dos planos}
Considere la región en el primer octante, acotada por las figuras de las ecuaciones
$y^2+z^2=9$, $z=x$ y $y=0$. Su gráfico se muestra a conti\-nuación:

\vskip0.2cm
\noindent
\begin{minipage}{\textwidth}
	\centering
	\captionof{figure}{Sólido determinado por $0 \leq y \leq \sqrt{9-z^2}$ \quad y \quad $0 \leq x \leq z$.}
	\label{fig:solido_yz_xz}
	\begin{minipage}{0.35\linewidth}
		\centering
		\vspace{-2mm}
		\includegraphics[width=\linewidth]{Fig/208.pdf}
	\end{minipage}
	\hspace{1em}
	\begin{minipage}{0.35\linewidth}
		\centering
		\vspace{-2mm}
		\includegraphics[width=\linewidth]{Fig/209.pdf}
	\end{minipage}
	\vspace{0.2cm}
	\fuentepropia
\end{minipage}


Enseguida se muestra la región proyectada sobre los planos coordenados.

\vskip0.2cm
\noindent
\begin{minipage}{\textwidth}
	\centering
	\captionof{figure}{Proyecciones de $y^2 + z^2 = 9$, $z = x$, $y = 0$ y sus intersecciones.}
	\label{fig:proyecciones_cilindro_plano}
	\begin{minipage}{0.29\linewidth}
		\centering
		\vspace{-2mm}
		\includegraphics[width=\linewidth]{Fig/210.pdf}
	\end{minipage}
	\hspace{1em}
	\begin{minipage}{0.29\linewidth}
	\centering
	\vspace{-2mm}
		\includegraphics[width=\linewidth]{Fig/211.pdf}
	\end{minipage}
	\hspace{1em}
	\begin{minipage}{0.29\linewidth}
		\centering
		\vspace{-2mm}
		\includegraphics[width=\linewidth]{Fig/212.pdf}
	\end{minipage}
	\vspace{0.2cm}
	
	\fuentepropia
\end{minipage}

\vskip0.2cm
Con ayuda de estas figura se pueden plantear las integrales necesarias para determinar el volumen del mencionado sólido.
\begin{multline*}
\int_{0}^{3}\int_{0}^{\sqrt{9-x^{2}}}\int_{x}^{\sqrt{9-y^{2}}}dzdydx =
\int_{0}^{3}\int_{0}^{\sqrt{9-y^{2}}}\int_{0}^{z}dxdzdy \\
=\int_{0}^{3}\int_{0}^{\sqrt{9-y^{2}}}\int_{x}^{\sqrt{9-y^{2}}}dzdxdy
\int_{0}^{3}\int_{0}^{\sqrt{9-z^{2}}}\int_{0}^{z}dxdydz \\
= \int_{0}^{3}\int_{0}^{z}\int_{0}^{\sqrt{9-z^{2}}}dydxdz = \int_{0}^{3}\int_{x}^{3}\int_{0}^{\sqrt{9-z^{2}}}dydzdx = 9
\end{multline*}
\vskip0.2cm
Tres de estas integrales se evalúan con \verb"Mathematica" usando las siguientes instrucciones:

\vskip0.2cm
\noindent
\begin{tcolorbox}[breakable,enhanced,title=\bf Instrucción en \verb"Mathematica", top=2pt,bottom=2pt,boxsep=0pt,before skip=2pt,after skip=0pt]
\begin{Verbatim}[formatcom=\letraverbatim]
Integrate[1,{z,0,3},{y,0,Sqrt[9-z^2]},{x,0,z}]
Integrate[1,{y,0,3},{x,0,Sqrt[9-y^2]},{z,x,
Sqrt[9-y^2]}]
Integrate[1,{x,0,3},{z,x,3},{y,0,Sqrt[9-z^2]}]
\end{Verbatim}
\end{tcolorbox}
\vspace{0.7em}
\noindent

Con el siguiente proceso también se halla este volumen.
\vskip0.2cm
%\begin{tcolorbox}[title=\bfseries Parametrizando el sólido]
Sobre el plano $yz,$ la región de integración se parametriza como
\[(0,3u\cos v,3u\sin v),\] para $u\in \lbrack 0,3],$ $v\in [0,\frac{\pi}{2}].$ Un punto sobre el plano $z=x,$ tiene la forma
$P(3u\sin v,3u\cos v,3u\sin v) ,$ mientras que sobre el plano $yz$ un punto arbitrario $Q$ tiene la forma $Q(0,3u\cos v,3u\sin v) .$
El sólido se parametriza formando segmentos rectilíneos puntos $P$ y $Q,$
esto define la función ${\bf r}$, simplificando la expresión $(1-w)P + wQ,$
\begin{eqnarray*}
{\bf r}(u,v,w) &=&\langle3 u w \sin (v),3 u \cos (v),3 u \sin (v)\rangle,
\end{eqnarray*}
con $u\in [0,1],$ $v\in [ 0,\frac{\pi }{2}],$ $w\in [0,1].$ Ahora se hallan los siguientes vectores
\begin{eqnarray*}
{\bf r}_{u}&=&\left\langle 3 w \sin (v),3 \cos (v),3 \sin (v)\right\rangle \\
{\bf r}_{v}&=&\left\langle 3 u w \cos (v),-3 u \sin (v),3 u \cos (v) \right\rangle \\
{\bf r}_{w}&=&\left\langle 3 u \sin (v),0,0 \right\rangle .
\end{eqnarray*}
Además $\left| {\bf r}_{u}\cdot \left( {\bf r}_{v}\times {\bf r}_{w}\right) \right| = 27u^{2}\sin(v),$ que determina el elemento diferencial de volumen, al integrarlo se halla el volumen del sólido.
\[\int_{0}^{1}\int_{0}^{\pi /2}\int_{0}^{1} 27u^{2}\sin(v) dudvdw=9.\]
%\end{tcolorbox}

Los cálculos mostrados antes se pueden efectuar brevemente, usando \verb"Mathematica",
mediante la ejecución de las siguientes instrucciones:

\vskip0.2cm
\noindent
\begin{tcolorbox}[breakable,enhanced,title=\bf Instrucción en \verb"Mathematica", top=2pt,bottom=2pt,boxsep=0pt,before skip=2pt,after skip=0pt]
\begin{Verbatim}[formatcom=\letraverbatim]
{y,z}={3*u*Cos[v],3*u*Sin[v]};
r=(1-w)*{0,y,z}+w*{z,y,z};
A=Dot[D[r,u],Cross[D[r,v],D[r,w]]];
Integrate[A,{u,0,1},{v,0,Pi/2},{w,0,1}]
\end{Verbatim}
\end{tcolorbox}
\vspace{0.7em}
\noindent

En coordenadas cilíndricas, escritas en  forma
$y=r \cos \theta$, $z=r \sin \theta$, con $\theta \in[0,\pi/2]$, $r \in [0,3];$
esta integral se evalúa como sigue:

$\int_0^{\pi/2} \int_0^3 \int_0^{r\sin\theta} r dxdr d\theta =9. $
El cálculo anterior se obtuvo en \verb"Mathematica" mediante la ejecución de la siguiente instrucción:

\begin{tcolorbox}[breakable,enhanced]
\begin{Verbatim}[formatcom=\letraverbatim]
Integrate[r,{\theta,0,Pi/2},{r,0,3},{x,0,r*Cos[\theta]}
\end{Verbatim}
\end{tcolorbox}


\subsection*{Un cilindro y dos planos II}
Las expresiones
$z=-y$, $x^2+y^2=x$ y $z=0$ corresponden a dos planos y un cilindro. Para $y<0$, estas expresiones acotan el sólido que se muestra a continuación.

\begin{figure}[H]
\centering
\caption{Sólido determinado por $x^2 + y^2 \leq x$, \quad para \quad $0 \leq z \leq -y$.}
\label{fig:solido_disco_inclinado}
\vspace{-2mm}
\includegraphics[width=4.0cm,height=3.5cm]{Img/12c.pdf}
\vspace{-2mm}
\includegraphics[height=4.5cm]{Img/12a.pdf}
\vspace{-2mm}
\includegraphics[height=4.5cm]{Img/12b.pdf}
\vspace{2mm}
\fuentepropia
\end{figure}


Estas figuras se obtienen ejecutando las siguientes instrucciones:

\vskip0.2cm
\noindent
\begin{tcolorbox}[breakable,enhanced,title=\bf Instrucción en \verb"Mathematica", top=2pt,bottom=2pt,boxsep=0pt,before skip=2pt,after skip=0pt]
\begin{Verbatim}[formatcom=\letraverbatim]
ContourPlot3D[{z==-y,x^2+y^2==x,z==0},{x,0,1},
{y,-1,1},{z,-1,1},Ticks->{{0,1},{0,1},{0,1}},
Mesh->False,ContourStyle->{Orange,LightGreen,Yellow}]

RegionPlot3D[0<z<-y&&x>x^2+y^2,{x,0,1},{y,-1/2,0},
{z,0,1/2},PlotPoints->90,Mesh->False,PlotStyle->Orange]
\end{Verbatim}
\end{tcolorbox}
\vspace{0.7em}
\noindent

Las expresiones $x^2+y^2\leq x,$ $0\leq z \leq -y$ determinan un triángulo en el plano $yz$ y un semicírculo en los planos $xz$ y $xy$.


\vskip0.2cm
\noindent
\begin{minipage}{\textwidth}
	\centering
	\captionof{figure}{Proyecciones del sólido en los planos coordenados}
	\label{fig:proyecciones_solido}
	\begin{minipage}{0.27\linewidth}
		\centering
		\includegraphics[width=\linewidth]{Fig/213.pdf}
	\end{minipage}
	\begin{minipage}{0.27\linewidth}
		\centering
		\includegraphics[width=\linewidth]{Fig/214.pdf}
	\end{minipage}
	\begin{minipage}{0.27\linewidth}
		\centering
		\includegraphics[width=\linewidth]{Fig/215.pdf}
	\end{minipage}
	
	\vspace{0.3em}
	\fuentepropia
\end{minipage}
\vskip0.2cm


El volumen del sólido está dado por cualquiera de las siguientes integrales.

\vspace{1ex} 
\noindent
\makebox[\linewidth][l]{
	\scriptsize % <--- REDUCCIÓN MÁXIMA DE LA FUENTE
	$\begin{aligned}
		\bigintsss_{0}^{1} \bigintsss_{-\sqrt{x-x^{2}}}^{0}\ds \bigintsss_{0}^{-y}dz\,dy\,dx &= \bigintsss_{-1/2}^{0}\bigintsss_{\frac{1}{2}-\sqrt{\frac{1}{4}-y^{2}}}^{\frac{1}{2}+ \sqrt{\frac{1}{4}-y^{2}}}\bigintsss_{0}^{-y}dz\,dx\,dy \\
		&= \bigintsss_{0}^{1}\bigintsss_{0}^{\sqrt{x-x^{2}}} \bigintsss_{-\sqrt{x-x^{2}}}^{-z}dy\,dz\,dx \\
		&= \bigintsss_{0}^{1/2}\bigintsss_{\frac{1}{2}-\sqrt{\frac{1}{4}-z^{2}}}^{\frac{1}{2}+ \sqrt{\frac{1}{4}-z^{2}}}\bigintsss_{-\sqrt{x-x^{2}}}^{-z}dy\,dx\,dz \\
		&= \bigintsss_{0}^{1/2}\bigintsss_{-1/2}^{-z}\bigintsss_{\frac{1}{2}-\sqrt{\frac{1}{4}-y^{2}}}^{\frac{1}{2}+\sqrt{\frac{1}{4}-y^{2}}}dx\,dy\,dz \\
		&= \bigintsss_{-1/2}^{0} \bigintsss_{0}^{-y}\bigintsss_{\frac{1}{2}-\sqrt{\frac{1}{4}-y^{2}}}^{\frac{1}{2}+\sqrt{\frac{1}{4}-y^{2}}}dx\,dz\,dy = \frac{1}{12}
	\end{aligned}$
}
\vspace{1ex}

Estas integrales se pueden evaluar en \verb"Mathematica" mediante el uso de las siguientes instrucciones:

\vskip0.2cm
\noindent
\begin{tcolorbox}[breakable,enhanced,title=\bf Instrucción en \verb"Mathematica", top=2pt,bottom=2pt,boxsep=0pt,before skip=2pt,after skip=0pt]
\begin{Verbatim}[formatcom=\letraverbatim]
Integrate[1,{x,0,1},{y,-Sqrt[x-x^2],0},{z,0,-y}]
Integrate[1,{y,-1/2,0},{x,1/2-Sqrt[1/4-y^2],1/2
+Sqrt[1/4-y^2]},{z,0,-y}]
Integrate[1,{x,0,1},{z,0,Sqrt[x-x^2]},
{y,-Sqrt[x-x^2],-z}]
NIntegrate[1,{z,0,1/2},{x,1/2-Sqrt[1/4-z^2],
1/2+Sqrt[1/4-z^2]},{y,-Sqrt[x-x^2],-z}]
NIntegrate[1,{z,0,1/2},{y,-1/2,-z},
{x,1/2-Sqrt[1/4-y^2],1/2+Sqrt[1/4-y^2]}]
NIntegrate[1,{y,-1/2,0},{z,0,-y},{x,1/2-Sqrt[1/4-y^2],
1/2+Sqrt[1/4-y^2]}]
\end{Verbatim}
\end{tcolorbox}
\vspace{0.7em}
\noindent


Si la densidad del sólido está dada por $\delta(x,y,z) =xy^{2},$ entonces la masa del sólido está dada por
\begin{multline*} \bigintsss_{0}^{1}\bigintsss_{-\sqrt{x-x^{2}}}^{0}\bigintsss_{0}^{-y}xy^{2}dzdydx \\ =\bigintsss_{-1/2}^{0}\bigintsss_{\frac{1}{2}-\sqrt{\frac{1}{4}-y^{2}}}^{\frac{1}{2}+\sqrt{\frac{1}{4}-y^{2}}}\bigintsss_{0}^{-y}xy^{2}dzdxdy= \frac{1}{240}.
\end{multline*}
% $\bigintsss_{0}^{1}\bigintsss_{0}^{\sqrt{x-x^{2}}}\bigintsss_{-\sqrt{x-x^{2}}}^{-z}xz^{2}dydzdx=\bigintsss_{0}^{1/2}\bigintsss_{\frac{1}{2}-\sqrt{\frac{1}{4}-z^{2}}}^{\frac{1}{2}+\sqrt{\frac{1}{4}-z^{2}}}\bigintsss_{-\sqrt{x-x^{2}}}^{-z}xz^{2}dydxdz= \frac{1}{720}$
%$\bigintsss_{0}^{1/2}\bigintsss_{-1/2}^{-z}\bigintsss_{\frac{1}{2}-\sqrt{\frac{1}{4}-y^{2}}}^{\frac{1}{2}+\sqrt{\frac{1}{4}-y^{2}}}xz^{2}dxdydz=\bigintsss_{-1/2}^{0}\bigintsss_{0}^{-y}\bigintsss_{\frac{1}{2}-\sqrt{\frac{1}{4}-y^{2}}}^{\frac{1}{2}+\sqrt{\frac{1}{4}-y^{2}}}xz^{2}dxdzdy= \frac{1}{720}$
El volumen y la masa del sólido también se pueden determinar con el siguiente procedimiento.
%\vskip0.1cm

\vskip0.3cm
%\begin{tcolorbox}[title=\bfseries Parametrizando el sólido]
A partir de la ecuación $x^{2}+y^{2}=x$, que equivale a
$\left( x-\frac{1}{2}\right) ^{2}+y^{2}=\frac{1}{4},$ que representa una circunferencia cuyo interior se puede parametrizar como $x=\frac{1}{2}+u\cos v$, $y=u\sin v$, para concluir que un punto en la base del sólido puede parametrizarse en la forma,
$P=\left( \frac{1}{2}+u\cos v,u\sin v,0\right) $ y un punto del sólido sobre el plano con ecuación $z=-y,$ es la forma
$Q=\left( \frac{1}{2}+u\cos v,u\sin v,-u\sin v\right).$
\vskip0.3cm
Se pueden formar segmentos rectilíneos que unan estos puntos y parame- \\ trizar así el sólido. Esto es $(1-w)Q+wP$
es por esta razón que se usara la función
\[{\bf r}\left( u,v,w\right) =\left\langle \frac{1}{2}+u\cos (v) ,u\sin
(v) , (w-1)u\sin (v) \right\rangle\]
para $u\in [0,1/2],\,v\in [-\pi,0],\,w\in [0,1],\, $
con lo anterior, obtenemos los siguientes vectores
\begin{eqnarray*}
{\bf r}_{u}&=&\langle \cos (v) ,\sin (v),\sin (v) \left( w-1\right) \rangle\\
{\bf r}_{v}&=&\langle -u\sin (v) ,u\cos (v) ,u\cos(v) \left( w-1\right) \rangle\\
{\bf r}_{w}&=& \langle 0,0,u\sin (v) \rangle
\end{eqnarray*}
El elemento diferencial de volumen es
$| {\bf r}_{u}\cdot \left( {\bf r}_{v}\times {\bf r}_{w}\right)| =-u^{2}\sin (v) ,$
entonces, el volumen del sólido esta dado por
\vspace{-4mm}
{\small
\begin{multline*}
\ds \int_{0}^{1}\ds \int_{-\pi }^{0}\ds \int_{0}^{1/2}\left| u^{2}\sin (v)
\right| dudvdw \\ = \int_{0}^{1}\ds \int_{-\pi }^{0} \int_{0}^{1/2} -u^{2}\sin (v)dudvdw=\frac{1}{12}.
\end{multline*}
}
%\end{tcolorbox}
	
Con las siguientes indicaciones se integra el término
$|{\bf r}_{u}\cdot \left( {\bf r}_{v}\times {\bf r}_{w}\right)|$, que con los límites adecuados nos produce su volumen.

\vspace{0.6em}
\noindent
\begin{tcolorbox}[breakable,enhanced,title=\bf Instrucción en \verb"Mathematica", top=2pt,bottom=2pt,boxsep=0pt,before skip=2pt,after skip=0pt]
\begin{Verbatim}[formatcom=\letraverbatim]
{x,y,z}={1/2+u*Cos[v],u*Sin[v],-y}
r=(1-w)*{x,y,0}+w*{x,y,z};
M=Abs[Dot[D[r,u],Cross[D[r,v],D[r,w]]]];
Integrate[M,{u,0,1/2},{v,-Pi,0},{w,0,1}]
\end{Verbatim}
\end{tcolorbox}
\vspace{0.7em}
\noindent

\begin{tcolorbox}
Si la densidad en cualquier punto $(x,y,z)$ de sólido es $xy^{2}$,
entonces, $\delta \left( u,v,w\right) =u^2 \sin ^2(v) \left(u \cos (v)+\frac{1}{2}\right)$. Además
{\small
\[\delta \left( u,v,w\right) | {\bf r}_{u}\cdot \left( {\bf r}_{v}\times {\bf r}_{w}\right)|=\left|-\frac{1}{2}u^4\sin^3(v)(2u\cos(v)+1)\right|\]}

por lo tanto, la masa del sólido esta dada por
$\int_{0}^{1}\int_{-\pi }^{0}\int_{0}^{1/2}\left|-\frac{1}{2}u^4\sin^3(v)(2u\cos(v)+1)\right| dudvdw=\frac{1}{240} $
\end{tcolorbox}


\subsection*{Tres planos y un cilindro parabólico II}
Las siguientes ecuaciones $z=2-y,$ $x= 4-y^2,$ $x=0,$ $y=0 $ definen un sólido en el espacio.

\vskip0.2cm
\noindent
\begin{minipage}{\textwidth}
	\centering
	\captionof{figure}{Sólido determinado por $0 \leq y + z \leq 2$ \quad y \quad $0 \leq x \leq 4 - y^2$.}
	\label{fig:solido_yz_xy}
	\begin{minipage}{0.29\linewidth}
		\centering
		\vspace{-4mm}
		\includegraphics[width=\linewidth]{Fig/216.pdf}
	\end{minipage}
	\hspace{1em}
	\begin{minipage}{0.29\linewidth}
		\centering
		\vspace{-4mm}
		\includegraphics[width=\linewidth]{Fig/217.pdf}
	\end{minipage}
	
	\vspace{1ex}
	\fuentepropia
\end{minipage}
\vskip0.2cm


Este sólido determina a su vez las siguientes regiones.

\vskip0.2cm
\noindent
\begin{minipage}{\textwidth}
	\centering
	\captionof{figure}{Proyección del sólido en los planos coordenados.}
	\label{fig:proyecciones_planes_coordenados}
	\begin{minipage}{0.28\linewidth}
		\centering
		\includegraphics[width=\linewidth]{Fig/218.pdf}
	\end{minipage}
	\hspace{1em}
	\begin{minipage}{0.28\linewidth}
		\centering
		\includegraphics[width=\linewidth]{Fig/219.pdf}
	\end{minipage}
	\hspace{1em}
	\begin{minipage}{0.28\linewidth}
		\centering
		\includegraphics[width=\linewidth]{Fig/220.pdf}
	\end{minipage}
	\vspace{0.2cm}
	\fuentepropia
\end{minipage}
\vskip0.2cm


Con ellas puede ser más sencillo ver los límites de integración al momento de plantear las integrales múltiples en los seis órdenes posibles, que son, respectivamente, las siguientes:
\begin{multline*}
\int_{0}^{4}\int_{0}^{\sqrt{4-x}}\int_{0}^{2-y}dzdydx= \int_{0}^{2}\int_{0}^{4-y^{2}}\int_{0}^{2-y}dzdxdy\\
=\int_{0}^{2}\int_{0}^{2-z}\int_{0}^{4-y^{2}}dxdydz= \int_{0}^{2}\int_{0}^{2-y}\int_{0}^{4-y^{2}}dxdzdy\\
=\int_{0}^{2}\int_{0}^{4z-z^{2}}\int_{0}^{2-z}dydxdz+\int_{0}^{2}%
\int_{4z-z^{2}}^{4}\int_{0}^{\sqrt{4-x}}dydxdz \\
= \int_{0}^{4}\int_{2-\sqrt{4-x}}^{2}\int_{0}^{2-z}dydzdx+\int_{0}^{4}%
\int_{0}^{2-\sqrt{4-x}}\int_{0}^{\sqrt{4-x}}dydzdx= \frac{20}{3}.
\end{multline*}

En \verb"Mathematica", las siguientes instrucciones evalúan estas integrales.

\vspace{0.6em}
\noindent
\begin{tcolorbox}[breakable,enhanced,title=\bf Instrucción en \verb"Mathematica", top=2pt,bottom=2pt,boxsep=0pt,before skip=2pt,after skip=0pt]
\begin{Verbatim}[formatcom=\letraverbatim]
Integrate[1,{x,0,4},{z,2-Sqrt[4-x],2},{y,0,2-z}]+
Integrate[1,{x,0,4},{z,0,2-Sqrt[4-x]},{y,0,Sqrt[4-x]}]
Integrate[1,{z,0,2},{x,0,4*z-z^2},{y,0,2-z}]
+Integrate[1,{z,0,2},{x,4*z-z^2,4},{y,0,Sqrt[4-x]}]
Integrate[1,{z,0,2},{y,0,2-z},{x,0,4-y^2}]
Integrate[1,{y,0,2},{z,0,2-y},{x,0,4-y^2}]
Integrate[1,{x,0,4},{y,0,Sqrt[4-x]},{z,0,2-y}]
Integrate[1,{y,0,2},{x,0,4-y^2},{z,0,2-y}]
\end{Verbatim}
\end{tcolorbox}
\vspace{0.7em}
\noindent


El volumen también puede hallarse mediante el siguiente procedimiento.
\vskip0.2cm


%\begin{tcolorbox}
Tomaremos como base del sólido al triángulo que se forma en el plano $yz.$ Un punto sobre el eje $y$ tiene la forma $(t,0) $ y otro punto sobre la recta $z=2-y$ posee la forma $(t,2-t),$ con $t \in [0,2].$
Uniendo estos puntos con un segmento rectilíneo, se paramétrica este triángulo, esto es $(1-u)(t,0) +u(t,2-t)=(t,u(2-t)).$
Como $y=t$, sobre el cilindro se satisface que $x=4-t^2$.

\vskip0.2cm
Con el fin de parametrizar el sólido, se formará la combinación convexa de la forma
\[(1-w) \langle 0,t,u( 2-t)\rangle +w\langle 4-t^{2},t,u( 2-t) \rangle,\]
con lo cual se define la función
${\bf r}( t,u,w) = \langle w( 4-t^{2}),t,2u-tu\rangle ,$
en donde $t\in \lbrack 0,2],$ $u\in \lbrack 0,1],$ $w\in \lbrack 0,1].$
Al derivar parcialmente se obtienen los siguientes vectores
\[{\bf r}_{t} = \langle -2wt,1,-u\rangle,\]
\[{\bf r}_{u} = \langle 0,0,2-t\rangle\]
\[{\bf r}_{w} = \langle 4-t^{2},0,0\rangle.\]
\vskip0.2cm
El triple producto vectorial es $|{\bf r}_{t}\cdot ({\bf r}_{u}\times {\bf r}_{v}) | = ( t+2)( t-2) ^{2},$
integrándolo nos proporciona el volumen del sólido
\[\int_{0}^{1}\int_{0}^{1}\int_{0}^{2}(t+2)(t-2)^{2}dtdudw= \frac{20}{3}.\]
%\end{tcolorbox}

Este cálculo se puede evaluar con las siguientes instrucciones:

\vspace{0.6em}
\noindent
\begin{tcolorbox}[breakable,enhanced,title=\bf Instrucción en \verb"Mathematica", top=2pt,bottom=2pt,boxsep=0pt,before skip=2pt,after skip=0pt]
\begin{Verbatim}[formatcom=\letraverbatim]
r=(1-w)*{0,t,u*(2-t)}+w*{4-t^2,t,u*(2-t)};
Integrate[Dot[D[r,w],Cross[D[r,t],D[r,u]]],{u,0,1},
{t,0,2},{w,0,1}]
\end{Verbatim}
\end{tcolorbox}
\vspace{0.7em}
\noindent


\subsection*{Tres planos y un cilindro parabólico III}
Evaluar {\small $\ds \iiint_{E}6xydV$,} en donde $E$ es el sólido acotado por las figuras de las ecuaciones $z=1+x+y,\,\, y=\sqrt{x},\, x=1,\,\,z=0,\,\, y=0.$ Los siguientes gráficos muestran la región en mención.

\vskip0.2cm
\noindent
\begin{minipage}{\textwidth}
	\centering
	\captionof{figure}{Sólido determinado por $0 \leq z \leq 1 + x + y$, \quad $y^2 \leq x \leq 1$, \quad $0 \leq y \leq 1$.}
	\label{fig:solido_parabolico_inclinado}
	\begin{minipage}{0.34\linewidth}
		\centering
		\vspace{-6mm}
		\includegraphics[width=\linewidth]{Fig/221.pdf}
	\end{minipage}
	\hspace{1em}
	\begin{minipage}{0.16\linewidth}
		\centering
		\vspace{-6mm}
		\includegraphics[width=\linewidth]{Img/9b.pdf}
	\end{minipage}
	\hspace{1em}
	\begin{minipage}{0.16\linewidth}
		\centering
		\vspace{-6mm}
		\includegraphics[width=\linewidth]{Img/9a.pdf}
	\end{minipage}
	\vspace{0.2cm}
	\fuentepropia
\end{minipage}
	


Las siguientes instrucciones permiten obtener las figuras anteriores:

\vspace{0.6em}
\noindent
\begin{tcolorbox}[breakable,enhanced,title=\bf Instrucción en \verb"Mathematica", top=2pt,bottom=2pt,boxsep=0pt,before skip=2pt,after skip=0pt]
\begin{Verbatim}[formatcom=\letraverbatim]
ContourPlot3D[{z==1+x+y,x==y^2,x==1},{x,0,1},{y,0,1},
{z,0,3},Mesh->False,Ticks->{{1},{0,1},{0,1,2}},
ContourStyle->{Orange,Blend[{Yellow,Pink,Green}]}]
RegionPlot3D[0<z<1+x+y&&y^2<x<1&&0<y<1,{x,0,1},
{y,0,1},{z,0,3},PlotStyle->Blend[{Yellow,Orange,Red}],
PlotPoints->80,Mesh->False]
\end{Verbatim}
\end{tcolorbox}
\vspace{0.7em}
\noindent


La intersección del plano y el cilindro parabólico está dada por la ecuación \linebreak
{\small $z=1+x+\sqrt{x}$}, de donde {\small $x=(-\frac{1}{2}+\sqrt{z-3/4})^2$}. Además, la otra posible sustitución nos lleva a la expresión $z=1+y+y^2$, lo cual equivale a
{\small $y=-\frac{1}{2}+\frac{1}{2}\sqrt{4z-3}$}.
\vskip0.2cm

%\vspace*{-1.5\baselineskip} % Sube el bloque sin empalmar el texto superior

\noindent
\begin{minipage}{\textwidth}
\captionsetup{font=footnotesize}
	\centering
	\captionof{figure}{Visualización del sólido y sus proyecciones para el cálculo de volumen.}
	\label{fig:volumen_solido_proyecciones}
	\begin{minipage}{0.25\linewidth}
		\centering
		\vspace{-2mm}
		\includegraphics[width=\linewidth,height=1\linewidth]{Fig/222.pdf}
	\end{minipage}
	\hspace{1em}
	\begin{minipage}{0.25\linewidth}
		\centering
		\vspace{-2mm}
		\includegraphics[width=\linewidth,height=1\linewidth]{Fig/223.pdf}
	\end{minipage}
	\hspace{1em}
	\begin{minipage}{0.25\linewidth}
		\centering
		\vspace{-2mm}
		\includegraphics[width=\linewidth,height=1\linewidth]{Fig/224.pdf}
	\end{minipage}
	\vspace{0.3cm}
	\fuentepropia
\end{minipage}
	
\vskip0.2cm
La proyección en el plano $xy$ es una región acotada por una porción de parábola y dos rectas; además $z$ recorre la región desde el plano $z=0$ hasta el plano $z=1+x+y$.
Con lo anterior, la integral triple está determinada por:
\[\int_{0}^{1}\int_{0}^{\sqrt{x}}\int_{0}^{1+x+y}6xydzdydx=
\int_{0}^{1}\int_{y^2}^{1}\int_{0}^{1+x+y}6xydzdxdy= \frac{65}{28}.\]
\vskip0.2cm
En el caso de plantear las integrales en el orden $dxdydz$ y $dydxdz$ puede ser necesario presentar cuatro integrales triples, los gráficos ilustran el porqué se da esta situación.

{\small
\begin{multline*} \int_{0}^{1}\int_{0}^{1+x}\int_{0}^{\sqrt{x}}6xydydzdx+\int_{0}^{1}\int_{1+x}^{1+x+\sqrt[2]{x}}\int_{z-1-x}^{\sqrt{x}}6xydydzdx \\
\bigintsss_{0}^{1}\bigintsss_{0}^{1}\bigintsss_{0}^{\sqrt{x}}6xydydxdz
+\bigintsss_{1}^2\bigintsss_{\left( -\frac{1}{2}+\sqrt{z-\frac{3}{4}}\right) ^2}^{z-1}\bigintsss_{z-1-x}^{\sqrt{x}}6xydydxdz\\
+\bigintsss_{1}^2\bigintsss_{z-1}^{1}\bigintsss_{0}^{\sqrt{x}}6xydydxdz +\bigintsss_{2}^{3}\bigintsss_{\left( -\frac{1}{2}+\sqrt{z-\frac{3}{4}}\right)
^2}^{1}\bigintsss_{z-1-x}^{\sqrt{x}}6xydydxdz= \frac{65}{28}.
\end{multline*}
}

Integrando con respecto a $x$, las integrales que se deben evaluar son
{\footnotesize
\begin{multline*}
\bigintsss_{0}^{1}\bigintsss_{0}^{1}\bigintsss_{y^2}^{1}6xydxdydz+\bigintsss_{1}^{3}\bigintsss_{-\frac{1}{2}+\frac{1}{2}\sqrt{4z-3}}^{1}\bigintsss_{y^2}^{1}6xydxdydz+ \\ \bigintsss_{1}^{2}\bigintsss_{0}^{-\frac{1}{2}+\frac{1}{2}\sqrt{4z-3}}\bigintsss_{z-1-y}^{1}6xydxdydz+\bigintsss_{2}^{3}\bigintsss_{z-2}^{-\frac{1}{2}+\frac{1}{2}\sqrt{4z-3}}\bigintsss_{z-1-y}^{1}6xydxdydz \\
=\bigintsss_{0}^{1}\bigintsss_{0}^{1+y+y^2}\bigintsss_{y^2}^{1}6xydxdzdy+\bigintsss_{0}^{1}\bigintsss_{1+y+y^2}^{2+y}\bigintsss_{z-1-y}^{1}6xydxdzdy= \frac{65}{28}.
\end{multline*}
}

\vskip0.3cm
Las últimas integrales se evalúan en \verb"Mathematica" mediante la ejecución de la siguiente instrucción:

\vspace{0.6em}
\noindent
\begin{tcolorbox}[breakable,enhanced,title=\bf Instrucción en \verb"Mathematica", top=2pt,bottom=2pt,boxsep=0pt,before skip=2pt,after skip=0pt]
\begin{Verbatim}[formatcom=\letraverbatim]
Integrate[6*x*y,{y,0,1},{z,0,1+y+y^2},{x,y^2,1}]+
Integrate[6*x*y,{y,0,1},{z,1+y+y^2,2+y},{x,z-1-y,1}]
Integrate[6*x*y,{z,0,1},{y,0,1},{x,y^2,1}]+
Integrate[6*x*y,{z,1,3},{y,-1/2+Sqrt[4*z-3]/2,1},
{x,y^2,1}]+Integrate[6*x*y,{z,1,2},{y,0,-1/2+
Sqrt[4*z-3]/2},{x,z-1-y,1}]+Integrate[6*x*y,
{z,2,3},{y,z-2,-1/2+Sqrt[4*z-3]/2},{x,z-1-y,1}]
\end{Verbatim}
\end{tcolorbox}
\vspace{0.7em}
\noindent


Ahora se muestra otra manera de evaluar la integral triple.
\vskip0.2cm
%\begin{tcolorbox}[title=\bfseries Parametrizando el sólido]
\noindent
\begin{minipage}{\textwidth}
	\centering
	\captionof{figure}{Sólido determinado por $0 \leq z \leq 1 + x + y$, \quad $y^2 \leq x \leq 1$, \quad $0 \leq y \leq 1$.}
	\label{fig:solido_parabolico_inclinado_225_227}
	\begin{minipage}{0.29\linewidth}
		\centering
		\vspace{-6mm}
		\includegraphics[width=\linewidth]{Fig/225.pdf}
	\end{minipage}
	\hspace{1em}
	\begin{minipage}{0.29\linewidth}
		\centering
		\vspace{-6mm}
		\includegraphics[width=\linewidth]{Fig/226.pdf}
	\end{minipage}
	\hspace{1em}
	\begin{minipage}{0.29\linewidth}
		\centering
		\vspace{-6mm}
		\includegraphics[width=\linewidth]{Fig/227.pdf}
	\end{minipage}
	
	\vspace{1ex}
	\fuentepropia
\end{minipage}
	
\vskip0.3cm
En el plano $xy,$ con dos puntos $A=(t^2,t)$ y $B=(t^2,0),$ se forma la combinación $(1-u)B+uA=\langle t^2,ut \rangle,$
que parametriza la base del sólido con $t\in [0,1]$ y $u \in [0,1].$
El sólido se parametriza con una combinación convexa de dos puntos: uno sobre el plano $xy$ y otro sobre el plano $z=1+x+y$, esto es,
$\left( 1-w\right) \langle t^{2},tu,0\rangle +w\langle t^{2},tu,1+t^{2}+tu\rangle,$
${\bf r} (t,u,w) =\left\langle t^{2},tu,w\left(1+t^{2}+tu\right) \right\rangle ,$
con $t \in [0,1],$ $u \in [0,1],$ $w\in [0,1].$
Derivando parcialmente se obtienen los siguientes vectores
\begin{eqnarray*}
{\bf r}_t &=& \left\langle 2t,u,w\left( 2t+u\right) \right\rangle \\
{\bf r}_u &=& \left\langle 0,t,wt\right\rangle \\
{\bf r}_w &=& \left\langle 0,0,1+t^{2}+tu\right\rangle
\end{eqnarray*}
El triple producto vectorial es $|{\bf r}_t\cdot({\bf r}_u\times {\bf r}_w)| = | 2t^{2}\left(1+t^{2}+tu\right) |.$
%\begin{eqnarray*}
%|{\bf r}_t \cdot ({\bf r}_t \times {\bf r}_w )| &=& | 2t^{2}\left(1+t^{2}+tu\right) |
%\end{eqnarray*}
Integrando $6xy=6t^{3}u$, con los límites establecidos  se obtiene la integral
\[\ds \iiint_{E}6xydV = \int_{0}^{1}\int_{0}^{1}\int_{0}^{1}\left( 6t^{3}u\right) \left|2t^{2}\left( 1+t^{2}+tu\right) \right| dtdudw= \frac{65}{28}.\]
%\end{tcolorbox}

\subsection*{Tres planos y un cilindro parabólico IV}

Hallar el volumen del sólido acotado por
$z=1-y, \,\, y=\sqrt{x},\,\, x=0,\,\, z=0$.
La figura de este objeto se puede generar con la siguiente instrucción:

\vspace{0.6em}
\noindent
\begin{tcolorbox}[breakable,enhanced,title=\bf Instrucción en \verb"Mathematica", top=2pt,bottom=2pt,boxsep=0pt,before skip=2pt,after skip=0pt]
\begin{Verbatim}[formatcom=\letraverbatim]
RegionPlot3D[0<z<1-y&&0<x<y^2&&0<y<1,{x,0,1},{y,0,1},
{z,0,1},PlotPoints->90,Mesh->False,PlotStyle->Orange]
\end{Verbatim}
\end{tcolorbox}
\vspace{0.7em}
\noindent

% %La figura de dichas expresiones y el sólido se esbozan a continuación. 1
% \begin{center}
% \begin{pspicture}(-2,-0.9)(3,2.2) % \psgrid
% \psset{algebraic,Alpha=60,Beta=15,xunit=1.5,yunit=1.4}
% \pstThreeDSquare[linecolor=black!60,fillcolor=blue!25,fillstyle=solid](-1,1,0)(0,-1,1)(3,0,0)
% %\pstThreeDCoor[xMin=-1.4,xMax=3,yMin=-0.8,yMax=2,zMin=0,zMax=1.7,linecolor=black,IIIDticks,IIIDlabels,spotX=90,spotY=90,spotZ=180]
% \parametricplotThreeD[yPlotpoints=30,drawStyle=xyLines,linewidth=0.3pt,linecolor=gray](-0.5,1.3)(-0,1.2){t^2|t|u}
% \parametricplotThreeD[yPlotpoints=30,drawStyle=xyLines,linewidth=0.3pt,linecolor=gray](-0,1.2)(-0.5,1.3){u^2|u|t}
% \parametricplotThreeD[linewidth=1.2pt,linecolor=black](0,1){t^2|t|1-t}
% \parametricplotThreeD[linewidth=1.2pt,linecolor=black](0,1){t^2|t|0}
% % \pstThreeDLine[linewidth=1.2pt,linecolor=black](0,1,0)(0,0,0)(0,0,1)(0,1,0)(1,1,0)
% \pstThreeDLine[linewidth=0.3pt,linecolor=blue,linestyle=dashed](2,0,1)(2,0,0)(2,1,0)
% \uput[0](0.2,1.35){$z=1-y$}
% \uput[0](-1.2,1.2){$x=y^2$}
% \parametricplotThreeD[linewidth=0.6pt,linecolor=black,linestyle=dashed](0,1){(t-1)^2|0|t}
% \pstThreeDLine[linewidth=0.6pt,linecolor=black,linestyle=dashed](1,0,0)(1,1,0)
% \end{pspicture}

% \begin{pspicture}(-1.5,-0.8)(2,2.2) % \psgrid
% \psset{algebraic,Alpha=40,Beta=15,yunit=1.75,xunit=1.5}
% \pstThreeDCoor[spotX=-90,xMin=-0.8,xMax=1.3,yMin=-0.4,yMax=2,zMin=0,zMax=1.3,spotZ=180,IIIDticks,%
% IIIDlabels]
% \pscustom[fillstyle=crosshatch,hatchwidth=0.2pt,hatchcolor=blue,opacity=0.4]{
% \parametricplotThreeD(0,1){t^2|t|1-t}
% \pstThreeDLine(1,1,0)(0,1,0)(0,0,1)
% }
% \pscustom[fillstyle=solid,fillcolor=lightgray!70,opacity=0.5]{
% \parametricplotThreeD(0,1){t^2|t|1-t}
% \parametricplotThreeD(1,0){t^2|t|0}
% \pstThreeDLine(0,0,0)(0,0,1)
% }
% \psset{linewidth=0.7pt}
% \pstThreeDLine(0,0,1)(0,1,0)(1,1,0)
% \pstThreeDLine[linestyle=dashed](1,1,0)(1,0,0)
% \parametricplotThreeD(0,1){t^2|t|1-t}
% \parametricplotThreeD[linestyle=dashed](0,1){t^2|t|0}
% \pstThreeDPut(0,-0.2,1){$1$}
% \pstThreeDPut(0.8,-0.3,0){$1$}
% \end{pspicture}
% \includegraphics[height=4.5cm]{Img/62a.pdf}
% \captionof{figure}{\small Sólido determinado por $0\leq z \leq 1-y,$ $0\leq x\leq y^2,$ $0\leq y\leq 1$}
% \end{center}

% \begin{center}
% \begin{pspicture}(-1.5,-0.8)(2,2.2) % \psgrid
% \psset{algebraic,Alpha=40,Beta=15,yunit=1.75,xunit=1.5}
% \pstThreeDCoor[spotX=-90,xMin=-0.8,xMax=1.3,yMin=-0.4,yMax=2,zMin=0,zMax=1.3,spotZ=180,IIIDticks,%
% IIIDlabels]
% \pscustom[fillstyle=crosshatch,hatchwidth=0.2pt,hatchcolor=blue,opacity=0.4]{
% \parametricplotThreeD(0,1){t^2|t|1-t}
% \pstThreeDLine(1,1,0)(0,1,0)(0,0,1)
% }
% \pscustom[fillstyle=solid,fillcolor=lightgray!70,opacity=0.5]{
% \parametricplotThreeD(0,1){t^2|t|1-t}
% \parametricplotThreeD(1,0){t^2|t|0}
% \pstThreeDLine(0,0,0)(0,0,1)
% }
% \psset{linewidth=0.7pt}
% \pstThreeDLine(0,0,1)(0,1,0)(1,1,0)
% \pstThreeDLine[linestyle=dashed](1,1,0)(1,0,0)
% \parametricplotThreeD(0,1){t^2|t|1-t}
% \parametricplotThreeD[linestyle=dashed](0,1){t^2|t|0}
% \pstThreeDPut(0,-0.2,1){$1$}
% \pstThreeDPut(0.8,-0.3,0){$1$}
% \end{pspicture}
% \includegraphics[height=4.5cm]{Img/62a.pdf}
% \captionof{figure}{Sólido determinado por $0\leq z \leq 1-y,$ $0\leq x\leq y^2,$ $0\leq y\leq 1.$}
% \end{center}
%\vspace{0.2cm}

\noindent
\begin{minipage}{\textwidth}
	\centering
	\captionof{figure}{Sólido determinado por $0\leq z \leq 1-y,$ $0\leq x\leq y^2,$ $0\leq y\leq 1.$}
	\begin{minipage}{0.30\linewidth}
		\centering
		\vspace{-15mm}
		\includegraphics[width=\linewidth]{Fig/303.pdf}
	\end{minipage}
	\hspace{0.2em}
	\begin{minipage}{0.30\linewidth}
		\centering
		\vspace{-15mm}
		\includegraphics[width=\linewidth]{Fig/302.pdf}
	\end{minipage}
	\hspace{0.2em}
	\begin{minipage}{0.20\linewidth}
		\centering
		\vspace{-15mm}
		\includegraphics[width=\linewidth]{Imgn/62a.pdf}
	\end{minipage}
	
	\vspace{-8mm}
	
	%\vspace{0.3em}
	\fuentepropia
\end{minipage}
\vskip0.2cm

En el plano $xz$, reemplazando $y=\sqrt{x}$ en la ecuación del plano $y=1-z,$
se obtiene $z=1-\sqrt{x},$ lo cual equivale a $x=(1-z)^2.$
En el plano $xy$ los límites se determinan con las expresiones $x=y^2,$ $y=1$ y $x=0$.
En el plano $yz$ la región proyectada es un triángulo acotado por los ejes coordenados y la recta $z=1-y$.
\vskip0.2cm

\noindent
\begin{minipage}{\textwidth}
	\centering
	\captionof{figure}{Proyecciones en los planos coordenados.}
	\begin{minipage}{0.23\linewidth}
		\centering
		\vspace{-10mm}
		\includegraphics[width=\linewidth]{Fig/304.pdf}
	\end{minipage}
	\hspace{0.2em}
	\begin{minipage}{0.23\linewidth}
		\centering
		\vspace{-10mm}
		\includegraphics[width=\linewidth]{Fig/305.pdf}
	\end{minipage}
	\hspace{0.2em}
	\begin{minipage}{0.23\linewidth}
		\centering
		\vspace{-10mm}
		\includegraphics[width=\linewidth]{Fig/306.pdf}
	\end{minipage}
	
	%\vspace{0.3em}
	
	\vspace{-5mm}
	\fuentepropia
\end{minipage}


Estos elementos ayudan a plantear las integrales que determinan el volumen del sólido
\begin{multline*}
\int_{0}^{1}\int_{\sqrt{x}}^{1}\int_{0}^{1-y}dzdydx=
\int_{0}^{1}\int_{0}^{y^2}\int_{0}^{1-y}dzdxdy \\
= \int_{0}^{1}\int_{0}^{1-y}\int_{0}^{y^2}dxdzdy= \int_{0}^{1}\int_{0}^{1-z}\int_{0}^{y^2}dxdydz \\
= \int_{0}^{1}\int_{0}^{\left( 1-z\right)^2}\int_{\sqrt{x}}^{1-z}dydxdz=
\int_{0}^{1}\int_{0}^{1-\sqrt{x}}\int_{\sqrt{x}}^{1-z}dydzdx= \frac{1}{12}.
\end{multline*}
Enseguida se presentan las indicaciones para evaluar estas integrales.

\vspace{0.6em}
\noindent
\begin{tcolorbox}[breakable,enhanced,title=\bf Instrucción en \verb"Mathematica", top=2pt,bottom=2pt,boxsep=0pt,before skip=2pt,after skip=0pt]
\begin{Verbatim}[formatcom=\letraverbatim]
Integrate[1,{z,0,1},{x,0,(1-z)^2},{y,Sqrt[x],1-z}]
Integrate[1,{x,0,1},{y,Sqrt[x],1},{z,0,1-y}]
Integrate[1,{y,0,1},{z,0,1-y},{x,0,y^2}]
\end{Verbatim}
\end{tcolorbox}
\vspace{0.7em}
\noindent

El mismo valor se puede conseguir siguiendo el siguiente proceso.


%\begin{tcolorbox}[enhanced,colback=white,colframe=black,title=\bfseries Parametrizando el sólido,fonttitle=\bfseries,boxrule=0.8pt,arc=2pt,left=4pt, right=4pt,top=4pt, bottom=4pt]
		
%\vspace*{-0.2em}
Sobre la región de integración del plano $xy$ con dos puntos $P(t^2,t)$ y 
$Q(t^2,1)$ se forma una combinación convexa $(1-u)P+uQ$, que describe la región sombreada. 
Simplificando, se obtiene $\langle t^2,\,u+(1-u)t\rangle$, con 
$t\in[0,1]$, $u\in[0,1].$
	
\vskip0.2cm
\noindent
\begin{minipage}{\textwidth}
	\centering
	\captionof{figure}{\small  $0 \le x \le y^2,\ 0 \le x \le 1.$}
	\begin{minipage}{0.36\linewidth}
		\centering
		\vspace{-18mm}
		\includegraphics[width=\linewidth]{Fig/307.pdf}
	\end{minipage}
	
	\vspace{-12mm}
	\fuentepropia
\end{minipage}
\vskip0.2cm
	
\vspace{0.5em}
\noindent
Sobre el plano $z=1-y$ se satisface que la coordenada $z$ tiene la forma 
$z=(1-t)(1-u).$ El sólido se parametriza mediante una combinación convexa de dos puntos de la forma 
$\langle t^2,\,u+(1-u)t,\,0\rangle$ y $\langle t^2,\,u+(1-u)t,\,(1-t)(1-u)\rangle$, 
lo cual permite definir la función
\[
\mathbf{r}(t,u,w)=\langle t^2,\,t+(1-u)t,\,(1-t)(u-1)w\rangle,
\]
con $t,u,w\in[0,1].$ Con esta función se obtienen los siguientes vectores:

\vspace{-8mm}
\begin{align*}
\mathbf{r}_t &= \langle 2t,\,1-u,\,-(1-u)w\rangle, &
\mathbf{r}_w &= \langle 0,\,0,\,(1-t)(1-u)\rangle,\\
\mathbf{r}_u &= \langle 0,\,1-t,\,-(1-t)w\rangle, &
\mathbf{r}_u\times\mathbf{r}_t &= \langle 0,\,-2tw(1-t),\,-2t(1-t)\rangle.
\end{align*}
	
\noindent
Con lo cual 
\[
|\mathbf{r}_w\cdot(\mathbf{r}_u\times\mathbf{r}_t)|
=|2t(1-t)^2(1-u)|,
\]
y su integración en los límites dados proporciona el volumen del sólido:
{\small
\[
\int_0^1\!\!\int_0^1\!\!\int_0^1
|2t(1-t)^2(1-u)|\,dt\,du\,dw
=\tfrac{1}{12}
\]
}	
%\end{tcolorbox}


Este procedimiento se realiza con la ayuda de las siguientes instrucciones:

\vspace{0.6em}
\noindent
\begin{tcolorbox}[breakable,enhanced,title=\bf Instrucción en \verb"Mathematica", top=2pt,bottom=2pt,boxsep=0pt,before skip=2pt,after skip=0pt]
\begin{Verbatim}[formatcom=\letraverbatim]
{x,y}=(1-u)*{t^2,t}+u*{t^2,1};
z=1-y;
R=Factor[(1-w)*{x,y,z}+w*{x,y,0}];
Integrate[Dot[D[R,w],Cross[D[R,u],D[R,t]]],
{w,0,1},{u,0,1},{t,0,1}]
\end{Verbatim}
\end{tcolorbox}
\vspace{0.7em}
\noindent

\subsection*{Un cilindro y dos planos III}
La ecuación $x^{2}+z^{2}=1$ determina un cilindro de radio $1$ con eje de simetría el eje $y$, esta superficie es cortada por los planos con ecuaciones $y=x$ y $y=2x.$

\vskip0.2cm
\noindent
\begin{minipage}{\textwidth}
	\centering
	\captionof{figure}{Superficies $x^{2} + z^{2} = 1$, $y = x$ y $y = 2x$, con $y \geq 0$.}
	\label{fig:cilindro_planos_ypositivo}
	\begin{minipage}{0.17\linewidth}
		\centering
		\includegraphics[width=\linewidth]{Img/27b.pdf}
	\end{minipage}
	\hspace{1em}
	\begin{minipage}{0.17\linewidth}
		\centering
		\includegraphics[width=\linewidth]{Img/27a.pdf}
	\end{minipage}
	\hspace{1em}
	\begin{minipage}{0.17\linewidth}
		\centering
		\includegraphics[width=\linewidth]{Img/27c.pdf}
	\end{minipage}
	\hspace{1em}
	\begin{minipage}{0.17\linewidth}
		\centering
		\includegraphics[width=\linewidth]{Img/27d.pdf}
	\end{minipage}
	
	\vspace{0.5em}
	\fuentepropia
\end{minipage}
\vskip0.2cm

Con las siguientes instrucciones se obtienen las figuras anteriores:

\vspace{0.6em}
\noindent
\begin{tcolorbox}[breakable,enhanced,title=\bf Instrucción en \verb"Mathematica", top=2pt,bottom=2pt,boxsep=0pt,before skip=2pt,after skip=0pt]
\begin{Verbatim}[formatcom=\letraverbatim]
ContourPlot3D[{1==x^2+z^2,y==x,y==2*x},{x,0,1},
{y,0,2},{z,-1,1},PlotPoints->50,Mesh->False,
ContourStyle->{Orange,Cyan,LightRed},
Ticks->{{0,1},{0,1,2},{-1,0}}]
RegionPlot3D[-Sqrt[1-x^2]<z<Sqrt[1-x^2]&&x<y
<2*x&&0<x<1,{x,0,1},{y,0,2},{z,-1,1},PlotPoints->90,
Mesh->False,Ticks->{{0,1},{0,1,2},{-1,0,1}}]
\end{Verbatim}
\end{tcolorbox}
\vspace{0.7em}
\noindent

Con el fin de plantear las integrales triples que permiten obtener el volumen del sólido presentamos las siguientes figuras.

\vskip0.2cm
\noindent
\begin{minipage}{\textwidth}
	\centering
	\captionof{figure}{Proyección del sólido en los planos coordenados.}
	\begin{minipage}{0.30\linewidth}
		\centering
		\vspace{-9mm}		\includegraphics[width=\linewidth]{Fig/308.pdf}
	\end{minipage}
	\hspace{1em}
	\begin{minipage}{0.30\linewidth}
		\centering
		\vspace{-9mm}
		\includegraphics[width=\linewidth]{Fig/309.pdf}
	\end{minipage}
	\hspace{1em}
	\begin{minipage}{0.30\linewidth}
		\centering
		\vspace{-9mm}
		\includegraphics[width=\linewidth]{Fig/310.pdf}
	\end{minipage}
	
	\vspace{-8mm}
	\parbox{0.95\linewidth}{\centering\fontsize{9pt}{10pt}\selectfont\textbf{Fuente:} elaboración propia}
\end{minipage}
\vskip0.2cm

% \begin{figure}[H]......2
% \centering

% \vspace{0.4cm}

% % Primer gráfico: plano xz
% \begin{pspicture}(-0.6,-1.4)(2.5,2.1)
% \psset{algebraic}
% \pscustom[fillstyle=vlines,hatchcolor=gray!40,hatchsep=1.4pt,hatchwidth=0.7pt,linestyle=none]{%
%   \parametricplot{-1.57}{1.57}{cos(t),sin(t)}%
% }
% \psaxes[linecolor=black]{->}(0,0)(-0.6,-1.5)(2,2.1)[$x$,-90][$z$,0]
% \parametricplot[linecolor=red!80!black,arrows=->]{-1.6}{1.9}{cos(t),sin(t)}
% \psline[linecolor=red!80!black,arrows=-](0,1)(0,-1)
% \rput(1.6,1.3){$x^2+z^2=1$}
% \end{pspicture}

% \vspace{0.5cm}

% % Segundo gráfico: plano yz
% \begin{pspicture}(-1.2,-1.4)(3.2,2.1)
% \psset{algebraic}
% \pscustom[fillstyle=vlines,hatchcolor=gray!40,hatchsep=1.4pt,hatchwidth=0.7pt,linestyle=none]{%
%   \parametricplot{-1.57}{1.57}{cos(t),sin(t)}%
% }
% \pscustom[fillstyle=solid,fillcolor=gray!40,linestyle=none]{%
%   \parametricplot{-1.57}{1.57}{cos(t),sin(t)}%
%   \parametricplot{1.57}{1.05}{2*cos(t),sin(t)}%
%   \psline(1,0)(1,0.87)%
%   \parametricplot{-1.05}{-1.57}{2*cos(t),sin(t)}%
% }
% \pscustom[fillstyle=hlines,hatchcolor=gray!40,linestyle=none,hatchsep=1.4pt,hatchwidth=0.7pt]{%
%   \parametricplot{-1.05}{1.05}{2*cos(t),sin(t)}%
% }
% \psaxes[linecolor=black]{->}(0,0)(-0.6,-1.2)(3,1.75)[$y$,-90][$z$,0]
% \parametricplot[linecolor=red!80!black,arrows=-]{-1.57}{1.57}{cos(t),sin(t)}
% \parametricplot[linecolor=red!80!black,arrows=->]{-1.57}{1.7}{2*cos(t),sin(t)}
% \psline[linecolor=red!80!black,arrows=-](0,1)(0,-1)
% \psline[linecolor=red!80!black,arrows=-,linewidth=0.6pt,linestyle=dashed](1,-0.87)(1,0.87)
% \psline[linecolor=black,arrows=->](0.65,-0.8)(1,-1.2)
% \rput(2,-1.2){$y^2+z^2=1$}
% \rput(1.6,1.3){$y^2+4z^2=4$}
% \end{pspicture}

% \vspace{0.5cm}

% % Tercer gráfico: plano xy
% \begin{pspicture}(-1,-0.6)(2,2.8)
% \psset{algebraic,xunit=1.2}
% \psaxes[linecolor=black,linewidth=0.6pt]{->}(0,0)(-0.5,-0.5)(2,2.5)[$x$,90][$y$,0]
% \psline[fillstyle=vlines,hatchcolor=gray!50,hatchsep=1pt](0,0)(1,1)(0.5,1)(0,0)
% \psline[fillstyle=hlines,hatchcolor=gray!50,hatchsep=1pt](1,1)(1,2)(0.5,1)
% \psline[linecolor=red!80!black,arrows=-,linewidth=0.6pt,linestyle=dashed](1,1)(0.5,1)
% \parametricplot[linecolor=red!80!black,arrows=->]{-0.25}{1.2}{t,t}
% \parametricplot[linecolor=red!80!black,arrows=->]{-0.25}{1.2}{t,2*t}
% \parametricplot[linecolor=red!80!black,arrows=-]{1}{2}{1,t}
% \rput(1.8,1.2){$y=x$}
% \rput(1.8,2.2){$y=2x$}
% \end{pspicture}

% \caption{Proyección del sólido en los planos coordenados.}
% \end{figure}

Con esta representación determinamos las integrales que nos aportan el volumen del sólido. Al integrar primero con respecto a $y$, esta recorre puntos desde el plano $y=x$ hasta el plano $y=2z$. Los límites para $x$ y $z$ se encuentran en la semicircunferencia unitaria en el plano $xz$.
\vskip0.2cm
\[\int_{-1}^{1}\int_{0}^{\sqrt{1-z^{2}}}\int_{x}^{2x}dydxdz= \int_{0}^{1}\int_{-\sqrt{1-x^{2}}}^{\sqrt{1-x^{2}}}\int_{x}^{2x}dydzdx= \frac{2}{3}\]
\vskip0.2cm
Al integrar con respecto a $z$, los planos $y=x$ y $y=2x$ son paralelos al eje $z$ por lo que los límites para $z$ van desde la parte negativa del cilindro a la parte positiva de este. Por tanto, la integral en este caso es:
\begin{multline*} \int_{0}^{1}\int_{x}^{2x}\int_{-\sqrt{1-x^{2}}}^{\sqrt{1-x^{2}}}dzdydx= \int_{0}^{1}\int_{\frac{y}{2}}^{y}\int_{-\sqrt{1-x^{2}}}^{\sqrt{1-x^{2}}}dzdxdy \\ +\int_{1}^2\int_{\frac{y}{2}}^{1}\int_{-\sqrt{1-x^{2}}}^{\sqrt{1-x^{2}}}dzdxdy= \frac{2}{3}. 
\end{multline*}

Integrando con respecto a $x$, la región se divide en varias partes, según como se observe. La proyección en el plano $yz$ ayuda a visualizar estos elementos:
{\small
\[\bigintsss_{-1}^{1}\bigintsss_{0}^{\sqrt{1-z^{2}}}\bigintsss_{\frac{y}{2}}^{y}dxdydz+\bigintsss_{-1}^{1}\bigintsss_{\sqrt{1-z^{2}}}^{\sqrt{4-4z^{2}}}\bigintsss_{\frac{y}{2}}^{\sqrt{1-z^{2}}}dxdydz= \frac{2}{3}.\]
}

En el orden $dxdzdy$ la evaluación respectiva es:

\vspace{-8mm}
{\small
\begin{multline*}
\bigintsss_{0}^{1}\bigintsss_{-\sqrt{1-y^{2}}}^{\sqrt{1-y^{2}}}\bigintsss_{\frac{y}{2}}^{y}dxdzdy+\bigintsss_{0}^{1}\bigintsss_{\sqrt{1-y^{2}}}^{\sqrt{1-\frac{y^{2}}{4}}}\bigintsss_{\frac{y}{2}}^{\sqrt{1-z^{2}}}dxdzdy +\\
\bigintsss_{0}^{1}\bigintsss_{-\sqrt{1-\frac{y^{2}}{4}}}^{-\sqrt{1-y^{2}}}\bigintsss_{\frac{y}{2}}^{\sqrt{1-z^{2}}}dxdzdy+\bigintsss_{1}^2\bigintsss_{-\sqrt{1-\frac{y^{2}}{4}}}^{\sqrt{1-\frac{y^{2}}{4}}}\bigintsss_{\frac{y}{2}}^{\sqrt{1-z^{2}}}dxdzdy= \frac{2}{3}.
\end{multline*}
}

Las integrales anteriores en los órdenes $dydxdz,$ $dzdxdy$ y $dxdzdy$, se evalúan en \verb"Mathematica"
ejecutando las siguientes instrucciones:

\vspace{0.6em}
\noindent
\begin{tcolorbox}[breakable,enhanced,title=\bf Instrucción en \verb"Mathematica", top=2pt,bottom=2pt,boxsep=0pt,before skip=2pt,after skip=0pt]
\begin{Verbatim}[formatcom=\letraverbatim]
Integrate[1,{z,-1,1},{x,0,Sqrt[1-z^2]},{y,x,2x}]
Integrate[1,{y,0,1},{x,y/2,y},{z,-Sqrt[1-x^2],
Sqrt[1-x^2]}]+Integrate[1,{y,1,2},{x,y/2,1},
{z,-Sqrt[1-x^2],Sqrt[1-x^2]}]
Integrate[1,{y,0,1},{z,-Sqrt[1-y^2],Sqrt[1-y^2]},
{x,y/2,y}]+Integrate[1,{y,0,1},{z,Sqrt[1-y^2],
Sqrt[1-y^2/4]},{x,y/2,Sqrt[1-z^2]}]+Integrate[1,
{y,0,1},{z,-Sqrt[1-y^2/4],-Sqrt[1-y^2]},{x,y/2,
Sqrt[1-z^2]}]+Integrate[1,{y,1,2},{z,-Sqrt[1-y^2/4],
Sqrt[1-y^2/4]},{x,y/2,Sqrt[1-z^2]}]
\end{Verbatim}
\end{tcolorbox}
\vspace{0.7em}
\noindent

También se puede hacer un cambio a coordenadas polares, $x=r\cos \theta $ y $z=r\sin \theta $,
de aquí se sigue lo siguiente:
{\small
\[\int_{-\pi /2}^{\pi /2}\int_{0}^{1}\int_{r\cos \theta }^{2r\cos \theta}rdydrd\theta = \frac{2}{3}.\]
}

\vspace{-2mm}
Ahora se ilustra otra posibilidad de obtener el mismo cálculo.

\vskip0.2cm
%\begin{tcolorbox}[enhanced,colback=white,colframe=black,title=\bfseries Parametrizando el sólido,fonttitle=\bfseries,boxrule=0.8pt,arc=2pt,left=4pt, right=4pt,top=4pt,bottom=4pt]


En el plano $xz$, la proyección del sólido es un semicírculo que se parametriza como:
{\small
\[ \la u\sin t, 0, u\cos t \ra, \quad t\in [-\pi /2, \pi/2], \quad u\in [0,1]. \]
}
Con un punto $P(u\sin t, u\sin t, u\cos t)$ sobre el plano $y=x$, y otro punto $Q(u\sin t, 2u\sin t, u\cos t)$ sobre el plano $y=2x.$

\vskip0.2cm
\noindent
\begin{minipage}{\textwidth}
	\centering
	\captionof{figure}{\small ${\bf r}(u,t)=\left\langle u\sin t,u\cos t\right\rangle$}
	\begin{minipage}{0.34\linewidth}
		\centering
		\vspace{-16mm}
		\includegraphics[width=\linewidth]{Fig/311.pdf}
	\end{minipage}
	
	\vspace{-10mm}
	\fuentepropia
\end{minipage}
\vskip0.2cm


\noindent
Se define una parametrización del sólido, que se obtiene mediante una combinación convexa
$\left( 1-v\right) \left\langle u\sin t,2u\sin t,u\cos t\right\rangle +v\left\langle u\sin t,u\sin t,u\cos t\right\rangle,$
entre dos puntos $P$ y $Q.$
$ {\bf r}(t,u,v) = \left\langle u\sin t,u\left( 2-v\right) \sin t,u\cos t\right\rangle ,$
en donde $t\in [ -\pi /2,\pi/2]$, $u\in [ 0,1]$ y $v\in [0,1].$

\vskip0.2cm
Esta función nos permite hallar los siguientes vectores:
%\vspace{-2mm}
{\small
\begin{eqnarray*}
{\bf r}_{t} &=& \left\langle u\cos t,u\left(2-v\right) \cos t,-u\sin t\right\rangle \\
{\bf r}_{u} &=& \left\langle \sin t,\left( 2-v\right) \sin t,\cos t\right\rangle \\
{\bf r}_{v} &=& \left\langle 0,-u\sin t,0\right\rangle.
\end{eqnarray*}
}

%\vspace{-2mm}
El producto vectorial triple esta dado por $\left|{\bf r}_{t}\cdot \left({\bf r}_{u}\times {\bf r}_{v}\right) \right| =\left| u^{2}\sin t\right|,$
que se integra con los límites establecidos y nos proporciona el volumen del sólido en mención
{\small
\[\int_{0}^{1}\int_{0}^{1}\int_{-\pi /2}^{\pi /2}\left| u^{2}\sin t\right|dtdudv= \frac{2}{3}.\]
}
%\end{tcolorbox}
%\vskip0.2cm
%\vspace{-2mm}
El cálculo anterior se obtiene con las siguientes instrucciones:

\vspace{0.6em}
\noindent
\begin{tcolorbox}[breakable,enhanced,title=\bf Instrucción en \verb"Mathematica", top=2pt,bottom=2pt,boxsep=0pt,before skip=2pt,after skip=0pt]
\begin{Verbatim}[formatcom=\letraverbatim]
{x,y,z}={u*Sin[t],u*Sin[t],u*Cos[t]};
r=(1-v)*{x,y,z}+v*{x,2*y,z};
A=Abs[Dot[D[r,v],Cross[D[r,u],D[r,t]]]];
Integrate[A,{u,0,1},{t,-Pi/2,Pi/2},{v,0,1}]
\end{Verbatim}
\end{tcolorbox}
\vspace{0.7em}
\noindent


\subsection*{Un paraboloide y cuatro planos} \index{Paraboloide}
El paraboloide con ecuación  $z = 4 - x^2 - y^2$ y el plano $x = 2-2y$ se intersecan en una región del espacio que a continuación se muestra.

\vskip0.2cm
\noindent
\begin{minipage}{\textwidth}
	\centering
	\captionof{figure}{Sólido definido por $0 \leq z \leq 4 - x^2 - y^2$, \quad $0 \leq x \leq 2 - 2y$, \quad $0 \leq y \leq 1$.}
	\label{fig:solido_parabolico_inclinado_44}
	\begin{minipage}{0.17\linewidth}
		\centering
		\includegraphics[width=\linewidth]{Img/44d.pdf}
	\end{minipage}
	\hspace{1em}
	\begin{minipage}{0.17\linewidth}
		\centering
		\includegraphics[width=\linewidth]{Img/44a.pdf}
	\end{minipage}
	\hspace{1em}
	\begin{minipage}{0.17\linewidth}
		\centering
		\includegraphics[width=\linewidth]{Img/44c.pdf}
	\end{minipage}
	
	\vspace{0.5em}
	\fuentepropia
\end{minipage}
\vskip0.2cm

Estas superficies determinan un sólido en el primer octante.
Las anteriores figuras se obtienen ejecutando las siguientes instrucciones:

%\vskip0.2cm
\vspace{0.6em}
\noindent
\begin{tcolorbox}[breakable,enhanced,title=\bf Instrucción en \verb"Mathematica", top=2pt,bottom=2pt,boxsep=0pt,before skip=2pt,after skip=0pt]
\begin{Verbatim}[formatcom=\letraverbatim]
ContourPlot3D[{z==4-x^2-y^2,x==2-2*y,y==0},{x,0,2},
{y,0,2},{z,0,4},Ticks->{{0,2},{2,4},{0,2,4}},
MeshStyle->Gray]
RegionPlot3D[0<z<4-x^2-y^2&&0<x<2-2*y&&0<y<1,{x,0,2},
{y,0,1},{z,0,4},PlotPoints->90,Mesh->False,
PlotStyle->RGBColor[0.7,0.2,0.1]]
ContourPlot3D[{z==4-x^2-y^2,x==2-2*y},{x,-2,2},{y,-2,
2},{z,0,4},ContourStyle->{Directive[Opacity[0.8]],
Directive[Cyan,Opacity[0.7]]},Mesh->False]
\end{Verbatim}
\end{tcolorbox}
\vspace{0.7em}
\noindent

Plantearemos ahora las integrales triples que se requieren evaluar para determinar el volumen de esta región. Para determinar la proyección del sólido en los planos coordenados se debe eliminar una de las variables adecuadamente escogidas. En el plano $xy,$ (que equivale a $z=0$),
la proyección está dada por las expresiones $x = 2 -2y\,$ y $\,x^2 + y^2=4$.
\vskip0.3cm

En el plano $yz,$ la expresión que determina la región de integración se obtiene eliminando $x$ de la ecuación del paraboloide, esto es 
$z =4-(2-2y)^2-y^2$ por lo tanto
$z = 8y-5y^2,$ entonces $y = \frac{4}{5}\pm \frac{1}{5}\sqrt{16-5z}.$


\vskip0.3cm
\noindent
\begin{minipage}{\textwidth}
	\centering
	\captionof{figure}{Proyección del sólido en dos planos coordenados}
	\begin{minipage}{0.37\linewidth}
		\centering
		\vspace{-19mm}
		\includegraphics[width=\linewidth]{Fig/312.pdf}
	\end{minipage}
	\hspace{1em} 
	\begin{minipage}{0.37\linewidth}
		\centering
		\vspace{-19mm}
		\includegraphics[width=\linewidth]{Fig/313.pdf}
	\end{minipage}
	
	\vspace{-14mm}
	\fuentepropia
\end{minipage}
\vskip0.3cm



Al integrar primera con respecto a $z,$ esta toma valores desde el plano $xy$ hasta el paraboloide y los límites de integración con respecto a $x$ e $y$ se pueden obtener del primero de los gráficos anteriores. Mientras que al integrar primero con respecto,
a $x,$ esta toma valores desde el plano $yz$ hasta el paraboloide, en un caso, y desde el plano $yz$ hasta el plano $x=2-2y$ en el otro caso. Por lo tanto, el volumen está dado por las siguientes integrales.

\vspace{-7mm}
{\small
\begin{multline*}
\bigintsss_{0}^{1}\bigintsss_{0}^{2-2y}\bigintsss_{0}^{4-x^2-y^2}dzdxdy= \bigintsss_{0}^2\bigintsss_{0}^{1-x/2}\bigintsss_{0}^{4-x^2-y^2}dzdydx \\
=\bigintsss_{0}^{1}\bigintsss_{8y-5y^2}^{4-y^2}\bigintsss_{0}^{\sqrt{4-z-y^2}}dxdzdy+\bigintsss_{0}^{1}\bigintsss_{0}^{8y-5y^2}\bigintsss_{0}^{2-2y}dxdzdy = \frac{19}{6}.
\end{multline*}
}

La integral en el orden $dxdydz$ requiere un poco de observación adicional.

La expresión $s = 8y - 5y^2$ determina una parábola en el plano $yz$, su valor máximo se puede hallar derivando, esto es
$z^\prime =8-10y=0,$ entonces $y=\frac{4}{5},$ por tanto el valor máximo de $z$ es.
$8 \left( \frac{4}{5} \right)- 5\left( \frac{4}{5} \right)^2 = \frac{16}{5}.$
\vskip0.3cm
Esta condición hace que la proyección en el plano $yz$ se divida en varias regiones que la siguiente figura ilustra.


\vskip0.2cm
\noindent
\begin{minipage}{\textwidth}
	\centering
	\captionof{figure}{Proyección en el plano $yz$.}
	\begin{minipage}{0.38\linewidth}
		\centering
		\vspace{-15mm}
		\includegraphics[width=\linewidth]{Fig/314.pdf}
	\end{minipage}
	
	\vspace{-10mm}
	\fuentepropia
\end{minipage}
%\vskip0.3cm


De la expresión $z=8y-5y^2$ se sigue que
$y=\frac{4 \pm \sqrt{16-5z}}{5}$ y la expresión $z= 4-y^2$
equivale $y=\sqrt{4-z},$ con los elementos anteriormente descritos, el volumen del sólido está dado por:
%\vspace{-2mm}
{\footnotesize
\begin{multline*}
\bigintsss_{0}^{\frac{16}{5}}\bigintsss_{0}^{\frac{4}{5}-\frac{1}{5}\sqrt{16-5z}}\bigintsss_{0}^{\sqrt{4-z-y^2}}dxdydz
+\bigintsss_{\frac{16}{5}}^{4}\bigintsss_{0}^{\sqrt{4-z}}\bigintsss_{0}^{\sqrt{4-z-y^2}}dxdydz \\
+\bigintsss_{3}^{\frac{16}{5}}\bigintsss_{\frac{4}{5}+\frac{1}{5}\sqrt{16-5z}}^{\sqrt{4-z}}\bigintsss_{0}^{\sqrt{4-z-y^2}}dxdydz +\bigintsss_{0}^{3}\bigintsss_{\frac{4}{5}-\frac{1}{5}\sqrt{16-5z}}^{1}\bigintsss_{0}^{2-2y}dxdydz\\
+\bigintsss_{3}^{\frac{16}{5}}\bigintsss_{\frac{4}{5}-\frac{1}{5}\sqrt{16-5z}}^{\frac{4}{5}+\frac{1}{5}\sqrt{16-5z}}\bigintsss_{0}^{2-2y}dxdydz= \frac{19}{6}
\end{multline*}
}

En \verb"Mathematica" estas integrales se evalúan con la siguiente instrucción:

\vspace{0.6em}
\noindent
\begin{tcolorbox}[breakable,enhanced,title=\bf Instrucción en \verb"Mathematica", top=2pt,bottom=2pt,boxsep=0pt,before skip=2pt,after skip=0pt]
\begin{Verbatim}[formatcom=\letraverbatim]
Integrate[1,{z,0,16/5},{y,0,(4-Sqrt[16-5*z])/5},
{x,0,Sqrt[4-z-y^2]}]+Integrate[1,{z,3,16/5},
{y,(4+Sqrt[16-5*z])/5,Sqrt[4-z]},{x,0,Sqrt[4-z-y^2]}]
+Integrate[1,{z,16/5,4},{y,0,Sqrt[4-z]},
{x,0,Sqrt[4-z-y^2]}]+Integrate[1,{z,0,3},
{y,(4-Sqrt[16-5*z])/5,1},{x,0,2-2*y}]
+Integrate[1,{z,3,16/5},{y,(4-Sqrt[16-5*z])/5,
(4+Sqrt[16-5*z])/5},{x,0,2-2*y}]
\end{Verbatim}
\end{tcolorbox}
\vspace{0.7em}
\noindent


En la expresión del paraboloide se elimina a $y,$ de modo que la curva de intersección se encuentre en el plano $xz.$ Esto es
\[z= 4- \left ( \frac{2-x}{2} \right ) ^2-x^2= 3+x- \frac{5}{4}x^2.\]
El valor máximo de $z,$ está dado por
$z^\prime= 1-5x/2$, que al igualar a cero nos lleva a que $x=2/5$, este valor está dado por
$z=3+\frac{2}{5}-\frac{5}{4}\left(\frac{2}{5}\right)^2=\frac{16}{5}.$
\vskip0.3cm
De la expresión $z=3+x- \frac{5}{4}x^2$
también se sigue que $x=\frac{2}{5} \pm \frac{2}{5} \sqrt{16-5z}.$ En los siguientes gráficos se ilustra porque en el orden $dydzdx$ se requiere evaluar dos integrales triples mientras que en el orden $dydxdz$ son necesarias cinco integrales triples.

\vskip0.3cm
\noindent
\begin{minipage}{\textwidth}
	\centering
	\captionof{figure}{Proyección en el plano $xz$}
	\vspace{-20mm}
	\begin{minipage}{0.38\linewidth}
		\centering
		\includegraphics[width=\linewidth]{Fig/315.pdf}
	\end{minipage}
	\hspace{1em}
	\begin{minipage}{0.38\linewidth}
		\centering
		\includegraphics[width=\linewidth]{Fig/316.pdf}
	\end{minipage}
	
	\vspace{-10mm}
	
	\vspace{0.3em}
	\fuentepropia
\end{minipage}
\vskip0.3cm


Apoyándonos en estas figuras presentamos las integrales triples respec\-tivas:
\vspace{-2mm}
{\small
\[\bigintsss_{0}^2\bigintsss_{0}^{3+x-\frac{5}{4}x^2}\bigintsss_{0}^{1-x/2}dydzdx+\bigintsss_{0}^2\bigintsss_{3+x-\frac{5}{4}x^2}^{4-x^2}\bigintsss_{0}^{\sqrt{4-x^2-z}}dydzdx= \frac{19}{6}\]
}

{\footnotesize
\begin{multline*}
\bigintsss_{0}^{3}\bigintsss_{0}^{\frac{2}{5}+\frac{2}{5}\sqrt{16-5z}}\bigintsss_{0}^{1-x/2}dydxdz +\bigintsss_{3}^{\frac{16}{5}}\bigintsss_{\frac{2}{5}-\frac{2}{5}\sqrt{16-5z}}^{\frac{2}{5}+\frac{2}{5}\sqrt{16-5z}}\bigintsss_{0}^{1-x/2}dydxdz \\
+\bigintsss_{3}^{\frac{16}{5}}\bigintsss_{0}^{\frac{2}{5}-\frac{2}{5}\sqrt{16-5z}}\bigintsss_{0}^{\sqrt{4-x^2-z}}dydxdz
+\bigintsss_{\frac{16}{5}}^{4}\bigintsss_{0}^{\sqrt{4-z}}\bigintsss_{0}^{\sqrt{4-x^2-z}}dydxdz\\
+\bigintsss_{0}^{\frac{16}{5}}\bigintsss_{\frac{2}{5}+\frac{2}{5}\sqrt{16-5z}}^{\sqrt{4-z}}\bigintsss_{0}^{\sqrt{4-x^2-z}}dydxdz = \frac{19}{6}.
\end{multline*}
}

Estas integrales se evalúan en \verb"Mathematica" con la siguiente instrucción:

\vspace{0.6em}
\noindent
\begin{tcolorbox}[breakable,enhanced,title=\bf Instrucción en \verb"Mathematica", top=2pt,bottom=2pt,boxsep=0pt,before skip=2pt,after skip=0pt]
\begin{Verbatim}[formatcom=\letraverbatim]
Integrate[1,{z,0,3},{x,0,(2+2*Sqrt[16-5z])/5},
{y,0,1-x/2}]+Integrate[1,{z,3,16/5},
{x,2*(1-Sqrt[16-5z])/5,2*(1+Sqrt[16-5z])/5},
{y,0,1-x/2}]+Integrate[1,{z,3,16/5},{x,0,
(2-2*Sqrt[16-5z])/5},{y,0,Sqrt[4-z-x^2]}]+
Integrate[1,{z,16/5,4},{x,0,Sqrt[4-z]},{y,0,
Sqrt[4-z-x^2]}]+Integrate[1,{z,0,16/5},{x,
(2+2*Sqrt[16-5z])/5,Sqrt[4-z]},{y,0,Sqrt[4-z-x^2]}]
\end{Verbatim}
\end{tcolorbox}
\vspace{0.7em}
\noindent


%El programa ofrece como solución $\frac{17}{6}+\frac{2\pi}{25}-\frac{2}{375}\left( 10 \pi-44 +5\arctan\left(\frac{44}{117}\right)\right)+ \frac{1}{375}\left(37-60\arctan\left(\frac{1}{2}\right)\right)$ que es equivalente a $19/6.$
Con el siguiente proceso también hallamos el volumen del sólido.
\vskip0.2cm
%\begin{tcolorbox}[title=\bfseries Parametrizando el sólido]
Sobre el plano $xy,$ el sólido se proyecta como un triángulo con vértices $(0,0),$ $(2,0)$ y $(0,1)$. El segmento que une el punto $(0,0)$ con $(2,0)$ se parametriza como $\left\langle 2t,0\right\rangle, $ con $t\in [ 0,1].$
Los segmentos que unen cualquier punto de este segmento con el punto
$(0,1),$ se parametriza como
$(1-u)(0,1)+u(2t,0)= \left(2ut,1-u\right) .$
Es decir, todo punto en este triángulo es de la forma
$\left\langle 2ut,1-u\right\rangle $, $t\in\lbrack 0,1]$, $u\in [ 0,1]$.
Un punto sobre el paraboloide posee coordenada $z,$ que está dada por
$z=4-x^{2}-y^{2}=4-(2ut)^2-(1-u)^2= 3+2u-u^{2}\left(4t^{2}+1\right) .$
\vskip0.2cm
Consideremos un punto $P$ sobre la base del sólido {\small $P\left( 2ut,1-u,0\right)$}
y un punto $Q$ sobre el paraboloide, {\small $Q\left( 2ut,1-u,3+2u-u^{2}\left( 4t^{2}+1\right) \right)$}.
Al formar el segmento que une estos dos puntos, se obtiene una parametrización del sólido:
\vspace{-2mm}
\begin{multline*}
{\bf r}\left( t,u,w\right) = (1-w)P +wQ \\
= \left\langle 2ut,1-u,w\left( 3+2u-u^{2}\left( 4t^{2}+1\right)
\right) \right\rangle, \end{multline*}
en donde $t\in [ 0,1],$ $u\in [ 0,1],$ $w\in [ 0,1].$

\vskip0.3cm
Con la función ${\bf r}$ hallamos los siguientes vectores:
\begin{eqnarray*}
{\bf r}_{t} &=& \left\langle 2u,0,-8wu^{2}t\right\rangle \\
{\bf r}_{u} &=& \left\langle 2t,-1,w\left( 2-2u\left( 4t^{2}+1\right) \right) \right\rangle \\
{\bf r}_{w} &=& \left\langle 0,0,3+2u-u^{2}\left( 4t^{2}+1\right) \right\rangle.
\end{eqnarray*}
El elemento diferencial de volumen o producto vectorial triple está dado por
$|{\bf r}_{t} \cdot ({\bf r}_{u} \times {\bf r}_{w})| =\left| -2u\left( 3+2u-\left(4t^{2}+1\right) u^{2}\right) \right| ,$
integrando con los límites establecidos, se halla el volumen deseado:
\[\int_{0}^{1}\int_{0}^{1}\int_{0}^{1}\left| -2u\left( 3+2u-\left(4t^{2}+1\right) u^{2}\right)  \frac{19}{6}\right|\]
%\end{tcolorbox}

Estos cálculos se hacen en \verb"Mathematica" ejecutando estas instrucciones:

\vspace{0.6em}
\noindent
\begin{tcolorbox}[breakable,enhanced,title=\bf Instrucción en \verb"Mathematica", top=2pt,bottom=2pt,boxsep=0pt,before skip=2pt,after skip=0pt]
\begin{Verbatim}[formatcom=\letraverbatim]
{x,y}=(1-u)*{2t,0}+u*{0,1};
r=(1-w)*{x,y,0}+w*{x,y,4-x^2-y^2};
A=Dot[D[r,t],Cross[D[r,u],D[r,w]]];
Integrate[A,{t,0,1},{u,0,1},{w,0,1}]
\end{Verbatim}
\end{tcolorbox}
\vspace{0.7em}
\noindent


\subsection*{Un cilindro y dos planos IV}
Determinar el volumen del sólido acotado por las figuras de las expresiones
$z=1-|y|$ y $z=x^2.$ La región del espacio se muestra a continuación:

\vskip0.2cm
\noindent
\begin{minipage}{\textwidth}
	\centering
	\captionof{figure}{Sólido determinado por $x^2 \leq z \leq 1 - |y|$.}
	\label{fig:solido_parabolico_absoluto_2}
	\begin{minipage}{0.23\linewidth}
		\centering
		\includegraphics[width=\linewidth]{Img/28a.pdf}
	\end{minipage}
	\hspace{1em} 
	\begin{minipage}{0.23\linewidth}
		\centering
		\includegraphics[width=\linewidth]{Img/28c.pdf}
	\end{minipage}
		\hspace{1em} 
	\begin{minipage}{0.23\linewidth}
		\centering
		\includegraphics[width=\linewidth]{Img/28d.pdf}
	\end{minipage}
	
	\vspace{0.3em}
	\fuentepropia
	\end{minipage}
\vskip0.3cm

Las figuras anteriores se consiguen usando las siguientes instrucciones:

\vspace{0.6em}
\noindent
\begin{tcolorbox}[breakable,enhanced,title=\bf Instrucción en \verb"Mathematica", top=2pt,bottom=2pt,boxsep=0pt,before skip=2pt,after skip=0pt]
\begin{Verbatim}[formatcom=\letraverbatim]
ContourPlot3D[{z==x^2,z==1-y,z==1+y},{x,-1,1},
{y,-1,1},{z,0,1},Ticks->{{-1,0,1},{-1,0,1},{0,1}},
PlotPoints->50,Mesh->False]
RegionPlot3D[z-1<y<1-z&&x^2<z<1,{x,-1,1},{y,-1,1},
{z,0,1},PlotPoints->90,Mesh->False,PlotStyle->Orange]
\end{Verbatim}
\end{tcolorbox}
\vspace{0.7em}
\noindent
\vskip0.2cm

De la ecuación $z=x^2$ al reemplazar a $y=\left| 1-z^2\right| $
se sigue que $z=1\pm y$ y también $z=x^2-1$ o $z=1-x^2$. Estas dos curvas determinan la región de integración sobre el plano $xz.$

\vskip0.2cm
\noindent
\begin{minipage}{\textwidth}
	\centering
	\captionof{figure}{Proyección del sólido en los planos coordenados.}
	\label{fig:proyeccion_solido_parabolico_absoluto_2}
	\begin{minipage}{0.30\linewidth}
		\centering
		\vspace{-8mm}
		\includegraphics[width=\linewidth]{Fig/317.pdf}
	\end{minipage}
	\hspace{1em} 
	\begin{minipage}{0.30\linewidth}
		\centering
		\vspace{-8mm}
		\includegraphics[width=\linewidth]{Fig/318.pdf}
	\end{minipage}
	\hspace{1em} 
	\begin{minipage}{0.30\linewidth}
		\centering
		\vspace{-8mm}
		\includegraphics[width=\linewidth]{Fig/319.pdf}
	\end{minipage}
	
	\vspace{-7mm}
	\fuentepropia
\end{minipage}
\vskip0.3cm

Con estos gráficos se plantean las integrales que determinan el volumen del sólido. Integrando primero con respecto a $z,$ esta recorre puntos desde el cilindro parabólico $z=x^2$ hasta los planos $z=1-y$ y $z=1+y.$
\begin{multline*}
\int_{-1}^{1}\int_{0}^{1-x^2}\int_{x^2}^{1-y}dzdydx+\int_{-1}^{1}\int_{x^2-1}^{0}\int_{x^2}^{1+y}dzdydx= \frac{16}{15}\\
\int_0^1\int_{-\sqrt{1-y}}^{\sqrt{1-y}}\int_{x^2}^{1-y}dzdxdy+\int_{-1}^{0}\int_{-\sqrt{1+y}}^{\sqrt{1+y}}\int_{x^2}^{1+y}dzdxdy=\frac{16}{15}.
\end{multline*}
Mientras que si se integra primero con respecto a $y$,
esta recorre los puntos que se hallan entre los planos $y=z-1$ y $y=1-z.$
$\int_{-1}^{1}\int_{x^2}^{1}\int_{z-1}^{1-z}dydzdx=\int_{0}^{1}\int_{-\sqrt{z}}^{\sqrt{z}}\int_{z-1}^{1-z}dydxdz=\frac{16}{15}.$
La variable $x$ recorre los puntos que van desde la parte negativa, hasta la parte positiva del cilindro $z=x^2$,
estos se describen como $x=\pm \sqrt{z}.$
\begin{multline*} = \int_{-1}^{0}\int_{0}^{1+y}\int_{-\sqrt{z}}^{\sqrt{z}}dxdzdy+ \int_{0}^{1}\int_{0}^{1-y}\int_{-\sqrt{z}}^{\sqrt{z}}dxdzdy \\
=\int_{0}^{1}\int_{z-1}^{1-z}\int_{-\sqrt{z}}^{\sqrt{z}}dxdydz= \frac{16}{15}.
\end{multline*}
En \verb"Mathematica", las integrales anteriores en los órdenes $dzdydx,$ $dydzdy$ y $dxdydz$,
se evalúan ejecutando las siguientes instrucciones:
\vskip0.2cm

\vspace{0.6em}
\noindent
\begin{tcolorbox}[breakable,enhanced,title=\bf Instrucción en \verb"Mathematica", top=2pt,bottom=2pt,boxsep=0pt,
before skip=2pt,after skip=0pt]
\begin{Verbatim}[formatcom=\letraverbatim]
Integrate[1,{x,-1,1},{y,0,1-x^2},{z,x^2,1-y}]
+Integrate[1,{x,-1,1},{y,x^2-1,0},{z,x^2,1+y}]
Integrate[1,{x,-1,1},{z,x^2,1},{y,z-1,1-z}]
Integrate[1,{z,0,1},{y,z-1,1-z},{x,-Sqrt[z],Sqrt[z]}]
\end{Verbatim}
\end{tcolorbox}
\vspace{0.7em}
\noindent

Este valor también se puede determinar de la siguiente manera.


\vskip0.2cm
%\begin{tcolorbox}[title=\bfseries Parametrizando el sólido]
	
En el plano $xz,$ la proyección del sólido se parametriza con
$(1-u)(t,t^2)+u(t,1),$ $t\in [-1,1],$ $u\in [0,1].$
Con un punto $P(t,t^2(1-u)+u-1,t^2(1-u)+u) $ sobre el plano $y=1+z$, y otro punto $Q(t,1-t^2(1-u)-u,t^2(1-u)+u) $ sobre el plano $z=1-y$.


\vskip0.2cm
\noindent
\begin{minipage}{\textwidth}
	\centering
	\captionof{figure} {$x^2 \leq y \leq 1.$}
	\begin{minipage}{0.32\linewidth}
		\centering
		\vspace{-10mm}
		\includegraphics[width=\linewidth]{Fig/319.pdf}
	\end{minipage}
	
	\vspace{-4mm}
	\fuentepropia
\end{minipage}
\vskip0.2cm


Con estos dos puntos, definimos la parametrización del sólido, que es una combinación convexa de la forma $(1-w)P+wQ$, esto es,
{\small
\begin{equation*}
{\bf r}\left( t,u,w\right) = \left\langle t,\left( t^2-1\right) \left( 1-u\right) \left( 1-2w\right),t^2\left( 1-u\right) +u\right\rangle,
\end{equation*}
}

en donde $t\in [ -1,1]$, $u\in [ 0,1]$, $w\in [ 0,1]$. Con la función
${\bf r}$ hallamos los siguientes vectores:
\vspace{-2mm}
\begin{eqnarray*}
{\bf r}_{t}&=&\left\langle 1,2t\left( 1-u\right) \left(1-2w\right) ,2t\left( 1-u\right) \right\rangle \\
{\bf r}_{u} &=& \left\langle 0,\left( 1-t^2\right) \left( 1-2w\right),1-t^2\right\rangle \\
{\bf r}_{w} &=& \left\langle 0,2\left( 1-t^2\right) \left( 1-u\right),0\right\rangle.
\end{eqnarray*}
El producto vectorial triple es
$\left| {\bf r}_{t}\cdot \left({\bf r}_{u}\times {\bf r}_{w}\right) \right| =
%\left| \left( 1,2t\left( 1-u\right) \left( 1-2w\right) ,2t\left( 1-u\right) \right) \cdot \left( \left( 0,\left( 1-t^2\right) \left( 1-2w\right) ,1-t^2\right)\times \left( 0,2\left( 1-t^2\right) \left( 1-u\right) ,0\right) \right)\right| =
2\left| \left( t^2-1\right) ^2\left( u-1\right) \right| .$
Este término se integra para encontrar el volumen del sólido.
\[\int_{0}^{1}\int_{0}^{1}\int_{-1}^{1}2\left| \left( t^2-1\right)^2\left( u-1\right) \right| dtdudw= \frac{16}{15}.\]
%\end{tcolorbox}

Estos cálculos se hacen en \verb"Mathematica" usando las siguientes instrucciones:

\vspace{0.6em}
\noindent
\begin{tcolorbox}[breakable,enhanced,title=\bf Instrucción en \verb"Mathematica", top=2pt,bottom=2pt,boxsep=0pt,before skip=2pt,after skip=0pt]
\begin{Verbatim}[formatcom=\letraverbatim]
{x,z}=(1-u)*{t,t^2}+u*{t,1};
s=(1-w)*{x,1-z,z}+w*{x,z-1,z};
A=Dot[D[s,t],Cross[D[s,u],D[s,w]]];
Integrate[Abs[A],{t,-1,1},{u,0,1},{w,0,1}]
\end{Verbatim}
\end{tcolorbox}
\vspace{0.7em}
\noindent

\subsection*{Dos cilindros y dos planos III}
Los planos con ecuaciones $z=2-y$ y $z=0$, junto con los cilindros con ecuaciones
$x^2+y^2=4$ y $x^2+y^2=2y,\,$ que denominaremos cilindro mayor y menor respectivamente, acotan el siguiente sólido.

\vskip0.2cm
\noindent
\begin{minipage}{\textwidth}
	\centering
	\captionof{figure}{Sólido determinado por $0 \leq z \leq 2 - y$, \quad $2y \leq x^2 + y^2 \leq 4$.}
	\label{fig:solido_corona_inclinada}
	\begin{minipage}{0.23\linewidth}
		\centering
		\vspace{-2mm}
		\includegraphics[width=\linewidth]{Img/31a.pdf}
	\end{minipage}
	\hspace{1em} 
	\begin{minipage}{0.22\linewidth}
		\centering
		\vspace{-2mm}
		\includegraphics[width=\linewidth]{Img/31b.pdf}
	\end{minipage}	
	
	\vspace{0.3em}
	\fuentepropia
\end{minipage}
\vskip0.2cm


Estas figuras se pueden generar con los siguientes comandos.

\vspace{0.6em}
\noindent
\begin{tcolorbox}[breakable,enhanced,title=\bf Instrucción en \verb"Mathematica", top=2pt,bottom=2pt,boxsep=0pt,before skip=2pt,after skip=0pt]
\begin{Verbatim}[formatcom=\letraverbatim]
RegionPlot3D[x^2+y^2<4&&x^2+y^2>2*y&&z>0&&z<2-y,
{x,-2,2},{y,-2,2},{z,0,4},PlotPoints->90,Mesh->False]
\end{Verbatim}
\end{tcolorbox}
\vspace{0.7em}
\noindent

Presentaremos ahora el cálculo del volumen del sólido de otras formas.

\subsubsection*{Integración en el orden $dxdydz$ o $dxdzdy$}

La proyección en el plano $yz$ del sólido es un triángulo. Además, $x$ recorre diferentes regiones, en una de ellas va desde la parte negativa del cilindro mayor a la parte positiva del mismo cilindro.

\vskip0.3cm
En otra región, el recorrido va desde la parte negativa del cilindro mayor a la parte negativa del cilindro menor y luego desde la parte positiva del cilindro menor a la parte positiva del cilindro mayor.


\vskip0.2cm
\noindent
\begin{minipage}{\textwidth}
	\centering
	\captionof{figure}{Proyección en el plano $yz$}
	\begin{minipage}{0.29\linewidth}
		\centering
		\vspace{-8mm}
		\includegraphics[width=\linewidth]{Fig/320.pdf}
	\end{minipage}
	
	\vspace{-4mm}
	\fuentepropia
\end{minipage}
\vskip0.2cm

Por esta razón, el volumen se determina evaluando las siguientes integrales:
\vspace{-8mm}

\begin{multline*}
V=\int_{-2}^{0}\int_{0}^{2-y}\int_{-\sqrt{4-y^2}}^{\sqrt{4-y^2}}dxdzdy +
\int_{0}^2\int_{0}^{2-y}\int_{-\sqrt{4-y^2}}^{-\sqrt{2y-y^2}}dxdzdy\\
+\int_{0}^2\int_{0}^{2-y}\int_{\sqrt{2y-y^2}}^{\sqrt{4-y^2}}dxdzdy \\
=\int_{0}^2\int_{-2}^{0}\int_{-\sqrt{4-y^2}}^{\sqrt{4-y^2}}dxdydz
+\int_{2}^{4}\int_{-2}^{2-z}\int_{-\sqrt{4-y^2}}^{\sqrt{4-y^2}}dxdydz+ \\
\int_{0}^2\int_{0}^{2-z}\int_{-\sqrt{4-y^2}}^{-\sqrt{2y-y^2}}dxdydz
+\int_{0}^2\int_{0}^{2-z}\int_{\sqrt{2y-y^2}}^{\sqrt{4-y^2}}dxdydz=7\pi.
\end{multline*}


Estas integrales se evaluaron usando las siguientes instrucciones:

\vspace{0.6em}
\noindent
\begin{tcolorbox}[breakable,enhanced,title=\bf Instrucción en \verb"Mathematica", top=2pt,bottom=2pt,boxsep=0pt,before skip=2pt,after skip=0pt]
\begin{Verbatim}[formatcom=\letraverbatim]
Integrate[1,{y,-2,0},{z,0,2-y},{x,-Sqrt[4-y^2],Sqrt[4
y^2]}]+Integrate[1,{y,0,2},{z,0,2-y},{x,-Sqrt[4-y^2],
Sqrt[2*y-y^2]}]+Integrate[1,{y,0,2},{z,0,2-y},{x,
Sqrt[2*y-y^2],Sqrt[4-y^2]}]
Integrate[1,{z,0,2},{y,-2,0},{x,-Sqrt[4-y^2],Sqrt[4
y^2]}]+Integrate[1,{z,2,4},{y,-2,2-z},{x,-Sqrt[4-
y^2],Sqrt[4-y^2]}]+Integrate[1,{z,0,2},{y,0,2-z},
{x,-Sqrt[4-y^2],-Sqrt[2*y-y^2]}]+Integrate[1,{z,0,2}
,{y,0,2-z},{x,Sqrt[2*y-y^2],Sqrt[4-y^2]}]
\end{Verbatim}
\end{tcolorbox}
\vspace{0.7em}
\noindent


\subsubsection*{Integración en el orden $dzdxdy$ y $dzdydx.$}

\vskip0.2cm
\noindent
\begin{minipage}{\textwidth}
	\centering
	\captionof{figure}{Integración del sólido en el plano $xy$}
	\begin{minipage}{0.45\linewidth}
		\centering
		\vspace{-18mm}
		\includegraphics[width=\linewidth]{Fig/321.pdf}
	\end{minipage}
	\hspace{1em}
	\begin{minipage}{0.35\linewidth}
		\centering
		\vspace{-18mm}
		\includegraphics[width=\linewidth]{Fig/322.pdf}
	\end{minipage}
	
	
	\vspace{-15mm}
	\fuentepropia
\end{minipage}
\vskip0.2cm


Este gráfico es útil para calcular el volumen mediante el planteamiento de las integrales triples necesarias, integrando primero con respecto a $z.$

\vspace{-4mm}
\begin{multline*}
\int_{0}^2\int_{-\sqrt{4-y^2}}^{-\sqrt{2y-y^2}}\int_{0}^{2-y}dzdxdy+\int_{0}^2\int_{\sqrt{2y-y^2}}^{\sqrt{4-y^2}}\int_{0}^{2-y}dzdxdy \\ +\int_{-2}^{0}\int_{-\sqrt{4-y^2}}^{\sqrt{4-y^2}}\int_{0}^{2-y}dzdxdy
=\int_{-2}^{-1}\int_{-\sqrt{4-x^2}}^{\sqrt{4-x^2}}\int_{0}^{2-y}dzdydx\\ + \int_{1}^2\int_{-\sqrt{4-x^2}}^{\sqrt{4-x^2}}\int_{0}^{2-y}dzdydx +\int_{-1}^{1}\int_{-\sqrt{4-x^2}}^{1-\sqrt{1-x^2}}\int_{0}^{2-y}dzdydx \\ +\int_{-1}^{1}\int_{1+\sqrt{1-x^2}}^{\sqrt{4-x^2}}\int_{0}^{2-y}dzdydx= 7\pi
\end{multline*}

\subsubsection*{Integración en el orden $dydxdz$ y $dydzdx$}

En estos casos, se trazan unas líneas auxiliares (en azul) sobre el sólido, con el propósito de evidenciar que el sólido se divide en diferentes regiones, lo cual dará origen a varias integrales triples.

\vskip0.2cm
\noindent
\begin{minipage}{\textwidth}
	\centering
	\captionof{figure}{Cono truncado oblicuo determinado por ${\bf r}(u,t)$}
	\begin{minipage}{0.40\linewidth}
		\centering
		\vspace{-18mm}
		\includegraphics[width=\linewidth]{Fig/323.pdf}
	\end{minipage}
	\hspace{1em} 
	\begin{minipage}{0.40\linewidth}
		\centering
		\vspace{-18mm}
		\includegraphics[width=\linewidth]{Fig/324.pdf}
	\end{minipage}
	
	\vspace{-8mm}
	\fuentepropia
\end{minipage}
\vskip0.2cm

Al reemplazar $y=2-z,$ de la ecuación del plano en las ecuaciones de los cilindros se obtiene
$x^2+\left( 2-z\right) ^2=4,$ por tanto $x^2+z^2=4z$ y $x^2+\left( 2-z\right) ^2=2\left( 2-z\right),$ entonces $x^2+z^2=2z.$
El sólido proyectado en el plano $xz$, junto con las líneas mencionadas antes, se observan en los siguientes gráficos.

\vskip0.3cm
\noindent
\begin{minipage}{\textwidth}
	\centering
	\captionof{figure}{Cono truncado oblicuo determinado por ${\bf r}(u,t)$}
	\begin{minipage}{0.35\linewidth}
		\centering
		\vspace{-10mm}
		\includegraphics[width=\linewidth]{Fig/325.pdf}
	\end{minipage}
	\hspace{1em}
	\begin{minipage}{0.35\linewidth}
		\centering
		\vspace{-10mm}
		\includegraphics[width=\linewidth]{Fig/326.pdf}
	\end{minipage}
	
	\vspace{-8mm}
	\fuentepropia
\end{minipage}
\vskip0.3cm

Las diferentes tonalidades y texturas sugieren que en cada una de esas regiones se debe plantear una integral triple.
\vskip0.2cm
Integrando primero con respecto a $y$, se debe observar que
$y$ recorre los puntos que van desde la parte negativa del cilindro de mayor radio al plano.
Por tal motivo, las integrales necesarias son las siguientes: 
\begin{eqnarray}
\int_{-2}^{-1}\int_{2-\sqrt{4-x^2}}^{2+\sqrt{4-x^2}}\int_{-\sqrt{4-x^2}}^{2-z}dydzdx+
\int_{-1}^{1}\int_{1+\sqrt{1-x^2}}^{2+\sqrt{4-x^2}}\int_{-\sqrt{4-x^2}}^{2-z}dydzdx \nonumber \\
+ \int_{1}^2\int_{2-\sqrt{4-x^2}}^{2+\sqrt{4-x^2}}\int_{-\sqrt{4-x^2}}^{2-z}dydzdx \approx 11.2156 \label{31a}
\end{eqnarray}
Cuando $y$ va desde el cilindro de mayor radio al mismo cilindro, son dos regiones que en el plano $xz,$ que se proyectan de color anaranjado 
\begin{eqnarray}
\int_{-2}^{-1}\int_{0}^{2-\sqrt{4-x^2}}\int_{-\sqrt{4-x^2}}^{\sqrt{4-x^2}}dydzdx
+\int_{1}^2\int_{0}^{2-\sqrt{4-x^2}}\int_{-\sqrt{4-x^2}}^{\sqrt{4-x^2}}dydzdx \nonumber \\
=\frac{16\pi}{3} -4\sqrt{3} -\frac{20}{3} \label{31b}
\end{eqnarray}
Si $y$ recorre la región desde la parte positiva del cilindro de menor radio a la parte positiva del cilindro grande se obtiene
\begin{equation}\label{31c}
\int_{-1}^{1}\int_{0}^{2-\sqrt{4-x^2}}\int_{1+\sqrt{1-x^2}}^{\sqrt{4-x^2}}dydzdx\approx 0.0445
\end{equation}
La base de este conjunto posee una región en la \emph{parte trasera} desde el cilindro de radio mayor al menor
\begin{equation}\label{31d}
\int_{-1}^{1}\int_{0}^{1-\sqrt{1-x^2}}\int_{-\sqrt{4-x^2}}^{1-\sqrt{1-x^2}}dydzdx\approx 0.9781
\end{equation}

Desde la parte positiva del cilindro peque\~{n}o al plano, cuya proyección en el plano $xz,$ corresponde a la parte en azul oscuro, tendremos la siguiente integral 
\begin{equation}\label{31e}
\int_{-1}^{1}\int_{2-\sqrt{4-x^2}}^{1-\sqrt{1-x^2}}\int_{1+\sqrt{1-x^2}}^{2-z}dydzdx\approx 0.0376
\end{equation}
\vskip0.2cm
Desde la parte negativa del cilindro grande a la parte negativa del cilindro peque\~{n}o, cuya proyección en el plano $xz$ es la región cuadriculada: 
\begin{equation}\label{31f}
\int_{-1}^{1}\int_{1-\sqrt{1-x^2}}^{1+\sqrt{1-x^2}}\int_{-\sqrt{4-x^2}}^{1-\sqrt{1-x^2}}dydzdx\approx 6.5551
\end{equation}
Sumando las expresiones (\ref{31a}), (\ref{31b}), (\ref{31c}), (\ref{31d}), (\ref{31e}) y (\ref{31f})
se obtiene un valor a aproximadamente $7\pi$.

\vskip0.2cm
\noindent
\begin{minipage}{\textwidth}
	\centering
	\captionof{figure}{Proyección en el plano $xz$.}
	\begin{minipage}{0.38\linewidth}
		\centering
		\vspace{-15mm}
		\includegraphics[width=\linewidth]{Fig/327.pdf}
	\end{minipage}
	
	\vspace{-8mm}
	\fuentepropia
\end{minipage}
\vskip0.2cm

\begin{comment}
Estos cálculos se hacen en \verb"Mathematica" como se muestra enseguida.
% \fontsize{9.5}{12}\selectfont
\begin{verbatim}
a=NIntegrate[1,{x,-2,-1},{z,2-Sqrt[4-x^2],2+
Sqrt[4-x^2]},{y,-Sqrt[4-x^2],2-z}]+Integrate[1,
{x,-1,1},{z,1+Sqrt[1-x^2],2+Sqrt[4-x^2]},
{y,-Sqrt[4-x^2],2-z}]+NIntegrate[1,{x,1,2},
{z,2-Sqrt[4-x^2],2+Sqrt[4-x^2]},{y,-Sqrt[4-x^2],2-z}]
b=Integrate[1,{x,-2,-1},{z,0,2-Sqrt[4-x^2]},
{y,-Sqrt[4-x^2],Sqrt[4-x^2]}]+Integrate[1,{x,1,2},
{z,0,2-Sqrt[4-x^2]},{y,-Sqrt[4-x^2],Sqrt[4-x^2]}]
c=NIntegrate[1,{x,-1,1},{z,0,2-Sqrt[4-x^2]},
{y,1+Sqrt[1-x^2],Sqrt[4-x^2]}]
d=NIntegrate[1,{x,-1,1},{z,0,1-Sqrt[1-x^2]},
{y,-Sqrt[4-x^2],1-Sqrt[1-x^2]}]
e=NIntegrate[1,{x,-1,1},{z,2-Sqrt[4-x^2],
1-Sqrt[1-x^2]},{y,1+Sqrt[1-x^2],2-z}]
f=NIntegrate[1,{x,-1,1},{z,1-Sqrt[1-x^2],
1+Sqrt[1-x^2]},{y,-Sqrt[4-x^2],1-Sqrt[1-x^2]}]
a+b+c+d+e+f
\end{verbatim}

\end{comment}


Planteando las integrales en el orden $dydxdz$, es necesario considerar varias regiones.
Una de ellas es la que recorre los puntos desde la parte negativa del cilindro de mayor radio al plano.
El volumen es el resultado de considerar varias integrales triples, las cuales presentamos a continuación:

\vspace{-7mm}
\begin{multline} \label{31k}
\int_{1}^2\int_{-1}^{-\sqrt{2z-z^2}}\int_{-\sqrt{4-x^2}}^{2-z}dydxdz \\
+\int_{1}^2\int_{\sqrt{2z-z^2}}^{1}\int_{-\sqrt{4-x^2}}^{2-z}dydxdz \approx 0. 8822
\end{multline}
\begin{eqnarray}
\int_{2-\sqrt{3}}^{2+\sqrt{3}}\int_{-\sqrt{4z-z^2}}^{-1}\int_{-\sqrt{4-x^2}}^{2-z}dydxdz&\approx& 3. 3333 \label{31g}
\end{eqnarray}
\begin{eqnarray}
\int_{2-\sqrt{3}}^{2+\sqrt{3}}\int_{1}^{\sqrt{4z-z^2}}\int_{-\sqrt{4-x^2}}^{2-z}dydxdz &\approx& 3. 3333 \label{31h} \\
\int_{2}^{2+\sqrt{3}}\int_{-1}^{1}\int_{-\sqrt{4-x^2}}^{2-z}dydxdz
&=& \frac{2}{3}\pi \sqrt{3} \label{31i} \\
\int_{2+\sqrt{3}}^{4}\int_{-\sqrt{4z-z^2}}^{\sqrt{4z-z^2}}\int_{-\sqrt{4-x^2}}^{2-z}dydxdz&\approx& 0.0391 \label{31j}
\end{eqnarray}
\vskip0.3cm
Desde el cilindro de mayor radio al mismo cilindro, hay dos regiones que en el plano $xz,$ se proyectan de color anaranjado.


%\vspace{-2mm}
\begin{eqnarray}
\int_{0}^{2-\sqrt{3}}\int_{-2}^{-1}\int_{-\sqrt{4-x^2}}^{\sqrt{4-x^2}}dydxdz
+\int_{2-\sqrt{3}}^2\int_{-2}^{-\sqrt{4z-z^2}}\int_{-\sqrt{4-x^2}}^{\sqrt{4-x^2}}dydxdz \nonumber \\
+\int_{0}^{2-\sqrt{3}}\int_{1}^2\int_{-\sqrt{4-x^2}}^{\sqrt{4-x^2}}dydxdz
+\int_{2-\sqrt{3}}^2\int_{\sqrt{4z-z^2}}^2\int_{-\sqrt{4-x^2}}^{\sqrt{4-x^2}}dydxdz \nonumber \\
= \frac{16\pi}{3}-\frac{20}{3}-4\sqrt{3} \label{31l}
\end{eqnarray}
\vskip0.3cm
De la parte positiva del cilindro peque\~{n}o al plano, cuya proyección en el plano $xz,$ corresponde a la parte en azul oscuro.

\begin{multline}
\int_{0}^{2-\sqrt{3}}\int_{-\sqrt{4z-z^2}}^{-\sqrt{2z-z^2}}\int_{1+\sqrt{1-x^2}}^{2-z}dydxdz + \\ \int_{2-\sqrt{3}}^{1}\int_{-1}^{-\sqrt{2z-z^2}}\int_{1+\sqrt{1-x^2}}^{2-z}dydxdz+ \int_{0}^{2-\sqrt{3}}\int_{\sqrt{2z-z^2}}^{\sqrt{4z-z^2}}\int_{1+\sqrt{1-x^2}}^{2-z}dydxdz \\
+\int_{2-\sqrt{3}}^{1}\int_{\sqrt{2z-z^2}}^{1}\int_{1+\sqrt{1-x^2}}^{2-z}dydxdz \approx 0.0376 \label{31m}
\end{multline}

\vskip0.4cm
Cuando $y$ recorre los puntos desde la parte positiva del cilindro pequeño a la parte positiva del mayor cilindro, la integral es:
\vspace{-8mm}
\begin{multline}
\int_{0}^{2-\sqrt{3}}\int_{-1}^{-\sqrt{4z-z^2}}\int_{1+\sqrt{1-x^2}}^{\sqrt{4-x^2}}dydxdz+ \\
\int_{0}^{2-\sqrt{3}}\int_{\sqrt{4z-z^2}}^{1}\int_{1+\sqrt{1-x^2}}^{\sqrt{4-x^2}}dydxdz \approx 0.0445 \label{31n}
\end{multline}

\vskip0.3cm
Existe una parte (la región en azul es parte de ella) que posee una \emph{región atras}, donde $y$ se recorre desde la parte negativa del cilindro mayor hasta la parte negativa el cilindro menor.
\vskip0.2cm
\begin{multline}
\int_{0}^{1}\int_{-1}^{-\sqrt{2z-z^2}}\int_{-\sqrt{4-x^2}}^{1-\sqrt{1-x^2}}dydxdz \\
+\int_{0}^{1}\int_{\sqrt{2z-z^2}}^{1}\int_{-\sqrt{4-x^2}}^{1-\sqrt{1-x^2}}dydxdz \approx 0.9781 \label{31o}
\end{multline}

\vskip0.4cm
Desde la parte negativa del cilindro grande hasta la parte negativa del cilindro pequeño, cuya proyección en el plano $xz$ es la región cuadriculada.
\vskip0.4cm
\begin{equation} \label{31p}
\int_{0}^2\int_{-\sqrt{2z-z^2}}^{\sqrt{2z-z^2}}\int_{-\sqrt{%
4-x^2}}^{1-\sqrt{1-x^2}}dydxdz \approx 6. 5551
\end{equation}
\vskip0.4cm
Sumando las expresiones (\ref{31k}), (\ref{31g}), (\ref{31h}), (\ref{31i}), (\ref{31j}), (\ref{31l}), (\ref{31m}), (\ref{31n}),
(\ref{31o}) y (\ref{31p}), se obtiene un valor aproximado a $7\pi.$

\vskip0.3cm
Con ayuda de \verb"Mathematica" se puede obtener estos valores, para tal fin se definen cada una de las integrales como $a$, $b$, $c$, $d$, $e$, $f$, $g$, $h$, $i$ y $j$.
La suma es aproximadamente $7\pi$.
\vskip0.3cm
\vspace{0.6em}
\noindent
\begin{tcolorbox}[breakable,enhanced,title=\bf Instrucción en \verb"Mathematica", top=2pt,bottom=2pt,boxsep=0pt,before skip=2pt,after skip=0pt]
\begin{Verbatim}[formatcom=\letraverbatim]
a=NIntegrate[1,{z,2-Sqrt[3],2+Sqrt[3]},{x,Sqrt[4*z-z^2],-1},
{y,-Sqrt[4-x^2],2-z}]
b=NIntegrate[1,{z,2-Sqrt[3],2+Sqrt[3]},{x,1,Sqrt[4*z-z^2]},
{y,-Sqrt[4-x^2],2-z}]
c=NIntegrate[1,{z,2,2+Sqrt[3]},{x,-1,1},{y,-Sqrt[4-x^2],2-z}]
d=NIntegrate[1,{z,2+Sqrt[3],4},{x,-Sqrt[4*z-z^2],
Sqrt[4*z-z^2]},{y,-Sqrt[4-x^2],2-z}]
e=NIntegrate[1,{z,1,2},{x,-1,-Sqrt[2*z-z^2]},
{y,-Sqrt[4-x^2],2-z}]+NIntegrate[1,{z,1,2},
{x,Sqrt[2*z-z^2],1},{y,-Sqrt[4-x^2],2-z}]
f=NIntegrate[1,{z,0,2-Sqrt[3]},{x,-2,-1},{y,
Sqrt[4-x^2],Sqrt[4-x^2]}]+NIntegrate[1,{z,
2-Sqrt[3],2},{x,-2,-Sqrt[4*z-z^2]},{y,
Sqrt[4-x^2],Sqrt[4-x^2]}]+NIntegrate[1,
{z,0,2-Sqrt[3]},{x,1,2},{y,-Sqrt[4-x^2],
Sqrt[4-x^2]}]+NIntegrate[1,{z,2-Sqrt[3],2},
{x,Sqrt[4*z-z^2],2},{y,-Sqrt[4-x^2],Sqrt[4-x^2]}]
g=Integrate[1,{z,0,2-Sqrt[3]},{x,-Sqrt[4*z-z^2],
Sqrt[2*z-z^2]},{y,1+Sqrt[1-x^2],2-z}]+Integrate[1,{z,2-
Sqrt[3],1},{x,-1,-Sqrt[2*z-z^2]},
{y,1+Sqrt[1-x^2],2-z}]+
Integrate[1,{z,0,2-Sqrt[3]},{x,Sqrt[2*z-z^2],
Sqrt[4*z-z^2]},
{y,1+Sqrt[1-x^2],2-z}]+Integrate[1,{z,2-Sqrt[3],1},
{x,Sqrt[2*z-z^2],1},{y,1+Sqrt[1-x^2],2-z}]
h=NIntegrate[1,{z,0,2-Sqrt[3]},{x,-1,-Sqrt[4*z-z^2]},
{y,1+Sqrt[1-x^2],Sqrt[4-x^2]}]+
NIntegrate[1,{z,0,2-Sqrt[3]},
{x,Sqrt[4*z-z^2],1},{y,1+Sqrt[1-x^2],Sqrt[4-x^2]}]
i=NIntegrate[1,{z,0,1},{x,-1,-Sqrt[2*z-z^2]},
{y,-Sqrt[4-x^2],
1-Sqrt[1-x^2]}]+NIntegrate[1,{z,0,1},
{x,Sqrt[2*z-z^2],1},
{y,-Sqrt[4-x^2],1-Sqrt[1-x^2]}]
j=NIntegrate[1,{z,0,2},{x,-Sqrt[2*z-z^2],Sqrt[2*z-z^2]},
{y,-Sqrt[4-x^2],1-Sqrt[1-x^2]}]

a+b+c+d+e+f+g+h+i+j
\end{Verbatim}
\end{tcolorbox}
\vspace{0.7em}
\noindent

\vskip0.2cm
Presentamos dos maneras de calcular el volumen del sólido, la primera de ellas es usando coordenadas cilíndricas, luego se parametriza el sólido.
\begin{align*}
\int_0^{\pi} \int_0^2 \int_{2 \sin \theta}^{2 - r \sin \theta} r\,dz\,dr\,d\theta 
+
\int_{\pi}^{2\pi} \int_0^2 \int_0^{2 - r \sin \theta} r\,dz\,dr\,d\theta &= 7\pi
\end{align*}


En coordenadas polares, los cilindros se representan como
{\small $r=2\sin \theta$} y {\small $r=2.$} Las parametrizaciones elegidas son
{\small $\langle 2\sin \theta \cos \theta , 2\sin \theta\sin \theta \rangle$} y {\small $\langle 2 \cos \theta , 2\sin \theta \rangle.$}
Para recorrer adecuadamente cada curva, iniciando en el punto $(0,2),$ se sustituye $\theta$ por $\theta + \frac{\pi}{2}.$

	
\vskip0.2cm
\noindent
\begin{minipage}{\textwidth}
	\centering
	\captionof{figure}{\textbf{$2y \leq x^2 + y^2 \leq 4.$}}
	\begin{minipage}{0.38\linewidth}
		\centering
		\vspace{-15mm}
		\includegraphics[width=\linewidth]{Fig/328.pdf}
	\end{minipage}
	
	\vspace{-10mm}
	\fuentepropia
\end{minipage}
\vskip0.2cm
	

Los puntos $P$ y $Q$ del gráfico, se determinan por medio de las funciones
$\mathbf{r}_1(\theta)= \langle -\sin(2\theta), 1+\cos(2\theta) \rangle,$
y $\mathbf{r}_2(\theta) = \langle -2 \sin(2\theta), 2\cos(2\theta)\rangle,$
res\-pectivamente.
Una combinación convexa de estos parametriza la región sombreada. Es decir,

\vspace{-8mm}
\[\mathbf{r}_3(\theta, u) = \langle (u-2) \sin(2\theta), u - (u-2) \cos(2\theta) \rangle, \theta \in [0, \pi], u \in [0, 1].\]

Con un punto $A((u-2)\sin(2\theta),u-(u-2)\cos(2\theta),0)$ el plano $z=0$ y otro punto
$B((u-2)\sin(2\theta),u-(u-2)\cos(2\theta),2-(u-(u-2)\cos(2\theta))),$
sobre el plano $z=2-y$ se forma la combinación convexa de la forma $(1-w)A+ wB,$ que describe el sólido, esto define la función:

\vspace{-8mm}
\[\mathbf{r}(\theta, u, w) = \left\langle (u-2) \sin(2\theta), u - (u-2) \cos(2\theta), 2w(2-u) \sin^2(\theta) \right\rangle,\]

en donde $\theta \in [0,\pi],\,u\in[0,1]$ y $w\in[0,1].$ 
Derivando parcialmente se obtienen los siguientes vectores

\vspace{-8mm}
{\footnotesize
\begin{align*}
{\bf r}_\theta &= \langle 2(u-1)\cos(2\theta),2(u-2)\sin(2\theta),2w(2-u)\sin(2\theta) \rangle \\
{\bf r}_u &= \langle \sin(2\theta),2\sin^2(\theta),-2w\sin^2(\theta) \rangle \\
{\bf r}_w &= \langle 0,0,2(2-u)\sin^2(\theta) \rangle \\
{\bf r}_{\theta} \times {\bf r}_u &= \langle 0, 4w(2 - u) \sin^2(t), 4 (2 - u) \sin^2(t) \rangle
\end{align*}
}

\vspace{-3mm}
Con lo cual, $|{\bf r}_w \cdot ({\bf r}_{\theta} \times {\bf r}_u)|= 8(u-2)^2\sin^4(\theta).$
Entonces, el volumen del sólido está dado por:

{\footnotesize
\[\int_0^1 \int_0^1 \int_0^\pi 8(u-2)^2\sin^4(\theta) d\theta du dw =7\pi\]
}
	
%\end{tcolorbox}

Este valor se obtiene en \verb"Mathematica" usando las siguientes instrucciones.


\vspace{0.6em}
\noindent
\begin{tcolorbox}[breakable,enhanced,title=\bf Instrucción en \verb"Mathematica", top=2pt,bottom=2pt,boxsep=0pt,before skip=2pt,after skip=0pt]
\begin{Verbatim}[formatcom=\letraverbatim]
{x,y,z}={(u-2)*Sin[2*t],u-(u-2)*Cos[2*t],2-y};
r=Simplify[(1-w)*{x,y,0}+w*{x,y,z}]
A=Simplify[Dot[D[r,w],Cross[D[r,t],D[r,u]]]]
Integrate[Abs[A],{w,0,1},{u,0,1},{t,0,Pi}]
\end{Verbatim}
\end{tcolorbox}
\vspace{0.7em}
\noindent

\section{Sólidos acotados por paraboloides} \index{Sólidos!con paraboloides}
\vskip0.5cm

\subsection*{Un paraboloide elíptico} \index{Paraboloide!elíptico}
Consideremos el paraboloide con vértice en el punto $(4,3,2)$ que abre hacia la parte negativa del eje $x$, determinado por la ecuación
$\frac{(z-2)^2}{4}+\frac{(y-3)^2}{9}=1-\frac{x}{4},$
esta superficie junto con el plano $x=0,$ acotan un sólido. La figura es la siguiente:

\vskip0.2cm
\noindent
\begin{minipage}{\textwidth}
	\centering
	\captionof{figure}{Sólido definido por $0 \leq x \leq 4 - (z - 2)^2 - \frac{4}{9}(y - 3)^2$.}
	\label{fig:solido_eliptico_desplazado_2}
	\begin{minipage}{0.20\linewidth}
		\centering
		\includegraphics[width=\linewidth]{Img/32a.pdf}
	\end{minipage}
	\hspace{1em} 
	\begin{minipage}{0.20\linewidth}
		\centering
		\includegraphics[width=\linewidth]{Img/32.pdf}
	\end{minipage}
	
	\vspace{0.3em}
	
	\vspace{0.3em}
	\fuentepropia
\end{minipage}

Las figuras anteriores se consiguen con los siguientes comandos.

\vspace{0.6em}
\noindent
\begin{tcolorbox}[breakable,enhanced,title=\bf Instrucción en \verb"Mathematica", top=2pt,bottom=2pt,boxsep=0pt,before skip=2pt,after skip=0pt]
\begin{Verbatim}[formatcom=\letraverbatim]
ContourPlot3D[(z-2)^2/4+(y-3)^2/9==1-x/4,{x,0,4},
{y,0,6},{z,0,4},PlotPoints->80,MeshStyle->Gray,
Mesh->5,ContourStyle->Directive[RGBColor[1,0.3,0.3],
Specularity[RGBColor[1,1,0],10]]]
\end{Verbatim}
\end{tcolorbox}
\vspace{0.7em}
\noindent

En el plano $yz$ $(x=0),$ se satisface que $\frac{\left( z-2\right) ^2}{4}+%
\frac{\left( y-3\right) ^2}{9}=1,$ esta ecuación representa una elipse con centro en el punto $(2,3),$ alargada en el sentido de las $y$, y puede parametrizarse en la forma
$y=3\sin t+3,\,\, z=2+2\cos t,\,\,t\in [0,2\pi ].$
\vskip0.3cm
La proyección en el plano $xy,$ se consigue cuando $z=2,$ y es representada por la ecuación
$\frac{\left( y-3\right) ^2}{9}=1-\frac{x}{4},$ que se puede reescrbir como $x=-\frac{4}{9}y^2+\frac{8}{3}y.$ Esta ecuación representa una parábola que abre hacia la parte negativa del eje $x$ y posee
vértice en el punto $\left( 4,3\right) .$ Una parametrización de esta curva es
$y=3t,\,\, x=-\frac{4}{9}\left( 3t\right) ^2+\frac{8}{3}(3t) = 8t-4t^2,$ para $t\in [0,2].$
\vskip0.3cm
Cuando $y=3,$ obtendremos la proyección del sólido en el plano $xz,$
esto es $\frac{\left( z-2\right) ^2}{4}=1-\frac{x}{4},$ de donde
$x=4(1-\frac{(z-2)^2}{4}) = 4z-z^2.$
Esta curva se puede parametrizar como $x=4t-t^2,\,\,z=t,$ para $t\in [0,4].$
\vskip0.3cm
Con los resultados anteriores se presentan ahora las figuras.

\vskip0.2cm
\noindent
\begin{minipage}{\textwidth}
	\centering
	\captionof{figure}{Proyecciones del sólido en los planos coordenados.}
	\label{fig:proyeccion_solido_eliptico_desplazado_2}
	\begin{minipage}{0.30\linewidth}
		\centering
		\vspace{-10mm}
		\includegraphics[width=\linewidth]{Fig/329.pdf}
	\end{minipage}
	\hspace{1em} 
	\begin{minipage}{0.30\linewidth}
		\centering
		\vspace{-10mm}
		\includegraphics[width=\linewidth]{Fig/330.pdf}
	\end{minipage}
	\hspace{1em} 
	\begin{minipage}{0.30\linewidth}
		\centering
		\vspace{-10mm}
		\includegraphics[width=\linewidth]{Fig/331.pdf}
	\end{minipage}
	
	\vspace{-8mm}
	\fuentepropia
\end{minipage}
\vskip0.2cm

% \begin{figure}[H]
% \centering
% % Primer gráfico
% \begin{pspicture}(-0.5,-0.4)(3.5,3) %\psgrid
% \psset{algebraic,unit=0.5}
% \psaxes[Dx=2,Dy=2,linecolor=gray]{->}(0,0)(-0.5,-1.2)(5,6.2)[$x$,-90][$z$,0]
% \pscustom[fillstyle=solid,fillcolor=gray!30,opacity=0.7,linestyle=none]{
%   \psparametricplot{0}{4}{4*t-t^2,t}
%   \psline(0,4)(0,0)
% }
% \psparametricplot{0}{4}{4*t-t^2,t}
% \rput(2.5,4.5){$\frac{(z-2)^2}{4}=1-\frac{x}{4}$}
% \end{pspicture}

% \vspace{0.5cm}

% % Segundo gráfico
% \begin{pspicture}(-0.5,-0.4)(3,3) % \psgrid
% \psset{algebraic,yunit=0.35,xunit=0.45}
% \psaxes[Dx=2,Dy=3,linecolor=gray]{->}(0,0)(-1.5,-1.2)(6.2,8)[$x$,-90][$y$,0]
% \pscustom[fillstyle=solid,fillcolor=gray!30,opacity=0.7,linestyle=none]{
%   \psparametricplot{0}{2}{8*t-4*t^2,3*t}
%   \psline(0,6)(0,0)
% }
% \psparametricplot{0}{2}{8*t-4*t^2,3*t}
% \rput(4.5,6){$\frac{(y-3)^2}{9}=1-\frac{x}{4}$}
% \end{pspicture}

% \vspace{0.5cm}

% % Tercer gráfico
% \begin{pspicture}(-1,-0.4)(3.5,3) % \psgrid
% \psset{algebraic,unit=0.5}
% \psaxes[Dx=2,Dy=2,linecolor=gray]{->}(0,0)(-0.5,-1.2)(7,6)[$y$,-90][$z$,0]
% \pscustom[fillstyle=solid,fillcolor=gray!30,opacity=0.7,linestyle=none]{
%   \psparametricplot{0}{\psPiTwo}{3+3*sin(t),2+2*cos(t)}
% }
% \psparametricplot{0}{\psPiTwo}{3+3*sin(t),2+2*cos(t)}
% \psline[linecolor=blue](5,3.5)(5,0.5)
% \rput(4,4.7){$\frac{(z-2)^2}{4}+\frac{(y-3)^2}{9}=1$}
% \rput(5,3.5){$\bullet$}
% \rput(5,0.5){$\bullet$}
% \rput(5,2.5){$\bullet$}
% \rput(4.5,0.8){$A$}
% \rput(4.5,3.2){$B$}
% \rput(5.5,2.5){$P$}
% \end{pspicture}
% \caption{Proyecciones del sólido en los planos coordenados.}
% \end{figure}

Estas figuras resultan útiles al momento de plantear las integrales triples que se requiere evaluar para calcular el volumen del sólido, veámoslo
\vskip0.3cm
A partir de la ecuación del paraboloide %$\frac{(z-2)^2}{4}+\frac{(y-3)^2}{9}=1-\frac{x}{4}$,
se obtiene {\footnotesize $x=4\left( 1-\frac{(z-2)^2}{4}-\frac{(y-3)^2}{9}\right),$} por tanto

\[
\frac{(z-2)^2}{4} + \frac{(y-3)^2}{9} = 1 - \frac{x}{4},
\]
se obtiene
{\small
\[
x = 4\left( 1 - \frac{(z-2)^2}{4} - \frac{(y-3)^2}{9} \right),
\]
}

por tanto
\[
\displaystyle
\int_{0}^{6}
\int_{2 - \frac{2}{3}\sqrt{6y - y^2}}^{2 + \frac{2}{3}\sqrt{6y - y^2}}
\int_{0}^{4 - (z - 2)^2 - \frac{4}{9}(y - 3)^2}
dx\,dz\,dy = 12\pi,
\]
también
\[
\displaystyle
\int_{0}^{4}
\int_{3 - \frac{3}{2}\sqrt{-z^2 + 4z}}^{3 + \frac{3}{2}\sqrt{-z^2 + 4z}}
\int_{0}^{4 - (z - 2)^2 - \frac{4}{9}(y - 3)^2}
dx\,dy\,dz = 12\pi.
\]

De la ecuación original, se sigue que
\[
(z - 2)^2 = 4\left( 1 - \frac{x}{4} - \frac{(y - 3)^2}{9} \right)
= \frac{1}{9}\left( 24y - 9x - 4y^2 \right),
\]
por tanto
\[
z = 2 \pm \frac{1}{3}\sqrt{-9x - 4y^2 + 24y}.
\]

Entonces, las integrales triples, integrando primero con respecto a \( z \), son:
\[
\displaystyle
\int_{0}^{4}
\int_{3 - \frac{3}{2}\sqrt{4 - x}}^{3 + \frac{3}{2}\sqrt{4 - x}}
\int_{2 - \frac{1}{3}\sqrt{-9x - 4y^2 + 24y}}^{2 + \frac{1}{3}\sqrt{-9x - 4y^2 + 24y}}
dz\,dy\,dx = 12\pi,
\]
\[
\displaystyle
\int_{0}^{6}
\int_{0}^{\frac{8}{3}y - \frac{4}{9}y^2}
\int_{2 - \frac{1}{3}\sqrt{-9x - 4y^2 + 24y}}^{2 + \frac{1}{3}\sqrt{-9x - 4y^2 + 24y}}
dz\,dx\,dy = 12\pi.
\]

De la misma forma,
\[
(y - 3)^2 = 9\left( 1 - \frac{x}{4} - \frac{(z - 2)^2}{4} \right)
= \frac{9}{4}(4z - x - z^2),
\]
entonces
\[
y = 3 \pm \frac{3}{2}\sqrt{4z - x - z^2}.
\]
Con lo cual, el volumen está dado por
{\small
\[\bigintss_{0}^{4}\bigintss_{2-\sqrt{4-x}}^{2+\sqrt{4-x}}\bigintss_{3-\frac{3}{2}\sqrt{4z-x-z^2} }^{3+\frac{3}{2}\sqrt{4z-x-z^2}}dydzdx= 12\pi\]
\[\bigintss_{0}^{4}\bigintss_{0}^{4z-z^2}\bigintss_{3-\frac{3}{2}\sqrt{4z-x-z^2}}^{3+\frac{3}{2}\sqrt{4z-x-z^2} }dydxdz=12\pi.\]
}

Algunas de estas integrales se evalúan con los siguientes comandos:

\vspace{0.6em}
\noindent
\begin{tcolorbox}[breakable,enhanced,title=\bf Instrucción en \verb"Mathematica", top=2pt,bottom=2pt,boxsep=0pt,before skip=2pt,after skip=0pt]
\begin{Verbatim}[formatcom=\letraverbatim]
Integrate[1,{x,0,4},{y,3-3*Sqrt[4-x]/2,3+
3*Sqrt[4-x]/2},{z,2-Sqrt[24*y-4y^2-9*x]/3,
2+Sqrt[24*y-4y^2-9*x]/3}]
Integrate[1,{y,0,6},{z,2-2Sqrt[6*y-y^2]/3,2
+2Sqrt[6*y-y^2]/3},{x,0,4-(z-2)^2-4*(y-3)^2/9}]
Integrate[1,{z,0,4},{x,0,4*z-z^2},{y,3-
3*Sqrt[4*z-x-z^2]/2,3+3*Sqrt[4*z-x-z^2]/2}]
Integrate[1,{z,0,4},{y,3-3/2*Sqrt[4z-z^2],3+
3/2*Sqrt[4z-z^2]},{x,0,4-(z-2)^2-4/9*(y-3)^2}]
\end{Verbatim}
\end{tcolorbox}
\vspace{0.7em}
\noindent


Se presenta ahora, otra manera de hallar el volumen.
\vskip0.2cm

%\begin{tcolorbox}[title=\bfseries Parametrizando el sólido]

La proyección del sólido en el plano $yz$ corresponde a una elipse cuyo interior se puede parametrizar como:
$y=3+3u\sin t,\ z=2+2u\cos t,$ en donde $u \in [0,1]$, $t \in [0,2\pi].$
Con dos puntos $A$ y $B$ sobre la elipse (con igual coordenada $y$) se escoge un punto $P$;
este punto se proyecta sobre el paraboloide. Con una combinación convexa entre $P$ y $Q$
se parametriza el sólido.


\vskip0.2cm
\noindent
\begin{minipage}{\textwidth}
	\centering
	\captionof{figure}{$(2uv, 2(1 - v + uv)).$}
	\label{fig:curva_parametrica_uv_3}
	\vspace{-18mm}
	\begin{minipage}{0.40\linewidth}
		\centering
		\includegraphics[width=\linewidth]{Fig/332.pdf}
	\end{minipage}
	
	\vspace{-16mm}
	\fuentepropia
\end{minipage}
\vskip0.2cm

Las coordenadas $x$ del punto $Q$ que está sobre el paraboloide satisfacen que
$x=4-4u^2.$ Por tanto, los puntos mencionados son
\[P\left(0,3+3u\sin t,2+2u\cos t\right)\] y \[Q\left( 4-4u^2,3+3u\sin t+3,2+2u\cos t\right).\]
\vskip0.2cm
El sólido se parametriza simplificando la expresión $(1-w)Q+wP$,
esto permite definir la función

\[{\bf r}(u,t,w)=\left\langle 4(1-w)(1-u^2),3+3u\sin t,2+2u\cos t\right\rangle,\]
\vskip0.2cm
con $\, u\in [0,1],\,t\in [0,2\pi ],\, w\in [0,1] $.
Con esta función, hallamos los siguientes vectores
\begin{eqnarray*}
{\bf r}_{t}&=& \langle 0,3u\cos t,-2u\sin t \rangle \\
{\bf r}_{u}&=& \langle 8\left( w-1\right) u,3\sin t,2\cos t\rangle \\
{\bf r}_{w}&=& \langle 4u^2-4,0,0\rangle \\
{\bf r}_{u}\times {\bf r}_{w} &=& \langle 0,8(1-u^2)\cos(t),12(u^2-1)\sin (t)\rangle.
\end{eqnarray*}
\vskip0.2cm
Entonces
$ \left|{\bf r}_{t}\cdot \left({\bf r}_{u}\times {\bf r}_{w}\right) \right| = |24u^{3}-24u|.$

Con lo cual, el volumen del sólido esta dado por

\[\int_{0}^{1}\int_{0}^{1}\int_{0}^{2\pi }\left| 24u^{3}-24u\right|dtdudw\\
= 12\pi.\]
%\end{tcolorbox}

\subsection*{Una esfera y un paraboloide} \index{Esfera}
La esfera con ecuación $x^2+y^2+z^2=2$ y el paraboloide
$z=x^2+y^2$ se intersecan en $z=1.$
Las siguientes figuras muestran esta situación y el sólido acotado por estas superficies.

\vskip0.3cm
\noindent
\begin{minipage}{\textwidth}
	\centering
	\captionof{figure}{Sólido definido por $x^2 + y^2 \leq z \leq \sqrt{2 - x^2 - y^2}$.}
	\label{fig:solido_paraboloide_esfera}
	\begin{minipage}{0.26\linewidth}
		\centering
		\includegraphics[width=\linewidth]{Img/81a.pdf}
	\end{minipage}
	\hspace{1em}
	\begin{minipage}{0.26\linewidth}
		\centering
		\includegraphics[width=\linewidth]{Img/81b.pdf}
	\end{minipage}
	
	\vspace{0.3em}
	\fuentepropia
\end{minipage}
\vskip0.3cm

Estas figuras se consiguen con las siguientes indicaciones:


\vspace{0.6em}
\noindent
\begin{tcolorbox}[breakable,enhanced,title=\bf Instrucción en \verb"Mathematica", top=2pt,bottom=2pt,boxsep=0pt,before skip=2pt,after skip=0pt]
\begin{Verbatim}[formatcom=\letraverbatim]
ContourPlot3D[{z==x^2+y^2,x^2+y^2+z^2==2},
{x,-1.5,1.5},{y,-1.5,1.5},{z,-1.5,1.5},Mesh->False,
ContourStyle->{Directive[Cyan,Opacity[0.9]],
Directive[Orange,Opacity[0.65]]},PlotPoints->60]
RegionPlot3D[x^2+y^2<z<Sqrt[2-x^2-y^2],{x,-1,1},
{y,-1,1},{z,0,1.5},PlotStyle->Orange,Mesh->False,
PlotPoints->80,Ticks->{{-1,0,1},{-1,0,1},{0,1}}]
\end{Verbatim}
\end{tcolorbox}
\vspace{0.7em}
\noindent

La proyección en el plano $xy$ está dada por $x^2+y^2=1,$ que corresponde a una circunferencia de radio $1$, mientras que en el plano $xz$ la proyección es una porción de parábola y una circunferencia de radio $\sqrt{2}$.

\vskip0.2cm
\noindent
\begin{minipage}{\textwidth}
	\centering
	\captionof{figure}{Regiones $x^2\leq z \leq \sqrt{2-x^2}$, $x^2+y^2\leq 1.$}
	\begin{minipage}{0.35\linewidth}
		\centering
		\vspace{-10mm}
		\includegraphics[width=\linewidth]{Fig/333.pdf}
	\end{minipage}
	\hspace{1em}
	\begin{minipage}{0.35\linewidth}
		\centering
		\vspace{-10mm}
		\includegraphics[width=\linewidth]{Fig/334.pdf}
	\end{minipage}
	
	\vspace{-8mm}
	\fuentepropia
\end{minipage}
\vskip0.2cm

El volumen del sólido acotado por estas superficies está dado por
\begin{multline*}
\int_{-1}^{1}\int_{-\sqrt{1-x^2}}^{\sqrt{1-x^2}}\int_{x^2+y^2}^{\sqrt{2-x^2-y^2}}dzdydx = \\ \int_{-1}^{1}\int_{-\sqrt{1-y^2}}^{\sqrt{1-y^2}}\int_{x^2+y^2}^{\sqrt{2-x^2-y^2}}dzdxdy= \frac{\pi}{6} ( 8\sqrt{2}-7)
\end{multline*}

Al integrar primero con respecto a $y$, el sólido queda dividido en dos porciones, en una de ellas $y$ toma valores desde la parte negativa de la esfera hasta la parte positiva de la misma; en la otra porción $y$ recorre los puntos desde la parte negativa del paraboloide hasta su parte positiva.
Por lo tanto, el volumen en este caso está dado por
%\vskip0.2cm
\[
\int_{-1}^{1} \int_{1}^{\sqrt{2 - x^2}} \int_{-\sqrt{2 - x^2 - z^2}}^{\sqrt{2 - x^2 - z^2}} dy\, dz\, dx
+ \int_{-1}^{1} \int_{x^2}^{1} \int_{-\sqrt{z - x^2}}^{\sqrt{z - x^2}} dy\, dz\, dx
\approx 2.259
\]
\vspace{-2mm}
{\small
\[
\int_{1}^{\sqrt{2}} \int_{-\sqrt{2 - z^2}}^{\sqrt{2 - z^2}} \int_{-\sqrt{2 - x^2 - z^2}}^{\sqrt{2 - x^2 - z^2}} dy\, dx\, dz
+ \int_{0}^{1} \int_{-\sqrt{z}}^{\sqrt{z}} \int_{-\sqrt{z - x^2}}^{\sqrt{z - x^2}} dy\, dx\, dz
= \frac{\pi}{6}(8\sqrt{2} - 7)
\]
}

\vskip0.2cm
Las anteriores integrales se evaluaron en \verb"Mathematica" usando las si\-guientes instrucciones:

\vspace{0.6em}
\noindent
\begin{tcolorbox}[breakable,enhanced,title=\bf Instrucción en \verb"Mathematica", top=2pt,bottom=2pt,boxsep=0pt,before skip=2pt,after skip=0pt]
\begin{Verbatim}[formatcom=\letraverbatim]
Integrate[1,{x,-1,1},{y,-Sqrt[1-x^2],Sqrt[1-x^2]},
{z,x^2+y^2,Sqrt[2-x^2-y^2]}]
Integrate[1,{y,-1,1},{x,-Sqrt[1-y^2],Sqrt[1-y^2]},
{z,x^2+y^2,Sqrt[2-x^2-y^2]}]
Integrate[1,{x,-1,1},{z,1,Sqrt[2-x^2]},
{y,-Sqrt[2-x^2-z^2],Sqrt[2-x^2-z^2]}]+Integrate[1,
{x,-1,1},{z,x^2,1},{y,-Sqrt[z-x^2],Sqrt[z-x^2]}]
Integrate[1,{z,1,Sqrt[2]},{x,-Sqrt[2-z^2],
Sqrt[2-z^2]},{y,-Sqrt[2-x^2-z^2],Sqrt[2-x^2-z^2]}]
+Integrate[1,{z,0,1},{x,-Sqrt[z],Sqrt[z]},
{y,-Sqrt[z-x^2],Sqrt[z-x^2]}]
\end{Verbatim}
\end{tcolorbox}
\vspace{0.7em}
\noindent


En el orden $dxdydz$ y $dxdzdy,$ la situación es muy similar al caso $dydzdx$,
por la simetría de la figura con respecto al plano
$xz$ o $yz.$ Con el siguiente procedimiento también se halla el volumen del sólido.

\vskip0.2cm
%\begin{tcolorbox}[title=\bfseries Parametrizando el sólido]
La proyección del sólido en el plano $xy$ en la forma $\left( u\cos v,u\sin v\right) ,$ con $u\in [0,1],$ $v\in [0,2\pi ],$ la coordenada $z$ sobre la esfera y el paraboloide están dadas, respectivamente, por $z =\sqrt{2-x^2-y^2}= \sqrt{2-u^2}$ y $z=x^2+y^2= u^2.$
El sólido se parametriza simplificando una combinación convexa de la forma $(1-w)P+wQ,$
en donde $P\left( u\cos v,u\sin v,\sqrt{2-u^2}\right)$ es un punto sobre la esfera y
$Q\left( u\cos v,u\sin v,u^2\right)$ es un punto sobre el paraboloide.
\vspace{-2mm}
\begin{eqnarray*}
{\bf r}\left( u,v,w\right) &=& \left\langle u\cos (v) ,u\sin(v),(1-w) \sqrt{2-u^2}+wu^2\right \rangle ,
\end{eqnarray*}
%\vskip0.2cm
con $u\in [0,1],$ $v\in [0,2\pi ],$ $u\in [0,1].$

\vskip0.2cm
Con esta función hallamos los vectores:
\begin{eqnarray*}
{\bf r}_{u} &=& \langle \cos (v) ,\sin (v) ,-\frac{\left( 1-w\right) u}{\sqrt{2-u^2}}+2wu\rangle \\
{\bf r}_{v}&=& \langle -u\sin (v) ,u\cos (v) ,0\rangle \\
{\bf r}_{w} &=& \langle 0,0,-\sqrt{2-u^2}+u^2\rangle.
\end{eqnarray*}
%\vskip0.2cm
El elemento diferencial de volumen es
$\left| {\bf r}_{w}\cdot \left( {\bf r}_{u}\times {\bf r}_{v}\right) \right| = \left| u^{3}-u\sqrt{2-u^2}\right|$
cuya integral nos proporciona el volumen del sólido:
\vskip0.2cm
\[\int_{0}^{1}\int_{0}^{1}\int_{0}^{2\pi }\left| u^{3}-u\sqrt{2-u^2}\right|
dvdudw= \frac{4}{3}\pi \sqrt{2}-\frac{7}{6}\pi.\]
%\end{tcolorbox}
\vskip0.2cm
En coordenadas cilíndricas este valor se obtiene evaluando la siguiente integral
\[\int_{0}^{2\pi }\int_{0}^{1}\int_{r^2}^{\sqrt{2-r^2}}rdzdrd\theta
= \frac{1}{6}\pi \left( 8\sqrt{2}-7\right).\]
Para usar coordenadas esféricas, la esfera posee ecuación $\rho =%
\sqrt{2}$ y el paraboloide $\rho \cos \phi =\left( \rho \sin \phi \cos
\theta \right) ^2+\left( \rho \sin \phi \sin \theta \right) ^2$
que se simplifica como
$\rho =\frac{\cos \phi }{\sin ^2\left( \phi \right)}= \cot \phi \csc \phi ,$
como la intersección se da en $z=1,$ con una circunferencia de radio $1,$ entonces
$\phi =\pi /4.$ Es por lo anterior, que el volumen está dado por:
\begin{multline*}
\int_{0}^{\frac{\pi}{4}}\int_{0}^{2\pi }\int_{0}^{\sqrt{2}}\rho ^2\sin \phi
d\rho d\theta d\phi \\
+\int_{\frac{\pi}{4}}^{\frac{\pi}{2}}\int_{0}^{2\pi }\int_{0}^{\cot \phi \csc \phi }\rho ^2\sin \phi d\rho d\theta d\phi = \frac{\pi}{6} \left( 8\sqrt{2}-7\right).
\end{multline*}

\subsection*{Un plano y un paraboloide}
El plano y el paraboloide con ecuaciones $z=6-y$ y $z=x^2+y^2$,
determinan el sólido que se presenta a continuación:

\vskip0.2cm
\begin{figure}[H]
\centering
\caption{Sólido determinado por $x^2 + y^2 \leq z \leq 6 - y$.}
\label{fig:solido_paraboloide_plano}
\includegraphics[height=4.5cm]{Img/14d.pdf}
\includegraphics[height=4.5cm]{Img/14a.pdf}
\includegraphics[height=4.5cm]{Img/14b.pdf}
\fuentepropia
\end{figure}


En \verb"Mathematica", este gráfico se consigue usando las siguientes instrucciones:

\vspace{0.6em}
\noindent
\begin{tcolorbox}[breakable,enhanced,title=\bf Instrucción en \verb"Mathematica", top=2pt,bottom=2pt,boxsep=0pt,before skip=2pt,after skip=0pt]
\begin{Verbatim}[formatcom=\letraverbatim]
ContourPlot3D[{z==6-y,z==x^2+y^2},{x,-4,4},{y,-4,4},
{z,0,10},Mesh->False,ContourStyle->{Orange,Cyan}]
RegionPlot3D[x^2+y^2<z<6-y&&-3<y<2,{x,-5/2,5/2},
{y,-3,2},{z,0,9},PlotPoints->90,Mesh->None]
\end{Verbatim}
\end{tcolorbox}
\vspace{0.7em}
\noindent

\vskip0.2cm
\noindent
\begin{minipage}{\textwidth}
	\centering
	\captionof{figure}{Proyecciones del sólido en los planos coordenados.}
	\begin{minipage}{0.30\linewidth}
		\centering
		\vspace{-11mm}
		\includegraphics[width=\linewidth]{Fig/335.pdf}
	\end{minipage}
	\hspace{1em} 
	\begin{minipage}{0.30\linewidth}
		\centering
		\vspace{-11mm}
		\includegraphics[width=\linewidth]{Fig/336.pdf}
	\end{minipage}
	\hspace{1em} 
	\begin{minipage}{0.30\linewidth}
		\centering
		\vspace{-11mm}
		\includegraphics[width=\linewidth]{Fig/337.pdf}
	\end{minipage}
	
	\vspace{-8mm}
	\fuentepropia
\end{minipage}
\vskip0.2cm

Al sustituir $z$ de la ecuación del plano en la ecuación del paraboloide, se sigue que $6-y=x^2 +y^2,$
con lo cual $x=\pm \sqrt{6-y^2-y}$ o también $y=-\frac{1}{2}\pm \frac{1}{2}\sqrt{25-4x^2}.$
Al reemplazar $y$ en la ecuación del paraboloide, se sigue que
$z=x^2 +y^2=x^2 +\left( 6-z\right)^2,$ esta expresión se simplifica como
$x^2+z^2-13z=-36$ y se puede reescribir como
$x^2+\left(z-\frac{13}{2}\right)^2=\frac{25}{4}.$
\vskip0.2cm
Las expresiones anteriores permiten graficar la proyección del sólido en los planos coordenados, con el propósito de plantear adecuadamente las integrales necesarias para determinar el volumen ($V$) u otros elementos relativos al sólido. En este caso,

\vspace{-6mm}
\begin{eqnarray*}
V&=& \int {-3}^2\int _{-\sqrt{6-y-y^2}}^{\sqrt{6-y-y^2}} \int _{x^2+y^2}^{6-y}dzdxdy = \int _{-3}^2\int _{y^2}^{6-y} \int _{-\sqrt{z-y^2}}^{\sqrt{z-y^2}}dxdzdy \\
&=&\int 4^9\int _{-\sqrt{z}}^{6-z} \int _{-\sqrt{z-y^2}}^{\sqrt{z-y^2}}dxdydz+\int _0^4\int_{-\sqrt{z}}^{\sqrt{z}}\int _{-\sqrt{z-y^2}}^{\sqrt{z-y^2}}dxdydz\\
&=& \int_{-\frac{5}{2}}^{\frac{5}{2}} \int_{-\frac{1}{2}-\sqrt{\frac{25}{4}-x^2}}^{-\frac{1}{2}+\sqrt{\frac{25}{4}-x^2}}\int_{x^2+y^2}^{6-y}dzdydx = \frac{625\pi}{32} \end{eqnarray*}
\vskip0.4cm
La proyección en  el plano $xz$ evidencia que en el orden de integración $dydxdz$ o $dydzdx$ puede ser necesario evaluar tres o cuatro integrales triples, según el caso.
En seguida se muestra una forma simplificada de obtener el volumen.
\begin{multline*}
\bi _{-\frac{5}{2}}^{\frac{5}{2}}\bi_{\frac{13}{2}-\sqrt{\frac{25}{4}-x^2}}^{\frac{13}{2}+\sqrt{\frac{25}{4}-x^2}}\bi_{-\sqrt{z-x^2}}^{6-z}dydzdx \\Z +\bi_{-\sqrt{6}}^{\sqrt{6}}\bi_{x^2}^{\frac{13}{2}-\sqrt{\frac{25}{4}-x^2}}\bi_{-\sqrt{z-x^2}}^{\sqrt{z-x^2}}dydzdx = \frac{625\pi}{32}
\end{multline*}

En \verb"Mathematica", algunas de estas integrales pueden evaluarse mediante las siguientes instrucciones:

\vspace{0.6em}
\noindent
\begin{tcolorbox}[breakable,enhanced,title=\bf Instrucción en \verb"Mathematica",top=2pt,bottom=2pt,boxsep=0pt,before skip=2pt,after skip=0pt]
\begin{Verbatim}[formatcom=\letraverbatim]
Integrate[1,{y,-3,2},{x,-Sqrt[6-y-y^2],Sqrt[6-y-y^2]},
{z,x^2+y^2,6-y}]
Integrate[1,{x,-5/2,5/2},{y,-1/2-Sqrt[25/4-x^2],
1/2+Sqrt[25/4-x^2]},{z,x^2+y^2,6-y}]
Integrate[1,{y,-3,2},{z,y^2,6-y},{x,-Sqrt[z-y^2],
Sqrt[z-y^2]}]
\end{Verbatim}
\end{tcolorbox}
\vspace{0.7em}
\noindent

Ahora se parametriza la región mencionada antes.

%\begin{tcolorbox}[title=\bfseries Parametrizando el sólido]
La base del sólido en el plano $xy$ está dada por
$x^2+y^2=6-y$ que equivale a $x^2+\left( y+\frac{1}{2}\right)^2=\frac{25}{4}.$
Esta región se parametriza al hacer $x=\frac{5}{2}u\cos(t)$, $y=-\frac{1}{2}+\frac{5}{2}u\sin(t)$
con $u\in [0,1],$ $t\in [0,2\pi].$ Un punto sobre el plano $z=6-y,$ tiene la forma
\vskip0.2cm
\[P\left(\frac{5}{2}\cos(t),-\frac{1}{2}+\frac{5}{2}\sin(t),6-(-\frac{1}{2}+\frac{5}{2}\sin(t))\right)\]
\vskip0.2cm
y un punto arbitrario sobre el paraboloide es de  forma
\vskip0.2cm
\[Q\left(\frac{5}{2}\cos(t),-\frac{1}{2}+\frac{5}{2}\sin(t),(\frac{5}{2}\cos(t))^2+(-\frac{1}{2}+\frac{5}{2}\sin(t))^2\right).\]
Formando segmentos rectilíneos que unan un punto del plano con un punto sobre el paraboloide se parametriza el sólido. Esto es, simplificando la expresión $(1-w)P+wQ,$ se define la función:
%\vskip0.2cm
\[{\bf r}\left(t,u,w\right)=\frac{5}{2} \left\langle u \cos (t),u\sin(t)-\frac{1}{5},\frac{5w}{2}(u^2-1)-u\sin(t)+\frac{13}{5} \right\rangle.\]
%\vskip0.2cm
con $u\in [0,1],$ $t\in [0,2\pi],$ $w\in [0,1],$
la cual parametriza el sólido. También se hallan los vectores:
{\small
\begin{eqnarray*}
{\bf r}_{t} &=& \frac{5}{2}\left\langle - u \sin (t), u \cos (t),- u \cos (t) \right\rangle \\
{\bf r}_{u} &=& \frac{5}{2}\left\langle \cos (t), \sin (t),5 u w-\sin (t)\right\rangle \\
{\bf r}_{w} &=& \frac{25}{4}\left\langle 0,0, \left(u^2-1\right) \right\rangle
\end{eqnarray*}
}
%\vskip0.2cm
El elemento diferencial de volumen es
$\left|{\bf r}_{t}\cdot\left({\bf r}_{u}\times {\bf r}_{w}\right)\right| =\frac{625}{16} u \left(1-u^2\right).$
Por último, evaluamos la integral que determina el volumen del sólido:
\[ \int_{0}^{1}\int_{0}^{2\pi}\int_{0}^{1} \left|{\bf r}_{t}\cdot\left({\bf r}_{u}\times {\bf r}_{w}\right)\right|dudtdw=\frac{625\pi}{32}\]
%\end{tcolorbox}
%\vskip0.2cm
El procedimiento descrito antes, se puede efectuar mediante la ejecución de las siguientes instrucciones: % La parametrización de la base se denomina \verb"A" y también define la función \verb"r" que parametriza el sólido. El producto triple se denomina \verb"B" que se integra en los límites establecidos.

\vspace{0.6em}
\noindent
\begin{tcolorbox}[breakable,enhanced,title=\bf Instrucción en \verb"Mathematica", top=2pt,bottom=2pt,boxsep=0pt,before skip=2pt,after skip=0pt]
\begin{Verbatim}[formatcom=\letraverbatim]
{x,y}={5/2*u*Cos[t],-1/2+5/2*u*Sin[t]};
r=Simplify[w*{x,y,x^2+y^2}+(1-w)*{x,y,6-y}];
B=Abs[Dot[D[r,t],Cross[D[r,u],D[r,w]]]]
Integrate[B,{t,0,2*Pi},{u,0,1},{w,0,1}]
\end{Verbatim}
\end{tcolorbox}
\vspace{0.7em}
\noindent

\subsection*{Un paraboloide, un cilindro y un plano}

Consideremos el sólido que se halla acotado bajo el paraboloide con ecuación $z=4-x^2-y^2$,
sobre el plano $z=0$ y fuera del cilindro con ecuación $x^2+y^2=2x$.
\vskip0.4cm
\vspace{-0.8mm}
\noindent
\begin{minipage}{\textwidth}
	\centering
	\captionof{figure}{Sólido determinado por $0\leq z\leq 4-x^2-y^2$, $2x \leq x^2+y^2$.}
	\label{fig:solido}
	\begin{minipage}{0.21\linewidth}
		\centering
		\vspace{-12mm}
		\includegraphics[width=\linewidth]{Img/3a.pdf}
	\end{minipage}
	\hspace{1em}
	\begin{minipage}{0.27\linewidth}
		\centering
		\vspace{-12mm}
		\includegraphics[width=\linewidth]{Fig/338.pdf}
	\end{minipage}
	\hspace{1em}
	\begin{minipage}{0.21\linewidth}
		\centering
		\vspace{-12mm}
		\includegraphics[width=\linewidth]{Img/3c.pdf}
	\end{minipage}
	
	\vspace{-8mm}
	\fuentepropia
\end{minipage}


\vskip0.2cm

Las instrucciones para obtener la figura anterior son las siguientes:

\vspace{0.6em}
\noindent
\begin{tcolorbox}[breakable,enhanced,title=\bf Instrucción en \verb"Mathematica", top=2pt,bottom=2pt,boxsep=0pt,before skip=2pt,after skip=0pt]
\begin{Verbatim}[formatcom=\letraverbatim]
RegionPlot3D[0<z<4-x^2-y^2&&x^2+y^2>2*x-0.05,
{x,-2,2},{y,-2,2},{z,0,4},Mesh->False,
PlotStyle->Orange,PlotPoints->90]
\end{Verbatim}
\end{tcolorbox}
\vspace{0.7em}
\noindent

La proyección de este sólido en el plano $xy$ es la siguiente:
\vskip0.2cm

\noindent
\begin{minipage}{\textwidth}
	\centering
	\captionof{figure}{Región $2x \leq x^2+y^2 \leq 4$.}
	\begin{minipage}{0.31\linewidth}
		\centering
		\vspace{-14mm}
		\includegraphics[width=\linewidth]{Fig/339.pdf}
	\end{minipage}
	\hspace{1em}
	\begin{minipage}{0.31\linewidth}
		\centering
		\vspace{-14mm}
		\includegraphics[width=\linewidth]{Fig/340.pdf}
	\end{minipage}
	
	\vspace{-8mm}
	\fuentepropia
\end{minipage}

\vspace{0.2cm} 

El relleno de esta región sugiere varios elementos: los límites de integración, en el orden $dzdydx$ se evalúa mediante tres integrales triples y en el orden $dzdxdy$ es necesario evaluar cuatro integrales triples. El volumen está dado por:

\vspace{-4mm}
{\small
\begin{multline*}
\int_{-2}^{0}\int_{-\sqrt{4-x^2}}^{\sqrt{4-x^2}}\int_{0}^{4-x^2-y^2}dzdydx+\int_{0}^2\int_{\sqrt{2x-x^2}}^{\sqrt{4-x^2}}\int_{0}^{4-x^2-y^2}dzdydx \\
+\int_{0}^2\int_{-\sqrt{4-x^2}}^{-\sqrt{2x-x^2}}\int_{0}^{4-x^2-y^2}dzdydx= \frac{11}{2}\pi
\end{multline*}
}
\vspace{-4mm}
{\small
\begin{multline*}
\int_{1}^2\int_{-\sqrt{4-y^2}}^{\sqrt{4-y^2}}\int_{0}^{4-x^2-y^2}dzdxdy
+\int_{-2}^{-1}\int_{-\sqrt{4-y^2}}^{\sqrt{4-y^2}}\int_{0}^{4-x^2-y^2}dzdxdy+ \\
\int_{-1}^{1}\int_{-\sqrt{4-y^2}}^{1-\sqrt{1-y^2}}\int_{0}^{4-x^2-y^2}dzdxdy
+\int_{-1}^{1}\int_{1+\sqrt{1-y^2}}^{\sqrt{4-y^2}}\int_{0}^{4-x^2-y^2}dzdxdy= \frac{11}{2}\pi
\end{multline*}
}

Estas integrales se evalúan en \verb"Mathematica" con las siguientes instrucciones

\vspace{0.6em}
\noindent
\begin{tcolorbox}[breakable,enhanced,title=\bf Instrucción en \verb"Mathematica", top=2pt,bottom=2pt,boxsep=0pt,before skip=2pt,after skip=0pt]
\begin{Verbatim}[formatcom=\letraverbatim]
Integrate[1,{x,-2,0},{y,-Sqrt[4-x^2],Sqrt[4-x^2]},
{z,0,4-x^2-y^2}]+Integrate[1,{x,0,2},{y,Sqrt[2*x-x^2],
Sqrt[4-x^2]},{z,0,4-x^2-y^2}]+Integrate[1,{x,0,2},
{y,-Sqrt[4-x^2],-Sqrt[2*x-x^2]},{z,0,4-x^2-y^2}]
Integrate[1,{y,1,2},{x,-Sqrt[4-y^2],Sqrt[4-y^2]},
{z,0,4-x^2-y^2}]+Integrate[1,{y,-2,-1},{x,-Sqrt[4-y^2],
Sqrt[4-y^2]},{z,0,4-x^2-y^2}]+Integrate[1,{y,-1,1},
{x,-Sqrt[4-y^2],1-Sqrt[1-y^2]},{z,0,4-x^2-y^2}]+
Integrate[1,{y,-1,1},{x,1+Sqrt[1-y^2],Sqrt[4-y^2]},
{z,0,4-x^2-y^2}]
\end{Verbatim}
\end{tcolorbox}
\vspace{0.7em}
\noindent

	
Integrando primero con respecto a $y,$ se puede plantear el volumen como la integral triple en donde $y$ va desde la parte negativa del paraboloide hasta la parte positiva del mismo y luego restar la parte correspondiente a la región que ocupa el cilindro, los límites para $x$ y $z$ se obtienen de la proyección del sólido en el plano $xz$.


\vskip0.3cm
\noindent
\begin{minipage}{\textwidth}
	\centering
	\captionof{figure}{$-\sqrt{4-z}\leq x \leq 2-z/2$}
	\begin{minipage}{0.39\linewidth}
		\centering
		\vspace{-17mm}
		\includegraphics[width=\linewidth]{Fig/341.pdf}
	\end{minipage}
	
	\vspace{-12mm}
	\fuentepropia
\end{minipage}
\vskip0.2cm


Con ayuda de este gráfico, el volumen se puede plantear mediante las siguientes integrales
\[\int_{0}^{4}\int_{-\sqrt{4-z}}^{\frac{4-z}{2}}\int_{-\sqrt{4-z-x^2}}^{\sqrt{4-z-x^2}}dydxdz
-\int_{0}^{4}\int_{0}^{\frac{4-z}{2}}\int_{-\sqrt{2x-x^2}}^{\sqrt{2x-x^2}}dydxdz \approx 17.2788.\]

\begin{multline*}
\int_{-2}^{0}\int_{0}^{4-x^2}\int_{-\sqrt{4-z-x^2}}^{\sqrt{4-z-x^2}}dydzdx
+\int_{0}^2\int_{0}^{4-2x}\int_{-\sqrt{4-z-x^2}}^{\sqrt{4-z-x^2}}dydzdx \\
-\int_{0}^2\int_{0}^{4-2x}\int_{-\sqrt{2x-x^2}}^{\sqrt{2x-x^2}}dydzdx = \frac{11}{2}\pi.
\end{multline*}
También se puede hacer uso de las coordenadas cilíndricas, cuyo planteamiento es el siguiente:
\[2 \int_{0}^{\pi /2}\int_{2\cos \theta}^2\int_{0}^{4-r^2}rdzdrd\theta +\int_{\pi /2}^{3\pi
/2}\int_{0}^2\int_{0}^{4-r^2}rdzdrd\theta = \frac{11}{2}\pi. \]

Las integrales anteriores se evaluaron con las siguientes instrucciones:

\vspace{0.6em}
\noindent
\begin{tcolorbox}[breakable,enhanced,title=\bf Instrucción en \verb"Mathematica", top=2pt,bottom=2pt,boxsep=0pt,before skip=2pt,after skip=0pt]
\begin{Verbatim}[formatcom=\letraverbatim]
Integrate[1,{z,0,4},{x,-Sqrt[4-z],(4-z)/2},
{y,-Sqrt[4-z-x^2],Sqrt[4-z-x^2]}]-Integrate[1,
{z,0,4},{x,0,(4-z)/2},{y,-Sqrt[2x-x^2],Sqrt[2x-x^2]}]
		
Integrate[1,{x,-2,0},{z,0,4-x^2},{y,-Sqrt[4-z-x^2],
Sqrt[4-z-x^2]}]+Integrate[1,{x,0,2},{z,0,4-2*x},
{y,-Sqrt[4-z-x^2],Sqrt[4-z-x^2]}]-Integrate[1,{x,0,2},
{z,0,4-2x},{y,-Sqrt[2x-x^2],Sqrt[2x-x^2]}]
		
2*Integrate[r,{\[Theta],0,Pi/2},{r,2*Cos[\[Theta]],2},
{z,0,4-r^2}]+Integrate[r,{\[Theta],Pi/2,3*Pi/2},
{r,0,2},{z,0,4-r^2}]
\end{Verbatim}
\end{tcolorbox}
\vspace{0.7em}
\noindent


\begin{ejer}
Plantear las integrales que se deben evaluar para determinar el volumen del sólido determinado por las expresiones $z\leq 4-x^2-y^2$, $x^2+y^2 \geq 2x$ y $z \geq 0$;
en las órdenes $dxdydz$ y $dxdzdy.$
\end{ejer}
\vskip0.2cm
El mismo valor obtenido antes se puede hallar de la siguiente manera. 
\vskip0.2cm

%\begin{tcolorbox}[title=\bfseries Parametrizando el sólido]
%\small
%\parskip=0pt
%\setlength{\belowdisplayskip}{0pt}
%\setlength{\abovedisplayskip}{0pt}
%\noindent
%\vspace{-20pt}

La proyección del sólido en el plano $xy$ permite ver que los puntos sobre la circunferencia interna puede representarse en la forma
$(1+\cos(t),\sin(t)),$ mientras que los puntos sobre la circunferencia externa son de la forma
$(2\cos(t),2\sin(t))$ para $t\in [0,2\pi]$.


\vskip0.2cm
\noindent
\begin{minipage}{\textwidth}
	\centering
	\captionof{figure}{Región $2x \leq x^2+y^2 \leq 4$}
	\label{fig:curva_parametrica_uv_3_2}
	\begin{minipage}{0.35\linewidth}
		\centering
		\vspace{-18mm}
		\includegraphics[width=\linewidth]{Fig/342.pdf}
	\end{minipage}
	
	\vspace{-10mm}
	\fuentepropia
\end{minipage}
\vskip0.2cm

Formando segmentos rectilíneos que unen las dos circunferencias, se parametriza la región entre las dos curvas, esto se logra simplificando
$\left( 1-u\right)\left(1+\cos(t) ,\sin(t)\right) +u\left(2\cos(t),2\sin(t)\right),$ para $u\in [0,1].$
Consideremos un punto
{\small $P\left( 1-u+(u+1)\cos (t),(1+u)\sin (t),0\right)$} en la base del sólido
y {\small $Q\left( 1-u+(u+1)\cos (t),(1+u)\sin (t),2(1-u^2)(1-\cos(t))\right)$} un \\ punto sobre el paraboloide.
Con lo anterior se pueden formar segmentos rectilíneos (paralelos al eje $z$) que unan la base del sólido con puntos sobre el paraboloide. Esto es $\left(1-v\right) P+vQ$; simplificando, obtendremos una parametrización para el sólido.
\begin{multline*}
{\bf r}(t,u,v) =\langle 1+(1+u) \cos t-u,( 1+u)\sin(t),
2v(1-u^2) ( 1-\cos t) \rangle,
\end{multline*}
\vskip0.2cm
donde $u\in [0,1],\,\, t\in [0,2\pi ],\,\, v\in [0,1].$ También puede verse que:
%\vspace{-4mm}
\begin{eqnarray*}
{\bf r}_{t}&=& \langle - \left( 1+u\right)\sin (t),\left( 1+u\right) \cos t,2v\left( 1-u^2\right) \sin t \rangle \\
{\bf r}_{u}&=& \langle -1+\cos t,\sin t,-4vu\left( 1-\cos t\right)\rangle\\
{\bf r}_{v}&=& \langle 0,0,2\left( 1-u^2\right) \left( 1-\cos t\right) \rangle.
\end{eqnarray*}
\vskip0.2cm
Con lo anterior,
${\bf r}_{u}\times {\bf r}_{v} = \big \langle 2 \left(1-u^2\right) \left( 1-\cos t\right)\sin t,2(1-\cos t) ^2(1-u^2) ,0\big \rangle .$
Además
\[\left|{\bf r}_{t}\cdot \left( {\bf r}_{u}\times {\bf r}_{v}\right)\right| = \left| 2\left( 1-\cos t\right)
^2\left( 1-u\right) \left( 1+u\right) ^2\right|,\]
\vskip0.2cm
con lo cual, el volumen del sólido esta dado por
\vskip0.2cm
\[\int_{0}^{1}\int_{0}^{1}\int_{0}^{2\pi }\left| 2\left( 1-\cos t\right)
^2\left( 1-u\right) \left( 1+u\right) ^2\right| dtdudv= \frac{11}{2}\pi.\]
%\end{tcolorbox}

\subsection*{Un paraboloide hueco
} \index{Paraboloide!hueco}
Consideraremos el sólido que se presenta en la sección \secref{p1}.

\vskip0.5cm
\noindent
\begin{minipage}{\textwidth}
	\centering
	\captionof{figure}{Región determinada por $2x \leq x^2 + y^2$ y $0 \leq z \leq 9 - x^2 - y^2$.}
	\label{fig:region_paraboloide_desplazada}
	\begin{minipage}{0.25\linewidth}
		\centering
		\includegraphics[width=\linewidth]{Img/79a.pdf}
	\end{minipage}
	\hspace{1em} 
	\begin{minipage}{0.25\linewidth}
		\centering
		\includegraphics[width=\linewidth]{Img/79b.pdf}
	\end{minipage}
	\hspace{1em} 
	\begin{minipage}{0.25\linewidth}
		\centering
		\includegraphics[width=\linewidth]{Img/79c.pdf}
	\end{minipage}
	
	\vspace{0.3em}
	\fuentepropia
\end{minipage}
\vskip0.2cm

La proyección del sólido en el plano $xy$ es una región acotada por las circunferencias
${\bf r}_1(t)= \langle 1+\cos(2t),\sin(2t)\rangle $
y ${\bf r}_2(t)= \langle 3\cos(2t),3\sin(2t)\rangle,$
que se parametriza por medio de una combinación de la forma $(1-u){\bf r}_1(t)+u{\bf r}_2(t),$ simplificando para $t\in [0,\pi]$, $u\in [0,1],$ se obtiene:
\[\langle 1-u+(1+2u)\cos (2t) ,(1+2u)\sin(2t)\rangle.\]


\vskip0.2cm
\noindent
\begin{minipage}{\textwidth}
	\centering
	\captionof{figure}{${\bf r}(t,u)$}
	\label{fig:curva_parametrica_uv_4}
	\begin{minipage}{0.34\linewidth}
		\centering
		\vspace{-18mm}
		\includegraphics[width=\linewidth]{Fig/343.pdf}
	\end{minipage}
	
	\vspace{-10mm}
	\fuentepropia
\end{minipage}
\vskip0.2cm


El sólido se parametriza mediante la función
{\small
\begin{eqnarray*}
{\bf r}\left( t,u,v\right) &=& \bigg\langle 1-u+\left( 1+2u\right) \cos (2t) ,\left( 1+2u\right)
\sin (2t) , \\ &&
\hskip2cm \left( 7-2u-5u^2+\left( 4u^2-2u-2\right) \cos(2t)\right) v\bigg\rangle,
\end{eqnarray*}
}

en donde $t\in [0,\pi]$, $u\in [0,1]$ y $v \in [0,1]$.
Con la función ${\bf r}$ hallamos los siguientes vectores:

\vspace{-4mm}
{\small
\begin{eqnarray*}
{\bf r}_{t} &=& \left\langle -2\left( 2u+1\right) \sin (2t)
,2\left( 2u+1\right) \cos (2t) ,-4v(2u^2-u-1)\sin (2t) \right\rangle \\
{\bf r}_{u} &=& \left\langle -1+2\cos (2t) ,2\sin \left(2t\right) ,v\left( -2-10u+\left( -2+8u\right) \cos (2t) \right) \right\rangle \\
{\bf r}_{v} &=& \left\langle 0,0,7-2u-5u^2+4\cos (2t)u^2-2\cos (2t) u-2\cos (2t) \right\rangle,
\end{eqnarray*}
}

\vskip0.2cm
luego,
$\left| {\bf r}_{v}\cdot ({\bf r}_{t}\times {\bf r}_{u}) \right| =\bigg| 2(\cos (2t)-2)(2u+1)(u-1) ((4u+2)\cos (2t) \linebreak -5u-7) \bigg| ,$
que corresponde al elemento diferencial de volumen, su integral en los
límites establecidos nos proporciona el volumen del sólido
\[\int_{0}^{1}\int_{0}^{1}\int_{0}^{\pi }\left| {\bf r}_{v}\cdot ({\bf r}_{t}\times {\bf r}_{u}) \right| dtdudv=33\pi.\]
\vskip0.2cm
Este cálculo se hizo en \verb"Mathematica" usando las siguientes instrucciones:

\vspace{0.6em}
\noindent
\begin{tcolorbox}[breakable,enhanced,title=\bf Instrucción en \verb"Mathematica", top=2pt,bottom=2pt,boxsep=0pt,before skip=2pt,after skip=0pt]
\begin{Verbatim}[formatcom=\letraverbatim]
{x,y}=(1-u)*{3*Cos[2*t],3*Sin[2*t]}+u*{1+Cos[2*t],
Sin[2*t]};
z=Simplify[9-x^2-y^2];
r={x,y,z*v}
A=Simplify[Dot[D[r,v],Cross[D[r,t],D[r,u]]]];
Integrate[A,{t,0,Pi},{u,0,1},{v,0,1}]
\end{Verbatim}
\end{tcolorbox}
\vspace{0.7em}
\noindent


En coordenadas rectangulares, el volumen está dado por
\begin{multline*} \int_{-3}^{3}\int_{-\sqrt{9-x^2}}^{\sqrt{9-x^2}}\left(
9-x^2-y^2\right) dydx- \\ 2\int_{0}^2\int_{0}^{\sqrt{2x-x^2}}\left(9-x^2-y^2\right)dydx= 33\pi,
\end{multline*}
y usando coordenadas polares, el planteamiento es el siguiente:
\[\int_{0}^{2\pi }\int_{0}^{3}\left( 9-r^2\right) rdrd\theta -\int_{0}^{\pi}\int_{0}^{2\cos \theta }\left( 9-r^2\right) rdrd\theta = 33\pi.\]
Estas integrales se evalúan con la siguiente rutina:

\vspace{0.6em}
\noindent
\begin{tcolorbox}[breakable,enhanced,title=\bf Instrucción en \verb"Mathematica", top=2pt,bottom=2pt,boxsep=0pt,before skip=2pt,after skip=0pt]
\begin{Verbatim}[formatcom=\letraverbatim]
Integrate[9-x^2-y^2,{x,-3,3},{y,-Sqrt[9-x^2],
Sqrt[9-x^2]}]-2*Integrate[9-x^2-y^2,{x,0,2},
{y,0,Sqrt[2*x-x^2]}]
Integrate[(9-r^2)*r,{t,0,2*Pi},{r,0,3}]-
Integrate[(9-r^2)*r,{t,0,Pi},{r,0,2*Cos[t]}]
\end{Verbatim}
\end{tcolorbox}
\vspace{0.7em}
\noindent


\subsection*{Un paraboloide y un cilindro parabólico}
Consideremos el sólido acotado por las figuras de las ecuaciones $z = x^2+y^2$ y
$z = 5 - y^2,$ que representan respectivamente un paraboloide y un cilindro parabólico; estas figuras se muestran enseguida.

\begin{figure}[H]
\centering
\caption{Sólido determinado por $x^2 + y^2 \leq z \leq 5 - y^2$.}
\label{fig:solido_paraboloide_plano_curvo}
\includegraphics[height=4.5cm]{Img/26a.pdf}
\includegraphics[height=4.5cm]{Img/26b.pdf}
\includegraphics[height=4.5cm]{Img/26c.pdf}
\fuentepropia
\end{figure}

Las instrucciones necesarias para obtenerlas son las siguientes.

\vspace{0.6em}
\noindent
\begin{tcolorbox}[breakable,enhanced,title=\bf Instrucción en \verb"Mathematica", top=2pt,bottom=2pt,boxsep=0pt,before skip=2pt,after skip=0pt]
\begin{Verbatim}[formatcom=\letraverbatim]
ContourPlot3D[{z==5-y^2,z==x^2+y^2},{x,-2.5,2.5},
{y,-2.5,2.5},{z,0,5},PlotPoints->90,
ContourStyle->{RGBColor[0.95,0.6,0.9],Cyan}]
RegionPlot3D[x^2+y^2<z<5-y^2&&-Sqrt[5-2*y^2]<x<
Sqrt[5-2*y^2],{x,-2.5,2.5},{y,-2.5,2.5},{z,0,5},
PlotPoints->90,Mesh->False]
\end{Verbatim}
\end{tcolorbox}
\vspace{0.7em}
\noindent


El sólido se proyecta en los planos coordenados como se muestra en las siguientes figuras.
\vskip0.2cm

\noindent
\begin{minipage}{\textwidth}
	\centering
	\captionof{figure}{Proyecciones del sólido en los planos coordenados.}
	\begin{minipage}{0.30\linewidth}
		\centering
		\vspace{-8mm}
		\includegraphics[width=\linewidth]{Fig/344.pdf}
	\end{minipage}
	\hspace{1em}
	\begin{minipage}{0.30\linewidth}
		\centering
		\vspace{-8mm}
		\includegraphics[width=\linewidth]{Fig/345.pdf}
	\end{minipage}
	\hspace{1em}
	\begin{minipage}{0.30\linewidth}
		\centering
		\vspace{-8mm}
		\includegraphics[width=\linewidth]{Fig/346.pdf}
	\end{minipage}
	
	\vspace{0.3em}
	
	\vspace{-8mm}
	\fuentepropia
\end{minipage}

\vskip0.2cm
Eliminando $z$ de las ecuaciones $z=x^2+y^2$ y $z=5-y^2$, se obtiene $x^2+2y^2=5;$
esta expresión en el plano $xy$ representa una elipse.
Para integrar primero con respecto a $z$ se debe tener en cuenta que $z$ recorre los puntos que se hallan entre el paraboloide y el cilindro, los límites en $y$ y $x$ se pueden obtener de la proyección en el plano $xy.$ Por tanto, el volumen del sólido está dado por las siguientes integrales:
\vspace{-2mm}
\begin{multline*} \int_{-\sqrt{5/2}}^{\sqrt{5/2}}\int_{-\sqrt{5-2y^2}}^{\sqrt{5-2y^2}}\int_{x^2+y^2}^{5-y^2}dzdxdy = \\ \int_{-\sqrt{5}}^{\sqrt{5}}\int_{-\frac{1}{2}\sqrt{10-2x^2}}^{\frac{1}{2}\sqrt{10-2x^2}}\int_{x^2+y^2}^{5-y^2}dzdydx= \frac{25\sqrt{2}\pi}{4}.
\end{multline*}
En el plano $xz,$ las curvas que describen la región de integración se determinan por las ecuaciones $2z=5+x^2,$ $z=x^2$ y $z=5.$
De la expresión $2z=5+x^2$ se sigue que $ x=\pm \sqrt{2z-5}.$ Por lo tanto, el volumen es
\[\int_{-\sqrt{5}}^{\sqrt{5}}\int_{x^2}^{\frac{5+x^2}{2}}\int_{-\sqrt{z-x^2}}^{\sqrt{z-x^2}}dydzdx+\int_{-\sqrt{5}}^{\sqrt{5}}\int_{\frac{5+x^2}{2}}^{5}\int_{-\sqrt{5-z}}^{\sqrt{5-z}}dydzdx= \frac{25}{4}\sqrt{2}\pi.\]
También se puede calcular el volumen de la siguiente manera
\begin{multline*}
\int_{5/2}^{5}\int_{-\sqrt{2z-5}}^{\sqrt{2z-5}}\int_{-\sqrt{5-z}}^{\sqrt{5-z}}dydxdz+\int_{5/2}^{5}\int_{-\sqrt{z}}^{-\sqrt{2z-5}}\int_{-\sqrt{z-x^2}}^{\sqrt{z-x^2}}dydxdz\\
+\int_{5/2}^{5}\int_{\sqrt{2z-5}}^{\sqrt{z}}\int_{-\sqrt{z-x^2}}^{\sqrt{z-x^2}}dydxdz+\int_{0}^{5/2}\int_{-\sqrt{z}}^{\sqrt{z}}\int_{-\sqrt{z-x^2}}^{\sqrt{z-x^2}}dydxdz\approx 27.768.
\end{multline*}

Al integrar primero con respecto a $x$ hay que tener presente que esta recorre los puntos que se encuentran entre la parte negativa del paraboloide (despejando $x$)
hasta la parte positiva de la misma superficie.
\vspace{-2mm}
\[
\displaystyle
\int_{0}^{5/2}
\int_{-\sqrt{z}}^{\sqrt{z}}
\int_{-\sqrt{z - y^2}}^{\sqrt{z - y^2}}
dx\,dy\,dz
+
\int_{5/2}^{5}
\int_{-\sqrt{5 - z}}^{\sqrt{5 - z}}
\int_{-\sqrt{z - y^2}}^{\sqrt{z - y^2}}
dx\,dy\,dz
\approx 27.768
\]

\[
\displaystyle
\int_{-\sqrt{5/2}}^{\sqrt{5/2}}
\int_{y^2}^{5 - y^2}
\int_{-\sqrt{z - y^2}}^{\sqrt{z - y^2}}
dx\,dz\,dy
= \frac{25}{4}\sqrt{2}\pi.
\]

\vskip0.2cm
Utilizando coordenadas cilíndricas las ecuaciones del paraboloide y el cilindro se reescriben como sigue
$z=r^2$ y $z=5-r^2\sin^2(\theta)$. La ecuación de la elipse de intersección de estas superficies está dada por
$ r^2 (3- \cos(2 \theta))=10.$
Es por esto que la integral que determina el volumen del sólido es
\[\int_0^{2\pi}\int_0^{\frac{\sqrt{10}}{\sqrt{3-\cos(2\theta)}}}\int_{r^2}^{5-r^2\sin ^2(\theta)}rdzdrd\theta=\frac{25\pi}{2\sqrt{2}}.\]
Algunas de estas integrales se evaluaron usando \verb"Mathematica" con las siguientes instrucciones:

\vspace{0.6em}
\noindent
\begin{tcolorbox}[breakable,enhanced,title=\bf Instrucción en \verb"Mathematica", top=2pt,bottom=2pt,boxsep=0pt,before skip=2pt,after skip=0pt]
\begin{Verbatim}[formatcom=\letraverbatim]
Integrate[1,{x,-Sqrt[5],Sqrt[5]},{y,-Sqrt[10-2*x^2]/2,
Sqrt[10-2*x^2]/2},{z,x^2+y^2,5-y^2}]
Integrate[1,{y,-Sqrt[5/2],Sqrt[5/2]},{x,-Sqrt[5-2y^2],
Sqrt[5-2y^2]},{z,x^2+y^2,5-y^2}]
Integrate[1,{y,-Sqrt[5/2],Sqrt[5/2]},{z,y^2,5-y^2},
{x,-Sqrt[z-y^2],Sqrt[z-y^2]}]
\end{Verbatim}
\end{tcolorbox}
\vspace{0.7em}
\noindent



\begin{comment}

\begin{verbatim}
Integrate[1,{x,-Sqrt[5],Sqrt[5]},{y,-Sqrt[10-2*x^2]/2,
Sqrt[10-2*x^2]/2},{z,x^2+y^2,5-y^2}]

Integrate[1,{y,-Sqrt[5/2],Sqrt[5/2]},{x,-Sqrt[5-2y^2],
Sqrt[5-2y^2]},{z,x^2+y^2,5-y^2}]

Integrate[1,{x,-Sqrt[5],Sqrt[5]},{z,x^2,(5+x^2)/2},
{y,-Sqrt[z-x^2],Sqrt[z-x^2]}]+Integrate[1,{x,-Sqrt[5],
Sqrt[5]},{z,(5+x^2)/2,5},{y,-Sqrt[5-z],Sqrt[5-z]}]

Integrate[1,{y,-Sqrt[5/2],Sqrt[5/2]},{z,y^2,5-y^2},
{x,-Sqrt[z-y^2],Sqrt[z-y^2]}]

Integrate[1,{z,5/2,5},{x,-Sqrt[2*z-5],Sqrt[2*z-5]},
{y,-Sqrt[5-z],Sqrt[5-z]}]+Integrate[1,{z,5/2,5},
{x,-Sqrt[z],-Sqrt[2z-5]},{y,-Sqrt[z-x^2],Sqrt[z-x^2]}]
+ Integrate[1,{z,5/2,5},{x,Sqrt[2z-5],Sqrt[z]},
{y,-Sqrt[z-x^2],Sqrt[z-x^2]}]+Integrate[1,{z,0,5/2},
{x,-Sqrt[z],Sqrt[z]},{y,-Sqrt[z-x^2],Sqrt[z-x^2]}]

Integrate[1,{z,0,5/2},{y,-Sqrt[z],Sqrt[z]},
{x,-Sqrt[z-y^2],Sqrt[z-y^2]}]+Integrate[1,{z,5/2,5},
{y,-Sqrt[5-z],Sqrt[5-z]},{x,-Sqrt[z-y^2],Sqrt[z-y^2]}]
\end{verbatim}

\verb"Integrate[r,{\[Theta],0,2*Pi},{r,0,Sqrt[10]/Sqrt[3-"
\verb"Cos[2\[Theta]]]},{z,r^2,5-r^2*Sin[\[Theta]]^2}]"
\end{comment}

\vskip0.2cm
Otra manera de hallar el volumen se presenta enseguida.
\vskip0.2cm

%\begin{tcolorbox}[title=\bfseries Parametrizando el sólido]
En el plano $xy$, la proyección del sólido se determina por
$x^2+2y^2=5,$ esta ecuación corresponde a elipse. Esta curva junto con su interior se parametrizan como {\small $x=\sqrt{5}u\cos t,$}
{\small $y=\sqrt{\frac{5}{2}}u\sin t),$} para {\small $t\in [0,2\pi ],$} {\small $u\in [0,1].$} 

\vskip0.2cm
La coordenada $z$ está dada por {\small $z = \frac{5}{2}u^2\cos ^2t+\frac{5}{2}u^2$} para el paraboloide , {\small $z=5-\frac{5}{2}u^2\sin ^2t$} para el cilindro.
Con un punto $P$ sobre el paraboloide y un punto $Q$ sobre el cilindro, se forma una combinación convexa de la forma $(1-w)P+wQ$ que parametriza el sólido, esto es,

\vspace{-5mm}
\[{\bf r}(t,u,w)= \langle \sqrt{5}u\cos t,\sqrt{\frac{5}{2}}u\sin t,\frac{5}{2}u^2( \cos ^2t+1) -5w(u^2-1)\rangle.\]

En donde $t \in [0,1],$ $u \in [0,1]$ y $w \in [0,1].$ Con esta función se hallan los siguientes vectores:
\vspace{-4mm}
\begin{eqnarray*}
{\bf r}_{t} & =& \langle -\sqrt{5}u\sin t,\frac{1}{2}\sqrt{10}u\cos t,-5u^2\cos t\sin t\rangle \\
{\bf r}_{u} &=& \langle \sqrt{5}\cos t,\frac{1}{2}\sqrt{10}\sin t,5u(\cos^2t+1)-10wu\rangle \\
{\bf r}_{w} &=& \langle0,0,-5u^2+5\rangle
\end{eqnarray*}


Con lo anterior, se calcula el producto vectorial triple que corresponde al elemento diferencial de volumen.
$| {\bf r}_{t}\cdot ( {\bf r}_{u}\times {\bf r}_{w}) | =\frac{25}{2}\sqrt{2} \left| u( u^2-1)\right|,$
que al integrarlo, nos produce el volumen del sólido mencionado
\[\frac{25\sqrt{2}}{2} \int_{0}^{1}\int_{0}^{1}\int_{0}^{2\pi } | u(u^2-1) | dtdudw= \frac{25}{4}\sqrt{2}\pi.\]
%\end{tcolorbox}

De manera más breve, \verb"Mathematica" permite hacer estos cálculos con la ejecución de las siguientes instrucciones:

\vspace{0.6em}
\noindent
\begin{tcolorbox}[breakable,enhanced,title=\bf Instrucción en \verb"Mathematica", top=2pt,bottom=2pt,boxsep=0pt,before skip=2pt,after skip=0pt]
\begin{Verbatim}[formatcom=\letraverbatim]
{x,y}={Sqrt[5]*u*Cos[t],Sqrt[5/2]*u*Sin[t]};
r=(1-w)*{x,y,5-y^2}+w*{x,y,x^2+y^2};
A=Dot[D[r,t],Cross[D[r,u],D[r,w]]];
Integrate[A,{t,0,2*Pi},{u,0,1},{w,0,1}]
\end{Verbatim}
\end{tcolorbox}
\vspace{0.7em}
\noindent


\subsection*{Dos paraboloides} \index{Paraboloide}
El siguiente ejercicio es tomado del libro \textcite{thomas_calculo_2015}.
Los paraboloides con ecuaciones $z=8-x^2 -y^2 $ y $z=x^2 +3y^2$ acotan un sólido en el espacio.


\noindent
\begin{minipage}{\textwidth}
	\centering
	\captionof{figure}{Sólido determinado por $x^2 + 3y^2 \leq z \leq 8 - x^2 - y^2$.}
	\label{fig:solido_paraboloides_elipticos}
	\begin{minipage}{0.23\linewidth}
		\centering
		\includegraphics[width=\linewidth]{Img/96a.pdf}
	\end{minipage}
	\hspace{1em}
	\begin{minipage}{0.23\linewidth}
		\centering
		\includegraphics[width=\linewidth]{Img/96b.pdf}
	\end{minipage}
	\hspace{1em}
	\begin{minipage}{0.22\linewidth}
		\centering
		\includegraphics[width=\linewidth]{Img/96c.pdf}
	\end{minipage}
	
	\vspace{0.3em}
	\fuentepropia
\end{minipage}
\vskip0.2cm

Estos gráficos se consiguen con la ayuda de las siguientes instrucciones:

\vspace{0.6em}
\noindent
\begin{tcolorbox}[breakable,enhanced,title=\bf Instrucción en \verb"Mathematica", top=2pt,bottom=2pt,boxsep=0pt,before skip=2pt,after skip=0pt]
\begin{Verbatim}[formatcom=\letraverbatim]
ContourPlot3D[{z==8-x^2-y^2,z==x^2+3*y^2},{x,-3,3},
{y,-3,3},{z,0,8},ContourStyle->{Directive[Cyan,
Opacity[0.8]],Directive[Orange,Opacity[0.9]]},
Mesh->False]
RegionPlot3D[x^2+3*y^2<z<8-x^2-y^2&&x^2+2y^2<4,
{x,-2,2},{y,-Sqrt[2],Sqrt[2]},{z,0,8},Mesh->False,
PlotPoints->90]
\end{Verbatim}
\end{tcolorbox}
\vspace{0.7em}
\noindent


La proyección en el plano $xz$ se obtiene eliminando $y$ de las ecuaciones de los paraboloides, esto es, al sumar las expresiones
$3z=24-3x^2 -3y^2$ y $z=x^2 +3y^2$ se obtiene $2z=12-x^2 $. Sumando las ecuaciones de los paraboloides, además\ de eliminar $x,$ se obtiene la ecuación que termina la proyección del sólido en el plano $yz,$ esto es $2z=8+2y^2 $ que equivale a $z=4+y^2 .$
\vskip0.2cm
Las siguientes figuras muestran la proyección del sólido en los planos $yz$ y $xz$; ellos ilustran el planteamiento de las integrales triples que determinan el volumen del sólido.

\vskip0.3cm
\noindent
\begin{minipage}{\textwidth}
	\centering
	\captionof{figure}{Cono truncado oblicuo determinado por ${\bf r}(u,t)$}
	\begin{minipage}{0.35\linewidth}
		\centering
		\vspace{-10mm}
		\includegraphics[width=\linewidth]{Fig/347.pdf}
	\end{minipage}
	\hspace{1em}
	\begin{minipage}{0.35\linewidth}
		\centering
		\vspace{-10mm}
		\includegraphics[width=\linewidth]{Fig/348.pdf}
	\end{minipage}

	\vspace{-9mm}
	\fuentepropia
\end{minipage}

\vskip0.3cm

Integrando primero con respecto a $y$, las integrales triples correspon\-dientes son:
\[\bi_{-2}^2 \bi_{6-\frac{1}{2}x^2 }^{8-x^2}\bi_{-\sqrt{8-z-x^2}}^{\sqrt{8-z-x^2}}dydzdx+\bi_{-2}^2 \bi_{x^2 }^{6-\frac{1}{2}x^2 }\bi_{-\sqrt{\frac{z-x^2 }{3}}}^{\sqrt{\frac{z-x^2 }{3}} }dydzdx= 8\sqrt{2}\pi\]

De la ecuación $2z=12-x^2 $, se sigue que $x=\pm \sqrt{12-2z}$
entonces
{\footnotesize
\begin{multline}
\bi_{4}^{6}\bi_{-\sqrt{8-z}}^{-\sqrt{12-2z}}\bi_{-\sqrt{8-z-x^2}}^{\sqrt{8-z-x^2}}dydxdz+\bi_{4}^{6}\int_{\sqrt{12-2z}}^{\sqrt{8-z}}\bi_{-\sqrt{8-z-x^2}}^{\sqrt{8-z-x^2}}dydxdz \\
+\bi_{6}^{8}\bi_{-\sqrt{8-z}}^{\sqrt{8-z}}\bi_{-\sqrt{8-z-x^2}}^{\sqrt{8-z-x^2}}dydxdz+\bi_{4}^{6}\int_{-\sqrt{12-2z}}^{\sqrt{12-2z}}\bi_{-\sqrt{\frac{z-x^2 }{3}} }^{\sqrt{\frac{z-x^2 }{3}}}dydxdz \\
+\bi_{0}^{4}\bi_{-\sqrt{z}}^{\sqrt{z}}\bi_{-\sqrt{\frac{z-x^2 }{3}}}^{\sqrt{\frac{z-x^2 }{3}} }dydxdz = 8\sqrt{2}\pi
\end{multline}
}
Estas integrales se evalúan en \verb"Mathematica" con las siguientes instrucciones:

\vspace{0.6em}
\noindent
\begin{tcolorbox}[breakable,enhanced,title=\bf Instrucción en \verb"Mathematica", top=2pt,bottom=2pt,boxsep=0pt,before skip=2pt,after skip=0pt]
\begin{Verbatim}[formatcom=\letraverbatim]
Integrate[1,{x,-2,2},{z,6-x^2/2,8-x^2},
{y,-Sqrt[8-z-x^2],Sqrt[8-z-x^2]}]+Integrate[1,
{x,-2,2},{z,x^2,6-x^2/2},{y,-Sqrt[(z-x^2)/3],
Sqrt[(z-x^2)/3]}]
Integrate[1,{z,4,6},{x,-Sqrt[8-z],-Sqrt[12-2z]},
{y,-Sqrt[8-z-x^2],Sqrt[8-z-x^2]}]+Integrate[1,{z,4,6},
{x,Sqrt[12-2z],Sqrt[8-z]},{y,-Sqrt[8-z-x^2],
Sqrt[8-z-x^2]}]+Integrate[1,{z,6,8},{x,-Sqrt[8-z],
Sqrt[8-z]},{y,-Sqrt[8-z-x^2],Sqrt[8-z-x^2]}]+
Integrate[1,{z,4,6},{x,-Sqrt[12-2z],Sqrt[12-2z]},
{y,-Sqrt[(z-x^2)/3],Sqrt[(z-x^2)/3]}]+Integrate[1,
{z,0,4},{x,-Sqrt[z],Sqrt[z]},{y,-Sqrt[(z-x^2)/3],
Sqrt[(z-x^2)/3]}]
\end{Verbatim}
\end{tcolorbox}
\vspace{0.7em}
\noindent


Integrando primero con respecto a $x$, las integrales triples que determinan el volumen del sólido son las siguientes:
{\footnotesize
\begin{multline} \bi_{-\sqrt{2}}^{\sqrt{2}}\bi_{4+y^2 }^{8-y^2}\bi_{-\sqrt{8-z-y^2}}^{\sqrt{8-z-y^2}}dxdzdy
+\bi_{-\sqrt{2}}^{\sqrt{2}}\bi_{3y^2 }^{4+y^2}\bi_{-\sqrt{z-3y^2 } }^{\sqrt{z-3y^2}}dxdzdy\\
= \bi_{6}^{8}\bi_{-\sqrt{8-z}}^{\sqrt{8-z}}\bi_{-\sqrt{8-z-y^2 }}^{\sqrt{8-z-y^2}}dxdydz+\bi_{4}^{6}\bi_{-\sqrt{z-4}}^{\sqrt{z-4}}\bi_{-\sqrt{8-z-y^2}}^{\sqrt{8-z-y^2 }}dxdydz \\
+\bi_{4}^{6}\int_{-\sqrt{z/3}}^{-\sqrt{z-4}}\bi_{-\sqrt{z-3y^2}}^{\sqrt{z-3y^2}}dxdydz+\bi_{4}^{6}\bi_{\sqrt{z-4}}^{\sqrt{z/3}}\bi_{-\sqrt{z-3y^2}}^{\sqrt{z-3y^2 }}dxdydz \\
+\bi_{0}^{4}\bi_{-\sqrt{z/3}}^{\sqrt{z/3}}\bi_{-\sqrt{z-3y^2 }}^{\sqrt{z-3y^2 } }dxdydz = 8\sqrt{2}\pi
\end{multline}
}
Con las siguientes instrucciones se pueden evaluar estas integrales.
\vskip0.2cm
\vspace{0.6em}
\noindent
\begin{tcolorbox}[breakable,enhanced,title=\bf Instrucción en \verb"Mathematica", top=2pt,bottom=2pt,boxsep=0pt,before skip=2pt,after skip=0pt]
\begin{Verbatim}[formatcom=\letraverbatim]
Integrate[1,{y,-Sqrt[2],Sqrt[2]},{z,4+y^2,8-y^2},
{x,-Sqrt[8-z-y^2],Sqrt[8-z-y^2]}]+Integrate[1,{y,
Sqrt[2],Sqrt[2]},{z,3*y^2,4+y^2},{x,-Sqrt[z-3y^2],
Sqrt[z-3y^2]}]
Integrate[1,{z,6,8},{y,-Sqrt[8-z],Sqrt[8-z]},{x,
Sqrt[8-z-y^2],Sqrt[8-z-y^2]}]+Integrate[1,{z,4,6},
{y,-Sqrt[z-4],Sqrt[z-4]},{x,-Sqrt[8-z-y^2],
Sqrt[8-z-y^2]}]+Integrate[1,{z,4,6},{y,-Sqrt[z/3],
Sqrt[z-4]},{x,-Sqrt[z-3y^2],Sqrt[z-3y^2]}]+
Integrate[1,{z,4,6},{y,Sqrt[z-4],Sqrt[z/3]},{x,
Sqrt[z-3y^2],Sqrt[z-3y^2]}]+Integrate[1,{z,0,4},{y,
Sqrt[z/3],Sqrt[z/3]},{x,-Sqrt[z-3y^2],Sqrt[z-3y^2]}]
\end{Verbatim}
\end{tcolorbox}
\vspace{0.7em}
\noindent



La proyección en el plano $xy$ está dada por
$ 8-x^2 -y^2 =x^2 +3y^2,$ que se simplifica en la forma $ 4=x^2 +2y^2;$
esta ecuación corresponde a una elipse cuya figura se muestra enseguida.

\vskip0.2cm
\noindent
\begin{minipage}{\textwidth}
	\centering
	\captionof{figure}{$x^2 + 2y^2 \leq 4$}
	\label{fig:curva_parametrica_uv_5}
	\begin{minipage}{0.36\linewidth}
		\centering
		\vspace{-18mm}
		\includegraphics[width=\linewidth]{Fig/349.pdf}
	\end{minipage}
	
	\vspace{-10mm}
	\fuentepropia
\end{minipage}
\vskip0.3cm


Con ayuda de esta figura se puede plantear las integrales que permiten hallar el volumen del sólido.
En los órdenes $dzdydx$ y $dzdxdy$, las integrales son las siguientes.

\vskip0.4cm
En los órdenes $dzdydx$ y $dzdxdy$, las integrales son las siguientes.

\begin{multline*}
\int_{-2}^2 \int_{-\frac{1}{2}\sqrt{8-2x^2 }}^{\frac{1}{2}\sqrt{8-2x^2 }}\int_{x^2 +3y^2 }^{8-x^2 -y^2 }dzdydx= \\
\int_{-\sqrt{2}}^{\sqrt{2}}\int_{-\sqrt{4-2y^2 }}^{\sqrt{4-2y^2 }}\int_{x^2 +3y^2 }^{8-x^2 -y^2 }dzdxdy= 8\sqrt{2}\pi . \end{multline*}
\vskip0.2cm
Estas integrales se evalúan con la ejecución de las siguientes instrucciones:

\vspace{0.6em}
\noindent
\begin{tcolorbox}[breakable,enhanced,title=\bf Instrucción en \verb"Mathematica", top=2pt,bottom=2pt,boxsep=0pt,before skip=2pt,after skip=0pt]
\begin{Verbatim}[formatcom=\letraverbatim]
Integrate[1,{x,-2,2},{y,-Sqrt[8-2x^2]/2,
Sqrt[8-2x^2]/2},{z,x^2+3*y^2,8-x^2-y^2}]
Integrate[1,{y,-Sqrt[2],Sqrt[2]},{x,-Sqrt[4-2*y^2],
Sqrt[4-2*y^2]},{z,x^2+3*y^2,8-x^2-y^2}]
\end{Verbatim}
\end{tcolorbox}
\vspace{0.7em}
\noindent

Otra alternativa para obtener el volumen del sólido es la siguiente.
\vskip0.3cm

% \begin{figure}[H]
% \centering
% \begin{pspicture}(-2.2,-1.2)(2.5,1.4) %\psgrid
% \psset{algebraic,unit=0.7}
% \pscustom[fillstyle=vlines,hatchcolor=gray,hatchsep=1.5pt,linestyle=none]{
%   \psparametricplot{0}{\psPiTwo}{2*cos(t)|sqrt(2)*sin(t)}
% }
% \psaxes[Dx=2,deltax=2,linecolor=black]{->}(0,0)(-2.5,-1.5)(3,1.8)[$x$,90][$y$,0]
% \psparametricplot[linecolor=black!60!red]{0}{\psPiTwo}{2*cos(t)|sqrt(2)*sin(t)}
% % \rput(2.5,1.3){$x^2+2y^2=4$}
% \end{pspicture}
% \caption{$x^2+2y^2 \leq 4$.}
% \end{figure}
% \par
% \justifying

%\begin{tcolorbox}[title=\bfseries Parametrizando el sólido]

La intersección de los paraboloides está dada por \( x^2+2y^2=4 \), que representa una elipse en el plano \( xy \). Al ser $x = 2u \sin(t)$, $y = \sqrt{2}u \cos(t)$, con $u \in [0,1]$, $t \in [0,2\pi],$ se describe esta región.


\vskip0.2cm
\noindent
\begin{minipage}{\textwidth}
	\centering
	\captionof{figure}{$x^2 + 2y^2 \leq 4$}
	\label{fig:curva_parametrica_uv_6}
	\begin{minipage}{0.38\linewidth}
		\centering
		\vspace{-18mm}
		\includegraphics[width=\linewidth]{Fig/349.pdf}
	\end{minipage}
	
	\vspace{-11mm}
	\fuentepropia
\end{minipage}
\vskip0.3cm


Con dos puntos:
\begin{align*}
P &= \left( 2u\sin(t),\ \sqrt{2}u\cos(t),\ 8 - u^2(3 - \cos(2t)) \right), \\
Q &= \left( 2u\sin(t),\ \sqrt{2}u\cos(t),\ u^2(5 + \cos(2t)) \right).
\end{align*}

Cada uno sobre un paraboloide. Una combinación convexa de la forma \( (1-w)P + wQ \) define la parametrización del sólido:
\[
\mathbf{r}(u,t,w) = \left\langle 2u\sin t,\ \sqrt{2}u\cos t,\ 8(1-w) + u^2(8w - 3 + \cos(2t)) \right\rangle.
\]

De esta función obtenemos los siguientes vectores:
\begin{align*}
\mathbf{r}_t &= \left\langle 2u\cos(t),\ -\sqrt{2}u\sin(t),\ -2u^2\sin(2t) \right\rangle \\
\mathbf{r}_u &= \left\langle 2\sin(t),\ \sqrt{2}\cos(t),\ 2u(8w - 3 + \cos(2t)) \right\rangle \\
\mathbf{r}_w &= \left\langle 0,\ 0,\ 8(u^2 - 1) \right\rangle.
\end{align*}

Por tanto,
\[
\left| \mathbf{r}_w \cdot (\mathbf{r}_u \times \mathbf{r}_t) \right| = 16\sqrt{2}u(1 - u^2).
\]

Entonces el volumen del sólido está dado por:
\vskip0.2cm
\[
\int_0^1 \int_0^1 \int_0^{2\pi} 16\sqrt{2}u(1 - u^2)\, dt\, du\, dw = 8\sqrt{2}\pi.
\]
%\end{tcolorbox}
\vskip0.2cm
El cálculo anterior resulta sencillo de evaluar en \verb"Mathematica" con ayuda de las siguientes instrucciones:


\vspace{0.6em}
\noindent
\begin{tcolorbox}[breakable,enhanced,title=\bf Instrucción en \verb"Mathematica", top=2pt,bottom=2pt,boxsep=0pt,before skip=2pt,after skip=0pt]
\begin{Verbatim}[formatcom=\letraverbatim]
{x,y}={2*u*Sin[t],Sqrt[2]*u*Cos[t]};
r=Simplify[(1-w)*{x,y,8-x^2-y^2}+w*{x,y,x^2+3*y^2}];
Integrate[Dot[D[r,w],Cross[D[r,u],D[r,t]]],{w,0,1},
{u,0,1},{t,0,2*Pi}]
\end{Verbatim}
\end{tcolorbox}
\vspace{0.7em}
\noindent


\subsection*{Dos paraboloides II}
Las ecuaciones $63x^2+252x-8+28z^2+56z+36y=0$y $9x^2+36x+16+4z^2+8z-12y=0$
representan dos paraboloides. Estas superficies acotan la región en el espacio que se muestra a continuación.

\begin{figure}[H]
\centering
\caption{Región acotada por los paraboloides.}
\label{fig:region_paraboloides_89}
\includegraphics[height=4.5cm]{Img/89b.pdf}
\includegraphics[height=4.5cm]{Img/89c.pdf}
\includegraphics[height=4.5cm]{Img/89d.pdf}
\fuentepropia
\end{figure}

Con las siguientes instrucciones se obtuvieron estos gráficos.

\vspace{0.6em}
\noindent
\begin{tcolorbox}[breakable,enhanced,title=\bf Instrucción en \verb"Mathematica", top=2pt,bottom=2pt,boxsep=0pt,before skip=2pt,after skip=0pt]
\begin{Verbatim}[formatcom=\letraverbatim]
ContourPlot3D[{(x+2)^2/4+(z+1)^2/9==(8/7-y/7),
36*(y/3+2/3)==9*(x+2)^2+4*(z+1)^2},{x,-6,2},
{y,-2,8},{z,-7,5},Mesh->False,PlotPoints->60,
ContourStyle->{Directive[Orange,Opacity[0.95]],
Directive[Cyan,Opacity[0.8]]}]
RegionPlot3D[3*(x+2)^2/4+(z+1)^2/3-2<=y<8-7*(x+2)^2/4-
7*(z+1)^2/9&&-Sqrt[9-9*(x+2)^2/4]-1<z<Sqrt[9-
9*(x+2)^2/4]-1&&-4<x<0,{x,-4,0},{y,-2,8},{z,-4,2},
PlotPoints->60,MeshStyle->RGBColor[0.4,0.5,0.4],
PlotStyle->Orange,BoxRatios->{2,2,1.5}]
\end{Verbatim}
\end{tcolorbox}
\vspace{0.7em}
\noindent


Las ecuaciones inicialmente consideradas, al factorizarlas, se reescriben como
$63(x+2)^2+28(z+1)^2=36(8-y)$ y $9\left( x+2\right) ^2+4\left( z+1\right) ^2=12(y+2).$
Eliminado $y$ se obtiene $9x^2+36x+4z^2+8z=-4,$ que corresponde a la proyección del sólido en el plano $xz$, esta ecuación equivale a
$\frac{\left( x+2\right) ^2}{4}+\frac{\left( z+1\right) ^2}{9}=1$
y representa una elipse con centro en el punto $(-2,-1)$,
alargada hacia el eje $z.$ De esta ecuación se sigue que
$x=-2\pm \sqrt{4-\frac{4\left( z+1\right) ^2}{9}} $ y $z=-1\pm
\sqrt{9-\frac{9\left( x+2\right) ^2}{4}}.$ 
\vskip 0.3cm
Si $z=-1,$ las ecuaciones se reducen a
$\frac{\left( x+2\right) ^2}{4}=\frac{8-y}{7}$ y $9\left( x+2\right) ^2=36\left( \frac{y+2}{3}\right) $
y representan dos parábolas en el plano $xy.$
Si $x=-2,$, estas ecuaciones se reducen a $\frac{\left( z+1\right) ^2}{9}=\frac{8-y}{7}$ y
$4\left( z+1\right) ^2=36\left( \frac{y+2}{3}\right) $
y representan dos parábolas en el plano $yz$.
\vskip 0.3cm
Por lo descrito anteriormente, el sólido en mención se proyecta sobre los planos coordenados como se muestra a continuación.
\vskip0.4cm

\noindent
\begin{minipage}{\textwidth}
	\centering
	\captionof{figure}{Región acotada por los paraboloides proyectada en los planos coordenados.}
	\begin{minipage}{0.30\linewidth}
		\centering
		\vspace{-10mm}	\includegraphics[width=\linewidth]{Fig/350.pdf}
	\end{minipage}
	\hspace{1em}
	\begin{minipage}{0.30\linewidth}
		\centering
		\vspace{-10mm}
		\includegraphics[width=\linewidth]{Fig/351.pdf}
	\end{minipage}
	\hspace{1em}
	\begin{minipage}{0.30\linewidth}
		\centering
		\vspace{-10mm}
		\includegraphics[width=\linewidth]{Fig/352.pdf}
	\end{minipage}
	\fuentepropia
\end{minipage}
\vskip0.3cm

Utilizando esta representación se plantean las integrales triples que determinan el volumen del sólido:
\vskip0.3cm
\begin{equation*}
\bigintsss_{-4}^{0}\bigintsss_{-1-\sqrt{9-\frac{9\left( x+2\right) ^2}{4}}}^{-1+\sqrt{9-\frac{9\left( x+2\right) ^2}{4}}}
\bigintsss_{\frac{3}{4}x^2+3x+\frac{4}{3}+\frac{1}{3}z^2+\frac{2}{3}z}^{\frac{2}{9}-\frac{7}{4}x^2-7x-\frac{7}{9}z^2-\frac{14}{9}z}
dydzdx = 30\pi
\end{equation*}

\begin{equation*}
\bigintsss_{-4}^2\bigintsss_{-2-\sqrt{4-\frac{4\left( z+1\right) ^2}{9}}}^{-2+\sqrt{4-\frac{4\left( z+1\right) ^2}{9}}}
\bigintsss_{\frac{3}{4}x^2+3x+\frac{4}{3}+\frac{1}{3}z^2+\frac{2}{3}z}^{\frac{2}{9}-\frac{7}{4}x^2-7x-\frac{7}{9}z^2-\frac{14}{9}z}
dydxdz = 30\pi.
\end{equation*}
\vskip0.4cm
Para integrar primero con respecto a $z$, se debe obtener $z$ de las expresiones originales y tener en cuenta que $z$ recorre los puntos desde la parte negativa de cada paraboloide hasta su parte positiva.

\begin{multline*}
\bigintsss_{-4}^{0}\bigintsss_{1}^{8-\frac{7}{4}\left( x+2\right) ^2}\bigintsss_{-1-3\sqrt{\frac{8-y}{7}-\frac{\left( x+2\right) ^2}{4}}}^{-1+3\sqrt{\frac{8-y}{7}-\frac{\left( x+2\right) ^2}{4}}}dzdydx \\
+\bigintsss_{-4}^{0}\bigintsss_{\frac{3}{4}\left( x+2\right) ^2-2}^{1}\bigintsss_{-1-3\sqrt{\frac{y+2}{3}-\frac{1}{4}\left(x+2\right) ^2}}^{-1+3\sqrt{\frac{y+2}{3}-\frac{1}{4}\left( x+2\right) ^2}}dzdydx= 30\pi.
\end{multline*}

Al integrar primero con respecto a $x,$ se debe presente que $x$ igual que en el caso anterior, recorre los puntos desde la parte negativa de cada paraboloide hasta la parte positiva del mismo, por tal razón las integrales en este caso son
\vspace{-2mm}
\begin{multline*}
\bigintsss_{-4}^2\bigintsss_{1}^{8-\frac{7\left( z+1\right) ^2}{9}}\bigintsss_{-2-2\sqrt{\frac{8-y}{7}-\frac{\left( z+1\right) ^2}{9}}}^{-2+2\sqrt{\frac{8-y}{7}-\frac{\left( z+1\right) ^2}{9}}}dxdydz \\
+\bigintsss_{-4}^2\bigintsss_{\frac{\left(z+1\right) ^2}{3}-2}^{1}\bigintsss_{-2-2\sqrt{\frac{1}{3}\left( y+2\right) -\frac{1}{9}\left( z+1\right) ^2}}^{-2+2\sqrt{\frac{1}{3}\left( y+2\right) -\frac{1}{9}\left( z+1\right) ^2}}dxdydz= 30\pi
\end{multline*}
Con las siguientes instrucciones se evaluó la integral anterior.

\vspace{0.6em}
\noindent
\begin{tcolorbox}[breakable,enhanced,title=\bf Instrucción en \verb"Mathematica", top=2pt,bottom=2pt,boxsep=0pt,before skip=2pt,after skip=0pt]
\begin{Verbatim}[formatcom=\letraverbatim]
Integrate[1,{z,-4,2},{y,1,8-7*(z+1)^2/9},{x,-2-
2*Sqrt[(8-y)/7-(z+1)^2/9],-2+2*Sqrt[(8-y)/7-
(z+1)^2/9]}]+Integrate[1,{z,-4,2},{y,-2+(z+1)^2/3,1},
{x,-2-Sqrt[(4/3)*(y+2)-4(z+1)^2/9],-2+Sqrt[(4/3)*(y+2)
4(z+1)^2/9]}]
\end{Verbatim}
\end{tcolorbox}
\vspace{0.7em}
\noindent


También puede usarse la siguiente opción para calcular el volumen.

\vskip0.3cm

%\begin{tcolorbox}
En el plano $xz$ la proyección del sólido es una elipse, cuyo interior se parametriza como $x=2u\cos(t)-2$, $z=3u\sin(t)-1.$ Sobre los parabo\-loides, se obtiene $y=3u^2-2$ y $y=8-7u^2;$
con dos puntos sobre estas superficies $P(-2+2u\cos (t) ,-1+3u\sin (t),3u^2-2)$ y
$Q(2u\cos (t)-2,3u\sin (t)-1,8-7u^2)$, se forma una combinación convexa de la forma $(1-w)P+wQ,$ que define la función que parametriza el sólido.
\vspace{-3mm}
\begin{equation*}
{\bf r}\left( t,u,w\right) =\langle 2u\cos t-2,3u\sin t-1,10w\left( 1-u^2\right) +3u^2-2 \rangle,
\end{equation*}
con $t\in [0,2\pi ],$ $u\in [0,1],$ $w\in [0,1];$
y obtenemos los siguientes vectores
\begin{eqnarray*}
{\bf r}_{t}&=& \langle -2u\sin t,3u\cos t,0\rangle \\
{\bf r}_{u}&=& \langle 2\cos t,3\sin t,6u-20uw\rangle\\
{\bf r}_{w}&=& \langle 0,0,10-10u^2\rangle.
\end{eqnarray*}
El elemento diferencial de volumen es
$\left\vert {\bf r}_{t}\cdot \left({\bf r}_{u}\times {\bf r}_{w}\right) \right\vert
=\left\vert 60u\left( u^2-1\right) \right\vert ,$
cuya integral proporciona el volumen del sólido mencionado
\vskip0.2cm
\[\int_{0}^{1}\int_{0}^{1}\int_{0}^{2\pi }\left\vert 60u\left( u^2-1\right)\right\vert dtdudw= 30\pi.\]
%\end{tcolorbox}

\section{Un hiperboloide y una esfera} \index{Esfera} \index{Hiperboloide}

Las ecuaciones $x^2+y^2+z^2=16$ y $x^2+z^2=9+y^2,$ representan una esfera de radio $4$
y un hiperboloide de un manto con eje de simetría el eje $y.$
La región interior a ambas superficies determinó el sólido que se muestra en seguida.

\vskip0.2cm
\noindent
\begin{minipage}{\textwidth}
	\centering
	\captionof{figure}{Región determinada por $x^2 + y^2 + z^2 \leq 16$ y $x^2 + z^2 \leq 9 + y^2$.}
	\label{fig:region_esferica_cilindrica}
	\begin{minipage}{0.20\linewidth}
		\centering
		\includegraphics[width=\linewidth]{Img/61a.pdf}
	\end{minipage}
	\hspace{1em} 
	\begin{minipage}{0.20\linewidth}
		\centering
		\includegraphics[width=\linewidth]{Img/61c.pdf}
	\end{minipage}
	\hspace{1em} 
	\begin{minipage}{0.20\linewidth}
		\centering
		\includegraphics[width=\linewidth]{Img/61b.pdf}
	\end{minipage}
	
	\vspace{0.3em}
	\fuentepropia
\end{minipage}

\vskip0.2cm
La intersección de la esfera y el hiperboloide se da en $y=\pm \frac{\sqrt{14}}{2}.$
En seguida se muestra la proyección del sólido en los planos $xz$ y $xy$.
Para la región que contiene al origen en el plano $xz$, $y$ recorre, por un lado, los puntos del sólido que van desde la parte negativa de la esfera hasta su parte positiva, y por otro lado, para la región entre las dos circunferencias, $y$
recorre los puntos entre el hiperboloide y la esfera; esto ocurre en dos regiones simétricas con respecto al plano $xz.$

\vskip0.2cm
\noindent
\begin{minipage}{\textwidth}
	\centering
	\captionof{figure}{Proyección en los planos coordenados.}
	\begin{minipage}{0.30\linewidth}
		\centering
		\vspace{-17mm}
		\includegraphics[width=\linewidth]{Fig/353.pdf}
	\end{minipage}
	\hspace{1em}
	\begin{minipage}{0.38\linewidth}
		\centering
		\vspace{-17mm}
		\includegraphics[width=\linewidth]{Fig/354.pdf}
	\end{minipage}
	
	\vspace{-14mm}
	\fuentepropia
\end{minipage}

% \begin{figure}[H]
% \centering
% % Primer gráfico
% \begin{pspicture}(-1.75,-2)(3,2) %\psgrid
% \psset{algebraic,unit=0.45}
% \pscustom[fillstyle=hlines,hatchcolor=gray!60,hatchsep=1.2pt,hatchwidth=0.5pt,opacity=0.7,linestyle=none]{
%   \psparametricplot{0}{\psPiTwo}{3*cos(t)|3*sin(t)}
% }
% \pscustom[fillstyle=solid,fillcolor=gray!60,opacity=0.7,linestyle=none]{
%   \psparametricplot{0}{\psPiTwo}{3*cos(t)|3*sin(t)}
%   \psline(3,0)(4,0)
%   \psparametricplot{\psPiTwo}{0}{4*cos(t)|4*sin(t)}
% }
% \psaxes[Dy=3,Dx=3,linecolor=black,labelFontSize=\scriptstyle]{->}(0,0)(-5,-5)(5.5,5.6)[$x$,90][$z$,0]
% \psparametricplot[linecolor=black!60!blue]{0}{\psPiTwo}{3*cos(t)|3*sin(t)}
% \psparametricplot[linecolor=black!60!blue]{0}{\psPiTwo}{4*cos(t)|4*sin(t)}
% \psline[linecolor=black!60!blue,linewidth=0.6pt,arrows=<->](-3,-3.5)(-3,3.75)
% \psline[linecolor=black!60!blue,linewidth=0.6pt,arrows=<->](3,-3.5)(3,3.75)
% \rput(0.3,0.8){$x^2+z^2=9$}
% \rput(4,4.5){$x^2+z^2=\frac{25}{2}$}
% \end{pspicture}

% \vspace{0.5cm}

% % Segundo gráfico
% \begin{pspicture}(-2.8,-2)(3,2) % \psgrid
% \psset{algebraic,unit=0.425}
% \psaxes[Dy=4,Dx=3,linecolor=black,labelFontSize=\scriptstyle]{->}(0,0)(-5,-5)(8,5.6)[$x$,90][$y$,0]
% \pscustom[fillstyle=hlines,hatchcolor=gray!60,hatchsep=1.2pt,hatchwidth=0.5pt,opacity=0.7,linestyle=none]{
%   \psparametricplot{-0.59}{0.59}{3*COSH(t)|3*SINH(t)}
%   \psparametricplot{0.49}{2.65}{4*cos(t)|4*sin(t)}
%   \psparametricplot{0.59}{-0.59}{-3*COSH(t)|3*SINH(t)}
%   \psparametricplot{3.63}{5.79}{4*cos(t)|4*sin(t)}
% }
% \psparametricplot[linecolor=black!60!blue]{0}{\psPiTwo}{4*cos(t)|4*sin(t)}
% \psparametricplot[linecolor=black!60!blue,arrows=->]{-1.1}{1.2}{3*COSH(t)|3*SINH(t)}
% \psparametricplot[linecolor=black!60!blue,arrows=->]{-1.1}{1.2}{-3*COSH(t)|3*SINH(t)}
% \psline[linecolor=black!60!blue,linewidth=0.6pt,arrows=->](-5,-1.87)(5.5,-1.87)
% \psline[linecolor=black!60!blue,linewidth=0.6pt,arrows=->](3.53,3.5)(3.53,-4)
% \rput(8.5,-2.1){$y=-\sqrt{\frac{7}{2}}$}
% \rput(3.75,-5){$x=\frac{5}{\sqrt{2}}$}
% \rput(7,5){$x^2-y^2=9$}
% % \rput(-0.4,2){$x^2+y^2=16$}
% \end{pspicture}
% \caption{Proyección en los planos coordenados.}
% \end{figure}
% \par
% \justifying

\vskip0.2cm
Sobre el plano $xz,$ se puede utilizar la simetría de la región, integrando la región para $z \geq 0,$ y se multiplica por $4$ a estas integrales. En este caso, el volumen del sólido está dado por:

\vspace{-4mm}
{\small
\begin{multline*}
\hskip-0.4cm \int_{-3}^{3}\int_{-\sqrt{9-x^2}}^{\sqrt{9-x^2}}\int_{-\sqrt{16-x^2-z^2}}^{\sqrt{16-x^2-z^2}}dydzdx + 4\bigg(\int_{-3}^{3}\int_{\sqrt{9-x^2}}^{\sqrt{\frac{25}{2}-x^2}}\int_{\sqrt{x^2+z^2-9}}^{\sqrt{16-x^2-z^2}}dydzdx+ \\
2\int_{3}^{\frac{5}{\sqrt{2}}}\int_{0}^{\sqrt{\frac{25}{2}-x^2}}\int_{\sqrt{x^2+z^2-9}}^{\sqrt{16-x^2-z^2}}dydzdx\bigg)= \frac{2\pi}{3}(128-7 \sqrt{14}).
\end{multline*}
}

%\vskip0.2cm
Si se integra primero con respecto a $z$, hay tres regiones en el espacio, en dos de ellas $z$ recorre los puntos desde la parte negativa de la esfera hasta su parte positiva; en la tercer región $z$ recorre los puntos que van desde la parte negativa del hiperboloide hasta la parte positiva del mismo, razón por la cual las integrales a evaluar en este caso son:
\vspace{-2mm}
\begin{multline*}
2\bi_{\sqrt{\frac{7}{2}}}^{4}\bi_{-\sqrt{16-y^2}}^{\sqrt{16-y^2}}\bi_{-\sqrt{16-x^2-y^2}}^{\sqrt{16-x^2-y^2}}dzdxdy+ \\
\bi_{-\sqrt{\frac{7}{2}}}^{\sqrt{\frac{7}{2}}}\bi_{-\sqrt{9+y^2}}^{\sqrt{9+y^2}}\bi_{-\sqrt{9+y^2-x^2}}^{\sqrt{9+y^2-x^2}}dzdxdy =\frac{2}{3}\pi\left( 128-7\sqrt{14}\right).
\end{multline*}
En coordenadas cilíndricas, el cálculo es el siguiente:
\begin{multline*}
\int_{0}^{2\pi }\int_{0}^{3}\int_{-\sqrt{16-r^2}}^{\sqrt{16-r^2}}rdydrd\theta \\ +2\int_{0}^{2\pi}\int_{3}^{\frac{5}{\sqrt{2}}}\int_{\sqrt{r^2-9}}^{\sqrt{16-r^2}}rdydrd\theta = \frac{2\pi}{3}(128-7 \sqrt{14}),
\end{multline*}
y su obtención en \verb"Mahematica" se consigue con la siguiente sentencia:

\vspace{0.6em}
\noindent
\begin{tcolorbox}[breakable,enhanced,title=\bf Instrucción en \verb"Mathematica", top=2pt,bottom=2pt,boxsep=0pt,before skip=2pt,after skip=0pt]
\begin{Verbatim}[formatcom=\letraverbatim]
Integrate[r,{\[Theta],0,2*Pi},{r,0,3},{y,-Sqrt[16-
r^2],Sqrt[16-r^2]}]+2*Integrate[r,{\[Theta],0,2*Pi},
{r,3,5/Sqrt[2]},{y,Sqrt[r^2-9],Sqrt[16-r^2]}]
\end{Verbatim}
\end{tcolorbox}
\vspace{0.7em}
\noindent


\begin{comment}
Estos gráficos se obtienen en \verb"Mathematica" con las siguientes instrucciones

\begin{verbatim}
ContourPlot3D[{x^2+z^2==9+y^2,x^2+y^2+z^2==16},{x,-5,5},
{y,-4,4},{z,-5,5},Mesh->False,Ticks->{{-3,0,3},{-3,0,3},
{-4,-2,2,4}},AxesLabel->{x,y,z},
ContourStyle->{Blend[{Yellow,Cyan},1/9],Orange}]
\end{verbatim}

\begin{verbatim}
RegionPlot3D[x^2+z^2<9+y^2&&x^2+y^2+z^2<16,{x,-4,4},{y,-4,4},
{z,-4,4},PlotPoints->90,Mesh->False,Ticks->{{-4,-2,0,2,4},
{-3,-1,1,3},{-3,-1,1,3}},PlotStyle->RGBColor[0.3,0.9,0.5],
AxesLabel->{x,y,z}]
\end{verbatim}


La intersección de la esfera y el hiperboloide se obtiene al sustituir $x^2+z^2=9+y^2$ en la ecuación $x^2+y^2+z^2=16,$
lo que genera la expresión $2y^2=7,$, es decir $y=\pm \frac{\sqrt{14}}{2}.$

% \begin{figure}
% \centering
% % Primer gráfico: proyección en plano xz
% \begin{pspicture}(-2,-1.4)(3,2) % \psgrid
% \psset{algebraic,unit=0.36}
% \pscustom[fillstyle=hlines,hatchcolor=gray!60,hatchsep=1.2pt,hatchwidth=0.5pt,opacity=0.7,linestyle=none]{
%   \psparametricplot{0}{\psPiTwo}{3*cos(t)|3*sin(t)}
% }
% \pscustom[fillstyle=solid,fillcolor=gray!60,opacity=0.7,linestyle=none]{
%   \psparametricplot{0}{\psPiTwo}{3*cos(t)|3*sin(t)}
%   \psline(3,0)(4,0)
%   \psparametricplot{\psPiTwo}{0}{4*cos(t)|4*sin(t)}
% }
% \psaxes[Dy=3,Dx=3,linecolor=black,labelFontSize=\scriptstyle]{->}(0,0)(-5,-5)(5.5,5.6)[$x$,90][$z$,0]
% \psparametricplot[linecolor=black!60!blue]{0}{\psPiTwo}{3*cos(t)|3*sin(t)}
% \psparametricplot[linecolor=black!60!blue]{0}{\psPiTwo}{4*cos(t)|4*sin(t)}
% \psline[linecolor=black!60!blue,linewidth=0.6pt,arrows=<->](-3,-3.5)(-3,3.75)
% \psline[linecolor=black!60!blue,linewidth=0.6pt,arrows=<->](3,-3.5)(3,3.75)
% \rput(0.3,0.8){$x^2+z^2=9$}
% \rput(4,4.5){$x^2+z^2=\frac{25}{2}$}
% \end{pspicture}

% \vspace{0.5cm}

% % Segundo gráfico: proyección en plano xy
% \begin{pspicture}(-2.5,-1.4)(3,2) % \psgrid
% \psset{algebraic,unit=0.35}
% \psaxes[Dy=4,Dx=3,linecolor=black,labelFontSize=\scriptstyle]{->}(0,0)(-5,-5)(8,5.6)[$x$,90][$y$,0]
% \pscustom[fillstyle=hlines,hatchcolor=gray!60,hatchsep=1.2pt,hatchwidth=0.5pt,opacity=0.7,linestyle=none]{
%   \psparametricplot{-0.59}{0.59}{3*COSH(t)|3*SINH(t)}
%   \psparametricplot{0.49}{2.65}{4*cos(t)|4*sin(t)}
%   \psparametricplot{0.59}{-0.59}{-3*COSH(t)|3*SINH(t)}
%   \psparametricplot{3.63}{5.79}{4*cos(t)|4*sin(t)}
% }
% \psparametricplot[linecolor=black!60!blue]{0}{\psPiTwo}{4*cos(t)|4*sin(t)}
% \psparametricplot[linecolor=black!60!blue,arrows=->]{-1.1}{1.2}{3*COSH(t)|3*SINH(t)}
% \psparametricplot[linecolor=black!60!blue,arrows=->]{-1.1}{1.2}{-3*COSH(t)|3*SINH(t)}
% \psline[linecolor=black!60!blue,linewidth=0.6pt,arrows=->](-5,-1.87)(5.5,-1.87)
% \psline[linecolor=black!60!blue,linewidth=0.6pt,arrows=->](3.53,3.5)(3.53,-4)
% \rput(7.75,-1.9){$y=-\sqrt{\frac{7}{2}}$}
% \rput(3.75,-5){$x=\frac{5}{\sqrt{2}}$}
% \rput(4.5,5){$x^2-y^2=9$}
% \rput(-0.4,2){$x^2+y^2=16$}
% \end{pspicture}
% \caption{Proyección en los planos coordenados.}
% \end{figure}



% \begin{figure}
% \centering
% \begin{pspicture}(-2.5,-1.8)(3,2.3) % \psgrid
% \psset{algebraic,unit=0.45}
% \pscustom[fillstyle=hlines,hatchcolor=gray!90,hatchsep=1.2pt,hatchwidth=0.5pt,linestyle=none]{
%   \psparametricplot{0}{\psPiTwo}{3*cos(t)|3*sin(t)}
% }
% \pscustom[fillstyle=solid,fillcolor=gray!60,opacity=0.7,linestyle=none]{
%   \psparametricplot{0}{\psPiTwo}{3*cos(t)|3*sin(t)}
%   \psline(3,0)(4,0)
%   \psparametricplot{\psPiTwo}{0}{5/sqrt(2)*cos(t)|5/sqrt(2)*sin(t)}
% }
% \psaxes[Dy=3,Dx=3,linecolor=black,labelFontSize=\scriptstyle]{->}(0,0)(-4.5,-4.5)(5.5,5.2)[$x$,90][$z$,0]
% \psparametricplot[linecolor=black!60!blue]{0}{\psPiTwo}{3*cos(t)|3*sin(t)}
% \psparametricplot[linecolor=black!60!blue]{0}{\psPiTwo}{5/sqrt(2)*cos(t)|5/sqrt(2)*sin(t)}
% %\rput(4,4.75){$x^2+z^2=\frac{25}{2}$}
% %\rput(5,3.5){$x^2+z^2=9$}
% \psline[linewidth=1.2pt,linestyle=dashed](-1.96,2.94)(0,0)
% % \psline[linewidth=0.6pt,arrows=->](2.2,2)(4,3)
% \rput(-1.96,2.94){$\bullet$}
% \rput(-1.66,2.496){$\bullet$}
% \rput(-1.8,1.65){$P$}
% \rput(-2.2,3.55){$Q$}
% \end{pspicture}
% \caption{$9 \leq x^2+z^2 \leq \frac{25}{2}$.}
% \label{figura168}
% \end{figure}

El eje $y$ es el eje de simetría del hiperboloide, por este motivo usaremos la proyección en el plano $xz,$ para construir la parametrización del sólido. El gráfico muestra la figura de curva de intersección $x^2+z^2=\frac{25}{2}$
y la traza del hiperboloide cuando $y=0.$
La región acotada por $x^2+z^2=9,$ se describe usando las expresiones $x=3u\sin v,$ $z=3u\cos v.$

Sobre el sólido, esta región se proyecta y determina los puntos que estan entre la parte negativa de la esfera hasta la parte positiva de la misma. Es decir, la expresión para $y,$ es
$y=\pm \sqrt{16-x^2-z^2}= \pm \sqrt{16-9u^2}$.
Es por lo anterior que la siguiente combinación convexa
$(1-w) \langle 3u\sin v,-\sqrt{16-9u^2},3u\cos v\rangle + w\langle 3u\sin v,\sqrt{16-9u^2},3u\cos v\rangle$
permite definir una función que parametriza el sólido
$ {\bf r} (u,v,w)= \left\langle 3u\sin v,\left(2w-1\right) \sqrt{16-9u^2},3u\cos v\right\rangle $
{\bf El volumen de la parte del sólido interior a la esfera y al cilindro} $x^2+z^2=9,$
se determina con el siguiente procedimiento. Con la función ${\bf r}$ hallamos los siguientes vectores
\begin{eqnarray*}
{\bf r}_{u}&=& \langle 3\sin v,\frac{9u(1-2w)}{\sqrt{16-9u^2}},3\cos v\rangle \\
{%\bf r}_{v}&=& \langle 3u\cos v,0,-3u\sin v \rangle \\
{\bf r}_{w}&=& \langle 0,2\sqrt{16-9u^2},0\rangle
\end{eqnarray*}
Entonces, $\left\vert {\bf r}_{w}\cdot \left({\bf r}_{u}\times {\bf r}_{v}\right) \right\vert =18u\sqrt{16-9u^2},$
Por tanto, el volumen de la mencionada parte del sólido es
\begin{equation} \label{hiperesfera1}
\int_{0}^{1}\int_{0}^{1}\int_{0}^{2\pi }2u\sqrt{16-u^2} dvdudw= \frac{4}{3}(64 - 7\sqrt{7})\pi
\end{equation}


\begin{tcolorbox}
Por otra parte, la región entre las dos circunferencias de la figura \ref{figura168}, se describe formando segmentos rectilíneos que unan dos puntos $P=\left(\frac{5}{\sqrt{2}}\sin v,\frac{5}{\sqrt{2}}\cos v\right)$
y $Q=\left(3 \sin v,3\cos v\right)$. Esto se consigue simplificando una combinación convexa de la forma
$(1-u)P +uQ$, entonces
\begin{eqnarray*}
x &=& \frac{1}{2} \left(-5\sqrt{2}u+6u+5\sqrt{2}\right)\sin(v) \\
z &=& \frac{1}{2}\left(-5\sqrt{2}u+6u+5\sqrt{2}\right)\cos (v)
\end{eqnarray*}
Las siguientes expresiones determinan, respectivamente, la coordenada $y$ sobre la esfera y el hiperboloide.
\begin{eqnarray*}
y_1&=&\sqrt{16-x^2-z^2}=\frac{1}{\sqrt{2}}\sqrt{\left(30 \sqrt{2}-43\right) u^2+\left(50-30 \sqrt{2}\right) u+7} \\
y_2 &=& \sqrt{x^2+z^2-9}=\frac{1}{\sqrt{2}}\sqrt{(1-u) \left(\left(30 \sqrt{2}-43\right) u+7\right)}
\end{eqnarray*}
La expresión $(1-w)y_1+wy_2$ determina la componente en $y$ de la función que parametriza el sólido.
Es por lo anterior que la siguiente función se utilizará como dicha parametrización.
\begin{multline*}
{\bf r}(u,v,w) = \bigg\langle \frac{1}{2} \left(-5\sqrt{2}u+6u+5\sqrt{2}\right)\sin(v),(1-w)y_1+wy_2,\\
\frac{1}{2}\left(-5\sqrt{2}u+6u+5\sqrt{2}\right)\cos (v) \bigg\rangle
\end{multline*}
Puede probarse (usando \verb"Mathematica", por ejemplo) que el elemento diferencial de volumen es
$|{\bf r_w}\cdot({\bf r_u}\times{\bf r_v})|=\bigg\vert\frac{1}{2}\left(30\sqrt{2}u-43u-15\sqrt{2}+25\right)(y_1-y_2)\bigg\vert.$
El volumen de la segunda región es
\begin{eqnarray}
\label{hiperesfera2}
2\int_0^1\int_0^1\int_0^{2\pi}|{\bf r_w}\cdot({\bf r_v}\times{\bf r_u})|dudvdw&=& \frac{14\sqrt{14}}{3}(\sqrt{2}-1)\pi \\
&\approx & 22.722 \nonumber
\end{eqnarray}
Nótese que la integral en (\ref{hiperesfera2}) se multiplicó por $2$ por la simetría del sólido con respecto al plano $xz.$
\end{tcolorbox}
El volumen del sólido es la suma de los resultados obtenidos por las expresiones (\ref{hiperesfera1}) y (\ref{hiperesfera2}).
\vskip0.2cm
Los cálculos anteriores se efectuaron en \verb"Mathematica" utilizando las siguientes instrucciones.

\begin{verbatim}
{x,z}={3*u*Sin[v],3*u*Cos[v]};
r=(1-w)*{x,-Sqrt[16-x^2-z^2],z}+w*{x,Sqrt[16-x^2-z^2],z};
Integrate[Dot[D[r,w],Cross[D[r,u],D[r,v]]],{w,0,1},
{u,0,1},{v,0,2*Pi}]
\end{verbatim}
%
\begin{verbatim}
{x1,z1}=5/Sqrt[2]*{Sin[v],Cos[v]};
{x2,z2}=3*{Sin[v],Cos[v]};
{x,z}=(1-u)*{x1,z1}+u*{x2,z2};
y1=Simplify[Sqrt[16-x^2-z^2]]
y2=Simplify[Sqrt[x^2+z^2-9]]
r={x,(1-w)*y1+w*y2,z};
A=Simplify[Dot[D[r,u],Cross[D[r,v],D[r,w]]]]
2*NIntegrate[A,{w,0,1},{u,0,1},{v,0,2*Pi}]
\end{verbatim}
\end{comment}

\begin{ejer}
Plantee y evalúe las integrales que determinan el volumen del sólido determinado por
\begin{enumerate}[font=\normalfont]
\item \textit{Las expresiones $63x^2+252x+28z^2+56z+36y=8$ y $9x^2+36x+16+4z^2+8z-12y=0$,
En las órdenes $dxdydz$ y $dzdxdy.$}
\item \textit{$x^2 + y^2 + z^2 \leq 16$ y $x^2 + z^2 \leq 9 + y^2$; en los órdenes $dydxdz$, $dzdydx$ $dxdydz$ y $dxdzdy.$}
Parametrice el sólido y obtenga el mismo valor por este método.
\end{enumerate}
\end{ejer}

\section{Un semicono y una semiesfera} \index{Semiesfera}
El semicono con ecuación $z=1-\sqrt{x^2+y^2}$ y la semiesfera determinada por la expresión
$z=-\sqrt{1-x^2-y^2}$, se intersecan en una circunferencia de radio $1,$ centrada en el origen.
La figura de este sólido se obtiene en \verb"Mathematica" con la siguiente instrucción:


\vspace{0.6em}
\noindent
\begin{tcolorbox}[breakable,enhanced,title=\bf Instrucción en \verb"Mathematica", top=2pt,bottom=2pt,boxsep=0pt,before skip=2pt,after skip=0pt]
\begin{Verbatim}[formatcom=\letraverbatim]
RegionPlot3D[x^2+y^2+z^2<=1&&z<1-Sqrt[x^2+y^2]&&x^2
+y^2<=1,{x,-1,1},{y,-1,1},{z,-1,1},PlotPoints->90,
Mesh->None]
ContourPlot3D[{x^2+y^2+z^2==1,(z-1)^2==x^2+y^2},
{x,-2,2},{y,-2,2},{z,-1,1},PlotPoints->90,Mesh->None,
BoxRatios->{1,1,1/2},ContourStyle->Opacity[0.85]]
\end{Verbatim}
\end{tcolorbox}
\vspace{0.7em}
\noindent
\vskip0.2cm

El sólido mencionado es el siguiente
\vskip0.2cm

\noindent
\begin{minipage}{\textwidth}
	\centering
	\captionof{figure}{Sólido determinado por $-\sqrt{1-x^2-y^2} \leq z\leq 1-\sqrt{x^2+y^2}.$}
	 \label{fig:solido_eps}
	\begin{minipage}{0.28\linewidth}
		\centering
		\includegraphics[width=\linewidth]{Img/13c.pdf}
	\end{minipage}
	\hspace{1em}
	\begin{minipage}{0.20\linewidth}
		\centering
		\includegraphics[width=\linewidth]{Img/13b.pdf}
	\end{minipage}
	\hspace{1em}
	\begin{minipage}{0.20\linewidth}
		\centering
		\includegraphics[width=\linewidth]{Img/13a.pdf}
	\end{minipage}
	
	\vspace{0.3em}
	\fuentepropia
\end{minipage}

\vskip0.2cm
\noindent
\begin{minipage}{\textwidth}
	\centering
	\captionof{figure}{Proyecciones del sólido en los planos coordenados.}
	 \label{fig:proyecciones}
	\begin{minipage}{0.32\linewidth}
		\centering
		\vspace{-10mm}
		\includegraphics[width=\linewidth]{Fig/355.pdf}
	\end{minipage}
	\hspace{1em}
	\begin{minipage}{0.32\linewidth}
		\centering
		\vspace{-10mm}
		\includegraphics[width=\linewidth]{Fig/356.pdf}
	\end{minipage}
	\begin{minipage}{0.28\linewidth}
		\centering
		\vspace{-10mm}
		\includegraphics[width=\linewidth]{Fig/357.pdf}
	\end{minipage}
	
	\vspace{-10mm}
	\fuentepropia
\end{minipage}

\vskip0.2cm
Con estas figuras se establecen las integrales triples, que determinan el volumen del sólido:
{\small
\[
\int_{-1}^{1} \int_{-\sqrt{1 - x^2}}^{\sqrt{1 - x^2}} \int_{-\sqrt{1 - x^2 - y^2}}^{1 - \sqrt{x^2 + y^2}} dz\,dy\,dx
=
\int_{-1}^{1} \int_{-\sqrt{1 - y^2}}^{\sqrt{1 - y^2}} \int_{-\sqrt{1 - x^2 - y^2}}^{1 - \sqrt{x^2 + y^2}} dz\,dx\,dy
= \pi.
\]
}

\vskip0.3cm
De la ecuación del cono se sigue que
{\small $y=\pm \sqrt{\left( 1-z\right) ^2-x^2}$} y con la ecuación de la semiesfera se obtiene {\small $y=\pm \sqrt{1-x^2-z^2}.$} Lo anterior permite obtener el volumen del sólido planteando la integral triple en el orden $dydxdz,$

$\ds \int_{0}^{1}\int_{z-1}^{1-z}\int_{-\sqrt{\left( 1-z\right)^2-x^2}}^{\sqrt{\left( 1-z\right) ^2-x^2}}dydxdz+\int_{-1}^{0}\int_{-\sqrt{1-z^2}}^{\sqrt{1-z^2}}\int_{-\sqrt{1-x^2-z^2}}^{\sqrt{1-x^2-z^2}}dydxdz=\pi. $
\vskip0.2cm
Esta integral se evalúa con \verb"Mathematica", ejecutando la siguiente instrucción.

\vspace{0.6em}
\noindent
\begin{tcolorbox}[breakable,enhanced,title=\bf Instrucción en \verb"Mathematica", top=2pt,bottom=2pt,boxsep=0pt,before skip=2pt,after skip=0pt]
\begin{Verbatim}[formatcom=\letraverbatim]
Integrate[1,{z,0,1},{x,z-1,1-z},{y,-Sqrt[(1-z)^2-x^2],
Sqrt[(1-z)^2-x^2]}]+Integrate[1,{z,-1,0},
{x,-Sqrt[1-z^2],Sqrt[1-z^2]},{y,-Sqrt[1-x^2-z^2],
Sqrt[1-z^2-x^2]}]
\end{Verbatim}
\end{tcolorbox}
\vspace{0.7em}
\noindent


En el caso del orden de integración $dydzdx,$, la primera de las integrales anteriores debe expresarse como la suma de dos integrales triples, es decir,
\vspace{-4mm}
{\small
\begin{multline*}
\int_{-1}^{0}\int_{0}^{1+x}\int_{-\sqrt{(1-z)^2-x^2}}^{\sqrt{(1-z)^2-x^2}}dydzdx+\int_{0}^{1}\int_{0}^{1-x}\int_{-\sqrt{\left( 1-z\right) ^2-x^2}}^{\sqrt{\left( 1-z\right) ^2-x^2}}dydzdx\\
+\int_{-1}^{1}\int_{-\sqrt{1-z^2}}^{\sqrt{1-z^2}}\int_{-\sqrt{1-x^2-z^2}}^{\sqrt{1-x^2-z^2}}dydzdx = \pi.
\end{multline*}
}
El sólido es simétrico con respecto al plano $xz$, las integrales en el orden $dydxdz$ y $dxdydz$ son idénticas.
En coordenadas cilíndricas este valor está dado por
{\small$\ds \int_{0}^{2\pi }\int_{0}^{1}\int_{-\sqrt{1-r^2}}^{1-r}rdzdrd\theta = \pi.$}
\vskip0.2cm

%\begin{tcolorbox}[title=\bfseries Parametrizando el sólido]
La intersección de las superficies está dada por
$x^2+y^2=1,$ que se parametriza con
$x= u\cos v,$ $y=u\sin v.$ Para la esfera $z= -\sqrt{1-u^2}$
y con el semicono se hace $z=1-u.$ Una combinación convexa de la forma
$\left( 1-w\right) \left( u\cos v,u\sin v,-\sqrt{1-u^2}\right) +w\left(u\cos v,u\sin v,1-u\right).$
Se define una parametrización del sólido mediante:
\vspace{-2mm}
\[
\mathbf{r}(u, v, w) =
\left\langle
u\cos v,\ 
u\sin v,\ 
\sqrt{1 - u^2}(w - 1) - w(u - 1)
\right\rangle.
\]

Con los parámetros en los siguientes rangos:

\[
u \in [0, 1], \quad v \in [0, 2\pi], \quad w \in [0, 1].
\]

Derivando parcialmente, se obtiene el elemento diferencial de volumen.
%\begin{eqnarray*} {\bf r}_{u}&=& \langle \cos(v),\sin(v),\frac{(1-w)u}{\sqrt{1-u^2}}-w\rangle \\
%{\bf r}_{v}&=&\langle -u\sin (v) ,u\cos (v) ,0\rangle \\ {\bf r}_{w}&=&\langle 0,0,\sqrt{1-u^2}-u+1\rangle . \end{eqnarray*}
$\left| {\bf r}_{w}\cdot \left( {\bf r}_{u}\times {\bf r}_{v}\right) \right| = u\left( \sqrt{1-u^2}-u+1\right),$
cuya integración nos proporciona el volumen del sólido,
%\vskip0.2cm
{\small
\[\int_{0}^{1}\int_{0}^{2\pi}\int_{0}^{1}u\left( \sqrt{1-u^2}-u+1\right) dudvdw= \pi.\]}
%\end{tcolorbox}

\section{Un elipsoide y un hiperboloide
} \index{Elipsoide} \index{Hiperboloide}

Las ecuaciones $x^2 + z^2 = 1 + y^2$ y $x^2 + y^2 + 4z^2 = 16$ representan un elipsoide y un hiperboloide de un manto; presentamos el sólido acotado por estas.


\noindent
\begin{minipage}{\textwidth}
	\centering
	\captionof{figure}{Sólido acotado por el elipsoide y el hiperboloide.}
	 \label{fig:elipsoide_hiperboloide}
	\begin{minipage}{0.20\linewidth}
		\centering
		\vspace{-2mm}
		\includegraphics[width=\linewidth]{Img/60c.pdf}
	\end{minipage}
	\hspace{1em}
	\begin{minipage}{0.20\linewidth}
		\centering
		\vspace{-2mm}
		\includegraphics[width=\linewidth]{Img/60a.pdf}
	\end{minipage}
	\hspace{1em}
	\begin{minipage}{0.20\linewidth}
		\centering
		\vspace{-2mm}
		\includegraphics[width=\linewidth]{Img/60aa.pdf}
	\end{minipage}
	
	\vspace{0.3em}
	\fuentepropia
\end{minipage}

\vskip0.2cm
El sólido se proyecta en los planos coordenados con el propósito de plantear las integrales que determinan su volumen.
\vskip0.2cm
\noindent
\begin{minipage}{\textwidth}
	\centering
	\captionof{figure}{El sólido proyectado en los planos $xy$ y $yz$.}
	\vspace{-15mm}
	\begin{minipage}{0.37\linewidth}
		\centering
		\vspace{-4mm}
		\includegraphics[width=\linewidth]{Fig/358.pdf}
	\end{minipage}
	\hspace{1em}
	\begin{minipage}{0.37\linewidth}
		\centering
		\vspace{-4mm}
		\includegraphics[width=\linewidth]{Fig/359.pdf}
	\end{minipage}
	
	%\vspace{0.3em}
	
	\vspace{-12mm}
	\fuentepropia
\end{minipage}
\vskip0.2cm


\begin{center}
	\begin{minipage}{0.40\linewidth} % Le doy un ancho adecuado para la figura
		\centering
		\captionof{figure}{Figura 6.X: Título descriptivo.} 
		\vspace{-19mm}
		\includegraphics[width=\linewidth]{Fig/360.pdf}
		\vspace{-1.2\baselineskip}
		\vspace{-2cm}
		\parbox{\linewidth}{\centering\fontsize{8pt}{10pt}\selectfont\textbf{Fuente:} elaboración propia.}
	\end{minipage}
\end{center}

\vskip0.2cm
\begin{ejer}
Plantee y evalúe las integrales que determinan el volumen del sólido determinado por las ecuaciones $x^2 + z^2 = 1 + y^2$ y $x^2 + y^2 + 4z^2 = 16.$
\end{ejer}


\section{Cilindros parabólicos} \index{Sólidos!cilindros parabólicos}
\vskip0.5cm
\subsection*{Dos cilindros parabólicos y un plano II}

Las ecuaciones $z=4-y^2$ y $y=2-x^2$ determinan dos cilindros parabólicos.
Estas superficies, junto con el plano $z=0$ acotan el sólido que se muestra a continuación, y del cual plantearemos las integrales que se requiere evaluar para determinar su volumen.

\vskip0.2cm
\noindent
\begin{minipage}{\textwidth}
	\centering
	\captionof{figure}{$0 \leq z \leq 4 - y^2$, \quad $y \leq 2 - x^2$.}
	\label{fig:proyeccion_xyz}
	\begin{minipage}{0.23\linewidth}
		\centering
		\includegraphics[width=\linewidth]{Img/34d.pdf}
	\end{minipage}
	\begin{minipage}{0.23\linewidth}
		\centering
		\includegraphics[width=\linewidth]{Img/34a.pdf}
	\end{minipage}
	\begin{minipage}{0.23\linewidth}
		\centering
		\includegraphics[width=\linewidth]{Img/34c.pdf}
	\end{minipage}
	
	\vspace{0.3em}
	\fuentepropia
\end{minipage}
\vskip0.2cm


Las figuras mostradas se obtienen con la ayuda de los siguientes comandos.

\vspace{0.6em}
\noindent
\begin{tcolorbox}[breakable,enhanced,title=\bf Instrucción en \verb"Mathematica", top=2pt,bottom=2pt,boxsep=0pt,before skip=2pt,after skip=0pt]
\begin{Verbatim}[formatcom=\letraverbatim]
ContourPlot3D[{z==4-y^2,2-x^2==y},{x,-2,2},{y,-2,2},
{z,0,4},ContourStyle->{Blend[{Yellow,Red}],
RGBColor[0.3,0.6,0.4]},Mesh->None,PlotPoints->60]
RegionPlot3D[0<z<4-y^2&&-Sqrt[2-y]<x<Sqrt[2-y]&&-2
<y<2,{x,-2,2},{y,-2,2},{z,0,4},PlotPoints->90,
Mesh->None,PlotStyle->Orange]
\end{Verbatim}
\end{tcolorbox}
\vspace{0.7em}
\noindent


Esta figura la proyectaremos sobre planos paralelos a los planos coordenados, con el fin de hacer sencilla la determinación de los límites de integración.
\vskip0.2cm
Integrando primero con respecto a $z,$ los puntos recorridos van desde $z=0$ hasta el cilindro $z=4-y^2$ y los límites para $x$ e $y$ se pueden obtener del segundo de los siguientes gráficos. Además, integrando primero con respecto a $x,$
se recorren puntos desde la parte negativa del cilindro $y = 2 - x^2$
hasta la parte positiva del mismo, pues el cilindro $z=4-y^2$ es paralelo al eje $x.$
\vskip0.2cm

\noindent
\begin{minipage}{\textwidth}
	\centering
	\captionof{figure}{Proyecciones del sólido en los planos $yz$ y $xy$.}
	\label{fig:proyecciones_yz_xy}
	\begin{minipage}{0.35\linewidth}
		\centering
		\vspace{-18mm}
		\includegraphics[width=\linewidth]{Fig/361.pdf}
	\end{minipage}
	\hspace{1em}
	\begin{minipage}{0.40\linewidth}
		\centering
		\vspace{-18mm}
		\includegraphics[width=\linewidth]{Fig/362.pdf}
	\end{minipage}
	
	
	\vspace{-15mm}
	\fuentepropia
\end{minipage}
\vskip0.2cm

Con la información observada en estos gráficos, las siguientes integrales triples determinan el volumen del sólido.

\begin{eqnarray*}
\int_{-2}^2\int_{-2}^{2-x^2}\int_{0}^{4-y^2}dzdydx=
\int_{-2}^2\int_{-\sqrt{2-y}}^{\sqrt{2-y}}\int_{0}^{4-y^2}dzdxdy &=& \frac{1024}{35} \\
\int_{0}^{4}\int_{-\sqrt{4-z}}^{\sqrt{4-z}}\int_{-\sqrt{2-y}}^{\sqrt{2-y}}dxdydz=
\int_{-2}^2\int_{0}^{4-y^2}\int_{-\sqrt{2-y}}^{\sqrt{2-y}}dxdzdy&=& \frac{1024}{35}.
\end{eqnarray*}
\vskip0.3cm


La proyección en el plano $xz$ merece varias consideraciones y la región queda dividida en dos subregiones; en una de ellas y recorre los puntos desde el cilindro
$z =4-y^2$ hasta el cilindro $y=2-x^2$ y en la otra subregion $y$ recorre puntos desde la parte negativa del cilindro $z=4-y^2$ hasta la parte positiva del mismo cilindro.


\vskip0.2cm
\noindent
\begin{minipage}{\textwidth}
	\centering
	\captionof{figure}{Proyección en el plano $xz$}
	\begin{minipage}{0.38\linewidth}
		\centering
		\vspace{-18mm}
		\includegraphics[width=\linewidth]{Fig/363.pdf}
	\end{minipage}
	
	\vspace{-14mm}
	\fuentepropia
\end{minipage}
\vskip0.2cm



Además, utilizando  la figura anterior, es pertinente considerar que, si en la ecuación, se sustituye en la otra ecuación, se obtiene lo siguiente:
\[z=4-y^2=4-(2-x^2)^2=4x^2-x^4.\]
De esta expresión se sigue que $\,\,2-x^2=\pm\sqrt{4-z},$ por lo tanto $2\pm \sqrt{4-z}=x^2$ y de ahí se obtiene las cuatro expresiones que determinan a $x,$ estas ecuaciones 
$ x = \pm \sqrt{2 \pm \sqrt{4-z}}.$
En la siguiente figura se ilustra esta situación.

\vskip0.2cm
\noindent
\begin{minipage}{\textwidth}
	\centering
	\captionof{figure}{Proyección del sólido en el plano $xz$.}
	\label{fig:proyeccion_xz}
	\begin{minipage}{0.60\linewidth}
		\centering
		\vspace{-43mm}
		\includegraphics[width=\linewidth]{Fig/364.pdf}
	\end{minipage}

	\vspace{-38mm}
	\fuentepropia
\end{minipage}
\vskip0.2cm

Ahora se presenta la integral triple faltante. Su valor fue necesario hallarlo de manera numérica utilizando un asistente computacional.

\[\int_{-\sqrt{2}}^{\sqrt{2}}\int_{4x^2-x^{4}}^{4}\int_{-\sqrt{4-z}}^{\sqrt{4-z}}dydzdx+\int_{-2}^2\int_{0}^{4x^2-x^{4}}\int_{-\sqrt{4-z}}^{2-x^2}dydzdx \approx 29.2571\]
\begin{multline*}
\int_{0}^{4}\int_{-\sqrt{2-\sqrt{4-z}}}^{\sqrt{2-\sqrt{4-z}}}\int_{-\sqrt{4-z}}^{\sqrt{4-z}}dydxdz+\int_{0}^{4}\int_{-\sqrt{2+\sqrt{4-z}}}^{-\sqrt{2-\sqrt{4-z}}}\int_{-\sqrt{4-z}}^{2-x^2}dydxdz\\
+\int_{0}^{4}\int_{\sqrt{2-\sqrt{4-z}}}^{\sqrt{2+\sqrt{4-z}}}\int_{-\sqrt{4-z}}^{2-x^2}dydxdz \approx 29.2571.
\end{multline*}

La evaluación numérica de las integrales anteriores se consigue ejecutando las siguientes instrucciones:

\vspace{0.6em}
\noindent
\begin{tcolorbox}[breakable,enhanced,title=\bf Instrucción en \verb"Mathematica", top=2pt,bottom=2pt,boxsep=0pt,before skip=2pt,after skip=0pt]
\begin{Verbatim}[formatcom=\letraverbatim]
NIntegrate[1,{x,-Sqrt[2],Sqrt[2]},{z,4*x^2-x^4,4},
{y,-Sqrt[4-z],Sqrt[4-z]}]+NIntegrate[1,{x,-2,2},
{z,0,4*x^2-x^4},{y,-Sqrt[4-z],2-x^2}]
NIntegrate[1,{z,0,4},{x,-Sqrt[2-Sqrt[4-z]],Sqrt[2-
Sqrt[4-z]]},{y,-Sqrt[4-z],Sqrt[4-z]}]+NIntegrate[1,
{z,0,4},{x,-Sqrt[2+Sqrt[4-z]],-Sqrt[2-Sqrt[4-z]]},
{y,-Sqrt[4-z],2-x^2}]+NIntegrate[1,{z,0,4},{x,Sqrt[2-
Sqrt[4-z]],Sqrt[2+Sqrt[4-z]]},{y,-Sqrt[4-z],2-x^2}]
\end{Verbatim}
\end{tcolorbox}
\vspace{0.7em}
\noindent


Otra manera de hallar el volumen del sólido es la siguiente.
%\vskip0.2cm

%\begin{tcolorbox}[title=\bfseries Parametrizando el sólido]
En el plano $xy$, la proyección del sólido se parametriza mediante con una combinación convexa de dos puntos
$A(t,-2)$ y $B(t,2-t^2)$; esto es $(1-u)A+uB=(t,(4-t^2)u-2).$
Con un punto $P\left(t,2-t^2-4u+ut^2,0\right) $ en la base del sólido y un punto $Q\left( t,(4-t^2)u-2,u(4-t^2)((t^2-4)u+4)\right)$
sobre el cilindro se forma la combinación $(1-w)P+wQ,$ que define la función
{\small\[{\bf r}(u,v,w) = \langle t,(4-t^2)u-2,-(t^2-4)uw((t^2-4)u+4)\rangle,\]}
que parametriza el sólido con $t\in [ -2,2]$, $u\in [0,1]$, $w\in [0,1].$

\vskip0.3cm
De esta función se obtienen los siguientes vectores
{\small
\begin{eqnarray*}
{\bf r}_{t}&=& \langle 1,2t(u-1),4wt(1-u)((u-1)t^2-4u+2) \rangle\\
{\bf r}_{u} &=& \langle 0,t^2-4,w(t^2-4) ((1-u)t^2+4u) +w(1-u)(t^2-4)^2 \rangle \\
{\bf r}_{w} &=& \langle 0,0,(u-1)( t^2-4) \left(t^2-ut^2+4u\right) \rangle.
\end{eqnarray*}
}
Con lo anterior, el elemento diferencial de volumen está dado por {\small \[\left| {\bf r}_{t}\cdot \left({\bf r}_{u}\times {\bf r}_{w}\right) \right|
=\left| \left( t^2-4\right) ^2\left( u-1\right) \left(t^2-ut^2+4u\right) \right|,\]}

entonces, el volumen del sólido esta dado por

\[\int_{0}^{1}\int_{0}^{1}\int_{-2}^2\left| \left( t^2-4\right)
^2\left( u-1\right) \left( t^2-ut^2+4u\right) \right|
dtdudw= \frac{1024}{35}.\]
%\end{tcolorbox}
\vskip0.2cm
Con las siguientes instrucciones, se realiza el procedimiento anterior:

\begin{tcolorbox}
\begin{Verbatim}[formatcom=\letraverbatim]
{x,y}=u*{t,-2}+(1-u)*{t,2-t^2};
r=Simplify[(1-w)*{x,y,0}+w*{x,y,4-y^2}]
A=Abs[Dot[D[r,t],Cross[D[r,u],D[r,w]]]];
Integrate[A,{t,-2,2},{u,0,1},{w,0,1}]
\end{Verbatim}
\end{tcolorbox}

\subsection*{Un cilindro parabólico y uno trigonométrico}

Hallar el volumen del sólido acotado por las figuras de las expresiones $z=1-x^2$ y $z=\cos(y)$, con
$y\in [0,2\pi]$. Estas superficies y el sólido que determinan se muestran enseguida.
\vskip0.2cm

\noindent
\begin{minipage}{\textwidth}
	\centering
	\captionof{figure}{$\cos y \leq z \leq 1 - x^2$, \quad $y \in [0, 2\pi]$.}
	\label{fig:region_cos_1mx2}
	\begin{minipage}{0.21\linewidth}
		\centering
		\includegraphics[width=\linewidth]{Img/84b.pdf}
	\end{minipage}
	\hspace{1em}
	\begin{minipage}{0.21\linewidth}
		\centering
		\includegraphics[width=\linewidth]{Img/84c.pdf}
	\end{minipage}
	\hspace{1em}
	\begin{minipage}{0.21\linewidth}
		\centering
		\includegraphics[width=\linewidth]{Img/84d.pdf}
	\end{minipage}
	
	\vspace{0.3em}
	\fuentepropia
\end{minipage}


\vskip0.2cm
Con las siguientes instrucciones se consiguen estas figuras:

\vspace{0.6em}
\noindent
\begin{tcolorbox}[breakable,enhanced,title=\bf Instrucción en \verb"Mathematica", top=2pt,bottom=2pt,boxsep=0pt,before skip=2pt,after skip=0pt]
\begin{Verbatim}[formatcom=\letraverbatim]
ContourPlot3D[{z==1-x^2,z==Cos[y]},{x,-Sqrt[2],
Sqrt[2]},{y,0,2*Pi},{z,-1,1},Mesh->False,
ContourStyle->{Orange,Green}]
RegionPlot3D[Cos[y]<z<1-x^2&&-2<x<2,{x,-Sqrt[2],
Sqrt[2]},{y,0,2*Pi},{z,-1,1},PlotPoints->30,
Mesh->False]
\end{Verbatim}
\end{tcolorbox}
\vspace{0.7em}
\noindent

\noindent
\begin{minipage}{\textwidth}
	\centering
	\captionof{figure}{Proyecciones del sólido en los planos coordenados.}
	\label{fig:proyecciones_coordenadas_2}
	\begin{minipage}{0.26\linewidth}
		\centering
		\vspace{-14mm}
		\includegraphics[width=\linewidth]{Fig/365.pdf}
	\end{minipage}
	\hspace{1em}
	\begin{minipage}{0.26\linewidth}
		\centering
		\vspace{-14mm}
		\includegraphics[width=\linewidth]{Fig/366.pdf}
	\end{minipage}
	\hspace{1em}
	\begin{minipage}{0.26\linewidth}
		\centering
		\vspace{-14mm}
		\includegraphics[width=\linewidth]{Fig/367.pdf}
	\end{minipage}
	
	\vspace{-8mm}
	\fuentepropia
\end{minipage}


Con estas figuras se plantean las integrales que determinan el volumen del sólido:

\vspace{-8mm}
\begin{multline*}
\int_{0}^{2\pi}\int_{-\sqrt{1-\cos y}}^{\sqrt{1-\cos y}}\int_{\cos y}^{1-x^2}dzdxdy= \int_{-\sqrt{2}}^{\sqrt{2}}\int_{\arccos \left( 1-x^2\right) }^{2\pi -\arccos \left( 1-x^2\right) }\int_{\cos y}^{1-x^2}dzdydx \\
\int_{-1}^{1}\int_{-\sqrt{1-z}}^{\sqrt{1-z}}\int_{\arccos z}^{2\pi -\arccos z}dydxdz= \int_{-\sqrt{2}}^{\sqrt{2}}\int_{-1}^{1-x^2}\int_{\arccos z}^{2\pi -\arccos z}dydzdx \\
\int_{0}^{2\pi }\int_{\cos y}^{1}\int_{-\sqrt{1-z}}^{\sqrt{1-z}}dxdzdy=\int_{-1}^{1}\int_{\arccos z}^{2\pi -\arccos z}\int_{-\sqrt{1-z}}^{\sqrt{1-z}}dxdydz = \frac{64\sqrt{2}}{9}.
\end{multline*}

Las anteriores integrales se evaluaron usando \verb"Mathematica";
el código requerido en tres de ellas, se presenta a continuación.
%\vskip0.2cm

\vspace{0.6em}
\noindent
\begin{tcolorbox}[breakable,enhanced,title=\bf Instrucción en \verb"Mathematica", top=2pt,bottom=2pt,boxsep=0pt,before skip=2pt,after skip=0pt]
\begin{Verbatim}[formatcom=\letraverbatim]
Integrate[1,{y,0,2*Pi},{x,-Sqrt[1-Cos[y]],
Sqrt[1-Cos[y]]},{z,Cos[y],1-x^2}]
Integrate[1,{z,-1,1},{x,-Sqrt[1-z],Sqrt[1-z]},
{y,ArcCos[z],2*Pi-ArcCos[z]}]
Integrate[1,{y,0,2*Pi},{z,Cos[y],1},{x,-Sqrt[1-z],
Sqrt[1-z]}]
\end{Verbatim}
\end{tcolorbox}
\vspace{0.7em}
\noindent

El volumen también se encuentra con el siguiente procedimiento.
\vskip0.2cm
% \begin{figure}[H]
% \centering
% \begin{pspicture}(-2,-0.6)(2,3.7) % \psgrid
% \psset{algebraic,xunit=0.9,yunit=0.4}
% \pscustom[fillstyle=solid,fillcolor=gray!40]{
%   \parametricplot{0}{\psPiTwo}{sqrt(1-cos(t)),t}
%   \parametricplot{\psPiTwo}{0}{-sqrt(1-cos(t)),t}
% }
% \psaxes[arrows=->,showorigin=false,dx=1.4,xLabels={-\sqrt{2},,,\sqrt{2}},yunit=\psPi,ytrigLabels,trigLabelBase=1](0,0)(-2,-0.3)(2.35,2.6)[$x$,90][$y$,180]
% \parametricplot[arrows=->,linewidth=1.2pt]{-0.3}{7}{sqrt(1-cos(t)),t}
% \parametricplot[arrows=->,linewidth=1.2pt]{-0.3}{7}{-sqrt(1-cos(t)),t}
% \rput{0}(-1.31,2.36){$\bullet$}
% \rput{0}(1.31,2.36){$\bullet$}
% \psline[linewidth=1pt,linecolor=black](-1.31,2.4)(1.31,2.4)
% \rput{0}(-1.65,2.5){$A$}
% \rput{0}(1.65,2.5){$B$}
% \end{pspicture}
% \caption{$\cos(y) = 1 - x^2$.}
% \label{fig:curva_cosyx2}
% \end{figure}
% \par
% \justifying
%\begin{tcolorbox}[title=\bfseries Parametrizando el sólido]
\noindent 
Con dos puntos $A\left( -\sqrt{1-\cos t},t\right) $
y $B\left( \sqrt{1-\cos t},t\right)$ en la proyección del sólido en el plano $xy$,
se forma la expresión
$(1-u)A+uB,$ que parametriza la región mencionada. % esto es $ \left((1-2u)\sqrt{1-\cos t},t\right).$
Sobre las dos superficies que determinan el sólido se satisface que
$z= 4u\left( 1-u\right) +\left( 1-2u\right)^2\cos t$ y $z=\cos (t).$


\vskip0.2cm
\noindent
\begin{minipage}{\textwidth}
	\centering
	\captionof{figure}{$\cos(y) = 1 - x^2$}
	\label{fig:curva_cosyx2}
	\begin{minipage}{0.25\linewidth}
		\centering
		\vspace{-8mm}
		\includegraphics[width=\linewidth]{Fig/368.pdf}
	\end{minipage}
	
	\vspace{-4mm}
	\fuentepropia
\end{minipage}
\vskip0.2cm


El sólido se parametriza mediante una función $\mathbf{r}$, que es combinación convexa con dos puntos de la forma $P\left((2u - 1)\sqrt{1 - \cos t}, t, \cos t\right)$ y $Q\left((1 - 2u)\sqrt{1 - \cos t}, 4u(1 - u) + (1 - 2u)^2 \cos t\right)$. Esto es
%\vskip0.2cm
{\small
\[
\begin{aligned}
	\mathbf{r}(t, u, w) = \left\langle (1 - 2w)(2u - 1)\sqrt{1 - \cos t}, t, 4uw(1 - u) + \right. \\
	\left. (1 - 4wu + 4wu^2) \cos t \right\rangle.
\end{aligned}
\]
}

%\vskip0.2cm
Con esta función hallamos los siguientes vectores
{\small
\begin{eqnarray*}
{\bf r}_{t} &=& \left\langle \frac{1}{2}\frac{\left(
1-2w\right) \left( 2u-1\right) }{\sqrt{1-\cos t}}\sin t,1,\left(
4wu-4wu^2-1\right) \sin t\right\rangle \\
{\bf r}_{u}&=& \left\langle 2\left( 1-2w\right) \sqrt{1-\cos
t},0,4w\left( 1-2u\right) \left( 1-\cos t\right) \right\rangle \\
{\bf r}_{w}&=& \left\langle 2\left( 1-2u\right) \sqrt{1-\cos t},0,4u\left( u-1\right) \left( \cos t-1\right) \right\rangle
\end{eqnarray*}
}

Además, $\left| {\bf r}_{w}\cdot ( {\bf r}_{t}\times {\bf r}_{u}) \right| =8(w+u(2w+1)(u-1))\sqrt{\left( 1-\cos t\right) ^{3}};$
el volu\-men del sólido esta dado por
\vspace{-2mm}
{\small
\[\int_{0}^{1}\int_{0}^{1}\int_{0}^{2\pi }8(w+u(2w+1)(u-1)) \sqrt{\left( 1-\cos t\right) ^{3}}%
dtdudw= \frac{64}{9}\sqrt{2}\]
}
%\end{tcolorbox}
%\vskip0.2cm
La rutina descrita antes se puede hacer en \verb"Mathematica" con ayuda de las siguientes instrucciones.

\vspace{0.6em}
\noindent
\begin{tcolorbox}[breakable,enhanced,title=\bf Instrucción en \verb"Mathematica", top=2pt,bottom=2pt,boxsep=0pt,before skip=2pt,after skip=0pt]
\begin{Verbatim}[formatcom=\letraverbatim]
{x,y,z}={(2*u-1)*Sqrt[1-Cos[t]],t,1-x^2};
r=(1-w)*{x,y,z}+w*{x,y,Cos[t]};
A=Dot[D[r,t],Cross[D[r,u],D[r,w]]];
Integrate[A,{t,0,2*Pi},{u,0,1},{w,0,1}]
\end{Verbatim}
\end{tcolorbox}
\vspace{0.7em}
\noindent


\subsection*{Dos cilindros parabólicos y un plano}
Los cilindros parabólicos con ecuaciones
$z= 4x-x^2,$ $z= 4y-y^2$ y el plano $z=0,$ acotan un sólido en el primer octante, su gráfico se muestra a continuación.

\vskip0.2cm
\noindent
\begin{minipage}{\textwidth}
	\centering
	\captionof{figure}{$0 \leq z \leq 4x - x^2$, \quad $0 \leq z \leq 4y - y^2$.}
	\label{fig:parabolas_z_2}
	\begin{minipage}{0.22\linewidth}
		\centering
		\includegraphics[width=\linewidth]{Img/10b.pdf}
	\end{minipage}
	\hspace{1em}
	\begin{minipage}{0.22\linewidth}
		\centering
		\includegraphics[width=\linewidth]{Img/10a.pdf}
	\end{minipage}
	
	\vspace{0.5mm}
	\fuentepropia
\end{minipage}
\vskip0.2cm

Las instrucciones para obtener los gráficos del sólido que se consideran en este ejemplo, son las siguientes:

\vspace{0.6em}
\noindent
\begin{tcolorbox}[breakable,enhanced,title=\bf Instrucción en \verb"Mathematica", top=2pt,bottom=2pt,boxsep=0pt,before skip=2pt,after skip=0pt]
\begin{Verbatim}[formatcom=\letraverbatim]
ContourPlot3D[{z==4*x-x^2,z==4*y-y^2},{x,0,4},{y,0,4},
{z,0,4},PlotPoints->80,ContourStyle->{Yellow,Green}]
RegionPlot3D[0<z<4*x-x^2&&0<z<4*y-y^2,{x,0,4},{y,0,4},
{z,0,4},PlotStyle->Blend[{Red,Orange,Gray}],
PlotPoints->90,Mesh->False]
\end{Verbatim}
\end{tcolorbox}
\vspace{0.7em}
\noindent


En la base de la región, es decir en el plano $xy$, se satisface que
$4x-x^2=4y-y^2$ de donde $(x-2)^2=(y-2)^2$, con lo cual $y-2= \pm(x-2)$, esto nos conduce a que $y=x$ o $y=4-x$. Esto divide la base en cuatro partes, por tal razón al integrarse en el orden $dzdydx$
se puede plantear una integral triple y multiplicarla por cuatro.

\vskip0.2cm
\noindent
\begin{minipage}{\textwidth}
	\centering
	\captionof{figure}{Proyección en el plano $xy$ y un cuarto del sólido.}
	\label{fig:proyeccion_xy_cuarto}
	\begin{minipage}{0.35\linewidth}
		\centering
		\vspace{-10mm}
		\includegraphics[width=\linewidth]{Fig/369.pdf}
	\end{minipage}
	\hspace{1em}
	\begin{minipage}{0.25\linewidth}
		\centering
		\vspace{-10mm}
		\includegraphics[width=\linewidth]{Img/10c.pdf}
	\end{minipage}
	
	\vspace{-4mm}
	\fuentepropia
\end{minipage}


\vskip0.2cm
La siguiente instrucción permite obtener la porción señalada:

\vspace{0.6em}
\noindent
\begin{tcolorbox}[breakable,enhanced,title=\bf Instrucción en \verb"Mathematica", top=2pt,bottom=2pt,boxsep=0pt,before skip=2pt,after skip=0pt]
\begin{Verbatim}[formatcom=\letraverbatim]
RegionPlot3D[0<z<2*x-x^2&&0<z<2*y-y^2&&2-x<y<x&&1<x<2,
{x,0,2},{y,0,2},{z,0,1},PlotPoints->90,Mesh->False,
PlotStyle->Orange]
\end{Verbatim}
\end{tcolorbox}
\vspace{0.7em}
\noindent


Cualquiera de las siguientes posibilidades permite calcular el volumen de la figura en mención.

\begin{itemize}
\item $4\ds \int_{2}^{4}\int_{4-x}^{x}\int_{0}^{4x-x^2}dzdydx=32$
\item $2\left( \ds \int_{0}^2\int_{x}^{4-x}\int_{0}^{4x-x^2}dzdydx+\int_{2}^{4}\int_{4-x}^{x}\int_{0}^{4x-x^2}dzdydx\right) =32$
\item $2\left( \ds \int_{0}^2\int_{y}^{4-y}\int_{0}^{4y-y^2}dzdxdy+\int_{2}^{4}\int_{4-y}^{y}\int_{0}^{4y-y^2}dzdxdy\right) =32$
\end{itemize}
\vskip0.2cm

\noindent
\begin{minipage}{\textwidth}
	\centering
	\captionof{figure}{Proyecciones en los planos $xz$ y $yz$.}
	\label{fig:proyecciones_xz_yz}
	\begin{minipage}{0.30\linewidth}
		\centering
		\vspace{-10mm}
		\includegraphics[width=\linewidth]{Fig/370.pdf}
	\end{minipage}
	\hspace{1em}
	\begin{minipage}{0.30\linewidth}
		\centering
		\vspace{-10mm}
		\includegraphics[width=\linewidth]{Fig/371.pdf}
	\end{minipage}
	
	\vspace{-8mm}
	\fuentepropia
\end{minipage}
\vskip0.2cm


Estas figuras permiten el planteamiento de las integrales triples en los órdenes restantes.
\[\int_{0}^{4}\int_{2-\sqrt{4-z}}^{2+\sqrt{4-z}}\int_{2-\sqrt{4-z}}^{2+\sqrt{4-z}}dxdydz= \int_{0}^{4}\int_{0}^{4y-y^2}\int_{2-\sqrt{4-z}}^{2+\sqrt{4-z}}dxdzdy= 32\]
\[\int_{0}^{4}\int_{2-\sqrt{4-z}}^{2+\sqrt{4-z}}\int_{2-\sqrt{4-z}}^{2+\sqrt{4-z}}dydxdz=\int_{0}^{4}\int_{0}^{4x-x^2}\int_{2-\sqrt{4-z}}^{2+\sqrt{4-z}}dydzdx=32\]
\vskip0.2cm
Algunos de estos cálculos pueden hacerse con las siguientes instrucciones:

\vspace{0.6em}
\noindent
\begin{tcolorbox}[breakable,enhanced,title=\bf Instrucción en \verb"Mathematica", top=2pt,bottom=2pt,boxsep=0pt,before skip=2pt,after skip=0pt]
\begin{Verbatim}[formatcom=\letraverbatim]
4*Integrate[1,{x,2,4},{y,4-x,x},{z,0,4*x-x^2}]
Integrate[1,{z,0,4},{x,2-Sqrt[4-z],2+Sqrt[4-z]},
{y,2-Sqrt[4-z],2+Sqrt[4-z]}]
\end{Verbatim}
\end{tcolorbox}
\vspace{0.7em}
\noindent


Este valor también puede obtenerse con el siguiente proceso.


\vskip0.2cm
%\begin{tcolorbox}[title=\bfseries Parametrizando el sólido]
\noindent

En el plano $xz$ se forma una combinación convexa entre dos puntos de la forma
$P(2-\sqrt{4-t}, t)$ y $Q(2 + \sqrt{4-t}, t)$, obteniendo $(2+(2u-1)\sqrt{4-t}, t)$ que parametriza la porción interior a la parábola sobre el eje $z$, con $t \in[0,4],$ $u \in[0,1].$


\vskip0.2cm
\noindent
\begin{minipage}{\textwidth}
	\centering
	\captionof{figure}{$0 \leq y \leq 4x - x^2$}
	\label{fig:region_y_vs_x}
	\begin{minipage}{0.32\linewidth}
		\centering
		\vspace{-15mm}
		\includegraphics[width=\linewidth]{Fig/372.pdf}
	\end{minipage}
	
	\vspace{-10mm}
	\fuentepropia
\end{minipage}
\vskip0.2cm



El sólido se parametriza simplificando una combinación convexa entre dos puntos de la forma
$(2+(2u-1)\sqrt{4-t},2-\sqrt{4-t},t)$ y $(2+(2u-1)\sqrt{4-t},2+\sqrt{4-t},t).$
Por tanto, la parametrización escogida para el sólido es
${\bf r}(u,t,w)= \langle 2+(2u-1)\sqrt{4-t},2+(2w-1)\sqrt{4-t},t \rangle, $
en donde $t\in [0,4],$ $u\in [0,1]$ y $w\in [0,1].$
De esta función se obtiene los siguientes vectores:
{\small
\begin{align*}
	\mathbf{r}_u &= \left\langle 2\sqrt{4-t},0,0 \right\rangle, \\
	\mathbf{r}_w &= \left\langle 0,2\sqrt{4-t},0 \right\rangle, \\
	\mathbf{r}_t &= \left\langle \frac{1-2u}{2\sqrt{4-t}}, \frac{1-2w}{2\sqrt{4-t}},1 \right\rangle.
\end{align*}
}
Entonces, ${\bf r}_w \cdot ({\bf r}_t \times {\bf r}_u)= 4(4-t)$, el volumen del sólido está dado por
%\vskip0.2cm
{\small
\[\int_0^1 \int_0^1 \int_0^4 4(4-t)dtdudw = 32.\]}
%\end{tcolorbox}
\vskip0.2cm
La rutina anteriormente descrita se puede obtener en \verb"Mathematica", ejecutando las siguientes instrucciones:

\vspace{0.6em}
\noindent
\begin{tcolorbox}[breakable,enhanced,title=\bf Instrucción en \verb"Mathematica", top=2pt,bottom=2pt,boxsep=0pt,before skip=2pt,after skip=0pt]
\begin{Verbatim}[formatcom=\letraverbatim]
r=Factor[(1-w){2+(2*u-1)*Sqrt[4-t],2-Sqrt[4-t],t}+
w{2+(2*u-1)*Sqrt[4-t],2+Sqrt[4-t],t}]
Integrate[Dot[D[r,w],Cross[D[r,t],D[r,u]]],{u,0,1},
{t,0,4},{w,0,1}]
\end{Verbatim}
\end{tcolorbox}
\vspace{0.7em}
\noindent


\subsection*{Un cilindro parabólico y un plano inclinado}
Consideremos el sólido acotado por el gráfico de las ecuaciones $y=4-x^2,$ $z=2-y/2,$ $y=0,$ $z=0$,
hallaremos su volumen de varias maneras.
\vskip0.2cm
\noindent
\begin{minipage}{\textwidth}
	\centering
	\captionof{figure}{$0 \leq z \leq 4x - x^2$, \quad $0 \leq z \leq 4y - y^2$.}
	\label{fig:parabolas_z}
	\begin{minipage}{0.25\linewidth}
		\centering
		\vspace{-15mm}
		\includegraphics[width=\linewidth]{Img/T_0a.pdf}
	\end{minipage}
	\hspace{1em}
	\begin{minipage}{0.35\linewidth}
		\centering
		\vspace{-15mm}
		\includegraphics[width=\linewidth]{Fig/373.pdf}
	\end{minipage}
	
	\vspace{-10mm}
	\fuentepropia
\end{minipage}
\vskip0.2cm


Esta región nos lleva a presentar las siguientes figuras.
\vskip0.2cm
\noindent
\begin{minipage}{\textwidth}
	\centering
	\captionof{figure}{Proyecciones en los planos coordenados.}
	\label{fig:proyecciones_coordenadas}
	\begin{minipage}{0.29\linewidth}
		\centering
		\vspace{-10mm}
		\includegraphics[width=\linewidth]{Fig/374.pdf}
	\end{minipage}
	\hspace{1em}
	\begin{minipage}{0.29\linewidth}
		\centering
		\vspace{-10mm}
		\includegraphics[width=\linewidth]{Fig/375.pdf}
	\end{minipage}
	\hspace{1em}
	\begin{minipage}{0.29\linewidth}
		\centering
		\vspace{-10mm}
		\includegraphics[width=\linewidth]{Fig/376.pdf}
	\end{minipage}
	
	\vspace{-5mm}
	\fuentepropia
\end{minipage}
\vskip0.2cm

Las integrales que se evalúan para determinar el volumen del sólido son:
\[\int_{-2}^2\int_{0}^{4-x^2}\int_{0}^{2-y/2}dzdydx = \int_{0}^{4}\int_{-\sqrt{4-y}}^{\sqrt{4-y}}\int_{0}^{2-y/2}dzdxdy= \frac{64}{5}\]
\[\int_{0}^{4}\int_{0}^{2-y/2}\int_{-\sqrt{4-y}}^{\sqrt{4-y}}dxdzdy= \int_{0}^2\int_{0}^{4-2z}\int_{-\sqrt{4-y}}^{\sqrt{4-y}}dxdydz= \frac{64}{5}.\]
\vskip0.2cm
La proyección del sólido en el plano $xz$ es un rectángulo dividido por la curva
$z=x^2/2$. En el orden $dydxdz$ es necesario plantear tres integrales triples
\begin{multline*} \int_{0}^2\int_{-\sqrt{2z}}^{\sqrt{2z}}\int_{0}^{4-2z}dydxdz
+\int_{0}^2\int_{-2}^{-\sqrt{2z}}\int_{0}^{4-x^2}dydxdz \\ +\int_{0}^2\int_{\sqrt{2z}}^2\int_{0}^{4-x^2}dydxdz= \frac{64}{5} 
\end{multline*}
\[\int_{-2}^2\int_{x^2/2}^2\int_{0}^{4-2z}dydzdx +\int_{-2}^2\int_0^{x^2/2}\int_{0}^{4-x^2}dydzdx=\frac{64}{5}.\]
Estas integrales se evalúan en \verb"Mathematica" con las siguientes indicaciones:

\vspace{0.6em}
\noindent
\begin{tcolorbox}[breakable,enhanced,title=\bf Instrucción en \verb"Mathematica", top=2pt,bottom=2pt,boxsep=0pt,before skip=2pt,after skip=0pt]
\begin{Verbatim}[formatcom=\letraverbatim]
Integrate[1,{z,0,2},{y,0,4-2*z},{x,-Sqrt[4-y],
Sqrt[4-y]}]
Integrate[1,{x,-2,2},{z,x^2/2,2},{y,0,4-2*z}]+
Integrate[1,{x,-2,2},{z,0,x^2/2},{y,0,4-x^2}]
\end{Verbatim}
\end{tcolorbox}
\vspace{0.7em}
\noindent


El valor obtenido puede ser calculado con el siguiente método.
\vskip0.2cm

%\begin{tcolorbox}[title=\bfseries Parametrizando el sólido]
La región sobre el plano $xy$ se parametriza uniendo dos puntos, uno sobre el eje $x$ que tiene la forma $(t,0) $ y el otro sobre la curva $y=4-x^2$ que posee la forma $(t,4-t^2),$ con $t\in [-2,2].$
Simplificando $( 1-u) ( t,0) +u( t,4-t^2)$ se sigue, por un lado, que $z=2+y/2= 2-2u+\frac{1}{2}ut^2.$

Por tanto, la función ${\bf r}$ es una combinación convexa entre un punto sobre el plano $xy$
y otros sobre el plano $z=2-y/2;$ simplificando
$(1-w) \langle t,u( 4-t^2),0\rangle +w\langle t,u( 4-t^2) ,2-2u+\frac{1}{2}ut^2\rangle$
se define
$ {\bf r}( t,u,w) = \langle t,4u-ut^2,\frac{1}{2}w( 4-4u+ut^2)\rangle ,$
con $t \in [-2,2], \, u \in [0,1],\, w \in [0,1].$ Además,
{\footnotesize
\begin{eqnarray*} {\bf r}_{t} &=& \langle 1,-2tu,wtu\rangle \\ {\bf r}_{u}&=& \langle 0,4-t^2,w(t^2-4)/2 \rangle \\
{\bf r}_{w}&=& \langle 0,0,2-2u+ut^2/2 \rangle. \end{eqnarray*}
}
%\vskip0.2cm

El triple producto vectorial
$|{\bf r}_t\cdot ( {\bf r}_u \times {\bf r}_w)|=\frac{1}{2}| ( 4-t^2) ( 4-4u+ut^2) | $
se integra para obtener el volumen del sólido
%\vskip0.2cm
{\footnotesize
\[\int_{0}^{1}\int_{0}^{1}\int_{-2}^2\frac{1}{2}\left|(4-t^2)(4-4u+ut^2)\right|dtdudw= \frac{64}{5}\]
}
%\end{tcolorbox}


\subsection*{Un cilindro parabólico y un plano inclinado II}


Hallar el volumen del sólido acotado por
las figuras de las ecuaciones $y=0,$ $y=2-z,$ $z=2-x^2,$ $z=0$. Esta región del espacio se muestra a continuación.
\vskip0.2cm

\noindent
\begin{minipage}{\textwidth}
	\centering
	\captionof{figure}{Visualización tridimensional del sólido limitado por $z=2-x^2$ y $z=2-y$.}
	\label{fig:solido_3D_parabolico}
	\begin{minipage}{0.35\linewidth}
		\centering
		\vspace{-15mm}
		\includegraphics[width=\linewidth]{Fig/378.pdf}
	\end{minipage}
	\hspace{1em}
	\begin{minipage}{0.35\linewidth}
		\centering
		\vspace{-15mm}
		\includegraphics[width=\linewidth]{Fig/379.pdf}
	\end{minipage}
	
	\vspace{-8mm}
	\fuentepropia
\end{minipage}
\vskip0.2cm


Ahora se proyecta el sólido sobre los planos coordenados.
\vskip0.2cm

\noindent
\begin{minipage}{\textwidth}
	\centering
	\captionof{figure}{Proyecciones del sólido en los planos coordenados.}
	\label{fig:proyecciones_sólido}
	\begin{minipage}{0.25\linewidth}
		\centering
		\vspace{-12mm}
		\includegraphics[width=\linewidth]{Fig/380.pdf}
	\end{minipage}
	\hspace{1em}
	\begin{minipage}{0.25\linewidth}
		\centering
		\vspace{-12mm}
		\includegraphics[width=\linewidth]{Fig/381.pdf}
	\end{minipage}
	\hspace{1em}
	\begin{minipage}{0.25\linewidth}
		\centering
		\vspace{-12mm}
		\includegraphics[width=\linewidth]{Fig/382.pdf}
	\end{minipage}
	
	\vspace{-8mm}
	\fuentepropia
\end{minipage}
\vskip0.2cm


Con estos elementos, el volumen en mención está dado por las siguientes integrales triples en los seis órdenes posibles y utilizando coordenadas cartesianas.
\vspace{-1mm}
\begin{eqnarray*}
\int_{-\sqrt{2}}^{\sqrt{2}}\int_{0}^{2-x^2}\int_{0}^{2-z}dydzdx=
\int_{0}^2\int_{-\sqrt{2-z}}^{\sqrt{2-z}}\int_{0}^{2-z}dydxdz&=& \frac{16}{5}\sqrt{2}\\
\int_{0}^2\int_{0}^{2-z}\int_{-\sqrt{2-z}}^{\sqrt{2-z}}dxdydz=
\int_{0}^2\int_{0}^{2-y}\int_{-\sqrt{2-z}}^{\sqrt{2-z}}dxdzdy&=& \frac{16}{5}\sqrt{2}
\end{eqnarray*}

Al integrar primero con respecto a $z$ se debe tener en cuenta que en el plano $xy$ la región de integración está
dividida en dos partes y dependiendo de la forma como se recorra tendremos dos o tres integrales triples.
\[\int_{-\sqrt{2}}^{\sqrt{2}}\int_{0}^{x^2}\int_{0}^{2-x^2}dzdydx + \int_{-\sqrt{2}}^{\sqrt{2}}\int_{x^2}^2\int_{0}^{2-y}dzdydx= \frac{16}{5}\sqrt{2}\]
\begin{multline*} \int_{0}^2\int_{-\sqrt{2}}^{-\sqrt{y}}\int_{0}^{2-x^2}dzdxdy+\int_{0}^2\int_{\sqrt{y}}^{\sqrt{2}}\int_{0}^{2-x^2}dzdxdy \\
+\int_{0}^2\int_{-\sqrt{y}}^{\sqrt{y}}\int_{0}^{2-y}dzdxdy= \frac{16}{5}%
\sqrt{2}.
\end{multline*}
Las siguientes son las indicaciones a usar en \verb"Mathematica" para evaluar estas últimas integrales.

\vspace{0.6em}
\noindent
\begin{tcolorbox}[breakable,enhanced,title=\bf Instrucción en \verb"Mathematica", top=2pt,bottom=2pt,boxsep=0pt,before skip=2pt,after skip=0pt]
\begin{Verbatim}[formatcom=\letraverbatim]
Integrate[1,{x,-Sqrt[2],Sqrt[2]},{y,0,x^2},{z,0,2-x^2}]
+Integrate[1,{x,-Sqrt[2],Sqrt[2]},{y,x^2,2},{z,0,2-y}]
Integrate[1,{y,0,2},{x,-Sqrt[2],-Sqrt[y]},{z,0,2-x^2}]
+Integrate[1,{y,0,2},{x,Sqrt[y],Sqrt[2]},{z,0,2-x^2}]+
Integrate[1,{y,0,2},{x,-Sqrt[y],Sqrt[y]},{z,0,2-y}]
\end{Verbatim}
\end{tcolorbox}
\vspace{0.7em}
\noindent


Con el siguiente procedimiento se muestra este cálculo de otra manera.

%\begin{tcolorbox}[title=\bfseries Parametrizando el sólido]
La proyección en el plano $xz$ se parametriza con $x=t,$
$z=( 2-t^2) u,$ por tanto un punto del sólido en el plano
$xz$ es de la forma $( t,0,2-t^2) $ y un punto sobre el plano $y=2-z$ es de la forma $( t,2-(2-t^2) u,(2-t^2) u).$
Simplificando la expresión $(1-w)\langle t,0,(2-t^2) u\rangle +w\langle t,2-( 2-t^2)u,( 2-t^2) u\rangle$
se parametriza el sólido, esto es
%\vskip0.2cm
\[{\bf r}( t,u,w) =\langle t,w(2-2u+ut^2) ,2u-ut^2\rangle,\]
\vskip0.2cm
con $t\in [ -\sqrt{2},\sqrt{2}],$ $u\in [0,1],$ $w\in
\lbrack 0,1].$ Derivando parcialmente se obtiene
\[{\bf r}_{t}= \langle 1,2wtu,-2tu\rangle,\] \[{\bf r}_{u}= \langle 0,w( -2+t^2),2-t^2\rangle,\] 
\[{\bf r}_{w}= \langle 0,2-2u+ut^2,0\rangle,\]
\[|{\bf r}_{t}\cdot({\bf r}_{u}\times {\bf r}_{w})|=|(2-t^2)(2u-ut^2-2)|.\]
Este último valor se integra y proporciona el volumen del sólido
\[\int_{0}^{1}\int_{0}^{1}\int_{-\sqrt{2}}^{\sqrt{2}}| ( 2-t^2) ( 2u-ut^2-2) |dtdudw= \frac{16}{5}\sqrt{2}.\]
%\end{tcolorbox}

El cálculo anterior se hace en \verb"Mathematica" utilizando las siguientes sentencias:

\vspace{0.6em}
\noindent
\begin{tcolorbox}[breakable,enhanced,title=\bf Instrucción en \verb"Mathematica", top=2pt,bottom=2pt,boxsep=0pt,before skip=2pt,after skip=0pt]
\begin{Verbatim}[formatcom=\letraverbatim]
r=(1-w)*{t,0,u*(2-t^2)}+w*{t,2-u*(2-t^2),u*(2-t^2)};
A=Simplify[Dot[D[r,t],Cross[D[r,u],D[r,w]]]]
Integrate[Abs[A],{t,-Sqrt[2],Sqrt[2]},{u,0,1},
{w,0,1}]
\end{Verbatim}
\end{tcolorbox}
\vspace{0.7em}
\noindent


\subsection*{Tres cilindros parabólicos}

Las expresiones $y=x^2$, $y=8-x^2$ y $z=4-x^2$ representan en el espacio tres cilindros parabólicos. Estas superficies y el plano $xy$ determinan un sólido que a continuación se muestra.



\begin{figure}[H]
\centering
\caption{$0 \leq z \leq 4 - x^2$, \quad $x^2 \leq y \leq 8 - x^2$.}
\label{fig:solido_parabolico}
\vspace{-3mm}
\includegraphics[height=4.5cm]{Img/25a.pdf} \hspace{0.2cm}
\includegraphics[height=4.5cm]{Img/25b.pdf} \hspace{0.2cm}
\includegraphics[height=4.5cm]{Img/25c.pdf}
\fuentepropia
\end{figure}
%\vskip0.2cm

Estas figuras se pueden generar usando las siguientes instrucciones:

\vspace{0.6em}
\noindent
\begin{tcolorbox}[breakable,enhanced,title=\bf Instrucción en \verb"Mathematica", top=2pt,bottom=2pt,boxsep=0pt,before skip=2pt,after skip=0pt]
\begin{Verbatim}[formatcom=\letraverbatim]
ContourPlot3D[{y==x^2,y==8-x^2,z==4-x^2},{x,-2,2},
{y,0,8},{z,0,4},PlotPoints->80]
RegionPlot3D[0<z<4-x^2&&x^2<y<8-x^2,{x,-2,2},{y,0,8},
{z,0,4},PlotPoints->70,Mesh->False]
\end{Verbatim}
\end{tcolorbox}
\vspace{0.7em}
\noindent


Las integrales sobre el sólido se plantean a partir de las siguientes figuras:
\vskip0.2cm

\noindent
\begin{minipage}{\textwidth}
	\centering
	\captionof{figure}{Proyección del sólido en los planos $xy$ y $xz$.}
	\label{fig:proyeccion_xy_xz_2}
	\begin{minipage}{0.36\linewidth}
		\centering
		\vspace{-15mm}
		\includegraphics[width=\linewidth]{Fig/383.pdf}
	\end{minipage}
	\hspace{1em} 
	\begin{minipage}{0.36\linewidth}
		\centering
		\vspace{-15mm}
		\includegraphics[width=\linewidth]{Fig/384.pdf}
	\end{minipage}
	
	\vspace{-10mm}
	\fuentepropia
\end{minipage}
\vskip0.2cm

Esta representación nos ayuda a determinar las integrales triples que determinan su volumen.
Al integrar primera con respecto a $z$, esta recorre los puntos desde el plano $z=0$
hasta el cilindro $z=4-x^2$ y los límites de integración para $x$ e $y$ se obtienen del primero de los gráficos anteriores. Entonces, el volumen del sólido está dado por:
{\small
\[\int_{-2}^2\int_{x^2}^{8-x^2}\int_{0}^{4-x^2}dzdydx= \frac{1024}{15}\]
\[\int_{0}^{4}\int_{-\sqrt{y}}^{\sqrt{y}}\int_{0}^{4-x^2}dzdxdy+
\int_{4}^{8}\int_{-\sqrt{8-y}}^{\sqrt{8-y}}\int_{0}^{4-x^2}dzdxdy= \frac{1024}{15}.\]
}
Si se integra con respecto a $y$ primero, se debe tener en cuenta que esta recorre los puntos desde el cilindro $y=x^2$ hasta el cilindro $y=8-x^2$, que equivalen respectivamente a
$x= \pm \sqrt{y}$ y $x= \pm \sqrt{8-z}$. Los límites de integración para $x$ y $z$ se obtienen del segundo de los gráficos anteriores.
\vspace{-1mm}
{\small
\[\int_{-2}^2\int_{0}^{4-x^2}\int_{x^2}^{8-x^2}dydzdx=\int_{0}^{4}\int_{-\sqrt{4-z}}^{\sqrt{4-z}}\int_{x^2}^{8-x^2}dydxdz= \frac{1024}{15}.\]
}


En el plano $yz$ el sólido se proyecta como un rectángulo. Al sustituir en $z=4-x^2$
las dos expresiones que involucran $x^2$, se sigue $z=4-x^2=4-y,$ por una parte, y $z=4-(8-y)=y-4,$
por la parte restante.


\vskip0.2cm
\noindent
\begin{minipage}{\textwidth}
	\centering
	\captionof{figure}{Proyección en el plano $yz$}
	\label{fig:proyeccion_yz_2}
	\begin{minipage}{0.40\linewidth}
		\centering
		\vspace{-23mm}
		\includegraphics[width=\linewidth]{Fig/385.pdf}
	\end{minipage}
	
	\vspace{-17mm}
	\fuentepropia
\end{minipage}
\vskip0.2cm



En este caso, el volumen del sólido está determinado por:
{\small
\begin{multline*}
\int_{0}^{4}\int_{0}^{4-z}\int_{-\sqrt{y}}^{\sqrt{y}}dxdydz+
\int_{0}^{4}\int_{4-z}^{z+4}\int_{-\sqrt{4-z}}^{\sqrt{4-z}}dxdydz\\
+\int_{0}^{4}\int_{z+4}^{8}\int_{-\sqrt{8-y}}^{\sqrt{8-y}}dxdydz= \frac{1024}{15}
\end{multline*}
}
y en el orden $dxdzdy,$ el planteamiento es el siguiente:
\vspace{-2mm}
{\small
\begin{multline*}
\int_{0}^{4}\int_{0}^{4-y}\int_{-\sqrt{y}}^{\sqrt{y}}dxdzdy+
\int_{0}^{4}\int_{4-y}^{4}\int_{-\sqrt{4-z}}^{\sqrt{4-z}}dxdzdy+ \\
\int_{4}^{8}\int_{y-4}^{4}\int_{-\sqrt{4-z}}^{\sqrt{4-z}}dxdzdy+\int_{4}^{8}%
\int_{0}^{y-4}\int_{-\sqrt{8-y}}^{\sqrt{8-y}}dxdzdy=\frac{1024}{15}.
\end{multline*}
}
\vskip0.2cm
Algunas de estas integrales se evalúan en \verb"Mathematica" mediante la ejecución de las siguientes instrucciones:


\vspace{0.6em}
\noindent
\begin{tcolorbox}[breakable,enhanced,title=\bf Instrucción en \verb"Mathematica", top=2pt,bottom=2pt,boxsep=0pt,before skip=2pt,after skip=0pt]
\begin{Verbatim}[formatcom=\letraverbatim]
Integrate[1,{x,-2,2},{z,0,4-x^2},{y,x^2,8-x^2}]
Integrate[1,{z,0,4},{x,-Sqrt[4-z],Sqrt[4-z]},
{y,x^2,8-x^2}]
Integrate[1,{y,0,4},{x,-Sqrt[y],Sqrt[y]},
{z,0,4-x^2}]+Integrate[1,{y,4,8},{x,-Sqrt[8-y],
Sqrt[8-y]},{z,0,4-x^2}]
Integrate[1,{z,0,4},{y,0,4-z},{x,-Sqrt[y],Sqrt[y]}]+
Integrate[1,{z,0,4},{y,4-z,z+4},{x,-Sqrt[4-z],
Sqrt[4-z]}]+Integrate[1,{z,0,4},{y,z+4,8},
{x,-Sqrt[8-y],Sqrt[8-y]}]
\end{Verbatim}
\end{tcolorbox}
\vspace{0.7em}
\noindent


%Integrate[1,{y,0,4},{z,0,4-y},{x,-Sqrt[y],Sqrt[y]}]+
%Integrate[1,{y,0,4},{z,4-y,4},{x,-Sqrt[4-z],Sqrt[4-z]}]+
%Integrate[1,{y,4,8},{z,y-4,4},{x,-Sqrt[4-z],Sqrt[4-z]}]+
%Integrate[1,{y,4,8},{z,0,y-4},{x,-Sqrt[8-y],Sqrt[8-y]}]

Otra manera de obtener el mismo valor es la descrita a continuación.


%\begin{tcolorbox}[title=\bfseries Parametrizando el sólido]
\noindent

En el plano $xy$, con dos puntos $P(t,t^2)$ y $Q(t,8-t^2)$
sobre las curvas $y=x^2$ y $y=8-x^2$.
La región mostrada se parametriza simplificando la expresión $(1-u)Q+uP,$ es decir,
$(t,8(1-u)+(2u-1)t^2),$ con $t \in [-2,2]$, $u \in [0,1].$


\vskip0.2cm
\noindent
\begin{minipage}{\textwidth}
	\centering
	\captionof{figure}{Región delimitada por $x^2 \leq y \leq 8 - x^2$}
	\label{fig:region_parabolica}
	\begin{minipage}{0.38\linewidth}
		\centering
		\vspace{-18mm}
		\includegraphics[width=\linewidth]{Fig/386.pdf}
	\end{minipage}
	
	\vspace{-10mm}
	\fuentepropia
\end{minipage}
\vskip0.2cm


Sobre el cilindro $z=4-x^2,$ se satisface que $z=4-t^2$
y los puntos con coordenadas
$R(t,8(1-u)+(2u-1)t^2,0)$ y $S(t,8(1-u)+(2u-1)t^2,4-t^2)$
se unen mediante un segmento rectilíneos $(1-w)R+wS,$ que permite definir la función que parametriza el sólido.
\begin{align*}
\mathbf{r}(t,u,w) &= \langle t,8(1-u)+(2u-1)t^2,w(4-t^2)\rangle,
\end{align*}
con $t\in [-2,2],\, u\in [0,1], \, w\in [0,1].$
Esta función nos permite hallar los siguientes vectores
\vspace{-4mm}
\begin{align*}
\mathbf{r}_{t}&=\langle 1,2(2u-1)t,-2wt\rangle \\
\mathbf{r}_{u}&=\langle 0,-8+2t^2,0\rangle\\
\mathbf{r}_{w}&=\langle 0,0,4-t^2\rangle.
\end{align*}
Con estos se halla el producto vectorial triple
$|\mathbf{r}_{t}\cdot(\mathbf{r}_{u}\times \mathbf{r}_{w})| =2(t^2-4)^2,$
que se integra con los límites establecidos y nos proporciona el volumen del sólido
\[\int_{0}^{1}\int_{0}^{1}\int_{-2}^22(t^2-4)^2dtdudw=\frac{1024}{15}.\]

%\end{tcolorbox}
%\vskip0.2cm
Este proceso se hace en \verb"Mathematica" ejecutando las siguientes instrucciones:


\vspace{0.6em}
\noindent
\begin{tcolorbox}[breakable,enhanced,title=\bf Instrucción en \verb"Mathematica", top=2pt,bottom=2pt,boxsep=0pt,before skip=2pt,after skip=0pt]
\begin{Verbatim}[formatcom=\letraverbatim]
R={x,y}=Factor[u*{t,t^2}+(1-u)*{t,8-t^2}]
r=Simplify[(1-w)*{x,y,0}+w*{x,y,4-x^2}]
Integrate[Abs[Dot[D[r,t],Cross[D[r,u],D[r,w]]]],
{t,-2,2},{u,0,1},{w,0,1}]
\end{Verbatim}
\end{tcolorbox}
\vspace{0.7em}
\noindent


\subsection*{Dos cilindros parabólicos}

Consideraremos el sólido acotado por las figuras de las expresiones $z=4x-x^2$, $y=4x-x^2$, $y=0$ y $z=0$.
Un bosquejo de dicha región del espacio se muestra en la siguiente figura.
Las ecuaciones cuadráticas corresponden a dos cilindros parabólicos.
%\vskip0.4cm
\vskip0.2cm
\noindent
\begin{minipage}{\textwidth}
	\centering
	\captionof{figure}{Visualización del sólido definido por $0 \leq z \leq 4x - x^2$, \quad $0 \leq y \leq 4x - x^2$.}
	\label{fig:solido_4x_x2}
	\vspace{-10pt}
	\begin{minipage}{0.33\linewidth}
		\centering
		\vspace{-13mm}
		\includegraphics[width=\linewidth]{Fig/387.pdf}
	\end{minipage}
	\hspace{1em}
	\begin{minipage}{0.24\linewidth}
		\centering
		\vspace{-13mm}
		\includegraphics[width=\linewidth]{Fig/388.pdf}
	\end{minipage}
	\hspace{1em}
	\begin{minipage}{0.24\linewidth}
		\centering
		\vspace{-13mm}
		\includegraphics[width=\linewidth]{Img/7b.pdf}
	\end{minipage}
	
	\vspace{-10mm}
	\fuentepropia
\end{minipage}
\vskip0.4cm

La línea roja corresponde a la curva de intersección de los dos cilindros parabólicos; en esta curva se usó la parametrización
$\overrightarrow{r}(t)= \langle t,4t-t^2,4t-t^2\rangle,$ $t\in [0,4].$
Esta línea se proyecta en el plano $yz$ por medio de la recta $z=y.$
Estas figuras se consiguen con las siguientes instrucciones:


\vspace{0.6em}
\noindent
\begin{tcolorbox}[breakable,enhanced,title=\bf Instrucción en \verb"Mathematica", top=2pt,bottom=2pt,boxsep=0pt,before skip=2pt,after skip=0pt]
\begin{Verbatim}[formatcom=\letraverbatim]
ContourPlot3D[{z==4*x-x^2,y==4*x-x^2},{x,0,4},
{y,0,4},{z,0,4}]
RegionPlot3D[0<z<4*x-x^2&&0<y<4*x-x^2,{x,0,4},
{y,0,4},
{z,0,4},PlotStyle->{Blend[{Red,Orange,Gray}]},
Mesh->False,PlotPoints->50]
\end{Verbatim}
\end{tcolorbox}
\vspace{0.7em}
\noindent



La proyección del sólido en el plano $yz$ corresponde a un cuadrado, mientras en los planos $xy$ y $xz$ coincide con sectores parabólicos.

\vskip0.2cm
\noindent
\begin{minipage}{\textwidth}
	\centering
	\captionof{figure}{Proyección del sólido.}
	\label{fig:proyeccion_solido}
	\begin{minipage}{0.30\linewidth}
		\centering
		\vspace{-10mm}
		\includegraphics[width=\linewidth]{Fig/389.pdf}
	\end{minipage}
	\hspace{1em}
	\begin{minipage}{0.30\linewidth}
		\centering
		\vspace{-10mm}
		\includegraphics[width=\linewidth]{Fig/390.pdf}
	\end{minipage}
	\hspace{1em}
	\begin{minipage}{0.25\linewidth}
		\centering
		\vspace{-10mm}
		\includegraphics[width=\linewidth]{Fig/391.pdf}
	\end{minipage}
	
	\vspace{-8mm}
	\fuentepropia
\end{minipage}
\vskip0.2cm


Con ayuda de los elementos anteriores se presentan las integrales triples necesarias para calcular el volumen del sólido en mención. \\[0.2cm]
$\ds \int_{0}^{4}\int_{0}^{z}\int_{2-\sqrt{4-z}}^{2+\sqrt{4-z}}dxdydz+\int_{0}^{4}\int_{z}^{4}\int_{2-\sqrt{4-y}}^{2+\sqrt{4-y}}dxdydz= \frac{512}{15}$\\[0.1cm]
$\ds \int_{0}^{4}\int_{y}^{4}\int_{2-\sqrt{4-z}}^{2+\sqrt{4-z}}dxdzdy+\int_{0}^{4}\int_{0}^{y}\int_{2-\sqrt{4-y}}^{2+\sqrt{4-y}}dxdzdy= \frac{512}{15}$\\[0.1cm]
$\ds \int_{0}^{4}\int_{0}^{4x-x^2}\int_{0}^{4x-x^2}dzdydx = \int_{0}^{4}\int_{2-\sqrt{4-y}}^{2+\sqrt{4-y}}\int_{0}^{4x-x^2}dzdxdy= \frac{512}{15}$\\[0.1cm]
$\ds \int_{0}^{4}\int_{0}^{4x-x^2}\int_{0}^{4x-x^2}dydzdx= \int_{0}^{4}\int_{2-\sqrt{4-z}}^{2+\sqrt{4-z}}\int_{0}^{4x-x^2}dydxdz= \frac{512}{15}.$\\[0.15cm]
\vskip0.2cm
Algunas de estas integrales se evalúan en \verb"Mathematica" con las siguientes instrucciones:

\vspace{0.6em}
\noindent
\begin{tcolorbox}[breakable,enhanced,title=\bf Instrucción en \verb"Mathematica", top=2pt,bottom=2pt,boxsep=0pt,before skip=2pt,after skip=0pt]
\begin{Verbatim}[formatcom=\letraverbatim]
Integrate[1,{x,0,4},{y,0,4*x-x^2},{z,0,4*x-x^2}]
Integrate[1,{x,0,4},{z,0,4*x-x^2},{y,0,4*x-x^2}]
Integrate[1,{z,0,4},{x,2-Sqrt[4-z],2+Sqrt[4-z]},
{y,0,4*x-x^2}]
Integrate[1,{z,0,4},{y,0,z},{x,2-Sqrt[4-z],2+
Sqrt[4-z]}]+Integrate[1,{z,0,4},{y,z,4},{x,2-
Sqrt[4-y],2+Sqrt[4-y]}]
Integrate[1,{y,0,4},{z,y,4},{x,2-Sqrt[4-z],
2+Sqrt[4-z]}]+Integrate[1,{y,0,4},{z,0,y},
{x,2-Sqrt[4-y],2+Sqrt[4-y]}]
\end{Verbatim}
\end{tcolorbox}
\vspace{0.7em}
\noindent


Otra manera de calcular el volumen del sólido es la siguiente.


%\begin{tcolorbox}[title=\bfseries Parametrizando el sólido]
\noindent

En la proyección sobre el plano $xy$,
consideremos un punto $P(t,0)$ y un punto $Q(t,4t-t^2)$,
con $t \in [0,4]$. El segmento que los une corresponde a la expresión
$\left( 1-v\right) \left( t,0\right) +v\left( t,4t-t^2\right),$
que se simplifica como $(t,vt(4-t)),$ $v \in[0,1].$


\vskip0.2cm
\noindent
\begin{minipage}{\textwidth}
	\centering
	\captionof{figure}{Región delimitada por $0 \leq y \leq 4x - x^2$}
	\label{fig:region_4x_x2}
	\begin{minipage}{0.37\linewidth}
		\centering
		\vspace{-18mm}
		\includegraphics[width=\linewidth]{Fig/392.pdf}
	\end{minipage}
	
	\vspace{-14mm}
	\fuentepropia
\end{minipage}
\vskip0.2cm


Un punto en la base del sólido tiene la forma
$\left( t,vt(4-t),0\right).$ El cilindro con ecuación $y=4x-x^2$ es paralelo al eje $z,$
sobre esta región, un punto tiene la forma $\left( t,vt(4-t),4t-t^2\right)$.
Formamos segmentos rectilíneos que unen dichos puntos, usaremos el parámetro $w.$ Por lo tanto, al simplificar
$\left( 1-w\right) \left( t,vt(4-t),0\right) +w\left( t,vt(4-t),4t-t^2\right)$
obtendremos la función
\[{\bf r}(u,v,w)=\left\langle t,vt(4-t),-wt\left(t -4\right) \right\rangle,\]
que parametriza el sólido con $t\in [0,4],\,v\in [0,1],\,w\in [0,1]. $
A partir de esta se pueden obtener 
{\small
\begin{gather*}
	\mathbf{r}_{t}= \left\langle 1,4v-2vt,4w-2wt\right\rangle \\
	\mathbf{r}_{v}= \left\langle 0,4t-t^2,0\right\rangle \\
	\mathbf{r}_{w}= \left\langle 0,0,t\left( 4-t\right) \right\rangle \\
	\mathbf{r}_{v}\times \mathbf{r}_{w}= \left\langle t^2\left( t-4\right) ^2,0,0\right\rangle.
\end{gather*}
}
%\vskip0.2cm
Con esto se obtiene
$\left| {\bf r}_{t}\cdot \left( {\bf r}_{v}\times {\bf r}_{w}\right) \right| = t^2(t-4) ^2,$
el volumen está dado por:
%\vskip0.2cm
{\small
\[\int_{0}^{4}\int_{0}^{1}\int_{0}^{1}t^2\left( t-4\right)^2dvdwdt= \frac{512}{15}.\]
}
%\end{tcolorbox}
\vskip0.2cm
El cálculo anterior se hizo en \verb"Mathematica" con las siguientes instrucciones:


\vspace{0.6em}
\noindent
\begin{tcolorbox}[breakable,enhanced,title=\bf Instrucción en \verb"Mathematica", top=2pt,bottom=2pt,boxsep=0pt,before skip=2pt,after skip=0pt]
\begin{Verbatim}[formatcom=\letraverbatim]
{x,y}=Simplify[u*{t,t^2}+(1-u)*{t,8-t^2}]
r=Simplify[(1-w)*{x,y,0}+w*{x,y,4-x^2}]
A=Abs[Dot[D[r,w],Cross[D[r,t],D[r,u]]]];
Integrate[A,{t,-2,2},{u,0,1},{w,0,1}]
\end{Verbatim}
\end{tcolorbox}
\vspace{0.7em}
\noindent


\subsection*{Un cilindro parabólico y un plano inclinado III}

Las expresiones $z=2-y$, $y=x^2$ y $z=0$,
determinan una región en el espacio, su gráfico se muestra a continuación:
\vskip0.2cm

\noindent
\begin{minipage}{\textwidth}
	\centering
	\captionof{figure}{Visualización del sólido definido por $0 \leq z \leq 2 - y$, \quad $x^2 \leq y \leq 2$.}
	\label{fig:solido_parabolico_2y}
	\begin{minipage}{0.42\linewidth}
		\centering
		\vspace{-24mm}
		\includegraphics[width=\linewidth]{Fig/393.pdf}
	\end{minipage}
	\hspace{1em}
	\begin{minipage}{0.42\linewidth}
		\centering
		\vspace{-24mm}
		\includegraphics[width=\linewidth]{Fig/394.pdf}
	\end{minipage}
	
	\vspace{-15mm}
	\fuentepropia
\end{minipage}
\vskip0.2cm


Sobre los planos coordenados, el sólido se proyecta como sigue:
\vskip0.2cm

\noindent
\begin{minipage}{\textwidth}
	\centering
	\captionof{figure}{Proyecciones sobre los planos coordenados.}
	\label{fig:proyecciones_coordenadas_final}
	\begin{minipage}{0.30\linewidth}
		\centering
		\vspace{-10mm}
		\includegraphics[width=\linewidth]{Fig/395.pdf}
	\end{minipage}
	\hspace{1em}
	\begin{minipage}{0.30\linewidth}
		\centering
		\vspace{-10mm}
		\includegraphics[width=\linewidth]{Fig/396.pdf}
	\end{minipage}
	\hspace{1em}
	\begin{minipage}{0.30\linewidth}
		\centering
		\vspace{-10mm}
		\includegraphics[width=\linewidth]{Fig/397.pdf}
	\end{minipage}
	
	\vspace{-8mm}
	\fuentepropia
\end{minipage}
\vskip0.2cm


Estos gráficos nos ayudan a plantear las integrales necesarias para calcular el volumen del sólido.

\begin{multline*} \int_{-\sqrt{2}}^{\sqrt{2}}\int_{x^2}^2\int_{0}^{2-y}dzdydx=
\int_{0}^2\int_{-\sqrt{y}}^{\sqrt{y}}\int_{0}^{2-y}dzdxdy \\
= \int_{-\sqrt{2}}^{\sqrt{2}}\int_{0}^{2-x^2}\int_{x^2}^{2-z}dydzdx = \frac{32\sqrt{2}}{15} \end{multline*}

\begin{multline*} \int_{0}^2\int_{-\sqrt{2-z}}^{\sqrt{2-z}}\int_{x^2}^{2-z}dydxdz =
\int_{0}^2\int_{0}^{2-z}\int_{-\sqrt{y}}^{\sqrt{y}}dxdydz \\
= \int_{0}^2\int_{0}^{2-y}\int_{-\sqrt{y}}^{\sqrt{y}}dxdzdy = \frac{32\sqrt{2}}{15}. \end{multline*}
Se pueden evaluar unas de estas integrales usando \verb"Mathematica" con ayuda de las siguientes sentencias.


\vspace{0.6em}
\noindent
\begin{tcolorbox}[breakable,enhanced,title=\bf Instrucción en \verb"Mathematica", top=2pt,bottom=2pt,boxsep=0pt,before skip=2pt,after skip=0pt]
\begin{Verbatim}[formatcom=\letraverbatim]
Integrate[1,{x,-Sqrt[2],Sqrt[2]},{y,x^2,2},{z,0,2-y}]
Integrate[1,{z,0,2},{y,0,2-z},{x,-Sqrt[y],Sqrt[y]}]
Integrate[1,{y,0,2},{z,0,2-y},{x,-Sqrt[y],Sqrt[y]}]
Integrate[1,{x,-Sqrt[2],Sqrt[2]},{y,x^2,2},{z,0,2-y}]
\end{Verbatim}
\end{tcolorbox}
\vspace{0.7em}
\noindent


El siguiente procedimiento es una alternativa al cálculo del volumen.

\vskip0.2cm
%\begin{tcolorbox}[title=\bfseries Parametrizando el sólido]
En el plano $xy,$, la base del sólido se parametriza formando una combinación convexa
$\left( 1-u\right) \left( t,t^2\right) +u\left(t,2\right) = \left( t,2u-t^2u+t^2\right),$
entre dos puntos de la forma $\left( t,t^2\right) $ y $\left( t,2\right).$
Sobre el plano $z=2-y$ se satisface que $z=(t^2-2)(u-1) .$ Por tanto, el sólido se parametriza con una combinación convexa de la forma
\[\left( 1-w\right) \left\langle t,2u-t^2u+t^2,0\right\rangle +w\left\langle t,2u-t^2u+t^2,\left(t^2-2\right) \left( u-1\right) \right\rangle,\]
que al simplificarse nos permiten definir la función
\[{\bf r}(t,u,w)=\left\langle t,2u-t^2u+t^2,w\left( t^2-2\right) \left( u-1\right) \right\rangle,\]
con $t\in [-\sqrt{2},\sqrt{2}],$ $u\in [0,1],$ $w\in [0,1].$ 
\vskip0.2cm
Derivando parcialmente se obtienen los siguientes vectores
\begin{gather*}
	\mathbf{r}_{t} = \left\langle 1,2t-2tu,2tw\left( u-1\right) \right\rangle, \\
	\mathbf{r}_{u}= \left\langle 0,2-t^2,w\left( t^2-2\right)\right\rangle, \\
	\mathbf{r}_{w} =\left\langle 0,0,\left( t^2-2\right) \left( u-1\right)\right\rangle.
\end{gather*}

De modo que $\left\vert {\bf r}_{w}\cdot \left({\bf r}_{t}\times {\bf r}_{u}\right) \right\vert
=\left\vert \left( t^2-2\right) ^2\left( 1-u\right)\right\vert .$
Por tanto, el volumen del sólido está dado por
%\vskip0.2cm
\[\int_{0}^{1}\int_{0}^{1}\int_{-\sqrt{2}}^{\sqrt{2}}\left\vert
\left( t^2-2\right) ^2\left( 1-u\right) \right\vert dtdudw= \frac{32}{15}\sqrt{2}.\]
%\vskip0.2cm
%\end{tcolorbox}

El procedimiento anterior se hace en \verb"Mathematica" ejecutando las siguientes instrucciones:


\vspace{0.6em}
\noindent
\begin{tcolorbox}[breakable,enhanced,title=\bf Instrucción en \verb"Mathematica", top=2pt,bottom=2pt,boxsep=0pt,before skip=2pt,after skip=0pt]
\begin{Verbatim}[formatcom=\letraverbatim]
R={x,y}=(1-u)*{t,t^2}+u{t,2};
r=Simplify[(1-w)*{x,y,0}+w*{x,y,2-y}]
Integrate[Dot[D[r,w],Cross[D[r,t],D[r,u]]],
{t,-Sqrt[2],Sqrt[2]},{u,0,1},{w,0,1}]
\end{Verbatim}
\end{tcolorbox}
\vspace{0.7em}
\noindent



\subsection*{Un cilindro parabólico y un plano oblicuo}

A continuación, se muestra el sólido acotado por las figuras de las expresiones:
$z=2x-x^2, \,\, y=2-x,\,\, x=1,\,\, z=0, \,\, y=0 $.
\vskip0.4cm


\noindent
\begin{minipage}{\textwidth}
	\centering
	\captionof{figure}{Visualización del sólido definido por $0 \leq z \leq 2x - x^2$, \quad $1 \leq x \leq 2 - y$.}
	\label{fig:solido_2x_x2}
	\begin{minipage}{0.37\linewidth}
		\centering
		\vspace{-10mm}
		\includegraphics[width=\linewidth]{Fig/398.pdf}
	\end{minipage}
	\hspace{1em}
	\begin{minipage}{0.37\linewidth}
		\centering
		\vspace{-10mm}
		\includegraphics[width=\linewidth]{Fig/399.pdf}
	\end{minipage}
	
	\vspace{-10mm}
	\fuentepropia
\end{minipage}
\vskip0.2cm

Este sólido proyectado determina las siguientes regiones.

\vskip0.2cm
\noindent
\begin{minipage}{\textwidth}
	\centering
	\captionof{figure}{Proyecciones en cada uno de los planos coordenados.}
	\label{fig:proyecciones_planes}
	\begin{minipage}{0.28\linewidth}
		\centering
		\vspace{-9mm}
		\includegraphics[width=\linewidth]{Fig/400.pdf}
	\end{minipage}
	\hspace{1em}
	\begin{minipage}{0.28\linewidth}
		\centering
		\vspace{-9mm}
		\includegraphics[width=\linewidth]{Fig/401.pdf}
	\end{minipage}
	\hspace{1em}
	\begin{minipage}{0.28\linewidth}
		\centering
		\vspace{-9mm}
		\includegraphics[width=\linewidth]{Fig/402.pdf}
	\end{minipage}
	
	\vspace{-4mm}
	\fuentepropia
\end{minipage}
\vskip0.2cm

El volumen del sólido está dado por
\[\int_{1}^2\int_{0}^{2-x}\int_{0}^{2x-x^2}dzdydx = \int_{0}^{1}\int_{1}^{2-y}\int_{0}^{2x-x^2}dzdxdy= \frac{5}{12}\]
\[\int_{1}^2\int_{0}^{2x-x^2}\int_{0}^{2-x}dydzdx=
\int_{0}^{1}\int_{1}^{1+\sqrt{1-z}}\int_{0}^{2-x}dydxdz= \frac{5}{12}.\]

La curva de intersección proyectada en el plano $yz,$ satisface que
$z=2y-y^2,$ con lo cual $y=1\pm \sqrt{1-z}.$ La figura de esta ecuación divide en dos partes  la proyección del sólido que es un cuadrado, esto nos lleva a plantear dos integrales en el orden $dxdzdy$ o $dxdydz$.
\begin{eqnarray*}
\int_{0}^{1}\int_{0}^{2y-y^2}\int_{1}^{2-y}dxdzdy+\int_{0}^{1}%
\int_{2y-y^2}^{1}\int_{1}^{1+\sqrt{1-z}}dxdzdy &=& \frac{5}{12} \\
\int_{0}^{1}\int_{1-\sqrt{1-z}}^{1}\int_{1}^{2-y}dxdydz+\int_{0}^{1} \int_{0}^{1-\sqrt{1-z}}\int_{1}^{1+\sqrt{1-z}}dxdydz &=& \frac{5}{12}.
\end{eqnarray*}
En \verb"Mathematica" se pueden evaluar algunas de estas integrales ejecutando las siguientes instrucciones.


\vspace{0.6em}
\noindent
\begin{tcolorbox}[breakable,enhanced,title=\bf Instrucción en \verb"Mathematica", top=2pt,bottom=2pt,boxsep=0pt,before skip=2pt,after skip=0pt]
\begin{Verbatim}[formatcom=\letraverbatim]
Integrate[1,{y,0,1},{x,1,2-y},{z,0,2x-x^2}]
Integrate[1,{z,0,1},{x,1,1+Sqrt[1-z]},{y,0,2-x}]
Integrate[1,{x,1,2},{y,0,2-x},{z,0,2*x-x^2}]
Integrate[1,{x,1,2},{z,0,2*x-x^2},{y,0,2-x}]
\end{Verbatim}
\end{tcolorbox}
\vspace{0.7em}
\noindent


La otra posibilidad de obtener el volumen en mención es la siguiente.
%\begin{tcolorbox}[title=\bfseries Parametrizando el sólido]
La base del sólido es un triángulo en el plano $xy$ que se parametriza primero formando segmentos que unen los puntos $(1,0,0) $ .
$\left(2,0,0\right),$ posteriormente se forman otros segmentos que unan con $\left( 1,1,0\right) ,$
esto es $(1-u)(1+t,0,0)+u(1,1,0)= ((1-u)(1+t)+u,u,0).$
Con lo anterior, la coordenada $z$ sobre el cilindro satisface que $z=w (t+1-ut)(tu-t+1).$
Es por lo anterior, que el sólido se parametriza como sigue
%\vskip0.2cm
\[{\bf r}(t,u,w) =\left\langle t+1-ut,u,w(t+1-tu)(tu-t+1)\right\rangle,\]
con $t\in [0,1],$ $u\in [0,1],$ $w\in [0,1].$ Con esta función se determinan los siguientes vectores:
\vspace{-2mm}
\begin{eqnarray*}
{\bf r}_{t}&=& \left\langle 1-u,0,-2tw(u-1)^2\right\rangle \\
{\bf r}_{u}&=&\left\langle -t,1,2wt^2(1-u) \right\rangle \\
{\bf r}_{w}&=& \left\langle 0,0,1-t^2(u-1)^2\right\rangle
\end{eqnarray*}

El producto vectorial triple está dado por

\[
\left| {\bf r}_{t}\cdot \left( {\bf r}_{u}\times {\bf r}_{w}\right) \right| =
\left| (u-1)\, (t (1-u)+1)\, (t (u-1)+1) \right|,
\]

cuya integración nos proporciona el volumen del sólido:

\[
\int_{0}^{1} \int_{0}^{1} \int_{0}^{1}
\left| (u-1)\, (t (1-u)+1)\, (t (u-1)+1) \right|
\, dt\, du\, dw = \frac{5}{12}.
\]
\vskip0.2cm
%\end{tcolorbox}

Con las siguientes sentencias se hace el procedimiento descrito anteriormente.


\vspace{0.6em}
\noindent
\begin{tcolorbox}[breakable,enhanced,title=\bf Instrucción en \verb"Mathematica", top=2pt,bottom=2pt,boxsep=0pt,before skip=2pt,after skip=0pt]
\begin{Verbatim}[formatcom=\letraverbatim]
{x,y}=(1-u)*{1+t,0}+u*{1,1};
r=Simplify[{x,y,w*(2*x-x^2)}]
A=Simplify[Dot[D[r,w],Cross[D[r,u],D[r,t]]]]
Integrate[Abs[A],{t,0,1},{u,0,1},{w,0,1}]
\end{Verbatim}
\end{tcolorbox}
\vspace{0.7em}
\noindent


\subsection*{Dos cilindros parabólicos II}
La región del espacio acotada por las figuras de las expresiones
$x=z^2$ y $x=4y-y^2$, determinan el sólido que a continuación se presenta.
\vskip0.3cm

\noindent
\begin{minipage}{\textwidth}
	\centering
	\captionof{figure}{$z^2 \leq x \leq 4y - y^2$, \quad $y^2 + z^2 \leq 4y$.}
	\label{fig:solido_curvo}
	\begin{minipage}{0.20\linewidth}
		\centering
		\includegraphics[width=\linewidth]{Imgn/53a.pdf}
	\end{minipage}
	\hspace{1em}
	\begin{minipage}{0.21\linewidth}
		\centering
		\includegraphics[width=\linewidth]{Img/53b.pdf}
	\end{minipage}
	\hspace{1em}
	\begin{minipage}{0.21\linewidth}
		\centering
		\includegraphics[width=\linewidth]{Img/53c.pdf}
	\end{minipage}
	\hspace{1em}
	\begin{minipage}{0.21\linewidth}
		\centering
		\includegraphics[width=\linewidth]{Img/53d.pdf}
	\end{minipage}
	
	\vspace{-2mm}
	\fuentepropia
\end{minipage}
\vskip0.3cm

Estos gráficos se obtienen con la ayuda de las siguientes expresiones.


\vspace{0.6em}
\noindent
\begin{tcolorbox}[breakable,enhanced,title=\bf Instrucción en \verb"Mathematica", top=2pt,bottom=2pt,boxsep=0pt,before skip=2pt,after skip=0pt]
\begin{Verbatim}[formatcom=\letraverbatim]
ContourPlot3D[{x==z^2,x==4*y-y^2},{x,0,4},{y,0,4},
{z,-2,2},PlotPoints->50,Mesh->False,Ticks->{{0,2,4},
{0,1,3},{0,2}},ContourStyle->{Blend[{Yellow,Pink,
Green}],Orange}]
RegionPlot3D[-Sqrt[x]<z<Sqrt[x]&&0<x<4*y-y^2&&0<y<4,
{x,0,4},{y,0,4},{z,-2,2},PlotPoints->80,Mesh->False,
PlotStyle->Directive[Orange,Specularity[White,40]]]
\end{Verbatim}
\end{tcolorbox}
\vspace{0.7em}
\noindent


Al igualar estas expresiones para determinar su proyección en el plano $yz,$ obtenemos $(y-2)^2+z^2=4,$
que corresponde a una circunferencia de radio $2$ centrada en el punto $(2,0).$
\vskip 0.2cm
En el plano $xz$ la proyección es la porción de parábola $x=z^2$ para $x\in[0,4];$ mientras que para el plano $xy$ la región de integración es la porción de parábola $x=4y-y^2$ para $x\in[0,4].$ En los siguientes gráficos se ilustran estas regiones.
\vskip0.4cm

\noindent
\begin{minipage}{\textwidth}
	\centering
	\captionof{figure}{Proyecciones en los planos coordenados.}
	\label{fig:proyecciones_curvas}
	\begin{minipage}{0.29\linewidth}
		\centering
		\vspace{-10mm}
		\includegraphics[width=\linewidth]{Fig/403.pdf}
	\end{minipage}
	\hspace{1em}
	\begin{minipage}{0.29\linewidth}
		\centering
		\vspace{-10mm}
		\includegraphics[width=\linewidth]{Fig/404.pdf}
	\end{minipage}
	\hspace{1em}
	\begin{minipage}{0.29\linewidth}
		\centering
		\vspace{-10mm}
		\includegraphics[width=\linewidth]{Fig/405.pdf}
	\end{minipage}
	
	\vspace{-8mm}
	\fuentepropia
\end{minipage}
\vskip0.2cm

Con ayuda de los gráficos anteriores se pueden plantear las integrales que determinan el volumen del sólido; en un caso conviene escribir la expresión $x=4y-y^2$ en la forma
$y=2\pm \sqrt{4-x}.$
{\small
\begin{eqnarray*}
\int_{0}^{4}\int_{-\sqrt{4y-y^2}}^{\sqrt{4y-y^2}}\int_{z^2}^{4y-y^2}dxdzdy= \int_{-2}^2\int_{2-\sqrt{4-z^2}}^{2+\sqrt{4-z^2}}\int_{z^2}^{4y-y^2}dxdydz &=& 8\pi \\
\int_{-2}^2\int_{z^2}^{4}\int_{2-\sqrt{4-x}}^{2+\sqrt{4-x}}dydxdz=\int_{0}^{4}\int_{-\sqrt{x}}^{\sqrt{x}}\int_{2-\sqrt{4-x}}^{2+\sqrt{4-x}}dydzdx &=& 8\pi \\
\int_{0}^{4}\int_{0}^{4y-y^2}\int_{-\sqrt{x}}^{\sqrt{x}}dzdxdy= \int_{0}^{4}\int_{2-\sqrt{4-x}}^{2+\sqrt{4-x}}\int_{-\sqrt{x}}^{\sqrt{x}}dzdydx &= & 8\pi
\end{eqnarray*}
}
Algunas de estas integrales se evalúan como sigue:


\vspace{0.6em}
\noindent
\begin{tcolorbox}[breakable,enhanced,title=\bf Instrucción en \verb"Mathematica", top=2pt,bottom=2pt,boxsep=0pt,before skip=2pt,after skip=0pt]
\begin{Verbatim}[formatcom=\letraverbatim]
Integrate[1,{y,0,4},{z,-Sqrt[4y-y^2],Sqrt[4y-y^2]},
{x,z^2,4y-y^2}]
Integrate[1,{z,-2,2},{y,2-Sqrt[4-z^2],2+Sqrt[4-z^2]},
{x,z^2,4y-y^2}]
Integrate[1,{y,0,4},{x,0,4y-y^2},{y,-Sqrt[x],Sqrt[x]}]
Integrate[1,{z,-2,2},{x,z^2,4},{y,2-Sqrt[4-x],
2+Sqrt[4-x]}]
\end{Verbatim}
\end{tcolorbox}
\vspace{0.7em}
\noindent

El mismo valor se puede obtener de la siguiente manera.

%\begin{tcolorbox}[title=\bfseries Parametrizando el sólido]
En el plano $xz$ dos puntos sobre la región de integración poseen coordenadas
$( t,-\sqrt{t})$ y $(t,\sqrt{t})$; con una combinación convexa entre ellos se describe completamente la región.
Es decir, simplificando
%\vskip0.2cm
\[( 1-u) \langle t,-\sqrt{t}\rangle +u\langle t,\sqrt{t}\rangle;\]
%\vskip0.2cm
para $t\in [0,4]$ y $u\in [0,1]$,
se obtiene $\langle t,\sqrt{t}(2u-1)\rangle\!.$ Por lo tanto, un punto sobre la región de integración posee coordenadas $( t,\sqrt{t}( 2u-1)),$
mientras que dos puntos $P$ y $Q$ sobre el cilindro poseen coordenadas
$( t,2-\sqrt{4-t},\sqrt{t}( 2u-1) ) $ y $( t,2+\sqrt{4-t},\sqrt{t}( 2u-1) ) .$ % 
Una combinación convexa de estos puntos parametriza el sólido, esto es
%\vskip0.2cm
{\small
\[{\bf r}(t,u,w) = \langle t,2+( 1-2w) \sqrt{4-t},( 2u-1) \sqrt{t}\rangle.\]
}
%\vskip0.2cm
Con esta función hallamos los siguientes vectores:
{\small
\begin{eqnarray*}
{\bf r}_t & =& \langle 1,-\frac{1}{2}\frac{1-2w}{\sqrt{4-t}},\frac{2u-1}{2\sqrt{t}}\rangle \\
{\bf r}_u & =& \langle 0,0,2\sqrt{t}\rangle \\
{\bf r}_w & =& \langle 0,-2\sqrt{4-t},0\rangle.
\end{eqnarray*}
}
Es sencillo ver que $|{\bf r}_t \cdot ({\bf r}_u \times {\bf r}_w )| = 4\sqrt{t}\sqrt{4-t}.$
Integrando este término obtenemos el volumen deseado.
%\vskip0.2cm
\[\int_{0}^{1}\int_{0}^{1}\int_{0}^{4}( 4\sqrt{t}\sqrt{4-t}) dtdudw =8\pi.\]
%\vskip0.2cm
%\end{tcolorbox}

Ejecutando las siguientes instrucciones se obtiene el volumen del sólido.


\vspace{0.6em}
\noindent
\begin{tcolorbox}[breakable,enhanced,title=\bf Instrucción en \verb"Mathematica", top=2pt,bottom=2pt,boxsep=0pt,before skip=2pt,after skip=0pt]
\begin{Verbatim}[formatcom=\letraverbatim]
r=Simplify[(1-w)*{t,2+Sqrt[4-t],Sqrt[t]*(2*u-1)}+
w*{t,2-Sqrt[4-t],Sqrt[t]*(2*u-1)}]
Integrate[Dot[D[r,t],Cross[D[r,u],D[r,w]]],{w,0,1},
{u,0,1},{t,0,4}]
\end{Verbatim}
\end{tcolorbox}
\vspace{0.7em}
\noindent

\subsection*{Dos cilindros parabólicos y dos planos}

Las expresiones $z=y^2$, $x=z^2/4$, $x=0$ y $y=2,$ determinan un sólido cuya figura se muestra a continuación.
\vskip0.2cm

\noindent
\begin{minipage}{\textwidth}
	\centering
	\captionof{figure}{Visualización del sólido definido por $0 \leq 4x \leq z^2$, \quad $0 \leq z \leq y^2$.}
	\label{fig:solido_xz_y2}
	\begin{minipage}{0.32\linewidth}
		\centering
		\vspace{-15mm}
		\includegraphics[width=\linewidth]{Fig/406.pdf}
	\end{minipage}
	\hspace{1em}
	\begin{minipage}{0.28\linewidth}
		\centering
		\vspace{-15mm}
		\includegraphics[width=\linewidth]{Fig/407.pdf}
	\end{minipage}
	
	\vspace{-8mm}
	\fuentepropia
\end{minipage}

\vskip0.2cm
En el plano $yz$ la proyección se determina por las ecuaciones $y=2$, $z=y^2,$ $z=0,$
y en el plano $xz$ la proyección la determinan las expresiones $4x=z^2,\,\, z=4$ y $x=0$.
Al sustituir las ecuaciones
$4x=z^2$ y $z=y^2$, en términos de $x$ e $y$ se puede representar la proyección en el plano
$xy.$ Es decir, $4x=z^2=(y^2)^2$ luego $4x=y^4,$ con lo anterior se obtiene las siguientes figuras.


\vskip0.2cm
\noindent
\begin{minipage}{\textwidth}
	\centering
	\captionof{figure}{Proyección sobre los planos coordenados.}
	\label{fig:proyecciones_planes_curvas}
	\begin{minipage}{0.25\linewidth}
		\centering
		\vspace{-10mm}
		\includegraphics[width=\linewidth]{Fig/408.pdf}
	\end{minipage}
	\hspace{1em}
	\begin{minipage}{0.25\linewidth}
		\centering
		\vspace{-10mm}
		\includegraphics[width=\linewidth]{Fig/409.pdf}
	\end{minipage}
	\hspace{1em}
	\begin{minipage}{0.25\linewidth}
		\centering
		\vspace{-10mm}
		\includegraphics[width=\linewidth]{Fig/410.pdf}
	\end{minipage}
	
	\vspace{-4mm}
	\fuentepropia
\end{minipage}
\vskip0.2cm

Con ayuda de estos gráficos se pueden plantear las seis integrales triples que permiten calcular el volumen del sólido en mención.
\vspace{-2mm}
\begin{multline*}
\int_{0}^2\int_{0}^{y^{4}/4}\int_{2\sqrt{x}}^{y^2}dzdxdy=\int_{0}^{4}\int_{\sqrt[4]{4x}}^2\int_{2\sqrt{x}}^{y^2}dzdydx \\
= \int_{0}^{4}\int_{0}^{z^2/4}\int_{\sqrt{z}}^2dydxdz=
\int_{0}^{4}\int_{2\sqrt{x}}^{4}\int_{\sqrt{z}}^2dydzdx \\= \int_{0}^{4}\int_{\sqrt{z}}^2\int_{0}^{z^2/4}dxdydz= \int_{0}^2\int_{0}^{y^2}\int_{0}^{z^2/4}dxdzdy= \frac{32}{21}.
\end{multline*}
Se presentan las instrucciones para evaluar cuatro de estas integrales.


\vspace{0.6em}
\noindent
\begin{tcolorbox}[breakable,enhanced,title=\bf Instrucción en \verb"Mathematica", top=2pt,bottom=2pt,boxsep=0pt,before skip=2pt,after skip=0pt]
\begin{Verbatim}[formatcom=\letraverbatim]
Integrate[1,{y,0,2},{z,0,y^2},{x,0,z^2/4}]
Integrate[1,{z,0,4},{y,Sqrt[z],2},{x,0,z^2/4}]
Integrate[1,{x,0,4},{z,Sqrt[4*x],4},{y,Sqrt[z],2}]
Integrate[1,{y,0,2},{x,0,y^4/4},{z,2*Sqrt[x],y^2}]
\end{Verbatim}
\end{tcolorbox}
\vspace{0.7em}
\noindent


Los siguientes gráficos ofrecen otras vistas del mismo sólido.

\begin{figure}[H]
\centering
\caption{Visualización del sólido definido por $0 \leq 4x \leq z^2$, \quad $0 \leq z \leq y^2$.}
\label{fig:solido_80_xyz}
\includegraphics[height=4.5cm]{Img/80c.pdf}
\includegraphics[height=4.5cm]{Img/80b.pdf}
\includegraphics[height=4.5cm]{Img/80a.pdf}
\fuentepropia
\end{figure}

Estos se consiguen ejecutando las siguientes instrucciones.


\vspace{0.6em}
\noindent
\begin{tcolorbox}[breakable,enhanced,title=\bf Instrucción en \verb"Mathematica", top=2pt,bottom=2pt,boxsep=0pt,before skip=2pt,after skip=0pt]
\begin{Verbatim}[formatcom=\letraverbatim]
RegionPlot3D[0<x<z^2/4&&Sqrt[z]<y<2,{x,0,4},{y,0,2},
{z,0,4},PlotPoints->90,PlotStyle->Orange,Boxed->False,
MeshStyle->Gray,AxesStyle->False,Axes->None]
\end{Verbatim}
\end{tcolorbox}
\vspace{0.7em}
\noindent


Una alternativa para obtener el volumen del sólido descrito es la siguiente.



%\begin{tcolorbox}[title=\bfseries Parametrizando el sólido]
\noindent

Consideremos dos puntos $P$ y $Q$ con coordenadas $(t,t^2)$ y $(2,t^2)$ respectivamente, en la proyección del sólido en el plano $yz,$
con $t \in [0,2]$. Simplificando una combinación convexa de estos se llena la región, es decir \linebreak
$(1-u)(t,t^2)+u(2,t^2)=( t+u(2-t),t^2) $
$t \in [0,2],\, u \in [0,1].$


\vskip0.3cm
\noindent
\begin{minipage}{\textwidth}
	\centering
	\captionof{figure}{Proyección de la parábola $z = y^2$ y puntos $P$, $Q$ sobre la recta $z = 2$}
	\label{fig:proyeccion_parabola_pq}
	\begin{minipage}{0.32\linewidth}
		\centering
		\vspace{-15mm}
		\includegraphics[width=\linewidth]{Fig/411.pdf}
	\end{minipage}
	
	\vspace{-8mm}
	\fuentepropia
\end{minipage}
\vskip0.2cm


Ahora, el punto $\left( 0,t+u\left( 2-t\right) ,t^2\right) $ está
en el plano $x=0,$ mientras que el punto
$\left( \frac{1}{4}t^{4},t+u\left( 2-t\right) ,t^2\right) $ yace en el cilindro $x=\frac{1}{4}z^2.$ Uniendo estos dos puntos con un segmento rectilíneo se parametriza el sólido, es decir simplificando la expresión
\vspace{-2mm}
\[\left( 1-w\right) \left (t^{4}/4,t+u(2-t),t^2\right) +w\left( 0,t+u(2-t) ,t^2\right),\] 
\[{\bf r}\left( t,u,w\right) = \left\langle \frac{1}{4}\left(1-w\right) t^{4},t+2u-tu,t^2\right\rangle,\]
en donde $t\in [0,2],$ $u\in [0,1],$ $w\in [0,1].$
Derivando parcialmente, se hallan los siguientes vectores

\vspace{-4mm}
{\small
\begin{align*}
	\mathbf{r}_t &= \frac{\partial}{\partial t} \left\langle \frac{1}{4}(1 - w) t^4, t + 2u - tu, t^2 \right\rangle = \left\langle (1 - w) t^3, 1 - u, 2t \right\rangle \\
	\mathbf{r}_u &= \frac{\partial}{\partial u} \left\langle \frac{1}{4}(1 - w) t^4, t + 2u - tu, t^2 \right\rangle = \left\langle 0, 2 - t, 0 \right\rangle \\
	\mathbf{r}_w &= \frac{\partial}{\partial w} \left\langle \frac{1}{4}(1 - w) t^4, t + 2u - tu, t^2 \right\rangle = \left\langle -\frac{1}{4}t^4, 0, 0 \right\rangle.
\end{align*}
}

\vskip0.2cm
Con ellos se calcula el triple producto vectorial
{\small
\[\left| {\bf r}_{t}\cdot \left( {\bf r}_{u}\times {\bf r}_{w}\right) \right| =
\frac{1}{2}\left| t^{5}\left( t-2\right) \right|,\]
}
que corresponde al elemento diferencial de volumen. Integrándolo, nos proporciona el volumen del sólido
%\vskip0.2cm
{\small
\[\int_{0}^{1}\int_{0}^{1}\int_0^2\frac{1}{2}\left|t^{5}(t-2)\right|dtdudw= \frac{32}{21}.\]
}
%\vskip0.2cm
%\end{tcolorbox}

Este valor se obtiene en \verb"Mathematica" ejecutando las siguientes instrucciones.
\vskip0.2cm
\begin{tcolorbox}[breakable,enhanced,title=\bf Instrucción en \verb"Mathematica", top=2pt,bottom=2pt,boxsep=0pt,before skip=2pt,after skip=0pt]
\begin{Verbatim}[formatcom=\letraverbatim]
r=(1-w)*{t^4/4,t+u*(2-t),t^2}+w*{0,t+u*(2-t),t^2}
Integrate[Dot[D[r,w],Cross[D[r,t],D[r,u]]],{t,0,2},
{u,0,1},{w,0,1}]
\end{Verbatim}
\end{tcolorbox}

\subsection*{Un cilindro parabólico y dos planos II}

Las ecuaciones $z=y^2$, $x=-1$ y $2z+x=3$ acotan un sólido en el espacio. Enseguida se representan estas expresiones junto con la región acotada.
\vskip0.4cm

\noindent
\begin{minipage}{\textwidth}
	\centering
	\captionof{figure}{Visualización del sólido definido por $-1 \leq x \leq 3 - 2z$, \quad $-\sqrt{z} \leq y \leq \sqrt{z}$.}
	\label{fig:solido_88_xyz}
	\begin{minipage}{0.22\linewidth}
		\centering
		\includegraphics[width=\linewidth]{Img/88a.pdf}
	\end{minipage}
	\hspace{1em} 
	\begin{minipage}{0.22\linewidth}
		\centering
		\includegraphics[width=\linewidth]{Img/88b.pdf}
	\end{minipage}
	\hspace{1em} 
	\begin{minipage}{0.22\linewidth}
		\centering
		\includegraphics[width=\linewidth]{Img/88d.pdf}
	\end{minipage}
	
	\vspace{0.3em}
	\fuentepropia
\end{minipage}
\vskip0.3cm

las figuras se obtienen con las siguientes instrucciones:

\vspace{0.6em}
\noindent
\begin{tcolorbox}[breakable,enhanced,title=\bf Instrucción en \verb"Mathematica", top=2pt,bottom=2pt,boxsep=0pt,before skip=2pt,after skip=0pt]
\begin{Verbatim}[formatcom=\letraverbatim]
ContourPlot3D[{z==y^2,x==-1,z==(3-x)/2},{x,-1,3},
{y,-1.5,1.5},{z,0,2},ContourStyle->{Orange,Cyan,
RGBColor[0.2,0.4,0.4]},Mesh->False,BoxRatios->{2,2,1}]
RegionPlot3D[y^2<=z<(3-x)/2&&-1<x<3,{x,-1,3},
{y,-1.5,1.5},{z,0,2},PlotPoints->90,Mesh->False,
BoxRatios->{2,2,1},PlotStyle->RGBColor[0.65,0.1,0.2]]
\end{Verbatim}
\end{tcolorbox}
\vspace{0.7em}
\noindent


Este sólido se proyecta sobre los planos coordenados.
\vskip0.2cm

\noindent
\begin{minipage}{\textwidth}
	\centering
	\captionof{figure}{Proyecciones del sólido definido por $-1 \leq x \leq 1 - 3z$, \quad $y^2 \leq z \leq 2$, \quad $-1 \leq x \leq 3 - 2y^2$.}
	\label{fig:proyecciones_sistema_curvo}
	\begin{minipage}{0.27\linewidth}
		\centering
		\vspace{-10mm}
		\includegraphics[width=\linewidth]{Fig/412.pdf}
	\end{minipage}
	\hspace{1em}
	\begin{minipage}{0.27\linewidth}
		\centering
		\vspace{-10mm}
		\includegraphics[width=\linewidth]{Fig/413.pdf}
	\end{minipage}
	\hspace{1em}
	\begin{minipage}{0.27\linewidth}
		\centering
		\vspace{-10mm}
		\includegraphics[width=\linewidth]{Fig/414.pdf}
	\end{minipage}
	
	\vspace{-4mm}
	\fuentepropia
\end{minipage}
\vskip0.2cm

En coordenadas rectangulares, los límites de las integrales que determinan el volumen se establecen con ayuda de estas figuras,
\begin{multline*}
\bigintsss_{0}^2\bigintsss_{-1}^{3-2z}\bigintsss_{-\sqrt{z}}^{\sqrt{z}}dydxdz = \bigintsss_{-1}^{3} \bigintsss_{0}^{\frac{3-x}{2}} \bigintsss_{-\sqrt{z}}^{\sqrt{z}}dydzdx \\
= \bigintsss_{-\sqrt{2}}^{\sqrt{2}}\bigintsss_{-1}^{3-2y^2}\bigintsss_{y^2}^{\frac{3-x}{2}}dzdxdy= \bigintsss_{-1}^{3}\bigintsss_{-\sqrt{\frac{3-x}{2}}}^{\sqrt{\frac{3-x}{2}}}\bigintsss_{y^2}^{\frac{3-x}{2}}dzdydx \\
= \bigintsss_{-\sqrt{2}}^{\sqrt{2}}\bigintsss_{y^2}^2 \bigintsss_{-1}^{3-2z}dxdzdy=\bigintsss_{0}^2\bigintsss_{-\sqrt{z}}^{\sqrt{z}}\bigintsss_{-1}^{3-2z}dxdydz= \frac{64}{15}\sqrt{2}
\end{multline*}

Cinco de estas integrales se evalúan en \verb"Mathematica" ejecutando las siguien\-tes instrucciones:


\vspace{0.6em}
\noindent
\begin{tcolorbox}[breakable,enhanced,title=\bf Instrucción en \verb"Mathematica", top=2pt,bottom=2pt,boxsep=0pt,before skip=2pt,after skip=0pt]
\begin{Verbatim}[formatcom=\letraverbatim]
Integrate[1,{y,-Sqrt[2],Sqrt[2]},{z,y^2,2},{x,-1,3-2z}]
Integrate[1,{y,-Sqrt[2],Sqrt[2]},{x,-1,3-2y^2},
{z,y^2,(3-x)/2}]
Integrate[z^2,{z,0,2},{x,-1,3-2z},{y,-Sqrt[z],Sqrt[z]}]
Integrate[1,{z,0,2},{y,-Sqrt[z],Sqrt[z]},{x,-1,3-2z}]
Integrate[1,{x,-1,3},{y,-Sqrt[(3-x)/2],Sqrt[(3-x)/2]},
{z,y^2,(3-x)/2}]
\end{Verbatim}
\end{tcolorbox}
\vspace{0.7em}
\noindent


El siguiente proceso ofrece otra posibilidad de obtener el mismo valor.

%\begin{tcolorbox}[title=\bfseries Parametrizando el sólido]
%\footnotesize 
%\parskip=0pt 
%\setlength{\belowdisplayskip}{0pt} 
%\setlength{\abovedisplayskip}{0pt} 

En el plano $xy,$ un punto $M$ sobre la curva $x=3-2y^2$ tiene la forma $(3-2t^2,t)$ y sobre la recta $x=-1$, un punto $N$
tiene la forma $(-1,t)$. Por tanto, la proyección del sólido se parametriza mediante una combinación convexa de $M$ y $N$,
esto es $(1-u)N+uM=\left( 2u\left( 2-t^2\right) -1,t\right) ,$ $t \in [-\sqrt{2},\sqrt{2}],$ $u \in [0,1].$
\vskip0.4cm

\begin{center}
	% 1. Ampliamos la minipage al ancho total para que el texto NO se comprima
	\begin{minipage}{1.0\linewidth} 
		\centering
		% La leyenda ahora tiene espacio para extenderse horizontalmente
		\captionof{figure}{Región delimitada por $-1 \leq x \leq 3 - 2y^2$}
		\label{fig:region_parabolica_x} 
		
		\vspace{-20mm} % Espacio estético entre leyenda y figura
		
		% 2. Mantenemos la figura pequeña (ajusta este valor a tu gusto)
		\includegraphics[width=0.45\linewidth]{Fig/415.pdf}
		
		\vspace{-18mm}
		% Fuente centrada respecto a la página
		\parbox{\linewidth}{\centering\fontsize{9pt}{10pt}\selectfont\textbf{Fuente:} elaboración propia.}
	\end{minipage}
\end{center}
\vskip0.4cm
Supongamos un punto $P\left( 2u\left( 2-t^2\right) -1,t,t^2\right) $ sobre el cilindro $z=y^2,$
y un punto $Q$ sobre el plano $2z+x=3,$ cuya coordenada $z$ satisface
%\vskip0.2cm
\[
z = \frac{3 - x}{2} = \frac{3}{2} - \frac{1}{2} \left(2u \left(2 - t^2\right) - 1\right) = u\left(t^2 - 2\right) + 2.
\]
%\vskip0.2cm
Por tanto, $Q$ es de la forma $Q( 2u\left( 2-t^2\right) -1,t,u\left( t^2-2\right)+2).$ Formando una combinación convexa entre $P$ y $Q$, se parametriza el sólido, esto es,
%\vskip0.2cm
\begin{equation*}
{\bf r}\left( t,u,w\right) = \langle 2u\left( 2-t^2\right)-1,t,t^2\left( 1-w+uw\right) -2w\left( u-1\right) \rangle.
\end{equation*}
%\vskip0.2cm
Con esta función hallamos los vectores
\begin{gather*}
	\mathbf{r}_{t} = \left\langle -4tu,1,2t\left( uw-w+1\right) \right\rangle, \\
	\mathbf{r}_{u} = \left\langle 4-2t^2,0,t^2w-2w\right\rangle, \\
	\mathbf{r}_{w} = \left\langle 0,0,\left( u-1\right) t^2-2u+2\right\rangle.
\end{gather*}
\vskip0.2cm
El elemento diferencial de volumen está dado por
\[\left\vert {\bf r}_{t}\cdot \left( {\bf r}_{u}\times {\bf r}_{w}\right) \right\vert
=\left\vert \left( 2t^2-4\right) \left( \left( u-1\right)t^2-2u+2\right) \right\vert,\]
cuya integral determina el volumen del sólido
%\vskip0.2cm
\[\int_{-\sqrt{2}}^{\sqrt{2}}\int_{0}^{1}\int_{0}^{1}\left\vert \left(2t^2-4\right) \left( \left( u-1\right) t^2-2u+2\right) \right\vert dwdudt= \frac{64}{15}\sqrt{2}.\]
%\vskip0.2cm
%\end{tcolorbox}

Con las siguientes instrucciones se hacen los cálculos anteriores.
%\vskip0.2cm

\vspace{0.6em}
\noindent
\begin{tcolorbox}[breakable,enhanced,title=\bf Instrucción en \verb"Mathematica", top=2pt,bottom=2pt,boxsep=0pt,before skip=2pt,after skip=0pt]
\begin{Verbatim}[formatcom=\letraverbatim]
r=Simplify[(1-w)*{2u*(2-t^2)-1,t,t^2}+w*{2u*(2-t^2)-1,
t,u*(t^2-2)+2}]
A=Dot[D[r,w],Cross[D[r,t],D[r,u]]]
Integrate[Abs[A],{w,0,1},{u,0,1},{t,-Sqrt[2],Sqrt[2]}]
\end{Verbatim}
\end{tcolorbox}
\vspace{0.7em}
\noindent


\subsection*{Un cilindro parabólico y dos planos}

Consideremos el sólido $E$, acotado por las figuras de las expresiones
$x=1-z^2$, $z=1-y$, $y=0$. Estas superficies y el sólido se muestran enseguida.

\vskip0.4cm
\noindent
\begin{minipage}{\textwidth}
	\centering
	\captionof{figure}{Visualización del sólido definido por $0 \leq x \leq 1 - z^2$, \quad $0 \leq y \leq 1 - z$.}
	\label{fig:solido_30_xyz}
	\begin{minipage}{0.24\linewidth}
		\centering
		\vspace{-15mm}
		\includegraphics[width=\linewidth]{Img/30a.pdf}
	\end{minipage}
	\hspace{1em}
	\begin{minipage}{0.40\linewidth}
		\centering
		\vspace{-15mm}
		\includegraphics[width=\linewidth]{Fig/416.pdf}
	\end{minipage}
	\hspace{1em}
	\begin{minipage}{0.24\linewidth}
		\centering
		\vspace{-15mm}
		\includegraphics[width=\linewidth]{Img/30d.pdf}
	\end{minipage}
	
	\vspace{-12mm}
	\parbox{\linewidth}{\centering\fontsize{9pt}{10pt}\selectfont\textbf{Fuente:} elaboración propia.}
\end{minipage}

\vskip0.4cm
Con las siguientes instrucciones se obtienen las figuras anteriores:


\vspace{0.6em}
\noindent
\begin{tcolorbox}[breakable,enhanced,title=\bf Instrucción en \verb"Mathematica", top=2pt,bottom=2pt,boxsep=0pt,before skip=2pt,after skip=0pt]
\begin{Verbatim}[formatcom=\letraverbatim]
ContourPlot3D[{x==1-z^2,z==1-y},{x,0,1},{y,0,2},
{z,-1,1},Mesh->False,ContourStyle->{Orange,Gray,
Opacity[0.85]},PlotPoints->80]
RegionPlot3D[0<x<1-z^2&&0<y<1-z,{x,0,1},{y,0,2},
{z,-1,1},PlotStyle->Blend[{Yellow,Orange,Gray,Red}],
Mesh->False,PlotPoints->90]
\end{Verbatim}
\end{tcolorbox}
\vspace{0.7em}
\noindent


Con el propósito de plantear las integrales triples que se requieren para calcular el volumen de $E$, se muestran ahora las proyecciones en los planos coordenados.
Para el caso del plano $xy$, se elimina la variable $z$ de las expresiones que determinan la intersección, es decir, se hace la sustitución $x=1-z^2=1-(1-y)^2=2y-y^2.$

\vskip0.3cm
\noindent
\begin{minipage}{\textwidth}
	\centering
	\captionof{figure}{Proyecciones del sólido definido por $0 \leq x \leq 1 - z^2$, \quad $0 \leq y \leq 1 - z$, \quad $-1 \leq z \leq 1$.}
	\label{fig:proyecciones_1menosz}
	\begin{minipage}{0.28\linewidth}
		\centering
		\vspace{-10mm}
		\includegraphics[width=\linewidth]{Fig/417.pdf}
	\end{minipage}
	\hspace{1em}
	\begin{minipage}{0.28\linewidth}
		\centering
		\vspace{-10mm}
		\includegraphics[width=\linewidth]{Fig/418.pdf}
	\end{minipage}
	\hspace{1em}
	\begin{minipage}{0.28\linewidth}
		\centering
		\vspace{-10mm}
		\includegraphics[width=\linewidth]{Fig/419.pdf}
	\end{minipage}
	
	\vspace{-6mm}
	\fuentepropia
\end{minipage}
\vskip0.3cm


Ahora se presentan las integrales que determinan el volumen del sólido $E$:
\begin{eqnarray*}
\ds \int_{-1}^{1}\int_{0}^{1-z}\int_{0}^{1-z^2}dxdydz=\ds\int_{0}^2\int_{-1}^{1-y}\int_{0}^{1-z^2}dxdzdy&=& \frac{4}{3}\\
\ds \int_{-1}^{1}\int_{0}^{1-z^2}\int_{0}^{1-z}dydxdz=\ds\int_{0}^{1}\int_{-\sqrt{1-x}}^{\sqrt{1-x}}\int_{0}^{1-z}dydzdx&=& \frac{4}{3}\\
\ds \int_{0}^{1}\int_{1-\sqrt{1-x}}^{1+\sqrt{1-x}}\int_{-\sqrt{1-x}}^{1-y}dzdydx+\int_{0}^{1}\int_{0}^{1-\sqrt{1-x}}\int_{-\sqrt{1-x}}^{\sqrt{1-x}}dzdydx&=& \frac{4}{3}\\
\ds \int_{0}^2\int_{0}^{2y-y^2}\int_{-\sqrt{1-x}}^{1-y}dzdxdy+\int_{0}^{1} \int_{2y-y^2}^{1}\int_{-\sqrt{1-x}}^{\sqrt{1-x}}dzdxdy&=& \frac{4}{3}
\end{eqnarray*}
Las integrales anteriores se evaluaron en \verb"Mathematica" mediante la ejecución de las siguientes instrucciones:


\vspace{0.6em}
\noindent
\begin{tcolorbox}[breakable,enhanced,title=\bf Instrucción en \verb"Mathematica", top=2pt,bottom=2pt,boxsep=0pt,before skip=2pt,after skip=0pt]
\begin{Verbatim}[formatcom=\letraverbatim]
Integrate[1,{z,-1,1},{y,0,1-z},{x,0,1-z^2}]
Integrate[1,{y,0,2},{z,-1,1-y},{x,0,1-z^2}]
Integrate[1,{z,-1,1},{x,0,1-z^2},{y,0,1-z}]
Integrate[1,{x,0,1},{z,-Sqrt[1-x],Sqrt[1-x]},{y,0,1-z}]
Integrate[1,{x,0,1},{y,1-Sqrt[1-x],1+Sqrt[1-x]},
{z,-Sqrt[1-x],1-y}]+Integrate[1,{x,0,1},
{y,0,1-Sqrt[1-x]},{z,-Sqrt[1-x],Sqrt[1-x]}]
Integrate[1,{y,0,2},{x,0,2*y-y^2},{z,-Sqrt[1-x],1-y}]+
Integrate[1,{y,0,1},{x,2*y-y^2,1},{z,-Sqrt[1-x],
Sqrt[1-x]}]
\end{Verbatim}
\end{tcolorbox}
\vspace{0.7em}
\noindent


Si agregamos la condición de que la densidad de $E$ en cualquier punto $(x,y,z)$ es
dada por $\delta(x,y,z)=x^2y$, entonces su masa estará dada por
\begin{eqnarray*}
\ds \int_{-1}^{1}\int_{0}^{1-z}\int_{0}^{1-z^2}x^2ydxdydz=\ds\int_{0}^2\int_{-1}^{1-y}\int_{0}^{1-z^2}x^2ydxdzdy&=&\frac{32}{189}
%\\ \ds \int_{-1}^{1}\int_{0}^{1-z^2}\int_{0}^{1-z}x^2ydydxdz= \ds \int_{0}^{1}\int_{-\sqrt{1-x}}^{\sqrt{1-x}}\int_{0}^{1-z}x^2ydydzdx&=& \frac{32}{189}\\ \ds \int_{0}^{1}\int_{1-\sqrt{1-x}}^{1+\sqrt{1-x}}\int_{-\sqrt{1-x}}^{1-y}x^2ydzdydx+\int_{0}^{1}\int_{0}^{1-\sqrt{1-x}}\int_{-\sqrt{1-x}}^{\sqrt{1-x}}x^2ydzdydx &=& \frac{32}{189} \\ \ds \int_{0}^2\int_{0}^{2y-y^2}\int_{-\sqrt{1-x}}^{1-y}x^2ydzdxdy+ \int_{0}^{1}\int_{2y-y^2}^{1}\int_{-\sqrt{1-x}}^{\sqrt{1-x}}x^2ydzdxdy&=& \frac{32}{189}
\end{eqnarray*}
El cálculo del volumen y la masa también se obtienen como sigue.




%\begin{tcolorbox}[title=\bfseries Parametrizando el sólido]
\noindent

En el plano $xz$, sobre la curva $x=1-z^2$ consideramos dos puntos {\small $P(u,\sqrt{1-u})$} y {\small $Q(u,-\sqrt{1-u});$}
Formando segmentos rectilíneos de la forma {\small $\left(1-v\right)P +vQ$,} con {\small $v\in [0,1];$} así, se simplifica como {\small $\left(u,\sqrt{1-u}\left(2v-1\right)\right),$} para {\small $u \in [0,1], \,\, v \in [0,1].$}


\vskip0.3cm
\noindent
\begin{minipage}{\textwidth}
	\centering
	\captionof{figure}{Región delimitada por $0 \leq x \leq 1 - z^2$}
	\label{fig:region_parabolica_xz}
	\begin{minipage}{0.38\linewidth}
		\centering
		\vspace{-18mm}
		\includegraphics[width=\linewidth]{Fig/420.pdf}
	\end{minipage}
	
	\vspace{-10mm}
	\fuentepropia
\end{minipage}
\vskip0.2cm


Formando segmentos rectilíneos desde la base del sólido al plano con ecuación $z=1-y$;
esto es, una combinación convexa de la forma
$(1-w)(u,0,\sqrt{1-u}( 2v-1))+w(u,1-\sqrt{1-u}\left( 2v-1\right) ,\sqrt{1-u}\left( 2v-1\right)),$
que al simplificarse nos permite definir la función:

%se lleva a la forma
%$ \left( u,w\left( 1+\sqrt{1-u}\right) -2wv\sqrt{1-u},\left( 2v-1\right) \sqrt{1-u}\right) , $
\vspace{-6mm}
\[{\bf r}(u,v,w)= \left\langle u,w\left( 1+\sqrt{1-u}\right) -2wv\sqrt{1-u},\left( 2v-1\right) \sqrt{1-u}\right\rangle,\]
%\vskip0.2cm
la cual parametriza el sólido, para $u\in [0,1],\,v\in [0,1],\,w\in [0,1].$
Con lo anterior se obtiene:
\begin{gather*}
	\mathbf{r}_{u}=\left\langle 1,\frac{w}{2}\frac{2v-1}{\sqrt{1-u}},\frac{1-2v}{2\sqrt{1-u}} \right\rangle, \\
	\mathbf{r}_{v} = \left\langle 0,-2w\sqrt{1-u},2\sqrt{1-u}\right\rangle, \\
	\mathbf{r}_{w} = \left\langle 0,1+\left( 1-2v\right) \sqrt{1-u},0\right\rangle.
\end{gather*}

%\vskip0.2cm
Además, $ \left| {\bf r}_{u}\cdot \left( {\bf r}_{v}\times {\bf r}_{w}\right) \right|
= \left| -2\sqrt{1-u}+2\left( 2v-1\right) \left( 1-u\right) \right| .$
\vskip0.3cm
El volumen de $E$ está dado por
\[\ds
\int_{0}^{1}\int_{0}^{1}\int_{0}^{1}\left| -2\sqrt{1-u}+2\left( 2v-1\right)
\left( 1-u\right) \right| dudvdw= \frac{4}{3}.\]
%\end{tcolorbox}
%\vskip0.2cm

Las siguientes instrucciones permiten evaluar los cálculos anteriores.

%\begin{tcolorbox}
Si la densidad del sólido está dada por {\footnotesize $\delta \left( x,y,z\right) =x^2y,$} entonces,
{\footnotesize $\delta \left( u,v,w\right) =u^2w\left( \left( 1+\sqrt{1-u}\right) -2v\sqrt{1-u}\right);$}
por lo tanto, la masa del sólido esta dada por
{\footnotesize $\ds \int_{0}^{1}\int_{0}^{1}\int_{0}^{1} 2 u^2 w\sqrt{1-u} \left((1-2v)\sqrt{1-u}+1\right)^2 dudvdw= \frac{32}{189}.$}
%\end{tcolorbox}
\vskip0.4cm
Las siguientes instrucciones permiten evaluar los cálculos anteriores.


\vspace{0.6em}
\noindent
\begin{tcolorbox}[breakable,enhanced,title=\bf Instrucción en \verb"Mathematica", top=2pt,bottom=2pt,boxsep=0pt,before skip=2pt,after skip=0pt]
\begin{Verbatim}[formatcom=\letraverbatim]
{x,z}=(1-v)*{u,Sqrt[1-u]}+v*{u,-Sqrt[1-u]};
r={x,y,z}=Simplify[(1-w)*{x,0,z}+w*{x,1-z,z}]
A=Abs[Dot[D[r,w],Cross[D[r,u],D[r,v]]]];
Integrate[A,{u,0,1},{v,0,1},{w,0,1}]
Integrate[x^2*y*A,{u,0,1},{v,0,1},{w,0,1}]
\end{Verbatim}
\end{tcolorbox}
\vspace{0.7em}
\noindent


\section{Sólidos acotados por cilindros} \index{Sólidos!con cilindros}
\vskip0.4cm
\subsection*{Dos cilindros y dos planos} \label{2cilindros2planos}
Las figuras de las ecuaciones $z=10-x-y^2/8$, $z=0$, $x^3=y$ y $x=0,$ acotan el sólido que se muestra a continuación.
\vskip0.2cm
\noindent
\begin{minipage}{\textwidth}
	\centering
	\captionof{figure}{Visualización del sólido definido por $0 \leq z \leq 10 - x - \frac{y^2}{8}$ \quad y \quad $x^3 \leq y \leq \sqrt{8(10 - x)}$.}
	\label{fig:solido_11_xyz}
	\begin{minipage}{0.20\linewidth}
		\centering
		\includegraphics[width=\linewidth]{Img/11d.pdf}
	\end{minipage}
	\hspace{1em} 
	\begin{minipage}{0.20\linewidth}
		\centering
		\includegraphics[width=\linewidth]{Img/11c.pdf}
	\end{minipage}
	\hspace{1em} 
	\begin{minipage}{0.20\linewidth}
		\centering
		\includegraphics[width=\linewidth]{Img/11b.pdf}
	\end{minipage}
	
	\vspace{2mm}
	\fuentepropia
\end{minipage}
\vskip0.2cm

Las siguientes son instrucciones que se requieren para obtener la figura:

\vspace{0.6em}
\noindent
\begin{tcolorbox}[breakable,enhanced,title=\bf Instrucción en \verb"Mathematica", top=2pt,bottom=2pt,boxsep=0pt,before skip=2pt,after skip=0pt]
\begin{Verbatim}[formatcom=\letraverbatim]
ContourPlot3D[{z==10-x-y^2/8,y==x^3,z==0,x==0},
{x,0,2},{y,0,9},{z,0,10},ContourStyle->{Orange,Cyan,
LightGray,Yellow},Mesh->False]
RegionPlot3D[0<=z<=10-x-y^2/8&&x^3<=y<=Sqrt[8*(10-x)]
&&0<=x<=2,{x,0,2},{y,0,9},{z,0,10},PlotPoints->90,
Mesh->False]
\end{Verbatim}
\end{tcolorbox}
\vspace{0.7em}
\noindent


Para hallar la curva de intersección de las superficies, se debe observar que al reemplazar
$y=x^{3},$ en la expresión $z=8-x-y^2/8,$ se obtiene $z=8-x-x^{6}/8$ o $z=8-\sqrt[3]{y}-y^2/8.$

\vskip0.2cm
\noindent
\begin{minipage}{\textwidth}
	\centering
	\captionof{figure}{Proyecciones del sólido en los planos coordenados.}
	\label{figura115}
	\begin{minipage}{0.27\linewidth}
		\centering
		\vspace{-10mm}
		\includegraphics[width=\linewidth]{Fig/421.pdf}
	\end{minipage}
	\hspace{1em}
	\begin{minipage}{0.27\linewidth}
		\centering
		\vspace{-10mm}
		\includegraphics[width=\linewidth]{Fig/422.pdf}
	\end{minipage}
	\hspace{1em}
	\begin{minipage}{0.27\linewidth}
		\centering
		\vspace{-10mm}
		\includegraphics[width=\linewidth]{Fig/423.pdf}
	\end{minipage}
	
	\vspace{-4mm}
	\fuentepropia
\end{minipage}
\vskip0.2cm

En el plano $yz$, la proyección del sólido se divide en dos regiones y en el plano $xy$ la región de integración se determina por las ecuaciones $y=x^3,$ $x+y^2/8=10$ y $x=0;$ el volumen del sólido está dado por
\vspace{-10mm}

{\small
\begin{multline*}
\int_0^2\int_{x^3}^{\sqrt{8 (10-x)}}\int_0^{10-x-\frac{y^2}{8}}dzdydx=
\bigintsss _0^8\bigintsss _0^{\sqrt[3]{y}}\bigintsss _0^{10-x-\frac{y^2}{8}}dzdxdy \\
+\bigintsss_8^{\sqrt{80}}\bigintsss_0^{10-\frac{y^2}{8}}\bigintsss_0^{10-x-\frac{y^2}{8}}dzdxdy=\frac{320 \sqrt{5}}{3}-\frac{2488}{15}
\end{multline*}
}

{\footnotesize
\[\bigintsss_0^2\bigintsss_0^{10-x-\frac{x^6}{8}}\bigintsss_{x^3}^{\sqrt{8(10-x-z)}}dydzdx=\frac{320\sqrt{5}}{3}-\frac{2488}{15}.\]
}

Al integrar primero con respecto a $x$ se debe observar que en una región esta recorre los puntos desde el plano $yz$ hasta el cilindro $x=\sqrt[3]{y},$ y en otro caso hasta la superficie {\small $x=10-z-y^2/8.$}
\vspace{-3mm}
{\footnotesize
\begin{multline*}
\bigintsss _0^8\bigintsss _0^{10-\sqrt[3]{y}-\frac{y^2}{8}}\bigintsss _0^{\sqrt[3]{y}}dxdzdy+\bigintsss _0^8\bigintsss_{10-\sqrt[3]{y}-\frac{y^2}{8}}^{10-\frac{y^2}{8}}\bigintsss _0^{10-z-\frac{y^2}{8}}dxdzdy\\ +\bigintsss_8^{\sqrt{80}}\bigintsss
_0^{10-\frac{y^2}{8}}\bigintsss _0^{10-z-\frac{y^2}{8}}dxdzdy = \frac{8}{15} \left(200 \sqrt{5}-311\right).
\end{multline*}
}

%\vskip0.2cm
\vspace{-2mm}
Estos cálculos se evalúan en \verb"Mathematica" con las siguientes instrucciones.


%\vspace{0.6em}

\noindent
\begin{tcolorbox}[breakable,enhanced,title=\bf Instrucción en \verb"Mathematica", top=2pt,bottom=2pt,boxsep=0pt,before skip=2pt,after skip=0pt]
\begin{Verbatim}[formatcom=\letraverbatim]
Integrate[1,{x,0,2},{y,x^3,Sqrt[8*(10-x)]},
{z,0,10-x-y^2/8}]
Integrate[1,{y,0,8},{x,0,y^(1/3)},
{z,0,10-x-y^2/8}]+Integrate[1,{y,8,Sqrt[80]},
{x,0,10-y^2/8},{z,0,10-x-y^2/8}]
Integrate[1,{x,0,2},{z,0,10-x-x^6/8},{y,x^3,
Sqrt[8*(10-x-z)]}]
Integrate[1,{y,0,8},{z,0,10-y^(1/3)-y^2/8},
{x,0,y^(1/3)}]+Integrate[1,{y,0,8},{z,10-y^(1/3)
y^2/8,10-y^2/8},{x,0,10-z-y^2/8}]+Integrate[1,
{y,8,Sqrt[80]},{z,0,10-y^2/8},{x,0,10-z-y^2/8}]
\end{Verbatim}
\end{tcolorbox}
\vspace{0.7em}
\noindent


\subsection*{Cálculo del centro de masa}

Si la densidad del sólido en cualquier punto {\small $(x,y,z)$} está dada por {\small $\delta(x,y,z)=xy$,} entonces la masa del sólido es inversa.

{\small
\[M=\int _0^2\int {x^3}^{\sqrt{8 (10-x)}}\int 0^{10-x-\frac{y^2}{8}}xydzdydx=\frac{13000}{63}.\]
}

Los momentos con respecto a los planos coordenados están dados por:

\vspace{-5mm}
{\footnotesize
\begin{align*}
	\mu_{yz} &= \int _0^2\int _{x^3}^{\sqrt{8 (10-x)}}\int _0^{10-x-\frac{y^2}{8}} x^2y\,dz\,dy\,dx \\
	&= \frac{9952}{45} \\[1ex]
	\mu_{xz} &= \int _0^2\int _{x^3}^{\sqrt{8 (10-x)}}\int _0^{10-x-\frac{y^2}{8}} xy^2\,dz\,dy\,dx \\
	&= \frac{5120 \left(18700 \sqrt{5}-34721\right)}{35343} \\[1ex]
	\mu_{xy} &= \int _0^2\int _{x^3}^{\sqrt{8 (10-x)}}\int _0^{10-x-\frac{y^2}{8}} xyz\,dz\,dy\,dx. \\
	&= \frac{177584}{315}
\end{align*}
}

\vskip0.2cm
Con los elementos anteriores se obtienen las coordenadas del centro de masa:
$(\overline{x},\overline{y},\overline{z})=\left(\frac{8708}{8125},\frac{512 \sqrt{5}}{39}-\frac{4444288}{182325},\frac{22198}{8125}\right)$.
Este proceso se hace en \verb"Mathematica" como sigue:


\vspace{0.6em}
\noindent
\begin{tcolorbox}[breakable,enhanced,title=\bf Instrucción en \verb"Mathematica", top=2pt,bottom=2pt,boxsep=0pt,before skip=2pt,after skip=0pt]
\begin{Verbatim}[formatcom=\letraverbatim]
M=Integrate[x*y,{x,0,2},{y,x^3,Sqrt[8*(10-x)]},
{z,0,10-x-y^2/8}];
Mxy=Integrate[x*y*z,{x,0,2},{y,x^3,Sqrt[8*(10-x)]},
{z,0,10-x-y^2/8}];
Mxz=Integrate[x*y^2,{x,0,2},{y,x^3,Sqrt[8*(10-x)]},
{z,0,10-x-y^2/8}];
Myz=Integrate[x^2*y,{x,0,2},{y,x^3,Sqrt[8*(10-x)]},
{z,0,10-x-y^2/8}];
Simplify[{Myz/M,Mxz/M,Mxy/M}]
\end{Verbatim}
\end{tcolorbox}
\vspace{0.7em}
\noindent

El volumen también se puede obtener con el siguiente procedimiento.

%\begin{tcolorbox}[title=\bfseries Parametrizando el sólido]
%\small
La proyección del sólido en el plano $xy$, ver figura (\ref{figura115}),
se parametriza considerando dos puntos
{\footnotesize $A(t,t^{3})$} y {\footnotesize $B\left( t,\sqrt{8\left( 10-t\right) }\right) $} para {\footnotesize $t\in\lbrack 0,2].$}
Simplificando la expresión {\footnotesize $(1-u)A+uB$}, se obtiene
{\footnotesize $\left\langle t,t^{3}(1-u)+u\sqrt{8\left(10-t\right) }\right\rangle.$}
\vskip0.2cm
La coordenada $z,$ sobre la superficie $z=10-x-\frac{y^{2}}{8},$
que abreviaremos como $\overline{z}$, satisface que
{\footnotesize $\overline{z}=10-t-\frac{1}{8}\left( t^{3}(1-u)+u\sqrt{8\left( 10-t\right) }\right) ^{2}.$}
\vskip0.2cm

Con dos puntos de la forma
{\small \[P\left( t,t^{3}(1-u)+u\sqrt{8\left( 10-t\right) },0\right) \]} y
{\small \[Q\left( t,t^{3}(1-u)+u\sqrt{8\left( 10-t\right) },\overline{z}\right)\]}
se forma la combinación $(1-v) P+vQ,$ que define la función
%\vskip0.2cm
{\small
\[{\bf r}(t,u,v)=\left\langle t,t^{3}(1-u)+u\sqrt{8(10-t) },v \,\overline{z} \right\rangle,\]}
%\vskip0.2cm
que parametriza el sólido, con $t\in [0,2],\,u\in [0,1],\,v\in [0,1].$
	% Ahora bien, %\begin{eqnarray*}
	%{\bf r}_{t} &=& \bigg\langle 1,3t^{2}(1-u) -\frac{\sqrt{2}u}{\sqrt{10-t}},\\ && v\left( \frac{1}{4}\left( t^{3}(1-u) +u\sqrt{8(10-t)}\right) \left( 3t^{2}(u-1) +\frac{\sqrt{2}u}{\sqrt{10-t}}\right) -1\right) \bigg\rangle \\
	%{\bf r}_{u} &=&\left\langle 0,\sqrt{8(10-t)}-t^{3},\frac{v}{4}\left( t^{3}(1-u) +u\sqrt{8(10-t)}\right) \left( t^{3}-\sqrt{8(10-t)}\right) \right\rangle \\ {\bf r}_{v} &=&\left\langle 0,0,10-\frac{1}{8}\left( t^{3}(u-1) -u\sqrt{8(10-t)}\right) ^{2}-t\right\rangle
	%\end{eqnarray*}
El elemento diferencial de volumen
$\left\vert {\bf r}_{t}\cdot \left({\bf r}_{u}\times {\bf r}_{v}\right) \right\vert$
esta dado por
\vspace{-4mm}
%\vskip0.2cm

{\footnotesize
\[\bigg\vert \dfrac{(u-1)(t^{3}-\sqrt{8(10-t)})}{8}\bigg(8(u+1)(10-t)
+t^6(u-1)-2t^{3}u\sqrt{8(10-t)}\bigg) \bigg\vert.\]
}

%\vskip0.2cm
Su integración determina el volumen del sólido
%\vskip0.2cm
{\footnotesize
\begin{equation*}
\int_{0}^{1}\int_{0}^{1}\int_{0}^{2}\left\vert {\bf r}_{t}\cdot \left({\bf r}_{u}\times {\bf r}_{v}\right) \right\vert dtdudv=\frac{320}{3}\sqrt{5}-\frac{2488}{15}.
\end{equation*}
}
%\vskip0.2cm
%\end{tcolorbox}


Con las siguientes indicaciones se pueden efectuar estos cálculos.
\vskip0.2cm

%\vspace{0.6em}
\noindent
\begin{tcolorbox}[breakable,enhanced,title=\bf Instrucción en \verb"Mathematica", top=2pt,bottom=2pt,boxsep=0pt,before skip=2pt,after skip=0pt]
\begin{Verbatim}[formatcom=\letraverbatim]
{x,y}=(1-u)*{t,t^3}+u*{t,Sqrt[8*(10-t)]};
r=Simplify[{x,y,w*(10-x-y^2/8)}];
A=Abs[Dot[D[r,w],Cross[D[r,u],D[r,t]]]];
Integrate[A,{t,0,2},{u,0,1},{w,0,1}]
\end{Verbatim}
\end{tcolorbox}
\vspace{0.7em}
\noindent



\begin{ejer}
Plantear y evaluar las integrales que determinan el volumen del sólido $E$ de la sección \secref{2cilindros2planos}
En los órdenes $dydxdz$ y $dxdydz.$
\end{ejer}

\subsection*{Dos cilindros}
Las ecuaciones $x^2+z^2=16$ y $2x^2+y^2=8x$ representan un cilindro circular y un cilindro elíptico. Estas superficies acotan en el primer octante, el sólido que se muestra enseguida.


\vskip0.4cm
\noindent
\begin{minipage}{\textwidth}
	\centering
	\captionof{figure}{Visualización del sólido definido por $x^2 + z^2 \leq 16$ \quad y \quad $2x^2 + y^2 \leq 8x$}
	\label{fig:solido_cilindro_elipse}
	\begin{minipage}{0.22\linewidth}
		\centering
		\includegraphics[width=\linewidth]{Img/29a.pdf}
	\end{minipage}
	\begin{minipage}{0.22\linewidth}
		\centering
		\includegraphics[width=\linewidth]{Img/29b.pdf}
	\end{minipage}
	\begin{minipage}{0.22\linewidth}
		\centering
		\includegraphics[width=\linewidth]{Img/29c.pdf}
	\end{minipage}
	\begin{minipage}{0.22\linewidth}
		\centering
		\includegraphics[width=\linewidth]{Img/29d.pdf}
	\end{minipage}
	
	\vspace{0.3em}
	\fuentepropia
\end{minipage}
\vskip0.4cm



En \verb"Mathematica", estas figuras se consiguen con las siguientes sentencias:

\vspace{0.6em}
\noindent
\begin{tcolorbox}[breakable,enhanced,title=\bf Instrucción en \verb"Mathematica", top=2pt,bottom=2pt,boxsep=0pt,before skip=2pt,after skip=0pt]
\begin{Verbatim}[formatcom=\letraverbatim]
ContourPlot3D[{16==x^2+z^2,2*x^2+y^2==8*x},{x,0,4},
{y,0,3},{z,0,4},PlotPoints->60,Ticks->{{0,2,4},
{0,1,2},{0,2,4}},Mesh->False,ContourStyle->{Orange,
Blend[{Yellow,Pink,Green}]}]
RegionPlot3D[0<z<Sqrt[16-x^2]&&0<y<Sqrt[8*x-2*x^2]&&0
<x<4,{x,0,4},{y,0,3},{z,0,4},PlotPoints->90,
PlotStyle->Directive[Orange,Specularity[White,40]],
Mesh->False]
\end{Verbatim}
\end{tcolorbox}
\vspace{0.7em}
\noindent


Para establecer con claridad los límites en las integrales triples respectivas, que determinan el volumen del sólido, se verifica que en el plano $xz$ la región es un cuarto del círculo centrado en el origen y radio $4.$
Entre tanto, en el plano $xy$, la expresión $2x^2+y^2=8x$ equivale a $2(x-2)^2+y^2=8$ y representa una elipse con centro en el punto $(2,0).$

\vskip0.4cm
\noindent
\begin{minipage}{\textwidth}
	\centering
	\captionof{figure}{Proyecciones en los planos coordenados.}
	\label{fig:proyecciones_cilindro_elipse_curva}
	\begin{minipage}{0.25\linewidth}
		\centering
		\vspace{-18mm}
		\includegraphics[width=\linewidth]{Fig/424.pdf}
	\end{minipage}
	\hspace{1em}
	\begin{minipage}{0.25\linewidth}
		\centering
		\vspace{-18mm}
		\includegraphics[width=\linewidth]{Fig/425.pdf}
	\end{minipage}
	\hspace{1em}
	\begin{minipage}{0.38\linewidth}
		\centering
		\vspace{-18mm}
		\includegraphics[width=\linewidth]{Fig/426.pdf}
	\end{minipage}
	
	\vspace{-12mm}
	\fuentepropia
\end{minipage}
\vskip0.2cm

Con lo anterior, presentamos las integrales triples según el orden de integración señalado.
\vskip0.2cm
Al integrar primero con respecto a $x,$ se debe tener en cuenta las dos regiones proyectadas en el plano $yz$, en un caso se parte de la parte negativa del cilindro elíptico y se llega hasta la parte positiva del mismo cilindro.
Mientras que en la otra región $x$ recorre los puntos que van desde la parte negativa del cilindro elíptico hasta el cilindro circular.
De la expresión $x^2+z^2=16$, se sigue que {\footnotesize $x=\pm \sqrt{16-z^2}.$} Del reemplazo en la otra ecuación se obtiene
{\footnotesize $16-z^2+2y^2=4\sqrt{16-z^2},$} que equivale a
{\footnotesize $\left( z^2-y^2\right) ^2=16\left( z^2+2y^2\right).$}
De esta expresión se obtiene {\footnotesize $z=\sqrt{8+\frac{y^2}{2}\pm 2\sqrt{16-2y^2}}.$}
% \begin{eqnarray*} z &=& \frac{1}{2}\sqrt{32+2y^2\pm 8\sqrt{16-2y^2}} =\sqrt{8+\frac{y^2}{2}\pm 2\sqrt{16-2y^2}} \end{eqnarray*}
Con esto, el volumen del sólido está dado por
\vspace{-4mm}

{\small
\begin{multline*}
\bigintsss_{0}^{\sqrt{8}}\bigintsss_{0}^{\sqrt{8+\frac{y^2}{2}-2\sqrt{16-2y^2}}}\bigintsss_{2-\frac{1}{2}\sqrt{16-2y^2}}^{2+\frac{1}{2}\sqrt{16-2y^2}}dxdzdy \\
+\bigintsss_{0}^{\sqrt{8}}\bigintsss_{\sqrt{8+\frac{y^2}{2}-\sqrt{64-8y^2}}}^{\sqrt{8+\frac{y^2}{2}+\sqrt{64-8y^2}}}\bigintsss_{2-\frac{1}{2}\sqrt{16-2y^2}}^{\sqrt{16-z^2}}dxdzdy = \frac{176}{3}-12\sqrt{2}\ln( 3+2\sqrt{2}).
\end{multline*}
}
\vskip0.2cm
De la expresión $\left( z^2-y^2\right) ^2=16\left( z^2+2y^2\right),$ se puede obtener \\
$y=\sqrt{2z^2-32\pm 8\sqrt{16-z^2}},$ entonces tenemos otra forma de hallar el mismo volumen.
\begin{multline*}
\bigintsss_{0}^{2\sqrt{3}}\bigintsss_{\sqrt{2z^2-32+8\sqrt{16-z^2}}}^{\sqrt{8}}\bigintsss_{2-\frac{1}{2}\sqrt{16-2y^2}}^{2+\frac{1}{2}\sqrt{16-2y^2}}dxdydz \\
+\bigintsss_{0}^{4}\bigintsss_{0}^{\sqrt{2z^2-32+8\sqrt{16-z^2}}}\bigintsss_{2-\frac{1}{2}\sqrt{16-2y^2}}^{\sqrt{16-z^2}}dxdydz \approx 28.7519.
\end{multline*}
\vskip0.2cm
Se debe observar que de la expresión $2( x-2) ^2+y^2=8$ se sigue que $x=2\pm \frac{1}{2}\sqrt{16-2y^2}$,
por tanto, el volumen del sólido esta dado por
\vspace{-1mm}

{\small
\[
\int_{\scriptscriptstyle 0}^{\scriptscriptstyle 4}\int_{\scriptscriptstyle 0}^{\scriptscriptstyle\sqrt{8x-2x^2}}\int_{\scriptscriptstyle 0}^{\scriptscriptstyle\sqrt{16-x^2}}dzdydx \! = \! \int_{\scriptscriptstyle 0}^{\scriptscriptstyle\sqrt{8}}\int_{\scriptscriptstyle 2-\frac{1}{2}\sqrt{16-2y^2}}^{\scriptscriptstyle 2+\frac{1}{2}\sqrt{16-2y^2}}\int_{\scriptscriptstyle 0}^{\scriptscriptstyle\sqrt{16-x^2}}dzdxdy \! \approx \! 28.75.
\] 
}
\\ 
Para el orden $dydxdz$, $y$ recorre los puntos desde el plano $xz$ hasta el cilindro elíptico; en este caso, las integrales correspondientes son las siguientes.

\vspace{-4mm}
{\small
\[\int_{0}^{4}\int_{0}^{\sqrt{16-z^2}}\int_{0}^{\sqrt{8x-2x^2}}dydxdz=\int_{0}^{4}\int_{0}^{\sqrt{16-x^2}}\int_{0}^{\sqrt{8x-2x^2}}dydzdx \approx 28.75.\]
}

\vskip0.2cm
Una de las integrales anteriores se evaluó numéricamente en \verb"Mathematica"
ejecutando la siguiente instrucción.

\vspace{0.6em}
\noindent
\begin{tcolorbox}[breakable,enhanced,title=\bf Instrucción en \verb"Mathematica", top=2pt,bottom=2pt,boxsep=0pt,before skip=2pt,after skip=0pt]
\begin{Verbatim}[formatcom=\letraverbatim]
NIntegrate[1,{y,0,Sqrt[8]},{z,0,Sqrt[8+y^2/2
Sqrt[64-8*y^2]]},{x,2-Sqrt[16-2*y^2]/2,2+
Sqrt[16-2*y^2]/2}]+NIntegrate[1,{y,0,Sqrt[8]},
{z,Sqrt[8+y^2/2-Sqrt[64-8*y^2]],Sqrt[8+y^2/2+
Sqrt[64-8*y^2]]},{x,2-Sqrt[16-2*y^2]/2,Sqrt[16-z^2]}]
\end{Verbatim}
\end{tcolorbox}
\vspace{0.7em}
\noindent


La otra alternativa para calcular el volumen se muestra a continuación.
\vskip0.2cm

%\begin{tcolorbox}[title=\bfseries Parametrizando el sólido]
\noindent
La proyección del sólido en el plano $xz$ se describe mediante las ecuaciones
$x=4u\cos(t),\,\, z=4u\sin(t),$ para $u\in [0,1]$
y $t\in [0,\pi/2 ].$ De estas expresiones se obtiene $y$ en la ecuación del cilindro parabólico, es decir $y = \sqrt{8x-2x^2} =4\sqrt{2}\sqrt{u(1-u\cos(t)) \cos(t)}$.

\vskip0.2cm
\noindent
\begin{minipage}{\textwidth}
	\centering
	\captionof{figure}{Visualización del sólido definido por $x^2 + z^2 = 16$, \quad $2x^2 + y^2 = 8x$}
	\label{fig:solido_cilindro_elipse_3D_2}
	\begin{minipage}{0.35\linewidth}
		\centering
		\vspace{-15mm}
		\includegraphics[width=\linewidth]{Fig/427.pdf}
	\end{minipage}
	
	\vspace{-8mm}
	\fuentepropia
\end{minipage}
\vskip0.2cm


%\vskip0.8cm

El sólido se parametriza considerando un punto
%\vskip0.2cm 
\[P(4u\cos(t),0,4u\sin(t).\]
%\vskip0.2cm 
en el plano $xz$
y otro punto sobre el cilindro elíptico
%\vskip0.2cm
\[Q(4u\cos(t),4\sqrt{2}\sqrt{u(1-u\cos(t))\cos(t)},4u\sin(t)).\]
%\vskip0.2cm
Simplificando la combinación convexa $(1-w)P+wQ$, se define la siguiente función:
%\vskip0.2cm
\[{\bf r}(t,u,w) = \langle 4u\cos(t),4w\sqrt{2}\sqrt{u(1-u\cos(t))\cos(t)},4u\sin(t)\rangle,\]
%\vskip0.2cm
con $u\in [0,1],$ $t\in [ \pi/2 ]$ y $w\in [0,1]$.
% Ahora hallamos los siguientes vectores \begin{eqnarray*}
%{\bf r}_{t} &=& \langle -4u\sin(t),2w\frac{\sqrt{2}u\sin(t)(2u\cos(t)-1)}{\sqrt{u(1-u\cos(t))\cos(t)}},4u\cos(t)\rangle \\
%{\bf r}_{u} &=&\langle 4\cos(t),2w\frac{\sqrt{2}\cos(t)(1-2u\cos(t))}{\sqrt{(u(1-u\cos(t))\cos t))}},4\sin(t)\rangle \\
%{\bf r}_{w} &=&\langle 0,4\sqrt{2}\sqrt{u(1-u\cos(t))\cos(t)},0 \rangle
%\end{eqnarray*}
% además ${\bf r}_{t}\times {\bf r}_{u}=\langle -8w\sqrt{2}\sqrt{u}\frac{2u\cos(t)-1}{\sqrt{(1-u\cos(t))\cos(t)}},16u,0\rangle $
El triple producto vectorial proporciona el elemento diferencial de volumen,
%\vskip0.2cm
\[|{\bf r}_{w}\cdot({\bf r}_{t}\times {\bf r}_{u})|=64\sqrt{2}u\sqrt{u(1-u\cos (t))\cos(t)},\]
%\vskip0.2cm
que se integra y nos proporciona el volumen del sólido
\[\int_{0}^{1}\int_{0}^{1}\int_{0}^{\pi/2}64\sqrt{2}u\sqrt{u(1-u\cos(t))\cos(t)}dtdudw \approx 28.752.\]
%\end{tcolorbox}
\vskip0.2cm
Este cálculo se obtiene de manera numérica en \verb"Mathematica" ejecutando las siguientes instrucciones:


\vspace{0.6em}
\noindent
\begin{tcolorbox}[breakable,enhanced,title=\bf Instrucción en \verb"Mathematica", top=2pt,bottom=2pt,boxsep=0pt,before skip=2pt,after skip=0pt]
\begin{Verbatim}[formatcom=\letraverbatim]
{x,z,y}={4*u*Cos[t],4*u*Sin[t],Sqrt[8*x-2*x^2]};
r=Simplify[(1-w)*{x,0,z}+w*{x,y,z}]
B=Simplify[Abs[Dot[D[r,w],Cross[D[r,t],D[r,u]]]]]
NIntegrate[B,{t,0,Pi/2},{u,0,1},{w,0,1}]
\end{Verbatim}
\end{tcolorbox}
\vspace{0.7em}
\noindent


\subsection*{Dos cilindros y un plano}
Hallar el volumen del sólido acotado por los gráficos de las ecuaciones.
$z=1-x^2,$ $x^2+y^2=1$ y $z=0$. Las siguientes figuras ilustran esta situación.


\vskip0.4cm
\noindent
\begin{minipage}{\textwidth}
	\centering
	\captionof{figure}{Visualización del sólido definido por $x^2 + y^2 \leq 1$ \quad y \quad $0 \leq z \leq 1 - x^2$.}
	\label{fig:solido_disco_parabolico}
	\begin{minipage}{0.24\linewidth}
		\centering
		\vspace{-8mm}
		\includegraphics[width=\linewidth]{Img/85c.pdf}
	\end{minipage}
	\hspace{1em}
	\begin{minipage}{0.28\linewidth}
		\centering
		\vspace{-8mm}
		\includegraphics[width=\linewidth]{Fig/428.pdf}
	\end{minipage}
	\hspace{1em}
	\begin{minipage}{0.18\linewidth}
		\centering
		\vspace{-8mm}
		\includegraphics[width=\linewidth]{Img/85b.pdf}
	\end{minipage}
	
	\vspace{-4mm}
	\fuentepropia
\end{minipage}
\vskip0.2cm

Estas figuras se obtienen con la ayuda de las siguientes instrucciones:


\vspace{0.6em}
\noindent
\begin{tcolorbox}[breakable,enhanced,title=\bf Instrucción en \verb"Mathematica", top=2pt,bottom=2pt,boxsep=0pt,before skip=2pt,after skip=0pt]
\begin{Verbatim}[formatcom=\letraverbatim]
ContourPlot3D[{z==1-x^2,x^2+y^2==1},{x,-1,1},{y,-1,1},
{z,0,1},Mesh->False,ContourStyle->{Cyan,Orange}]
A=RegionPlot3D[0<z<1-x^2&&x^2+y^2<1,{x,-1,1},{y,-1,1},
{z,0,1},PlotStyle->RGBColor[0.8,0.2,0.2],Mesh->False,
PlotPoints->90,Axes->False,Boxed->False];
B=ParametricPlot3D[{{Cos[t],Sin[t],1-Cos[t]^2},
{Cos[t],Sin[t],0}},{t,0,2Pi},PlotStyle->{Gray,Gray}];
Show[A,B]
\end{Verbatim}
\end{tcolorbox}
\vspace{0.7em}
\noindent


\noindent
\begin{minipage}{\textwidth}
	\centering
	\captionof{figure}{Proyección del sólido en los planos coordenados.}
	\label{fig:proyeccion_parabolico_disco}
	\begin{minipage}{0.29\linewidth}
		\centering
		\vspace{-14mm}
		\includegraphics[width=\linewidth]{Fig/429.pdf}
	\end{minipage}
	\hspace{1em}
	\begin{minipage}{0.29\linewidth}
		\centering
		\vspace{-14mm}
		\includegraphics[width=\linewidth]{Fig/430.pdf}
	\end{minipage}
	\begin{minipage}{0.29\linewidth}
		\centering
		\vspace{-14mm}
		\includegraphics[width=\linewidth]{Fig/431.pdf}
	\end{minipage}
	
	\vspace{-7mm}
	\fuentepropia
\end{minipage}


\vskip0.2cm
En el plano $xy$ la ecuación $x^2+y^2=1$ determina una proyección del sólido, esto ayuda a plantear las integrales que permiten obtener el volumen.
\vspace{-2mm}
\[\int_{-1}^{1}\int_{-\sqrt{1-x^2}}^{\sqrt{1-x^2}}\int_{0}^{1-x^2}dzdydx= \int_{-1}^{1}\int_{-\sqrt{1-y^2}}^{\sqrt{1-y^2}}\int_{0}^{1-x^2}dzdxdy= \frac{3}{4}\pi.\]
%\vskip0.2cm
En el plano $yz,$ se satisface que $1-z+y^2=1,$ es decir $z=y^2$ cuya figura es una parábola. Por tal motivo, el volumen está dado por:
\begin{multline*}
\int_{0}^{1}\int_{-\sqrt{z}}^{\sqrt{z}}\int_{-\sqrt{1-z}}^{\sqrt{1-z}}dxdydz +\int_{0}^{1}\int_{-1}^{-\sqrt{z}}\int_{-\sqrt{1-y^2}}^{\sqrt{1-y^2}}dxdydz \\ +\int_{0}^{1}\int_{\sqrt{z}}^{1}\int_{-\sqrt{1-y^2}}^{\sqrt{1-y^2}}dxdydz
=\int_{-1}^1\int_{y^2}^1\int_{-\sqrt{1-z}}^{\sqrt{1-z}}dxdzdy \\ +\int_{-1}^{1}\int_0^{y^2}\int_{-\sqrt{1-y^2}}^{\sqrt{1-y^2}}dxdzdy=\frac{3}{4}\pi.
\end{multline*}

% \begin{multline*} \int_{0}^{1}\int_{-\sqrt{z}}^{\sqrt{z}}\int_{-\sqrt{1-z}}^{\sqrt{1-z}}dxdydz+\int_{0}^{1}\int_{-1}^{-\sqrt{z}}\int_{-\sqrt{1-y^2}}^{\sqrt{1-y^2}}%dxdydz \\ +\int_{0}^{1}\int_{\sqrt{z}}^{1}\int_{-\sqrt{1-y^2}}^{\sqrt{1-y^2}}dxdydz \approx 2.3562 \end{multline*}
En el plano $xz$, la proyección es una región acotada por una porción de parábola y en el plano $x$,
el volumen está dado por
%\vskip0.2cm
\[\int_{-1}^{1}\int_{0}^{1-x^2}\int_{-\sqrt{1-x^2}}^{\sqrt{1-x^2}}dydzdx = \int_{0}^{1}\int_{-\sqrt{1-z}}^{\sqrt{1-z}}\int_{-\sqrt{1-x^2}}^{\sqrt{1-x^2}}dydxdz \approx 2.3562.\]
%\vskip0.2cm
En coordenadas cilíndricas, esta integral se evalúa como sigue:
%\vskip0.2cm
\[\int_0^{2\pi} \int_0^1 \int_0^{1-r^2\cos^2 \theta}, r dzdrd \theta =\frac{3\pi}{4}.\]
%\vskip0.2cm
Algunas de estas integrales se evalúan mediante el uso de las siguientes instrucciones:

\vspace{0.6em}
\noindent
\begin{tcolorbox}[breakable,enhanced,title=\bf Instrucción en \verb"Mathematica", top=2pt,bottom=2pt,boxsep=0pt,before skip=2pt,after skip=0pt]
\begin{Verbatim}[formatcom=\letraverbatim]
Integrate[1,{x,-1,1},{y,-Sqrt[1-x^2],Sqrt[1-x^2]},
{z,0,1-x^2}]
Integrate[1,{y,-1,1},{x,-Sqrt[1-y^2],Sqrt[1-y^2]},
{z,0,1-x^2}]
Integrate[1,{y,-1,1},{z,y^2,1},{x,-Sqrt[1-z],
Sqrt[1-z]}]+Integrate[1,{y,-1,1},{z,0,y^2},
{x,-Sqrt[1-y^2],Sqrt[1-y^2]}]
Integrate[1,{x,-1,1},{z,0,1-x^2},{y,-Sqrt[1-x^2],
Sqrt[1-x^2]}]
Integrate[1,{z,0,1},{x,-Sqrt[1-z],Sqrt[1-z]},
{y,-Sqrt[1-x^2],Sqrt[1-x^2]}]
\end{Verbatim}
\end{tcolorbox}
\vspace{0.7em}
\noindent


El volumen también se puede obtener de la siguiente manera.

%\begin{tcolorbox}[title=\bfseries Parametrizando el sólido]
La base del sólido es un círculo unitario que se parametriza {\small $\left(u\cos \theta ,u\sin \theta ,0\right),$} sobre el cilindro parabólico se satisface que $z=1-x^2=1-u^2\cos ^2\theta .$ El sólido se parametriza formando una combinación convexa entre un punto en la base y otro punto sobre el cilindro parabólico, esto es
$ (1-w)(u\cos \theta ,u\sin \theta ,0) +w( u\cos \theta ,u\sin \theta ,1-u^2\cos ^2\theta);$
al simplificarse, nos permite definir la función
%\vskip0.2cm
{\small
\[{\bf r}\left( u,\theta ,w\right) =\langle u\cos \theta , u\sin \theta ,w(1-u^2\cos ^2\theta )\rangle,\]
}
%\vskip0.2cm
donde $\theta \in [0,2\pi]$, $u \in [0,1]$ y $w \in [0,1]$.

\vskip0.3cm
Con esta función hallamos los siguientes vectores
\begin{eqnarray*}
{\bf r}_{u} &=& \left\langle \cos \theta,\sin \theta, -2uw\cos ^2\theta \right\rangle \\
{\bf r}_{\theta } &=& \left\langle -u\sin \theta , u\cos \theta , 2u^2w\cos \theta \sin \theta \right\rangle \\
{\bf r}_{w} &=& \left\langle 0 , 0 , 1-u^2\cos ^2\theta \right\rangle
\end{eqnarray*}
El elemento diferencial de volumen está dado por

\[\left\vert {\bf r}_{u}\cdot \left( {\bf r}_{\theta }\times {\bf r}_{w}\right) \right\vert =\left\vert u\left( 1-u^2\cos ^2\theta \right) \right\vert,\]

entonces, el volumen del sólido esta dado por
\vspace{-2mm}
{\small
\[\int_{0}^{1}\int_{0}^{1}\int_{0}^{2\pi }u\left( 1-u^2\cos ^2\theta
\right) d\theta dudw= \frac{3\pi}{4}.\]
}
%\end{tcolorbox}

Este cálculo se hace en \verb"Mathematica" ejecutando las siguientes instrucciones.



\vspace{0.6em}
\noindent
\begin{tcolorbox}[breakable,enhanced,title=\bf Instrucción en \verb"Mathematica", top=2pt,bottom=2pt,boxsep=0pt,before skip=2pt,after skip=0pt]
\begin{Verbatim}[formatcom=\letraverbatim]
{x,y,z}={u*Cos[t],u*Sin[t],1-x^2};
r=(1-w)*{x,y,0}+w*{x,y,z}
A=Simplify[Abs[Dot[D[r,t],Cross[D[r,u],D[r,w]]]]]
[A,{t,0,2*Pi},{u,0,1},{w,0,1}]
\end{Verbatim}
\end{tcolorbox}
\vspace{0.7em}
\noindent


\subsection*{Tres cilindros y un plano}
Las ecuaciones $z=4-x^2,\,\, 2x=y^2$ y $4x=y^4-2y^2$ para $y \geq 0,$ acotan un sólido.
Las figuras de las expresiones que lo determinan y el sólido se muestran a continuación.




\vskip0.2cm
\noindent
\begin{minipage}{\textwidth}
	\centering
	\captionof{figure}{Visualización del sólido definido por $0 \leq z \leq 4 - x^2$ \quad y \quad $y^4 - 2y^2 \leq 4x \leq 2y^2$}
	\label{fig:solido_curvas_parabolicas}
	\begin{minipage}{0.20\linewidth}
		\centering
		\includegraphics[width=\linewidth]{Img/17c.pdf}
	\end{minipage}
	\begin{minipage}{0.20\linewidth}
		\centering
		\includegraphics[width=\linewidth]{Img/17d.pdf}
	\end{minipage}
	\begin{minipage}{0.20\linewidth}
		\centering
		\includegraphics[width=\linewidth]{Img/17b.pdf}
	\end{minipage}
	
	
	\vspace{0.3em}
	\fuentepropia
\end{minipage}
\vskip0.2cm

También se presentan las instrucciones que se requieren para graficarlo.


\vspace{0.6em}
\noindent
\begin{tcolorbox}[breakable,enhanced,title=\bf Instrucción en \verb"Mathematica", top=2pt,bottom=2pt,boxsep=0pt,before skip=2pt,after skip=0pt]
\begin{Verbatim}[formatcom=\letraverbatim]
ContourPlot3D[{z==4-x^2,2*x==y^2,4*x==y^4-2*y^2},
{x,-1,2},{y,0,2},{z,0,4},ContourStyle->{Orange,
Blend[{Yellow,Pink,Green}],LightGray},Mesh->None,
PlotPoints->60]
RegionPlot3D[0<z<4-x^2&&(y^4-2*y^2)/4<x<y^2/2,
{x,0,2},
{y,0,2},{z,0,4},PlotStyle->Blend[{Yellow,Orange,Red}],
PlotPoints->90,Mesh->None]
\end{Verbatim}
\end{tcolorbox}
\vspace{0.7em}
\noindent


La ecuación $z = 4-x^2$ en el plano $xz$ permite ver la proyección en dicho plano de la región de integración. Lo mismo sucede con las ecuaciones $4x = y^4 -2y^2$ y $2x = y^2$, que se dibujan en el plano $xy.$ En este caso $s = 4 - x^2$ equivale a $x = \pm\sqrt{4-z},$ mientras que $2x = y^2$ equivale a $y = \pm \sqrt{2x}$. La expresión $4x =y^4- 2y^2$ nos lleva a que $y = \pm \sqrt{1\pm \sqrt{1+ 4x}},$ es decir hay cuatro funciones que determinan la curva, con $x$ como variable independiente, es por esta razón que el gráfico posee diferentes colores para enfatizar dicha situación.


\vskip0.2cm
\noindent
\begin{minipage}{\textwidth}
	\centering
	\captionof{figure}{Proyecciones del sólido definido por $y^4 - 2y^2 \leq 4x \leq 2y^2$, \quad $0 \leq z \leq 4 - x^2$.}
	\label{fig:proyecciones_curvas_parabolicas}
	\begin{minipage}{0.58\linewidth}
		\centering
		\vspace{-32mm}
		\includegraphics[width=\linewidth]{Fig/432.pdf}
	\end{minipage}
	\hspace{1em}
	\begin{minipage}{0.32\linewidth}
		\centering
		\vspace{-32mm}
		\includegraphics[width=\linewidth]{Fig/433.pdf}
	\end{minipage}
	
	\vspace{-27mm}
	\fuentepropia
\end{minipage}


%\vskip0.4cm
Con ayuda de estos gráficos, podemos plantear las integrales triples que determinan el volumen del sólido.
\vspace{-1mm}
{\small
\begin{multline*}
\hskip-0.4cm \bigintss_{-\frac{1}{4}}^{0}\bigintss_{\sqrt{1-\sqrt{1+4x}}}^{\sqrt{1+\sqrt{1+4x}}}\bigintss_{0}^{4-x^2}dzdydx+\bigintss_{0}^2\bigintss_{\sqrt{2x}}^{\sqrt{1+\sqrt{1+4x}}}\bigintss_{0}^{4-x^2}dzdydx \\
\hskip-0.4cm \bigintss_{-\frac{1}{4}}^{0}\bigintss_{0}^{4-x^2}\bigintss_{\sqrt{1-\sqrt{1+4x}}}^{\sqrt{1+\sqrt{1+4x}}}dydzdx+\bigintss_{0}^2\bigintss_{0}^{4-x^2}\bigintss_{\sqrt{2x}}^{\sqrt{1+\sqrt{1+4x}}}dydzdx \\
= \bigintss_{0}^2 \bigintss_{\frac{y^{4}-2y^2}{4}}^{\frac{y^2}{2}}\bigintss_{0}^{4-x^2}dzdxdy= \frac{169472}{45045}.
\end{multline*}
}

En el plano $yz$ existen tres regiones en las que se divide la proyección del sólido; una de ellas es muy pequeña, está sombreada en negro y está
acotada por la recta $z=4$ y la curva $z=4- (y^4 -2y^2)^2/16,$ para $y \in [0, \sqrt{2}].$
\vskip0.2cm

\noindent
\begin{minipage}{\textwidth}
	\centering
	\captionof{figure}{Proyección del sólido en el plano $yz$.}
	\label{fig:proyeccion_yz_curvas}
	\begin{minipage}{0.45\linewidth}
		\centering
		\vspace{-20mm}
		\includegraphics[width=\linewidth]{Fig/434.pdf}
	\end{minipage}
	\hspace{1em}
	\begin{minipage}{0.45\linewidth}
		\centering
		\vspace{-20mm}
		\includegraphics[width=\linewidth]{Fig/435.pdf}
	\end{minipage}
	
	\vspace{-14mm}
	\parbox{\linewidth}{\centering\fontsize{8pt}{10pt}\selectfont\textbf{Fuente:} elaboración propia.}
\end{minipage}

\vskip0.2cm
En este caso, el volumen está dado por:

\vspace{-8mm}
\begin{multline*}
\bigintss_{0}^2\bigintss_{0}^{4-\frac{y^{4}}{4}}\bigintss_{\frac{y^{4}-2y^2}{4}}^{\frac{y^2}{2}}dxdzdy+\bigintss_{0}^2\bigintss_{4-\frac{y^{4}}{4}}^{4-\frac{\left(y^{4}-2y^2\right) ^2}{16}}\bigintss_{\frac{y^{4}-2y^2}{4}}^{\sqrt{4-z}}dxdzdy \\
+\bigintss_{0}^{\sqrt{2}}\bigintss_{4-\frac{\left( y^{4}-2y^2\right) ^2}{16}}^{4}\bigintss_{-\sqrt{4-z}}^{\sqrt{4-z}}dxdzdy \approx 3.7623.
\end{multline*}

La evaluación numérica de la integral anterior se puede obtener ejecutando las siguientes instrucciones.



\vspace{0.6em}
\noindent
\begin{tcolorbox}[breakable,enhanced,title=\bf Instrucción en \verb"Mathematica", top=2pt,bottom=2pt,boxsep=0pt,before skip=2pt,after skip=0pt]
\begin{Verbatim}[formatcom=\letraverbatim]
NIntegrate[1,{y,0,2},{z,0,4-y^4/4},{x,(y^4-2*y^2)/4,
y^2/2}]+NIntegrate[1,{y,0,2},{z,4-y^4/4,4-(y^4-
2*y^2)^2/16},{x,(y^4-2*y^2)/4,Sqrt[4-z]}]+NIntegrate[1,
{y,0,Sqrt[2]},{z,4-(y^4-2*y^2)^2/16,4},{x,-Sqrt[4-z],
Sqrt[4-z]}]
\end{Verbatim}
\end{tcolorbox}
\vspace{0.7em}
\noindent


También se puede obtener dicho valor como se muestra enseguida.



%\begin{tcolorbox}[title=\bfseries Parametrizando el sólido]
\noindent
Sobre el plano $xy,$ las ecuaciones $4x=y^{4}-2y^2$ y $2x=y^2,$ se parametrizan como
$\overrightarrow{\alpha}(t)=\left\langle \frac{1}{2}t^2,t\right\rangle $
y $\overrightarrow{\beta}(t)=\left\langle \frac{1}{4}\left(t^{4}-2t^2\right),t\right\rangle, $ $t\in [0,2].$

\vskip0.2cm
\noindent
\begin{minipage}{\textwidth}
	\centering
	\captionof{figure}{Región delimitada por $y^4 - 2y^2 \leq 4x \leq 2y^2$}
	\label{fig:solido_cilindro_elipse_3D}
	\begin{minipage}{0.34\linewidth}
		\centering
		\vspace{-15mm}
		\includegraphics[width=\linewidth]{Fig/436.pdf}
	\end{minipage}
	
	\vspace{-12mm}
	\fuentepropia
\end{minipage}
\vskip0.2cm


Una combinación convexa de la forma $(1-u)\overrightarrow{\alpha}(t) +u\overrightarrow{\beta}(t)$
parametriza la base del sólido, esto es
%\vskip0.2cm
{\small
\[\left\langle \frac{1}{4}t^2\left( \left( t^2-4\right) u+2\right) ,t\right\rangle\!.\]
}

La coordenada $z$ de un punto sobre el cilindro $z=4-x^2$ está dada por $z=4-\frac{t^{4}}{16}\left( \left( t^2-4\right) u+2\right) ^2.$
Por tanto, un punto $Q$ sobre el cilindro es de la forma
{\small $Q\left( \frac{1}{4}t^2((t^2-4)u+2),t,4-\frac{t^{4}}{16}\left( \left( t^2-4\right) u+2\right) ^2\right)\!,$}
con un punto en la base $P\left( \frac{1}{4}t^2\left( \left( t^2-4\right) u+2\right) ,t,0\right)\!,$
se forma la combinación convexa $(1-w)P+wQ$ que parametriza el sólido; esto define la siguiente función
\[
\mathbf{r} (t,u,w) = \left\langle \tfrac{t^2}{4}((t^2-4)u+2) ,t,w\left(4-\tfrac{1}{16}t^4((t^2-4)u+2)^2\right)\right\rangle\!;
\] \\
con $t\in [0,2],$ $u\in [0,1],$ $w\in [0,1].$
% luego hallamos los siguientes vectores. \begin{eqnarray*}
%{\bf r}_{t} &=& \langle t\left( \left( t^2-2\right) u+1\right) ,1,\frac{1}{2}wt^{3}\left( 4u-t^2u-2\right)\left(\left(t^2-2\right) u+1\right)\rangle \\
%{\bf r}_{u} &=& \langle \frac{1}{4}t^2\left( t^2-4\right) ,0,-\frac{1}{8}wt^{4}\left( t^2-4\right) \left( \left( t^2-4\right) u+2\right)\rangle \\
%{\bf r}_{w}&=& \langle 0,0,\frac{1}{16}\left( 4-t^2\right) \left(t^2\left( \left( t^2-4\right) u+2\right) +8\right) \left(t^2u+2\right) \rangle
%\end{eqnarray*}
El elemento diferencial de volumen $|{\bf r}_{t}\cdot ({\bf r}_{u}\times {\bf r}_{w})|$
es $\frac{t^2}{64}\left(t^2-4\right) ^2\left( t^2u+2\right) \left( t^2\left(
\left( t^2-4\right) u+2\right) \right)\!,$
término que integramos para conseguir el volumen del sólido
\[\int_{0}^{1}\int_{0}^{1}\int_{0}^2 |{\bf r}_{t}\cdot ({\bf r}_{u}\times {\bf r}_{w})| dtdudw= \frac{169472}{45045}.\]
%\end{tcolorbox}

Lo anterior se puede evaluar ejecutando las siguientes instrucciones:

\vspace{0.6em}
\noindent
\begin{tcolorbox}[breakable,enhanced,title=\bf Instrucción en \verb"Mathematica", top=2pt,bottom=2pt,boxsep=0pt,before skip=2pt,after skip=0pt]
\begin{Verbatim}[formatcom=\letraverbatim]
{x,y,z}=(1-u){t^2/2,t,0}+u{(t^4-2*t^2)/4,t,0};
r=Simplify[(1-w)*{x,y,0}+w*{x,y,4-x^2}]
A=Simplify[Dot[D[r,t],Cross[D[r,u],D[r,w]]]]
Integrate[A,{t,0,2},{u,0,1},{w,0,1}]
\end{Verbatim}
\end{tcolorbox}
\vspace{0.7em}
\noindent


\begin{ejer}
Plantear y evaluar las integrales triples en los órdenes $dxdydz$ y $dydxdz$,
necesarias para hallar el volumen del sólido acotado por las figuras de las expresiones $z = 4- x^2,$
$2x=y^2\,$ y $\,\, 4x = y^4 -2y^2$ para $y \geq 0.$
\end{ejer}

\subsection*{Tres cilindros rectos}

Los cilindros con ecuaciones
$x^2+y^2=2x$, $x^2+y^2=4x$ y $x^2+z^2=16$ generan el siguiente sólido.

\begin{figure}[H]
\centering
\caption{Visualización del sólido definido por $2x \leq x^2 + y^2 \leq 4x$ \quad y \quad $x^2 + z^2 \leq 16$.}
\label{fig:solido_anillo_cilindro}
\includegraphics[height=4.5cm]{Img/83a.pdf}\hskip0.25cm
\includegraphics[height=4.5cm]{Img/83b.pdf}\hskip0.25cm
\includegraphics[height=4.5cm]{Img/83c.pdf}
\parbox{\linewidth}{\centering\fontsize{9pt}{10pt}\selectfont\textbf{Fuente:} elaboración propia.}
\end{figure}

Estas figuras se consiguen ejecutando las siguientes instrucciones, se incluye una animación que muestra el recorrido de estas curvas en el plano $xy;$ esto se usa para obtener la parametrización de esta región.


\vspace{0.6em}
\noindent
\begin{tcolorbox}[breakable,enhanced,title=\bf Instrucción en \verb"Mathematica", top=2pt,bottom=2pt,boxsep=0pt,before skip=2pt,after skip=0pt]
\begin{Verbatim}[formatcom=\letraverbatim]
Manipulate[ParametricPlot[{{Cos[t]*Cos[t],
Cos[t]*Sin[t]},{2*Cos[t]*Cos[t],2*Cos[t]*Sin[t]}},
{t,-Pi/2,k},PlotRange->{{0,2},{-2,2}}],
{k,-Pi/2+Pi/40,Pi/2}]
ContourPlot3D[{16==x^2+z^2,x^2+y^2==2*x,x^2+y^2==4*x},
{x,-4,4},{y,-4,4},{z,-4,4},PlotPoints->90,Mesh->False,
Ticks->{{-4,-2,0,2,4},{-3,0,3},{-4,-2,0,2,4}},
ContourStyle->{Orange,LightGray,Cyan}]
RegionPlot3D[x^2+z^2<16&&x^2+y^2>2*x&&x^2+y^2<4*x&&0<
x<4,{x,0,4},{y,-2,2},{z,-4,4},Mesh->False]
\end{Verbatim}
\end{tcolorbox}
\vspace{0.7em}
\noindent


Con la proyección del sólido en el plano $xy$ se hallan los límites de integración en los
órdenes $dzdydx$ y $dzdxdy$; se muestra esta región, las diferentes tonalidades de sombreado determinan las integrales necesarias para su evaluación.
\vskip0.2cm


\noindent
\begin{minipage}{\textwidth}
	\centering
	\captionof{figure}{Región delimitada por $2x \leq x^2 + y^2 \leq 4x$.}
	\label{fig:region_corona_desplazada}
	\begin{minipage}{0.30\linewidth}
		\centering
		\vspace{-10mm}
		\includegraphics[width=\linewidth]{Fig/437.pdf}
	\end{minipage}
	\hspace{1em}
	\begin{minipage}{0.30\linewidth}
		\centering
		\vspace{-10mm}
		\includegraphics[width=\linewidth]{Fig/438.pdf}
	\end{minipage}
	
	\vspace{-8mm}
	\parbox{\linewidth}{\centering\fontsize{9pt}{10pt}\selectfont\textbf{Fuente:} elaboración propia.}
\end{minipage}
\vskip0.2cm

En coordenadas rectangulares, el volumen del mencionado sólido se puede determinar por medio de las siguientes integrales.
{\small
\begin{multline*}
\int_{-1}^{1}\int_{2-\sqrt{4-y^2}}^{1-\sqrt{1-y^2}}\int_{-\sqrt{16-x^2}}^{\sqrt{16-x^2}}dzdxdy+\int_{-1}^{1}\int_{1+\sqrt{1-y^2}}^{2+\sqrt{4-y^2}}\int_{-\sqrt{16-x^2}}^{\sqrt{16-x^2}}dzdxdy+\\
\int_{-2}^{-1}\int_{2-\sqrt{4-y^2}}^{2+\sqrt{4-y^2}}\int_{-\sqrt{16-x^2}}^{\sqrt{16-x^2}}dzdxdy+\int_{1}^2\int_{2-\sqrt{4-y^2}}^{2+\sqrt{4-y^2}}\int_{-\sqrt{16-x^2}}^{\sqrt{16-x^2}}dzdxdy \approx 57.207.
\end{multline*}
}
\vspace{-5mm}
{\small
\begin{multline*}
\int_{0}^2\int_{-\sqrt{4x-x^2}}^{-\sqrt{2x-x^2}}\int_{-\sqrt{16-x^2}}^{\sqrt{16-x^2}}dzdydx+\int_{0}^2\int_{\sqrt{2x-x^2}}^{\sqrt{4x-x^2}}\int_{-\sqrt{16-x^2}}^{\sqrt{16-x^2}}dzdydx\\
+\int_{2}^{4}\int_{-\sqrt{4x-x^2}}^{\sqrt{4x-x^2}}\int_{-\sqrt{16-x^2}}^{\sqrt{16-x^2}}dzdydx \approx 57.207
\end{multline*}
}

Estas integrales se evaluaron numéricamente en \verb"Mathematica";
se presentan las instrucciones para el orden $dzdydx$.


\vspace{0.6em}
\noindent
\begin{tcolorbox}[breakable,enhanced,title=\bf Instrucción en \verb"Mathematica", top=2pt,bottom=2pt,boxsep=0pt,before skip=2pt,after skip=0pt]
\begin{Verbatim}[formatcom=\letraverbatim]
NIntegrate[1,{x,0,2},{y,-Sqrt[4*x-x^2],-Sqrt[2*x-x^2]},
{z,-Sqrt[16-x^2],Sqrt[16-x^2]}]+NIntegrate[1,{x,0,2},
{y,Sqrt[2*x-x^2],Sqrt[4*x-x^2]},{z,-Sqrt[16-x^2],
Sqrt[16-x^2]}]+NIntegrate[1,{x,2,4},{y,-Sqrt[4*x-x^2],
Sqrt[4*x-x^2]},{z,-Sqrt[16-x^2],Sqrt[16-x^2]}]
\end{Verbatim}
\end{tcolorbox}
\vspace{0.7em}
\noindent

Haciendo un cambio a coordenadas cilíndricas el volumen esta dado por:
%\vskip0.2cm 
\[
\int_{0}^{\pi} \int_{2 \cos \theta}^{4 \cos \theta} \int_{-\sqrt{16-(r \cos \theta)^{2}}}^{\sqrt{16-(r \cos \theta)^{2}}} r\,dz\,dr\,d\theta \approx 57.207.
\]
%\vskip0.2cm


El mismo valor se puede obtener con el siguiente procedimiento.


%\begin{tcolorbox}[title=\bfseries Parametrizando el sólido]
\noindent
En el plano $xy$, un punto $P$ sobre el círculo interno tiene la forma
{\small
\[\left( 2\cos (\theta) \cos (\theta) ,2\cos (\theta) \sin (\theta) \right) = (2 \cos^2 (\theta), \sin(2\theta)),\]}
y un punto $Q$ sobre el círculo mayor posee la forma
{\small
\[\left( 4\cos (\theta) \cos (\theta) ,4\cos \left(
\theta \right) \sin (\theta) \right) = (4 \cos^2 (\theta), 2\sin(2\theta)).\]}

%\vspace{1ex} 

%\noindent
%\vskip0.4cm
\vspace{-2mm}
La combinación convexa $\left( 1-u\right)P +uQ$ se simplifica como
$ (2( 1+u)\cos^2(\theta), \linebreak
( 1+u)\sin (2\theta) ),$ y parametriza la región acotada por estas dos curvas. Sobre el cilindro $x^2+z^2 = 16$ se satisface que $z = \pm 2\sqrt{4-(1+u) ^2\cos ^{4}(\theta) }.$ Sobre el cilindro circular, consideremos dos puntos $P_1$ y $Q_1.$


\vskip0.3cm
\noindent
\begin{minipage}{\textwidth}
	\centering
	\captionof{figure}{Región delimitada por $2x \leq x^2 + y^2 \leq 4x$}
	\label{fig:corona_desplazada_xy}
	\begin{minipage}{0.37\linewidth}
		\centering
		\vspace{-15mm}
		\includegraphics[width=\linewidth]{Fig/439.pdf}
	\end{minipage}
	
	\vspace{-10mm}
	\fuentepropia
\end{minipage}
\vskip0.3cm

{\small
\[P_1= \big(2(1+u) \cos^2(\theta),\left( 1+u\right) \sin \left( 2\theta \right) ,-2\sqrt{4-(1+u)^2\cos ^{4}(\theta) }\big)\]  \[Q_1=\big(2(1+u) \cos ^2(\theta) ,\left( 1+u\right) \sin \left( 2\theta \right) ,2\sqrt{4-\left( 1+u\right) ^2\cos ^{4}(\theta) }\big);\]
}
%\vspace{7mm}
simplificando la expresión
$\left(1-w\right)P_1+wQ_1 $ se obtiene la función.
%\vspace{-4mm}
{\small
\[\mathbf{r}(\theta,u,w) = \left\langle 2(1+u)\cos^2\theta, (1+u)\sin(2\theta), 2(2w-1)\sqrt{4-(1+u)^2\cos^4\theta}\right\rangle\]
}
que parametriza el sólido, con $\theta \in [0,2\pi],$ $u \in [0,1]$ y $w \in [0,1].$
%% ${\bf r}_{\theta } = \left\langle -4\cos (\theta) \left(1+u\right) \sin (\theta) ,2\left( 1+u\right) \cos \left(2\theta \right) ,\frac{4\left( 2w-1\right) \left( 1+u\right) ^2\cos^{3}(\theta) \sin (\theta) }{\sqrt{4-\left(1+u\right) ^2\cos ^{4}(\theta) }}\right\rangle $\\
%%${\bf r}_{u} = \left\langle 2\cos ^2(\theta),\sin \left( 2\theta \right) ,\frac{-2\left( 2w-1\right) \left( 1+u\right)\cos ^{4}(\theta) }{\sqrt{4-\left( 1+u\right) ^2\cos^{4}(\theta) }} \right\rangle $ \\
%%${\bf r}_{w} = \left\langle 0,0,4\sqrt{\left( 4-\left( 1+u\right)^2\cos ^{4}(\theta) \right) }\right\rangle .$ \\
Con esta función se obtiene el elemento diferencial de volumen
%\vskip0.2cm
{\small
\[\left| {\bf r}_{w}\cdot \left({\bf r}_{\theta }\times {\bf r}_{u}\right) \right| =16\left| \cos ^2(\theta) \left( 1+u\right) \sqrt{4-\left( 1+u\right) ^2\cos ^{4}(\theta) }\right|,\]
}
%\vskip0.2cm
cuya integral es el volumen del sólido que se considera.
{\small
\[\int_{0}^{\pi }\int_{0}^{1}\int_{0}^{1}\left| {\bf r}_{w}\cdot \left({\bf r}_{\theta }\times {\bf r}_{u}\right) \right| dudwd\theta \approx 57.207.\]}
%\end{tcolorbox}
%\vskip0.2cm
Enseguida se presenta una forma de hacer parametrización de la base del sólido y los cálculos que permiten evaluar numéricamente la integral.


\vspace{0.6em}
\noindent
\begin{tcolorbox}[breakable,enhanced,title=\bf Instrucción en \verb"Mathematica", top=2pt,bottom=2pt,boxsep=0pt,before skip=2pt,after skip=0pt]
\begin{Verbatim}[formatcom=\letraverbatim]
{x,y}=Simplify[(1-u)*{2*Cos[t]*Cos[t],2*Cos[t]*Sin[t]}
+u*{4*Cos[t]*Cos[t],4*Cos[t]*Sin[t]}];
z=Simplify[(1-w)*Sqrt[16-x^2]+w*(-Sqrt[16-x^2])];
A=Dot[D[{x,y,z},t],Cross[D[{x,y,z},u],D[{x,y,z},w]]]
NIntegrate[A,{t,0,Pi},{u,0,1},{w,0,1}]
\end{Verbatim}
\end{tcolorbox}
\vspace{0.7em}
\noindent


Haciendo un cambio a coordenadas cilíndricas, el volumen está dado por:
\vspace{-2mm}
\[\int_{0}^{\pi }\int_{2\cos \theta }^{4\cos \theta }\int_{-\sqrt{16-\left(r\cos \theta \right)^2}}^{\sqrt{16-(r\cos\theta)^2}}rdzdrd\theta \approx 57.207.\]


\subsection*{Tres cilindros}
Las figuras de las ecuaciones $y^2 +2z^2=4,$ $2x=y^2$ y $4x=y^4-2y^2$ se muestran a continuación:


\vskip0.2cm
\noindent
\begin{minipage}{\textwidth}
	\centering
	\captionof{figure}{Tres cilindros}
	\label{fig:tres_cilindros}
	\begin{minipage}{0.20\linewidth}
		\centering
		\includegraphics[width=\linewidth]{Img/18d.pdf}
	\end{minipage}
	\begin{minipage}{0.20\linewidth}
		\centering
		\includegraphics[width=\linewidth]{Img/18c.pdf}
	\end{minipage}
	\begin{minipage}{0.20\linewidth}
		\centering
		\includegraphics[width=\linewidth]{Img/18e.pdf}
	\end{minipage}
	
	\vspace{0.3em}
	\fuentepropia
\end{minipage}
\vskip0.2cm



Estas superficies con $y\geq 0$ determinan la región que se muestra enseguida.

\vskip0.2cm
\noindent
\begin{minipage}{\textwidth}
	\centering
	\captionof{figure}{Sólido determinado por los tres cilindros.}
	\begin{minipage}{0.18\linewidth}
		\centering
		\includegraphics[width=\linewidth]{Img/18a.pdf}
	\end{minipage}
	\hspace{1em}
	\begin{minipage}{0.18\linewidth}
		\centering
		\includegraphics[width=\linewidth]{Img/18b.pdf}
	\end{minipage}
	
	\vspace{0.3em}
	\fuentepropia
\end{minipage}

\vskip0.2cm
Estas figuras se obtienen con ayuda de \verb"Mathematica", ejecutando los siguientes comandos.



\vspace{0.6em}
\noindent
\begin{tcolorbox}[breakable,enhanced,title=\bf Instrucción en \verb"Mathematica", top=2pt,bottom=2pt,boxsep=0pt,before skip=2pt,after skip=0pt]
\begin{Verbatim}[formatcom=\letraverbatim]
ContourPlot3D[{y^2+2*z^2==4,2*x==y^2,4*x==y^4-2*y^2},
{x,-1,2},{y,0,2},{z,-2,2},,PlotPoints->60,Mesh->False,
ContourStyle->{LightPurple,Orange,Blend[{Green,
Orange,Red}]}]
RegionPlot3D[-Sqrt[4-y^2]<z<Sqrt[4-y^2]&&(y^4-2*y^2)/4
<x<y^2/2,{x,-1,2},{y,0,2},{z,-2,2},PlotPoints->90,
Mesh->None,PlotStyle->Blend[{Green,Orange,Red}]]
\end{Verbatim}
\end{tcolorbox}
\vspace{0.7em}
\noindent


Con las siguientes tres instrucciones además de la figura de la superficie se obtienen las curvas de intersección.


\vspace{0.6em}
\noindent
\begin{tcolorbox}[breakable,enhanced,title=\bf Instrucción en \verb"Mathematica", top=2pt,bottom=2pt,boxsep=0pt,before skip=2pt,after skip=0pt]
\begin{Verbatim}[formatcom=\letraverbatim]
a=ContourPlot3D[{y^2+2*z^2==4,2*x==y^2,4*x==y^4
2*y^2},{x,-0.5,2},{y,0,2},{z,-2,2},PlotPoints->80,
Mesh->False,ContourStyle->{Directive[LightPurple,
Opacity[0.9]],Directive[Orange,Opacity[0.75]],
Directive[Blend[{Green,Orange,Red}],Opacity[0.75]],
Opacity[0.75]}];
b=ParametricPlot3D[{{2-t^2,Sqrt[4-2*t^2],t},
{(t^2-1)*(t^2-2),Sqrt[4-2*t^2],t}},{t,-Sqrt[2],
Sqrt[2]},PlotPoints->80,PlotStyle->{Directive[Blue,
Thickness->0.012],Directive[Red,Thickness->0.012]}];
Show[a,b]
\end{Verbatim}
\end{tcolorbox}
\vspace{0.7em}
\noindent


La expresión $y^2 +2z^2 = 4$ equivale a $y = \pm \sqrt{4-2z^2}$ o $z = \pm \sqrt{\frac{4-y^2}{2}} $ y de la expresión $4x = y^4 -2y^2$ podemos plantear cuatro funciones con $y$ como variable dependiente de $x$, esto es $y = \pm \sqrt{1\pm \sqrt{1+4x}}.$ \vskip0.3cm
En los planos $yz$ y $xy$ el sólido se proyecta como se muestra a continuación.


\vskip0.3cm
\noindent
\begin{minipage}{\textwidth}
	\centering
	\captionof{figure}{Proyección del sólido en los planos $yz$ y $xy$.}
	\label{fig:proyeccion_yz_xy}
	\begin{minipage}{0.35\linewidth}
		\centering
		\vspace{-15mm}
		\includegraphics[width=\linewidth]{Fig/440.pdf}
	\end{minipage}
	\hspace{1em}
	\begin{minipage}{0.35\linewidth}
		\centering
		\vspace{-15mm}
		\includegraphics[width=\linewidth]{Fig/441.pdf}
	\end{minipage}	
	
	\vspace{-10mm}
	\fuentepropia
\end{minipage}
\vskip0.2cm


Con las figuras anteriores planteamos las integrales triples que determinan el volumen del sólido, integrando primero con respecto a $x$ o $z$.
\vspace{-1mm}
{\small
\begin{multline*}
\hskip-0.3cm \bigintsss_{-\frac{1}{4}}^{0}\bigintsss_{\sqrt{1-\sqrt{1+4x}}}^{\sqrt{1+\sqrt{1+4x}}}\bigintsss_{-\sqrt{\frac{4-y^2}{2}}}^{\sqrt{\frac{4-y^2}{2}}}dzdydx+\bigintsss_{0}^2\bigintsss_{\sqrt{2x}}^{\sqrt{1+\sqrt{1+4x}}}\bigintsss_{-\sqrt{\frac{4-y^2}{2}}}^{\sqrt{\frac{4-y^2}{2}}}dzdydx\\
=\bigintsss_{-\sqrt{2}}^{\sqrt{2}}\bigintsss_{0}^{\sqrt{4-2z^2}}\bigintsss_{\frac{y^{4}-2y^2}{4}}^{\frac{y^2}{2}}dxdydz=\bigintsss_{0}^2\bigintsss_{-\sqrt{\frac{4-y^2}{2}}}^{\sqrt{\frac{4-y^2}{2}}}\bigintsss_{\frac{y^{4}-2y^2}{4}}^{\frac{y^2}{2}}dxdzdy \\
= \bigintsss_{0}^2\bigintsss_{\frac{y^{4}-2y^2}{4}}^{ \frac{y^2}{2}}\bigintsss_{-\sqrt{\frac{4-y^2}{2}}}^{\sqrt{\frac{4-y^2}{2}}}dzdxdy= \frac{\sqrt{2}\pi}{2}.
\end{multline*}
}

Dos de estas integrales se obtienen con ayuda de las siguientes sentencias.


\vspace{0.6em}
\noindent
\begin{tcolorbox}[breakable,enhanced,title=\bf Instrucción en \verb"Mathematica", top=2pt,bottom=2pt,boxsep=0pt,before skip=2pt,after skip=0pt]
\begin{Verbatim}[formatcom=\letraverbatim]
Integrate[1,{z,-Sqrt[2],Sqrt[2]},{y,0,Sqrt[4-2z^2]},
{x,(y^4-2y^2)/4,y^2/2}]
Integrate[1,{y,0,2},{x,(y^4-2y^2)/4,y^2/2},
{z,-Sqrt[(4-y^2)/2],Sqrt[(4-y^2)/2]}]
\end{Verbatim}
\end{tcolorbox}
\vspace{0.7em}
\noindent


De la ecuación del cilindro se sigue que $y^2= 4-2z^2$, al reemplazar esta expresión en las otras dos ecuaciones se obtienen las expresiones que definen la proyección de la figura en el plano $xz$. Esto es, para $2x = y^2 \to x = 2- z^2$ y de la ecuación
$4x=y^{4}-2y^2$ se sigue que $4x=\left( 4-2z^2\right) ^2-2\left( 4-2z^2\right),$ entonces
$x= (z^2-1)(z^2-2).$
\vskip0.2cm
Para integrar en el orden $dydxdz$ nos apoyamos en la siguiente figura.
Debemos tener en cuenta que en una región, $y$ recorre los puntos desde la curva con ecuación
$y = a $ hasta $y =b;$
en la otra región $y$ recorre puntos desde $y=\sqrt{2x}$ hasta
$y = \sqrt{4 -2z^2}$ y en la tercer región $y$ toma valores desde $y = \sqrt{2x}$ hasta
$y =b.$
\vskip0.4cm
Por la simetría de la figura con respecto al plano $xy$, la segunda gráfica de la figura \ref{fig:proyeccion_xzx} de las siguientes figuras se ha dividido en seis regiones que corresponden, respectivamente, a las seis integrales que se enuncian; luego de plantearlas basta multiplicar por $2.$

\vskip0.3cm
\noindent
\begin{minipage}{\textwidth}
	\centering
	\vspace{3mm}
	\captionof{figure}{Proyección del sólido en el plano $xz$.}
	\label{fig:proyeccion_xzx}
	\begin{minipage}{0.45\linewidth}
		\centering
		\vspace{-15mm}
		\includegraphics[width=\linewidth]{Fig/442.pdf}
	\end{minipage}
	\hspace{1em}
	\begin{minipage}{0.40\linewidth}
		\centering
		\vspace{-15mm}
		\includegraphics[width=\linewidth]{Fig/443.pdf}
	\end{minipage}
	
	\vspace{-12mm}
	\fuentepropia
\end{minipage}
\vskip0.5cm

Las regiones se enumeran en el mismo orden en que se presentan las siguien\-tes integrales, para abreviar usaremos la notación {\small $a= \sqrt{1-\sqrt{1+4x}}$} y {\small $b=\sqrt{1+\sqrt{1+4x}}.$}


{\small
\begin{multline*}
\hskip-0.3cm \bigintsss_{1}^{\sqrt{2}}\bigintsss_{(z^2-1)(z^2-2)}^{0} \bigintsss_{ a }^{\sqrt{4-2z^2}}dydxdz
+\bigintsss_{1}^{\sqrt{\frac{3}{2}}} \bigintsss_{-\frac{1}{4}}^{(z^2-1)(z^2-2)}\bigintsss_{a}^{b}dydxdz \\
+\bigintsss_{0}^{1}\bigintsss_{-\frac{1}{4}}^{0}\bigintsss_{ a }^{b}dydxdz+\bigintsss_{1}^{\sqrt{2}}\bigintsss_{0}^{2-z^2}\bigintsss_{\sqrt{2x}}^{\sqrt{4-2z^2}}dydxdz + \\
\hskip-0.3cm \bigintsss_{0}^{1}\bigintsss_{(z^2-1)(z^2-2)}^{2-z^2} \bigintsss_{\sqrt{2x}}^{\sqrt{4-2z^2}}dydxdz
+\bigintsss_{0}^{1}\bigintsss_{0}^{(z^2-1)(z^2-2)} \bigintsss_{\sqrt{2x}}^{b}dydxdz \approx \frac{\pi}{\sqrt{2}}.
\end{multline*}
}

\vskip0.2cm
Las integrales anteriores se pueden evaluar numéricamente mediante la ejecución de las siguientes instrucciones:


\vspace{0.6em}
\noindent
\begin{tcolorbox}[breakable,enhanced,title=\bf Instrucción en \verb"Mathematica", top=2pt,bottom=2pt,boxsep=0pt,before skip=2pt,after skip=0pt]
\begin{Verbatim}[formatcom=\letraverbatim]
2*(NIntegrate[1,{z,1,Sqrt[2]},{x,(z^2-1)*(z^2-2),0},
{y,Sqrt[1-Sqrt[1+4*x]],Sqrt[4-2*z^2]}]+NIntegrate[1,
{z,1,Sqrt[3/2]},{x,-1/4,(z^2-1)*(z^2-2)},{y,Sqrt[1
Sqrt[1+4*x]],Sqrt[1+Sqrt[1+4*x]]}]+NIntegrate[1,
{z,0,1},{x,-1/4,0},{y,Sqrt[1-Sqrt[1+4*x]],Sqrt[1
+Sqrt[1+4*x]]}]+Integrate[1,{z,1,Sqrt[2]},{x,0,
2-z^2},{y,Sqrt[2*x],Sqrt[4-2*z^2]}]+NIntegrate[1,
{z,0,1},{x,(z^2-1)*(z^2-2),2-z^2},{y,Sqrt[2*x],
Sqrt[4-2*z^2]}]+NIntegrate[1,{z,0,1},{x,0,
(2-z^2)(1-z^2)},{y,Sqrt[2*x],Sqrt[1+Sqrt[1+4*x]]}])
\end{Verbatim}
\end{tcolorbox}
\vspace{0.7em}
\noindent

Para plantear la integral triple en el orden $dydzdx,$ hay que obtener $z$ de la ecuación $x = (z^2-1)(z^2-2),$ esto es
$z = \pm \frac{1}{2} \sqrt{6 \pm 2 \sqrt{1+4x}};$
el signo depende de la región del plano que se considere. Usando la simetría de la figura, las integrales triples que se requieren evaluar para determinar el volumen del sólido son:

{\footnotesize
\begin{multline} \label{trescilindros1}
2\Bigg( \bigintsss_{-\frac{1}{4}}^{0} \bigintsss_{0}^{\frac{1}{2}\sqrt{6-2\sqrt{1+4x}}}\bigintsss_{ a }^{b}dydzdx+\\
\bigintsss_{-\frac{1}{4}}^0 \bigintsss_{\frac{1}{2}\sqrt{6-2\sqrt{1+4x}}}^{\frac{1}{2}\sqrt{6+2\sqrt{1+4x}}}\bigintsss_{ a }^{\sqrt{4-2z^2}}dydzdx+ \bigintsss_0^2 \bigintsss_{\frac{1}{2}\sqrt{6-2\sqrt{1+4x}}}^{\sqrt{2-x}}\bigintsss_{\sqrt{2x}}^{\sqrt{4-2z^2}}dydzdx
\\ + \bigintsss_{0}^2\bigintsss_{0}^{\frac{1}{2}\sqrt{6-2\sqrt{1+4x}}} \bigintsss_{\sqrt{2x}}^{b}dydzdx\Bigg) \approx \frac{\pi}{\sqrt{2}} \end{multline}
}

\vskip0.2cm
El volumen también se puede determinar de la siguiente manera.

%\begin{tcolorbox}[title=\bfseries Parametrizando el sólido]
En el plano $xy$ las curvas que determinan la proyección del sólido se parametrizan como
{\small $\overrightarrow{\alpha }(t)=\left\langle \frac{1}{4}\left(
t^{4}-2t^2\right) ,t\right\rangle$} y {\small $\overrightarrow{\beta }(t)=\left\langle \frac{1}{2}t^2,t\right\rangle;$} con $t\in [0,2].$
\vskip0.2cm
Con la expresión {\small $\left( 1-u\right) \overrightarrow{\alpha }(t)+u\overrightarrow{\beta }(t),$}
se parametriza la región entre estas curvas, esto es {\small $\left\langle \frac{1}{4}t^2\left( \left(1-u\right)t^2-2+4u\right) ,t\right\rangle.$}
La expresión {\small $z=\pm \sqrt{\frac{4-y^2}{2}},$} se reescribe como {\small $z=\pm \sqrt{2-\frac{1}{2}t^2}.$}
Dos puntos sobre el cilindro $y^2+2z^2=4$ tienen la forma
\[P\left( \frac{1}{4}t^2\left( \left( 1-u\right) t^2-2+4u\right) ,t,-\sqrt{2-\frac{1}{2}t^2}\right)\]
y \[Q\left( \frac{1}{4}t^2\left( \left(1-u\right) t^2-2+4u\right) ,t,\sqrt{2-\frac{1}{2}t^2}\right).\]
El sólido se parametriza simplificando la expresión $(1-w)P+wQ,$ que determina la siguiente función
\vskip0.2cm
\begin{equation*}
{\bf r}\left( t,u,w\right) = \left\langle \frac{1}{4}t^2\left( \left( 1-u\right)
t^2-2+4u\right) ,t,\left( w-\frac{1}{2}\right) \sqrt{8-2t^2}\right\rangle,
\end{equation*}
\vskip0.2cm

en donde $t\in [0,2],$ $u\in [0,1]$ y $w\in [0,1]$. %\vskip0.2cm
% Derivando parcialmente, hallamos los siguientes vectores
%% \begin{eqnarray*} {\bf r}_{t} &=& \left\langle \frac{1}{2}t\left( \left(1-u\right) t^2-2+4u\right) +\frac{1}{2}t^{3}\left( 1-u\right) ,1,\frac{t\left( 1-2w\right) }{\sqrt{8-2t^2}} \right\rangle \\ {\bf r}_{u} &=& \left\langle \frac{1}{4}t^2\left( 4-t^2\right) ,0,0\right\rangle \\
%%{\bf r}_{w}&=& \left\langle 0,0,\sqrt{8-2t^2}\right\rangle \end{eqnarray*}
Con esta función se obtiene el elemento diferencial de volumen
%\vskip0.2cm
\[\left| {\bf r}_{t}\cdot \left({\bf r}_{u}\times {\bf r}_{w}\right) \right| =\left| -\frac{1}{4}t^2\left( 4-t^2\right) \sqrt{8-2t^2}\right|,\]
%\vskip0.2cm 
cuya integración proporciona el volumen del sólido
%\vskip0.2cm
\[\int_{0}^{1}\int_{0}^{1}\int_{0}^2\left| -\frac{1}{4}t^2\left( 4-t^2\right) \sqrt{8-2t^2}\right|dtdudw= \frac{1}{2}\sqrt{2}\pi.\]
%\vskip0.2cm
%\end{tcolorbox}

Estos cálculos se realizan en \verb"Mathematica" con las siguientes instrucciones:


\vspace{0.6em}
\noindent
\begin{tcolorbox}[breakable,enhanced,title=\bf Instrucción en \verb"Mathematica", top=2pt,bottom=2pt,boxsep=0pt,before skip=2pt,after skip=0pt]
\begin{Verbatim}[formatcom=\letraverbatim]
r={x,y}=(1-u)*{(t^4-2*t^2)/4,t}+u*{t^2/2,t};
s=(1-w)*{x,y,-Sqrt[2-y^2/2]}+w*{x,y,Sqrt[2-y^2/2]};
A=Dot[D[s,t],Cross[D[s,u],D[s,w]]];
Integrate[Abs[A],{t,0,2},{u,0,1},{w,0,1}]
\end{Verbatim}
\end{tcolorbox}
\vspace{0.7em}
\noindent


Las integrales de la expresión (\ref{trescilindros1}) se evalúan ejecutando las siguientes instrucciones:


\vspace{0.6em}
\noindent
\begin{tcolorbox}[breakable,enhanced,title=\bf Instrucción en \verb"Mathematica", top=2pt,bottom=2pt,boxsep=0pt,before skip=2pt,after skip=0pt]
\begin{Verbatim}[formatcom=\letraverbatim]
(NIntegrate[1,{x,-1/4,0},{z,0,Sqrt[6-2Sqrt[1+4x]]/2},
{y,Sqrt[1-Sqrt[1+4x]],Sqrt[1+Sqrt[1+4x]]}]+
NIntegrate[1,{x,-1/4,0},{z,Sqrt[6-2Sqrt[1+4x]]/2,
Sqrt[6+2Sqrt[1+4x]]/2},{y,Sqrt[1-Sqrt[1+4x]],
Sqrt[4-2z^2]}]+NIntegrate[1,{x,0,2},{z,
Sqrt[6-2Sqrt[1+4x]]/2,Sqrt[2-x]},{y,Sqrt[2x],
Sqrt[4-2z^2]}]+NIntegrate[1,{x,0,2},{z,0,Sqrt[6-
2Sqrt[1+4x]]/2},{y,Sqrt[2x],Sqrt[1+Sqrt[1+4x]]}])*2
\end{Verbatim}
\end{tcolorbox}
\vspace{0.7em}
\noindent


\subsection*{Un cilindro y dos planos}

El cilindro elíptico con ecuaciones $4x^2+z^2=4$ y los planos $z=3+2y$ y
$z=5-x-3y$ acotan el sólido que se muestra a continuación:

\vskip0.2cm

\noindent
\begin{minipage}{\textwidth}
	\centering
	\captionof{figure}{Sólido definido por $4x^2 + z^2 \leq 4$ \quad y \quad $\dfrac{z - 3}{2} \leq y \leq \dfrac{5 - x - z}{3}$.}
	\label{fig:solido_cilindro_planos}
	\begin{minipage}{0.20\linewidth}
		\centering
		\includegraphics[width=\linewidth]{Img/55a.pdf}
	\end{minipage}
	\hspace{1em}
	\begin{minipage}{0.20\linewidth}
		\centering
		\includegraphics[width=\linewidth]{Img/55b.pdf}
	\end{minipage}
	\hspace{1em}
	\begin{minipage}{0.20\linewidth}
		\centering
		\includegraphics[width=\linewidth]{Img/55c.pdf}
	\end{minipage}
	
	\vspace{0.3em}
	\fuentepropia
\end{minipage}
\vskip0.2cm


Con las siguientes instrucciones se consiguen las figuras presentadas.

%\vskip0.4cm
\vspace{0.6em}
\noindent
\begin{tcolorbox}[breakable,enhanced,title=\bf Instrucción en \verb"Mathematica", top=2pt,bottom=2pt,boxsep=0pt,before skip=2pt,after skip=0pt]
\begin{Verbatim}[formatcom=\letraverbatim]
ContourPlot3D[{4*x^2+z^2==4,z==3+2*y,z==5-x-3*y},
{x,-2,2},{y,-3,4},{z,-2,2.5},Mesh->False,
ContourStyle->{Orange,RGBColor[0.2,0.6,0.4],Cyan}]
RegionPlot3D[(z-3)/2<y<(5-x-z)/5&&-Sqrt[1-z^2/4]<x<
Sqrt[1-z^2/4]&&-2<x<2,{x,-1,1},{y,-3,2},{z,-2,2},
Mesh->False,PlotStyle->{Blend[{Orange,Yellow},2/3]},
PlotPoints->90,Ticks->{{-1,0},{-2,0,1},{-1,0,1,2}}]
\end{Verbatim}
\end{tcolorbox}
\vspace{0.7em}
\noindent



En el plano $xz$ el sólido se proyecta como una elipse y los límites de integración para $y,$ van desde un plano al otro.
\vskip0.2cm
La proyección en el plano $xy,$ se determina reemplazando $z$ en la ecuación del cilindro $4x^2+z^2=4,$ por el respectivo término de $z$ en las ecuaciones de los planos, de tal manera que la expresión resultante solo dependa de $x$ y $y.$

\noindent
\begin{minipage}{\textwidth}
	\centering
	\captionof{figure}{Proyecciones del sólido en los planos $xy$ y $xz$.}
	\label{fig:proyeccion_xy_xz}
	\begin{minipage}{0.40\linewidth}
		\centering
		\vspace{-17mm}
		\includegraphics[width=\linewidth]{Fig/444.pdf}
	\end{minipage}
	\hspace{1em}
	\begin{minipage}{0.35\linewidth}
		\centering
		\vspace{-17mm}
		\includegraphics[width=\linewidth]{Fig/445.pdf}
	\end{minipage}
	
	\vspace{-12mm}
	\fuentepropia
\end{minipage}
\vskip0.2cm


Con lo anterior podemos plantear las integrales que determinan el volumen del sólido en los órdenes $dydxdz$ y $dydzdx,$
{\small
\begin{multline*} \bigintss_{-1}^{1}\bigintss_{-\sqrt{4-4x^2}}^{\sqrt{4-4x^2}}\bigintss_{\frac{z-3}{2}}^{\frac{5-x-z}{3}}dydzdx \\ =\bigintss_{-2}^2\bigintss_{-\sqrt{1-z^2/4}}^{\sqrt{1-z^2/4}}\bigintss_{\frac{z-3}{2}}^{\frac{5-x-z}{3}}dydxdz= \frac{19}{3}\pi.
\end{multline*}
}
%\vskip0.2cm
Con la ecuación del plano $z =3+2y$ se reemplaza $z$ en la ecuación del cilindro y se sigue que $4x^2+(3+2y)^2=4,$ que equivale a $4x^2+4(y+\frac{3}{2})^2=4;$
esta ecuación representa una circunferencia centrada en el punto $(0,-\frac{3}{2})$ de radio $1.$
Ahora, para el plano $z = 5-x-3y,$ la sustitución nos lleva a la ecuación
$4x^2 + (5-x-3y)^2 = 4,$ que corresponde a una elipse rotada en el plano $xy.$ De la expresión $4x^2+( 5-x-3y) ^2=4$, se obtiene que
$x=1-\frac{3}{5}y\pm \frac{2}{5}\sqrt{30y-20-9y^2},$
también
$y=\frac{5}{3}-\frac{1}{3}x\pm \frac{2}{3}\sqrt{1-x^2}.$
%\vskip0.4cm
De la expresión $4x^2+4( y+\frac{3}{2}) ^2=4$,
se obtiene que $x=\pm \frac{1}{2}\sqrt{ 4-( 3+2y) ^2}$
$y=-\frac{3}{2}\pm \sqrt{1-x^2}.$
Es de anotar, que en los extremos los valores de $z$ son menores en el cilindro que en los planos.
Con la información anterior podemos plantear las integrales que determinan el volumen del sólido en mención, en el orden $dzdydx$.

\vspace{-6mm}
{\small
\begin{multline*}
\bigintss_{-1}^{1}\bigintss_{-\frac{3}{2}-\sqrt{1-x^2}}^{-\frac{3}{2}+\sqrt{1-x^2}}\bigintss_{-2\sqrt{1-x^2}}^{3+2y}dzdydx \\ +\bigintss_{-1}^{1}\bigintss_{-\frac{3}{2}+\sqrt{1-x^2}}^{\frac{5}{3}-\frac{1}{3}x-\frac{2}{3}\sqrt{1-x^2}}\bigintss_{-\sqrt{4-4x^2}}^{\sqrt{4-4x^2}}dzdydx \\
+\bigintss_{-1}^{1}\bigintss_{\frac{5}{3}-\frac{1}{3}x-\frac{2}{3}\sqrt{1-x^2}}^{\frac{5}{3}-\frac{1}{3}x+\frac{2}{3}\sqrt{%
1-x^2}}\bigintss_{-\sqrt{4-4x^2}}^{5-x-3y}dzdydx= \frac{19}{3}\pi.
\end{multline*}
}
%\vskip0.2cm
Sobre el plano $yz,$ se hacen varias consideraciones, pues la proyección del sólido en este plano incluye una elipse rotada; además, de la expresión $4(5-z-3y)^2+z^2=4$ se sigue que
{\small
\[z=-\frac{12}{5}y+4\pm \frac{2}{5}\sqrt{5-9\left( y-\frac{5}{3}\right) ^2}.\]
}
La integral en el orden $dxdzdy$ corresponde a la suma de cinco integrales, esto se debe a que $x$ recorre los puntos entre la parte negativa del cilindro y su parte positiva; sin embargo, en otra región $x$ va desde la parte negativa del cilindro al plano $z=5-x-3y;$ esta última región se proyecta en el plano $yz$ como una elipse rotada.


\vskip0.3cm
\noindent
\begin{minipage}{\textwidth}
	\centering
	\captionof{figure}{Proyección del sólido en el plano $yz$}
	\label{fig:proyeccion_yz}
	\begin{minipage}{0.75\linewidth}
		\centering
		\vspace{-55mm}
		\includegraphics[width=\linewidth]{Fig/446.pdf}
	\end{minipage}
	
	\vspace{-50mm}
	\fuentepropia
\end{minipage}
\vskip0.3cm


Con lo anterior, y considerando las abreviaturas.
$a(y)=-\frac{12}{5}y+4+\frac{2}{5}\sqrt{5-9\left( y-\frac{5}{3}\right) ^2}$ y $b(y)=-\frac{12}{5}y+4-\frac{2}{5}\sqrt{5-9\left( y-\frac{5}{3}\right) ^2},$ el planteamiento es el siguiente:

\vspace{-6mm}
{\small
\begin{multline*}
\bigintsss_{-\frac{5}{2}}^{-\frac{1}{2}}\bigintsss_{-2}^{3+2y}\bigintsss_{-\sqrt{1-z^2/4}}^{\sqrt{1-z^2/4}}dxdzdy+\bigintsss_{-\frac{1}{2}}^{\frac{5-\sqrt{5}}{3}}\bigintsss_{-2}^2\bigintsss_{-\sqrt{1-z^2/4}}^{\sqrt{1-z^2/4}}dxdzdy\\
+\bigintsss_{\frac{5-\sqrt{5}}{3}}^{\frac{7}{3}}\bigintsss_{-2}^{a\left( y\right) }\bigintsss_{-\sqrt{1-z^2/4}}^{\sqrt{1-z^2/4}}dxdzdy+\bigintsss_{\frac{5-\sqrt{5}}{3}}^{1}\bigintsss_{b\left( y\right) }^2\bigintsss_{-\sqrt{1-z^2/4}}^{\sqrt{1-z^2/4}}dxdzdy\\
+\bigintsss_{\frac{5-\sqrt{5}}{3}}^{\frac{5+\sqrt{5}}{3}}\bigintsss_{a\left(y\right) }^{b\left( y\right) }\bigintsss_{-\sqrt{1-z^2/4}}^{5-z-3y}dxdzdy \approx 19.897.
\end{multline*}
}
\vskip0.2cm
Usando coordenadas cilíndricas, $x=r\cos \theta ,$ $z=r\sin \theta ,$ la ecuación del cilindro toma la forma
$4( r\cos \theta ) ^2+( r\sin \theta ) ^2=4$, que se simplifica
%\vskip0.2cm
{\small
\[r=\frac{2}{\sqrt{4-3\sin ^2\theta }}.\]
}
%\vskip0.2cm
Además, las ecuaciones de los planos toman la forma
$y=\frac{1}{2}r\sin \theta -\frac{3}{2}$ y
$y=-\frac{2}{3}r\sin \theta +\frac{5}{3}.$
Por lo tanto, el volumen del sólido está determinado por
\[\bigintsss_{0}^{2\pi }\bigintsss_{0}^{\frac{2}{\sqrt{4-3\sin ^2\theta }}}\bigintsss_{\frac{1}{2}r\sin \theta -\frac{3}{2}}^{-\frac{2}{3}r\sin\theta +\frac{5}{3}}rdydrd\theta = \frac{19}{3}\pi.\]
%\vskip0.2cm
Algunas de estas integrales se pueden evaluar numéricamente usando las siguientes instrucciones:


\vspace{0.6em}
\noindent
\begin{tcolorbox}[breakable,enhanced,title=\bf Instrucción en \verb"Mathematica", top=2pt,bottom=2pt,boxsep=0pt,before skip=2pt,after skip=0pt]
\begin{Verbatim}[formatcom=\letraverbatim]
a=Sqrt[1-z^2/4];
NIntegrate[1,{y,-5/2,-1/2},{z,-2,3+2*y},
{x,-a,a}]+
NIntegrate[1,{y,-1/2,(5-Sqrt[5])/3},{z,-2,2},{x,-a,a}]
+NIntegrate[1,{y,(5-Sqrt[5])/3,7/3},{z,-2,-12*y/5+4
2*Sqrt[30*y-20-9y^2]/5},{x,-a,a}]+NIntegrate[1,
{y,(5-Sqrt[5])/3,1},{z,-12*y/5+4+2*Sqrt[30*y-20-
9y^2]/5,2},{x,-a,a}]+NIntegrate[1,{y,(5-Sqrt[5])/3,
(5+Sqrt[5])/3},{z,-12*y/5+4-2*Sqrt[30*y-20-9y^2]/5,
12*y/5+4+2*Sqrt[30*y-20-9y^2]/5},{x,-a,5-z-3*y}]
\end{Verbatim}
\end{tcolorbox}
\vspace{0.7em}
\noindent



Este valor también se obtiene con el siguiente procedimiento:
\vskip0.2cm

%\begin{tcolorbox}[title=\bfseries Parametrizando el sólido]
%\footnotesize 
%\parskip=0pt 
%\setlength{\belowdisplayskip}{0pt}
%\setlength{\abovedisplayskip}{0pt} 
		
En el plano $xz$ el cilindro se proyecta como una elipse que se puede parametrizar en la forma
$x=\cos(t),$ $z=2\sin (t),$ para $t\in [0,2\pi ].$
La región interior a esta curva se parametriza de la siguiente manera: $x = u \cos(t),$ $z = 2u\sin(t),$ para
$t\in[0, 2\pi]$, $u\in[0, 1].$ La coordenada $y$ se puede determinar reemplazando en la ecuación de los planos.
De esta manera, para $z=3+2y$ se sigue $y = u\sin(t)-\frac{3}{2},$, mientras que en
$z = 5-x-3y,$ se obtiene que $y = -\frac{2}{3}u\sin(t)+\frac{5}{3}-\frac{1}{3}u \cos(t);$ por lo tanto, es natural que se definan las funciones:
\vspace{-2mm}
{\small
\begin{eqnarray*}
\overrightarrow{\alpha}(t,u)&=&\langle u \cos(t), u\sin(t)-\frac{3}{2},2u\sin(t)\rangle \\
\overrightarrow{\beta}(t,u) &=& \langle u \cos(t),-\frac{2}{3}u\sin(t)+\frac{5}{3}-\frac{1}{3}u \cos(t) ,2u\sin(t)\rangle.
\end{eqnarray*}
}
%\vskip0.2cm

Como parametrizaciones de la región sobre cada uno de los planos que acotan el cilindro. Con las funciones
$\overrightarrow{\alpha}$ y $\overrightarrow{\beta}$ se forma una combinación convexa que permite parametrizar el sólido. Es decir, al simplificar:
$(1-w)\overrightarrow{\alpha}(t,u)+w\overrightarrow{\beta}(t,u)$
se obtiene
\begin{multline*}
{\bf r}( t,u,w) = \big \langle u\cos (t),u\sin(t)-\frac{3}{2}-\frac{5}{3}wu\sin (t)+\frac{19}{6}w- \\ \frac{1}{3}wu\cos (t),2u\sin (t) \big \rangle, 
\end{multline*}
%\vskip0.2cm
en donde $t\in[0, 2\pi],\, u\in[0, 1],, w\in[0, 1], w\in[0, 1].$ También se obtienen los siguientes vectores:
{\small
\begin{eqnarray*}
{\bf r}_{t} &=& \langle -u\sin (t),u\cos (t)-\frac{5}{3}wu\cos (t)+\frac{1}{3}wu\sin (t),2u\cos (t)\rangle \\
{\bf r}_{u}&=& \langle \cos (t),\sin (t)-\frac{5}{3}w\sin (t)-\frac{1}{3}w\cos (t),2\sin (t)\rangle \\
{\bf r}_{w}&=& \langle 0,-\frac{5}{3}u\sin (t)+\frac{19}{6}-\frac{1}{3}u\cos (t),0\rangle.
\end{eqnarray*}
}
%\vskip0.2cm
Formando el producto vectorial triple, se obtiene
%\vskip0.2cm
\[| {\bf r}_{t}\cdot ( {\bf r}_{u}\times {\bf r}_{w}) | = \left| -\frac{10}{3}u^2\sin (t)+\frac{19}{3}u-\frac{2}{3}u^2\cos (t)\right |,\]
%\vskip0.2cm
que se integra con los límites definidos antes y nos proporciona el volumen del sólido
%\vskip0.2cm
\[\int_{0}^{1}\int_{0}^{1}\int_{0}^{2\pi }\left | -\frac{10}{3}u^2\sin (t)
+\frac{19}{3}u-\frac{2}{3}u^2\cos (t)\right | dtdudw= \frac{19}{3}\pi.\]
%\end{tcolorbox}
\vskip0.2cm
\begin{ejer}
Para el sólido acotado por las figuras de las ecuaciones $4x^2 + z^2 = 4,$ $z = 3 + 2y$ y $z = 5 -x-3y.$ Plantee y evalúe numéricamente las integrales triples en los órdenes $dzdxdy$ y $dxdydz.$
\end{ejer}

\section{Tres superficies y un plano}

El sólido acotado por las figuras de las expresiones $y=z^2$, $z=x^2$, $z=x^3$ y $y=0$,
se muestra a continuación. También se aprecian algunas líneas que corresponden a las proyecciones en los planos coordenados de las líneas de intersección de la superficie que lo conforman.

\vskip0.2cm
\noindent
\begin{minipage}{\textwidth}
	\centering
	\captionof{figure}{Visualización tridimensional del sólido definido por $0 \leq y \leq z^2$ \quad y \quad $x^2 \leq z \leq x^2$.}
	\label{fig:solido_parametrico_xyz}
	\begin{minipage}{0.28\linewidth}
		\centering
		\vspace{-12mm}
		\includegraphics[width=\linewidth]{Fig/447.pdf}
	\end{minipage}
	\hspace{1em}
	\begin{minipage}{0.28\linewidth}
		\centering
		\vspace{-12mm}
		\includegraphics[width=\linewidth]{Fig/448.pdf}
	\end{minipage}
	
	\vspace{-8mm}
	\fuentepropia
\end{minipage}
\vskip0.2cm


Las proyecciones del sólido en los planos coordenados son útiles para establecer los límites y evaluar la integral triple sobre dicha región. Enseguida se presentan las regiones correspondientes a los planos $xy$ y $yz$.

\vskip0.2cm
\noindent
\begin{minipage}{\textwidth}
	\centering
	\captionof{figure}{Proyecciones del sólido en los planos $xy$ y $yz$, definido por $x^6 \leq y \leq x^4$ \quad y \quad $0 \leq y \leq z^2$.}
	\label{fig:proyeccion_xy_yz}
	\begin{minipage}{0.30\linewidth}
		\centering
		\vspace{-17mm}
		\includegraphics[width=\linewidth]{Fig/449.pdf}
	\end{minipage}
	\hspace{1em}
	\begin{minipage}{0.30\linewidth}
		\centering
		\vspace{-17mm}
		\includegraphics[width=\linewidth]{Fig/450.pdf}
	\end{minipage}
	
	\vspace{-13mm}
	\fuentepropia
\end{minipage}


\vskip0.2cm
Utilizando la información anterior, el volumen del sólido está dado por las siguientes integrales
\vskip0.2cm
{\small
$\ds \int_{0}^{1}\int_{0}^{x^{6}}\int_{x^{3}}^{x^{2}}dzdydx+\int_{0}^{1}\int_{x^{6}}^{x^{4}}\int_{\sqrt{y}}^{x^{2}}dzdydx= \frac{1}{70}$
}

{\small
$\ds \int_{0}^{1}\int_{\root{6}\of{y}}^{1}\int_{x^{3}}^{x^{2}}dzdxdy+\int_{0}^{1}\int_{\root{4}\of{y}}^{\root{6}\of{y}}\int_{\sqrt{y}}^{x^{2}}dzdxdy= \frac{1}{70}$
}

{\small
$\ds \int_{0}^{1}\int_{0}^{z^{2}}\int_{\root{2}\of{z}}^{\root{3}\of{z}}dxdydz= \int_{0}^{1}\int_{\sqrt{y}}^{1}\int_{\root{2}\of{z}}^{\root{3}\of{z}}dxdzdy= \frac{1}{70}.$
}
\vskip0.4cm
En el plano $xz$ el sólido se proyecta como se muestra enseguida:

\vskip0.2cm
\noindent
\begin{minipage}{\textwidth}
	\centering
	\captionof{figure}{Visualización del sólido definido por $0 \leq y \leq z^2$ \quad y \quad $x^3 \leq z \leq x^2$.}
	\label{fig:solido_xz_yz}
	\begin{minipage}{0.30\linewidth}
		\centering
		\vspace{-13mm}
		\includegraphics[width=\linewidth]{Fig/451.pdf}
	\end{minipage}
	\hspace{1em}
	\begin{minipage}{0.30\linewidth}
		\centering
		\vspace{-13mm}
		\includegraphics[width=\linewidth]{Fig/452.pdf}
	\end{minipage}
	
	\vspace{-8mm}
	\fuentepropia
\end{minipage}


\vskip0.2cm
Las integrales triples que determinan el volumen del sólido en los órdenes $dydxdz$ y $dydzdx$
se presentan a continuación:

\[\ds \int_{0}^{1}\int_{\root{2}\of{z}}^{\root{3}\of{z}}\int_{0}^{z^{2}}dydxdz =\int_{0}^{1}\int_{x^{3}}^{x^{2}}\int_{0}^{z^{2}}dydzdx= \frac{1}{70}.\]
\vskip0.2cm
Las integrales presentadas antes se pueden evaluar en \verb"Mathematica" mediante la ejecución de las siguientes sentencias.


\vspace{0.6em}
\noindent
\begin{tcolorbox}[breakable,enhanced,title=\bf Instrucción en \verb"Mathematica", top=2pt,bottom=2pt,boxsep=0pt,before skip=2pt,after skip=0pt]
\begin{Verbatim}[formatcom=\letraverbatim]
Integrate[1,{x,0,1},{y,0,x^6},{z,x^3,x^2}]+
Integrate[1,{x,0,1},{y,x^6,x^4},{z,Sqrt[y],x^2}]
Integrate[1,{y,0,1},{x,y^(1/6),1},{z,x^3,x^2}]+
Integrate[1,{z,0,1},{y,0,z^2},{z,Sqrt[z],z^(1/3)}]
Integrate[1,{y,0,1},
{z,Sqrt[y],1},{x,Sqrt[z],z^(1/3)}]
Integrate[1,{z,0,1},{x,Sqrt[z],z^(1/3)},{y,0,z^2}]
Integrate[1,{x,0,1},{z,x^3,x^2},{y,0,z^2}]
\end{Verbatim}
\end{tcolorbox}
\vspace{0.7em}
\noindent


El volumen de esta figura también se puede hallar de la siguiente manera.


%\begin{tcolorbox}[title=\bfseries Parametrizando el sólido]
%\noindent
\vskip0.3cm
Sobre el sólido, consideremos los puntos
$A=(\sqrt[2]{u},0,u),$ $B=(\sqrt[3]{u},0,u),$ $C=(\sqrt[2]{u},u^{2},u)$ y $D=(\sqrt[3]{u},u^{2},u),$
para $u \in [0,1].$
Los segmentos $\overline{AB}$ y $\overline{CD}$ se parametrizan respectivamente como:
$\langle(1-v)\sqrt{u}(1-v)+v\sqrt[3]{u},0,u\rangle$ y
$\langle(1-v)\sqrt{u}+v\sqrt[3]{u},u^{2},u\rangle,$ $v \in [0,1].$
 

\vskip0.2cm
\noindent
\begin{minipage}{\textwidth}
	\centering
	\captionof{figure}{Visualización tridimensional del sólido definido por $x^3 \leq z \leq x^2$}
	\label{fig:solido_x3_x2}
	\begin{minipage}{0.36\linewidth}
		\centering
		\includegraphics[width=\linewidth]{Fig/9.pdf}
	\end{minipage}
	
	\vspace{0.3em}
	\fuentepropia
\end{minipage}
\vskip0.2cm


El sólido se recorre completamente uniendo un punto del segmento {\small $\overline{AB}$} con un punto del segmento {\small $\overline{CD},$} esto se consigue simplificando \\
{\small $\left( 1-w\right) \langle(1-v)\sqrt{u}(1-v)+v\sqrt[3]{u},0,u\rangle +w\langle(1-v)\sqrt{u}+v\sqrt[3]{u},u^{2},u\rangle$}
con lo cual, la función que parametriza el sólido es
%\vskip0.2cm
\vspace{-2mm}
\begin{equation} \label{32a}
{\bf r}\left( u,v,w\right) =\left\langle \sqrt{u}-\sqrt{u}v+v\sqrt[3]{u},wu^{2},u \right\rangle,
\end{equation}
%\vskip0.2cm
para $u\in [0,1],$ $v\in [0,1],w\in [0,1].$
También hallamos los vectores
{\small
\begin{eqnarray*}
{\bf r}_{u} &=& \left\langle \frac{1-v}{2\sqrt{u}}+\frac{1}{3}\frac{v}{\left(\sqrt[3]{u}\right) ^{2}},2wu,1\right\rangle \\
{\bf r}_{v} &=& \left\langle-\sqrt{u}+\sqrt[3]{u},0,0\right\rangle \\
{\bf r}_{w} &=& \left\langle 0,u^{2},0\right\rangle.
\end{eqnarray*}
}
%\vskip0.2cm
Además, $\left| {\bf r}_{u}\cdot \left( {\bf r}_{v}\times {\bf r}_{w}\right) \right| =
\left| \left( -\sqrt{u}+\sqrt[3]{u}\right) u^{2}\right|,$
por lo tanto, el volumen del sólido es
\[\int_{0}^{1}\int_{0}^{1}\int_{0}^{1}\left|\left(-\sqrt{u}+\sqrt[3]{u}\right) u^{2}\right|dudvdw= \frac{1}{70}.\]
%\end{tcolorbox}
%\vskip0.2cm
Si la densidad del sólido en mención en cualquier punto está dada por
$\delta \left( x,y,z\right) =xy^{2}z^{3},$ entonces la masa $M$ es:
%\vskip0.2cm
\[M= \ds \int_{0}^{1}\int_{x^{3}}^{x^{2}}\int_{0}^{z^{2}}xy^{2}z^{3}dydzdx=\ds \int_{0}^{1}\int_{\sqrt[2]{z}}^{\sqrt[3]{z}}\int_{0}^{z^{2}}xy^{2}z^{3}dydxdz= \frac{1}{2112}.\]
%\vskip0.2cm
Usando la parametrización escogida (\ref{32a}) se obtiene
%\vskip0.2cm
\[\delta \left( u,v,w\right) =\left( \sqrt{u}-\sqrt{u}v+v\sqrt[3]{u}\right) \left( wu^{2}\right) ^{2}\left( u\right) ^{3}.\]
%\vskip0.2cm
Entonces, la masa del sólido está dada por
%\vskip0.2cm
\[
\resizebox{\linewidth}{!}{
	$M=\int_{0}^{1}\int_{0}^{1}\int_{0}^{1} \left( \sqrt{u}(1-v)+v\sqrt[3]{u} \right) \left( wu^{2}\right) ^{2}u ^3 \left|\left( -\sqrt{u}+\sqrt[3]{u}\right) u^{2}\right| \,du\,dv$
}.
\]
%\vskip0.2cm
Los momentos con respecto a los planos coordenados están determinados por

\vspace{-10mm}
{\small
\begin{eqnarray*}
\mu_{yz} &=&\int _0^1\int _{x^3}^{x^2}\int _0^{z^2}x^2 y^2 z^3dydzdx =\frac{1}{2277} \\
\mu_{xz} &=& \int _0^1\int _{x^3}^{x^2}\int _0^{z^2}x y^3 z^3dydzdx =\frac{1}{3952} \\
\mu_{xy} &=& \int _0^1\int _{x^3}^{x^2}\int _0^{z^2}x y^2 z^4dydzdx=\frac{1}{2520}.
\end{eqnarray*}
}
\vskip0.2cm
Por lo tanto, las coordenadas $(\overline{x},\overline{y},\overline{z})$ del centro de masa del sólido están dadas por
$\left(\frac{\mu{yz}}{M}, \frac{\mu{xz}}{M},\frac{\mu{xy}}{M}\right)=\left(\frac{64}{69}, \frac{132}{247},\frac{88}{105}\right).$
\vskip0.2cm
\subsection*{Intersección de cilindros perpendiculares}
Consideremos el sólido acotado por dos cilindros de igual radio con ejes de simetría perpendicular. En particular, se puede suponer que poseen ecuaciones $x^{2}+z^{2}=r^{2}$ y
$x^{2}+y^{2}=r^{2},$ con $r>0$. 

%\vskip0.2cm
\vskip0.3cm
\noindent
\begin{minipage}{\textwidth}
	\centering
	\captionof{figure}{Sólido definido por $x^2 + z^2 \leq r^2$ \quad y \quad $x^2 + y^2 \leq r^2$.}
	\label{fig:bicilindro}
	\begin{minipage}{0.20\linewidth}
		\centering
		\includegraphics[width=\linewidth]{Img/51a.pdf}
	\end{minipage}
	\hspace{1em}
	\begin{minipage}{0.23\linewidth}
		\centering
		\includegraphics[width=\linewidth]{Img/51c.pdf}
	\end{minipage}
	\hspace{1em}
	\begin{minipage}{0.23\linewidth}
		\centering
		\includegraphics[width=\linewidth]{Img/51d.pdf}
	\end{minipage}
	
	\vspace{0.3em}
	\fuentepropia
\end{minipage}

\vskip0.3cm
Estos gráficos se obtienen ejecutando las siguientes instrucciones, en ellas se consideró $r=1$.


\vspace{0.6em}
\noindent
\begin{tcolorbox}[breakable,enhanced,title=\bf Instrucción en \verb"Mathematica", top=2pt,bottom=2pt,boxsep=0pt,before skip=2pt,after skip=0pt]
\begin{Verbatim}[formatcom=\letraverbatim]
ContourPlot3D[{x^2+y^2==1,x^2+z^2==1},{x,-1,1},
{y,-1,1},{z,-1,1},Mesh->False,PlotPoints->70,
ContourStyle->{Blend[{Cyan,Yellow},1/5],Yellow}]
RegionPlot3D[-Sqrt[1-y^2]<z<Sqrt[1-y^2]&&-Sqrt[1-
x^2]<y<Sqrt[1-x^2]&&-1<x<1,{x,-1,1},{y,-1,1},
{z,-1,1},PlotStyle->RGBColor[0.8,0.3,0.3],
PlotPoints->90,Mesh->False]
\end{Verbatim}
\end{tcolorbox}
\vspace{0.7em}
\noindent



En el plano $xy$, el sólido se proyecta como una circunferencia de radio $r,$ y en el plano
$yz$ el gráfico correspondiente es un cuadrado. Eliminado $x$ de las ecuaciones
$x^2+y^2=r^2$ y $x^2+z^2=r^2$ se obtienen las ecuaciones de las rectas $z = \pm y$,
razón por la cual, el cuadrado se divide en cuatro regiones.


\noindent
\begin{minipage}{\textwidth}
	\centering
	\captionof{figure}{Proyecciones del sólido definido por $x^2 + y^2 \leq r^2$, \quad $|x| \leq 1$ \quad y \quad $|z| \leq 1$.}
	\label{fig:cilindro_recortado}
	\begin{minipage}{0.25\linewidth}
		\centering
		\vspace{-10mm}
		\includegraphics[width=\linewidth]{Fig/454.pdf}
	\end{minipage}
	\hspace{1em}
	\begin{minipage}{0.25\linewidth}
		\centering
		\vspace{-10mm}
		\includegraphics[width=\linewidth]{Fig/455.pdf}
	\end{minipage}
	
	\vspace{-8mm}
	\fuentepropia
\end{minipage}
\vskip0.2cm

Estas figuras son útiles para plantear las integrales que determinan el volumen del sólido.
Integrando primero con respecto a $z$, esta recorre los puntos desde la parte negativa del cilindro con ecuación $x^2+z^2=r^2$ hasta su parte positiva y los límites en $x$ e $y$ se determinan con la circunferencia mostrada antes.
%\vskip0.2cm
\[\int_{-r}^{r}\int_{-\sqrt{r^{2}-x^{2}}}^{\sqrt{r^{2}-x^{2}}}\int_{-\sqrt{r^{2}-x^{2}}}^{\sqrt{r^{2}-x^{2}}}dzdydx=\int_{-r}^{r}\int_{-\sqrt{r^{2}-y^{2}}}^{\sqrt{r^{2}-y^{2}}}\int_{-\sqrt{r^{2}-x^{2}}}^{\sqrt{r^{2}-x^{2}}}dzdxdy= \frac{16}{3}r^{3}.\]
\vskip0.2cm
Para el planteamiento en  orden $dxdydz$, $x$ recorre los puntos entre la parte negativa del cilindro $x^2+z^2=r^2$ hasta su parte positiva y los límites para $y$ y $z$
se determinan con el cuadrado.
%\vskip0.2cm
\begin{multline*}
\int_{0}^{r}\int_{-z}^{z}\int_{-\sqrt{r^{2}-z^{2}}}^{\sqrt{r^{2}-z^{2}}}dxdydz+\int_{-r}^{0}\int_{z}^{-z}\int_{-\sqrt{r^{2}-z^{2}}}^{\sqrt{r^{2}-z^{2}}}dxdydz \\
+\int_{0}^{r}\int_{z}^{r}\int_{-\sqrt{r^{2}-y^{2}}}^{\sqrt{r^{2}-y^{2}}}dxdydz +\int_{-r}^{0}\int_{-z}^{r}\int_{-\sqrt{r^{2}-y^{2}}}^{\sqrt{r^{2}-y^{2}}}dxdydz\\
+\int_{0}^{r}\int_{-r}^{-z}\int_{-\sqrt{r^{2}-y^{2}}}^{\sqrt{r^{2}-y^{2}}}dxdydz+\int_{-r}^{0}\int_{-r}^{z}\int_{-\sqrt{r^{2}-y^{2}}}^{\sqrt{r^{2}-y^{2}}}dxdydz=\frac{16}{3}r^{3}.
\end{multline*}
\vskip0.2cm
Por la simetría de la figura, el planteamiento en los órdenes $dxdzdy$, $dydzdx$ y $dydxdz$ es muy similar al descrito en los casos ya vistos.

Estas integrales se evaluaron en \Verb"Mathematica" usando el comando \\ \Verb"Assuming",
debido a que el programa aporta una solución que depende de la naturaleza de $r$
(positiva, negativa, real o compleja). Esto se logra con las siguientes instrucciones.


\vspace{0.6em}
\noindent
\begin{tcolorbox}[breakable,enhanced,title=\bf Instrucción en \verb"Mathematica", top=2pt,bottom=2pt,boxsep=0pt,before skip=2pt,after skip=0pt]
\begin{Verbatim}[formatcom=\letraverbatim]
Integrate[1,{x,-r,r},{y,-Sqrt[r^2-x^2],Sqrt[r^2-x^2]},
{z,-Sqrt[r^2-x^2],Sqrt[r^2-x^2]}]
Assuming[r\[Element]Reals&&r>0,Integrate[1,{y,-r,r},
{x,-Sqrt[r^2-y^2],Sqrt[r^2-y^2]},{z,-Sqrt[r^2-x^2],
Sqrt[r^2-x^2]}]]

Assuming[r\[Element]Reals&&r>0,Integrate[1,{z,0,r},
{y,-z,z},{x,-Sqrt[r^2-z^2],Sqrt[r^2-z^2]}]+
Integrate[1,{z,-r,0},{y,z,-z},{x,-Sqrt[r^2-z^2],
Sqrt[r^2-z^2]}]+Integrate[1,{z,0,r},{y,z,r},
{x,-Sqrt[r^2-y^2],Sqrt[r^2-y^2]}]+Integrate[1,
{z,-r,0},{y,-z,r},{x,-Sqrt[r^2-y^2],Sqrt[r^2-y^2]}]
+Integrate[1,{z,0,r},{y,-r,-z},{x,-Sqrt[r^2-y^2],
Sqrt[r^2-y^2]}]+Integrate[1,{z,-r,0},{y,-r,z},
{x,-Sqrt[r^2-y^2],Sqrt[r^2-y^2]}]]
\end{Verbatim}
\end{tcolorbox}
\vspace{0.7em}
\noindent


En coordenadas cilíndricas, considerando el radio $a,$, el volumen se determina por la siguiente integral
%\vskip0.2cm
\[\int 0^{2 \pi }\int _0^a\int _{-\sqrt{a^2-r^2 \cos ^2(\theta)}}^{\sqrt{a^2-r^2\cos^2(\theta)}}rdzdrd\theta=\frac{16}{3}a^{3}.\]
%\vskip0.2cm
Esta se evalúa como sigue:


\vspace{0.6em}
\noindent
\begin{tcolorbox}[breakable,enhanced,title=\bf Instrucción en \verb"Mathematica", top=2pt,bottom=2pt,boxsep=0pt,before skip=2pt,after skip=0pt]
\begin{Verbatim}[formatcom=\letraverbatim]
Integrate[r,{\[Theta],0,2*Pi},{r,0,a},{z,-Sqrt[a^2
r^2*Cos[\[Theta]]^2],Sqrt[a^2-r^2*Cos[\[Theta]]^2]}]
\end{Verbatim}
\end{tcolorbox}
\vspace{0.7em}
\noindent


También se puede obtener el volumen usando el siguiente proceso.
\vskip0.3cm
%\begin{tcolorbox}[title=\bfseries Parametrizando el sólido]
El sólido se proyecta en el plano $xy$ como un círculo unitario; esta región se describe con las ecuaciones $x=u\cos v$, $y=u\sin v$, con $u \in [0,r]$, $v\in [0,2\pi]$.
Sobre el cilindro $x^2+z^2=r^2$ se satisface que $z =\pm \sqrt{r^2-u^2 \cos ^2v}.$
Por tanto, el sólido se parametriza mediante una combinación convexa de dos puntos de la forma
%\vskip0.2cm
\[P(u \cos (v),u \sin (v),-\sqrt{r^2-u^2 \cos ^2v})\] y \[Q(u \cos (v),u \sin (v),\sqrt{r^2-u^2 \cos ^2v}),\]
esto es, $(1-w)P+wQ.$

Simplificando esta expresión, se define la función
%\vskip0.2cm
\[{\boldsymbol \alpha}(u,v,w)= \langle u \cos (v),u \sin (v),(2w-1)\sqrt{r^2-u^2 \cos ^2(v)} \rangle,\]
%\vskip0.2cm
en donde $u\in [0,r]$, $v\in [0,2\pi]$ y $w\in [0,1].$
Derivando parcialmente se obtienen los siguientes vectores
\vspace{-3mm}
{\small
\begin{eqnarray*}
{\boldsymbol \alpha}_u &=& \langle \cos(v), \sin(v), \frac{u(1-2w)\cos^2(v)}{\sqrt{r^2-u^2 \cos^2(v)}} \rangle \\
{\boldsymbol \alpha}_v &=& \langle-u \sin(v), u\cos(v), \frac{u^2(2w-1)\sin(v)\cos(v)}{\sqrt{r^2-u^2 \cos^2v}}\rangle \\
{\boldsymbol \alpha}_w &=& \langle 0, 0, 2\sqrt{r^2-u^2 \cos^2v}\rangle.
\end{eqnarray*}
}
%\vskip0.2cm
Entonces,
{\small $| {\boldsymbol \alpha}_w \cdot ({\boldsymbol \alpha}_u \times {\boldsymbol \alpha}_v)| = 2u\sqrt{r^2-u^2 \cos^2(v)},$}
con lo cual el volumen del mencionado esta dado por
%\vskip0.2cm
{\small
\[\int_0^1 \int_0^{2\pi} \int_0^r 2u\sqrt{r^2-u^2 \cos^2(v)} du dv dw= \frac{16}{3}r^3\]
}
%\end{tcolorbox}
\vskip0.2cm
En \verb"Mathematica" el cálculo anterior se realiza con las siguientes instrucciones:


\vspace{0.6em}
\noindent
\begin{tcolorbox}[breakable,enhanced,title=\bf Instrucción en \verb"Mathematica", top=2pt,bottom=2pt,boxsep=0pt,before skip=2pt,after skip=0pt]
\begin{Verbatim}[formatcom=\letraverbatim]
{x,y,z}={u*Cos[v],u*Sin[v],Sqrt[r^2-x^2]};
R=Simplify[(1-w)*{x,y,-z}+w*{x,y,z}];
A=Assuming[r\[Element]Reals&&r>0,
Simplify[Dot[D[R,w],Cross[D[R,u],D[R,v]]]]]
Integrate[A,{v,0,2Pi},{u,0,r},{w,0,1}]
\end{Verbatim}
\end{tcolorbox}
\vspace{0.7em}
\noindent


Si la intersección se hace con tres cilindros del mismo radio, cuyos ejes de simetría son perpendiculares entre sí. La situación es un poco similar al procedimiento presentado cuando se trata de dos cilindros.
Consideremos los cilindros con ecuaciones
$x^2+y^2=r^2$, $y^2+z^2=r^2,$ y $x^2+z^2=r^2.$
Con las siguientes instrucciones se obtienen los gráficos presentados, en ellos se usó $r=1$.


\vspace{0.6em}
\noindent
\begin{tcolorbox}[breakable,enhanced,title=\bf Instrucción en \verb"Mathematica", top=2pt,bottom=2pt,boxsep=0pt,before skip=2pt,after skip=0pt]
\begin{Verbatim}[formatcom=\letraverbatim]
ContourPlot3D[{x^2+y^2==1,y^2+z^2==1,x^2+z^2==1},
{x,-1,1},{y,-1,1},{z,-1,1},Mesh->False,
ContourStyle->{Blend[{Cyan,Yellow},2/5],
RGBColor[0.9,0.8,0.6],Orange}]
RegionPlot3D[-Sqrt[1-y^2]<z<Sqrt[1-y^2]&&-Sqrt[1-
x^2]<y<Sqrt[1-x^2]&&-1<x<1,{x,-1,1},{y,-1,1},
{z,-1,1},PlotStyle->RGBColor[0.8,0.3,0.3],
PlotPoints->90,Mesh->False]
\end{Verbatim}
\end{tcolorbox}
\vspace{0.7em}
\noindent


El producto es el siguiente:

\vskip0.2cm
\noindent
\begin{minipage}{\textwidth}
	\centering
	\captionof{figure}{Intersección de tres cilindros: $x^2 + y^2 \leq r^2$, \quad $y^2 + z^2 \leq r^2$, \quad $x^2 + z^2 \leq r^2$}
	\label{fig:tricilindro}
	\begin{minipage}{0.21\linewidth}
		\centering
		\includegraphics[width=\linewidth]{Img/52a.pdf}
	\end{minipage}
	\begin{minipage}{0.25\linewidth}
		\centering
		\includegraphics[width=\linewidth]{Img/52b.pdf}
	\end{minipage}
	
	
	\vspace{0.3em}
	\fuentepropia
\end{minipage}
\vskip0.2cm





\noindent

La simetría de la figura con respecto a los planos coordenados, permite hacer el planteamiento en el primer octante o en la mitad de este.
La siguiente figura muestra la proyección del sólido en el plano $yz$.


\vskip0.2cm
\noindent
\begin{minipage}{\textwidth}
	\centering
	\captionof{figure}{Región definida por $y^2 + z^2 \leq r^2$, \quad $x > 0$ \quad y \quad $y > 0$}
	\label{fig:cuarto_disco_yz}
	\begin{minipage}{0.31\linewidth}
		\centering
		\vspace{-14mm}
		\includegraphics[width=\linewidth]{Fig/456.pdf}
	\end{minipage}
	
	\vspace{-9mm}
	\fuentepropia
\end{minipage}
\vskip0.2cm


Presentamos las integrales que permiten calcular el volumen, integrando primero con respecto a $z.$
%\vskip0.2cm
\vspace{-2mm}
{\small
\begin{multline*}
8\bigg( \int_{0}^{r/\sqrt{2}}\int_{0}^{x}\int_{0}^{\sqrt{r^{2}-x^{2}}}dzdydx+\int_{r/\sqrt{2}}^{r}\int_{0}^{\sqrt{r^{2}-x^{2}}}\int_{0}^{\sqrt{r^{2}-x^{2}}}dzdydx \\
+\int_{0}^{r/\sqrt{2}}\int_{x}^{\sqrt{r^{2}-x^{2}}}\int_{0}^{\sqrt{r^{2}-y^{2}}}dzdydx\bigg) = 8\left( 2-\sqrt{2}\right) r^{3} \\
8\bigg( \int_{0}^{r/\sqrt{2}}\int_{y}^{\sqrt{r^{2}-y^{2}}}\int_{0}^{\sqrt{r^{2}-x^{2}}}dzdxdy+\int_{0}^{r/\sqrt{2}}\int_{0}^{y}\int_{0}^{\sqrt{r^{2}-y^{2}}}dzdxdy \\
+\int_{r/\sqrt{2}}^{r}\int_{0}^{\sqrt{r^{2}-y^{2}}}\int_{0}^{\sqrt{r^{2}-y^{2}}}dzdxdy\bigg) = 8\left( 2-\sqrt{2}\right) r^{3}.
\end{multline*}
}
%\vskip0.2cm
Para evaluar las integrales anteriores, se ejecutaron estas instrucciones:


\vspace{0.6em}
\noindent
\begin{tcolorbox}[breakable,enhanced,title=\bf Instrucción en \verb"Mathematica", top=2pt,bottom=2pt,boxsep=0pt,before skip=2pt,after skip=0pt]
\begin{Verbatim}[formatcom=\letraverbatim]
Assuming[r\[Element]Reals&&r>0,8*(Integrate[1,{x,0,
r/Sqrt[2]},{y,0,x},{z,0,Sqrt[r^2-x^2]}]+Integrate[1,
{x,r/Sqrt[2],r},{y,0,Sqrt[r^2-x^2]},{z,0,Sqrt[r^2-
x^2]}]+Integrate[1,{x,0,r/Sqrt[2]},{y,x,Sqrt[r^2
x^2]},{z,0,Sqrt[r^2-y^2]}])]
Assuming[r\[Element]Reals&&r>0,8*(Integrate[1,{y,0,
r/Sqrt[2]},{x,y,Sqrt[r^2-y^2]},{z,0,Sqrt[r^2-x^2]}]+
Integrate[1,{y,0,r/Sqrt[2]},{x,0,y},{z,0,Sqrt[r^2-
y^2]}]+Integrate[1,{y,r/Sqrt[2],r},{x,0,Sqrt[r^2-
y^2]},{z,0,Sqrt[r^2-y^2]}])]
\end{Verbatim}
\end{tcolorbox}
\vspace{0.7em}
\noindent



Un cambio a coordenadas cilíndricas, usando $a$ como el radio de los cilindros y observando solamente la mitad del primer octante, nos lleva al mismo valor
\[16\int_{0}^{\pi/4}\int_{0}^{a}\int_{0}^{\sqrt{a^{2}-r^{2}\cos^{2}\theta}}rdzdrd\theta=8(2-\sqrt{2})a^3.\]

\section{Dos esferas} \index{Esfera}
Las ecuaciones $x^2+y^2+(z-1)^2=5$ y $x^2+y^2+z^2=4$
representan dos esferas que se intersectan en $z=0$. Consideraremos el sólido acotado por los hemisferios negativos en $z$.

\begin{figure}[H]
\centering
\caption{Región definida por $2\rho \cos \varphi + 4 \leq \rho^2 \leq 4$.}
\label{fig:region_polar_desplazada}
\includegraphics[height=4.5cm]{Img/90b.pdf}
\hspace{0.2cm}
\includegraphics[height=4.5cm]{Img/90a.pdf}
\fuentepropia
\end{figure}

Estos gráficos se obtienen con las siguientes instrucciones:


\vspace{0.6em}
\noindent
\begin{tcolorbox}[breakable,enhanced,title=\bf Instrucción en \verb"Mathematica", top=2pt,bottom=2pt,boxsep=0pt,before skip=2pt,after skip=0pt]
\begin{Verbatim}[formatcom=\letraverbatim]
ContourPlot3D[{x^2+y^2+(z-1)^2==5,x^2+y^2+z^2==4},
{x,-3,3},{y,-3,3},{z,-3,4},Mesh->False,
ContourStyle->{Directive[Cyan,Opacity[0.8]],
Directive[Orange,Opacity[0.9]]},PlotPoints->60]
\end{Verbatim}
\end{tcolorbox}
\vspace{0.7em}
\noindent


En coordenadas rectangulares, el volumen está dado por la integral triple en el orden $dzdydx:$
%\vskip0.2cm
\[\int_{-2}^{2}\int_{-\sqrt{4-x^{2}}}^{\sqrt{4-x^{2}}}\int_{-\sqrt{4-x^{2}-y^{2}}}^{1-\sqrt{5-x^{2}-y^{2}}}dzdydx=\frac{10}{3} \left(3-\sqrt{5}\right) \pi\]
%\vskip0.2cm
La proyección del sólido en el plano $xz$ se muestra a continuación.
\vskip0.2cm

\noindent
\begin{minipage}{\textwidth}
	\centering
	\captionof{figure}{Región definida por $2z + 4 \leq x^2 + z^2 \leq 4$, \quad $x^2 + (z - 1)^2 \leq 4$.}
	\label{fig:region_corona_desplazada_2}
	\begin{minipage}{0.35\linewidth}
		\centering
		\vspace{-20mm}
		\includegraphics[width=\linewidth]{Fig/457.pdf}
	\end{minipage}
	
	
	\vspace{-18mm}
	\fuentepropia
\end{minipage}

\vskip0.3cm
Con esta figura se plantea la integral que determina el volumen del sólido en el orden $dydzdx$, se usa la abreviatura
$a=\sqrt{5-x^{2}-(z-1)^{2}}.$
{\small
\begin{multline*}
\hskip-0.2cm \int_{-2}^{2}\int_{1-\sqrt{5-x^{2}}}^{0}\int_{-\sqrt{4-x^{2}-z^{2}}}^{-a}dydzdx +\int_{-2}^{2}\int_{1-\sqrt{5-x^{2}}}^{0}\int_{a}^{\sqrt{4-x^{2}-z^{2}}}dydzdx \\
+\int_{-2}^{2}\int_{-\sqrt{4-x^{2}}}^{1-\sqrt{5-x^{2}}}\int_{-\sqrt{4-x^{2}-z^{2}}}^{\sqrt{4-x^{2}-z^{2}}}dydzdx \approx 7.999.
\end{multline*}
}
%\vskip0.2cm
Por la simetría de la figura con respecto al plano $yz$, la integral en el orden $dxdzdy$ es similar.
Las integrales anteriores se evaluaron usando \verb"Mathematica" ejecutando las siguientes instrucciones:


\vspace{0.6em}
\noindent
\begin{tcolorbox}[breakable,enhanced,title=\bf Instrucción en \verb"Mathematica", top=2pt,bottom=2pt,boxsep=0pt,before skip=2pt,after skip=0pt]
\begin{Verbatim}[formatcom=\letraverbatim]
Integrate[1,{x,-2,2},{y,-Sqrt[4-x^2],Sqrt[4-x^2]},
{z,-Sqrt[4-x^2-y^2],1-Sqrt[5-x^2-y^2]}]
Integrate[r,{\[Theta],0,2Pi},{r,0,2},
{z,-Sqrt[4-r^2],1-Sqrt[5-r^2]}]
NIntegrate[2,{x,-2,2},{z,1-Sqrt[5-x^2],0},
{y,-Sqrt[4-x^2-z^2],-Sqrt[5-x^2-(z-1)^2]}]
+NIntegrate[2,{x,-2,2},{z,1-Sqrt[5-x^2],0},
{y,Sqrt[5-x^2-(z-1)^2],Sqrt[4-x^2-z^2]}]
+NIntegrate[1,{x,-2,2},{z,-Sqrt[4-x^2],1-Sqrt[5-x^2]},
{y,-Sqrt[4-x^2-z^2],Sqrt[4-x^2-z^2]}]
\end{Verbatim}
\end{tcolorbox}
\vspace{0.7em}
\noindent


En coordenadas cilíndricas, esta integral se escribe como
%\vskip0.2cm
\[\int_{0}^{2\pi }\int_{0}^{2}\int_{-\sqrt{4-r^{2}}}^{1-\sqrt{5-r^{2}}}rdzdrd\theta=\frac{10}{3}\pi\left(3-\sqrt{5}\right).\]
%\vskip0.2cm
En coordenadas esféricas, la ecuación $x^{2}+y^{2}+z^{2}=4$ se escribe como $\rho =2$ y la ecuación
$x^{2}+y^{2}+\left( z-1\right) ^{2}=5$ se transforma en $\rho ^{2}-2\rho \cos \phi=4,$ de donde
$\rho =\cos \phi +\sqrt{\cos ^{2}\phi +4}$.
La intersección entre las esferas se determina por $\phi = \pi /2.$
Entonces, el volumen del sólido esta dado por
%\vskip0.2cm
\[\int_{\pi /2}^{\pi }\int_{0}^{2\pi }\int_{\cos \phi +\sqrt{\cos ^{2}\phi +4}}^{2}\rho ^{2}\sin \phi d\rho d\theta d\phi = \frac{10}{3}\pi \left( 3-\sqrt{5}\right).\]
%\vskip0.2cm
La evaluación de esta integral se puede hacer como sigue.


\vspace{0.6em}
\noindent
\begin{tcolorbox}[breakable,enhanced,title=\bf Instrucción en \verb"Mathematica", top=2pt,bottom=2pt,boxsep=0pt,before skip=2pt,after skip=0pt]
\begin{Verbatim}[formatcom=\letraverbatim]
Integrate[\[Rho]^2*Sin[\[Phi]],{\[Phi],Pi/2,Pi},
{\[Theta],0,2*Pi},{\[Rho],Cos[\[Phi]]+Sqrt[4+
Cos[\[Phi]]^2],2}]
\end{Verbatim}
\end{tcolorbox}
\vspace{0.7em}
\noindent


El siguiente procedimiento permite obtener el mismo valor.

%\begin{tcolorbox}[title=\bfseries Parametrizando el sólido]
La curva de intersección se determina por $z=0$ y está dada por $x^{2}+y^{2}=4,$, cuyo interior se parametriza como
$x=u\cos v,$ $y=u\sin v,$ con $u\in [0,2],$ $v\in [0,2\pi ]$.
La coordenada $z$, en las dos esferas, es
$z= -\sqrt{4-u^{2}}$, $z= 1-\sqrt{5-u^{2}}$, respectivamente.
El sólido se parametriza formando una combinación convexa entre dos puntos de las esferas con iguales coordenadas $x$ e $y.$ Esto es,
$(1-w) \left\langle u\cos v,u\sin v,-\sqrt{4-u^{2}}\right\rangle +w\left\langle u\cos v,u\sin v,1-\sqrt{5-u^{2}}\right\rangle$
que se simplifica y define la función
{\small
\[
\resizebox{\linewidth}{!}{
	$\mathbf{r}(u,v,w)=\left\langle u\cos v,u\sin v,w\left(1-\sqrt{5-u^{2}}\right)-(1-w)\sqrt{4-u^{2}}\right\rangle.$
}
\]
}
Esta función nos permite obtener los siguientes vectores y el elemento diferencial de volumen
{\small
\begin{eqnarray*}
\mathbf{r} &=& \left\langle \cos v, \sin v, \frac{1}{1+u^2} \sqrt{4-u^2} \right\rangle + \frac{u}{\sqrt{5-u^2}} \mathbf{r} \\
\mathbf{r}_u &=& \left\langle \text{(completa aquí)} \right\rangle \\
\mathbf{r}_v &=& \left\langle -\sin v, \cos v, 0 \right\rangle \\
\mathbf{r}_u &=& \left\langle 0, 0, \sqrt{4-u^2} \right\rangle \\
dV &=& \left| \mathbf{r}_u \times \mathbf{r}_v \right| = \left| \sqrt{4-u^2} \right|.
\end{eqnarray*}
}
Entonces, el volumen del sólido está dado por
\vspace{-2mm}
\[\int_{0}^{1}\int_{0}^{2\pi }\int_{0}^{2}\left\vert u\left( \sqrt{5-u^{2}}-
\sqrt{4-u^{2}}-1\right)\right\vert dudvdw=\frac{10(3-\sqrt{5})}{3}\pi.\]
%\end{tcolorbox}

\section{Sólidos con cilindros cardióidicos} \index{Cardioide}

En esta sección se consideran cilindros cuyos cortes perpendiculares al eje de simetría son cardioides.

\index{Cilindro cardioidico}
\subsection*{Cilindros cardióidicos y el plano z=2-x/2}

En el plano, las ecuaciones $r = 1-\cos(t)$ y $r = 2- 2\cos(t)$ representan dos cardioides, mientras que en el espacio las mismas ecuaciones representan dos cilindros cardióidecos.



\vskip0.3cm
\noindent
\begin{minipage}{\textwidth}
	\centering
	\captionof{figure}{Región delimitada por $r = 1 - \cos(t)$, $r = 2 - 2\cos(t)$, \quad y \quad $z = 2 - \frac{x}{2}$.}
	\label{fig:cardioide_revolucion}
	\begin{minipage}{0.28\linewidth}
		\centering
		\vspace{-10mm}
		\includegraphics[width=\linewidth]{Fig/458.pdf}
	\end{minipage}
	\hspace{1em}
	\begin{minipage}{0.25\linewidth}
		\centering
		\vspace{-10mm}
		\includegraphics[width=\linewidth]{Img/36a.pdf}
	\end{minipage}
	
	\vspace{-8mm}
	\fuentepropia
\end{minipage}


\vskip0.3cm
Consideremos el sólido acotado por las figuras de las expresiones $z = 2-x/2,$
$z=0$ y los cilindros $r = 1- \cos(t)$ y $r = 2-2\cos(t).$




\vskip0.2cm
\noindent
\begin{minipage}{\textwidth}
	\centering
	\captionof{figure}{Sólido definido por $1 - \cos(t) \leq r \leq 2 - 2\cos(t)$ \quad y \quad $0 \leq z \leq 2 - \frac{x}{2}$}
	\label{fig:cardioide_z_inclinado}
	\begin{minipage}{0.22\linewidth}
		\centering
		\includegraphics[width=\linewidth]{Img/36c.pdf}
	\end{minipage}
	\begin{minipage}{0.22\linewidth}
		\centering
		\includegraphics[width=\linewidth]{Img/36d.pdf}
	\end{minipage}
	\begin{minipage}{0.22\linewidth}
		\centering
		\includegraphics[width=\linewidth]{Img/36e.pdf}
	\end{minipage}
	
	\vspace{0.3em}
	\fuentepropia
\end{minipage}
\vskip0.2cm



Las figuras de las cardioides en una animación que muestra el recorrido de esta curva se obtienen ejecutando las siguientes instrucciones:


\vspace{0.6em}
\noindent
\begin{tcolorbox}[breakable,enhanced,title=\bf Instrucción en \verb"Mathematica", top=2pt,bottom=2pt,boxsep=0pt,before skip=2pt,after skip=0pt]
\begin{Verbatim}[formatcom=\letraverbatim]
Manipulate[ParametricPlot[{{(1-Cos[t])*Cos[t],
(1-Cos[t])*Sin[t]},{2*(1-Cos[t])*Cos[t],
2*(1-Cos[t])*Sin[t]}},{t,0,k},PlotRange->{{-4,1},
{-3,3}}],{k,Pi/20,2*Pi}]
\end{Verbatim}
\end{tcolorbox}
\vspace{0.7em}
\noindent



Las tres superficies básicas que intervienen, se consiguen con las siguientes instrucciones:



\vspace{0.6em}
\noindent
\begin{tcolorbox}[breakable,enhanced,title=\bf Instrucción en \verb"Mathematica", top=2pt,bottom=2pt,boxsep=0pt,before skip=2pt,after skip=0pt]
\begin{Verbatim}[formatcom=\letraverbatim]
a=ParametricPlot3D[{{(1-Cos[t])*Cos[t],
(1-Cos[t])*Sin[t],u},{(2-2*Cos[t])*Cos[t],
(2-2*Cos[t])*Sin[t],u}},{u,0,4},{t,0,2*Pi},
Mesh->False,PlotPoints->80,PlotStyle->{Orange,
RGBColor[0.7,0.1,0.1]}];
b=ContourPlot3D[{z==2-x/2},{x,-5,5},{y,-5,5},{z,-2,4},
Mesh->False,ContourStyle->RGBColor[0.2,0.4,0.3]];
Show[a,b]
\end{Verbatim}
\end{tcolorbox}
\vspace{0.7em}
\noindent


Además, la parte del plano que se halla entre los dos cilindros cardioidicos se consigue con estas sentencias:


\vspace{0.6em}
\noindent
\begin{tcolorbox}[breakable,enhanced,title=\bf Instrucción en \verb"Mathematica", top=2pt,bottom=2pt,boxsep=0pt,before skip=2pt,after skip=0pt]
\begin{Verbatim}[formatcom=\letraverbatim]
ParametricPlot3D[{(1+u)*(1-Cos[t])*Cos[t],
(1+u)*(1-Cos[t])*Sin[t],2-(1+u)*(1-Cos[t])*Cos[t]/2},
{u,0,1},{t,0,2*Pi},Mesh->False,
PlotStyle->RGBColor[0.6,0.3,0.3],PlotPoints->80]
\end{Verbatim}
\end{tcolorbox}
\vspace{0.7em}
\noindent



La figura del sólido se consigue utilizando coordenadas cilíndricas y ejecutando las siguientes tres instrucciones:



\vspace{0.6em}
\noindent
\begin{tcolorbox}[breakable,enhanced,title=\bf Instrucción en \verb"Mathematica", top=2pt,bottom=2pt,boxsep=0pt,before skip=2pt,after skip=0pt]
\begin{Verbatim}[formatcom=\letraverbatim]
{r,t}={Sqrt[x^2+y^2],ArcTan[x,y]};
RegionPlot3D[1-Cos[t]<r<2-2*Cos[t]&&0<z<2-x/2,
{x,-4,2},{y,-4,4},{z,0,4},PlotPoints->90,Mesh->False]
\end{Verbatim}
\end{tcolorbox}
\vspace{0.7em}
\noindent



En el plano $xy,$ los cilindros cardióidicos se pueden parametrizar, haciendo
$x = (1 -\cos(t))\cos(t),$ $y = (1 - \cos(t)) \sin(t),\,t \in[0, 2\pi];$ esto para el cilindro menor. Mientras que $x = (2 - 2\cos(t))\cos(t),$ $y = (2 - 2\cos(t))\sin(t)$ se usarán para el cilindro mayor.
Lo anterior, nos permite definir las funciones vectoriales que representan las curvas de intersección de los cilindros con el plano $xy$ ($z = 0$), que son
%\vskip0.2cm
\begin{eqnarray*}
\overrightarrow{\alpha} (t) &=& \langle(1-\cos(t))\cos(t), (1-\cos(t))\sin(t), 0\rangle ,\,\, t \in[0, 2\pi] \\
\overrightarrow{\beta} (t)& = &\langle(2-2\cos(t))\cos(t), (2-2\cos(t))\sin(t), 0\rangle ,\,\, t \in[0, 2\pi.]
\end{eqnarray*}
%\vskip0.2cm
Con estas funciones formamos la combinación convexa $(1-u)\overrightarrow{\alpha}(t)+u\overrightarrow{\beta}(t)$que determina la región contenida entre las dos curvas (en el plano $z=0$). Esto es:
%\vskip0.2cm
{\small
\begin{equation}\label{36a}
\overrightarrow{r}_1(t,u)=(1+u)\left\langle 2 \sin ^2\left(\frac{t}{2}\right) \cos (t),\sin (t) (1-\cos (t)),0 \right\rangle,
\end{equation}
}
%\vskip0.2cm
para $u\in[0, 1]$ y $t \in[0, 2\pi]$; proyectando esta región sobre el plano $z=2-x/2$, obtenemos la función
{\small
\begin{eqnarray}
\overrightarrow{r}_2(t,u)=\bigg\langle 2 (u+1) \sin ^2\left(\frac{t}{2}\right) \cos(t),(u+1) \sin (t) (1-\cos (t)),\nonumber\\
2-(u+1) \sin^2\left(\frac{t}{2}\right) \cos (t) \bigg\rangle \label{36b}
\end{eqnarray}
}

\vskip0.2cm
\noindent
\begin{minipage}{\textwidth}
	\centering
	\captionof{figure}{Sólido definido por $1 - \cos t \leq r \leq 2 - 2\cos t$, \quad $0 \leq z \leq 2 - \frac{x}{2}$.}
	\label{fig:solido_cardioide_inclinado}
	\begin{minipage}{0.35\linewidth}
		\centering
		\vspace{-15mm}
		\includegraphics[width=\linewidth]{Fig/459.pdf}
	\end{minipage}
	\hspace{1em}
	\begin{minipage}{0.35\linewidth}
		\centering
		\vspace{-15mm}
		\includegraphics[width=\linewidth]{Fig/460.pdf}
	\end{minipage}
	
	\vspace{-8mm}
	\fuentepropia
\end{minipage}


\vskip0.2cm
Un punto $Q,$ sobre la base del sólido, tendrá la forma de la expresión \ref{36a}, mientras que un punto $P$ del sólido sobre el plano $z=2-x/2$,
tiene la forma \ref{36b}. Con el fin de parametrizar la región en el espacio, formamos una combinación convexa de los puntos $P$ y $Q$, de la forma
$(1-w)\overrightarrow{r}_1(t,u)+w\overrightarrow{r}_2(t,u).$ Esto define la función
\begin{multline*}
{\bf r}(t,u,w) = \bigg \langle 2 (u+1) \sin ^2\left(\frac{t}{2}\right)\cos(t),-(u+1) \sin (t) (\cos (t)-1),\\
w \left(2-(u+1)\sin ^2\left(\frac{t}{2}\right) \cos(t)\right)\bigg \rangle,
\end{multline*}
con $t\in [0,2\pi]$, $u\in [0,1]$, $w\in [0,1].$
Ahora obtenemos los vectores
{\small
	\[
	\mathbf{r}_{t} = (1+u)\left\langle \sin (2 t)-\sin (t),2\sin^2\left(\tfrac{t}{2}\right) (2\cos(t)+1),
	\frac{1}{2} w (\sin (t)-\sin (2 t)) \right\rangle,
	\]
	\[
	\mathbf{r}_{u} = \left\langle2 \sin ^2\left(\tfrac{t}{2}\right) \cos (t),-\sin(t)(\cos (t)+1),-w \sin^2\left(\tfrac{t}{2}\right) \cos (t) \right\rangle,
	\]
	\[
	\mathbf{r}_{w} = \left\langle 0,0,2-(u+1) \sin ^2\left(\tfrac{t}{2}\right) \cos(t)\right\rangle.
	\]
}
Además,
{\small
\[\left| {\bf r}_{t}\cdot \left( {\bf r}_{u}\times {\bf r}_{w}\right) \right| =
\left|(u+1) \sin^4\left(\frac{t}{2}\right) ((u+1)(\cos (2 t)-2 \cos (t))+u+9) \right|.\]
}
Por último, el volumen está dado por:
\vskip0.2cm
\[
\int_{0}^{1} \int_{0}^{1} \int_{0}^{2\pi}
\left| \mathbf{r}_{t} \cdot \left( \mathbf{r}_{u} \times \mathbf{r}_{w} \right) \right|
\, dt\, du\, dw = \frac{107}{8}\pi.
\]
\vskip0.2cm
La parametrización y el cálculo anterior, se pueden hacer en \verb"Mathematica" con ayuda de las siguientes instrucciones


\vspace{0.6em}
\noindent
\begin{tcolorbox}[breakable,enhanced,title=\bf Instrucción en \verb"Mathematica", top=2pt,bottom=2pt,boxsep=0pt,before skip=2pt,after skip=0pt]
\begin{Verbatim}[formatcom=\letraverbatim]
{x1,y1}={Cos[t]*(1-Cos[t]),Sin[t]*(1-Cos[t])};
r=(1-u)*{x1,y1,0}+u*{x2,y2,0};
s=(1-u)*{x1,y1,2-x1/2}+u*{2*x1,2*y1,2-x1};
v=(1-w)*s+w*r;
A=Dot[D[v,t],Cross[D[v,u],D[v,w]]];
Integrate[Abs[A],{t,0,2*Pi},{u,0,1},{w,0,1}]
\end{Verbatim}
\end{tcolorbox}
\vspace{0.7em}
\noindent


Otra manera de realizar dicho cálculo es el uso de coordenadas cilíndricas, de esta manera el volumen está dado por:
\[\int_{0}^{2\pi }\int_{1-\cos\theta}^{2-2\cos\theta}\int_{0}^{2-\frac{1}{2}r\cos\theta}rdzdrd\theta = \frac{107}{8}\pi.\]

\subsection*{Cilindros cardioidicos y el plano z=2-y/2}

Las ecuaciones $r = 1 - \cos(t)$ y $r = 2 -2 \cos(t)$ representan dos cardioides en el plano, mientras que en el espacio las mismas ecuaciones representan dos cilindros cardioides.
\vskip0.2cm

\noindent
\begin{minipage}{\textwidth}
	\centering
	\captionof{figure}{Región delimitada por $r = 1 - \cos(t)$, $r = 2 - 2\cos(t)$, \quad y \quad $z = 2 - \frac{y}{2}$.}
	\label{fig:cardioide_y_inclinado}
	\begin{minipage}{0.26\linewidth}
		\centering
		\vspace{-10mm}
		\includegraphics[width=\linewidth]{Fig/461.pdf}
	\end{minipage}
	\hspace{1em}
	\begin{minipage}{0.24\linewidth}
		\centering
		\vspace{-10mm}
		\includegraphics[width=\linewidth]{Img/35e.pdf}
	\end{minipage}
	
	\vspace{-5mm}
	\parbox{\linewidth}{\centering\fontsize{9pt}{10pt}\selectfont\textbf{Fuente:} elaboración propia.}
\end{minipage}
\vskip0.2cm


Consideremos el sólido acotado por las figuras de las expresiones $z = 2 -y/2,$
$z=0$ y los cilindros $r=1-\cos(t)$ y $r=2- 2 \cos(t).$

\begin{figure}[H]
\centering
\caption{Sólido definido por $1 - \cos(t) \leq r \leq 2 - 2\cos(t)$ \quad y \quad $0 \leq z \leq 2 - \frac{y}{2}$.}
\label{fig:cardioide_y_inclinado_2} \includegraphics[height=4.5cm]{Img/35a.pdf}
\hspace{0.3cm}
\includegraphics[height=4.5cm]{Img/35b.pdf}
\hspace{0.3cm}
\includegraphics[height=4.5cm]{Img/35c.pdf}
\parbox{\linewidth}{\centering\fontsize{8pt}{10pt}\selectfont\textbf{Fuente:} elaboración propia.}
\end{figure}

Las figuras se obtienen en coordenadas cilíndricas, definiendo $r$ y $t=\theta$.


\vspace{0.6em}
\noindent
\begin{tcolorbox}[breakable,enhanced,title=\bf Instrucción en \verb"Mathematica", top=2pt,bottom=2pt,boxsep=0pt,before skip=2pt,after skip=0pt]
\begin{Verbatim}[formatcom=\letraverbatim]
{r,t}={Sqrt[x^2+y^2],ArcTan[x,y]};
RegionPlot3D[1-Cos[t]<r<2-2*Cos[t]&&0<z<2-y/2,
{x,-4,2},{y,-4,4},{z,0,4},PlotPoints->90,Mesh->False]
\end{Verbatim}
\end{tcolorbox}
\vspace{0.7em}
\noindent


De igual manera, se pueden usar las siguientes instrucciones:


\vspace{0.6em}
\noindent
\begin{tcolorbox}[breakable,enhanced,title=\bf Instrucción en \verb"Mathematica", top=2pt,bottom=2pt,boxsep=0pt,before skip=2pt,after skip=0pt]
\begin{Verbatim}[formatcom=\letraverbatim]
a=ParametricPlot3D[{{(1-Cos[t])*Cos[t],
(1-Cos[t])*Sin[t],u},{(2-2*Cos[t])*Cos[t],
(2-2*Cos[t])*Sin[t],u}},{u,0,3.5},{t,0,2*Pi},
PlotStyle->{Orange,RGBColor[0.7,0.1,0.1]},
Mesh->False,PlotPoints->80];
b=ContourPlot3D[{z==2-y/2},{x,-5,2},{y,-4,4},{z,-2,
3.5},Mesh->False,ContourStyle->RGBColor[0.2,0.4,0.3]];
Show[a,b]
\end{Verbatim}
\end{tcolorbox}
\vspace{0.7em}
\noindent


El cilindro menor se parametriza por $x = (1-\cos(t)) \cos(t),$ $y = (1-\cos(t))\sin(t)$
y para el cilindro mayor se hace $x = (2- 2\cos(t)) \cos(t),$
$y = (2 -2\cos(t))\sin(t)$, con $t \in [0, 2\pi]$. Las siguientes son las funciones vectoriales que representan las curvas de intersección de los cilindros con el plano $xy$ $(z=0)$:
%\vskip0.2cm
{\small
\begin{eqnarray*}
\overrightarrow{\alpha }(t)&=&\langle ( 1-\cos ( t)) \cos(t),( 1-\cos(t)) \sin(t),0\rangle \\
\overrightarrow{\beta }(t)&=&\langle ( 2-2\cos(t))\cos(t),( 2-2\cos(t)) \sin(t),0\rangle,
\end{eqnarray*}
}
%\vskip0.2cm
con $t\in [0,2\pi ].$ Posteriormente, formamos una combinación convexa de ellas que parametriza la región contenida entre las dos curvas.
%\vskip0.2cm
{\small
\begin{eqnarray}
\overrightarrow{r}_1(t,u) &=& ( 1-u) \overrightarrow{\alpha }(t)+u\overrightarrow{\beta } (t)\nonumber \\
&=& (u+1)\langle 2 \sin ^2\left(\frac{t}{2}\right)\cos(t),\sin(t)(1-\cos(t)),0\rangle, \label{625x}
\end{eqnarray}
}
%\vskip0.2cm
para $u\in [0,1]$ y $t\in [0,2\pi ].$
Proyectando esta región sobre el plano $z=2-y/2$, esta superficie se parametriza como
%\vskip0.2cm
{\small
\begin{eqnarray}
\overrightarrow{r}_2(t,u) &=& \big \langle 2(u+1)\sin^2\left(\frac{t}{2}\right)\cos(t),(u+1)\sin(t)(1-\cos(t)), \nonumber\\
&& \hskip3.25 cm \frac{1}{2}(u+1)\sin(t)(\cos(t)-1)+2 \big \rangle \label{626x}
\end{eqnarray}
}
\vskip0.2cm
Cuya figura es la siguiente
\vskip0.2cm
\noindent
\begin{minipage}{\textwidth}
	\centering
	\captionof{figure}{Superficie $r_2(t,u)$ definida sobre región cardioidal con altura $z = 2 - \frac{r \sin(t)}{2}$.}
	\label{fig:superficie_cardioide_inclinada}
	\begin{minipage}{0.26\linewidth}
		\centering
		\vspace{-12mm}
		\includegraphics[width=\linewidth]{Img/35d.pdf}
	\end{minipage}
	\hspace{1em}
	\begin{minipage}{0.26\linewidth}
		\centering
		\vspace{-12mm}
		\includegraphics[width=\linewidth]{Fig/462.pdf}
	\end{minipage}
	
	\vspace{-8mm}
	\fuentepropia
\end{minipage}

\vskip0.2cm
Un punto $Q$ sobre la base del sólido tendrá la forma de la expresión (\ref{625x}), mientras que un punto $P$ del sólido sobre el plano $z=2-y/2$, tiene la forma \ref{626x}.
\vskip0.2cm

\noindent
\begin{minipage}{\textwidth}
	\centering
	\captionof{figure}{Sólido definido por $1 - \cos(t) \leq r \leq 2 - 2\cos(t)$ \quad y \quad $0 \leq z \leq 2 - \frac{y}{2}$.}
	\label{fig:solido_cardioide_y_inclinado}
	\begin{minipage}{0.35\linewidth}
		\centering
		\vspace{-12mm}
		\includegraphics[width=\linewidth]{Fig/463.pdf}
	\end{minipage}
	\hspace{1em}
	\begin{minipage}{0.35\linewidth}
		\centering
		\vspace{-12mm}
		\includegraphics[width=\linewidth]{Fig/464.pdf}
	\end{minipage}
	
	\vspace{-8mm}
	\fuentepropia
\end{minipage}


\vskip0.2cm

Con el fin de parametrizar la región en el espacio formamos una combinación convexa de los puntos $P$ y $Q$ de la forma $(1-w)\overrightarrow{r}_1(t,u)+w \overrightarrow{r}_2(t,u),$ que al simplificarse define la función

\vspace{-8mm}
\begin{multline}
{\bf r}( t,u,w) =\bigg\langle 2 (u+1) \sin ^2\left(\frac{t}{2}\right)\cos(t),(u+1)\sin (t) (1-\cos(t)), \\
w \left(\frac{1}{2} (u+1) \sin (t) (\cos (t)-1)+2\right)\bigg \rangle
\end{multline}
De la cual se obtienen los siguientes vectores
\begin{equation*}
	\begin{aligned}
		{\bf r}_{t} &= (u+1)\Big \langle (\sin(2t)-\sin(t)), \\
		&\qquad 2\sin^2\left(\frac{t}{2}\right)(2\cos (t)+1), w \sin ^2\left(\frac{t}{2}\right) (2 \cos (t)+1) \Big \rangle \\
		{\bf r}_{u} &= \Big\langle 2\sin^2\left(\frac{t}{2}\right)\cos(t),\sin(t)(-(\cos(t)-1)),\\
		&\qquad \frac{1}{2}w\sin(t)(\cos(t)-1)\Big\rangle \\
		{\bf r}_{w} &= \Big\langle 0,0,\frac{1}{2} (u+1) \sin (t) (\cos (t)-1)+2 \Big\rangle.
	\end{aligned}
\end{equation*}
Con los cuales se obtiene el elemento diferencial de volumen
$|{\bf r}_{t}\cdot({\bf r}_{u}\times{\bf r}_{w})|=\big|-(u+1)\sin^4\left(\frac{t}{2}\right)(8-2(u+1)\sin(t)+(u+1)\sin(2t)) \big| .$
Finalmente, el volumen del sólido está dado por
\vskip0.2cm
{\small
\[\int_{0}^{1}\int_{0}^{1}\int_{0}^{2\pi}\bigg|-(u+1)\sin^4\left(\frac{t}{2}\right)(8-(u+1)(2\sin(t)-\sin(2t)))\bigg| dtdudw= 9\pi\]
}
\vskip0.2cm
Este valor también se puede obtener utilizando coordenadas cilíndricas.
\vskip0.2cm
\[\int_{0}^{2\pi } \int_{1-\cos \theta }^{2-2\cos \theta }\int_{0}^{2-\frac{r}{2}\sin \theta }rdzdrd\theta = 9\pi.\]
\vskip0.2cm
Los cálculos anteriores se pueden realizar en \verb"Mathematica" ejecutando las siguientes sentencias.



\vspace{0.6em}
\noindent
\begin{tcolorbox}[breakable,enhanced,title=\bf Instrucción en \verb"Mathematica", top=2pt,bottom=2pt,boxsep=0pt,before skip=2pt,after skip=0pt]
\begin{Verbatim}[formatcom=\letraverbatim]
r1={Cos[t]*(1-Cos[t]),Sin[t]*(1-Cos[t])};
r2={Cos[t]*2*(1-Cos[t]),2*Sin[t]*(1-Cos[t])};
r={x,y}=Simplify[(1-u)*r1+u*r2];
R=Simplify[(1-w)*{x,y,0}+w*{x,y,2-y/2}];
A=Simplify[Dot[D[R,t],Cross[D[R,u],D[R,w]]]]
Integrate[Abs[A],{t,0,2*Pi},{u,0,1},{w,0,1}]
\end{Verbatim}
\end{tcolorbox}
\vspace{0.7em}
\noindent


\subsection*{Dos cilindros cardioidicos y un paraboloide}
El paraboloide con ecuación $4z = x^2 + y^2$, el plano $xy$ y los cilindros cuya base son las cardioiodes con ecuaciones $r=1-\cos t$ y $r=2 - 2 \cos t,$ acotan un sólido.




\vskip0.2cm
\noindent
\begin{minipage}{\textwidth}
	\centering
	\captionof{figure}{Sólido definido por $r = 1 - \cos t$, \quad $r = 2 - 2\cos t$, y $4z = x^2 + y^2$}
	\label{fig:solido_cardioide_paraboloide}
	\begin{minipage}{0.26\linewidth}
		\centering
		\includegraphics[width=\linewidth]{Img/37a.pdf}
	\end{minipage}
	\begin{minipage}{0.26\linewidth}
		\centering
		\includegraphics[width=\linewidth]{Img/37b.pdf}
	\end{minipage}
	\begin{minipage}{0.26\linewidth}
		\centering
		\includegraphics[width=\linewidth]{Img/37c.pdf}
	\end{minipage}
	
	\vspace{0.3em}
	\fuentepropia
\end{minipage}
\vskip0.2cm


Con las siguientes indicaciones se consiguen las figuras presentadas



\vspace{0.6em}
\noindent
\begin{tcolorbox}[breakable,enhanced,title=\bf Instrucción en \verb"Mathematica", top=2pt,bottom=2pt,boxsep=0pt,before skip=2pt,after skip=0pt]
\begin{Verbatim}[formatcom=\letraverbatim]
a=ParametricPlot3D[{{(1-Cos[t])*Cos[t],
(1-Cos[t])*Sin[t],u},{(2-2*Cos[t])*Cos[t],
(2-2*Cos[t])*Sin[t],u}},{u,0,4},{t,0,2*Pi},
PlotStyle->{Orange,RGBColor[0.7,0.1,0.1]},
Mesh->False,PlotPoints->80];
b=ContourPlot3D[{4*z==x^2+y^2},{x,-4,4},{y,-4,4},
{z,0,4},ContourStyle->RGBColor[0.4,0.5,0.3],
Mesh->False];
Show[b,a]
\end{Verbatim}
\end{tcolorbox}
\vspace{0.7em}
\noindent




\vskip0.2cm
\noindent
\begin{minipage}{\textwidth}
	\centering
	\captionof{figure}{Sólido definido por $1 - \cos(t) \leq r \leq 2 - 2\cos(t)$ \quad y \quad $0 \leq 4z \leq r^2$}
	\label{fig:cardioide_paraboloide}
	\begin{minipage}{0.25\linewidth}
		\centering
		\includegraphics[width=\linewidth]{Img/37d.pdf}
	\end{minipage}
	\begin{minipage}{0.25\linewidth}
		\centering
		\includegraphics[width=\linewidth]{Img/37f.pdf}
	\end{minipage}
	\begin{minipage}{0.25\linewidth}
		\centering
		\includegraphics[width=\linewidth]{Img/37e.pdf}
	\end{minipage}
	
	\vspace{0.3em}
	\fuentepropia
\end{minipage}
\vskip0.2cm



El sólido se consigue definiendo las coordenadas polares y ejecutando las siguientes instrucciones:


\vspace{0.6em}
\noindent
\begin{tcolorbox}[breakable,enhanced,title=\bf Instrucción en \verb"Mathematica", top=2pt,bottom=2pt,boxsep=0pt,before skip=2pt,after skip=0pt]
\begin{Verbatim}[formatcom=\letraverbatim]
{r,t}={Sqrt[x^2+y^2],ArcTan[x,y]};
RegionPlot3D[1-Cos[t]<r<2-2*Cos[t]&&0<z<(x^2+y^2)/4,
{x,-4,2},{y,-4,4},{z,0,4.2},PlotPoints->90,Mesh->False]
\end{Verbatim}
\end{tcolorbox}
\vspace{0.7em}
\noindent


La base del sólido es igual a la región definida por la fórmula (\ref{625x}). Sobre el paraboloide la coordenada $z$ del sólido es $z = (u+1)^2 \sin ^4\left(\frac{t}{2}\right),$ esto permite definir la función
{\small
\begin{multline} \label{628x}
\overrightarrow{r}_2(t,u)=(u+1)\bigg \langle 2\sin^2\left(\frac{t}{2}\right)\cos(t),\sin(t)(1-\cos(t)),\\ (u+1)\sin^4\left(\frac{t}{2}\right) \bigg\rangle
\end{multline}
}



\noindent
\begin{minipage}{\textwidth}
	\centering
	\captionof{figure}{Sólido definido por $1 - \cos(t) \leq r \leq 2 - 2\cos(t)$ \quad y \quad $0 \leq 4z \leq r^2$.}
	\label{fig:solido_cardioide_paraboloide_2}
	\begin{minipage}{0.33\linewidth}
		\centering
		\vspace{-18mm}
		\includegraphics[width=\linewidth]{Fig/464.pdf}
	\end{minipage}
	\hspace{1em}
	\begin{minipage}{0.33\linewidth}
		\centering
		\vspace{-18mm}
		\includegraphics[width=\linewidth]{Fig/465.pdf}
	\end{minipage}
	
	\vspace{-10mm}
	\fuentepropia
\end{minipage}
\vskip0.3cm


Un punto $Q$ sobre la base del sólido tendrá coordenadas determinadas por la expresión \ref{625x}
y un punto $P$ sobre el paraboloide tendrá la forma dada por la expresión \ref{628x};
con ellos se puede formar una combinación convexa que parametriza el sólido. Esto permite definir la función
\begin{equation*}
	\mathbf{r}(t,u,w) = (u+1)\left\langle 2\sin^2\left(\tfrac{t}{2}\right)\cos(t),\sin(t)(1-\cos(t)),(u+1)w\sin^4\left(\tfrac{t}{2}\right)\right\rangle,
\end{equation*}
en donde $t\in [0,2\pi],$ $u \in [0,1]$, $w \in [0,1].$
Con esta función se pueden obtener los siguientes vectores
{\footnotesize
\begin{align*}
	{\bf r}_{t} &= (u+1)\left \langle\sin(2t)-\sin(t), 2\sin^2\left(\frac{t}{2}\right)(2\cos (t)+1), w(u+1) \sin ^2\left(\frac{t}{2}\right) \sin (t) \right \rangle \\
	{\bf r}_{u} &=\left \langle 2 \sin ^2\left(\frac{t}{2}\right) \cos (t), -\sin (t)(\cos (t)+1), 2(u+1) w \sin ^4\left(\frac{t}{2}\right) \right \rangle \\
	{\bf r}_{w} &=\left \langle 0,0,(u+1)^2 \sin ^4\left(\frac{t}{2}\right) \right \rangle.
\end{align*}
}
Finalmente, el volumen está dado por
\vskip0.2cm
\[\int_{0}^{1}\int_{0}^{1}\int_{0}^{2\pi }\left | -4 (u+1)^3 \sin ^8\left(\frac{t}{2}\right)\right|dtdudw= \frac{525}{64}\pi.\]
\vskip0.2cm
Utilizando coordenadas cilíndricas, también se puede obtener este valor
\vskip0.2cm
\[\int_{0}^{2\pi}\int_{1-\cos\theta}^{2-2\cos\theta}\int_{0}^{\frac{1}{4}r^{2}}rdzdrd\theta=\frac{525}{64}\pi.\]
\vskip0.2cm
Las siguientes instrucciones permiten construir la parametrización del sólido, hallar el elemento diferencial de volumen y evaluar la integral que proporciona su volumen:

\vspace{0.6em}
\noindent
\begin{tcolorbox}[breakable,enhanced,title=\bf Instrucción en \verb"Mathematica", top=2pt,bottom=2pt,boxsep=0pt,before skip=2pt,after skip=0pt]
\begin{Verbatim}[formatcom=\letraverbatim]
r1={Cos[t]*(1-Cos[t]),Sin[t]*(1-Cos[t])};
r2={Cos[t]*2*(1-Cos[t]),2*Sin[t]*(1-Cos[t])};
{x,y}=Simplify[(1-u)*r1+u*r2];
R=Simplify[(1-w)*{x,y,0}+w*{x,y,(x^2+y^2)/4}];
A=Simplify[Dot[D[R,t],Cross[D[R,u],D[R,w]]]]
Integrate[Abs[A],{t,0,2*Pi},{u,0,1},{w,0,1}]
\end{Verbatim}
\end{tcolorbox}
\vspace{0.7em}
\noindent


\subsection*{Un cilindro cardióidico y dos planos}

Consideremos el sólido acotado por cilindro cardióidico con ecuación $r=2-2\cos (t)$ y los planos con ecuaciones $z=0$ y $z=1-y/3.$




\vskip0.2cm
\noindent
\begin{minipage}{\textwidth}
	\centering
	\captionof{figure}{Sólido definido por $r = 2 - 2\cos(t)$ \quad y \quad $z = 1 - \frac{y}{3}$}
	\label{fig:cardioide_plano_inclinado}
	\begin{minipage}{0.25\linewidth}
		\centering
		\includegraphics[width=\linewidth]{Img/38a.pdf}
	\end{minipage}
	\begin{minipage}{0.25\linewidth}
		\centering
		\includegraphics[width=\linewidth]{Img/38b.pdf}
	\end{minipage}
	\begin{minipage}{0.25\linewidth}
		\centering
		\includegraphics[width=\linewidth]{Img/38c.pdf}
	\end{minipage}
	
	\vspace{0.3em}
	\fuentepropia
\end{minipage}
\vskip0.2cm


Con las siguientes instrucciones se consiguen las figuras presentadas:


\vspace{0.6em}
\noindent
\begin{tcolorbox}[breakable,enhanced,title=\bf Instrucción en \verb"Mathematica", top=2pt,bottom=2pt,boxsep=0pt,before skip=2pt,after skip=0pt]
\begin{Verbatim}[formatcom=\letraverbatim]
a=ParametricPlot3D[{(2-2*Cos[t])*Cos[t],
(2-2*Cos[t])*Sin[t],u},{u,0,2},{t,0,2*Pi},
Mesh->False,PlotStyle->Orange,PlotPoints->90];
b=ContourPlot3D[z==1-y/3,{x,-4,2},{y,-3,3},{z,0,2},
Mesh->False,ContourStyle->RGBColor[0.4,0.5,0.3],
PlotPoints->80];
Show[b,a]
{r,t}={Sqrt[x^2+y^2],ArcTan[x,y]};
RegionPlot3D[r<2-2*Cos[t]&&0<z<1-y/3,{x,-4,2},
{y,-3,3},{z,0,2},PlotPoints->90,Mesh->False]
\end{Verbatim}
\end{tcolorbox}
\vspace{0.7em}
\noindent



A partir de la ecuación $r=2-2\cos(t)$ se sigue que, en el plano $xy$, la cardioide se parametriza
$x=2u\cos(t)(1-\cos(t))$ y $y=2u\sin(t)(1-\cos(t)).$
Consideremos un punto $Q(2u\cos(t)(1-\cos(t)),2u\sin(t)(1-\cos(t)),0),$ sobre la base del sólido y un punto
$P\big(2u\cos(t)(1-\cos(t)), \linebreak 2u\sin(t)(1-\cos(t)),1+\frac{2}{3}u\sin(t)(\cos(t)-1)\big),$
sobre el plano $z=1-y/3,$ para $u\in [0,1]$, $t\in [0,2\pi].$
Una combinación convexa de la forma $(1-w)P+wQ$
parametriza el sólido, esto es
\begin{eqnarray}
{\bf r}\left(t,u,w\right)=\big\langle 2u(1-\cos(t))\cos(t),2u\sin(t)(1-\cos(t)),\nonumber\\
\frac{1}{3}(1-w)(3-2u\sin(t)+u\sin(2t)) \big\rangle, \label{38a}
\end{eqnarray}
para $t \in [0,2\pi ]$, $u\in [0,1]$, $w\in [0,1]$.

\vskip0.2cm
\noindent
\begin{minipage}{\textwidth}
	\centering
	\captionof{figure}{Sólido definido por $0 \leq r \leq 2(1 - \cos t)$ \quad y \quad $0 \leq 3z \leq 3 - y$.}
	\label{fig:cardioide_plano_inclinado_final}
	\begin{minipage}{0.35\linewidth}
		\centering
		\vspace{-18mm}
		\includegraphics[width=\linewidth]{Fig/467.pdf}
	\end{minipage}
	\hspace{1em}
	\begin{minipage}{0.35\linewidth}
		\centering
		\vspace{-18mm}
		\includegraphics[width=\linewidth]{Fig/468.pdf}
	\end{minipage}
	
	\vspace{-10mm}
	\fuentepropia
\end{minipage}


\vskip0.2cm
Derivando parcialmente la función \ref{38a} se obtienen los siguientes vectores
{\small
\begin{flalign*}
	&{\bf r}_t = \Big \langle 2u\sin(t)(2\cos(t)-1), &\\
	&\qquad 4u\sin^2\left(\frac{t}{2}\right)(2\cos(t)+1), \frac{2}{3}u(w-1)(\cos(t)-\cos(2t)) \Big \rangle ,&\\
	& {\bf r}_{u} = \Big\langle 2(1-\cos(t))\cos(t), 2\sin(t)(1-\cos(t)),&\\
	&\qquad \frac{1}{3}(1-w)(\sin(2t)-2\sin(t)) \Big\rangle \quad \text{y}&\\
	&{\bf r}_{w}=\Big\langle 0,0,\frac{1}{3}(2u\sin(t)-2u\sin(t)\cos(t)-3) \Big\rangle .&
\end{flalign*}
}
\vskip0.2cm
Esto implica que {\small
$\left| {\bf r}_t \cdot \left( {\bf r}_{u}\times {\bf r}_{w}\right) \right| =
\left | \frac{16}{3}u\sin^4\left(\frac{t}{2}\right)(3-2u\sin(t)+u\sin(2t)) \right|.$
}
Por tanto, el volumen del sólido está dado por
{\small\[\int_{0}^{1}\int_{0}^{1}\int_{0}^{2\pi }\left( \frac{4}{3}u\left( \cos (t) -1\right) ^{2}\left( 2u\sin (t) \left( \cos (t) -1\right) +3\right)\right) dtdudw= 6\pi.\]
}
En \verb"Mathematica" el proceso anterior se ejecuta como sigue.


\vspace{0.6em}
\noindent
\begin{tcolorbox}[breakable,enhanced,title=\bf Instrucción en \verb"Mathematica", top=2pt,bottom=2pt,boxsep=0pt,before skip=2pt,after skip=0pt]
\begin{Verbatim}[formatcom=\letraverbatim]
P={u*(2-2*Cos[t])*Cos[t],u*(2-2*Cos[t])*Sin[t],
1-(u*(2-2*Cos[t])*Sin[t])/3};
Q={u*(2-2*Cos[t])*Cos[t],u*(2-2*Cos[t])*Sin[t],0};
r=Simplify[(1-w)*P+w*Q];
A=Simplify[Dot[D[r,t],Cross[D[r,u],D[r,w]]]];
Integrate[A,{w,0,1},{u,0,1},{t,0,2*Pi}]
\end{Verbatim}
\end{tcolorbox}
\vspace{0.7em}
\noindent


Para ver el resultado en cada instrucción, el usuario debe, además de ejecutar cada sentencia, eliminar el punto y coma (;) en algunas líneas.
En coordenadas cilíndricas, el volumen está dado por la siguiente integral
\vskip0.2cm
\[\int_{0}^{2\pi }\int_{0}^{2-2\cos \theta }\int_{0}^{1-\frac{1}{3} r\sin \theta }rdzdrd\theta = 6\pi.\]


\subsection*{Un cilindro cardióidico y un paraboloide}

Consideremos el parabóloide con ecuación $3z=9-(x+2)^{2}-(y+1)^{2}$ y el cilindro cardióidico
$r=1-\cos (t)$.
Esta región se presenta en la sección \secref{cilcardiopar}. \vskip0.2cm
En dos dimensiones, para $t\in [0,2\pi], $ la cardioide se parametriza como
{\small
\[\vec{r}(t)=\left\langle (1-\cos(t))\cos(t),(1-\cos(t))\sin(t)\right\rangle.\]
}

\begin{figure}[H]
\centering
\caption{Sólido definido por $r = 1 - \cos t$ \quad y \quad $3z = 4 - 4x - x^2 - 2y - y^2$.}
\label{fig:cardioide_superficie_parabolica}
\includegraphics[width=4cm,height=3cm]{Img/58a.pdf}\hskip0.5cm
\includegraphics[width=4cm,height=3cm]{Img/58b.pdf}
\fuentepropia
\end{figure}

La región interior a la cardioide en el plano \textit{xy} se parametriza como
\begin{align}
x(u,t)&= u(1-\cos (t))\cos (t) \label{38b} \\
y(u,t)&= u(1-\cos (t))\sin (t) \label{38c}
\end{align}
% ${\bf r}_1(t,u)=\left\langle u(1-\cos (t))\cos (t),u(1-\cos (t))\sin (t),0 \right\rangle ,$
La coordenada $z$ sobre el paraboloide satisface que,
\begin{equation}
	\resizebox{0.85\textwidth}{!}{%
		$z(u,t)= \frac{1}{3}\left(9-(u(1-\cos (t))\cos (t)+2)^{2} -(u(1-\cos (t))\sin(t)+1)^{2}\right)$
	}
	\label{38d}
\end{equation}

Este término en la ecuación \ref{38d} se reescribe como
{\footnotesize
\[\frac{1}{6}\left( (4u-u^2)\cos(2t) + (4u^2-8u)\cos(t) -4u\sin(t)+2u\sin(2t)-3u^2+4u+8 \right).\]
}


\noindent
\begin{minipage}{\textwidth}
	\centering
	\captionof{figure}{Región sobre el paraboloide interior al cilindro.}
	\label{fig:paraboloide_interior_cilindro}
	\begin{minipage}{0.35\linewidth}
		\centering
		\vspace{-15mm}
		\includegraphics[width=\linewidth]{Fig/469.pdf}
	\end{minipage}
	\hspace{1em} 
	\begin{minipage}{0.35\linewidth}
		\centering
		\vspace{-15mm}
		\includegraphics[width=\linewidth]{Fig/470.pdf}
	\end{minipage}
	
	\vspace{-10mm}
	\fuentepropia
\end{minipage}


\vskip0.2cm
La base del sólido se describe por medio de la función 
\vskip0.2cm
\[{\bf r}_1(u,t)=\langle x(u,t),y(u,t),0\rangle.\]
\vskip0.2cm
% \begin{equation} {\bf r}_1(u,t)=\langle x(u,t),y(u,t),0\rangle. \end{equation}
Utilizando las expresiones \ref{38b}, \ref{38c} y \ref{38d} se usará la función
%\vskip0.2cm
\begin{equation}
{\bf r}_2(u,t)=\langle x(u,t),y(u,t),z(u,t)\rangle,
\end{equation}
%\vskip0.2cm
con $u\in [0,1],$ $t\in [0,2\pi ],$ como una parametrización de la superficie del paraboloide dentro del cilindro cardióidico.
El sólido se parametriza con $w\in [0,1],$ mediante una combinación convexa entre un punto sobre el paraboloide y un punto sobre el plano $xy$, esto es $(1-w){\bf r}_1(u,t)+w{\bf r}_2(u,t)$.
Con lo cual, el sólido queda determinado por medio de la función
\begin{multline*}
{\bf r} (t,u,w) =\bigg\langle u(1-\cos(t))\cos(t),u\sin(t)(1-\cos(t)),\frac{1}{6}w\bigg(2u\sin(2t) \\
4u\sin(t)+4(u-2) u\cos(t)-(u-4)u\cos(2t)-3u^2+4u+8\bigg) \bigg\rangle.
\end{multline*}
Puede mostrarse que
\begin{multline*}
\left| {\bf r}_{t}\cdot \left( {\bf r}_{u}\times {\bf r}_{w}\right) \right|
=\bigg | \frac{1}{6}u(\cos(t)-1)^2\bigg(8+4u-3u^2-u((8-4u)\cos(t)+\\ (u-4)\cos(2t)+ 4\sin(t)-2\sin(2t))\bigg)\bigg |,
\end{multline*}
entonces,
\[\int_{0}^{1}\int_{0}^{1}\int_{0}^{2\pi}\left| {\bf r}_{t}\cdot \left( {\bf r}_{u}\times {\bf r}_{w}\right) \right|dtdudw= \frac{47}{16}\pi.\]

%Al expresar la ecuación del parabolide en coordenadas polares,
%$z = \frac{1}{3}\left( 9-(r\cos \theta +2)^{2}-(r\sin \theta +1)^{2}\right)=\frac{4}{3}-\frac{4}{3}r\cos \theta -\frac{1}{3}r^{2}-\frac{2}{3}r\sin \theta ,$
El volumen también puede hallarse evaluando la siguiente integral:
\[
\bigintss_{0}^{2\pi}
\bigintss_{0}^{1 - \cos \theta}
\bigintss_{0}^{\frac{4}{3} - \frac{4}{3}r\cos \theta - \frac{1}{3}r^{2} - \frac{2}{3}r\sin \theta}
r\, dz\, dr\, d\theta = \frac{47}{16}\pi
\]

Los cálculos anteriores se pueden obtener en \verb"Mathematica", definiendo $x$ e
$y,$ como la parametrización de la cardioide. Luego definir la parametrización $r$ del sólido, hallar el producto vectorial triple se denomina $A$ y se integra con los límites descritos anteriormente.


\vspace{0.6em}
\noindent
\begin{tcolorbox}[breakable,enhanced,title=\bf Instrucción en \verb"Mathematica", top=2pt,bottom=2pt,boxsep=0pt,before skip=2pt,after skip=0pt]
\begin{Verbatim}[formatcom=\letraverbatim]
{x,y}={u*(1-Cos[t])*Cos[t],u*(1-Cos[t])*Sin[t]};
r=(1-w)*{x,y,0}+w*{x,y,(9-(x+2)^2-(y+1)^2)/3};
A=Simplify[Dot[D[r,t],Cross[D[r,u],D[r,w]]]];
Integrate[Abs[A],{t,0,2*Pi},{u,0,1},{w,0,1}]
Integrate[r,{\[Theta],0,2*Pi},{r,0,1-Cos[\[Theta]]},
{z,0,4/3-4/3*r*Cos[\[Theta]]
r^2/3-2/3*r*Sin[\[Theta]]}]
\end{Verbatim}
\end{tcolorbox}
\vspace{0.7em}
\noindent


\section{Sólidos acotados por otras superficies}
%\vskip0.9cm
\subsection*{Superficie exponencial y  paraboloide}

Las superficies definidas por las ecuaciones $z=x^{2}+y^{2}$ y $z=e^{1-x^{2}-y^{2}}$ el sólido que se muestra a continuación:

\begin{figure}[H]
\centering
\caption{Sólido definido por $x^{2}+y^{2} \leq z \leq e^{1-x^{2}-y^{2}}$.}
\label{fig:paraboloide_exponencial}
\includegraphics[height=4.5cm]{Img/54a.pdf} \hspace{0.2cm}
\includegraphics[height=4.5cm]{Img/54b.pdf} \hspace{0.2cm}
\includegraphics[height=4.5cm]{Img/54c.pdf}
\fuentepropia
\end{figure}

Estos gráficos se obtienen por medio de las siguientes instrucciones.


\vspace{0.6em}
\noindent
\begin{tcolorbox}[breakable,enhanced,title=\bf Instrucción en \verb"Mathematica", top=2pt,bottom=2pt,boxsep=0pt,before skip=2pt,after skip=0pt]
\begin{Verbatim}[formatcom=\letraverbatim]
RegionPlot3D[x^2+y^2<z<Exp[1-x^2-y^2]&&-Sqrt[1-x^2]<y<
Sqrt[1-x^2]&&-1<x<1,{x,-1,1},{y,-1,1},{z,0,3},
PlotPoints->90,PlotStyle->{Blend[{Orange,Red},1/4]},
Mesh->False]
ContourPlot3D[{z==x^2+y^2,z==Exp[1-x^2-y^2]},{x,-2,2},
{y,-2,2},{z,0,2.8},ContourStyle->{Directive[Orange,
Opacity[0.7]],Directive[RGBColor[0.6,0.1,0.1],
Opacity[0.6]]},Mesh->False,BoxRatios->{2,2,1}]
\end{Verbatim}
\end{tcolorbox}
\vspace{0.7em}
\noindent


La intersección de las superficies se da en $z=1,$
por lo tanto, en el plano $xy,$ la intersección es el conjunto de puntos que satisface la ecuación $x^{2}+y^{2}=1.$


\noindent
\begin{minipage}{\textwidth}
	\centering
	\captionof{figure}{Proyección del sólido en los planos $xz$ y $xy$.}
	\label{fig:proyecciones_xz_xy}
	\begin{minipage}{0.30\linewidth}
		\centering
		\vspace{-15mm}
		\includegraphics[width=\linewidth]{Fig/471.pdf}
	\end{minipage}
	\hspace{1em}
	\begin{minipage}{0.30\linewidth}
		\centering
		\vspace{-15mm}
		\includegraphics[width=\linewidth]{Fig/472.pdf}
	\end{minipage}
	
	\vspace{-10mm}
	\fuentepropia
\end{minipage}


\vskip0.2cm
Ahora presentamos las integrales que se requieren evaluar para hallar el volumen del mencionado sólido.
%\vskip0.2cm
{\small
\[\int_{-1}^{1}\int_{-\sqrt{1-x^{2}}}^{\sqrt{1-x^{2}}}\int_{x^{2}+y^{2}}^{e^{1-x^{2}-y^{2}}}dzdydx= \int_{-1}^{1}\int_{-\sqrt{1-y^{2}}}^{\sqrt{1-y^{2}}}\int_{x^{2}+y^{2}}^{e^{1-x^{2}-y^{2}}}dzdxdy \approx 3.8273.\]
}
%\vskip0.2cm

Al despejar $y,$ de las expresiones que determinan las superficies e integrar primero con respecto a $y,$ se obtiene por un lado que {\small $y=\pm \sqrt{1-x^{2}-\ln z}$} y {\small $y=\pm \sqrt{z-x^{2}},$} por el otro. Entonces, el volumen está dado por
%\vskip0.2cm
{\small
\[\int_{0}^{1}\int_{-\sqrt{z}}^{\sqrt{z}}\int_{-\sqrt{z-x^{2}}}^{\sqrt{z-x^{2}}}dydxdz+\int_{1}^{e}\int_{-\sqrt{1-\ln z}}^{\sqrt{1-\ln z}}\int_{-\sqrt{1-x^{2}-\ln z}}^{\sqrt{1-x^{2}-\ln z}}dydxdz=\frac{1}{2}\pi ( 2e-3)\]
}
\[\int_{-1}^{1}\int_{x^{2}}^{1}\int_{-\sqrt{z-x^{2}}}^{\sqrt{z-x^{2}}}dydzdx+\int_{-1}^{1}\int_{1}^{e^{1-x^{2}}}\int_{-\sqrt{1-x^{2}-\ln z}}^{\sqrt{1-x^{2}-\ln z}}dydzdx \approx 3.8273.\]
%\vskip0.2cm
Por la simetría de la figura, al intercambiar \(x\) con \(y\) en las ecuaciones que determinan las superficies, es evidente que se obtiene:
%\vskip0.2cm

{\small
\[
\int_{0}^{1} \int_{-\sqrt{z}}^{\sqrt{z}} \int_{-\sqrt{z - y^{2}}}^{\sqrt{z - y^{2}}} dx\, dy\, dz
+ \int_{1}^{e} \int_{-\sqrt{1 - \ln z}}^{\sqrt{1 - \ln z}} \int_{-\sqrt{1 - y^{2} - \ln z}}^{\sqrt{1 - y^{2} - \ln z}} dx\, dy\, dz
= \pi e - \frac{3}{2}\pi
\]
}
\[
\int_{-1}^{1} \int_{y^{2}}^{1} \int_{-\sqrt{z - y^{2}}}^{\sqrt{z - y^{2}}} dx\, dz\, dy
+ \int_{-1}^{1} \int_{1}^{e^{1 - y^{2}}} \int_{-\sqrt{1 - y^{2} - \ln z}}^{\sqrt{1 - y^{2} - \ln z}} dx\, dz\, dy
\approx 3.8273.
\]
\vskip0.2cm
Algunas de estas integrales se evaluaron usando \verb|Mathematica|, las indicaciones son las siguientes:


\vspace{0.6em}
\noindent
\begin{tcolorbox}[breakable,enhanced,title=\bf Instrucción en \verb"Mathematica", top=2pt,bottom=2pt,boxsep=0pt,before skip=2pt,after skip=0pt]
\begin{Verbatim}[formatcom=\letraverbatim]
NIntegrate[1,{x,-1,1},{y,-Sqrt[1-x^2],Sqrt[1-x^2]},
{z,x^2+y^2,Exp[1-x^2-y^2]}]
Integrate[1,{y,-1,1},{x,-Sqrt[1-y^2],Sqrt[1-y^2]},
{z,x^2+y^2,Exp[1-x^2-y^2]}]
Integrate[1,{z,0,1},{x,-Sqrt[z],Sqrt[z]},
{y,-Sqrt[z-x^2],Sqrt[z-x^2]}]+Integrate[1,
{z,1,Exp[1]},{x,-Sqrt[1-Log[z]],Sqrt[1-Log[z]]},
{y,-Sqrt[1-x^2-Log[z]],Sqrt[1-x^2-Log[z]]}]
NIntegrate[1,{y,-1,1},{z,y^2,1},{x,-Sqrt[z-y^2],
Sqrt[z-y^2]}]+NIntegrate[1,{y,-1,1},{z,1,Exp[1-y^2]},
{x,-Sqrt[1-Log[z]-y^2],Sqrt[1-Log[z]-y^2]}]
\end{Verbatim}
\end{tcolorbox}
\vspace{0.7em}
\noindent
\vskip0.2cm


Utilizando el siguiente proceso también aporta el mismo valor.

%\begin{tcolorbox}[title=\bfseries Parametrizando el sólido]
	
La región de integración en el plano $xy$ es el círculo unitario, un punto en este conjunto se determina por {\small $(u\cos t,u\sin t),$} para $u\in [0,1]$ y {\small $t\in [0,2\pi].$}
Un punto sobre el paraboloide posee coordenadas {\small $P( u\cos t,u\sin t,u^{2}),$}
y un punto sobre la superficie exponencial tiene la forma
{\small $Q(u\cos t,u\sin t,e^{1-u^{2}}).$} Con una combinación convexa de estos puntos se parametriza el sólido en mención.
\begin{multline*}
{\bf r}( u,t,w) = (1-w) Q +wP = \langle u\cos t,u\sin t,e^{1-u^{2}}( 1-w) +wu^{2}\rangle,
\end{multline*}
con $u\in [0,1]$, $w\in [0,1]$ y $t\in [0,2\pi ].$

%\vskip0.2cm
De esta función se obtiene
\begin{eqnarray*}
{\bf r}_{t} &=& \langle -u\sin t,u\cos t,0\rangle \\
{\bf r}_{u} &=&\langle \cos t,\sin t,-2( 1-w)
ue^{1-u^{2}}+2wu\rangle \\
{\bf r}_{w}&=& \langle 0,0,u^{2}-e^{1-u^{2}}\rangle \\
|{\bf r}_{t}\cdot ({\bf r}_{u}\times {\bf r}_{w})| &=& |u\exp(1-u^{2})-u^{3}|=ue^{1-u^{2}}-u^{3}.
\end{eqnarray*}
%\vskip0.2cm
El volumen del sólido se halla evaluando la siguiente integral
%\vskip0.2cm
\[\int_{0}^{1}\int_{0}^{1}\int_{0}^{2\pi }( ue^{1-u^{2}}-u^{3})
dtdwdu= \pi e-\frac{3}{2}\pi.\]
%\vskip0.2cm
%\end{tcolorbox}

Con el cambio a coordenadas cilíndricas, la integral respectiva es
\[\int_{0}^{2\pi }\int_{0}^{1}\int_{r^{2}}^{e^{1-r^{2}}}rdzdrd\theta
= \frac{1}{2}\pi ( 2e-3).\]
%\vskip0.2cm
Con las siguientes instrucciones se obtuvieron los cálculos anteriores.


\vspace{0.6em}
\noindent
\begin{tcolorbox}[breakable,enhanced,title=\bf Instrucción en \verb"Mathematica", top=2pt,bottom=2pt,boxsep=0pt,before skip=2pt,after skip=0pt]
\begin{Verbatim}[formatcom=\letraverbatim]
{x,y}={u*Cos[t],u*Sin[t]}
r=Simplify[(1-w)*{x,y,x^2+y^2}+w*{x,y,Exp[1-x^2-y^2]}]
A=Simplify[Dot[D[r,w],Cross[D[r,u],D[r,t]]]]
Integrate[A,{u,0,1},{t,0,2*Pi},{w,0,1}]
Integrate[r,{\[Theta],0,2*Pi},{r,0,1},
{z,r^2,Exp[1-r^2]}]
\end{Verbatim}
\end{tcolorbox}
\vspace{0.7em}
\noindent


\subsection*{Una semiesfera y un cilindro} \index{Semiesfera}
Hallar el volumen del sólido determinado por las siguientes expresiones.
$x^2 +y^2 +z^2 \leq 9,$ $ y^2 +z^2 \geq 1,$ $z \leq 0.$ La figura se forma cuando una esfera de radio $3$ se perfora por un cilindro de radio $1$ cuyo eje de simetría es una recta que pasa por el centro de la esfera; luego, un plano que contiene al eje de simetría del cilindro divide esta región del espacio en dos porciones iguales y se toma una de estas.

\vskip0.2cm
\noindent
\begin{minipage}{\textwidth}
	\centering
	\captionof{figure}{Sólido definido por $1 \leq y^2 + z^2 \leq 9 - x^2$ \quad y \quad $z \geq 0$.}
	\label{fig:cilindro_parabolico_superior}
	\begin{minipage}{0.25\linewidth}
		\centering
		\includegraphics[width=\linewidth]{Img/56c.pdf}
	\end{minipage}
	\hspace{1em}
	\begin{minipage}{0.30\linewidth}
		\centering
		\includegraphics[width=\linewidth]{Img/56d.pdf}
	\end{minipage}
	
	\vspace{0.3em}
	\fuentepropia
\end{minipage}

\vskip0.2cm
Las siguientes instrucciones permiten obtener el gráfico del sólido.


\vspace{0.6em}
\noindent
\begin{tcolorbox}[breakable,enhanced,title=\bf Instrucción en \verb"Mathematica", top=2pt,bottom=2pt,boxsep=0pt,before skip=2pt,after skip=0pt]
\begin{Verbatim}[formatcom=\letraverbatim]
RegionPlot3D[x^2+y^2+z^2<9&&y^2+z^2>1&&-3<z<0,
{x,-3,3},{y,-3,3},{z,-3,2},PlotPoints->90,Mesh->False,
PlotStyle->Orange,Ticks->{{-4,-2,0,2,4},{-3,-1,1,3},
{-2,0,2}},AxesLabel->{x,y,z}]
\end{Verbatim}
\end{tcolorbox}
\vspace{0.7em}
\noindent


Con el fin de plantear las integrales que determinan el volumen del sólido, se presentan las siguientes figuras.

\vskip0.2cm
\noindent
\begin{minipage}{\textwidth}
	\centering
	\captionof{figure}{Proyección del sólido en los planos coordenados.}
	\label{figsemicirc}
	\begin{minipage}{0.28\linewidth}
		\centering
		\vspace{-10mm}
		\includegraphics[width=\linewidth]{Fig/473.pdf}
	\end{minipage}
	\hspace{1em}
	\begin{minipage}{0.28\linewidth}
		\centering
		\vspace{-10mm}
		\includegraphics[width=\linewidth]{Fig/474.pdf}
	\end{minipage}
	\hspace{1em}
	\begin{minipage}{0.28\linewidth}
		\centering
		\vspace{-10mm}
		\includegraphics[width=\linewidth]{Fig/475.pdf}
	\end{minipage}
	
	\vspace{-4mm}
	\fuentepropia
\end{minipage}


\vskip0.2cm
El volumen del sólido que se considera está dado por las siguientes integrales.
En el orden $dxdzdy$ se debe observar que $x$ recorre puntos desde la parte negativa de la esfera hasta su parte positiva.
%\vskip0.2cm
\vspace{-2mm}
\begin{multline*}
\int_{-3}^{-1}\int_{-\sqrt{9-y^2 }}^{0}\int_{-\sqrt{9-y^2 -z^2 }}^{\sqrt{9-y^2 -z^2 }}dxdzdy+\int_{-1}^{1}\int_{-\sqrt{9-y^2 }}^{-\sqrt{1-y^2 }}\int_{-\sqrt{9-y^2 -z^2 }}^{\sqrt{9-y^2 -z^2 }}dxdzdy\\
+\int_{1}^{3}\int_{-\sqrt{9-y^2 }}^{0}\int_{-\sqrt{9-y^2 -z^2 }}^{\sqrt{9-y^2 -z^2 }}dxdzdy = \frac{32}{3}\sqrt{2}\pi
\end{multline*}
%\vskip0.2cm
\vspace{-7mm}
\begin{multline*}
\int_{-1}^{0}\int_{-\sqrt{9-z^2 }}^{-\sqrt{1-z^2 }}\int_{-\sqrt{9-y^2 -z^2 }}^{\sqrt{9-y^2 -z^2 }}dxdydz+\int_{-3}^{-1}\int_{-\sqrt{9-z^2 }}^{\sqrt{9-z^2 }}\int_{-\sqrt{9-y^2 -z^2 }}^{\sqrt{9-y^2 -z^2}}dxdydz\\
+\int_{-1}^{0}\int_{\sqrt{1-z^2 }}^{\sqrt{9-z^2 }}\int_{-\sqrt{9-y^2 -z^2 }}^{\sqrt{9-y^2 -z^2 }}dxdydz= \frac{32}{3}\sqrt{2}\pi
\end{multline*}

%\vskip0.2cm

En el plano $xz$ hay una porción en la región de integración que corresponde a un rectángulo y existe una región en el plano $xz$ que no contiene puntos $x$ en $[-3,\sqrt{8}]$ y $[-3,\sqrt{8}],$ entre otros.
%\vskip0.2cm
\begin{multline*}
\int_{-1}^{0}\int_{-\sqrt{8}}^{\sqrt{8}}\int_{-\sqrt{9-x^{2}-z^{2}}}^{-\sqrt{1-z^{2}}}dydxdz
+\int_{-1}^{0}\int_{-\sqrt{8}}^{\sqrt{8}}\int_{\sqrt{1-z^{2}}}^{\sqrt{9-x^{2}-z^{2}}}dydxdz\\
+\int_{-3}^{-1}\int_{-\sqrt{9-z^{2}}}^{\sqrt{9-z^{2}}}\int_{-\sqrt{9-x^{2}-z^{2}}}^{\sqrt{9-x^{2}-z^{2}}}dydxdz = \frac{32}{3}\sqrt{2}\pi
\end{multline*}
%\vskip0.2cm
\vspace{-7mm}
\begin{multline*}
\int_{-\sqrt{8}}^{\sqrt{8}} \int_{-1}^{0}\int_{-\sqrt{9-x^{2}-z^{2}}}^{-\sqrt{1-z^{2}}}dydzdx
+\int_{-\sqrt{8}}^{\sqrt{8}}\int_{-1}^{0}\int_{\sqrt{1-z^{2}}}^{\sqrt{9-x^{2}-z^{2}}}dydzdx \\
+\int_{-\sqrt{8}}^{\sqrt{8}}\int_{-\sqrt{9-x^{2}}}^{-1}\int_{-\sqrt{9-x^{2}-z^{2}}}^{\sqrt{9-x^{2}-z^{2}}}dydzdx = \frac{32}{3}\sqrt{2}\pi.
\end{multline*}
%\vskip0.2cm
Integrando primero, con respecto a $z,$ existe en el plano $xy$ una región donde $z$ recorre los puntos desde la parte negativa de la esfera al plano
$z=0$ y otra región donde $z$ recorre los puntos de la parte negativa de la esfera hasta la parte negativa del cilindro con ecuación $y^2+z^2=1.$
\vspace{-2mm}
{\small
\begin{multline*}
\int_{1}^{3}\int_{-\sqrt{9-y^2}}^{\sqrt{9-y^2}}\int_{-\sqrt{9-x^2 -y^2 }}^{0}dzdxdy+\int_{-3}^{-1}\int_{-\sqrt{9-y^2 }}^{\sqrt{9-y^2 }}\int_{-\sqrt{9-x^2 -y^2 }}^{0}dzdxdy\\
+\int_{-1}^{1}\int_{-\sqrt{8}}^{\sqrt{8}}\int_{-\sqrt{9-x^2 -y^2 }}^{-\sqrt{1-y^2 }}dzdxdy= \frac{32}{3}\sqrt{2}\pi
\end{multline*}
}
\vspace{-9mm}
{\small
\begin{multline*}
\int_{-\sqrt{8}}^{\sqrt{8}}\int_{-\sqrt{9-x^2 }}^{-1}\int_{-\sqrt{9-x^2 -y^2 }}^{0}dzdydx+\int_{-\sqrt{8}}^{\sqrt{8}}\int_{1}^{\sqrt{9-x^2 }}\int_{-\sqrt{9-x^2 -y^2 }}^{0}dzdydx\\
+\int_{-\sqrt{8}}^{\sqrt{8}}\int_{-1}^{1}\int_{-\sqrt{9-x^2 -y^2 }}^{-\sqrt{1-y^2 }}dzdydx= \frac{32}{3}\sqrt{2}\pi.
\end{multline*}
}
%\vskip0.2cm
Al hacer un cambio a coordenadas cilíndricas
$y = r\cos \theta,\, z = r\sin \theta,$, el volumen está dado por.
%\vskip0.2cm
{\small
\[\int_{\pi}^{2\pi}\int_{1}^{3}\int_{-\sqrt{9-r^2}}^{\sqrt{9-r^2}}rdxdrd\theta=\frac{32}{3}\sqrt{2}\pi.\]
}
%\vskip0.2cm
Algunos de los cálculos anteriores se pueden evaluar usando \verb"Mathematica"
mediante la ejecución de las siguientes instrucciones.



\vspace{0.6em}
\noindent
\begin{tcolorbox}[breakable,enhanced,title=\bf Instrucción en \verb"Mathematica", top=2pt,bottom=2pt,boxsep=0pt,before skip=2pt,after skip=0pt]
\begin{Verbatim}[formatcom=\letraverbatim]
Integrate[r,{t,Pi,2*Pi},{r,1,3},{x,-Sqrt[9-r^2],
Sqrt[9-r^2]}]
Integrate[1,{y,-3,-1},{z,-Sqrt[9-y^2],0},
{x,-Sqrt[9-y^2-z^2],Sqrt[9-y^2-z^2]}]+Integrate[1,
{y,-1,1},{z,-Sqrt[9-y^2],-Sqrt[1-y^2]},{x,
Sqrt[9-y^2-z^2],Sqrt[9-y^2-z^2]}]+Integrate[1,
{y,1,3},{z,-Sqrt[9-y^2],0},{x,-Sqrt[9-y^2-z^2],
Sqrt[9-y^2-z^2]}]

Integrate[1,{z,-1,0},{x,-Sqrt[8],Sqrt[8]},
{y,-Sqrt[9-x^2-z^2],-Sqrt[1-z^2]}]+Integrate[1,
{z,-1,0},{x,-Sqrt[8],Sqrt[8]},{y,Sqrt[1-z^2],
Sqrt[9-x^2-z^2]}]+Integrate[1,{z,-3,-1},{x,
Sqrt[9-z^2],Sqrt[9-z^2]},{y,-Sqrt[9-x^2-z^2],
Sqrt[9-x^2-z^2]}]
NIntegrate[1,{x,-Sqrt[8],Sqrt[8]},{y,-Sqrt[9-x^2],-1},
{z,-Sqrt[9-x^2-z^2],0}]+NIntegrate[1,{x,-Sqrt[8],
Sqrt[8]},{y,1,Sqrt[9-x^2]},{z,-Sqrt[9-x^2-z^2],0}]
+NIntegrate[1,{x,-Sqrt[8],Sqrt[8]},{y,-1,1},
{z,-Sqrt[9-x^2-z^2],-Sqrt[1-y^2]}]
\end{Verbatim}
\end{tcolorbox}
\vspace{0.7em}
\noindent


El volumen también se obtiene mediante el siguiente procedimiento.

%\begin{tcolorbox}[title=\bfseries Parametrizando el sólido]
Consideremos la proyección del sólido en el plano $yz,$ (véase la figura \ref{figsemicirc}). Un punto sobre la semicircunferencia de radio $1,$ posee coordenadas
$( \cos( t) ,\sin ( t) ) $ y un punto sobre la semicircunferencia de radio $3,$ es de la forma
$( 3\cos ( t),3\sin ( t) ) ,$ con $t \in [ \pi ,2\pi ].$
Por lo tanto, una combinación convexa de estos dos puntos parametriza la región en mención:
$( 1-u) ( \cos ( t) ,\sin ( t) )+u( 3\cos ( t) ,3\sin ( t) ) $
que se simplifica como $ ( ( 1+2u) \cos t,( 1+2u) \sin t) .$
\vskip0.2cm
Sobre el sólido, $x$ recorre los puntos que van desde la parte negativa de la esfera hasta la parte positiva de la misma. Dos puntos arbitrarios
$P$ y $Q$ en estos dos lugares tienen coordenadas
$P=(-2\sqrt{( u+2)( 1-u) },(1+2u)\cos t,(1+2u)\sin t), \,\, \text{y} $
$Q=(2\sqrt{( u+2) ( 1-u) },(1+2u)\cos t,(1+2u)\sin t).$
\vskip0.2cm
Es posible parametrizar el sólido mediante una combinación convexa de la forma $(1-w)P+wQ$, esto define la función ${\bf r}.$
%\vskip0.2cm
\vspace{-2mm}
{\small
\begin{multline*}
{\bf r}(t,u,w) = \big \langle 2( 2w-1) \sqrt{( u+2)( 1-u) },\cos t+2u\cos t, \\ \sin t+2u\sin t \big \rangle,
\end{multline*}
}
en donde $t \in [0,2\pi],$ $u \in [0,1]$ y $w \in [0,1].$
Al derivar parcialmente con respecto a $t$, $u$ y $w$ hallamos respectivamente, los siguientes vectores:
%\vskip0.2cm
\vspace{-2mm}
{\small
\begin{eqnarray*}
{\bf r}_{t} &=& \langle 0,-\sin t-2u\sin t,\cos t+2u\cos t\rangle \\
{\bf r}_{u} &=& \left \langle \frac{( 1-2w)( 1+2u)}{\sqrt{( u+2) ( 1-u) }},2\cos t,2\sin t\right \rangle \\
{\bf r}_{w} &=& ( 4\sqrt{( u+2) ( 1-u) },0,0).
\end{eqnarray*}
}
%\vskip0.2cm
\vspace{-2mm}
Su triple producto vectorial triple es
{\small
\[|{\bf r}_{t}\cdot({\bf r}_{u}\times{\bf r}_{w})|=8(1+2u)\sqrt{(u+2)(1-u)};\]
}
%\vskip0.2cm
al integrarlo con los límites definidos, nos lleva al volumen del
sólido
%\vskip0.2cm
{\small
\[\int_{0}^{1}\int_{0}^{1}\int_{\pi}^{2\pi}8(1+2u)\sqrt{(u+2)(1-u)}\,dtdudw=\frac{32}{3}\sqrt{2}\pi.\]
}
%\end{tcolorbox}

\subsection*{Sólidos cardióidicos} \index{Cardioide} \index{Sólidos!cardioidicos}

En esta sección consideraremos el sólido generado por la rotación de una cardioide alrededor de una recta exterior a ella.
Compararemos los valores obtenidos parametrizando el sólido y comparándolos con el valor obtenido aplicando el teorema de Pappus o la conocida formulación para superficies de revolución. \index{Teorema de Pappus!área superficial}

\index{Toro cardioidico!volumen}
\subsection*{Toro cardióidico generado por $\bm{\mathsf{r=1-\sin t}}$}

Para $0\leq t \leq 2\pi,$, la cardioide con ecuación $r=1-\sin t$ su área está dada por:
\[\int_{0}^{2\pi }\int_{0}^{1-\sin \theta}rdrd\theta = \frac{3}{2}\pi.\]
El volumen del sólido generado cuando la curva parametrizada en la forma {\small $\langle x(t),y(t)\rangle,$} con {\small $t \in [a,b]$}, gira alrededor del eje $y,$
está dado por {\small $\ds \pi \int_{a}^{b}\left( x(t)\right) ^{2}y^{\prime }(t)dt.$}
El volumen está dado por {\small $\ds\pi \int_{a}^{b}\left( y(t)\right) ^{2}x^{\prime }(t)dt,$}
si el giro se hace alrededor del eje $x.$
En el caso de la cardioide $r=1-\sin \theta$ que se desplazada tres unidades sobre el eje polar, esta curva se puede parametrizar como {\small $x(t)=3+\left( 1-\sin (t) \right) \cos (t),$}
{\small $y(t)=\left( 1-\sin (t) \right) \sin (t) ,$} {\small $0\leq t \leq 2\pi.$} Entonces el volumen $V$ del sólido que se obtiene al girar la cardiode sobre la recta $\theta = \pi/2,$ está dado por
{\small \[V=\pi \int_{0}^{2\pi }\left( 3+\left( 1-\sin (t) \right) \cos(t) \right) ^{2}(\cos t-2\cos t\sin t)dt= 9\pi ^{2}.\]
}
\vspace{-6mm}
%\vskip0.2cm
\begin{figure}[H]
\centering
\caption{Toro cardióidico generado por $r = 1 - \sin \theta$.}
\label{fig:toro_cardioidico}
\includegraphics[height=4.5cm]{Img/63d.pdf} \hspace{0.4cm}
\includegraphics[height=4.5cm]{Img/63e.pdf}
\fuentepropia
\end{figure}

Las coordenadas del centro de masa, $(\overline{x},\overline{y})$ están determinadas por
\[\frac{2}{3\pi }\int_{0}^{2\pi }\int_{0}^{1-\sin \theta }\left( 3+r\cos \theta \right) rdrd\theta = 3,\]
%\vskip0.2cm
\[\frac{2}{3\pi }\int_{0}^{2\pi }\int_{0}^{1-\sin \theta }\left( r\sin \theta \right) rdrd\theta = -\frac{5}{6}.\]
El volumen de este sólido, usando el teorema de Pappus, es \index{Teorema de Pappus!volumen}
$2\pi \left( \frac{3}{2}\pi \right) \left( 3\right) = 9\pi ^{2}.$

% \begin{eqnarray*} \index{Centro de masa} \overline{x} &=& \frac{2}{3\pi }\int_{0}^{2\pi }\int_{0}^{1-\sin \theta }\left( 3+r\cos \theta \right) rdrd\theta = 3 \\ \overline{y}&=& \frac{2}{3\pi }\int_{0}^{2\pi }\int_{0}^{1-\sin \theta }\left( r\sin \theta \right) rdrd\theta = -\frac{5}{6}
% \end{eqnarray*}

%\begin{tcolorbox}[title=\bfseries Parametrizando el sólido]
%\footnotesize 
%\parskip=0pt 
%\setlength{\belowdisplayskip}{0pt} 
%\setlength{\abovedisplayskip}{0pt} 
\vskip0.2cm
La cardioide se rellena usando la parametrización
%\vskip0.2cm
\[\langle 3+u\left( 1-\sin (t) \right) \cos (t),u\left( 1-\sin (t) \right) \sin (t) \rangle,\]
%\vskip0.2cm
con $t\in [0,2\pi]$ y $u\in [0,1]$.
Agregando a la primera componente $\cos(v)$ y $\sin(v)$, se obtiene una parametrización de la región rotada,
\begin{multline*}
{\bf r}(t,u,v)=\bigg \langle\left ( 3+u\left( 1-\sin (t) \right) \cos (t) \right) \cos (v) , \\ \left( 3+u\left( 1-\sin \left(
t\right) \right) \cos (t) \right) \sin (v) ,u\left(1-\sin (t) \right) \sin (t) \bigg \rangle,
\end{multline*}
\vskip0.2cm
con $t\in [0,2\pi ],\,\, u\in [0,1]$ y $v\in [0,2\pi].$
Derivando parcialmente se obtienen los vectores: 
\vskip0.2cm
${\bf r}_{u} = ( 1-\sin t) \langle \cos t\cos v, \cos t\sin v, \sin t\rangle ,$ \\
${\bf r}_{v} = (3+u(1-\sin t) \cos t) \langle -
\sin v,\cos t) \cos v,0\rangle ,$ \\
${\bf r}_{t} = \bigg \langle -u( \cos \left( 2t\right) +\sin t) \cos
v,-u\left( \cos \left( 2t\right) +\sin t\right) \sin v, \\
u\left( 1-2\sin t\right) \cos t \bigg \rangle.$
\vskip0.2cm
Con ellos, se obtiene el elemento diferencial de volumen
$|{\bf r}_{t}\cdot \left( {\bf r}_{u}\times {\bf r}_{v}\right)| =
|u\left( 6-3\cos ^{2}t+u\left(\sin t-3\right) \cos ^{3}t+4u(1-\sin t)\cos t-6\sin t\right)|,$
cuya integración nos proporciona el volumen del sólido
%\vskip0.2cm
\[\int_{0}^{2\pi }\int_{0}^{1}\int_{0}^{2\pi} |{\bf r}_{t}\cdot \left( {\bf r}_{u}\times {\bf r}_{v}\right)|
dtdudv= 9\pi ^{2}.\]
%\end{tcolorbox}

Si la densidad del solido es $\delta \left( x,y,z\right) =z^{2},$ entonces la masa del sólido es:
%\vskip0.2cm
\begin{equation*}
\int_{0}^{2\pi }\int_{0}^{1}\int_{0}^{2\pi}\left(u\left(1-\sin(t)\right)\sin(t)\right)^{2}|{\bf r}_{t}\cdot\left({\bf r}_{u}\times {\bf r}_{v}\right)|dtdudv= \frac{147}{16}\pi ^{2}.
\end{equation*}
\vskip0.2cm
Los cálculos anteriores se evaluan con las siguientes instrucciones: %en \verb"Mathematica" con ayuda de la siguiente rutina


\vspace{0.6em}
\noindent
\begin{tcolorbox}[breakable,enhanced,title=\bf Instrucción \verb"Mathematica", top=2pt,bottom=2pt,boxsep=0pt,before skip=2pt,after skip=0pt]
\begin{Verbatim}[formatcom=\letraverbatim]
{x,y,z}={3+u*(1-Sin[t])*Cos[t],(1-u*Sin[t])*Sin[t],y};
r={x*Cos[v],x*Sin[v],z};
A=Simplify[Dot[D[r,t],Cross[D[r,u],D[r,v]]]]
Integrate[A,{t,0,2*Pi},{u,0,1},{v,0,2*Pi}]
Integrate[z^2*A,{t,0,2*Pi},{u,0,1},{v,0,2*Pi}]
\end{Verbatim}
\end{tcolorbox}
\vspace{0.7em}
\noindent



\subsection*{Toro cardióidico generado por $\bm{\mathsf{r=1-\cos t}}$}

Consideraremos la cardioide con ecuación $r=1-\cos(t),$ $0\leq t \leq 2\pi.$
En el plano polar, si esta curva se desplaza tres unidades sobre el eje polar, y luego rota alrededor de una recta paralela a
$\theta=\pi/2$, se obtiene el sólido $E$.

\begin{figure}[H]
\centering
\caption{Toro cardioidico generado por $r = 1 - \cos \theta$.}
\label{fig:toro_cardioidico_cos}
\includegraphics[height=4.5cm]{Img/59f.pdf}
\hspace{0.4cm}
\includegraphics[height=4.5cm]{Img/59e.pdf}
\fuentepropia
\end{figure}

El área de la cardioide es
%\vskip0.2cm
\[\int_{0}^{2\pi }\int_{0}^{1-\cos \theta }rdrd\theta = \frac{3}{2}\pi,\]
%\vskip0.2cm
y las coordenadas del centro de masa están dadas por
%\vskip0.2cm
\begin{eqnarray*}
\overline{x} &=& \frac{2}{3\pi}\int_{0}^{2\pi}\int_{0}^{1-\cos \theta }\left( 3+r\cos \theta \right) rdrd\theta =\frac{13}{6} \\
\overline{y} &=& \frac{2}{3\pi}\int_0^{2\pi}\int_0^{1-\cos\theta}(r\sin\theta)rdrd\theta = 0.
\end{eqnarray*}
\vskip0.2cm
De acuerdo con el teorema de Pappus, el volumen del toro cardioidico generado por
$r=1-\cos t$, es:
\[2\pi \left( \frac{13}{6}\right) \frac{3}{2}\pi =\frac{13}{2}\pi ^{2}.\]
\vskip0.2cm
La cardioide se parametriza en el plano $xy,$ como
\[x(t)=3+\left( 1-\cos (t) \right) \cos (t),\quad y(t)=\left( 1-\cos (t) \right) \sin (t),\quad t\in[0,2\pi].\]

El volumen $V,$ del sólido generado por la revolución de esta región sobre la recta $\theta = \pi/2,$ está dado por
\[V=\pi \int_{0}^{2\pi }\left( 3+\left( 1-\cos (t) \right) \cos (t) \right) ^{2}(\cos t-2\cos ^{2}t+1)dt =
\frac{13}{2}\pi ^{2}.\]

%\begin{tcolorbox}[title=\bfseries Parametrizando el sólido]
En el plano $xy$, para rellenar la cardioide se usa la parametrización
${\bf r}_1(t,u)=\langle 3+u\left( 1-\cos (t) \right) \cos (t),u\left( 1-\cos (t) \right) \sin (t) \rangle ,$ con $t\in [0,2\pi ],\,\, u\in [0,1].$
Si esta región gira sobre la recta $x=3$, se genera un toro cardioidico, cualquier punto sobre este sólido tiene la forma determinada por la si\-guiente función
\begin{multline*}
{\bf r}(t,u,v)=\bigg \langle \left( 3+u\left( 1-\cos (t) \right) \cos \left(
t\right) \right) \cos (v) , \\
\left( 3+u\left( 1-\cos \left(t\right) \right) \cos (t) \right) \sin (v) , u\left(1-\cos (t) \right) \sin (t) \bigg \rangle,
\end{multline*}
con $t\in [0,2\pi ],\,\,v\in [0,2\pi ],\,\, u\in [0,1].$
Al derivar parcialmente con respecto a $t$, $u$ y $v$ se obtiene el elemento diferencial de volumen:
\begin{multline*}
\left|{\bf r}_{t}\cdot \left( {\bf r}_{u}\times {\bf r}_{v}\right)\right| 
= \Big|(1+u)\cos^2(t)\left(6 - u + 2u\cos(t) - u\cos(2t)\right) \\
\times \sin^2\left( \frac{t}{2} \right)\Big|
\end{multline*}
Esto se integra y nos proporciona el volumen del toro cardióidico generado por la curva $r = 1 - \cos t$:
\[
\int_{0}^{2\pi }\int_{0}^{2\pi }\int_{0}^{1}
\left|{\bf r}_{t}\cdot \left( {\bf r}_{u}\times {\bf r}_{v}\right)\right| \, dudvdt
= \frac{13}{2}\pi ^{2}.
\]
%\end{tcolorbox}

En \verb"Mathematica" se pueden efectuar estos cálculos ejecutando las si\-guientes instrucciones:



\vspace{0.6em}
\noindent
\begin{tcolorbox}[breakable,enhanced,title=\bf Instrucción \verb"Mathematica", top=2pt,bottom=2pt,boxsep=0pt,before skip=2pt,after skip=0pt]
\begin{Verbatim}[formatcom=\letraverbatim]
R=(3+u*(1-Cos[t])*Cos[t]);
r={x,y,z}={R*Cos[v],R*Sin[v],u*(1-Cos[t])*Sin[t]};
A=Simplify[Dot[D[r,t],Cross[D[r,u],D[r,v]]]];
Integrate[Abs[A],{t,0,2*Pi},{v,0,2*Pi},{u,0,1}]
\end{Verbatim}
\end{tcolorbox}
\vspace{0.7em}
\noindent



\subsection*{Un sólido cardióidico} \index{Cardioide}

En el plano $xy$, el interior de la cardioide $r=1-\cos(t)$
se parametriza por medio de las expresiones
$x(t,u)=u\left( 1-\cos (t)\right)\cos(t)$, $y(t,u)=u\left( 1-\cos (t) \right) \sin (t), $
para $u \in [0,1]$ y $t \in [0,2\pi].$
\vskip0.2cm
Construiremos una superficie que en el plano $xy$ coincida con la cardioide mencionada.
La ecuación $r=1-\cos(t)$ corresponde en coordenadas rectangulares a la expresión
$(x^2+y^2)=(x^2+y^2+x)^2$; es por esto que al hacer
$z=(x^2+y^2)-(x^2+y^2+x)^2,$ se satisface la igualdad en $z=0$.
Sustituyéndose $x$ e $y$ en la expresión anterior se obtiene:
%\vskip0.2cm
\[z(t,u)= 8 (u-1) u^2 \sin ^6\left(\frac{t}{2}\right)((u-1)\cos (t)-u-1).\]
%\vskip0.2cm
La superficie parametrizada por las ecuaciones $x(t,u),$ $y(t,u)$ $z(t,u)$ se muestra a continuación:


\vskip0.2cm
\noindent
\begin{minipage}{\textwidth}
	\centering
	\captionof{figure}{Superficie definida por $z = (x^2 + y^2) - (x^2 + y^2 + x)^2$}
	\label{fig:superficie_compuesta}
	\begin{minipage}{0.26\linewidth}
		\centering
		\includegraphics[width=\linewidth]{Img/65a.pdf}
	\end{minipage}
	\begin{minipage}{0.26\linewidth}
		\centering
		\includegraphics[width=\linewidth]{Img/65b.pdf}
	\end{minipage}
	
	\vspace{0.3em}
	\fuentepropia
\end{minipage}
\vskip0.2cm


El volumen del sólido acotado por esta superficie y el plano $z=0$
se hallan definiendo la función ${\bf r}\left( t,u,v\right),$ que es una combinación convexa de dos puntos de la forma $(x(t,u),y(t,u),0)$ y $(x(t,u),y(t,u),z(t,u)).$ Esto es:
\begin{multline*}
{\bf r}\left( t,u,v\right) = \bigg \langle u(1-\cos(t))\cos(t),u\sin(t)(1-\cos(t)), \\
8(u-1)u^2v\sin ^6\left(\frac{t}{2}\right)((u-1)\cos(t)-u-1 \bigg \rangle,
\end{multline*}
\vskip0.2cm
en donde \( t \in [0,2\pi] \), \( u \in [0,1] \) y \( v \in [0,1] \).
\vskip0.2cm
Con esta función, hallamos el producto vectorial triple:
%\vskip0.2cm
\[
\left| \mathbf{r}_{t} \cdot \left( \mathbf{r}_{u} \times \mathbf{r}_{v} \right) \right| =
\left| 32 (1 - u) u^3 \sin^{10}\left( \frac{t}{2} \right) \left( (u - 1) \cos(t) - u - 1 \right) \right|,
\]
\vskip0.2cm
que se integra y nos proporciona el volumen del sólido en mención:
{\small
\[
\int_{0}^{1} \int_{0}^{1} \int_{0}^{2\pi}
\left| 32 (1 - u) u^3 \sin^{10}\left( \frac{t}{2} \right) \left( (u - 1) \cos(t) - u - 1 \right) \right|
\, dt\, du\, dv = \frac{35}{32}\pi.
\]
}

El volumen se puede plantear en coordenadas cilíndricas como sigue:
%\vskip0.2cm
\[
\int_{0}^{2\pi} \int_{0}^{1 - \cos \theta}
\int_{0}^{r^{2} \sin^{2}(\theta) - r^{4} - 2r^{3} \cos(\theta)}
r\, dz\, dr\, d\theta = \frac{35}{32}\pi.
\]
%\vskip0.2cm
Las figuras y las integrales se consiguen con las siguientes instrucciones:


\vspace{0.6em}
\noindent
\begin{tcolorbox}[breakable,enhanced,title=\bf Instrucción \verb"Mathematica", top=2pt,bottom=2pt,boxsep=0pt,before skip=2pt,after skip=0pt]
\begin{Verbatim}[formatcom=\letraverbatim]
{x,y}={u*(1-Cos[t])*Cos[t],u*(1-Cos[t])*Sin[t]};
z=(x^2+y^2)-(x^2+y^2+x)^2;
R=(1-v)*{x,y,0}+v*{x,y,z};
ParametricPlot3D[{x,y,z},{t,0,2*Pi},{u,0,1},
Ticks->{{-2,-1,0},{-1,0,1},{0,2}},MeshStyle->Gray,
PlotPoints->80,PlotStyle->Orange,PlotRange->{{-2,1/2},
{-3/2,3/2},{0,2}}]
A=Simplify[Abs[Dot[D[R,t],Cross[D[R,u],D[R,v]]]]];
Integrate[A,{t,0,2*Pi},{u,0,1},{v,0,1}]
Clear[r,z,t]
Integrate[r,{t,0,2*Pi},{r,0,1-Cos[t]},
{z,0,r^2-(r^2+r*Cos[t])^2}]
\end{Verbatim}
\end{tcolorbox}
\vspace{0.7em}
\noindent



\subsection*{Un cardioide revolucionado
} \index{Cardioide}


La cardioide escrita en forma polar como $r=1-\cos \theta,$ se parametriza como $x=\left( 1-\cos (t) \right) \cos (t) $
y $y=\left( 1-\cos (t) \right) \sin (t) ,$ para $t\in [0,2\pi].$
Sea $E$ el sólido obtenido cuando esta curva gira alrededor del eje polar.

\vskip0.2cm
\noindent
\begin{minipage}{\textwidth}
	\centering
	\captionof{figure}{Sólido $E$ generado por la rotación de $r = 1 - \cos \theta$.}
	\label{fig:toro_cardioidico_E}
	\begin{minipage}{0.25\linewidth}
		\centering
		\includegraphics[width=\linewidth]{Img/66b.pdf}
	\end{minipage}
	\hspace{1em}
	\begin{minipage}{0.25\linewidth}
		\centering
		\includegraphics[width=\linewidth]{Img/66a.pdf}
	\end{minipage}
	
	\vspace{0.3em}
	\fuentepropia
\end{minipage}

\vskip0.2cm
Esta figura se genera en \verb"Mathematica" usando los siguientes comandos:


\vspace{0.6em}
\noindent
\begin{tcolorbox}[breakable,enhanced,title=\bf Instrucción \verb"Mathematica", top=2pt,bottom=2pt,boxsep=0pt,before skip=2pt,after skip=0pt]
\begin{Verbatim}[formatcom=\letraverbatim]
{x,y}={(1-Cos[t])*Cos[t],(1-Cos[t])*Sin[t]};
r={x,y*Cos[v],y*Sin[v]};
ParametricPlot3D[r,{t,0,Pi},{v,0,2*Pi},MeshStyle->Gray,
PlotStyle->Orange,PlotPoints->60,Mesh->{20,15}]
\end{Verbatim}
\end{tcolorbox}
\vspace{0.7em}
\noindent


El volumen de $E$, usando el teorema de Pappus, está dado por
%\vskip0.2cm
{\small
\begin{eqnarray*}
V&=&\pi \int_{0}^{\pi }\left(- \left( 1-\cos (t) \right) \sin
(t) \right) ^{2}\frac{d}{dt}\left( \left( 1-\cos (t)
\right) \cos (t) \right) dt \\
&=& \pi \int_{0}^{\pi }\left( \sin (t) \left( \cos (t)
+1\right) \left( 1-2\cos (t) \right) \left( \cos (t)
1\right) ^{3} \right) dt = \frac{8}{3}\pi
\end{eqnarray*}
}
%\vskip0.2cm
Otra manera de obtener este valor es la siguiente.

%\begin{tcolorbox}[title=\bfseries Parametrizando el sólido]
%\vskip0.2cm
Una parametrización de la región interior a la cardioide es
\[x=u\left( 1-\cos (t) \right) \cos (t),\]
%\vskip0.2cm
\[y=u\left( 1-\cos (t) \right) \sin (t),\]
en donde $t\in [0,\pi]$ y $u\in [0,1].$
Si esta región gira sobre el eje polar, agregando $\cos(v)$ y $\sin(v)$ a $y$, con $v \in [0,2\pi]$, se obtiene una parametrización de $E$, esto es:
%\vskip0.2cm
\[{\bf r}(t,u,v)=u(1-\cos (t))\left\langle \cos (t) ,\sin (t) \cos (v) ,\sin(t) \sin (v) \right \rangle.\]
%\vskip0.2cm
El producto vectorial triple es
$\left| {\bf r}_{t}\cdot \left( {\bf r}_{u}\times {\bf r}_{v}\right) \right| = \left|
8 u^2 \sin ^6\left(\frac{t}{2}\right) \sin (t)\right| ,$
Entonces, el volumen de $E$ está dado por
%\vskip0.2cm
\[\int_{0}^{2\pi }\int_{0}^{1}\int_{0}^{\pi }\left| u^{2}\sin (t)
\left( \cos (t) -1\right) ^{3}\right| dtdudv= \frac{8}{3}\pi.\]
%\vskip0.2cm
%\end{tcolorbox}
	
	
Estos cálculos se pueden hacer en \verb"Mathematica" con ayuda de las siguientes instrucciones:


\vspace{0.6em}
\noindent
\begin{tcolorbox}[breakable,enhanced,title=\bf Instrucción \verb"Mathematica", top=2pt,bottom=2pt,boxsep=0pt,before skip=2pt,after skip=0pt]
\begin{Verbatim}[formatcom=\letraverbatim]
r=u*(1-Cos[t])*{Cos[t],Sin[t]*Cos[v],Sin[t]*Sin[v]};
A=Simplify[Dot[D[r,t],Cross[D[r,u],D[r,v]]]]
Integrate[Abs[A],{t,0,Pi},{v,0,2Pi},{u,0,1}]
\end{Verbatim}
\end{tcolorbox}
\vspace{0.7em}
\noindent
%\clearpage


\subsection*{Ejercicios propuestos para terminar}
En esta sección se muestran varios sólidos con las instrucciones requeridas para obtener su gráfico y
al menos una posibilidad de obtener su volumen mediante integrales múltiples.
El lector deberá, como ejercicio, obtener el mismo valor a través del planteamiento de integrales triples en otros 
ordenes de integración o planteando cambios de variables.
\vskip0.2cm
\begin{enumerate}[leftmargin=0cm,itemindent=0.5cm,labelwidth=\itemindent,labelsep=0cm,align=left,label={{\bf \arabic*)}}]
	
\item En \cite{stewart_calculo_2018}, sección 12.1, se plantea el siguiente ejercicio:
Describa y trace un sólido con las siguientes propiedades: cuando es iluminado por rayos paralelos al eje $z,$ su sombra es un disco circular. Si los rayos son paralelos al eje $y,$ su sombra es un cuadrado. Si los rayos son paralelos al eje $x,$
su sombra es un triángulo isósceles. Presentamos un sólido con estas características.

\vskip0.2cm

La proyección en el plano $xy$ corresponde a una círculo, cuya frontera se parametriza en la forma
$x= \cos(t)$, $y= \sin(t)$ para $t\in [0,2\pi].$ El diámetro en la base es $2,$ por tanto consideraremos un segmento paralelo al eje $y$ que sea perpendicular y simétrico al eje $z$. Un punto en la base tendrá la forma $P(\cos(t),\sin(t),0),$ mientras que un punto en la parte superior posee la forma $Q(0,\sin(t),2)$. Formando segmentos rectilíneos entre los puntos anteriores, es decir simplificando $(1-u)P+uQ,$
se obtiene una función que parametriza dicha superficie:
\vspace{-2mm}
\begin{eqnarray*}
{\bf r}(t,u) &=& \langle (1-u)\cos(t),\sin(t),2u \rangle, \,\, t\in [0,2\pi],\, u\in [0,1].
\end{eqnarray*}

\vskip0.2cm
\noindent
\begin{minipage}{\textwidth}
	\centering
	\captionof{figure}{\small Sólido cuyas proyecciones son un rectángulo, un círculo o un triángulo}
	\begin{minipage}{0.21\linewidth}
		\centering
		\includegraphics[width=\linewidth]{Imgn/T_20a.pdf}
	\end{minipage}
	\hspace{1em}
	\begin{minipage}{0.21\linewidth}
		\centering
		\includegraphics[width=\linewidth]{Imgn/T_20b.pdf}
	\end{minipage}
	\hspace{1em}
	\begin{minipage}{0.21\linewidth}
		\centering
		\includegraphics[width=\linewidth]{Imgn/T_20c.pdf}
	\end{minipage}
	\vspace{0.3em}
	\fuentepropia
\end{minipage}

\vskip0.2cm		
Estas figuras se consiguen con las siguientes instrucciones:

\vspace{0.6em}
\noindent
\begin{tcolorbox}[breakable,enhanced,title=\bf Instrucción \verb"Mathematica", top=2pt,bottom=2pt,boxsep=0pt,before skip=2pt,after skip=0pt]
\begin{Verbatim}[formatcom=\letraverbatim]
R=(1-u)*{0,Sin[t],2}+u*{Cos[t],Sin[t],0};
ParametricPlot3D[R,{t,0,2*Pi},{u,0,1},PlotPoints->80,
Ticks->{{-1,0,1},{-1,0,1},{1,2}},MeshStyle->Black]
\end{Verbatim}
\end{tcolorbox}
	
\vskip0.4cm	
El sólido se parametriza rellenando la base de sólido que corresponde a un círculo y formando segmentos rectilíneos hasta el segmento superior. Es decir,
%\vskip0.2cm
\vspace{-2mm}
\begin{eqnarray*}
{\bf r}(t,u,w) &=&(1-w)(0,u\sin(t),2)+w(u\cos(t),u\sin(t),0) \\
&=& \langle uw\cos(t),u\sin(t),2(1-w) \rangle, 
\end{eqnarray*}
%\vskip0.2cm
con $t\in [0,2\pi],\, u\in [0,1],\, w\in [0,1].$
El elemento diferencial de volumen está dado por $|{\bf r}_w \cdot ({\bf r}_t \times {\bf r}_u) |=2uw,$ por lo tanto el volumen está dado por
%\vskip0.2cm

$$\int_0^1 \int_0^1 \int_0^{2\pi} 2 uw dt du dw = \pi.$$
%\vskip0.2cm

Otro sólido con los requerimientos hechos es el que satisface las expresiones
$0\leq z \leq 1-|x|$, $x^2+y^2\leq 1.$ Obviamente su volumen no es el mismo al obtenido anteriormente.
\vskip0.3cm
	
\item El sólido acotado por las superficies de los los paraboloides con ecuaciones $2z=6-2x^2-y^2$ y $z+3=x^2+y^2,$ se muestra a continuación.

\vskip0.3cm
\noindent
\begin{minipage}{\textwidth}
	\centering
	\captionof{figure}{\small $2x^2+2y^2-6\leq 2z\leq 6-2x^2-y^2$ y $,$}
	\begin{minipage}{0.23\linewidth}
		\centering
		\includegraphics[width=\linewidth]{Imgn/T_8a.pdf}
	\end{minipage}
	\hspace{1em}
	\begin{minipage}{0.25\linewidth}
		\centering
		\includegraphics[width=\linewidth]{Imgn/T_8e.pdf}
	\end{minipage}
	
	\vspace{0.3em}
	\fuentepropia
\end{minipage}

\vskip0.2cm
Ahora presentamos las instrucciones necesarias para conseguir esta figura.


\vspace{0.6em}
\noindent
\begin{tcolorbox}[breakable,enhanced,title=\bf Instrucción \verb"Mathematica", top=2pt,bottom=2pt,boxsep=0pt,before skip=2pt,after skip=0pt]
\begin{Verbatim}[formatcom=\letraverbatim]
RegionPlot3D[z<(6-2*x^2-y^2)/2&&z>x^2+y^2-3,{x,-2,2},
{y,-2,2},{z,-3,3},PlotPoints->90,Mesh->False];
\end{Verbatim}
\end{tcolorbox}

\vskip0.4cm	
Al despejar $z$ de las dos ecuaciones e igualar estas expresiones se obtiene $4x^{2}+3y^{2}=12$. Esta expresión corresponde a la intersección de las superficies consideradas y representa, en el plano
$xy$, a una elipse alargada en sentido de las $y$ centrada en el origen. Por tanto, el volumen del sólido está dado por
{\small
\begin{equation} \label{ejercicio1} 
\bigintss_{-\sqrt{3}}^{\sqrt{3}}\bigintss_{-\frac{2}{\sqrt{3}}\sqrt{3-x^{2}}}^{\frac{2}{\sqrt{3}}\sqrt{3-x^{2}}}\bigintss_{x^{2}+y^{2}-3}^{3-x^{2}-\frac{1}{2}y^{2}}dzdydx= 6\sqrt{3}\pi .
\end{equation}
}
\vskip0.2cm
Otra manera de hallar este valor es la siguiente. La elipse y su interior se parametriza como $x=\sqrt{3}u\cos \theta$, $y=2u\sin \theta ,$ con $u \in [0,1]$, $\theta \in [0,2\pi].$ Reemplazando estas expresiones en las ecuaciones de los paraboloides se sigue que

\vspace{-12mm}
\begin{eqnarray*}
z &=& 3-x^{2}-\frac{1}{2}y^{2}= 3-\frac{5}{2}u^{2}-\frac{1}{2}u^{2}\cos 2\theta \\
z&=&x^2+y^2-3= \frac{7}{2}u^{2}-\frac{1}{2}u^{2}\cos 2\theta -3.
\end{eqnarray*}

Sobre cada uno de los paraboloides consideremos dos puntos con coordenadas
$Q(\sqrt{3}u\cos \theta ,2u\sin \theta ,\frac{7}{2}u^{2}-\frac{1}{2}u^{2}\cos 2\theta-3)$
y $P(\sqrt{3}u\cos \theta,2u\sin \theta,\\ 3-\frac{5}{2}u^{2}-\frac{1}{2}u^{2}\cos 2\theta).$  
Una combinación convexa de estos puntos permite parametrizar el sólido, es decir, simplificando la expresión $(1-w)P+wQ$ se define la función
%\vskip0.2cm
{\small
\[{\bf r}(u,\theta ,w)=\left\langle\sqrt{3}u\cos\theta,2u\sin\theta,6u^{2}w-\frac{1}{2}u^{2}\cos2\theta -6w-\frac{5}{2}u^{2}+3 \right\rangle,\]}
%\vskip0.2cm

en donde $u \in [0,1]$, $\theta \in [0,2\pi] ,$ $w\in [0,1].$ 
Con esta función hallamos los siguientes vectores
\vspace{-4mm}
\begin{eqnarray*}
{\bf r}_{\theta}&=& \left\langle -\sqrt{3}u\sin \theta ,2u\cos \theta ,u^{2}\sin 2\theta \right\rangle\\
{\bf r}_{u} &=& \left\langle \sqrt{3}\cos \theta ,2\sin \theta ,12uw-u\cos 2\theta -5u\right\rangle \\
{\bf r}_{w} &=& \left\langle 0,0,6u^{2}-6\right\rangle \\
{\bf r}_{\theta }\times {\bf r}_{u}&=& \left\langle 12u^{2}\left( 2w-1\right)
\cos \theta ,4\sqrt{3}u^{2}\left( 3w-1\right) \sin \theta ,-2\sqrt{3}u \right\rangle.
\end{eqnarray*}

El elemento diferencial de volumen es $\left\vert {\bf r}_{w}\cdot \left({\bf r}_{\theta}\times {\bf r}_{u}\right) \right\vert= 12\sqrt{3}u\left( 1-u^{2}\right) ,$ cuya integración con los límites establecidos nos proporciona el volumen del sólido $$\int_{0}^{1}\int_{0}^{1}\int_{0}^{2\pi }12\sqrt{3}u\left( 1-u^{2}\right)
d\theta dudw= 6\sqrt{3}\pi.$$
\vskip0.2cm
	
\item Consideremos el sólido determinado por las expresiones $x^2+y^2+z^2\leq 9,$ $1+\cos\theta \leq r\leq 2+\cos\theta$. Esto es, la región del espacio acotada por dos cilindros cardióidicos e interior la esfera de radio $3$ centrada en el origen.
\vskip0.2cm
\noindent
\begin{minipage}{\textwidth}
	\centering
	\captionof{figure}{\small $x^2+y^2+z^2\leq 9$ \& $1+\cos\theta\leq r\leq 2+\cos\theta$}
	\begin{minipage}{0.22\linewidth}
		\centering
		\includegraphics[width=\linewidth]{Imgn/T_7a}
	\end{minipage}
	\hspace{1em}
	\begin{minipage}{0.22\linewidth}
		\centering
		\includegraphics[width=\linewidth]{Imgn/T_7b.pdf}
	\end{minipage}
	\hspace{1em}
	\begin{minipage}{0.22\linewidth}
		\centering
		\includegraphics[width=\linewidth]{Imgn/T_7c.pdf}
	\end{minipage}
	
	\vspace{0.3em}
	\fuentepropia
\end{minipage}

\vskip0.2cm
Esta figura y su proyección en el plano $xy$ se obtiene con estas instrucciones:
\vskip0.2cm	


\vspace{0.6em}
\noindent
\begin{tcolorbox}[breakable,enhanced,title=\bf Instrucción \verb"Mathematica", top=2pt,bottom=2pt,boxsep=0pt,before skip=2pt,after skip=0pt]
\begin{Verbatim}[formatcom=\letraverbatim]
PolarPlot[{1+Cos[t],2+Cos[t]},{t,0,2Pi}]
RegionPlot3D[x^2+y^2+z^2<9&&(x^2+y^2-x)^2>x^2+y^2&&
(x^2+y^2-x)^2<3*(x^2+y^2),{x,-1.2,3.2},{y,-2.5,2.5},
{z,-3,3},PlotPoints->90,Mesh->False]
\end{Verbatim}
\end{tcolorbox}

\vskip0.2cm
%\clearpage
La base del sólido se muestra a continuación:

\vskip0.2cm	
\noindent
\begin{minipage}{\textwidth}
	\centering
	\captionof{figure}{\small $1+\cos\theta\leq r\leq 2+\cos\theta$}
	\begin{minipage}{0.30\linewidth}
		\centering
		\vspace{-12mm}
		\includegraphics[width=\linewidth]{Imgn/T_600.pdf}
	\end{minipage}
	
	\vspace{-8mm}
	\fuentepropia
\end{minipage}
\vskip0.3cm

Sobre la proyección del sólido en el plano $xy$ consideremos dos puntos,
$P=((1+\cos(t))\cos(t),(1+\cos(t))\sin(t))$ y $Q=((2+\cos(t))\cos(t),(2+\cos(t))\sin(t)),$ 
uno sobre cada curva. La región se parametriza al formar una combinación convexa de la forma $(1-u)P+uQ$ entre estos puntos, es decir $(\cos(t)(1+u+\cos(t)),(1+u+\cos(t))\sin(t)).$

\vskip0.2cm 

Las siguientes instrucciones muestran en \verb"Mathematica" cómo hacer el procedimiento completo y además permiten hallar su área.
\vskip0.2cm	


\vspace{0.6em}
\noindent
\begin{tcolorbox}[breakable,enhanced,title=\bf Instrucción \verb"Mathematica", top=2pt,bottom=2pt,boxsep=0pt,before skip=2pt,after skip=0pt]	
\begin{Verbatim}[formatcom=\letraverbatim]
r={x,y}=Simplify[(1-u)*{(1+Cos[t])*Cos[t],
(1+Cos[t])*Sin[t]}+u*{(2+Cos[t])*Cos[t],
(2+Cos[t])*Sin[t]}];
w=(1-v)*{x,y,-Sqrt[9-x^2-y^2]}+v*{x,y,
Sqrt[9-x^2-y^2]};
A=Abs[Simplify[Dot[D[w,t],Cross[D[w,u],D[w,v]]]]];
NIntegrate[A,{t,0,2Pi},{u,0,1},{v,0,1}]
\end{Verbatim}
\end{tcolorbox}

\vskip0.4cm
Utilizando coordenadas cilíndricas, el volumen del mencionado sólido está dado por la siguiente integral:
%\vskip0.2cm
$$\int_0^{2\pi} \int_{1+\cos \theta}^{2+\cos \theta}\int_{-\sqrt{9-r^2}}^{\sqrt{9-r^2}}r dzdrd\theta \approx 40.7091.$$
\vskip0.2cm
La evaluación numérica de esta integral se consigue ejecutando la siguiente instrucción:


\vskip0.2cm
\vspace{0.6em}
\noindent
\begin{tcolorbox}[breakable,enhanced,title=\bf Instrucción \verb"Mathematica", top=2pt,bottom=2pt,boxsep=0pt,before skip=2pt,after skip=0pt]
\begin{Verbatim}[formatcom=\letraverbatim]
NIntegrate[r,{t,0,2Pi},{r,1+Cos[t],2+Cos[t]},
{z,-Sqrt[9-r^2],Sqrt[9-r^2]}]
\end{Verbatim}
\end{tcolorbox}
\vskip0.4cm	
	
\item El sólido acotado por las superficies del cilindro $x^2+y^2=4$ y la esfera $x^2+y^2+z^2=9$ se muestran a continuación.
\vskip0.2cm	

\noindent
\begin{minipage}{\textwidth}
	\centering
	\captionof{figure}{\small $x^2+y^2 \leq 4$ \& $6x^2+y^2+z^2\leq 9.$}
	\begin{minipage}{0.20\linewidth}
		\centering
		\includegraphics[width=\linewidth]{Imgn/T_9a.pdf}
	\end{minipage}
	\hspace{1em} 
	\begin{minipage}{0.18\linewidth}
		\centering
		\includegraphics[width=\linewidth]{Imgn/T_9d.pdf}
	\end{minipage}
	
	\vspace{0.3em}
	\fuentepropia
\end{minipage}

\vskip0.2cm
Las siguientes instrucciones permiten obtenerlo:


\vspace{0.6em}
\noindent
\begin{tcolorbox}[breakable,enhanced,title=\bf Instrucción \verb"Mathematica", top=2pt,bottom=2pt,boxsep=0pt,before skip=2pt,after skip=0pt]	
\begin{Verbatim}[formatcom=\letraverbatim]
RegionPlot3D[x^2+y^2<4&&x^2+y^2+z^2<=9,{x,-2,2},
{y,-2,2},{z,-3.5,3.5},PlotPoints->90,Mesh->False]
\end{Verbatim}
\end{tcolorbox}

\vskip0.4cm
Usando coordenadas rectangulares, se plantea la integral que representa el volumen de dicha región.
%\vskip0.2cm
$$\bigintsss_{-2}^2 \bigintsss_{-\sqrt{4-x^{2}}}^{\sqrt{4-x^{2}}}\bigintsss_{-\sqrt{9-x^{2}-y^{2}}}^{\sqrt{9-x^{2}-y^{2}}}dzdydx \approx 66.265. $$
%\vskip0.2cm
En el plano $xy$ la proyección del cilindro y su interior se parametriza como $x=2u\cos \theta, $
$y=2u\sin \theta $ con $u \in [0,1]$ y $\theta \in [0,2\pi];$ sobre la esfera la expresión correspondiente es
%\vskip0.2cm
\[z=\pm \sqrt{9-\left( 2u\cos \theta \right)^2 -\left( 2u\sin \theta \right) ^2 }= \pm \sqrt{9-4u^{2}}.\]
%\vskip0.2cm
De modo que el sólido se parametriza mediante una combinación convexa de dos puntos sobre la esfera dentro del cilindro, esto es:
%\vskip0.2cm
{\small
\[\left( 1-w\right) \left( 2u\cos \theta ,2u\sin
\theta,\sqrt{9-4u^{2}}\right) +w\left( 2u\cos \theta ,2u\sin \theta,-\sqrt{9-4u^{2}}\right).\]
}
%\vskip0.2cm

Lo que nos lleva a presentar la función
%\vskip0.2cm
\begin{eqnarray*}
{\bf r}\left( u,\theta ,w\right) &=& \left\langle 2u\cos \theta ,2u\sin \theta ,\left( 1-2w\right) \sqrt{9-4u^{2}}\right\rangle 
\end{eqnarray*}
%\vskip0.2cm

con $u \in [0,1]$, $\theta \in [0,2\pi],$ $w \in [0,2\pi];$ 
con esta se hallan los siguientes vectores

\vspace{-10mm}
\begin{eqnarray*}
{\bf r}_{\theta } &=& \left\langle -2u\sin \theta ,2u\cos \theta ,0\right\rangle \\
{\bf r}_{u} &=& \left\langle 2\cos \theta ,2\sin \theta ,\frac{4u\left( 2w-1\right) }{\sqrt{9-4u^{2}}}%
\right\rangle \\
{\bf r}_{w}&=& \left\langle 0,0,-2\sqrt{9-4u^{2}}\right\rangle. 
\end{eqnarray*}
El elemento diferencial de volumen es
%\vskip0.2cm
\begin{eqnarray*}
\left\vert {\bf r}_{w}\cdot \left({\bf r}_{u}\times {\bf r}_{\theta }\right) \right\vert 
=8u\sqrt{9-4u^{2}},
\end{eqnarray*}
%\vskip0.2cm
por lo tanto el volumen es
\vskip0.2cm
\begin{equation} \label{ejercicio2}
\int_{0}^{1}\int_{0}^{1}\int_{0}^{2\pi }8u\sqrt{9-4u^{2}}d\theta dudw= \frac{4}{3}\pi \left(27-5\sqrt{5}\right) 
\end{equation} 
\vskip0.2cm

\item El sólido que contiene a los puntos interiores, tanto a la esfera con ecuación $x^2+y^2+z^2=9$ como al cilindro cardióidico
$r=1+\cos \theta,$ se muestra a continuación.

\vskip0.2cm	
\noindent
\begin{minipage}{\textwidth}
	\centering
	\captionof{figure}{\small $x^2+y^2+z^2\leq 9$ \& $r\leq 1+\cos \theta.$}
	\begin{minipage}{0.20\linewidth}
		\centering
		\includegraphics[width=\linewidth]{Imgn/T_12a}
	\end{minipage}
	\hspace{1em}
	\begin{minipage}{0.20\linewidth}
		\centering
		\includegraphics[width=\linewidth]{Imgn/T_12b}
	\end{minipage}
	
	\vspace{0.3em}
	\fuentepropia
\end{minipage}

\vskip0.2cm	
Las instrucciones necesarias para conseguirlo son las siguientes:
\vskip0.2cm


\vspace{0.6em}
\noindent
\begin{tcolorbox}[breakable,enhanced,title=\bf Instrucción en \verb"Mathematica", top=2pt,bottom=2pt,boxsep=0pt,before skip=2pt,after skip=0pt]
\begin{Verbatim}[formatcom=\letraverbatim]
RegionPlot3D[(x^2+y^2-x)^2<x^2+y^2&&x^2+y^2+z^2<9,
{x,-1/2,2},{y,-3/2,3/2},{z,-3,3},Mesh->False,
PlotPoints->80]
\end{Verbatim}
\end{tcolorbox}

\vskip0.4cm
Usando coordenadas cilíndricas se puede plantear la integral que determina el volumen de esta región.
%\vskip0.2cm
\begin{multline} \label{ejercicio3}
\int_{0}^{2\pi }\int_{0}^{1+\cos \theta}\int_{-\sqrt{9-r^{2}}}^{\sqrt{9-r^{2}}}rdzdrd\theta \\ 
= 2\int_{0}^{2\pi }\int_{0}^{1+\cos \theta }r\sqrt{9-r^{2}}\,drd\theta \approx 25.8179 
\end{multline}
\vskip0.2cm
	
\item El siguiente es la figura del sólido acotado por el paraboloide $x^2+y^2=3-z$ y el paraboloide hiperbólico $x^2-y^2=3z+6.$

\vskip0.2cm
\noindent
\begin{minipage}{\textwidth}
	\centering
	\captionof{figure}{\small $x^2+y^2\geq 3-z$ \& $x^2-y^2\leq 3z+6$}
	%{\label{ejercicio4a}
	\begin{minipage}{0.20\linewidth}
		\centering
		\includegraphics[width=\linewidth]{Imgn/T_10a.pdf}
	\end{minipage}
	\hspace{1em}
	\begin{minipage}{0.20\linewidth}
		\centering
		\includegraphics[width=\linewidth]{Imgn/T_10b.pdf}
	\end{minipage}
	
	\vspace{0.3em}
	\fuentepropia
\end{minipage}

\vskip0.2cm
Con las siguientes instrucciones se consigue esta figura:


\vspace{0.6em}
\noindent
\begin{tcolorbox}[breakable,enhanced,title=\bf Instrucción \verb"Mathematica", top=2pt,bottom=2pt,boxsep=0pt,before skip=2pt,after skip=0pt]	
\begin{Verbatim}[formatcom=\letraverbatim]
RegionPlot3D[x^2+y^2<3-z&&x^2-y^2<3*z+6,{x,-2,2},
{y,-3,3},{z,-5,3},PlotPoints->90,Mesh->False]
\end{Verbatim}
\end{tcolorbox}

\vskip0.4cm	
	
La intersección de estas dos superficies se obtiene al igualar las expresiones para $z$, esto es, $$3-x^{2}-y^{2}=\frac{x^{2}-y^{2}-6}{3},$$
la cual se simplifica como $4x^2+2y^2=15$. Esta ecuación corresponde a una elipse en el plano $xy$ centrada en el origen y alargada en sentido de las $y.$ Por tanto, el volumen del sólido está dado por
\vskip0.2cm
{\small
\begin{equation} \label{ejercicio4}
\bigintss_{-\frac{\sqrt{15}}{2}}^{\frac{\sqrt{15}}{2}}\bigintss_{-\frac{1}{2}\sqrt{2} \sqrt{15-4x^{2}}}^{\frac{1}{2}\sqrt{2}\sqrt{15-4x^{2}}}\bigintss_{\frac{x^{2}-y^{2}-6}{3}}^{3-x^{2}-y^{2}}dzdydx= \frac{75}{8}\sqrt{2}\pi
\end{equation}
}
\vskip0.2cm

Otra posibilidad para hallar este valor es parametrizar la elipse mediante las expresiones
{\small $x=\frac{\sqrt{15}}{2}u\cos \theta $}, {\small $y=\sqrt{\frac{15}{2}}u\sin \theta $} con {\small $u \in [0,1]$} y {\small $\theta \in [0,2\pi].$}

Luego, reemplazar en las ecuaciones del paraboloide y el paraboloide hiperbólico, lo cual produce, respectivamente:	
%\vskip0.2cm
{\small
\begin{multline*}
z = 3-\left( \frac{\sqrt{15}}{2}u\cos \theta \right) ^2 -\left( \sqrt{\frac{15}{2}}u\sin \theta \right) ^2  = \frac{15}{8}u^{2}(\cos 2\theta -3)+3 \\
z =  \frac{1}{3}\left( \left( \frac{\sqrt{15}}{2}u\cos \theta \right)^2-\left( \sqrt{\frac{15}{2}}u\sin \theta \right) ^2 -6\right) =  \frac{15}{8}u^{2}(\cos 2\theta -1)-2.
\end{multline*}
}
%\vskip0.2cm
Con estas expresiones tendremos dos puntos de la forma 

{\small
\begin{eqnarray*}
P& =&\left( \frac{\sqrt{15}}{2}u\cos \theta,\sqrt{\frac{15}{2}}u\sin \theta , \frac{15}{8}u^{2}(\cos 2\theta -3)+3\right) \\
Q&=&\left( \frac{\sqrt{15}}{2}u\cos \theta,\sqrt{\frac{15}{2}}u\sin \theta,\frac{15}{8}u^{2}\cos 2\theta-\frac{5}{8}u^{2}-2\right)\!.
\end{eqnarray*}
}

\vskip0.2cm
Una combinación convexa de ellos, $(1-w)P+wQ,$ determina una parametrización del sólido
%\vskip0.2cm
{\small
\[{\bf r}(u,\theta,w) = \left\langle \frac{\sqrt{15}}{2}u\cos v,\sqrt{\frac{15}{2}}u\sin v,5w\left( u^{2}-1\right) +\frac{15}{8}u^{2}\left(\cos2v-3\right) +3 \right\rangle.\]
}

De esta función se obtienen los siguientes vectores:
{\small
\begin{align*}
{\bf r}_{\theta } &= \left\langle -\frac{1}{2}\sqrt{15}u\sin \theta ,\frac{1}{2}%
\sqrt{2}\sqrt{15}u\cos \theta ,-\frac{15}{4}u^{2}\sin 2\theta \right\rangle \\ 
{\bf r}_{u} &= \left\langle \frac{1}{2}\sqrt{15}\cos \theta ,\frac{1}{2}\sqrt{2} \sqrt{15}\sin \theta ,\frac{15}{4}u\left( \cos 2\theta -3\right) +10uw \right\rangle \\
{\bf r}_{w} &= \left\langle 0,0,5u^{2}-5\right\rangle \\
{\bf r}_{\theta }\times {\bf r}_{u} &= \left\langle \frac{5\sqrt{30}}{4} u^{2} (4w-3) \cos \theta ,\frac{5\sqrt{15}}{2}u^{2}(2w-3) \sin \theta ,-\frac{15\sqrt{2}}{4}u \right\rangle.
\end{align*}
}
\vskip0.2cm

El elemento diferencial de volumen es
$\left\vert {\bf r}_{w}\cdot \left({\bf r}_{\theta }\times {\bf r}_{u}\right) \right\vert
=  \left\vert \frac{75}{4}\sqrt{2}\left( u-u^{3}\right) \right\vert ,$
cuya integración nos proporciona el volumen deseado
\vskip0.2cm
$$\int_{0}^{1}\int_{0}^{1}\int_{0}^{2\pi }\left\vert \frac{75}{4}\sqrt{2}\left( u-u^{3}\right) \right\vert d\theta dudw
= \frac{75}{8}\sqrt{2}\pi .$$
\vskip0.2cm
	
\item El sólido cuya región, descrita en coordenadas esféricas, comprende los puntos que satisfacen la expresión $\rho \leq 1+\cos \phi,$ se presenta a continua\-ción:
\vskip0.2cm
\noindent
\begin{minipage}{\textwidth}
	\centering
	\captionof{figure}{\small Cardioide de revolución}
	\begin{minipage}{0.20\linewidth}
		\centering
		\includegraphics[width=\linewidth]{Imgn/T_13a.pdf}
	\end{minipage}
	\hspace{1em}
	\begin{minipage}{0.20\linewidth}
		\centering
		\includegraphics[width=\linewidth]{Imgn/T_13d.pdf}
	\end{minipage}
	
	\vspace{0.3em}
	\fuentepropia
\end{minipage}
\vskip0.2cm

Las siguientes instrucciones permiten obtener la figura:


\vspace{0.6em}
\noindent
\begin{tcolorbox}[breakable,enhanced,title=\bf Instrucción en \verb"Mathematica", top=2pt,bottom=2pt,boxsep=0pt,before skip=2pt,after skip=0pt]	
\begin{Verbatim}[formatcom=\letraverbatim]
RegionPlot3D[(x^2+y^2+z^2-x)^2<x^2+y^2+z^2,{x,-1,2},
{y,-2,2},{z,-3/2,3/2},PlotPoints->30,Mesh->False,
BoxRatios->{1,1,1}]
\end{Verbatim}
\end{tcolorbox}
	
\vskip0.4cm	
En coordenadas esféricas el volumen del sólido está dado por

\begin{equation}
\label{ejercicio5}
\int_{0}^{\pi }\int_{0}^{2\pi }\int_{0}^{1+\cos \phi }\rho ^2 \sin \phi
d\rho d\theta d\phi =\frac{8}{3}\pi.
\end{equation}
	
Sustituyendo $\rho =1+\cos \phi,$ en las fórmulas que definen las coordenadas esféricas, se obtienen las siguientes expresiones: 
$x=u\left( 1+\cos \phi \right) \sin \phi \sin \theta ,$ $y=u\left( 1+\cos \phi \right) \sin \phi \cos \theta ,$  $z=u\left( 1+\cos \phi \right) \cos \phi .$
Por tanto, se define la función 
%\vskip0.2cm
$${\bf r}(u,\theta,\phi)=(1+\cos \phi) \left\langle u \sin
\phi \cos \theta ,u\sin \phi \sin \theta ,u \cos \phi \right\rangle $$
\vskip0.2cm
que parametriza el mencionado sólido con $u\in[0,1]$, $\phi \in[0,\pi]$, $\theta\in[0,2\pi]$. 
Con la función ${\bf r}$ se hallan los siguientes vectores
\noindent
{\small
	\begin{gather*}
		{\bf r}_{u} = (1+\cos \phi) \langle \cos \theta \sin \phi, \sin \theta \sin \phi, \cos \phi \rangle \\[6pt]
		{\bf r}_{\phi} = u \langle (\cos \phi + \cos (2\phi)) \cos \theta, (\cos \phi + \cos (2\phi)) \sin \theta, -(2\cos \phi + 1) \sin \phi \rangle \\[6pt]
		{\bf r}_{\theta} = \langle -u(\cos \phi + 1) \sin \theta \sin \phi, u(\cos \phi + 1) \cos \theta \sin \phi, 0 \rangle 
	\end{gather*}
}
\vskip0.2cm

Además {\footnotesize ${\bf r}_{u}\times {\bf r}_{\phi }= \left\langle  -u\left( \cos \phi +1\right) ^2 \sin \theta ,u\left( \cos \phi +1\right) ^2 \cos \theta ,0 \right\rangle.$}
%\linebreak
El elemento diferencial de volumen es {\footnotesize $\left\vert {\bf r}_{\theta }\cdot \left({\bf r}_{u}\times {\bf r}_{\phi }\right) \right\vert =u^{2}\left( \cos \phi +1\right) ^{3}\sin \phi ;$} con lo anterior, se plantea la integral que determina el volumen del sólido
%\vskip0.2cm 
{\small
$$\int_{0}^{1}\int_{0}^{2\pi }\int_{0}^{\pi }u^{2}\left( \cos \phi +1\right) ^{3}\sin \phi d\phi d\theta du= \frac{8}{3}\pi.$$
}
%\vskip0.2cm

\item El sólido acotado por la superficie del semicono $y=\sqrt{x^2+z^2}$ y la esfera $x^2+y^2+z^2=9,$
se muestran enseguida. \label{ejercicio7a}
\vskip0.2cm

\noindent
\begin{minipage}{\textwidth}
	\centering
	\captionof{figure}{\small $\sqrt{x^2+z^2}\leq y$ \& $x^2+y^2+z^2\leq 9$}
	\begin{minipage}{0.21\linewidth}
		\centering
		\includegraphics[width=\linewidth]{Imgn/T_14a.pdf}
	\end{minipage}
	\hspace{1em}
	\begin{minipage}{0.21\linewidth}
		\centering
		\includegraphics[width=\linewidth]{Imgn/T_14b.pdf}
	\end{minipage}
	
	\vspace{0.3em}
	\fuentepropia
\end{minipage}

\vskip0.2cm	
Estas figuras se obtienen con ayuda de las siguientes instrucciones:


\vspace{0.6em}
\noindent
\begin{tcolorbox}[breakable,enhanced,title=\bf Instrucción en \verb"Mathematica", top=2pt,bottom=2pt,boxsep=0pt,before skip=2pt,after skip=0pt]	
\begin{Verbatim}[formatcom=\letraverbatim]
RegionPlot3D[y>Sqrt[x^2+z^2]&&x^2+y^2+z^2<9,
{x,-2.5,2.5},{y,0,3.2},{z,-2.5,2.5},PlotPoints->90,
Mesh->False,BoxRatios->{1,1,1}]
\end{Verbatim}
\end{tcolorbox}

\vskip0.4cm
	
La intersección de las dos superficies está dada por la expresión $x^{2}+z^{2}=\frac{9}{2}.$ Además, sobre el sólido $y$ recorre los puntos que van desde el semicono hasta la esfera. Por tanto, el volumen se determina evaluando la siguiente integral:
$$\bigintsss_{-\frac{3}{\sqrt{2}}}^{\frac{3}{\sqrt{2}}}\bigintsss_{-\sqrt{\frac{9}{2}-x^{2}}}^{\sqrt{\frac{9}{2}-x^{2}}}\bigintsss_{\sqrt{x^{2}+z^{2}}}^{\sqrt{9-x^{2}-z^{2}}}dydzdx = 9 (2-\sqrt{2}) \pi.$$
%\vskip0.2cm

Utilizando coordenadas esféricas la integral correspondiente es
%\vskip0.2cm
\begin{equation} \label{ejercicio7}
\int_{0}^{2\pi }\int_{0}^{\frac{3}{\sqrt{2}}}\int_{r}^{\sqrt{9-r^{2}}}rdydrd\theta = 9\pi \left( 2-\sqrt{2}\right).
\end{equation}
%\vskip0.2cm

También se puede calcular el volumen considerando un punto sobre el cono {\small $P\left(u\cos v,u,u\sin v\right)$} y otro punto {\small $Q\left( u\cos v,\sqrt{9-u^{2}},u\sin v\right)$} sobre la esfera. De esta manera de forma una combinación convexa entre estos puntos y se parametriza el sólido. Esto es
% ${\bf r}(u,v,w) =(1-w) P +wQ  = \left\langle u\cos v,u+w\sqrt{9-u^{2}}-uw,u\sin v\right\rangle $
\vspace{-2mm}
{\small
\begin{eqnarray*} 
{\bf r}(u,v,w) &=&(1-w) P +wQ  = \left\langle u\cos v,u+w\sqrt{9-u^{2}}-uw,u\sin v\right\rangle,
\end{eqnarray*}
}
%\vskip0.2cm
con {\small $u\in \lbrack 0,\frac{3}{\sqrt{2}}],$} {\small $v\in \lbrack 0,2\pi ]$} y {\small $w\in \lbrack 0,1].$} De esta función se obtienen los siguien\-tes vectores: {\footnotesize ${\bf r}_{u} =\left\langle \cos v,1-w-\frac{uw}{\sqrt{9-u^{2}}},\sin v\right\rangle ,$} {\footnotesize ${\bf r}_{v} = \left\langle -u\sin v,0,u\cos v\right\rangle$} y {\footnotesize ${\bf r}_{w} = \left\langle 0,\sqrt{9-u^{2}}-u,0\right\rangle.$} El elemento diferencial de volumen es
{\footnotesize $\left\vert {\bf r}_{u}\cdot \left({\bf r}_{v}\times {\bf r}_{w}\right)\right\vert 
\\
=\left\vert u^{2}-u\sqrt{9-u^{2}}\right\vert.$} Con lo anterior, el volumen está dado por
\vskip0.2cm

{\small
$$\int_{0}^{1}\int_{0}^{2\pi }\int_{0}^{\frac{3}{\sqrt{2}}}\left\vert u^{2}-u \sqrt{9-u^{2}}\right\vert 
dudvdw= 9\pi \left( 2-\sqrt{2}\right).$$
}

\vskip0.2cm
	
\item El hiperboloide de un manto con ecuación $2x^2+3y^2=z^2+7$ y el elipsoide $2x^2+y^2+z^2=9$ acotan el sólido que se muestra a continuación. Esta región se presenta en detalle en la página \pageref{elipsoidehiperboloide}.

\vskip0.3cm
\noindent
\begin{minipage}{\textwidth}
	\centering
	\captionof{figure}{\small $2x^2+3y^2\leq z^2+7$ \& $2x^2+y^2+z^2\leq 9$}
	\begin{minipage}{0.23\linewidth}
		\centering
		\includegraphics[width=\linewidth]{Imgn/T_11a.pdf}
	\end{minipage}
	\hspace{1em}
	\begin{minipage}{0.23\linewidth}
		\centering
		\includegraphics[width=\linewidth]{Imgn/T_11b.pdf}
	\end{minipage}
	
	\vspace{0.3em}
	\fuentepropia
\end{minipage}

\vskip0.3cm
Esta figura se consigue ejecutando las siguientes instrucciones:
%\vskip0.2cm


\vspace{0.6em}
\noindent
\begin{tcolorbox}[breakable,enhanced,title=\bf Instrucción en \verb"Mathematica", top=2pt,bottom=2pt,boxsep=0pt,before skip=2pt,after skip=0pt]
\begin{Verbatim}[formatcom=\letraverbatim]
RegionPlot3D[2*x^2+3*y^2<z^2+7&&2*x^2+y^2+z^2<9,
{x,-2,2},{y,-2,2},{z,-3,3.2},PlotPoints->90,
Mesh->False];
\end{Verbatim}
\end{tcolorbox}

\vskip0.4cm	
En el plano $xy$, la proyección de la parte más delgada del hiperboloide posee ecuación $2x^{2}+3y^{2}=7.$ Eliminando $z$ de las ecuaciones iniciales se obtiene $x^{2}+y^{2}=4,$ que corresponde a la proyección en el plano $xy$ de la intersección de las superficies. 
\vskip0.4cm
Hay una región en la cual los l\'{\i}mites de integración para $z$ van desde el elipsoide hasta el mismo elipsoide, mientras que en la región restante, los límites para $z$ van desde el hiperboloide hasta el elipsoide. Por la simetr\'{\i}a de la figura se puede establecer que el volumen está dado por

\vspace{-4mm}
{\small
\begin{multline}
\bigintsss_{0}^{\sqrt{\frac{7}{2}}}\bigintsss_{0}^{\sqrt{\frac{7-2x^{2}}{3}}}\bigintsss_{0}^{\sqrt{9-2x^{2}-y^{2}}}8dzdydx \\ 
+ \bigintsss_{0}^{\sqrt{\frac{7}{2}}}\bigintsss_{\sqrt{\frac{7-2x^{2}}{3}}}^{\sqrt{4-x^{2}}}\bigintsss_{\sqrt{2x^{2}+3y^{2}-7}}^{\sqrt{9-2x^{2}-y^{2}}}8dzdydx\\
+\bigintsss_{\sqrt{\frac{7}{2}}}^2 \bigintsss_{0}^{\sqrt{4-x^{2}}}\bigintsss_{\sqrt{2x^{2}+3y^{2}-7}}^{\sqrt{9-2x^{2}-y^{2}}}8dzdydx \approx 51.8616 \label{ejercicio6}
\end{multline}
}
%\vskip0.2cm	
En el plano $xz$ la proyección del sólido genera una elipse y una hipérbola; mientras que la curva de intersección proyectada en el plano $xz$ se obtiene eliminando $y$ de las ecuaciones iniciales, esto se reduce a la fórmula $x^2+z^2=5$. El volumen del sólido se determina integrando primero con respecto a $y$ sobre dos regiones,
en una de ellas $y$ recorre puntos que van desde el elipsoide hasta la misma superficie y en la otra región $y$ se recorre los puntos desde el hiperboloide hasta el mismo. Con lo anterior, el planteamiento es el siguiente:
{\small
\begin{multline} 
\hskip-0.4cm \bigintsss_{0}^2 \bigintsss_{\sqrt{5-x^{2}}}^{\sqrt{9-2x^{2}}}\bigintsss_{0}^{\sqrt{9-2x^{2}-z^{2}}}8dydzdx + 
\bigintsss_{\sqrt{\frac{7}{2}}}^2 \bigintsss_{\sqrt{2x^{2}-7}}^{\sqrt{5-x^{2}}}\bigintsss_{0}^{\sqrt{\frac{z^{2}-2x^{2}+7}{3}}}8dydzdx  \\
+\bigintsss_{0}^{\sqrt{\frac{7}{2}}}\bigintsss_{0}^{\sqrt{5-x^{2}}}\bigintsss_{0}^{\sqrt{\frac{z^{2}-2x^{2}+7}{3}}}8dydzdx \approx 51.8616
\end{multline}
}
%\vskip0.2cm

La evaluación numérica de las anteriores integrales se obtuvo mediante la 
ejecución de las siguientes instrucciones:
%\vskip0.2cm


\vspace{0.6em}
\noindent
\begin{tcolorbox}[breakable,enhanced,title=\bf Instrucción en \verb"Mathematica", top=2pt,bottom=2pt,boxsep=0pt,before skip=2pt,after skip=0pt]
\begin{Verbatim}[formatcom=\letraverbatim]
NIntegrate[8,{x,0,Sqrt[7/2]},{y,0,Sqrt[(7-2*x^2)/3],
{z,0,Sqrt[9-2*x^2-y^2]}]+NIntegrate[8,{x,0,Sqrt[7/2]},
{y,Sqrt[(7-2*x^2)/3],Sqrt[4-x^2]},{z,Sqrt[2*x^2
+3*y^2-7],Sqrt[9-2*x^2-y^2]}]+NIntegrate[8,
{x,Sqrt[7/2],2},{y,0,Sqrt[4-x^2]},{z,Sqrt[2*x^2+
3*y^2-7],Sqrt[9-2*x^2-y^2]}]
\end{Verbatim}
\end{tcolorbox}
\vskip0.4cm	
%NIntegrate[8,{x,0,2},{z,Sqrt[5-x^2],Sqrt[9-2*x^2]},
%{y,0,Sqrt[9-2*x^2-z^2]}]+NIntegrate[8,{x,0,Sqrt[7/2]},	%{z,0,Sqrt[5-x^2]},{y,0,Sqrt[(z^2-2*x^2+7)/3]}]+
%NIntegrate[8,{x,Sqrt[7/2],2},{z,Sqrt[2*x^2-7],
%Sqrt[5-x^2]},{y,0,Sqrt[(z^2-2*x^2+7)/3]}]

\item El sólido acotado por las superficies del paraboloide $4y=x^2+z^2$ y la esfera $x^2+y^2+z^2=12$ 
se muestran a continuación: \label{ejercicio8a}

\vskip0.2cm
\noindent
\begin{minipage}{\textwidth}
	\centering
	\captionof{figure}{\small $4y\geq x^2+z^2$ \& $x^2+y^2+z^2\leq 12.$}
	\begin{minipage}{0.23\linewidth}
		\centering
		\includegraphics[width=\linewidth]{Imgn/T_15a.pdf}
	\end{minipage}
	\hspace{1em}
	\begin{minipage}{0.23\linewidth}
		\centering
		\includegraphics[width=\linewidth]{Imgn/T_15b.pdf}
	\end{minipage}
	
	\vspace{0.3em}
	\fuentepropia
\end{minipage}

\vskip0.2cm
Esta figura se consigue ejecutando las siguientes instrucciones:


\vspace{0.6em}
\noindent
\begin{tcolorbox}[breakable,enhanced,title=\bf Instrucción en \verb"Mathematica", top=2pt,bottom=2pt,boxsep=0pt,before skip=2pt,after skip=0pt]	
\begin{Verbatim}[formatcom=\letraverbatim]
RegionPlot3D[4*y>x^2+z^2&&x^2+y^2+z^2<12,{x,-3,3},
{y,0,3.7},{z,-3,3},PlotPoints->90,Mesh->False,
BoxRatios->{1,1,1}]
\end{Verbatim}
\end{tcolorbox}
\vskip0.4cm
	
Eliminando $x^2+z^2$ de las dos ecuaciones, obtenemos que $y^{2}+4y=12$, cuya solución es $y=2$ y $y=-6;$ por la ubicación de la región en el espacio consideramos el caso $y=2.$ Esto implica que en el plano $xz,$ la proyección del sólido está determinada por la ecuación $x^{2}+z^{2}=8;$ esto corresponde a un c\'{\i}rculo
centrado en el origen de radio $4$. Por tanto, utilizando coordenadas rectangulares e integrando primero con respecto a $y,$ la integral que determina el volumen es
\vskip0.2cm
$$\int _{-\sqrt{8}}^{\sqrt{8}}\int _{-\sqrt{8-x^2}}^{\sqrt{8-x^2}}\int _{\frac{1}{4}
\left(x^2+y^2\right)}^{\sqrt{12-x^2-y^2}}dzdydx= \frac{8}{3}\pi(6\sqrt{3}-5)$$
\vskip0.2cm
Esta integral se puede evaluar al ejecutar la siguiente instrucción:


\vspace{0.6em}
\noindent
\begin{tcolorbox}[breakable,enhanced,title=\bf Instrucción en \verb"Mathematica", top=2pt,bottom=2pt,boxsep=0pt,before skip=2pt,after skip=0pt]	
\begin{Verbatim}[formatcom=\letraverbatim]
Integrate[1,{x,-Sqrt[8],Sqrt[8]},{y,-Sqrt[8-x^2],
Sqrt[8-x^2]},{z,(x^2+y^2)/4,Sqrt[12-x^2-y^2]}]
\end{Verbatim}
\end{tcolorbox}

\vskip0.4cm	
Al hacer un cambio a coordenadas cilíndricas, el volumen está dado por
\vskip0.2cm
\begin{equation} \label{ejercicio8}
\int_0^{\sqrt{8}}\int_{0}^{2\pi}\int_{\frac{r^{2}}{4}}^{\sqrt{12-r^{2}}}rdyd\theta dr=\frac{8}{3}\pi\left(6\sqrt{3}-5\right)\!. 
\end{equation}
%\vskip0.2cm
Por otra parte, la ecuación $x^{2}+z^{2}=8$ sugiere usar la parametrización $x=\sqrt{8}u\cos v,$ $z=\sqrt{8}u\sin v$, en donde $u \in [0,1],$ $v \in [0,2\pi]$. Reemplazando $x$ y $z$ en la expresión de la esfera y del paraboloide se obtienen expresiones correspondientes para $y$,
es decir $y=\sqrt{12-x^{2}-z^{2}}= 2\sqrt{3-2u^{2}},$ $y=\frac{x^{2}+z^{2}}{4}= 2u^{2}.$ 
\vskip0.2cm
De esta manera, tendremos dos puntos sobre las superficies mencionadas. Al hacer una combinación convexa de estos puntos se parametriza el sólido
{\footnotesize
\[\left( 1-w\right) \left( \sqrt{8}u\cos v,2u^{2},\sqrt{8}u\sin v\right) \linebreak 	+w\left( \sqrt{8}u\cos v,2\sqrt{3-2u^{2}},\sqrt{8}u\sin v\right)\!;\]
}
por esta razón se considera la función
%\vskip0.2cm
$$ {\bf r}(u,v,w)= \langle 2\sqrt{2}u\cos v,2w\sqrt{3-2u^{2}}-2u^{2}w+2u^{2},2\sqrt{2}u\sin v \rangle,$$
\vskip0.2cm
con $u \in [0,1],$ $v \in [0,2\pi],$ $w \in [0,1].$ De aquí se obtienen los siguientes vectores:
{\small
\begin{eqnarray*}
{\bf r}_{u} &=& \left\langle 2\sqrt{2}\cos v,4u-4uw-\frac{4uw}{\sqrt{3-2u^{2}}},2\sqrt{2}\sin v \right\rangle  \\
{\bf r}_{v} &=& \left\langle  -2\sqrt{2}u\sin v,0,2\sqrt{2}u\cos v\right\rangle  \\
{\bf r}_{w} &=& \left\langle  0,2\sqrt{3-2u^{2}}-2u^{2},0\right\rangle.
\end{eqnarray*}
}
%\vskip0.2cm

Por lo descrito antes, el elemento diferencial de volumen está dado por $\left\vert {\bf r}_{w} \cdot({\bf r}_{u} \times {\bf r}_{v}) \right\vert = \left\vert 16u\left( u^{2}-\sqrt{3-2u^{2}}\right) \right\vert,$ que al integrarlo con los límites establecidos, produce el volumen del mencionado sólido:
%\vskip0.2cm 
$$\int_{0}^{1}\int_{0}^{1}\int_{0}^{2\pi }\left\vert 16u\left( u^{2}-\sqrt{3-2u^{2}}\right) \right\vert dvdudw= \frac{8}{3}\pi \left( 6\sqrt{3}-5\right)\!.$$
%\vskip0.3cm

\item Hallar el volumen del sólido acotado por el semicono con ecuación 
$x=\sqrt{y^2+z^2}$ y el paraboloide $y^2+z^2+4x=12.$

\vskip0.4cm
\noindent
\begin{minipage}{\textwidth}
	\centering
	\captionof{figure}{\small $r\leq x \leq 3-r^2/4$}
	\begin{minipage}{0.21\linewidth}
		\centering
		\includegraphics[width=\linewidth]{Imgn/T_17a.pdf}
	\end{minipage}
	\hspace{1em}
	\begin{minipage}{0.21\linewidth}
		\centering
		\includegraphics[width=\linewidth]{Imgn/T_17b.pdf}
	\end{minipage}
	
	\vspace{0.3em}
	\fuentepropia
\end{minipage}

\vskip0.2cm
Las siguientes instrucciones producen la figura:
%\vskip0.2cm	


\vspace{0.6em}
\noindent
\begin{tcolorbox}[breakable,enhanced,title=\bf Instrucción en \verb"Mathematica", top=2pt,bottom=2pt,boxsep=0pt,before skip=2pt,after skip=0pt]
\begin{Verbatim}[formatcom=\letraverbatim]
RegionPlot3D[x>Sqrt[y^2+z^2]&&y^2+z^2+4*x<12,
{x,0,3.25},{y,-2,2},{z,-2,2},PlotPoints->90,
Mesh->False]
\end{Verbatim}
\end{tcolorbox}


\vskip0.4cm	
Estas superficies se intersectan en una región del plano $yz$
determinada por la ecuación $y^{2}+z^{2}=4.$
Se debe tener presente que $x$ recorre los puntos desde el semicono hasta el paraboloide. También se presenta la integral correspondiente al cambio a coordenadas cilíndricas; por tanto, 
el volumen del sólido está dado por
%\vskip0.2cm 
\vspace{-2mm}
\begin{multline} \label{ejercicio10}
\int_{-2}^2 \int_{-\sqrt{4-y^2}}^{\sqrt{4-y^2}}\int_{\sqrt{y^2+z^{2}}}^{\frac{12-y^{2}-z^{2}}{4}}dxdzdy \\
=\int_{0}^2 \int_{0}^{2\pi }\int_{r}^{3-\frac{r^{2}}{4}}rdxd\theta dr= \frac{14}{3}\pi.
\end{multline}
%\vskip0.3cm
%$$\int_{0}^2 \int_{0}^{2\pi }\int_{r}^{3-\frac{r^{2}}{4}}rdxd\theta dr= \frac{14}{3}\pi .$$
\item Presentamos ahora el sólido acotado por el paraboloide $x+3=y^2+z^2$ 
y el paraboloide hiperbólico $y^2-z^2=9x-3.$

\vskip0.2cm
\noindent
\begin{minipage}{\textwidth}
	\centering
	\captionof{figure}{\small $3-y^2-z^2\leq x$ \& $y^2-z^2\leq 9x-3$}
	\begin{minipage}{0.20\linewidth}
		\centering
		\includegraphics[width=\linewidth]{Imgn/T_16a.pdf}
	\end{minipage}
	\hspace{1em}
	\begin{minipage}{0.20\linewidth}
		\centering
		\includegraphics[width=\linewidth]{Imgn/T_16b.pdf}
	\end{minipage}
	\hspace{1em}
	\begin{minipage}{0.20\linewidth}
		\centering
		\includegraphics[width=\linewidth]{Imgn/T_16c.pdf}
	\end{minipage}
	
	\vspace{0.3em}
	\fuentepropia
\end{minipage}

\vskip0.2cm
Ejecutando las siguientes instrucciones, se consigue la siguiente figura:	
\vskip0.2cm


%\vspace{0.6em}
\noindent
\begin{tcolorbox}[breakable,enhanced,title=\bf Instrucción en \verb"Mathematica", top=2pt,bottom=2pt,boxsep=0pt,before skip=2pt,after skip=0pt]
\begin{Verbatim}[formatcom=\letraverbatim]
RegionPlot3D[x+3>y^2+z^2&&y^2-z^2>9*x-3,{x,-3.5,1.5,
{y,-2,2},{z,-2.2,2.2},PlotPoints->90,Mesh->False,
BoxRatios->{1,1,1}]
\end{Verbatim}
\end{tcolorbox}

\vskip0.4cm
	
Eliminando $x$ de las ecuaciones del paraboloide y del paraboloide hiperbólico se obtiene $4y^{2}+5z^{2}=15,$ que corresponde a una elipse
en el plano $yz;$ esta ecuación define la región de integración en dicho plano. Con esta información se plantea la integral que representa el volumen del
sólido,
%\vskip0.2cm
$$\bigintss_{-\frac{\sqrt{15}}{2}}^{\frac{\sqrt{15}}{2}}\bigintss_{-\sqrt{3-\frac{4y^{2}}{5}}}^{\sqrt{3-\frac{4y^{2}}{5}}}\bigintss_{y^{2}+z^{2}-3}^{\frac{y^{2}-z^{2}+3}{9}}dxdzdy= \frac{5}{2}\sqrt{5}\pi.$$
\vskip0.2cm	
La región del plano $xz$ determinada por la ecuación $4y^{2}+5z^{2}=15,$ se puede representar mediante las expresiones $y=\frac{\sqrt{15}}{2}u\cos v,$ $z=\sqrt{3}u\sin v$, con $u \in [0,1]$ y $v \in [0,2\pi].$ La expresión correspondiente para $x$ en las dos superficies que acotan el sólido son: 
%\vskip0.2cm
\begin{eqnarray*}
x &=&y^{2}+z^{2}-3= \frac{3}{8}u^{2}\cos 2v+\frac{27}{8}u^{2}-3\\
x&=&\frac{y^{2}-z^{2}+3}{9}= \frac{3}{8}u^{2}\cos 2v+\frac{1}{24}u^{2}+\frac{1}{3}
\end{eqnarray*} 
%\vskip0.2cm
Con lo mencionado antes, 
se pueden considerar dos puntos: uno sobre el paraboloide {\small $Q\left( \frac{3}{8}u^{2}\cos 2v+\frac{27}{8}u^{2}-3,\frac{\sqrt{15}}{2}u\cos v,\sqrt{3}u\sin v\right)$} y otro sobre el paraboloide hiperbólico
{\small $P\left(\frac{3}{8}u^{2}\cos2v+\frac{1}{24}u^{2}+\frac{1}{3},\frac{\sqrt{15}}{2}u\cos v,\sqrt{3}u\sin v\right);$}
Con ellos se forma una combinación convexa {\small $(1-w) P +wQ,$ con $w \in [0,1],$} que al simplificarse define la función que parametriza el sólido,
%\vskip0.2cm 
{\small
\begin{multline*} 
{\bf r}(u,v,w)= \bigg\langle \frac{10}{3}w(u^{2}-1)+\frac{1}{8}u^{2}\left(3\cos 2v+\frac{1}{3}\right)+\frac{1}{3}, \\
\frac{\sqrt{15}}{2}u\cos v,\sqrt{3}u\sin v \bigg \rangle.
\end{multline*}
}
%\vskip0.2cm
Derivando parcialmente hallamos los siguientes vectores
%\vskip0.2cm
{\small
\begin{eqnarray*}
{\bf r}_{u} &=&\left\langle \frac{1}{12}u+\frac{3}{4}u\cos 2v+\frac{20}{3}uw,\frac{1}{2}\sqrt{15}\cos v,\sqrt{3}\sin v\right\rangle\\
{\bf r}_{v}&=&\left\langle -\frac{3}{4}u^{2}\sin 2v,-\frac{1}{2}\sqrt{15}u\sin v,\sqrt{3}u\cos v\right\rangle  \\
{\bf r}_{w}&=&\left\langle \frac{10}{3}u^{2}-\frac{10}{3},0,0\right\rangle  \\
{\bf r}_{w}\times {\bf r}_{u}&=& \frac{10}{3} \left\langle 0,-\sqrt{3}\left( u^{2}-1\right) \sin v,\frac{1}{2}\sqrt{15}\left( u^{2}-1\right) \cos v \right\rangle.
\end{eqnarray*}
}
\vskip0.2cm
Por último, se evalúa el producto vectorial triple, que determina el ele\-mento diferencial de volumen
$\left\vert {\bf r}_{v}\cdot \left({\bf r}_{w}\times {\bf r}_{u}\right) \right\vert= \left\vert 5\sqrt{5}u\left( u^{2}-1\right) \right\vert.$ Así, se integra con los límites establecidos y obtendremos el volumen deseado
\vskip0.2cm
\begin{equation}\label{ejercicio9}
\int_{0}^{1}\int_{0}^{1}\int_{0}^{2\pi }\left\vert 5\sqrt{5}u\left(u^{2}-1\right) \right\vert dvdudw= \frac{5}{2}\sqrt{5}\pi.
\end{equation}
\vskip0.2cm

\item Hallar el volumen del sólido que se obtiene de perforar la esfera centrada en el origen de radio $3$ con el cilindro con ecuación $x^2+y^2=2x+y.$

\vskip0.2cm
\noindent
\begin{minipage}{\textwidth}
	\centering
	\captionof{figure}{\small $2x+y \leq x^2+y^2$ \& $x^2+y^2+z^2\leq 9.$}
	\begin{minipage}{0.23\linewidth}
		\centering
		\includegraphics[width=\linewidth]{Imgn/T_19a.pdf}
	\end{minipage}
	\hspace{1em}
	\begin{minipage}{0.20\linewidth}
		\centering
		\includegraphics[width=\linewidth]{Imgn/T_19b.pdf}
	\end{minipage}
	
	\vspace{0.3em}
	\fuentepropia
\end{minipage}

\vskip0.2cm
Ejecutando las siguientes instrucciones se consigue la figura:
\vskip0.2cm	



%\vspace{0.6em}
\noindent
\begin{tcolorbox}[breakable,enhanced,title=\bf Instrucción en \verb"Mathematica", top=2pt,bottom=2pt,boxsep=0pt,before skip=2pt,after skip=0pt]
\begin{Verbatim}[formatcom=\letraverbatim]
RegionPlot3D[2*x+y<x^2+y^2<9&&x^2+y^2+z^2<=9,{x,-3,3},
{y,-3,3},{z,-3,3},PlotPoints->90,Mesh->False]
\end{Verbatim}
\end{tcolorbox}


\vskip0.2cm	
Al integrar en coordenadas rectangulares con el orden $dzdydx$ es necesario evaluar las cuatro integrales triples que a continuación se presentan:
%\vskip0.2cm
\begin{multline*} 
\bigintss_{-3}^{1-\frac{\sqrt{5}}{2}}\bigintss_{-\sqrt{9-x^{2}}}^{\sqrt{9-x^{2}}}\bigintss_{-\sqrt{9-x^{2}-y^{2}}}^{\sqrt{9-x^{2}-y^{2}}}dzdydx+ \\
\bigintsss_{1-\frac{\sqrt{5}}{2}}^{1+\frac{\sqrt{5}}{2}}\bigintsss_{-\sqrt{9-x^{2}}}^{\frac{1}{2}-\sqrt{\frac{5}{4}-\left( x-1\right) ^2 }}\bigintsss_{-\sqrt{9-x^{2}-y^{2}}}^{\sqrt{9-x^{2}-y^{2}}}dzdydx +\\
\bigintsss_{1-\frac{\sqrt{5}}{2}}^{1+\frac{\sqrt{5}}{2}}\bigintsss_{\frac{1}{2}+\sqrt{\frac{5}{4}-\left( x-1\right) ^2 }}^{\sqrt{9-x^{2}}}\bigintsss_{-\sqrt{9-x^{2}-y^{2}}}^{\sqrt{9-x^{2}-y^{2}}}dzdydx\\
+\bigintsss_{1+\frac{\sqrt{5}}{2}}^{3}\bigintsss_{-\sqrt{9-x^{2}}}^{\sqrt{9-x^{2}}}\bigintsss_{-\sqrt{9-x^{2}-y^{2}}}^{\sqrt{9-x^{2}-y^{2}}}dzdydx \approx 92.2262.
\end{multline*}
%\vskip0.2cm
	
Estas integrales se evaluaron numéricamente en \verb"Mathematica":
\vskip0.2cm


%\vspace{0.6em}
\noindent
\begin{tcolorbox}[breakable,enhanced,title=\bf Instrucción en \verb"Mathematica", top=2pt,bottom=2pt,boxsep=0pt,before skip=2pt,after skip=0pt]
\begin{Verbatim}[formatcom=\letraverbatim]
Integrate[1,{x,-3,1-Sqrt[5]/2},{y,-Sqrt[9-x^2],
Sqrt[9-x^2]},{z,-Sqrt[9-x^2-y^2],Sqrt[9-x^2-y^2]}]
+NIntegrate[1,{x,1-Sqrt[5]/2,1+Sqrt[5]/2},
{y,-Sqrt[9-x^2],1/2-Sqrt[5/4-(x-1)^2]},{z,
-Sqrt[9-x^2-y^2],Sqrt[9-x^2-y^2]}]+Integrate[1,
{x,1+Sqrt[5]/2,3},{y,-Sqrt[9-x^2],Sqrt[9-x^2]},
{z,-Sqrt[9-x^2-y^2],Sqrt[9-x^2-y^2]}]+NIntegrate[1,
{x,1-Sqrt[5]/2,1+Sqrt[5]/2},{y,1/2+Sqrt[5/4-(x-1)^2],
Sqrt[9-x^2]},{z,-Sqrt[9-x^2-y^2],Sqrt[9-x^2-y^2]}]
\end{Verbatim}
\end{tcolorbox}

\vskip0.4cm	
El volumen del sólido también puede calcularse como la diferencia entre el volumen de la esfera y
el volumen de la región acotada por la esfera e interior al cilindro, esto es,
%\vskip0.2cm
$$36\pi -\bigintss_{1-\frac{\sqrt{5}}{2}}^{1+\frac{\sqrt{5}}{2}}\bigintss_{\frac{1}{2}- \sqrt{\frac{5}{4}-(x-1)^2 }}^{\frac{1}{2}+\sqrt{\frac{5}{4}-(x-1) ^2 }}\bigintss_{-\sqrt{9-x^{2}-y^{2}}}^{\sqrt{9-x^{2}-y^{2}}}dzdydx \approx 92.2262.$$
%\vskip0.2cm
	
\item El sólido acotado por las figuras de las ecuaciones $z=4-y^{2},$ $z=4-x,$ $z=0,$ $x=0,$ 
se perfora por el cilindro $\left( z-2\right) ^{2}+y^{2}=1.$ Esta región se presentó en la sección \secref{solidohueco3} de la página \pageref{solidohueco3}; esta se muestra en las siguientes figuras.

\vskip0.2cm
\noindent
\begin{minipage}{\textwidth}
	\centering
	\captionof{figure}{\small $1\leq (z-2)^{2}+y^{2}$ \& $0\leq z\leq 4-y^2$}
	\begin{minipage}{0.18\linewidth}
		\centering
		\vspace{-4mm}
		\includegraphics[width=\linewidth]{Imgn/5b.pdf}
	\end{minipage}
	\hspace{1em}
	\begin{minipage}{0.18\linewidth}
		\centering
		\vspace{-4mm}
		\includegraphics[width=\linewidth]{Imgn/5a.pdf}
	\end{minipage}
	\hspace{1em}
	\begin{minipage}{0.19\linewidth}
		\centering
		\vspace{-4mm}
		\includegraphics[width=\linewidth]{Imgn/5c.pdf}
	\end{minipage}
	
	\vspace{0.3em}
	\fuentepropia
\end{minipage}
\vskip0.2cm

El volumen puede obtenerse como la diferencia entre el volumen del sólido sin perforar y el volumen de la parte cilíndrica. Esto es,
\vspace{-2mm}
$$ \int _{-2}^2\int _0^{4-y^2}\int _0^{4-z}dxdzdy-\int _{-1}^1\int _{2-\sqrt{1-y^2}}^{2+\sqrt{1-y^2}}\int _0^{4-z}dxdzdy = \frac{128}{5}-2\pi.$$
%\vskip0.2cm


%\vspace{0.6em}
\noindent
\begin{tcolorbox}[breakable,enhanced,title=\bf Instrucción en \verb"Mathematica", top=2pt,bottom=2pt,boxsep=0pt,before skip=2pt,after skip=0pt]
\begin{Verbatim}[formatcom=\letraverbatim]
Integrate[1,{y,-2,2},{z,0,4-y^2},{x,0,4-z}]
-Integrate[1,{y,-1,1},{z,2-Sqrt[1-y^2],2+
Sqrt[1-y^2]},{x,0,4-z}]
\end{Verbatim}
\end{tcolorbox}

\vskip0.4cm

El mencionado sólido se parametrizó mediante la expresión (\ref{solidohueco2}).
Con las siguientes instrucciones se define la parametrización, se halla el elemento diferencial de 
volumen y se evalúan las integrales con los límites establecidos:
\vskip0.2cm

%\vspace{0.6em}
\noindent
\begin{tcolorbox}[breakable,enhanced,title=\bf Instrucción en \verb"Mathematica", top=2pt,bottom=2pt,boxsep=0pt,before skip=2pt,after skip=0pt]	
\begin{Verbatim}[formatcom=\letraverbatim]
r1={0,(1-u)*Cos[t]+u*(2-4*t/Pi),(1-u)*(2+Sin[t])
+16*t*u/Pi*(1-t/Pi)};
r2={0,(1-u)*Cos[t]+u*(4*(t-Pi)/Pi-2),(1-u)*(2+Sin[t])};
r3={4-(1-u)*(2+Sin[t])-16*t*u/Pi*(1-t/Pi),(1-u)*Cos[t]
+u*(2-4*t/Pi),(1-u)*(2+Sin[t])+16*t*u/Pi*(1-t/Pi)};
r4={4-((1-u)*(2+Sin[t])),(1-u)*Cos[t]
+u*(4*(t-Pi)/Pi-2),(1-u)*(2+Sin[t])};
{R1,R2}={(1-w)*r1+w*r3,(1-w)*r2+w*r4};
Integrate[Dot[D[R1,w],Cross[D[R1,u],D[R1,t]]],
{t,0,Pi},{u,0,1},{w,0,1}]+Integrate[Dot[D[R2,w],
Cross[D[R2,u],D[R2,t]]],{t,Pi,2*Pi},{u,0,1},{w,0,1}]
\end{Verbatim}
\end{tcolorbox}

\end{enumerate}
\vskip0.4cm

\begin{ejer}
Determine el volumen del sólido acotado por las figuras de las expresiones señaladas.
Hágalo de diversas maneras, incluyendo la parametrización de la región. Se agrega la figura y las instrucciones para conseguirlo.
\vskip0.2cm

\begin{enumerate}[leftmargin=1.2em, itemindent=0pt,labelwidth=\labelwidth,labelsep=0.5em]
\item $x^2+y^2=4+z$, $4z=(4-x^2-y^2)e^{1-x^2-y^2}$. 
\normalfont
\vskip0.2cm


%\vspace{0.6em}
\noindent
\begin{tcolorbox}[breakable,enhanced,title=\bf Instrucción en \verb"Mathematica", top=2pt,bottom=2pt,boxsep=0pt,before skip=2pt,after skip=0pt]
\begin{Verbatim}[formatcom=\letraverbatim]
RegionPlot3D[x^2+y^2-4<=z&&z<=(4-x^2-y^2)*Exp[1
-x^2-y^2],{x,-2,2},{y,-2,2},{z,-4,12},
PlotPoints->90,Mesh->False]
\end{Verbatim}
\end{tcolorbox}

\vskip0.4cm

\noindent
\begin{minipage}{\textwidth}
	\centering
	\captionof{figure}{\small $x^2+y^2-4\leq z \leq (4-x^2-y^2)e^{1-x^2-y^2}$}
	\begin{minipage}{0.22\linewidth}
		\centering
		\includegraphics[width=\linewidth]{Imgn/T_42a.pdf}
	\end{minipage}
	\hspace{1em}
	\begin{minipage}{0.22\linewidth}
		\centering
		\includegraphics[width=\linewidth]{Imgn/T_42b.pdf}
	\end{minipage}
	\hspace{1em}
	\begin{minipage}{0.22\linewidth}
		\centering
		\includegraphics[width=\linewidth]{Imgn/T_42c.pdf}
	\end{minipage}
	
	\vspace{0.3em}
	\fuentepropia
\end{minipage}
\vskip0.2cm
		
\item $z=x^2+y^2,$ $4-z=3(x^2+y^2)^3.$
\vskip0.2cm 
\normalfont


%\vspace{0.6em}
\noindent
\begin{tcolorbox}[breakable,enhanced,title=\bf Instrucción en \verb"Mathematica", top=2pt,bottom=2pt,boxsep=0pt,before skip=2pt,after skip=0pt]
\begin{Verbatim}[formatcom=\letraverbatim]
RegionPlot3D[x^2+y^2<=z<=4-3*(x^2+y^2)^3&&x^2
+y^2<=1,{x,-1,1},{y,-1,1},{z,0,4},Mesh->False,
PlotPoints->90,PlotStyle->RGBColor[0.3,0.6,0.3]]
\end{Verbatim}
\end{tcolorbox}

\vskip0.2cm

\noindent
\begin{minipage}{\textwidth}
	\centering
	\captionof{figure}{\small $x^2+y^2\leq z\leq 4-3(x^2+y^2)^3$}
	\begin{minipage}{0.22\linewidth}
		\centering
		\includegraphics[width=\linewidth]{Imgn/T_40a.pdf}
	\end{minipage}
	\hspace{1em}
	\begin{minipage}{0.22\linewidth}
		\centering
		\includegraphics[width=\linewidth]{Imgn/T_40b.pdf}
	\end{minipage}
	\hspace{1em}
	\begin{minipage}{0.22\linewidth}
		\centering
		\includegraphics[width=\linewidth]{Imgn/T_40c.pdf}
	\end{minipage}
	
	\vspace{0.3em}
	\fuentepropia
\end{minipage}
\vskip0.2cm
		
\item $z=4-x^2-y^2,$ $z=3(x^2+y^2)^6.$ 

\vskip0.2cm
\noindent
\begin{minipage}{\textwidth}
	\centering
	\captionof{figure}{\small $3(x^2+y^2)^6 \leq z \leq 4-x^2-y^2$} 
	\begin{minipage}{0.20\linewidth}
		\centering
		\includegraphics[width=\linewidth]{Imgn/T_41a.pdf}
	\end{minipage}
	\hspace{1em}
	\begin{minipage}{0.20\linewidth}
		\centering
		\includegraphics[width=\linewidth]{Imgn/T_41b.pdf}
	\end{minipage}
	
	\vspace{0.3em}
	\fuentepropia
\end{minipage}
\vskip0.2cm
 
\normalfont


%\vspace{0.6em}
\noindent
\begin{tcolorbox}[breakable,enhanced,title=\bf Instrucción en \verb"Mathematica", top=2pt,bottom=2pt,boxsep=0pt,before skip=2pt,after skip=0pt]
\begin{Verbatim}[formatcom=\letraverbatim]
RegionPlot3D[3*(x^2+y^2)^6<=z<=4-x^2-y^2&&x^2
+y^2<=1,{x,-1,1},{y,-1,1},{z,0,4},PlotPoints->90,
Mesh->False,PlotStyle->RGBColor[0.6,0.2,0.2]]
\end{Verbatim}
\end{tcolorbox}

\vskip0.2cm
\end{enumerate}
\end{ejer}
   
\chapter{Integración de campos vectoriales} \chapter{Integración de campos vectoriales}\label{chap:stokes}

Los campos vectoriales se presentaron con algunas generalidades en la sección \secref{funcionesenelespacio}. \label{integracioncamposvectoriales}
Ahora nos ocuparemos de algunos resultados relativos a la integración de estas funciones.

\section{Integrales de línea en campos vectoriales}

Algunos libros como \cite{stewart_calculo_2018} o \cite{thomas_calculo_2015} motivan la definición de integrales sobre una curva mediante conceptos físicos como el trabajo. En tal sentido, una fuerza cons\-tante ${\bf F}$ mueve una partícula sobre el segmento que une al punto $A$ con el punto $B$. Si $W$ es el trabajo realizado, entonces $W={\bf F}\cdot \overrightarrow{AB}$.

Abordaremos inicialmente el caso bidimensional.
Si una fuerza no constante está determinada por ${\bf F}(x,y)= \langle P(x,y),Q(x,y)\rangle $
y el movimiento no es rectilíneo, sino que se hace a través de una trayectoria la cual está dada
${\bf r}(t) =\left\langle f(t),g(t) \right\rangle $, con $t\in [a,b]$, siendo $f$ y $g$ funciones continuas en $[a,b]$.
En cualquier punto de la curva $C$ la fuerza está dada por ${\bf F}({\bf r}(t))$.

\begin{figure}[H]
  \centering
  \caption{Curva en $\mathbb{R}^2$}
  \includegraphics[width=0.8\textwidth]{Fig/170.pdf}
  \fuentepropia
\end{figure}

Consideremos una partición del intervalo $[a,b]$, esto es, un conjunto finito de puntos
$\{t_0, t_1, t_2, \dots, t_{n-1}, t_n\}$, con $a = t_0 < t_1 < t_2 < \cdots < t_{n-1} < t_n = b$.

Entonces, el vector $\overrightarrow{A_{i-1}A_i}$ está dado por
\[
  \overrightarrow{A_{i-1}A_i} = \mathbf{r}(t_i) - \mathbf{r}(t_{i-1}) =
  \left\langle f(t_i) - f(t_{i-1}),\; g(t_i) - g(t_{i-1}) \right\rangle\!.
\]

Por el teorema del valor medio, existen $t_i^{\ast} \in [t_{i-1}, t_i]$ y
$t_i^{\ast\ast} \in [t_{i-1}, t_i]$, tales que
\[
  \overrightarrow{A_{i-1}A_i} =
  \left\langle f'(t_i^{\ast}),\; g'(t_i^{\ast\ast}) \right\rangle \, \Delta t_i.
\]

Para cada $i$, consideremos $\mathbf{F}_i$ como una aproximación al vector fuerza
$\mathbf{F}(\mathbf{r}(t))$ sobre la curva $C$ desde el punto $A_{i-1}$ hasta $A_i$. Por lo tanto, el trabajo realizado por $\mathbf{F}$ para trasladar una partícula desde $A_{i-1}$ hasta $A_i$ está dado por:
{\small
  \begin{multline*}
    \Delta W_i =
    \left\langle
    P\left( \left\langle f(t_i^{\ast}),\; g(t_i^{\ast}) \right\rangle \right),
    Q\left( \left\langle f(t_i^{\ast\ast}),\; g(t_i^{\ast\ast}) \right\rangle \right)
    \right\rangle \cdot
    \left\langle f'(t_i^{\ast}),\; g'(t_i^{\ast\ast}) \right\rangle \, \Delta t_i
  \end{multline*}
}

Una aproximación al trabajo total para desplazar una partícula sobre la curva $C$ desde $\mathbf{r}(a)$ hasta $\mathbf{r}(b)$ está dada por:
\[
  W \approx \sum_{i=1}^{n} \Delta W_i.
\]

Si $n \to \infty$, entonces $\Delta t_i \to 0$, por lo tanto:
\begin{eqnarray}
  W &=& \lim_{n\rightarrow\infty}\sum_{i=1}^{n}\Delta W_{i} \nonumber \\
  &=& \int_{a}^{b}\left\langle P\left( \left\langle f(t) ,g\left(t\right) \right\rangle \right) ,Q\left( \left\langle f(t)
  ,g(t) \right\rangle \right) \right\rangle \cdot \left\langle f^{\prime }(t) ,g^{\prime }(t) \right\rangle dt \label{integraldelinea2}
\end{eqnarray}
Simplificando la ecuación \ref{integraldelinea2} se precisa la expresión que determina el siguiente concepto.

\begin{defi}
  Si $C$ es una curva suave en el plano parametrizada por la función ${\bf r}(t)=\langle f(t),g(t)\rangle$, con $t \in [a,b]$
  y ${\bf F}(x,y)= \langle P(x,y),Q(x,y)\rangle $ es un campo vectorial con funciones componentes continuas, se define la integral del campo vectorial ${\bf F}$ sobre la curva $C$ como
  \begin{equation} \label{integraldelinea3}
    \int_C {\bf F} \cdot d {\bf r} =\int_{a}^{b}{\bf F}({\bf r}(t)) \cdot {\bf r}^{\prime }(t) dt
  \end{equation} \index{Integrales de línea}
\end{defi}

Este concepto se puede extender a campos vectoriales tridimensionales. Para una función vectorial
{\small ${\bf r}(t)=\langle f(t),g(t),h(t)\rangle$,} con {\small $t \in [a,b]$;}
que parametriza a una curva suave $C$ y {\footnotesize ${\bf F}(x,y,z)= \langle P(x,y,z),Q(x,y,z), R(x,y,z)\rangle$} es un campo vectorial con funciones componentes continuas, entonces la integral del campo vectorial ${\bf F}$ sobre la curva $C$ está dada por
\begin{equation}\label{integraldelinea7}
  \scalebox{0.82}{%
    $\displaystyle \int_C {\bf F} \cdot d {\bf r} =\int_{a}^{b}{\bf F}(\langle f(t),g(t),h(t)\rangle) \cdot \langle f'(t),g'(t),h'(t)\rangle dt =\int_{a}^{b}{\bf F}({\bf r}(t)) \cdot {\bf r}^{\prime }(t)$
  }
\end{equation}

\begin{ejemplo} \label{integraldelinea}
  Evaluar la integral de línea de ${\bf F}\left( x,y,z\right) =\left\langle xy,yz,-x^{2}\right\rangle $ sobre la curva que une los puntos $\left(-1,1,3\right) $ con
  $\left( 1,1,1\right)$, considerando las siguientes parametrizaciones:

  \begin{itemize}
    \item ${\bf r}(t) =\left\langle t^{2},t,2-t\right\rangle $, con $t\in[-1,1]$. En este caso,\\ ${\bf F}\left({\bf r}(t) \right) \cdot {\bf r}^{\prime }(t)= \left\langle t^{3},t\left( 2-t\right) ,-t^{4}\right\rangle \cdot \left\langle 2t,1,-1\right\rangle = 3t^{4}-t^{2}+2t$.
      \begin{eqnarray*}
        \int_{C}{\bf F}\cdot d{\bf r} &=&\int_{-1}^{1}{\bf F}\left({\bf r}(t) \right) \cdot {\bf r}^{\prime }(t)
        = \int_{-1}^{1}(3t^{4}-t^{2}+2t) dt= \frac{8}{15}
      \end{eqnarray*}

    \item ${\bf r}\left( r\right) =\left\langle 2t-1,1,3-2t\right\rangle $, con $t\in [0,1];$ Para este caso ${\bf F}\left({\bf r}(t) \right) \cdot {\bf r}^{\prime }(t)= \left\langle 2t-1,3-2t,-\left( 2t-1\right) ^{2}\right\rangle \cdot \left\langle 2,0,-2\right\rangle = 8t^{2}-4t$
      \begin{eqnarray*}
        \int_{C}{\bf F}\cdot d{\bf r}&=& \int_{0}^{1}{\bf F}\left({\bf r}(t) \right) \cdot {\bf r}^{\prime }(t) dt = \int_{0}^{1} (8t^{2}-4t) dt = \frac{2}{3}
      \end{eqnarray*}
  \end{itemize}
  Nótese que, para el mismo campo vectorial, la curva une los mismos puntos y el valor de la integral es distinto.
\end{ejemplo}

Los cálculos anteriores se evalúan en \verb"Mathematica" ejecutando los siguien\-tes comandos:

\begin{tcolorbox}[title=\bf Instrucción \verb"Mathematica"]\begin{Verbatim}[formatcom=\letraverbatim]
r={x,y,z}={t^2,t,2-t}
Integrate[Dot[{x*y,y*z,-x^2},D[r,t]],{t,-1,1}]
r={x,y,z}={2*t-1,1,3-2*t}
Integrate[Dot[{x*y,y*z,-x^2},D[r,t]],{t,0,1}]
\end{Verbatim}
\end{tcolorbox}

\subsection{Teorema fundamental para integrales de línea}
De acuerdo con \cite{acosta_gempeler_doce_2020},
para un campo escalar $\phi$ con dominio $D \subseteq \mathbb{R}^3$ y $C$ una curva en
$D$ que se parametriza mediante una función vectorial continua ${\bf r} : [a, b]\to D$. Para calcular la circulación del campo vectorial ${\bf F}=\nabla \phi$ a lo largo de $C$, se sigue que
\begin{eqnarray*}
\int_{C}F\cdot d{\bf r}&=&\int_{C}\nabla\phi\cdot dr=\int_{a}^{b}\nabla \phi\left({\bf r}(t)\right)\cdot {\bf r}^{\prime}(t) dt\\
&=& \int_{a}^{b} \frac{d}{dt} \left(\phi \left({\bf r}(t)\right) \right) dt= \phi \left({\bf r}(b)\right) -\phi\left({\bf r}\left( a\right) \right)\!.
\end{eqnarray*}
Por lo tanto, si ${\bf F}$ es un campo gradiente, entonces la circulación de ${\bf F}$ a lo largo de toda curva dentro del dominio de ${\bf F}$ solo depende de los puntos extremos de dicha curva.

\begin{ejemplo}
Evaluar la integral de línea de {\small ${\bf F}\left( x,y,z\right) =\left\langle y^{2}-2xz,2xy,-x^{2}\right\rangle $}
sobre la curva que une los puntos $\left( -1,1,3\right) $ con $\left( 1,1,1\right)$,
considerando las siguien\-tes parametrizaciones.

\begin{itemize}
  \item ${\bf r}(t) =\left\langle t,t^{2},2-t\right\rangle $, con $t\in [-1,1]$
    \begin{eqnarray*}
      \int_{C}{\bf F}\cdot dr &=& \int_{-1}^{1}{\bf F}\left({\bf r}(t) \right) \cdot {\bf r}^{\prime }(t) dt\\
      &=& \int_{-1}^{1}\left\langle t^{4}-2t\left( 2-t\right) ,2t^{3},-t^{2}\right\rangle \cdot \left\langle 1,2t,-1\right\rangle dt \\
      &=&\int_{-1}^{1}\left(5t^{4}+3t^{2}-4t\right) dt= 4
    \end{eqnarray*}

  \item ${\bf r}(t) =\left\langle 2t-1,1,3-2t\right\rangle $, con $t\in [0,1]$
    \begin{eqnarray*}
      \int_{C}{\bf F}\cdot d{\bf r}&=&\int_{0}^{1}{\bf F}\left({\bf r}(t) \right) \cdot{\bf r}^{\prime }(t) dt \\
      %&=& \int_{0}^{1}\left\langle 1-2\left( 2t-1\right) \left( 3-2t\right) ,2\left(2t-1\right) ,-\left( 2t-1\right) ^{2}\right\rangle \cdot \left\langle 2,0,-2\right\rangle dt \\
      &=& \int_{0}^{1}\left( 24t^{2}-40t+16\right)dt= 4.
    \end{eqnarray*}
\end{itemize}

Nótese que ${\bf F}$ es conservativo con función potencial $\phi(x,y,z)=xy^2-x^2z$, con lo cual
$\int_{C}{\bf F}\cdot d{\bf r}= \phi(1,1,1)-\phi(-1,1,3)=4$.
Compárese con los resultados del ejemplo (\ref{integraldelinea}).
\end{ejemplo}

En \verb"Mathematica" se pueden evaluar los cálculos anteriores ejecutando los siguientes comandos:

\begin{tcolorbox}[title=\bf Instrucción \verb"Mathematica"]\begin{Verbatim}[formatcom=\letraverbatim]
r={x,y,z}={t,t^2,2-t}
Integrate[Dot[{y^2-2*x*z,2*x*y,-x^2},D[r,t]],{t,-1,1}]
r={x,y,z}={2*t-1,1,3-2*t}
Integrate[Dot[{y^2-2*x*z,2*x*y,-x^2},D[r,t]],{t,0,1}]
\end{Verbatim}
\end{tcolorbox}

La función potencial se obtiene en \verb"Mathematica" con el siguiente proceso.

\begin{tcolorbox}[title=\bf Instrucción \verb"Mathematica"]\begin{Verbatim}[formatcom=\letraverbatim]
{P,Q,R}={y^2-2*x*z,2*x*y,-x^2};
DSolve[{D[u[x,y,z],x]==P,D[u[x,y,z],y]==Q,
D[u[x,y,z],z]==R},u[x,y,z],x,y,z]
\end{Verbatim}
\end{tcolorbox}

\begin{ejemplo}
Evaluar la integral de línea de {\small ${\bf F}\left( x,y,z\right) =\left\langle 4 x^3 y,x^4+3 y^2 z,y^3 \right\rangle$}
sobre la curva $\alpha$ parametrizada por ${\bf r}(t)=\langle 2 t-1,t+2,9 t^2 \rangle$ con $t \in [-1,2]$.

El campo es conservativo, la función potencial es $\phi(x,y,z)=x^4 y+y^3z+C$, con lo cual
\[\int_\alpha{\bf F} \cdot d{\bf r}=\phi({\bf r}(2))-\phi({\bf r}(-1))=\phi(3,4,36)-\phi(-3,1,9)=2538.\]
\end{ejemplo}

Con las siguientes instrucciones se evalúa la integral de línea.

\begin{tcolorbox}[title=\bf Instrucción \verb"Mathematica"]\begin{Verbatim}[formatcom=\letraverbatim]
F={4x^3y,x^4+3y^2z,y^3};
r[t_]:={2t-1,2+t,3*3*t^2}
{x,y,z}=r[t];
Integrate[Dot[F,D[r[t],t]],{t,-1,2}]
\end{Verbatim}
\end{tcolorbox}

Con los siguientes comandos se halla la función potencial y se utiliza el teorema fundamental para integrales de línea.

\begin{tcolorbox}[title=\bf Instrucción \verb"Mathematica"]\begin{Verbatim}[formatcom=\letraverbatim]
Clear[x,y,z,r,t]
{P,Q,R}={4x^3y,x^4+3y^2z,y^3};
\[Phi][x_,y_,z_]=u[x,y,z]/.First@DSolve[{D[u[x,y,z],
x]==P,D[u[x,y,z],y]==Q,D[u[x,y,z],z]==R},
u[x,y,z],x,y,z]
\[Phi][3,4,36]-\[Phi][-3,1,9\]
\end{Verbatim}
\end{tcolorbox}

En la ecuación \ref{integraldelinea3} al considerar el vector tangente unitario y $s$ es la longitud de la curva, haciendo el cambio ${\bf T}(t)= \frac{{\bf r}^\prime(t)}{\Vert {\bf r}^\prime(t) \Vert} = \frac{d{\bf r}}{ds}ds$,
se obtiene una expresión equivalente a la fórmula \ref{integraldelinea7} en términos del vector tangente unitario.

\begin{figure}[H]
\centering
\caption{Vectores ${\bf F}(x,y,z)$, ${\bf T}$ y ${\bf n}$}
\includegraphics[width=0.8\textwidth]{Fig/171.png}
\fuentepropia
\end{figure}

{\small
\begin{equation} \label{integraldelinea4}
\int_C {\bf F} \cdot d{\bf r}=\int_C{\bf F}({\bf r}(t))\cdot\frac{{\bf r}^\prime(t)}{\Vert{\bf r}^\prime(t)\Vert}ds
=\int_C{\bf F}({\bf r}(t)) \cdot {\bf T}(t)ds=\int_C{\bf F}\cdot {\bf T}ds
\end{equation}
}

\subsection{Integrales de flujo y de circulación para campos de velocidad}

En la ecuación \ref{integraldelinea4}, según \cite{thomas_calculo_2015},
si ${\bf F}$ representa el campo de velocidad de un fluido que fluye a través de una región en el espacio, la integral de ${\bf F}\cdot {\bf T}$ a lo largo de una curva $C$ nos proporciona el \emph{flujo} del fluido a lo largo de la curva, o la \emph{circulación} alrededor de $C$.

\begin{defi}
Si $C$ es una curva suave parametrizada por ${\bf r}(t)$ y sea ${\bf F}(x,y)$ un campo vectorial cuyas funciones componentes son continuas. La circulación de ${\bf F}$ a lo largo de la curva $C$
es \[\int_C{\bf F} \cdot {\bf T} ds.\] \index{Circulación}
\end{defi}

En el plano $xy$ consideremos una curva cerrada simple $C$, la razón con que un fluido entra o sale de una región acotada por $C$ se determina de la siguiente manera.

\begin{defi}
Si $C$ es una curva cerrada, simple y suave en el dominio de un campo vectorial ${\bf F}(x,y) = P(x, y)\vi + Q(x, y)\vj$, cuyas funciones componentes son continuas, el flujo de {\bf F} a través de $C$ es \[\text{Flujo de ${\bf F}$ a través de $C$:}\int_C{\bf F} \cdot {\bf n} ds\]
En donde {\bf n} denota el vector normal unitario a la curva $C$ que apunta hacia afuera.
\end{defi}

Una exposición más detallada se encuentra en \cite{thomas_calculo_2015}.

Al término ${\bf F}(x(t),y(t)) \cdot {\bf T}(t)$ se le denomina la \emph{densidad de circulación} del campo ${\bf F}$ en el punto $(x(t),y(t))$ de la curva $C$. La \emph{densidad de flujo}
del campo ${\bf F}$, en el punto $(x(t),y(t))$ de la curva $C$, se define como la componente normal de ${\bf F}$, es decir,
${\bf F}(x(t),y(t)) \cdot {\bf n}(t)$. \index{Densidad!de circulación} \index{Densidad!de flujo}
De modo que integrando la densidad de circulación sobre la curva se obtiene la circulación e integrando la densidad de flujo se obtiene el flujo.

Si ${\bf T}$ y ${\bf n}$ denotan los vectores tangente y normal unitario a la curva $\alpha$
en el punto $(x(t),y(t))$ entonces,
${\bf T}(x,y)= \frac{1}{\sqrt{(x^\prime(t))^2+y^\prime(t))^2}}\la x^\prime(t), y^\prime(t)\ra$ y
${\bf n}(x,y)= \frac{1}{\sqrt{(x^\prime(t))^2+y^\prime(t))^2}}\la -y^\prime(t), x^\prime(t)\ra$.

\subsection{Densidad rotacional y densidad de expansión}
En \cite{acosta_gempeler_doce_2020}se define la densidad rotacional como la circulación por unidad de área, y la densidad de expansión es el flujo por unidad de área. Analizaremos estos conceptos en campos vectoriales bidimensionales. Posteriormente, se extienden a campos tridimensionales.

Consideremos ${\bf F} (x,y)=P(x,y)\vi+Q(x,y)\vj$ un campo de velocidad del flujo de un fluido en el plano tal que las primeras derivadas parciales de $P(x,y) $ y $Q(x,y) $ son continuas en cada punto del rectángulo $R$.

\begin{figure}[H]
\centering
\caption{La razón del flujo a través de cada lado en la dirección indicada}
\includegraphics[width=0.8\textwidth]{Fig/172.pdf}
\fuentepropia
\end{figure}

Según \cite{thomas_calculo_2015}, se desea medir la tasa en un punto del giro de una rueda con paletas, cuyo eje es perpendicular al plano, en un fluido que se mueve en una región plana. Esto nos da una idea de la forma en que circula un fluido con respecto a los ejes, localizados en diferentes puntos, perpendiculares a la región. En ocasiones, los físicos se refieren a esto como la densidad de circulación de un campo vectorial {\bf F}  en un punto.

Supongamos un rectángulo $R$ con lados paralelos a los ejes coordenados y con longitudes $\Delta x$ y $\Delta y$. Si las funciones componentes
$P(x,y) $ y $Q(x,y) $ preservan el signo en una pequeña región del plano que contiene al rectángulo $R$.
Para este procedimiento supondremos que los componentes de ${\bf F} $ son positivas.
La razón de circulación de ${\bf F} $ alrededor de la frontera de $R$ es la suma de las razones de flujo a lo largo de los lados en dirección tangencial. Para la orilla inferior, la razón de flujo es aproximadamente el producto de la densidad de circulación de ${\bf F}$ en el punto medio para la longitud del segmento, es decir,
${\bf F} (x,y)\cdot\vi\Delta x=P(x,y)\Delta x$.

Las razones de flujo son positivas o negativas, dependiendo de las componentes de
${\bf F} $. Aproximamos la razón de circulación neta alrededor de la frontera rectangular de $R$ sumando las razones de flujo a lo largo de los cuatro lados, como lo indican los siguientes productos punto.

\begin{tabular}{rl}
Arriba &  ${\bf F} \left( x,y+\Delta y\right) \cdot (-\vi) \Delta x=-P\left( x,y+\Delta y\right) \Delta x$ \\
Abajo & ${\bf F} (x,y) \cdot \vi\Delta x=P(x,y) \Delta x$ \\
Lado derecho & ${\bf F} ( x+\Delta x,y) \cdot \vj\Delta y=Q\left(x+\Delta x,y\right) \Delta y$ \\
Lado izquierdo & ${\bf F} (x,y) \cdot \left( -\vj\right) \Delta y=-Q(x,y) \Delta y$.
\end{tabular}

Al sumar estos elementos y aplicar el teorema del valor medio, se sigue que

\begin{multline}
Q( x+\Delta x,y) \Delta y-Q(x,y) \Delta y-P\left(x,y+\Delta y\right) \Delta x+P(x,y) \Delta x \\
=(Q( x+\Delta x,y) -Q(x,y)) \Delta y-\left(P\left( x,y+\Delta y\right) -P(x,y) \right) \Delta x\\
\approx \left( \frac{\partial Q}{\partial x}(x^{\ast },y)\Delta x\right) \Delta y-\left( \frac{\partial P}{\partial y}\left(x,y^{\ast }\right) \Delta y\right) \Delta x \
\end{multline}
Para algún $x^* \in (x,x+\Delta x)$ e $y^* \in (y,y+\Delta y)$.
Es por esto que la razón de circulación alrededor del rectángulo $R$ es

\[\left( \frac{\partial Q}{\partial x}(x^{\ast },y) \Delta x\right) \Delta y-\left( \frac{\partial P}{\partial y}\left( x,y^{\ast }\right) \Delta y\right) \Delta x\]

Al dividir  el área de $R$ se obtiene una estimación de la razón de circulación por unidad de área o densidad rotacional para el rectángulo

\[\frac{\left( \frac{\partial Q}{\partial x}(x^{\ast },y) \Delta x\right) \Delta y-\left( \frac{\partial P}{\partial y}\left( x,y^{\ast}\right) \Delta y\right) \Delta x}{\Delta x\Delta y}=\frac{\partial Q}{\partial x}(x^{\ast },y) -\frac{\partial P}{\partial y}\left(x,y^{\ast }\right)\!. \]

Lo anterior motiva la siguiente definición.

\begin{defi}
Sea ${\bf F} (x,y)=P(x,y) \vi+Q(x,y) \vj$ un campo vectorial cuyas funciones componentes poseen primeras derivadas parciales continuas.
{\bf La densidad rotacional} del campo vectorial ${\bf F} (x,y)$ en el punto $(x,y)$ está dada por
\begin{equation} \label{integraldelinea5}
\frac{\partial Q}{\partial x}(x,y) -\frac{\partial P}{\partial y}(x,y)
\end{equation}
Este término se denomina la componente $\vk$ del rotacional y se denota
$({\rm rot}{\bf F} )\cdot \vk$.
\index{Densidad!rotacional}
\end{defi}

De manera similar al caso anterior, supongamos que las componentes del campo preservan el signo en toda la pequeña del rectángulo $R$; se desea determinar la razón con la que el fluido sale de $R$.

\begin{figure}[H]
\centering
\caption{La razón del flujo a través de cada lado en la dirección indicada}
\includegraphics[width=0.8\textwidth]{Fig/173.pdf}
\fuentepropia
\end{figure}

De manera aproximada, la razón del flujo neto a través de la frontera del rectángulo $R$ corresponde a la suma de las razones de flujo a través de los cuatro lados de $R$. En el lado inferior, el flujo es aproximadamente el producto de la densidad de flujo de ${\bf F}$ en el punto medio por la longitud del segmento, es decir
${\bf F} (x,y) \cdot \left( -\vj\right) \Delta x$.
Por lo anterior, se consideran los siguientes productos punto para el flujo de salida

\begin{tabular}{rl}
Arriba & ${\bf F} \left( x,y+\Delta y\right) \cdot \mathbf{j}\Delta x=Q\left(x,y+\Delta y\right) \Delta x$ \\
Abajo & ${\bf F} (x,y) \cdot \left( -\vj\right) \Delta x=-Q\left(x,y\right) \Delta x$ \\
Lado derecho & ${\bf F} ( x+\Delta x,y) \cdot \vi\Delta y=P\left(x+\Delta x,y\right) \Delta y$ \\
Lado izquierdo & ${\bf F} (x,y) \cdot \left( -\vi\right) \Delta y=-P(x,y) \Delta y$
\end{tabular}

Al sumar estos elementos y aplicar el teorema del valor medio, se sigue que
\begin{multline}
P( x+\Delta x,y) \Delta y-P(x,y) \Delta y+Q\left(x,y+\Delta y\right) \Delta x-Q(x,y) \Delta x \\
= \left( P( x+\Delta x,y) -P(x,y) \right)\Delta y+\left( Q\left( x,y+\Delta y\right) -Q(x,y) \right)\Delta x \\
\approx \left( \frac{\partial P}{\partial x}(x^{\ast },y)\Delta x\right) \Delta y+\left( \frac{\partial Q}{\partial y}\left(x,y^{\ast }\right) \Delta y\right) \Delta x
\end{multline}
Para algún $x^* \in (x,x+\Delta x)$ y $y^* \in (y,y+\Delta y)$.
Por lo tanto, el flujo a través de la frontera del rectángulo $R$ es aproximadamente
\[\frac{\partial P}{\partial x}(x^*,y) \Delta x\Delta y+\frac{\partial Q}{\partial y}(x,y^*) \Delta y\Delta x.\]
Al dividir  el área del rectángulo $R$, se establece una estimación del flujo total por unidad de área o densidad de flujo para el rectángulo $R$. Lo anterior motiva la siguiente definición.

\begin{defi}
Si ${\bf F} (x,y)=P(x,y) \vi+Q(x,y) \vj$ es un campo vectorial cuyas funciones componentes poseen primeras derivadas parciales continuas.
La divergencia o \textit{densidad de expansión} del campo vectorial ${\bf F}$ en el punto $(x,y)$.
\begin{equation} \label{integraldelinea6}
{\rm div}{\bf F} =\frac{\partial P}{\partial x}(x,y) +\frac{\partial Q}{\partial y}(x,y) .
\end{equation}
\index{Densidad!de expansión} \index{Divergencia}
\end{defi}

La densidad rotacional y la densidad de expansión para campos vectoriales tridimensionales ${\bf F}(x,y,z) =P(x,y,z)\vi +Q(x,y,z)\vj+R(x,y,z)\vk$,
se pueden generalizar a partir de las ecuaciones (\ref{integraldelinea5})  y (\ref{integraldelinea6});
estos elementos corresponden al rotacional y la divergencia, respectivamente.
Estos elementos se presentaron en la sección \secref{divergenciayrotacional}.
\index{Rotacional} \index{Divergencia}

\section{Teorema de Green} \label{teoremadegreen}
En \cite{thomas_calculo_2015}, se presentan los siguientes resultados en dos formas; recordemos que ${\bf T}$ y ${\bf n}$
denotan los vectores tangente unitario y normal unitario a la curva $C$.
Una exposición completa se puede consultar en \cite{acosta_gempeler_doce_2020}.
\begin{teorema}[\bfseries Teorema de Green $-$ circulación rotacional o forma tangencial]
Sea $C$ una curva suave por tramos, cerrada y simple que encierra una región $R$ en el plano. Sea ${\bf F} (x,y)=P(x,y)\vi+Q(x,y)\vj$ un campo vectorial donde $P(x,y)$ y $Q(x,y)$ tienen primeras derivadas parciales continuas en una región abierta que contiene a $R$. Entonces, la circulación en contra de las manecillas del reloj de ${\bf F}$ alrededor de $C$ es la doble integral de $({\rm rot}{\bf F})\cdot \vk$
sobre la región $R$. \index{Región abierta} \index{Teorema!de Green}
\begin{equation}
\int_C {\bf F}\cdot{\bf T}ds=\int_C Pdx+Qdy=\iint_R \left(\frac{\partial Q}{\partial x}-\frac{\partial P}{\partial y}\right)dxdy
\end{equation}
\end{teorema}

\begin{teorema}[\bfseries Teorema de Green $-$ divergencia de flujo o forma normal]
Sea $C$ una curva suave por tramos, cerrada y simple que encierra una región R en el plano.
Sea ${\bf F} (x,y)=P(x,y)\vi+Q(x,y)\vj$ un campo vectorial donde $P(x,y)$ y $Q(x,y)$ tienen primeras derivadas parciales continuas en una región abierta que contiene a $R$. Entonces, el flujo hacia fuera de ${\bf F}$ a través de $C$ es igual a la doble integral de ${\rm div}{\bf F}$ sobre la región $R$ encerrada por $C$. \index{Región abierta}
\begin{equation}
\int_C{\bf F}\cdot{\bf n}ds=\int_C Pdy-Qdx=\iint_R \left(\frac{\partial P}{\partial x}+\frac{\partial Q}{\partial y}\right)dxdy
\end{equation}
\end{teorema}

Ilustremos lo anterior.

\begin{ejemplo}
\textls[-10]{Calcular la circulación en sentido antihorario  y el flujo hacia afuera del campo
${\bf F}(x,y)=(x^{2}y+3x)\vi+(2x^{3}-y)\vj$ y la curva $C:$ $x^{2}+2y^{2}=2$.}
\end{ejemplo}

\begin{parrafoitemize}
\item Cálculo de la circulación a través del teorema de Green con la integral doble
\[\frac{\partial Q}{\partial x}-\frac{\partial P}{\partial y}=\frac{\partial }{\partial x}\left( 2x^{3}-y\right)-
\frac{\partial }{\partial y}\left( x^{2}y+3x\right) = 5x^{2}\]
por lo tanto
\[\iint\limits_{R}\left( \frac{\partial Q}{\partial x}-\frac{\partial P}{\partial y}\right) dA=\int_{-1}^{1}\int_{-\sqrt{2-2y^{2}} }^{\sqrt{2-2y^{2}} }5x^{2}dxdy= \frac{5}{2}\sqrt{2}\pi\]
\item Cálculo de la circulación a través de la integral de línea. En este caso, la curva se parametriza como
${\bf r}(t) =\left\langle \sqrt{2}\cos t,\sin t\right\rangle $, $t\in \lbrack 0,2\pi ]$.
Por lo tanto ${\bf F}({\bf r}(t))=\left\langle \left( \sin 2t+3\sqrt{2}\right) \cos t,4\sqrt{2}\cos ^{3}t-\sin t\right\rangle$, además
\[{\bf F}({\bf r}(t))\cdot{\bf r}^{\prime}(t)=2\sqrt{2}\cos(2t)-\frac{7}{2}\sin(2t)+\frac{3}{4}\sqrt{2}\cos(4t)+\frac{5}{4}\sqrt{2}.\] Entonces la circulación está dada por
\begin{multline*} \int_{0}^{2\pi}{\bf F}\left({\bf r}(t)\right)\cdot{\bf r}^{\prime}(t) dt = \\ \int_{0}^{2\pi}\left( 2\sqrt{2}\cos 2t-\frac{7}{2}\sin 2t+\frac{3}{4}\sqrt{2}\cos 4t+\frac{5}{4}\sqrt{2}\right) dt
= \frac{5}{2}\sqrt{2}\pi
\end{multline*}
Este cálculo se efectúa en \verb"Mathematica" con los siguientes comandos:

\begin{tcolorbox}[title=\bf Instrucción \verb"Mathematica"]\begin{Verbatim}[formatcom=\letraverbatim]
r={x,y}={Sqrt[2]*Cos[t],Sin[t]};
F={x^2*y+3*x,2*x^3-y};
Integrate[Dot[F,D[r,t]],{t,0,2*Pi}]
\end{Verbatim}
\end{tcolorbox}

\item Para el flujo hacia afuera a través de la integral doble
\[\frac{\partial P}{\partial x}+\frac{\partial Q}{\partial y}=\frac{\partial }{\partial x}\left( x^{2}y+3x\right) +\frac{\partial }{\partial y}\left(2x^{3}-y\right) = 2xy+2.\]
entonces el flujo es
\[\iint\limits_{R}\left(\frac{\partial P}{\partial x}+\frac{\partial Q}{\partial y}\right)dA = \int_{-1}^{1}\int_{-\sqrt{2-2y^{2}}}^{\sqrt{2-2y^{2}} }\left( 2xy+2\right) dxdy= 2\sqrt{2}\pi.\]

\item Flujo hacia afuera a través de la integral de línea
$Pdy-Qdx=$ $\left( 10\sin t\cos ^{3}t+4\sqrt{2}\cos ^{2}t-\sqrt{2}\right) dt;$
por lo tanto,
\[\int_{C}Pdx-Qdy=\int_{0}^{2\pi }\left( 10\sin t\cos ^{3}t+4\sqrt{2}\cos ^{2}t-\sqrt{2}\right) dt= 2\sqrt{2}\pi\]
\end{parrafoitemize}

\begin{ejemplo}
\textls[-10]{Considere el campo vectorial
\[{\bf F}(x,y)=\left\langle x^{2}-\frac{x}{1+y^{2}},e^{x}+\arctan y\right\rangle\!,\]
calcule el flujo hacia afuera de ${\bf F}$ a través de la cardioide $r=a(1+\cos \theta )$, $a>0$.}
% La curva se parametriza por medio de la función ${\bf r}\left( \theta \right) =\left\langle a(1+\cos\theta )\cos \theta ,a(1+\cos \theta )\sin \theta \right\rangle\!,$ $\theta \in [0,2\pi].$
En este caso,
\[\frac{\partial P}{\partial x}+\frac{\partial Q}{\partial y}=\frac{\partial }{\partial x}\left( x^{2}-\frac{x}{1+y^{2}}\right) +\frac{\partial }{\partial y}\left( e^{x}+\arctan y\right) = 2x;\]
por lo tanto, el flujo está dado por
\[\iint\limits_{R}\left( \frac{\partial P}{\partial x}+\frac{\partial Q}{\partial y}\right) dA=\iint\limits_{R}2xdA=\int_{0}^{2\pi}\int_{0}^{a(1+\cos \theta )} 2r\cos \theta  rdrd\theta= \frac{5}{2}\pi a^{3}.\]

\end{ejemplo}
Este cálculo se efectúa en \verb"Mathematica" con los siguientes comandos:

\begin{tcolorbox}[title=\bf Instrucción \verb"Mathematica"]\begin{Verbatim}[formatcom=\letraverbatim]
{P,Q}={x^2-x/(y^2+1),Exp[x]+ArcTan[y]};
A=D[P,x]+D[Q,y];
{x,y}={r*Cos[\[Theta]],r*Sin[\[Theta]]};
Integrate[A*r,{\[Theta],0,2*Pi},
{r,0,a*(1+Cos[\[Theta]])}]
\end{Verbatim}
\end{tcolorbox}

\section{Superficies orientadas}

Según \textcite{marsden_calculo_2013}, una superficie parametrizada es una función ${\bf r}:D\subset \mathbb{R}^{2}\rightarrow \mathbb{R}^{3}$, en donde $D$ es algún dominio en $\mathbb{R}^{2};$ la superficie $\sigma $ correspondiente a la función ${\bf r}$ es su imagen $\sigma ={\bf r}\left( D\right) $, que se escribe como
\[{\bf r}\left( u,v\right) =\left\langle x\left( u,v\right) ,y\left( u,v\right),z\left( u,v\right) \right\rangle\!.\]
Si las funciones $x(u,v) $, $y(u,v) $ y $z(u,v)$ son diferenciables, se dice que ${\bf r}$ es diferenciable. Si ${\bf r}$ es diferenciable en $\left( u_{0},v_{0}\right) $, al fijar $u$, se obtiene que la función ${\bf r}\left( u_{0},v\right) $ posee como imagen una curva en $\mathbb{R}^{3}$.
Los vectores ${\bf r}_{u}$ y ${\bf r}_{v}$ son tangentes a esta curva en $\left(u_{0},v_{0}\right) ;$
de modo que ${\bf r}_{u}\times {\bf r}_{v}$ es normal a la superficie en $\left( u_{0},v_{0}\right) $.
Se dice que la superficie es \emph{suave} \index{Superficie!suave} si ${\bf r}_{u}\times {\bf r}_{v}\neq {\bf 0}$, en todo punto de $S$ (esto puede depender de la parametrización elegida).

La superficie $S$ se dice \emph{orientada} \index{Superficie!orientada} si consta de dos lados (uno exterior o positivo y el otro interior o negativo) y en todo punto, hay dos vectores normales unitarios ${\bf N}_{1}$ y ${\bf N}_{2}$, con ${\bf N}_{2}=-{\bf N}_{1}$.

Ilustraremos lo anterior con unos ejemplos. En estos casos se presentan algunas superficies con una parametrización determinada (${\bf r}(u,v)$) y unos vectores del campo vectorial ${\bf r}_{u}\times {\bf r}_{v}$
(no necesariamente normalizados), se muestran pocos de ellos para diferenciarlos entre sí.

\begin{parrafoitemize}

\item La función ${\bf r}(u,t)= \langle \cos(t)\sin(u),\sin(t)\sin(u),\cos(u) \rangle $, para $t \in [0,2\pi]$ y $u \in [0,\pi]$, parametriza la esfera unitaria, que es una superficie suave y orientada.
%

\begin{figure}[H]
\centering
\caption{Campo vectorial ${\bf r}_u \times {\bf r}_v$ sobre la esfera}
\includegraphics[width=0.8\textwidth]{Fig/174.pdf}
\fuentepropia
\end{figure}

\item En las siguientes figuras se ilustra la superficie $z=xy/4$ y el campo vectorial $\langle y/4,x/4,1 \rangle$
en algunos puntos.

\begin{figure}[H]
\centering
\caption{Campo vectorial ${\bf r}_u \times {\bf r}_v$ sobre la superficie $4z=xy$}
\includegraphics[height=4.5cm]{Fig/175.pdf}
\fuentepropia
\end{figure}

Usando la parametrización ${\bf r}(x,y)=\langle x,y,xy/4\rangle$, $x \in [-2,2]$,
$y \in [-2,2]$.

\begin{figure}[H]
\centering
\caption{Campo vectorial  ${\bf r}_u \times {\bf r}_v$}
\includegraphics[width=0.8\textwidth]{Fig/176.pdf}
\fuentepropia
\end{figure}

Es sencillo ver que ${\bf r}_x \times {\bf r}_y \neq \langle 0,0,0\rangle$, para todo $x,y$,
es decir que dicha superficie es suave, también puede probarse que esta superficie es orientada.

\item La \emph{cinta de Mobius} \index{Cinta de Mobius} es un ejemplo de una superficie no orientada.
\index{Superficie!no orientada}
Para $t \in [0,2\pi]$ y $u\in[-1,1]$, esta región se puede parametrizar en la forma
{\small
\[{\bf r}(u,t)= \langle \cos(t)(2+u\cos(t/2)/2),\sin(t)(2+u\cos(t/2)),u\sin(t/2)/2 \rangle.\]
}

\begin{figure}[H]
\centering
\caption{Cinta de Mobius}
\includegraphics[width=0.8\textwidth]{Fig/177.pdf}
\fuentepropia
\end{figure}

\item El hiperboloide de un manto con ecuación $x^2-y^2=1+z^2$ se parametriza en la forma
${\bf r}(u,v) = \langle \cosh(u)\cos(t), \cosh(u)\sin(t),\sinh(u) \rangle$, con $t \in [0,2\pi]$, $u \in \mathbb{R}$.
Puede comprobarse que esta superficie es orientada y suave. Se muestran algunos vectores del campo vectorial ${\bf r}_t \times {\bf r}_u$ que son ortogonales a la superficie. \index{Hiperboloide}

\begin{figure}[H]
\centering
\caption{Campo vectorial ${\bf r}_x \times {\bf r}_y$ sobre el hiperboloide}
\includegraphics[width=0.8\textwidth]{Fig/178.pdf}
\fuentepropia
\end{figure}

\item La ecuación $z= \frac{-1}{10}(x^2+y^2) \cos(2 \sqrt{x^2+y^2})$, representa una superficie que se parametriza por $x=u\cos(t)$, $y=u\sin(t)$ y $z=-\frac{u^2}{10}\cos(2u)$, en donde $t\in [0,2\pi]$ y
$u\in [0,\infty)$. Esta superficie está orientada y es suave para $(x,y,z) \neq (0,0,0)$.
%\clearpage

\begin{figure}[H]
\centering
\caption{Campo vectorial ${\bf r}_u \times {\bf r}_v$}
\includegraphics[width=0.8\textwidth]{Fig/179.pdf}
\fuentepropia
\end{figure}

\item \emph{La botella de Klein}, \index{Botella de Klein}: esta región es un ejemplo de superficie no orientada. \index{Superficie!no orientada}
Para $u\in [0, 2\pi]$ y $v\in [0, \pi]$. La siguiente función vectorial determina dicha región:
\begin{multline*}
\overrightarrow{r} (u,v)= \big \langle 6\cos(v)(1+\sin(v))-4(1-\cos(v)/2)\cos(u),\\
16\sin(v), 4(1-\cos(v)/2)\sin(u) \big \rangle.
\end{multline*}

\begin{figure}[H]
\centering
\caption{La botella de Klein}
\includegraphics[width=0.5\textwidth]{Fig/180.pdf}
\fuentepropia
\end{figure}

\item La porción del paraboloide con ecuación $z=1-\frac{x^2+y^2}{4} $, $z\geq 0$
se parametriza como  ${\bf r}(t,u) = \langle u\cos(t),u\sin(t),1-\frac{u^2}{4}\rangle $, $t\in [0,2\pi]$ y $u \in [0,2]$.

\begin{figure}[H]
\centering
\caption{Campo vectorial ${\bf r}_u \times {\bf r}_v$ sobre el paraboloide}
\includegraphics[width=0.8\textwidth]{Fig/181.pdf}
\fuentepropia
\end{figure}

\item La función ${\bf r}(u,t)=\langle (3+\sin(t))\cos(u),(3+\sin(t))\sin(u),\cos(t) \rangle $, con $u,t \in [0,2\pi]$,
parametriza un toro; puede probarse que esta es una superficie suave y orientada.

\begin{figure}[H]
\centering
\caption{Campo vectorial ${\bf r}_u \times {\bf r}_v$ sobre el toro}
\includegraphics[width=0.8\textwidth]{Fig/182.pdf}
\fuentepropia
\end{figure}

\end{parrafoitemize}

\section{Dos teoremas fundamentales}
Los siguientes dos resultados son importantes para la realización de los ejercicios presentados en este capítulo; el enunciado es tomado de \cite{thomas_calculo_2015}.
\begin{teorema}[\bfseries Teorema de Stokes] \index{Teorema!de Stokes}
Sea $S$ una superficie orientada suave por partes que tiene como frontera una curva suave por partes $C$. Sea ${\bf F}(x,y,z)= P(x,y,z)\vi+ Q(x,y,z)\vj + R(x,y,z)\vk$
un campo vectorial cuyos componentes tienen primeras derivadas parciales continuas sobre una región abierta que contiene a $S$. Entonces, la circulación de ${\bf F}$ alrededor de $C$ en la dirección contraria a las manecillas del reloj con respecto al vector unitario ${\bf n}$ normal a la superficie es igual a la integral de \[\int_C {\bf F} \cdot d{\bf r}=\iint_S \rot{F}\cdot {\bf n} d\sigma\]
\index{Región abierta}
\end{teorema}

La prueba se puede conseguir en \cite{apostol_calculus_1999}.
\begin{teorema}[\bfseries Teorema de la divergencia] \index{Teorema!de la divergencia}
Sea ${\bf F}(x,y,z)= P(x,y,z)\vi+ Q(x,y,z)\vj + R(x,y,z)\vk$
un campo vectorial cuyos componentes tienen primeras derivadas parciales, y sea $S$ una superficie cerrada, orientada, suave por partes. El flujo de ${\bf F}$ a través de $S$ en la dirección del campo unitario normal exterior ${\bf n}$ es igual a la integral de ${\rm div\,{\bf F}}$ sobre la región $D$ encerrada por la superficie: \index{Flujo}
\[\iint_S {\bf F}\cdot {\bf n} d\sigma = \iiint_D {\rm div\,{\bf F}} dV.\]
\end{teorema}
La prueba se puede conseguir en \cite{apostol_calculus_1999} o \cite{marsden_calculo_2013}.

\section{Ejercicios que usan el teorema de Stokes}
Ilustraremos el uso de estos resultados en algunas situaciones particulares; en ellas se comprueba el cálculo señalado de más de una forma.
\subsection{Superficie sobre un paraboloide}

Para el campo vectorial ${\bf F}(x,y,z)=\la y^3,z^3,x^3 \ra$,
evaluar {\footnotesize $\ds \iint_S {\rm rot\,}{\bf F}\cdot d{\bf S}$} sobre la superficie del paraboloide
$x=3-\left( y-2\right) ^{2}-\left(z-1\right) ^{2}$
acotada por el cilindro $\left( y-2\right) ^{2}+\left(
z-2\right) ^{2}=4$. Esta región se consideró en la página \pageref{curvas2}.

\begin{figure}[H]
\centering
\caption{Superficie sobre el paraboloide interior al cilindro}
\includegraphics[width=0.4\textwidth]{Img/70b.pdf} \hspace{5mm}
\includegraphics[width=0.4\textwidth]{Img/T_22a.pdf}

\includegraphics[width=0.4\textwidth]{Img/T_22b.pdf}
\fuentepropia
\end{figure}

La figura se obtiene con las siguientes instrucciones:

\begin{tcolorbox}[title=\bf Instrucción \verb"Mathematica"]\begin{Verbatim}[formatcom=\letraverbatim]
{y,z,x}={2+2*u*Cos[v],2+2*u*Sin[v],(y-2)^2+(z-1)^2-3};
a=ParametricPlot3D[{x,y,z},{u,0,1},{v,0,2Pi}];
b=SliceVectorPlot3D[{1,-4+2y,-2+2z},x==-3+(y-2)^2+
(z-1)^2,{x,-3,7},{y,0,4},{z,0,4.2},VectorScale->{0.08,
Scaled[0.4]},RegionFunction->Function[{x,y,z},(y-2)^2+
(z-2)^2<4],VectorStyle->Black,VectorPoints->18,
PlotStyle->Opacity[0]];
c=ParametricPlot3D[{2+4*Sin[t],2+2*Cos[t],2+2*Sin[t]},
{t,0,2*Pi},PlotStyle->{Blend[{Red,Black},1/2],
Thickness->0.01},PlotPoints->80];
Show[b,a,c]
\end{Verbatim}
\end{tcolorbox}

Con el campo vectorial ${\bf F}$ se verifica que
${\rm rot\,}{\bf F}(x,y,z)=\la -3z^2,-3 x^2,-3 y^2\ra$.
La proyección de la superficie en el plano $yz$ se parametriza mediante las expresiones
$y = 2+2u \cos (v)$, $z = 2+2u \sin (v)$. Reemplazando en la expresión que define el paraboloide se sigue que $x = -3+ (y-2)^2+(z-1)^2=-2+4u^2+4u\sin(v)$.
En resumen, para $u\in [0,1],\, v\in [0,2\pi]$, la siguiente función parametriza la superficie sobre el paraboloide con la restricción mencionada.
\[{\bf r}(u,v)=\la -2+4u^2+4u\sin(v),2+2u\cos(v),2+2u\sin(v)\ra.\]

De esta función se obtiene
${\bf r}_u \times {\bf r}_v=\la u,-16u^2\cos(v),-8u(1+2u\sin(v)) \ra$
\begin{multline*}
{\rm rot\,}{\bf F}({\bf r}(u,v)) \cdot ({\bf r}_u \times {\bf r}_v) = 48u\bigg(2(1+u\cos(v))^2(1+2u\sin(v))\\ -(1+u\sin(v))^2
+u\cos(v)(4u^2-2+4u\sin(v))^2 \bigg ).
\end{multline*}
Con lo cual,
\begin{equation} \label{stokes1}
\iint_S {\rm rot\,}{\bf F}\cdot d{\bf S}=\int_0^1\int_0^{2\pi}{\rm rot\,}{\bf F}({\bf r}(u,v))\cdot({\bf r}_u\times{\bf r}_v) dvdu = 60 \pi
\end{equation}
Este cálculo puede efectuarse de manera simplificada mediante la ejecución de las siguientes instrucciones:

\begin{tcolorbox}[title=\bf Instrucción \verb"Mathematica"]\begin{Verbatim}[formatcom=\letraverbatim]
G=Curl[{y^3,z^3,x^3},{x,y,z}];
{y,z,x}={2+2*u*Cos[v],2+2*u*Sin[v],y^2-4*y+z^2-2*z+2};
r={x,y,z};
Integrate[Dot[G,Cross[D[r,u],D[r,v]]],
{v,0,2*Pi},{u,0,1}]
\end{Verbatim}
\end{tcolorbox}

La curva de intersección del paraboloide y el cilindro se determina reduciendo las ecuaciones originales a
$x = 2z - 2$. Con las expresiones $y = 2 + 2\cos(t)$, $z = 2 + 2\sin(t)$, se sigue que la función vectorial que describe dicha curva es

{\small
\[{\bf r}(t) = \left\langle 2 + 4\sin(t),\, 2 + 2\cos(t),\, 2 + 2\sin(t) \right\rangle\!,\; t \in [0, 2\pi].\]}

Entonces,

{\small
\[
{\bf F}({\bf r}(t)) =
\left\langle
(2 + 2\cos(t))^3,\;
(2 + 2\sin(t))^3,\;
(2 + 4\sin(t))^3
\right\rangle
\]
}
y también
{\small
\[
{\bf F}({\bf r}(t)) \cdot {\bf r}'(t) =
16\cos(t)\left(2(1 + \cos(t))^3 + (1 + 2\sin(t))^3\right)
+ 16\sin(t)(1 + \sin(t))^3.
\]
}
Por lo tanto
\begin{equation} \label{stokes2}
\int_0 ^{2\pi}{\bf F}({\bf r}(t)) \cdot{\bf r}'(t) dt = 60 \pi
\end{equation}
De acuerdo con el teorema de Stokes, los resultados obtenidos en \ref{stokes1} y \ref{stokes2} son iguales.
Este valor se obtiene con la ayuda de las siguientes instrucciones:

\begin{tcolorbox}[title=\bf Instrucción \verb"Mathematica"]\begin{Verbatim}[formatcom=\letraverbatim]
r={y,z,x}={2+2*Cos[t],2+2*Sin[t],2z-2};
Integrate[Dot[{y^3,z^3,x^3},D[r,t]],{t,0,2*Pi}]
\end{Verbatim}
\end{tcolorbox}

\subsection{Superficie de un paraboloide}

Evaluar $\ds \iint {\rm rot\,}{\bf F}\cdot d{\bf S}$, si ${\bf F\,}(x,y,z)= \la y-xz,y^2,z \ra $ y $S$ es la superficie sobre el paraboloide $2z=x^2+y^2$ acotada por el plano $-x+y+z=1$.

\begin{figure}[H]
\centering
\caption{$2z\leq x^2+y^2$ \& $z \leq 1+x-y$}
\includegraphics[width=0.4\textwidth]{Img/72a.pdf} \hspace{5mm}
\includegraphics[width=0.4\textwidth]{Img/T_26a.pdf}

\includegraphics[width=0.4\textwidth]{Img/T_26b.pdf}
\fuentepropia
\end{figure}

Con las siguientes instrucciones se obtiene la figura en \verb"Mathematica":

\begin{tcolorbox}[title=\bf Instrucción \verb"Mathematica"]\begin{Verbatim}[formatcom=\letraverbatim]
a=SliceVectorPlot3D[{2*x,2*y,-2},{2*z==x^2+y^2},
{x,-1,3},{y,-3,1},{z,0,6},RegionFunction->Function[{x,
y,z},(x-1)^2+(y+1)^2<=4],VectorStyle->Black,
VectorScale->{0.1,Scaled[0.2],Automatic},
VectorPoints->15];
b=ParametricPlot3D[{1+2*Cos[t],-1+2*Sin[t],
3+2*Cos[t]-2*Sin[t]},{t,0,2Pi},PlotStyle->Red];
Show[a,b]
\end{Verbatim}
\end{tcolorbox}

Las superficies mencionadas se trabajaron en la página \pageref{curvas1}.
La curva de intersección se parametrizó por medio de la función

\[{\bf r}(t)= \la 1+ 2\cos(t), -1+ 2\sin(t), 3+2\cos(t)-2\sin(t) \ra, \, t \in [0,2 \pi].\]
Entonces, ${\bf F}({\bf r}(t))\cdot{\bf r}'(t)= 2(2\cos(t)+\cos^2(t)-\sin(t))(2\sin(t)-1)$,
por lo tanto,

\begin{equation}\label{stokes3}
\int_C {\bf F} \cdot d{\bf r}=\int_0^{2\pi}2(2\cos(t)+\cos^2(t)-\sin(t))(2\sin(t)-1) dt = -8\pi.
\end{equation}
Del campo ${\bf F}(x,y,z)$ se sigue que ${\rm rot\,}{\bf F}(x,y,z)= \la 0,-x,-1\ra$.

La región sobre el paraboloide acotada por la curva $C$, se puede parametrizar mediante las expresiones
$x=1+2u\cos(t)$,  $y =-1+2u\sin(t)$.
La expresión para $z$ se obtiene de la ecuación del paraboloide, es decir,

\[z = (x^2 + y^2)/2=1+2u^2+2u\cos(t)-2u\sin(t);\] en donde $u\in [0,1]$,
$t\in [0,2\pi]$.
Con la función

\[{\bf r}(u,t)= \la 1+2u\cos(t),-1+2u\sin(t), 1+2u^2+2u\cos(t)-2u\sin(t) \ra,\]
se sigue que

{\small
\begin{eqnarray*}
{\bf r}_u &=& \la 2\cos(t), 2\sin(t), 2(2u +\cos(t)-\sin(t))\ra \\
{\bf r}_t &=& \la -2 u \sin(t), 2u\cos(t),-2u(\cos(t) + \sin(t)) \ra\\
{\bf r}_u \times {\bf r}_t &=& \la -4u(1+2u\cos(t)),4u(1-2u\sin(t)), 4 u  \ra.
\end{eqnarray*}
}

De modo que
${\rm rot\,}{\bf F} \cdot ({\bf r}_u \times {\bf r}_t) = 8u(-1-u\cos(t)+u\sin(t)+u^2\sin(2t)) $, entonces

{\small
\begin{multline*}
\iint_S {\rm rot\,}{\bf F}\cdot d{\bf S}=
\int_0^1\int_0^{2\pi}8u(-1-u\cos(t)+u\sin(t)+u^2\sin(2t))dtdu= -8\pi.
\end{multline*}
}

Este cálculo se obtiene como sigue, y es el mismo valor obtenido en \ref{stokes3}.

\begin{tcolorbox}[title=\bf Instrucción \verb"Mathematica"]\begin{Verbatim}[formatcom=\letraverbatim]
G=Curl[{y-x*z,y^2,z};,{x,y,z}]
r={x,y,z}={1+2*u*Cos[t],-1+2*u*Sin[t],(x^2+y^2)/2};
Integrate[Dot[G,Cross[D[r,u],D[r,t]]],
{t,0,2*Pi},{u,0,1}]
\end{Verbatim}
\end{tcolorbox}

\subsection{Superficie sobre una esfera}

Sea $S$ la superficie sobre la esfera con ecuación $x^{2}+y^{2}+z^{2}=2x$, con
$x\leq \frac{3}{2}$.
%\clearpage

\begin{figure}[H]
\centering
\caption{Superficie $S$ sobre una esfera}
\includegraphics[width=0.4\textwidth]{Fig/183.pdf} \hspace{5mm}
\includegraphics[width=0.4\textwidth]{Fig/184.pdf}
\fuentepropia
\end{figure}

Consideremos ${\bf F}\left(x,y,z\right)=\left\langle -y,y-z,-x\right\rangle $. Evaluar
{\small
\[\displaystyle \iint_{S}{\rm rot}\,{\bf F}\cdot dS.\]}

La superficie $S$ es una parte de la esfera de radio $1$ con centro en el punto  $(1,0,0)$, por lo tanto, usaremos coordenadas esféricas en la forma $x=\rho \cos \phi $, $y=\rho \cos \theta \sin \phi $,
$z=\rho \sin \theta \sin \phi $ para construir una parame- \\ trización de $S$. Tendremos en cuenta que la ecuación $x^{2}+y^{2}+z^{2}=2x$ se reescribe como $\rho =2\cos \phi $,
de modo que esta superficie $S$ se parametriza como
$x=2\cos ^{2}\phi $, $y=2\cos \theta \cos \phi \sin \phi $, $z=2\cos \phi
\sin \theta \sin \phi $. O de manera más simple:

{\small
\begin{align*}
\mathbf{u}(\theta, \phi) &= \left \langle 2 \cos^2 \phi, 2 \cos \theta \cos \phi \sin \phi, 2 \cos \phi \sin \theta \sin \phi \right \rangle \\
&= \left \langle 1 + \cos (2\phi), \cos \theta \sin (2\phi), \sin \theta \sin (2\phi) \right \rangle,
\end{align*}
}
con $\theta \in \lbrack 0,2\pi ]$, $\phi \in \lbrack \pi /6,\pi /2]$.
De esta función se obtiene

{\small
\begin{align*}
\mathbf{u}_{\theta} &= \left \langle 0, -\sin 2\phi \sin \theta, \sin 2\phi \cos \theta \right \rangle \\
\mathbf{u}_{\phi} &= \left \langle -2 \sin 2\phi, 2 \cos 2\phi \cos \theta, 2 \cos 2\phi \sin \theta \right \rangle \\
\mathbf{u}_{\theta} \times \mathbf{u}_{\phi} &= \left \langle -2 \cos 2\phi \sin 2\phi, -2 \sin^2 (2\phi) \cos \theta, -2 \sin^2 (2\phi) \sin \theta \right \rangle \\
&= \left \langle -\sin (4\phi), -2 \sin^2 (2\phi) \cos \theta, -2 \sin^2 (2\phi) \sin \theta \right \rangle.
\end{align*}
}

Además,

${\rm rot}\,{\bf F}=\nabla \times \left\langle -y,y-z,-x\right\rangle =
\left\langle 1,1,1\right\rangle $
y
\begin{eqnarray*}
{\bf F}({\bf u}(\theta,\phi))\cdot ({\bf u}_{\theta }\times {\bf u}_{\phi }) &=& -\sin 4\phi -2\left( \cos \theta +\sin \theta \right) \sin^{2}\left( 2\phi \right)\!.
\end{eqnarray*}
Con lo anterior, podemos evaluar la integral

{\small
\begin{multline*}
\iint_S {\rm rot}\, {\bf F} \cdot d{\bf S} = -\int_{\pi /6}^{\pi /2}\int_{0}^{2\pi }\left( \sin 4\phi +2\left( \cos
\theta +\sin \theta \right) \sin ^{2}\left( 2\phi \right) \right) d\theta d\phi = \frac{3\pi}{4}.
\end{multline*}
}

En \Verb"Mathematica" este proceso se hace con las siguientes instrucciones:

\begin{tcolorbox}[title=\bf Instrucción \verb"Mathematica"]\begin{Verbatim}[formatcom=\letraverbatim]
r[\[Theta]_,\[Phi]_]:={1+Cos[2\[Phi]],
Cos[\[Theta]]*Sin[2\[Phi]],Sin[\[Theta]]*Sin[2\[Phi]]}
ParametricPlot3D[r[\[Theta],\[Phi]],{\[Theta],0,2*Pi},
{\[Phi],Pi/6,Pi/2}]
G=Curl[{-y,y-z,-x},{x,y,z}];
{x,y,z}={r[\[Theta],\[Phi]][[1]],r[\[Theta],
\[Phi]][[2]],r[\[Theta],\[Phi]][[3]]};
A=Simplify[Dot[G,Cross[D[r[\[Theta],\[Phi]],
\[Theta]],D[r[\[Theta],\[Phi]],\[Phi]]]]]
Integrate[A,{\[Theta],0,2*Pi},{\[Phi],Pi/6,Pi/2}]
\end{Verbatim}
\end{tcolorbox}

La figura de la mencionada superficie y el campo vectorial ${\rm rot}\, {\bf F}$ se muestran a continuación:

\begin{figure}[H]
\centering
\caption{Campo vectorial ${\bf F}$ sobre la superficie $S$}
\includegraphics[width=0.4\textwidth]{Img/T_18a.pdf} \hspace{5mm}
\includegraphics[width=0.4\textwidth]{Img/T_18b.pdf}
\fuentepropia
\end{figure}

Las figuras anteriores se obtienen ejecutando las siguientes instrucciones:

\begin{tcolorbox}[title=\bf Instrucción \verb"Mathematica"]\begin{Verbatim}[formatcom=\letraverbatim]
a=SliceVectorPlot3D[{1,1,1},{x^2+y^2+z^2==2*x},
{x,0,3/2},{y,-1,1},{z,-1,1},VectorPoints->12,
VectorStyle->Black,VectorScale->{0.12,Scaled[0.2],
Automatic},Ticks->{{0,1},{-1,0,1},{-1,0,1}},
PlotStyle->Opacity[0]];
b=ParametricPlot3D[{1+Sqrt[3]/2,Cos[\[Theta]]/2,
Sin[\[Theta]]*Sin[2*Pi/6]},{\[Theta],0,2Pi},
PlotStyle->Blue];
c=ParametricPlot3D[{1+Cos[2*\[Phi]],Cos[\[Theta]]*
Sin[2*\[Phi]],Sin[\[Theta]]*Sin[2*\[Phi]]},
{\[Theta],0,2Pi},{\[Phi],Pi/6,Pi/2},Mesh->False];
Show[a,b,c]
\end{Verbatim}
\end{tcolorbox}

La frontera de la superficie $S$ es la circunferencia $C$ determinada por la ecuación
$\vec{\alpha}(t) =\left\langle \frac{3}{2},\frac{\sqrt{3}}{2}\cos t,\frac{\sqrt{3}}{2}\sin t\right\rangle $,
$t\in \lbrack 0,2\pi ];$ entonces,
\begin{eqnarray*}
{\bf F}(\ve{\alpha}(t))\cdot d\ve{\alpha}
& =& \frac{3}{4}\left ( \sin ^{2}t-\cos t\sin t-\sqrt{3}\cos t \right );
\end{eqnarray*}
entonces
\[\int_{C}{\bf F}\cdot d\alpha=\int_{0}^{2\pi}\frac{3}{4}\left(\sin^{2}t-\cos t\sin t-\sqrt{3}\cos t\right)dt=\frac{3}{4}\pi. \]
El cálculo anterior se obtiene al ejecutar las siguientes instrucciones:

%

\begin{tcolorbox}
\begin{Verbatim}[formatcom=\letraverbatim]
{x,y,z}={3/2,Sqrt[3]/2*Cos[t],Sqrt[3]/2*Sin[t]};
F[x_,y_,z_]:={-y,y-z,-x}
Dot[F[x,y,z],D[{x,y,z},t]]
Integrate[Dot[F[x,y,z],D[{x,y,z},t]],{t,0,2*Pi}]
\end{Verbatim}
\end{tcolorbox}

\subsection{Superficie sobre un elipsoide} \index{Elipsoide}

Evaluar $\ds \iint_{S}{\rm rot}\,{\bf F}\cdot d{\bf S}$, en donde $S$ es la parte del elipsoide
$x^{2}+4y^{2}+z^{2}=16$ con $z \geq 0$ y ${\bf F}\left(x,y,z\right) =e^{x}\vi+\cos \left( yz\right)\vj+z\vk$.
%\clearpage

\begin{figure}[H]
\centering
\caption{Superficie sobre el elipsoide}
\includegraphics[width=0.8\textwidth]{Fig/185.pdf}
\fuentepropia
\end{figure}

Utilizaremos la parametrización
{\small
${\bf r}\left( u,v\right) =\left\langle 4\cos u\sin v,2\sin u\sin v,4\cos v\right\rangle $,} con
{\small $u\in \lbrack 0,2\pi ]$, $v\in \lbrack 0,\pi/2]$.}

Entonces,
\begin{eqnarray*}
{\bf r}_{u}\left( u,v\right) &=& \left\langle -4\sin u\sin v,2\cos u\sin v,0\right\rangle \\
{\bf r}_{v}\left( u,v\right) &=&\left\langle 4\cos u\cos v,2\cos v\sin u,-4\sin v\right\rangle \\
{\bf r}_{u}\times {\bf r}_{v}&=&\left\langle -8\cos u\sin ^{2}v,-16\sin u\sin ^{2}v,-8\cos v\sin v\right\rangle\!.
\end{eqnarray*}
Con los elementos anteriores es sencillo ver que ${\rm rot}\,{\bf F}= \left\langle y\sin yz,0,0\right\rangle$, además
${\rm rot}\,{\bf F}({\bf r}(u,v)) \cdot ({\bf r}_{u}\times {\bf F}_{v})= -8\sin(2u)\sin^3(v)\sin(4\sin(u)\sin(2v)), $
por lo tanto,
{\footnotesize
\begin{multline*}
\iint_{S}{\rm rot}\,{\bf F}\cdot d{\bf S}=
\int_{0}^{\pi /2}\int_{0}^{2\pi}\left( -8\sin(2u)\sin^3(v)\sin(4\sin(u)\sin(2v))\right)dudv= 0.
\end{multline*}
}

Con las siguientes instrucciones se obtienen, además de la figura, los cálculos anteriores.

\begin{tcolorbox}[title=\bf Instrucción \verb"Mathematica"]\begin{Verbatim}[formatcom=\letraverbatim]
Clear[x,y,z,F,G,r]
F[x_,y_,z_]:={Exp[x],Cos[y*z],z}
G=Curl[F[x,y,z],{x,y,z}]
r[u_,v_]:={4*Cos[u]*Sin[v],2*Sin[u]*Sin[v],4*Cos[v]}
ParametricPlot3D[r[u,v],{u,0,2*Pi},{v,0,Pi/2}]
{x,y,z}=r[u,v];
A=Simplify[Dot[G,Cross[D[r[u,v],u],D[r[u,v],v]]]]
NIntegrate[A,{u,0,2*Pi},{v,0,Pi/2}]
\end{Verbatim}
\end{tcolorbox}

La frontera de $S$ es la curva $C$ determinada por la función vectorial
$\ve{\alpha}(t) =\left\langle 4\cos t,2\sin t,0\right\rangle$, $t\in \lbrack 0,2\pi ]$.
El mismo cálculo se puede evaluar usando una integral de línea, en este caso
${\bf F}(\ve{\alpha}(t))\cdot \ve{\alpha}^{\prime}(t)=\la e^{4\cos t},\cos(0),0\ra\cdot\la -4\sin t,2\cos t,0\ra$,
por lo tanto,

{\small
\[\int_{C}{\bf F}\cdot d\ve{\alpha} = \int_{0}^{2\pi}{\bf F}(\ve{\alpha}(t))\cdot
\ve{\alpha}^{\prime }\left( t\right) dt = \int_{0}^{2\pi}\left( 2\cos t-4e^{4\cos t}\sin t\right) dt= 0.\]}

%

\begin{tcolorbox}[title=\bf Instrucción \verb"Mathematica"]\begin{Verbatim}[formatcom=\letraverbatim]
{x,y,z}={4*Cos[t],2*Sin[t],0};
F[x_,y_,z_]:={Exp[x],Cos[y*z],z}
Integrate[Dot[F[x,y,z],D[{x,y,z},t]],{t,0,2*Pi}]
\end{Verbatim}
\end{tcolorbox}

\subsection{Flujo sobre una superficie}

Calcular el flujo del campo vectorial ${\bf F}(x,y,z)= \la 2x,2y,z\ra$ sobre la superficie $S$
determinada por la ecuación $z=3-x^2-y^4+2y^2$.

\begin{figure}[H]
\centering
\caption{Campo vectorial ${\bf F}(x,y,z)=\la 2x,2y,z\ra$ sobre la superficie $S$}
\includegraphics[width=0.4\textwidth]{Img/T_4a.pdf} \hspace{5mm}
\includegraphics[width=0.4\textwidth]{Img/T_4b.pdf}
\fuentepropia
\end{figure}

Las siguientes instrucciones permiten obtener la figura que se muestra:

\begin{tcolorbox}[title=\bf Instrucción \verb"Mathematica"]\begin{Verbatim}[formatcom=\letraverbatim]
a=ContourPlot3D[3-x^2-y^4+2*y^2==z,{x,-2,2},
{y,-2,2},{z,0,4.2},Mesh->False,PlotPoints->60];
Show[a,SliceVectorPlot3D[{2*x,2*y,z},{3-x^2-y^4+
2*y^2==z},{x,-2,2},{y,-2,2},{z,0,4.4},
VectorStyle->Black,VectorScale->{0.12,Scaled[0.3],
Automatic},VectorPoints->12,PlotStyle->Opacity[0]]]
\end{Verbatim}
\end{tcolorbox}

La ecuación $z=3-x^{2}-y^{4}+2y^{2}$ determina la superficie $S$
que intersecta al plano $xy$, en una curva caracterizada por la expresión
$4=x^{2}+\left( y^{2}-1\right) ^{2}$. Esta región se muestra enseguida:

\begin{figure}[H]
\centering
\caption{$z=3-x^2-y^4+2y^2$ y su proyección en el plano $xy$}
\includegraphics[width=0.4\textwidth]{Fig/186.pdf} \hspace{5mm}
\includegraphics[width=0.4\textwidth]{Fig/187.pdf}
\fuentepropia
\end{figure}

Es claro que si $x=0$ entonces {\footnotesize $y=\pm \sqrt{3}$,} esta región se puede parametrizar formando segmentos rectilíneos entre dos puntos de la forma $A$ y $B$, que son de forma {\footnotesize $A=\left( -\sqrt{3-t^{4}+2t^{2}},t\right)$,} {\footnotesize $B=\left( \sqrt{3-t^{4}+2t^{2}},t\right)$.}
Una combinación convexa entre $A$ y $B$, permite parametrizar la proyección de la superficie sobre el plano $xy$,
esto es {\footnotesize $\left( 1-u\right)A+uB=\left(\left(2u-1\right)\sqrt{2t^{2}-t^{4}+3},t\right)$.}
La coordenada $z$ sobre la superficie, se obtiene al reemplazar $x=(2u-1)\sqrt{2t^{2}-t^{4}+3}$ y $y=t$, en
$z=3-x^{2}-y^{4}+2y^{2}$ obteniendo $z=4u(1-u)(3-t^{4}+2t^{2})$.
Por lo anterior, la superficie se describe por medio de la función
\[{\bf r}(t,u) =\la (2u-1) \sqrt{2t^{2}-t^{4}+3},t,4u(1-u)\left(3-t^{4}+2t^{2}\right)\ra,\]
con $t\in \lbrack -\sqrt{3},\sqrt{3}]$, $u\in \lbrack 0,1]$.
Con lo cual,
\[{\bf F}\left({\bf r}\left(t,u\right) \right) =\left\langle 2\left( 2u-1\right) \sqrt{2t^{2}-t^{4}+3},2t,4u\left( 1-u\right) \left( 3-t^{4}+2t^{2}\right)\right\rangle\!.\]
De la función $\bf  {r}$ obtenemos los siguientes vectores:
\begin{eqnarray*}
{\bf r}_{t}&=& \left\langle -\frac{1}{2}\left( 4t^{3}-4t\right) \frac{2u-1}{\sqrt{3-t^{4}+2t^{2}}},1,4u\left( 4t^{3}-4t\right) \left( u-1\right) \right\rangle \\
{\bf r}_{u}&=& \left\langle 2\sqrt{2t^{2}-t^{4}+3},0,4\left( t^{2}-3\right)
\left( t^{2}+1\right) \left( 2u-1\right)  \right\rangle\!.
\end{eqnarray*}
Para tener la orientación de la superficie hacia afuera se hace el producto
{\footnotesize
\[
{\bf r}_{u} \times {\bf r}_{t} =
\left\langle
4(1 - 2u)(t^{2} + 1)(t^{2} - 3),\;
\frac{-8t(t^6 - 3t^4 - t^2 + 3)}{\sqrt{3 - t^{4} + 2t^{2}}},\;
2\sqrt{3 - t^{4} + 2t^{2}}
\right\rangle\!.
\]
}

El término ${\bf F}\left({\bf r}(u,v)\right) \cdot \left({\bf r}_{u} \times {\bf r}_{v}\right)$ se puede simplificar de la siguiente mane\-ra:

{\small
\begin{multline*}
8\sqrt{2t^{2} - t^{4} + 3} \left(
3 + t^2(t^2u + t^2 - 6u) + {} \right.
\left. 3u^2(2t^{2} - t^{4} + 3) - 9u
\right),
\end{multline*}
}

de modo que el flujo de ${\bf F}$ sobre $S$ está dado por:

\[
\int_{-\sqrt{3}}^{\sqrt{3}} \int_{0}^{1}
{\\bf F}\left({\bf r}(u,v)\right) \cdot \left({\bf r}_{u} \times {\bf r}_{v}\right)
\, du\, dt \approx 136.205.
\]

Los cálculos anteriores se pueden obtener en \verb|Mathematica| con las siguientes instrucciones; con estos, además de hallar la expresión que define la proyección en el plano $xy$, se construye la parametrización mostrada antes.

\begin{tcolorbox}[title=\bf Instrucción \verb"Mathematica"]\begin{Verbatim}[formatcom=\letraverbatim]
Solve[3-x^2-y^4+2y^2==0,x]
{x1,x2,y}={-Sqrt[3+2t^2-t^4],Sqrt[3+2t^2-t^4],t};
x=(1-u)*x1+u*x2;
z=Factor[3-x^2-t^4+2t^2];
r={x,y,z};
A=Factor[Dot[{2*x,2*y,z},Cross[D[r,u],D[r,t]]]];
NIntegrate[A,{u,0,1},{t,-Sqrt[3],Sqrt[3]}]
\end{Verbatim}
\end{tcolorbox}

\subsection{Una región acotada por dos cilindros parabólicos}

Evalúe $\ds \iint_S \rot{F} \cdot d{\bf S}$ en donde ${\bf F}(x,y,z)=\left\langle y^{2},yz-x,xz-y\right\rangle $ y %% $
$S$ es la superficie acotada por las figuras de las expresiones $x=z^2$ y $x=4y-y^2$.
Esta superficie se presenta en la página \pageref{curvas3}.

En la figura~\ref{curvas3} se muestra el sólido y la representación de los campos vectoriales $\la -1,0,-2z\ra$ y  $\la 1,-4+2y,0\ra$, que determinan la orientación hacia afuera de la superficie.

\begin{figure}[H]
\centering
\caption{$x = z^2$ con $(y-2)^2+z^2 \leq 4$}
\label{curvas3}
\includegraphics[width=0.4\textwidth]{Img/53e.pdf} \hspace{5mm}
\includegraphics[width=0.4\textwidth]{Img/53h.pdf}

\includegraphics[width=0.4\textwidth]{Img/53g.pdf}
\fuentepropia
\end{figure}

Estas figuras se consiguen ejecutando las siguientes instrucciones:

\begin{tcolorbox}[title=\bf Instrucción \verb"Mathematica"]\begin{Verbatim}[formatcom=\letraverbatim]
a=ParametricPlot3D[{4*Sin[t]^2,2+2*Cos[t],2*Sin[t]},
{t,0,2*Pi},PlotStyle->{Black,Thickness[0.01]}];
b=SliceVectorPlot3D[{-1,0,2*z},{x==z^2},{x,0,4.5},
{y,0,4},{z,-2.5,2.5},RegionFunction->Function[{x,y,z},
(y-2)^2+z^2<=4],VectorScale->{0.14,Scaled[0.2],
Automatic},VectorPoints->18,VectorStyle->Black];
Show[b,a]
c=ParametricPlot3D[{4*Sin[t]^2,2+2*Cos[t],2*Sin[t]},
{t,0,2*Pi},PlotStyle->{Black,Thickness[0.01]}];
d=SliceVectorPlot3D[{1,-4+2*y,0},{x==4*y-y^2},
{x,-0.5,4},{y,-0.5,5},{z,-2,2},
RegionFunction->Function[{x,y,z},(y-2)^2+z^2<=4],
VectorScale->{0.2,Scaled[0.2],Automatic},
VectorStyle->Black,VectorPoints->20];
Show[d,c]
\end{Verbatim}
\end{tcolorbox}

Las ecuaciones $x=z^{2}$ y $x=4y-y^{2}$ representan dos cilindros parabólicos, su intersección proyectada en el plano $yz$ está dada por la ecuación $\left( y-2\right) ^{2}+z^{2}=4$, la cual se parametriza en la forma $y=2+2u\cos v$, $z=2u\sin v$, con $u\in \lbrack 0,1]$, $v\in \lbrack 0,2\pi ]$. Haremos el cálculo sobre cada región por separado.

\subsubsection{Integración sobre el cilindro \texorpdfstring{$x=z^{2}$}{x=z^2}}

Del campo vectorial se  sigue que $\rot{F}= \left\langle -y-1,-z,-2y-1\right\rangle $.
Además, la superficie en mención se parametriza por medio de la función
\[{\bf r}\left( u,v\right) =\left\langle 4u^{2}\sin ^{2}v,2+2u\cos v,2u\sin v\right\rangle\!,\, u\in \lbrack 0,1],\,v\in \lbrack 0,2\pi ].\] Por lo tanto,
$\rot{F}\left({\bf r}\left(u,v\right)\right) = \left\langle -2u\cos v-3,-2u\sin v,-4u\cos v-5\right\rangle $.
De la función ${\bf r}$ se obtienen los siguientes vectores

\begin{align*}
{\bf r}_{u} &= \left\langle 8u\sin ^{2}v,2\cos v,2\sin v\right\rangle \\
{\bf r}_{v} &=\left\langle 8u^{2}\cos v\sin v,-2u\sin v,2u\cos v\right\rangle \\
{\bf r}_{u}\times {\bf r}_{v} &=\left\langle 4u,0,-16u^{2}\sin v\right\rangle\!.
\end{align*}
Además, $\rot{F}\left({\bf r}(u,v)\right)\cdot\left({\bf r}_{u}\times{\bf r}_{v}\right)=-4 u(4u^2\cos(2v)-4u^2+4 u \cos (v)+5)$, entonces,
\begin{equation}\label{stokes6}
\int_{0}^{2\pi }\int_{0}^{1}\left( -4 u(4u^2\cos(2v)-4u^2+4 u \cos (v)+5)\right) dudv= -12\pi
\end{equation}
Las instrucciones requeridas en \verb"Mathematica" se presentan a continuación:

\begin{tcolorbox}[title=\bf Instrucción \verb"Mathematica"]\begin{Verbatim}[formatcom=\letraverbatim]
G=Curl[{y^2,y*z-x,x*z-y},{x,y,z}]
r={y,z,x}={2+2*u*Cos[v],2*u*Sin[v],z^2}
Integrate[Dot[G,Cross[D[r,u],D[r,v]]],
{u,0,1},{v,0,2Pi}]
\end{Verbatim}
\end{tcolorbox}

\subsubsection{Integración sobre el cilindro \texorpdfstring{$x=4y-y^{2}$}{x=4y-y^2}}

Sobre esta región se satisface que
$y= (2+2u\cos v)$ y $z=2u\sin v$ por lo tanto $x=4-2u^{2}\left( 1+\cos 2v\right);$
por lo tanto, la superficie se parametriza utilizando la función
\[{\bf r}\left( u,v\right) =\left\langle 4-2u^{2}\left( 1+\cos 2v\right) ,2+2u\cos v,2u\sin v\right\rangle\!,\]
con $u\in[0,1],\, v\in [0,2\pi]$.
Por lo tanto, $\rot{F}\left({\bf r}\left( u,v\right)\right) =\big \langle -2u\cos v-3,-2u\sin v,-4u\cos v-5\big \rangle$.
De la función ${\bf r}$ se sigue que
\begin{eqnarray*}
{\bf r}_{u} &=& \left\langle -4u\left( \cos 2v+1\right) ,2\cos v,2\sin v\right\rangle \\
{\bf r}_{v}&=& \left\langle 4u^{2}\sin 2v,-2u\sin v,2u\cos v\right\rangle \\
{\bf r}_{u}\times {\bf r}_{v}&=& \left\langle 4u,16u^{2}\cos v,0\right\rangle;
\end{eqnarray*}
además \[\rot{F}\left({\bf r}(u,v) \right)\cdot\left({\bf r}_{u}\times {\bf r}_{v}\right)= -4u\left( 4u^{2}\sin 2v+2u\cos v+3\right),\]
entonces,
\begin{equation}\label{stokes7}
\int_{0}^{2\pi }\int_{0}^{1}-4u\left( 4u^{2}\sin 2v+2u\cos v+3\right)dudv= -12\pi
\end{equation}
Se presentan las instrucciones que se requieren en \verb"Mathematica".

\begin{tcolorbox}[title=\bf Instrucción \verb"Mathematica"]\begin{Verbatim}[formatcom=\letraverbatim]
F={y^2,y*z-x,x*z-y};
G=Curl[{y^2,y*z-x,x*z-y},{x,y,z}];
{y,z,x}={2+2*u*Cos[v],2*u*Sin[v],4*y-y^2}
Integrate[Dot[G,Cross[D[{x,y,z},u],D[{x,y,z},v]]],
{u,0,1},{v,0,2Pi}]
\end{Verbatim}
\end{tcolorbox}

\subsubsection{Integración sobre la curva de intersección}

En la página \pageref{curvas3}, la curva de intersección se parametrizó  por medio de la función ${\bf r}(t)=\left\langle 4\sin ^{2}t,2+2\cos t,2\sin t\right\rangle$, $t \in [0,2\pi]$.
De esta expresión se sigue que
${\bf r}^{\prime}(t)  = \left\langle 8\cos t\sin t,-2\sin t,2\cos t\right\rangle$. Además,
{\small
\[{\bf F}\left({\bf r}(t) \right) =  \big \langle 4(\cos t+1) ^{2},4(\cos t+1)(\cos t+\sin t-1), 2(3\sin t-\cos t-\sin 3t-1) \big \rangle.\]
}

Por lo tanto,
{\small
\[{\bf F}\left({\bf r}\left(t\right) \right) \cdot {\bf r}^{\prime}(t) = -8\cos ^{2}\left( \frac{t}{2}\right) (1+\cos(t)-\cos(2t)-3\sin(t)-5\sin(2t)-\sin(3t));\]
}
entonces,
\begin{equation}\label{stokes8}
\int_{0}^{2\pi}{\bf F}\left({\bf r}\left(t\right) \right) \cdot {\bf r}^{\prime}(t)dt
= -12\pi
\end{equation}
De acuerdo con el teorema de Stokes, se buscará la igualdad de los planteamientos que conducen a los resultados obtenidos en \ref{stokes6}, \ref{stokes7} y \ref{stokes8}.
En \verb"Mathematica" esta integral se evalúa usando las siguientes instrucciones:

\begin{tcolorbox}[title=\bf Instrucción \verb"Mathematica"]\begin{Verbatim}[formatcom=\letraverbatim]
F={y^2,y*z-x,x*z-y};
{y,z,x}={2+2*Cos[t],2*Sin[t],z^2};
Integrate[Simplify[Dot[F,D[{x,y,z},t]]],{t,0,2*Pi}]
\end{Verbatim}
\end{tcolorbox}

\index{Flujo}
En las siguientes secciones se calcula el flujo de un campo vectorial sobre ciertas superficies.
Este valor se halla por diversos métodos.

\subsection{Una superficie de tipo exponencial}

Consideremos el campo vectorial ${\bf F}(x,y,z) = \la x-y,x+y,2z \ra$, comprobar el teorema de Stokes sobre la superficie con ecuación $z=(x+y^2)e^{1-x^2-y^2}$ para $x^2+y^2\leq 4$.
%\clearpage
%\vspace*{-1.6em}

\begin{figure}[H]
\centering
\caption{$z=(x+y^2)e^{1-x^2-y^2}$}
\includegraphics[width=0.4\textwidth]{Img/T_33a.pdf} \hspace{5mm}
\includegraphics[width=0.4\textwidth]{Img/T_33c.pdf}

\includegraphics[width=0.4\textwidth]{Img/T_33b.pdf}
\fuentepropia
\end{figure}

Con las siguientes orientaciones se consiguen estas figuras; en ellos se hace uso de la instrucción:
\verb"ListSliceVectorPlot3D". \index{ListSliceVectorPlot3D}

\begin{tcolorbox}[title=\bf Instrucción \verb"Mathematica"]\begin{Verbatim}[formatcom=\letraverbatim]
datos=Table[{x-y,x+y,2*z},{z,-2,2},{y,-2,2},{x,-2,2}];
a=ContourPlot3D[z==(x+y^2)*Exp[1-x^2-y^2],{x,-2,2},
{y,-2,2},{z,-4/3,4/3},MeshStyle->Gray,
BoxRatios->{1,1,0.75}];
b=ListSliceVectorPlot3D[datos,ImplicitRegion[z==(x+
y^2)*Exp[1-x^2-y^2],{x,y,z}],DataRange->{{-2,2},
{-2,2},{-4/3,3/2}},VectorPoints->15,VectorStyle
->{Thickness->0.005,Black,Thickness->0.004},
PlotPoints->60,PlotStyle->{Opacity[0]}];
c=ParametricPlot3D[{2*Cos[t],2*Sin[t],2*Exp[-3]
*(Cos[t]+2Sin[t]^2)},{t,0,2Pi},PlotStyle->Red];
Show[a,b,PlotRange->{{-2,2},{-2,2},{-4/3,3/2}}]
\end{Verbatim}
\end{tcolorbox}

En el plano $xy$, la región determinada por la desigualdad $x^2+y^2\leq 4$ se describe por medio de las expresiones
$x=2u\cos(t)$, $y=2u\sin(t)$, con lo cual \[z=(x+y^2)e^{1-x^2-y^2}=2ue^{1-4u^2}(\cos(t)+2u\sin^2(t)).\]
Por lo tanto, la función
\[{\bf r}(t,u) = \la 2u\cos(t),2u\sin(t) ,2ue^{1-4u^2}(\cos(t)+2u\sin^2(t)))\ra,\]
con $t\in [0,2\pi]$, $u\in [0,1]$,  parametriza la superficie en mención.
Puede verse que
\begin{multline*}
{\bf r}_u \times {\bf r}_t = \bigg \langle 4 e^{1-4 u^2} u \left(4 u^3 \cos (t)-4 u^3
\cos (3 t)+4 u^2 \cos (2 t)+4 u^2-1\right),\\
16e^{1-4 u^2} u^2 \sin (t) \left(2 u^2 \cos (2t)-2 u \cos (t)-2 u^2+1\right),4 u\bigg \rangle;
\end{multline*}
además,
$\rot{F}\cdot ({\bf r}_u \times {\bf r}_t)= 8u$, entonces
\begin{equation} \label{stokes9}
\iint_S \rot{F}\cdot {\bf n} d \sigma= \int_0^1\int_0^{2\pi} 8u dtdu =8\pi.
\end{equation}
Este cálculo se hace en \verb"Mathematica" con las siguientes instrucciones:

\begin{tcolorbox}[title=\bf Instrucción \verb"Mathematica"]\begin{Verbatim}[formatcom=\letraverbatim]
G=Curl[{x-y,x+y,2*z},{x,y,z}]
r={x,y,z}={2*u*Cos[t],2*u*Sin[t],
(x+y^2)*Exp[1-x^2-y^2]};
Integrate[Dot[G,Cross[D[r,u],D[r,t]]],{u,0,1},
{t,0,2Pi}]
\end{Verbatim}
\end{tcolorbox}

Para evaluar la integral de línea sobre la frontera de la superficie en mención, se puede usar la función
\[{\bf r}(t) = \la 2\cos(t),2\sin(t) ,2e^{-3}(\cos(t)+2\sin^2(t)))\ra,\]
con $t\in [0,2\pi]$ como parametrización de dicha curva. Es claro que
\[{\bf F}({\bf r}(t))= \la 2\cos(t)-2\sin(t),2\sin (t)+2\cos(t),4e^{-3}\left(2\sin ^2(t)+\cos(t)\right)\ra;\]
además,
${\bf F}({\bf r}(t)) \cdot {\bf r}^\prime (t)= 4 + 4e^{-6}(3\sin(2t)-\sin(t)+3\sin(3t)-2\sin(4t))$.

Por lo tanto,
\begin{equation} \label{stokes10}
\scalebox{0.85}{%
$\displaystyle \int_C {\bf F} \cdot d{\bf r} = \int_0^{2\pi}(4 + 4e^{-6}(3\sin(2t)-\sin(t)+3\sin(3t)-2\sin(4t))) dt = 8\pi$
}
\end{equation}

Es evidente la igualdad de los resultados obtenidos en \ref{stokes9} y \ref{stokes10}.
Las siguientes instrucciones permiten realizar este procedimiento:

\begin{tcolorbox}[title=\bf Instrucción \verb"Mathematica"]\begin{Verbatim}[formatcom=\letraverbatim]
{x,y}={2*Cos[t],2*Sin[t]};
{x,y,z}={2*Cos[t],2*Sin[t],(x+y^2)*Exp[1-x^2-y^2]};
Integrate[Dot[{x-y,x+y,2*z},D[{x,y,z},t]],{t,0,2Pi}]
\end{Verbatim}
\end{tcolorbox}

\subsection{Flujo sobre la superficie \texorpdfstring{$\bm{\mathsf{z=\ln(x^2+y^2)}}$}{z=ln(x²+y²)}}

Calcular el flujo del campo vectorial ${\bf F}(x,y,z)= \la yz,-xz,-xy \ra$ sobre la superficie determinada por las expresiones.
$z=\ln(x^2+y^2)$, con $\frac{1}{9} \leq x^2+y^2 \leq 9$.

\begin{figure}[H]
\centering
\caption{Campo vectorial ${\bf F}$ sobre la superficie $z=\ln (x^2+y^2)$}
\includegraphics[width=0.4\textwidth]{Img/T_5a.pdf} \hspace{5mm}
\includegraphics[width=0.4\textwidth]{Img/T_5b.pdf}
\fuentepropia
\end{figure}

La figura se consigue ejecutando las siguientes instrucciones:

\begin{tcolorbox}[title=\bf Instrucción \verb"Mathematica"]\begin{Verbatim}[formatcom=\letraverbatim]
SliceVectorPlot3D[{y*z,-x*z,-x*y},{z==Log[x^2+y^2]},
{x,-3,3},{y,-3,3},{z,-1,2},VectorPoints->14,
RegionFunction->Function[{x,y,z},0.1<x^2+y^2<9],
VectorStyle->Black,VectorScale->{0.16,Scaled[0.22],
Automatic}]
\end{Verbatim}
\end{tcolorbox}

La superficie en mención se parametriza mediante la función
\[{\bf r}\left( u,v\right) =\left\langle u\cos v,u\sin v,\ln u^{2}\right\rangle\!,\,\, v\in [0,2\pi], \,\, u\in \left[1/3,3\right ].\] Con esta función se hallan los siguientes vectores
{\small
\begin{align*}
{\bf r}_{u} &= \frac{\partial}{\partial u}\left \langle u\cos v, u\sin v, \ln u^2 \right \rangle = \left \langle \cos v, \sin v, \frac{2}{u} \right \rangle \\
{\bf r}_{v} &= \frac{\partial}{\partial v}\left \langle u\cos v, u\sin v, \ln u^2 \right \rangle = \left \langle -u\sin v, u\cos v, 0 \right \rangle \\
{\bf r}_{u} \times {\bf r}_{v} &= \left \langle \cos v, \sin v, \frac{2}{u} \right \rangle \times \left \langle -u\sin v, u\cos v, 0 \right \rangle = \left \langle -2\cos v, -2\sin v, u \right \rangle.
\end{align*}
}
Además, del campo vectorial considerado se obtiene
%${\bf F}\left( x,y,z\right) =\left\langle yz,-xz,-xy\right\rangle $
\begin{eqnarray*}
{\bf F}({\bf r}(u,v)) %&=&\left\langle (u\sin v) \ln u^{2},-(u\cos v) \ln u^{2},-(u\cos v)(u\sin v)\right\rangle \\
&=& \left\langle u\sin v\ln u^{2},-u\cos v\ln u^{2},-u^{2}\cos v\sin v \right\rangle
\end{eqnarray*}
y también ${\bf F}\left({\bf r}\left( u,v\right) \right) \cdot \left({\bf r}_{u}\times {\bf r}_{v}\right) = -u^{3}\cos v\sin v$,
por lo tanto, el flujo de ${\bf F}$ sobre la superficie en mención es
\[\int_{0}^{2\pi }\int_{1/3}^{3}-u^{3}\cos v\sin vdudv= 0.\]

\subsection{Flujo sobre la superficie \texorpdfstring{$z=xe^{1-x^2-y^2}$}{z=xe a la potencia de 1-x2-y2}}

Calcular el flujo del campo ${\bf F}(x,y,z) =\la x^{2},-y,1+z\ra $ sobre la superficie $S$ determinada por la expresión $z=xe^{1-x^2-y^2}$, $|x|\leq 2$, $|y|\leq 2$.

\begin{figure}[H]
\centering
\caption{Campo vectorial ${\bf F}$ sobre la superficie $z=xe^{1-x^2-y^2}$}
\includegraphics[width=0.4\textwidth]{Img/T_6b.pdf} \hspace{5mm}
\includegraphics[width=0.4\textwidth]{Img/T_6a.pdf}
\fuentepropia
\end{figure}

Con las siguientes instrucciones se obtiene la figura de la superficie en mención:

\begin{tcolorbox}[title=\bf Instrucción \verb"Mathematica"]\begin{Verbatim}[formatcom=\letraverbatim]
a=ContourPlot3D[z==x*Exp[1-x^2-y^2],{x,-2,2},{y,-2,2},
{z,-1.2,1.2},Mesh->False,ContourStyle->Orange];
Show[a,SliceVectorPlot3D[{x^2,-y,1+z},
{x*Exp[1-x^2-y^2]==z},{x,-2,2},{y,-2,2},
{z,-1.4,1.4},PlotStyle->{Opacity[0]},
VectorPoints->20,VectorScale->{0.15,Scaled[0.25],
Automatic},VectorStyle->Black],BoxRatios->{1,1,0.5}]
\end{Verbatim}
\end{tcolorbox}

La superficie que se considera se puede parametrizar de la siguiente mane\-ra
${\bf u}\left(x,y\right) =\la x,y,xe^{1-x^{2}-y^{2}}\ra $, con $x\in[-2,2]$, $y\in[-2,2]$.
De esta función se obtienen los siguientes vectores:
\begin{eqnarray*}
{\bf u}_{x} &=& = \la 1,0,\left( 1-2x^{2}\right) e^{-x^{2}-y^{2}+1}\ra \\
{\bf u}_{y} &=&  \la 0,1,-2xye^{-x^{2}-y^{2}+1}\ra \\
{\bf u}_{x}\times {\bf u}_{y} &=& \la e^{-x^{2}-y^{2}+1}(2x^{2}-1),2xye^{-x^{2}-y^{2}+1},1 \ra.
\end{eqnarray*}
También se puede verificar que
${\bf F}\left({\bf u}\left( x,y\right) \right) =\la x^{2},-y,1+xe^{1-x^{2}-y^{2}}\ra$
y
\[{\bf F}\left({\bf u}\left( x,y\right) \right) \cdot \left({\bf u}_{x}\times {\bf u}_{y}\right)
= x\left( 2x^{3}-x-2y^{2}+1\right)e^{-x^{2}-y^{2}+1}+1.\]
Entonces, el flujo de ${\bf F}$ a través de $S$ está dado por
{\small
\[\int_{-2}^{2}\int_{-2}^{2}\left( x\left( 2x^{3}-x-2y^{2}+1\right)e^{-x^{2}-y^{2}+1}+1\right) dydx \approx 22.706.\]
}

El cálculo anterior es posible realizarlo en \verb"Mathamatica" ejecutando las siguientes instrucciones:

\begin{tcolorbox}[title=\bf Instrucción \verb"Mathematica"]\begin{Verbatim}[formatcom=\letraverbatim]
u={x,y,x*Exp[1-x^2-y^2]};
NIntegrate[Dot[{x^2,-y,1+x*Exp[1-x^2-y^2]},
Cross[D[u,x],D[u,y]]],{x,-2,2},{y,-2,2}]
\end{Verbatim}
\end{tcolorbox}

\section{Cálculo de flujo utilizando el teorema de la divergencia}

Presentamos ejemplos en los que se utiliza el teorema de la divergencia y se compara con la integración sobre la superficie.

\subsection{Flujo sobre la superficie de un elipsoide} \index{Elipsoide}
Hallar el flujo del campo vectorial ${\bf F}(x,y,z)= \la xz,y,x \ra$ sobre la superficie del elipsoide con ecuación $x^2+4y^2+4z^2=4$. La siguiente figura ilustra esta situación.
\vskip2cm

\begin{figure}[H]
\centering
\caption{Campo vectorial ${\bf F}(x,y,z)= \la xz,y,x \ra$ sobre el elipsoide}
\includegraphics[width=0.4\textwidth]{Img/T_1a.pdf} \hspace{5mm}
\includegraphics[width=0.4\textwidth]{Img/T_1b.pdf}
\fuentepropia
\end{figure}

Para obtener esta figura, es necesario compilar las siguientes instrucciones:

\begin{tcolorbox}[title=\bf Instrucción \verb"Mathematica"]\begin{Verbatim}[formatcom=\letraverbatim]
SliceVectorPlot3D[{x*z,y,x},{x^2+4*y^2+4*z^2==4},
{x,-2,2},{y,-1,1},{z,-1,1},VectorPoints->12,
VectorStyle->Black,VectorScale->{0.15,
Scaled[0.2],Automatic},BoxRatios->{1,1,1/2},
Ticks->{{-2,0,2},{-2,0,2},{-1,0,1}}]
\end{Verbatim}
\end{tcolorbox}

La ecuación del elipsoide se reescribe en la forma $\frac{x^{2}}{4}+y^{2}+z^{2}=1$,
y se puede parametrizar multiplicando por $2$ la componente $x$ de la expresión correspondiente a las ecuaciones en coordenadas esféricas, que describen la esfera unitaria. Es por lo anterior que utilizaremos la función
\[{\bf r}\left( u,v\right) =\left\langle 2\cos u\sin v,\sin u\sin v,\cos
v\right\rangle\!,\,u\in \lbrack 0,2\pi ],\, v\in \lbrack 0,\pi ].\]
Al derivar parcialmente, obtenemos los siguientes vectores

\begin{eqnarray*}
{\bf r}_{u} &=& \left\langle -2\sin u\sin v,\cos u\sin v,0\right\rangle  \\
{\bf r}_{v} &=& \left\langle 2\cos u\cos v,\cos v\sin u,-\sin v\right\rangle \\
{\bf r}_{u}\times {\bf r}_{v}&=&  \left\langle \cos u\sin ^{2}v,2\sin u\sin^{2}v,2\sin v\cos v\right\rangle\!.
\end{eqnarray*}

Y con ellos obtenemos las siguientes expresiones:
{\small
\begin{align*}
{\bf F}\left( {\bf r}(u,v) \right) &=
\left\langle
2\cos u \sin v \cos v,\;
\sin u \sin v,\;
2\cos u \sin v
\right\rangle\!, \\
{\bf F}({\bf r}(u,v)) \cdot ({\bf r}_{u} \times {\bf r}_{v}) &=
2\left(
\sin^{2} u \sin v + (\cos u \cos v)(2 + \cos u \sin v)
\right) \sin^{2} v.
\end{align*}
}

Entonces, el flujo de ${\bf F}$ a través de $S$ está dado por:
\[
\int_{0}^{\pi} \int_{0}^{2\pi}
{\bf F}\left( {\bf r}(u,v) \right) \cdot \left({\bf r}_{u} \times {\bf r}_{v}\right)
\, du\, dv = \frac{8}{3}\pi.
\]

Los cálculos anteriores se pueden hacer en \verb|Mathematica| ejecutando las siguientes sentencias:

\begin{tcolorbox}[title=\bf Instrucción \verb"Mathematica"]\begin{Verbatim}[formatcom=\letraverbatim]
r={x,y,z}={2*Cos[u]*Sin[v],Sin[u]*Sin[v],Cos[v]};
A=Simplify[Dot[{x*z,y,x},Cross[D[r,v],D[r,u]]]]
Integrate[A,{u,0,2*Pi},{v,0,Pi}]
\end{Verbatim}
\end{tcolorbox}

Otra posibilidad de hacer este cálculo es mediante el uso del teorema de la divergencia.
Para este propósito se necesita calcular la divergencia del campo vectorial ${\bf F}$, esto es
\[{\rm div}{\bf F}=\nabla \cdot \langle xz,y,x\rangle = z+1.\]
En este caso, el flujo está dado por
{\small
\[\bigintss_{-2}^{2}\bigintss_{-\sqrt{1-\frac{x^{2}}{4}}}^{\sqrt{1-\frac{x^{2}}{4}}}\bigintss_{-\sqrt{1-\frac{x^{2}}{4}-y^{2}}}^{\sqrt{1-\frac{x^{2}}{4}-y^{2}}}\left( z+1\right)dzdydx= \frac{8}{3}\pi.\]
}

Este cálculo se evaluó en \verb"Mathematica" con la siguiente instrucción:

\begin{tcolorbox}[title=\bf Instrucción \verb"Mathematica"]\begin{Verbatim}[formatcom=\letraverbatim]
Integrate[z+1,{x,-2,2},{y,-Sqrt[1-x^2/4],
Sqrt[1-x^2/4]},{z,-Sqrt[1-x^2/4-y^2],
Sqrt[1-x^2/4-y^2]}]
\end{Verbatim}
\end{tcolorbox}

\subsection{Flujo sobre la superficie de un paraboloide}

Calcular el flujo del campo vectorial ${\bf F}\left( x,y,z\right) =\left\langle xz,x(y+1),z(y+1)\right\rangle$,
sobre la superficie del paraboloide $z=4-x^{2}-y^{2}$ y el disco sobre el plano $z=0$, que resulta de proyectar esta región sobre dicho plano.

\begin{figure}[H]
\centering
\caption{Campo vectorial ${\bf F}(x,y,z)= \la xz,x(y+1),z(y+1) \ra$ sobre el paraboloide}
\includegraphics[width=0.4\textwidth]{Img/T_2a.pdf} \hspace{5mm}
\includegraphics[width=0.4\textwidth]{Img/T_2b.pdf}
\fuentepropia
\end{figure}

La figura anterior se consigue con la ayuda de las siguientes instrucciones:

\begin{tcolorbox}[title=\bf Instrucción \verb"Mathematica"]\begin{Verbatim}[formatcom=\letraverbatim]
SliceVectorPlot3D[{x*z,x*(y+1),z*(y+1)},{4-x^2-
y^2==z},{x,-2,2},{y,-2,2},{z,0,4},VectorPoints->12,
VectorStyle->Black,VectorScale->{0.16,Scaled[0.22],
Automatic},Ticks->{{-2,0,2},{-2,0,2},{0,2,4}}]
\end{Verbatim}
\end{tcolorbox}

La proyección en el plano $xy$ del paraboloide corresponde a un círculo centrado en el origen del radio $2$.
Para calcular el flujo sobre el disco que forma la base del sólido, se parametriza esta superficie en la forma
\[{\bf r}\left(u,v\right) = \left\langle u\cos v,u\sin v,0\right\rangle ,\, u\in \lbrack 0,2],\,v\in \lbrack 0,2\pi ].\]
Es sencillo ver que
${\bf F}\left({\bf r}\left(u,v\right)\right)=\left\langle 0,u\left(u\sin v+1\right)\cos v,0\right\rangle;$
con la función ${\bf r}$ se obtienen los vectores:
\[{\bf r}_{u} = \left\langle \cos v, \sin v, 0 \right\rangle\] y
\[{\bf r}_{v} = \left\langle -u\sin v, u\cos v, 0 \right\rangle\!.\]
Para conseguir que la orientación sea hacia afuera, se evalúa:
\[
{\bf r}_{v} \times {\bf r}_{u} =
\left\langle -u\sin v, u\cos v, 0 \right\rangle \times
\left\langle \cos v, \sin v, 0 \right\rangle =
\left\langle 0, 0, -u \right\rangle\!.
\]

Es claro que
\[
{\\ F}\left({\bf r}(u,v)\right) \cdot \left({\bf r}_{u} \times {\bf r}_{v}\right) = 0;
\]
por lo tanto, el flujo sobre la superficie en mención es
\begin{equation} \label{stokes4}
\int_{0}^{2\pi }\int_{0}^{2}{\bf F}\left({\bf r}\left(u,v\right)\right) \cdot \left({\bf r}_{u}\times {\bf r}_{v}\right) dudv =0
\end{equation}

La superficie sobre el paraboloide se puede parametrizar en la forma
${\bf r}\left(u,v\right) =\left\langle u\cos v,u\sin v,4-u^{2}\right\rangle\!,\, u\in \lbrack 0,2],\, v\in \lbrack 0,2\pi ]$.
Con esta función hallamos los siguientes vectores:
\begin{eqnarray*}
{\bf r}_{u}&=& \left\langle \cos v,\sin v,-2u\right\rangle \\
{\bf r}_{v} &=& \left\langle -u\sin v,u\cos v,0\right\rangle \\
{\bf r}_{u}\times {\bf r}_{v} &=& \left\langle 2u^{2}\cos v,2u^{2}\sin v,u\right\rangle\!.
\end{eqnarray*}
Sustituyéndose adecuadamente, el término ${\bf F}\left({\bf r}\left(u,v\right)\right) \cdot \left({\bf r}_{u}\times {\bf r}_{v}\right)$ se simplifica como
$\big(4u-u^{3}+u^{2}((4-u^{2}) \sin v+(8u-2u^{3})\cos^{2}v +u(u\sin v+1) \sin (2v))\big)$.
Integrando esta expresión obtendremos el flujo de ${\bf F}$ a través de $S$,
\begin{equation} \label{stokes5}
\int_{0}^{2\pi }\int_{0}^{2}{\bf F}\left({\bf r}\left(u,v\right)\right) \cdot \left({\bf r}_{u}\times {\bf r}_{v}\right) dudv= \frac{56}{3}\pi
\end{equation}
Los procedimientos que conducen a la integral (\ref{stokes5}) se logran con la ejecución de las siguientes instrucciones:

\begin{tcolorbox}[title=\bf Instrucción \verb"Mathematica"]\begin{Verbatim}[formatcom=\letraverbatim]
r={x,y,z}={u*Cos[v],u*Sin[v],4-u^2};
Integrate[Dot[{x*z,x*(y+1),z*(y+1)},Cross[D[r,v],
D[r,u]]],{u,0,2*Pi},{v,0,Pi}]
\end{Verbatim}
\end{tcolorbox}

La otra alternativa para calcular el flujo es usar el teorema de la divergencia, como ${\rm div}{\bf F}=\nabla \cdot \left\langle xz,x\left( y+1\right) ,z\left( y+1\right) \right\rangle = x+y+z+1$.
Entonces, el flujo de ${\bf F}$ sobre la superficie de $S$ está dado por
\[\int_{-2}^{2}\int_{-\sqrt{4-x^{2}}}^{\sqrt{4-x^{2}}}\int_{0}^{4-x^{2}-y^{2}}\left( x+y+z+1\right) dzdydx= \frac{56}{3}\pi.\]
Este valor es el producto de sumar los resultados obtenidos en \ref{stokes4} y \ref{stokes5}.
Usando \verb"Mathematica" las siguientes sentencias realizan este cálculo.

\begin{tcolorbox}[title=\bf Instrucción \verb"Mathematica"]\begin{Verbatim}[formatcom=\letraverbatim]
F={x*z,x*(y+1),z*(y+1)};
Integrate[Div[F,{x,y,z}],{x,-2,2},{y,-Sqrt[4-x^2],
Sqrt[4-x^2]},{z,0,4-x^2-y^2}]
\end{Verbatim}
\end{tcolorbox}

\subsection{Flujo sobre la superficie de un paraboloide y un plano}

Consideremos el sólido $E$ determinado por el paraboloide con ecuación $9y=3-9x^{2}-4z^{2}+4z$ y el plano $4z-9y=9$. Calcular el flujo neto del campo vectorial
{\small ${\bf F}\left( x,y,z\right) =\la xz,x(y+1),z(y+1) \ra$,} a través de la superficie de $E$. La proyección de esta región en el plano $xz$ está dada por la ecuación  $9x^{2}+4z^{2}=12$. Esta región se puede parametrizar en la forma
$x=\frac{2}{\sqrt{3}}u\cos v$, $z=\sqrt{3}u\sin v$, en donde  $u \in [0,1]$, $v \in [0,2\pi]$.
Las siguientes figuras muestran el sólido.
%\clearpage
%

\begin{figure}[H]
\centering
\caption{Superficies $9y=3-9x^{2}-4z^{2}+4z$ y $4z-9y=9$}
\includegraphics[width=0.4\textwidth]{Img/T_T1.pdf} \hspace{5mm}
\includegraphics[width=0.4\textwidth]{Img/T_T2.pdf}
\fuentepropia
\end{figure}

Las siguientes son las instrucciones para obtener las figuras.

\begin{tcolorbox}[title=\bf Instrucción \verb"Mathematica"]\begin{Verbatim}[formatcom=\letraverbatim]
ContourPlot3D[{9*y==3-9*x^2-4*z^2+4*z,4z-9y==9},
{x,-2,2},{y,-2,1},{z,-2.5,3},Mesh->False,
ContourStyle->{Cyan,Orange}]
{x,z}={2/Sqrt[3]*u*Cos[t],Sqrt[3]*u*Sin[t]};
a=ParametricPlot3D[{x,(3-9*x^2-4*z^2+4*z)/9,z},
{u,0,1},{t,0,2Pi},PlotStyle->Orange,PlotPoints->80];
b=ParametricPlot3D[{x,(4*z-9)/9,z},{u,0,1},{t,0,2Pi},
PlotStyle->Cyan,PlotPoints->80];
Show[a,b,BoxRatios->{1,1,1}]
\end{Verbatim}
\end{tcolorbox}

La siguiente es una representación del campo vectorial que actúa sobre $E$.

\begin{figure}[H]
\centering
\caption{Campo vectorial ${\bf F}(x,y,z)= \la xz,x(y+1),z(y+1) \ra$ sobre la superficie $E$}
\includegraphics[width=0.4\textwidth]{Img/T_21a.pdf} \hspace{5mm}
\includegraphics[width=0.4\textwidth]{Img/T_21b.pdf}

\includegraphics[width=0.4\textwidth]{Img/T_21c.pdf}
\fuentepropia
\end{figure}

Esta figura se consigue con la ayuda de las siguientes instrucciones:

\begin{tcolorbox}[title=\bf Instrucción \verb"Mathematica"]\begin{Verbatim}[formatcom=\letraverbatim]
a=RegionPlot3D[(4*z-9)/9<=y<=(3-9*x^2-4*z^2+4*z)/9
&&9*x^2+4*z^2<=12,{x,-2,2},{y,-2,1},{z,-2,2},
PlotPoints->90,Mesh->False];
b=SliceVectorPlot3D[{x*z,x*(y+1),z(y+1)},
{4*z-9*y==9,9*y==3-9*x^2-4*z^2+4*z},{x,-2,2},
{y,-2,1},{z,-2,2},RegionFunction->Function[{x,y,z},
9*x^2+4*z^2<12],VectorPoints->20,VectorStyle->Black,
PlotPoints->60,VectorScale->{0.3,Scaled[0.4],
Automatic},PlotStyle->Opacity[0]]
Show[a,b]
\end{Verbatim}
\end{tcolorbox}

\subsubsection{Integral sobre la porción del plano}

Sobre el plano la coordenada $y$ está dada por
$y=\frac{4}{9}z-1=\frac{4}{9}\sqrt{3}u\sin v-1$,
de modo que una parametrización de la región sobre la porción del plano acotada por el paraboloide es la siguiente:

${\bf r}\left( u,v\right) =\left\langle \frac{2}{\sqrt{3}}u\cos v,\frac{4}{9}\sqrt{3}u\sin v-1,\sqrt{3}u\sin v\right\rangle$,
$u\in \lbrack 0,1]$, $v\in\lbrack 0,2\pi ];$

con esta función hallamos los siguientes vectores:

{\small
\begin{eqnarray*}
{\bf r}_{u} &=&\left\langle \frac{2}{3}\sqrt{3}\cos v,\frac{4}{9}\sqrt{3}\sin v,\sqrt{3}\sin v\right\rangle \\
{\bf r}_{v}&=&\left\langle -\frac{2}{3}\sqrt{3}u\sin v,\frac{4}{9}\sqrt{3}u\cos v,\sqrt{3}u\cos v\right\rangle \\
{\bf r}_{u}\times {\bf r}_{v}&=& % \left\langle \frac{2}{3}\sqrt{3}\cos v,\frac{4}{9}\sqrt{3}\sin v,\sqrt{3}\sin v\right\rangle \times \left\langle -\frac{2}{3}\sqrt{3}u\sin v,\frac{4}{9}\sqrt{3}u\cos v,\sqrt{3}u\cos v\right\rangle \\
\left\langle 0,-2u,\frac{8}{9}u\right\rangle\!.
\end{eqnarray*}
}

También pueden comprobarse las siguientes igualdades,
{\small
\begin{eqnarray*}
{\bf F}\left( {\bf r}\left( u,v\right) \right) &=& \left\langle 2u^{2}\cos v\sin v,\frac{8}{9}u^{2}\cos v\sin v,\frac{4}{3}u^{2}\sin ^{2}v \right\rangle  \\
{\bf F}\left( {\bf r}\left( u,v\right) \right) \cdot \left( {\bf r}_{u}\times {\bf r}_{v}\right)
&=& %\left\langle 2u^{2}\cos v\sin v,\frac{8}{9}u^{2}\cos v\sin v,\frac{4}{3} u^{2}\sin ^{2}v \right\rangle\cdot \left\langle 0,-2u,\frac{8}{9}u\right\rangle =
\frac{32}{27}u^{3}\sin ^{2}v-\frac{16}{9}u^{3}\sin v\cos v
\end{eqnarray*}
}

De tal modo que el flujo sobre la porción del plano es
\begin{equation} \label{T1}
\scalebox{0.88}{%
$\displaystyle \iint_{S_1} {\bf F} \cdot d\boldsymbol{\sigma} = \int_{0}^{2\pi }\int_{0}^{1}\left(\frac{32}{27}u^{3}\sin ^{2}v-\frac{16}{9}u^{3}\sin v\cos v \right) dudv = \frac{8}{27}\pi$
}
\end{equation}

El procedimiento mostrado antes se puede efectuar en \verb"Mathematica"
ejecutando las siguientes sentencias:

\begin{tcolorbox}[title=\bf Instrucción \verb"Mathematica"]\begin{Verbatim}[formatcom=\letraverbatim]
{x,y}={2/Sqrt[3]*u*Cos[v],4/9*Sqrt[3]*u*Sin[v]-1}
z=Sqrt[3]*u*Sin[v];
F={x*z,x*(y+1),z*(y+1)};
A=Dot[F,Simplify[Cross[D[{x,y,z},u],D[{x,y,z},v]]]];
Integrate[A,{u,0,1},{v,0,2Pi}]
\end{Verbatim}
\end{tcolorbox}

\subsubsection{Integral sobre la porción del paraboloide}

Con los elementos anteriores, la coordenada $y$ sobre el paraboloide está dada por
$y = \frac{1}{9}\left( 3-9x^{2}-4z^{2}+4z\right)= \frac{4}{9}\sqrt{3}u\sin v+\frac{1}{3}-\frac{4}{3}u^{2}$,
de modo que una parametrización de la región sobre el paraboloide acotada por la intersección con el plano, es la siguiente

{\footnotesize
\[{\bf r}\left( u,v\right) =\left\langle \frac{2}{\sqrt{3}}u\cos v,\frac{4}{9}\sqrt{3}u\sin v+\frac{1}{3}-\frac{4}{3}u^{2},\sqrt{3}u\sin v\right\rangle\!,u\in \lbrack 0,1],\, v\in \lbrack 0,2\pi ].\]
}

Esta función nos permite hallar los siguientes vectores:
\begin{eqnarray*}
{\bf r}_{u} &=& \left\langle \frac{2}{3}\sqrt{3}\cos v,\frac{4}{9}\sqrt{3}\sin v-\frac{8}{3}u,\sqrt{3}\sin v\right\rangle \\
{\bf r}_{v} &=& \left\langle -\frac{2}{3}\sqrt{3}u\sin v,\frac{4}{9}\sqrt{3}u\cos v,\sqrt{3}u\cos v\right\rangle \\
{\bf r}_{v}\times {\bf r}_{u} %&=&\left\langle \frac{2}{3}\sqrt{3}\cos v,\frac{4}{9}\sqrt{3} \sin v-\frac{8}{3}u,\sqrt{3}\sin v\right\rangle \times \left\langle -\frac{2}{3}\sqrt{3}u\sin v,\frac{4}{9}\sqrt{3}u\cos v,\sqrt{3}u\cos v\right\rangle \\
&=&  \left\langle \frac{8}{3}\sqrt{3}u^{2}\cos v,2u,\frac{8}{9}u+\frac{16}{9}\sqrt{3}u^{2}\sin v\right\rangle\!.
\end{eqnarray*}
Además,
\begin{multline*}
{\bf F}\left( {\bf r}\left( u,v\right) \right) = \bigg\langle 2u^{2}\sin v\cos v,
\frac{2}{3}\sqrt{3}u\left( \frac{4}{9}\sqrt{3}u\sin v+\frac{4}{3}- \frac{4}{3}u^{2}\right) \cos v, \\
\sqrt{3}u\left( \frac{4}{9}\sqrt{3}u\sin v+\frac{4}{3}-\frac{4}{3}u^{2}\right) \sin v\bigg\rangle
\end{multline*}

{\small
\begin{multline*}
{\bf F}\left({\bf r}(u,v) \right) \cdot \left( {\bf r}_{v}\times {\bf r}_{u}\right) =
\frac{16}{3}\sqrt{3}u^{4}\sin v\cos ^{2}v+ \\
\left( \frac{16}{27}\sqrt{3}u^{2}\cos v - \frac{4}{9}\sqrt{3}u\sin v \left( \frac{8}{9}u-\frac{16}{9}\sqrt{3}u^{2}\sin v\right) \right) \left( \sqrt{3}u\sin v+3-3u^{2}\right)\!.
\end{multline*}
}

De lo anterior se sigue que
\begin{equation} \label{T2}
\iint_{S_2}{\bf F}\left( {\bf r}\left( u,v\right) \right) \cdot \left( {\bf r}_{u}\times {\bf r}_{v}\right)dudv=\frac{8}{27}\pi
\end{equation}
El procedimiento anterior se puede hacer en \verb"Mathematica" con ayuda de las siguientes instrucciones:

\begin{tcolorbox}[title=\bf Instrucción \verb"Mathematica"]\begin{Verbatim}[formatcom=\letraverbatim]
{x,z}={2/Sqrt[3]*u*Cos[v],Sqrt[3]*u*Sin[v]};
y=Simplify[(3-9x^2-4*z^2+4z)/9]
{r,F}={{x,y,z},{x*z,x*(y+1),z*(y+1)}};
A=Dot[F,Simplify[Cross[D[r,v],D[r,u]]]];
Integrate[A,{u,0,1},{v,0,2Pi}]
\end{Verbatim}
\end{tcolorbox}

\subsubsection{Usando el teorema de la divergencia}

Es claro que ${\rm div}{\bf F}=z+x+y+1$, por lo tanto, el flujo de ${\bf F}$ a través de $S$ está dado por
\[\bigintsss_{-\frac{2}{\sqrt{3}}}^{\frac{2}{\sqrt{3}}}\bigintsss_{-\frac{1}{2}\sqrt{12-9x^{2}}}^{\frac{1}{2}\sqrt{12-9x^{2}}}\bigintsss_{\frac{4}{9}z-1}^{\frac{4}{9}z+ \frac{1}{3}-x^{2}-\frac{4}{9}z^{2}}\left( z+x+y+1\right) dydzdx=\frac{16}{27}\pi,\]
que coincide con la suma de las expresiones \ref{T1} y \ref{T2}.
Este cálculo se consigue ejecutando las siguientes instrucciones:

\begin{tcolorbox}[title=\bf Instrucción \verb"Mathematica"]\begin{Verbatim}[formatcom=\letraverbatim]
F={x*z,x*(y+1),z*(y+1)}
Integrate[Div[F,{x,y,z}],{x,-2/Sqrt[3],2/Sqrt[3]},
{z,-Sqrt[12-9*x^2]/2,Sqrt[12-9*x^2]/2},{y,4z/9-1,
(3+4z-9*x^2-4z^2)/9}]
\end{Verbatim}
\end{tcolorbox}

\subsection{Flujo sobre una esfera}

Se calcula el flujo del campo vectorial ${\bf F}(x,y,z)= \la y^2,-xy,z \ra$ sobre la esfera con ecuación $x^2+y^2+z^2=4$.
El gráfico se obtiene con las siguientes instrucciones:

\begin{tcolorbox}[title=\bf Instrucción \verb"Mathematica"]\begin{Verbatim}[formatcom=\letraverbatim]
SliceVectorPlot3D[{y^2,-x*y,z},{4==x^2+y^2+z^2},
{x,-2,2},{y,-2,2},{z,-2,2},VectorPoints->10,
VectorPoints->12,VectorStyle->Black,
VectorScale->{0.2,Scaled[0.2],Automatic}]
\end{Verbatim}
\end{tcolorbox}

\begin{figure}[H]
\centering
\caption{Campo vectorial ${\bf F}(x,y,z)= \la y^2,-xy,z \ra$ sobre una esfera}
\includegraphics[width=0.4\textwidth]{Img/T_3a.pdf} \hspace{5mm}
\includegraphics[width=0.4\textwidth]{Img/T_3b.pdf}
\fuentepropia
\end{figure}

La esfera puede describirse utilizando la función
\[{\bf r}\left( u,v\right)=\left\langle 2\cos u\sin v,2\sin u\sin v,2\cos v\right\rangle ,\, u\in
\lbrack 0,2\pi ],\,v\in \lbrack 0,2\pi ].\]
Con esta función hallamos los siguientes vectores:
\begin{eqnarray*}
{\bf r}_{u}&=&\left\langle -2\sin u\sin v,2\cos u\sin v,0\right\rangle \\
{\bf r}_{v}&=& \left\langle 2\cos u\cos v,2\cos v\sin u,-2\sin v\right\rangle \\
{\bf r}_{v}\times {\bf r}_{u}&=&\la 4\cos u\sin ^{2}v,4\sin u\sin ^{2}v,4\cos v\sin v \ra.
\end{eqnarray*}
Con estos elementos se puede establecer la siguiente relación:
\[
\mathbf{F}\left( \mathbf{r}\left( u,v\right) \right) = \left\langle 4\sin^{2}u\sin^{2}v,\ -4\cos u\sin u\sin^{2}v,\ 2\cos v \right\rangle
\]

Entonces,
\[
\mathbf{F}\left( \mathbf{r}\left( u,v\right) \right) \cdot \left( \mathbf{r}_{u} \times \mathbf{r}_{v} \right) = -8\cos^{2}v\sin v.
\]

Por lo tanto, el flujo de \(\mathbf{F}\) a través de la superficie de la esfera está dado por:
\[
\int_{0}^{\pi} \int_{0}^{2\pi} 8\cos^{2}v\sin v\ dudv = \frac{32}{3}\pi.
\]

El proceso anterior se puede hacer en \verb"Mathematica" utilizando las siguien\-tes instrucciones:

\begin{tcolorbox}[title=\bf Instrucción \verb"Mathematica"]\begin{Verbatim}[formatcom=\letraverbatim]
F={y^2,-x*y,z};
r={x,y,z}={2*Cos[u]*Sin[v],2*Sin[u]*Sin[v],2*Cos[v]};
Integrate[Dot[Cross[D[r,v],D[r,u]],F],{u,0,2*Pi},
{v,0,Pi}]
\end{Verbatim}
\end{tcolorbox}

El flujo también se puede hallar usando el teorema de la divergencia; en ese caso, el cálculo es el siguiente
{\small
\[\iiint_E {\rm div} \,{\bf F} dV= \int _0^{\pi}\int _0^{2\pi}\int _0^{2} (1-\rho \sin (\phi))\rho^2\sin(\phi)d\rho d\theta d\phi = \frac{32}{3} \pi.\]}
Esta integral se evaluó en \verb"Mathematica" usando la siguiente instrucción:

\begin{tcolorbox}[title=\bf Instrucción \verb"Mathematica"]\begin{Verbatim}[formatcom=\letraverbatim]
Integrate[(1-\[Rho]*Sin[\[Phi]]*Cos[\[Theta]])*
\[Rho]^2*Sin[\[Phi]],{\[Rho],0,2},{\[Phi],0,Pi},
{\[Theta],0,2*Pi}]
\end{Verbatim}
\end{tcolorbox}

\section{Cálculo de flujo sobre superficies cerradas}

En esta sección se presentan algunos ejemplos en que se calcula el flujo sobre superficies cerradas y se compara el valor obtenido usando el teorema de la divergencia, cuando es posible hacerlo.
\begin{parrafonumerado}
\item Hallar el flujo neto del campo vectorial.

${\bf F}(x,y,z) = \la y^2,x^2,z \ra $ sobre la superficie cerrada del sólido acotado por los paraboloides con ecuaciones:

%$z+3=x^2+y^2$ y $2z=6-2x^2-y^2.$

\begin{gather*}
z + 3 = x^2 + y^2 \quad \text{y} \quad 2z = 6 - 2x^2 - y^2
\end{gather*}

\begin{figure}[H]
\centering
\caption{Campo vectorial ${\bf F}(x,y,z)=\la y^2,x^2,z\ra $ sobre una superficie}
\includegraphics[width=0.4\textwidth]{Img/T_8b.pdf} \hspace{5mm}
\includegraphics[width=0.4\textwidth]{Img/T_8d.pdf}
\fuentepropia
\end{figure}

Esta representación del campo vectorial se obtiene ejecutando las siguientes instrucciones:

\begin{tcolorbox}[title=\bf Instrucción \verb"Mathematica"]\begin{Verbatim}[formatcom=\letraverbatim]
a=RegionPlot3D[z<(6-2*x^2-y^2)/2&&z>x^2+y^2-3,
{x,-2,2},{y,-2,2},{z,-3,3},PlotPoints->90,
Mesh->False];
b=SliceVectorPlot3D[{y^2,x^2,z},{2*z==6-2*x^2-y^2,
z==x^2+y^2-3},{x,-2,2},{y,-2,2},{z,-3,3},
VectorPoints->20,VectorScale->{0.15,Scaled[0.25],
Automatic},RegionFunction->Function[{x,y,z},4*x^2
+3*y^2<=12],VectorStyle->Black,PlotStyle->Opacity[0]];
Show[a,b]
\end{Verbatim}
\end{tcolorbox}

De acuerdo con el ejercicio \ref{ejercicio1} de la página \pageref{ejercicio1}, la curva de intersección corresponde a la ecuación $4x^{2}+3y^{2}=12$, que representa una elipse.
Esta región y su interior se pueden representar en el plano $xy$ por medio de las ecuaciones
$x=\sqrt{3}u\cos v$, $y=2u\sin v$, con $u\in \lbrack 0,1]$,
$v\in \lbrack 0,2\pi ]$.
Calcularemos el flujo sobre las dos superficies que limitan el sólido.

\subsubsection{Flujo sobre el paraboloide \texorpdfstring{$\bm{\mathsf{2z=6-2x^2-y^2}}$}{2z=6-2x^2-y^2}}

En este caso, una función que parametriza esta región es:
{\small
\[
{\\\alpha}(u,v) =
\left\langle
\sqrt{3}u\cos v,\;
2u\sin v,\;
3 - \frac{1}{2}u^{2}\left(5 + \cos 2v\right)
\right\rangle\!,
u \in [0,1],\; v \in [0, 2\pi].
\]
}

Con esta función se obtienen los siguientes vectores:
\begin{eqnarray*}
{\bm\alpha}_{u}&=&\left\langle \sqrt{3}\cos v,2\sin v,-u\left( \cos 2v+5\right)
\right\rangle \\
{\bm\alpha}_{v}&=&\left\langle -\sqrt{3}u\sin v,2u\cos v,u^{2}\sin 2v\right\rangle \\
{\bm\alpha}_{u}\times {\bm\alpha}_{v}&=&\left\langle 12u^{2}\cos v,4\sqrt{3}u^{2}\sin v,2\sqrt{3}u\right\rangle
\end{eqnarray*}
${\bf F}\left({\bm \alpha}(u,v)\right)=\left\langle 4u^2\sin^2 v,3u^2\cos^2 v,3-\frac{1}{2}u^{2}\left(5+\cos 2v\right)\right\rangle $, entonces \\
${\bf F}\left({\bm\alpha}(u,v)\right) \cdot \left({\bm\alpha}_{u}\times {\bm\alpha} _{v}\right)
=48u^4\sin^2v\cos v+ \sqrt{3}u(6-u^2(\cos 2v+5)+12u^3\cos^2 v \sin v)$.
Con lo cual, el flujo sobre la superficie está dado por
\begin{equation}\label{flujo1}
\int_{0}^{1}\int_{0}^{2\pi}{\bf F}\left({\bm\alpha}(u,v)\right) \cdot \left({\bm\alpha}_{u}\times {\bm\alpha} _{v}\right)dvdu=\frac{7}{2}\sqrt{3}\pi
\end{equation}
El proceso descrito antes se obtiene con las siguientes instrucciones:

\begin{tcolorbox}[title=\bf Instrucción \verb"Mathematica"]\begin{Verbatim}[formatcom=\letraverbatim]
r={x,y,z}={Sqrt[3]*u*Cos[v],2*u*Sin[v],
(6-2x^2-y^2)/2};
Integrate[Dot[Cross[D[r,u],D[r,v]],{y^2,x^2,z}],
{u,0,1},{v,0,2*Pi}]
\end{Verbatim}
\end{tcolorbox}

\subsubsection{Flujo sobre el paraboloide \texorpdfstring{$\bm{\mathsf{z+3=x^2+y^2}}$}{z+3=x^2+y^2}}

Esta superficie se puede parametrizar mediante la siguiente función:
{\small
\[{\bm \beta} (u,v)=\left\langle \sqrt{3}u\cos v,2u\sin v,\frac{1}{2}u^{2}\left(7-\cos 2v\right)-3\right\rangle\!,
u\in \lbrack 0,1],\, v\in \lbrack 0,2\pi].\]
}

En este sentido,
\[{\bf F}({\bm \beta}(u,v))=\big \langle 4u^2\sin^2v,3u^2\cos^2v, \linebreak
\frac{1}{2}u^{2}(7-\cos 2v)-3\big \rangle .\]
Derivando parcialmente se hallan los siguientes vectores:
{\small
\begin{eqnarray*}
{\bm \beta}_{u}&=&\left\langle \sqrt{3}\cos v,2\sin v,u\left( 7-\cos 2v\right)\right\rangle \\
{\bm \beta}_{v}&=&\left\langle -\sqrt{3}u\sin v,2u\cos v,u^{2}\sin 2v\right\rangle \\
{\bm \beta}_{v}\times {\bm \beta}_{u}&=&\left\langle 12u^{2}\cos v,8\sqrt{3}u^{2}\sin v,-2\sqrt{3}u\right\rangle\!.
\end{eqnarray*}
}

Además,
{\footnotesize
\begin{multline*}
{\bf F}\left({\bm \beta}(u,v)\right) \cdot \left({\bm \beta}_{v}\times {\bm \beta}_{u}\right)
= 12u^{4}(2\sin v+\sqrt{3}\cos v)\sin(2v)+\sqrt{3}u(u^{2}(\cos(2v)-7) +6).
\end{multline*}
}
Con lo cual, el flujo sobre esta región está dado por
{\small
\begin{equation}\label{flujo2}
\int_{0}^{1}\int_{0}^{2\pi}{\bf F}\left({\bm \beta}(u,v)\right) \cdot \left({\bm \beta}_{v}\times {\bm \beta}_{u}\right)dvdu=\frac{5}{2}\sqrt{3}\pi
\end{equation}
}

El flujo sobre la superficie del sólido se determina sumando las expresiones \ref{flujo1} y \ref{flujo2}.
Dado que la ${\rm div}{\bf F}=1$, aplicando el teorema de la divergencia, el valor obtenido coincide con el volumen del sólido en la fórmula \ref{ejercicio1} de la página \pageref{ejercicio1}.

Con las siguientes instrucciones se obtiene el cálculo anterior:

\begin{tcolorbox}[title=\bf Instrucción \verb"Mathematica"]\begin{Verbatim}[formatcom=\letraverbatim]
r={x,y,z}={Sqrt[3]*u*Cos[v],2*u*Sin[v],x^2+y^2-3};
Integrate[Dot[Cross[D[r,v],D[r,u]],{y^2,x^2,z}],
{u,0,1},{v,0,2*Pi}]
\end{Verbatim}
\end{tcolorbox}

\item Hallar el flujo neto sobre el campo vectorial
${\bf F}(x,y,z) = \la y,-x,z \ra $ sobre la superficie frontera del sólido acotado por el cilindro $x^2+y^2=4$ y la esfera $x^2+y^2+z^2=9$.

\begin{figure}[H]
\centering
\caption{${\bf F}(x,y,z) = \la y,-x,z \ra$ sobre la superficie $x^2+y^2\leq 4$ \& $x^2+y^2+z^2\leq 9$}
\includegraphics[width=0.4\textwidth]{Img/T_9b.pdf} \hspace{5mm}
\includegraphics[width=0.4\textwidth]{Img/T_9c.pdf}
\fuentepropia
\end{figure}

La figura presentada se obtiene ejecutando las siguientes instrucciones:

\begin{tcolorbox}[title=\bf Instrucción \verb"Mathematica"]\begin{Verbatim}[formatcom=\letraverbatim]
a=RegionPlot3D[x^2+y^2<4&&x^2+y^2+z^2<=9,{x,-2,2},
{y,-2,2},{z,-3.5,3.5},PlotPoints->90,Mesh->False]
b=SliceVectorPlot3D[{y,-x,z},{x^2+y^2==4,x^2+y^2
+z^2==9},{x,-2,2},{y,-2,2},{z,-3.4,3.5},
VectorPoints->16,RegionFunction->Function[{x,y,z},
x^2+y^2<=4&&x^2+y^2+z^2<=9],VectorStyle->Black,
VectorScale->{0.15,Scaled[0.2],Automatic},
PlotStyle->Opacity[0]];
Show[a,b]
\end{Verbatim}
\end{tcolorbox}

La intersección de las superficies satisface la relación $z=\pm \sqrt{5}$.
El interior del cilindro se parametriza mediante las expresiones $x=2u\cos v$, $y=2u\sin v$,
para $u\in \lbrack 0,1]$, $v\in \lbrack 0,2\pi ]$. Para este sólido se consideran los siguientes dos casos.

\subsubsection{Flujo sobre la superficie de la esfera}
Utilizaremos la función
${\bf r}\left( u,v\right) =\langle 2u\cos v,2u\sin v,\sqrt{9-4u^{2}}\rangle $,
como parametrización de la esfera, con lo cual ${\bf F}({\bf r}(u,v))=\langle 2u\sin v,-2u\cos v, \linebreak \sqrt{9-4u^{2}}\rangle $.
Derivando parcialmente se obtienen los siguientes vectores:

\begin{eqnarray*}
{\bf r}_{u} &=& \left\langle 2\cos v,2\sin v,\frac{-4u}{\sqrt{9-4u^{2}}}\right\rangle \\
{\bf r}_{v}&=&\left\langle -2u\sin v,2u\cos v,0\right\rangle \\
{\bf r}_{u}\times {\bf r}_{v}&=&\left\langle \frac{8u^{2}\cos v}{\sqrt{9-4u^{2}}},\frac{8u^{2}\sin v}{\sqrt{9-4u^{2}}},4u\right\rangle\!.
\end{eqnarray*}

Con lo cual, {\small ${\bf F}\left({\bf r}\left(u,v\right)\right) \cdot \left({\bf r}_{u}\times {\bf r}_{v}\right)= 4u\sqrt{9-4u^{2}}$.}
Al ser dos regiones iguales (sobre la esfera) multiplicamos por 2 y haremos un solo cálculo, es decir, el flujo sobre la porción de la esfera que limita el sólido está dado por

\begin{equation}\label{flujo3}
2\int_{0}^{1}\int_{0}^{2\pi }4u\sqrt{9-4u^{2}}dvdu= \frac{4}{3}\pi \left( 27-5\sqrt{5}\right)
\end{equation}

El cálculo anterior se obtiene en \verb"Mathematica" con las siguientes instrucciones:

\begin{tcolorbox}[title=\bf Instrucción \verb"Mathematica"]\begin{Verbatim}[formatcom=\letraverbatim]
r={x,y,z}={2*u*Cos[v],2*u*Sin[v],Sqrt[9-x^2-y^2]};
Integrate[Dot[Cross[D[r,u],D[r,v]],{y,-x,z}],
{u,0,1},{v,0,2*Pi}]
\end{Verbatim}
\end{tcolorbox}

\subsubsection{Flujo sobre la superficie del cilindro}

En este caso la parametrización escogida es
${\bf w}\left( u,v\right) =\left\langle 2\cos u,2\sin u,v\right\rangle $, para $ u\in \lbrack 0,2\pi ]$, $v\in \lbrack -\sqrt{5},\sqrt{5}]$.
Con esta función hallamos los siguientes vectores:
\begin{eqnarray*}
{\bf w}_{u}&=&\left\langle -2\sin u,2\cos u,0\right\rangle \\
{\bf w}_{v}&=&\left\langle 0,0,1\right\rangle \\
{\bf w}_{u}\times {\bf w}_{v}&=&\left\langle 2\cos u,2\sin u,0\right\rangle\!.
\end{eqnarray*}
Además, ${\bf F}\left( {\bf w}\left( u,v\right) \right) =\left\langle 2\sin u,-2\cos u,v\right\rangle $
y ${\bf F}\left({\bf w}\left( u,v\right) \right) \cdot \left({\bf w}_{u}\times {\bf w}_{v}\right)= 0$,

con lo cual el flujo sobre la superficie del cilindro esta dado por
\begin{equation}\label{flujo4}
\int_S {\bf F}\cdot d{\bf S}=0
\end{equation}

El flujo neto es la suma de los resultados \ref{flujo3} y \ref{flujo4}. Como ${\rm div}\,{\bf F}=1$ en este caso, el valor hallado coincide con el volumen del sólido encontrado en la expresión \ref{ejercicio4} de la página \pageref{ejercicio4}.

\item Hallar el flujo neto sobre el campo vectorial ${\bf F}(x,y,z) = \la z,y,x \ra $
sobre la superficie frontera del sólido acotada por el hiperboloide de un manto con ecuaciones
$2x^2+3y^2=z^2+7$ y el elipsoide $2x^2+y^2+z^2=9$.

\label{elipsoidehiperboloide}

\begin{figure}[H]
\centering
\caption{Campo vectorial ${\bf F}(x,y,z)=\la z,y,x\ra$}
\includegraphics[width=0.4\textwidth]{Img/T_11c.pdf} \hspace{5mm}
\includegraphics[width=0.4\textwidth]{Img/T_11d.pdf}
\fuentepropia
\end{figure}

Con las siguientes instrucciones se consigue esta figura.

\begin{tcolorbox}[title=\bf Instrucción \verb"Mathematica"]\begin{Verbatim}[formatcom=\letraverbatim]
a=RegionPlot3D[2*x^2+3*y^2<z^2+7&&2*x^2+y^2+z^2<9,
{x,-2,2},{y,-2,2},{z,-3,3},PlotPoints->90,Mesh->False];
b=SliceVectorPlot3D[{z,y,x},{2*x^2+3*y^2==z^2+7,
2*x^2+y^2+z^2==9},{x,-Sqrt[15/4],Sqrt[15/4]},
{y,-Sqrt[15/2],Sqrt[15/2]},{z,-4,3},VectorPoints->20,
VectorStyle->Black,RegionFunction->Function[{x,y,z},
x^2+y^2<=4],VectorScale->{0.07,0.2,None},
PlotStyle->Opacity[0]];
Show[a,b]
\end{Verbatim}
\end{tcolorbox}

Calcularemos el flujo sobre la superficie que acota el sólido, es decir, la porción del elipsoide y, posteriormente, la parte del hiperboloide.

\subsubsection{Flujo sobre la porción del elipsoide}

Con las ecuaciones que definen el elipsoide y el hiperboloide de un manto se elimina $z$ y se obtiene la expresión
$x^2+y^2=4;$ la región que esta ecuación representa se parametriza mediante el uso de las expresiones
$x=2u\cos v$ y $y=2u \sin v$, con $u \in [0,2\pi]$ y $u\in [0,1]$.
El término que define a $z$ está dado por
\[z =\sqrt{9-2x^2-y^2} = \sqrt{9-6u^2-2u^2 \cos(2v)}.\]
Entonces; la función
\[{\bf r}(u,v) = \left \langle 2u\cos v,2u \sin v ,\sqrt{9-6u^2-2u^2 \cos(2v)} \right\rangle\]
parametriza la región de interés sobre el elipsoide. Con esta función se obtienen los siguientes vectores:
{\small
\begin{eqnarray*}
{\bf r}_u &=& \left \langle 2\cos v,2 \sin v , -\frac{2u(3+2\cos v)}{\sqrt{9-6u^2-2u^2 \cos(2v)}} \right\rangle \\
{\bf r}_v &=& \left \langle -2u\sin v,2u \cos v , \frac{2u^2 \sin v}{\sqrt{9-6u^2-2u^2 \cos(2v)}} \right\rangle \\
{\bf r}_u \times {\bf r}_v &=& \left \langle \frac{16u^2 \cos v}{\sqrt{9-6u^2-2u^2 \cos(2v)}}, \frac{8u^2 \sin v}{\sqrt{9-6u^2-2u^2 \cos(2v)}}, 4u \right\rangle\!.
\end{eqnarray*}
}

Además
${\bf F}({\bf r}(u,v) ) = \left \langle \sqrt{9-6u^2-2u^2 \cos(2v)},2u \sin v , 2u\cos v \right\rangle \!$,
con lo cual
{\small
\[{\bf F}({\bf r}(u,v) ) \cdot ({\bf r}_u \times {\bf r}_v)=\frac{ 8u \left (u-u \cos (2v) +3 \cos v \, \sqrt{9-2u^2(3+\cos(2v))}\right )}{\sqrt{9-6u^2-2u^2 \cos(2v)}};\]
}

por lo tanto, el flujo sobre la porción del elipsoide está dado por

\begin{equation} \label{ejercicio6e}
\int_0^1 \int_0^{2\pi} {\bf F}({\bf r}(u,v) ) \cdot ({\bf r}_u \times {\bf r}_v) dvdu \approx 10.844
\end{equation}
Este cálculo se hizo en \verb"Mathematica" ejecutando las siguientes instrucciones:

\begin{tcolorbox}[title=\bf Instrucción \verb"Mathematica"]\begin{Verbatim}[formatcom=\letraverbatim]
r={x,y,z}={2*u*Cos[v],2*u*Sin[v],Sqrt[9-2*x^2-y^2]};
NIntegrate[2*Dot[{z,y,x},Cross[D[r,u],D[r,v]]],
{v,0,2*Pi},{u,0,1}]
\end{Verbatim}
\end{tcolorbox}

Antes de presentar el siguiente procedimiento, nótese que la función $f(u)= (b-a)u+a$ satisface que $f(0)=a$ y $f(1)=b$.
De esta manera $f$ aplica el intervalo $[0,1]$ en el intervalo $[a,b]$. A su vez la función inversa
$g(v)= \frac{v-a}{b-a}$ aplica el intervalo $[a,b]$ en el intervalo $[0,1]$.
Usaremos estas ideas para parametrizar la región acotada por la porción de la circunferencia y la hipérbola que se des\-cribe a continuación.

\subsubsection{Parametrizando la región sobre el hiperboloide}

Al eliminar $y$ de las ecuaciones originales se obtiene $x^2+z^2=5$,
esta expresión representa, en el plano $xz$, una circunferencia de radio $\sqrt{5}$ que se parametriza como
$x=\sqrt{5}\cos t$, $z=\sqrt{5}\sin t$.
El hiperboloide proyectado en el plano $xz$ corresponde a un una hipérbola con ecuación $2x^2-z^2=7$.
Esta región se puede describir definiendo $x=\sqrt{\frac{7}{2}}\cosh (t),\,\, z=\sqrt{7}\sinh (t)$.
Las anteriores dos curvas se intersectan en los puntos $(\pm 2, \pm 1)$.

Entonces la región sobre la circunferencia para $x\in [-2,2]$ se describe por medio de la función
\[\overrightarrow{\alpha}(t)=\langle \sqrt{5}\cos t, \sqrt{5}\sin t \rangle, \,\,
t \in \left[ \arccos \left( \frac{2}{\sqrt{5}}\right),\pi -\arccos \left( \frac{2}{\sqrt{5}}\right) \right].\]

La región sobre la hipérbola desde el punto $(2,-1)$ al punto $(2,1)$ se describe por medio de la función
\[\overrightarrow{\beta}(t) = \left \langle \sqrt{\frac{7}{2}}\cosh (t), \sqrt{7}\sinh (t) \right \rangle, \,t \in \left[ -t_0 ,t_0 \right].\]
Se usó la abreviatura $t_0={\rm arccosh}\left( \sqrt{\frac{2}{7}}\right)$.
Estas curvas se muestran a continuación:

\begin{figure}[H]
\centering
\caption{Parametrización de la proyección del sólido en plano $xz$}
\begin{minipage}[c]{0.39\linewidth}
\centering

\includegraphics[width=\linewidth]{Fig/188.pdf}
\end{minipage}
\hspace{1em}
\begin{minipage}[c]{0.32\linewidth}
\centering

\includegraphics[width=\linewidth]{Fig/189.pdf}
\end{minipage}

\fuentepropia
\end{figure}

Para abreviar la escritura, usaremos las siguientes constantes:
{\small
\[a= \arccos \left( \frac{2}{\sqrt{5}}\right),\,\, b=\pi -\arccos \left( \frac{2}{\sqrt{5}}\right)\]
\[c= -{\rm arccosh}\left( \frac{2\sqrt{2}}{\sqrt{7}}\right) ,\, d={\rm arccosh}\left( \frac{2\sqrt{2}}{\sqrt{7}}\right)\!.\]}

Al redefinir las funciones que describen las curvas en mención con el sentido deseado y unificar el intervalo de definición a $t \in [0,1]$,

{\small
\begin{eqnarray*}
\overrightarrow{\alpha}(t) &=& \left \langle \sqrt{7/2}\cosh((d-c)t+c),\sqrt{7}\sinh((d-c)t+c)\right \rangle,\\
\overrightarrow{\beta}(t) &=&\left \langle \sqrt{5}\cos((b-a)(1-t)+a),\sqrt{5}\sin((b-a)(1-t)+a) \right \rangle.
\end{eqnarray*}
}

Formando una combinación convexa de estas expresiones se parametriza la región sombreada en el gráfico.
Con las siguientes instrucciones se dibujan las curvas mencionadas y se consigue una animación en \verb"Mathematica" que permite ver el recorrido de la región que se considera.

\begin{tcolorbox}[title=\bf Instrucción \verb"Mathematica"]\begin{Verbatim}[formatcom=\letraverbatim]
{a,b}={ArcCos[2/Sqrt[5]],Pi-a};
{c,d}={-ArcCosh[2*Sqrt[2/7]],c};
R[w_]:={Sqrt[5]*Cos[(b-a)*(1-w)+a],
Sqrt[5]*Sin[(b-a)*(1-w)+a]};
r[w_]:={Sqrt[7/2]*Cosh[(d-c)*w+c],
Sqrt[7]*Sinh[(d-c)*w+c]};
Animate[ParametricPlot[{R[u],(1-u)*r[t]+u*R[t],r[u]},
{u,0,1},PlotRange->{{-2,3},{-2,3}},PlotStyle->{Red,
Black,Red}],{t,0,1}]
\end{Verbatim}
\end{tcolorbox}

La función obtenida al hacer esta construcción es la siguiente:

\begin{multline}
{\bf r}(t,u)=(1-u)\left(\sqrt{7/2}\cosh((d-c)t+c),\sqrt{7}\sinh((d-c)t+c)\right )+ \\
u\left (\sqrt{5}\cos((b-a)(1-t)+a),\sqrt{5}\sin((b-a)(1-t)+a)\right) \label{ejercicio6ec}
\end{multline}
En donde $t \in [0,1],\,\, u \in [0,1]$. La expresión correspondiente a $y$ se consigue al sustituir $x$ y $z$
de la función anterior en la expresión que define el hiperboloide de un manto. La región obtenida corresponde a la cuarta parte de la superficie mencionada.

\begin{figure}[H]
\centering
\caption{Parametrización de la superficie del sólido}
\includegraphics[width=0.4\textwidth]{Img/T_11e.pdf} \hspace{5mm}
\includegraphics[width=0.4\textwidth]{Img/T_11f.pdf}
\fuentepropia
\end{figure}

\begin{tcolorbox}[title=\bf Instrucción \verb"Mathematica"]\begin{Verbatim}[formatcom=\letraverbatim]
r={x,y,z}={2*u*Cos[v],2*u*Sin[v],Sqrt[9-2*x^2-y^2]}
A1=ParametricPlot3D[{{x,y,z},{x,y,-z}},{u,0,1},
{v,0,2*Pi},PlotStyle->{RGBColor[0.8,0.2,0.2]},
MeshStyle->Black,PlotPoints->90];
{a,b}={ArcCos[2/Sqrt[5]],Pi-ArcCos[2/Sqrt[5]]};
{c,d}={-ArcCosh[2*Sqrt[2/7]],ArcCosh[2*Sqrt[2/7]]};
R=(1-u)*{Sqrt[7/2]*Cosh[(d-c)*t+c],Sqrt[7]*Sinh[(d-
c)*t+c]}+u*{Sqrt[5]*Cos[(b-a)*(1-t)+a],
Sqrt[5]*Sin[(b-a)*(1-t)+a]};
{x,z,y}={R[[1]],R[[2]],Sqrt[(z^2+7-2*x^2)/3]};
A2=ParametricPlot3D[{{x,y,z},{x,-y,z},-{x,y,z},
-{x,-y,z}},{u,0,1},{t,0,1},PlotStyle->{Cyan,Orange,
Cyan,Orange},PlotPoints->60];
Show[A2,A1,BoxRatios->{1,1,1}]
\end{Verbatim}
\end{tcolorbox}

Utilizando la expresión \ref{ejercicio6ec}, se establece que el flujo sobre la superficie del hiperboloide $H$,
está dado por la expresión
\begin{equation} \label{ejercicio6h}
4 \iint_H {\bf F}({\bf r}(t,u) ) \cdot ({\bf r}_u \times {\bf r}_v) dt dt \approx 41.0176
\end{equation}
Este valor se obtuvo con \verb"Mathematica" con estas instrucciones:

\begin{tcolorbox}[title=\bf Instrucción \verb"Mathematica"]\begin{Verbatim}[formatcom=\letraverbatim]
{a,b,d,c}={ArcCos[2/Sqrt[5]],Pi-a,
ArcCosh[2*Sqrt[2/7]],-d};
R=(1-u)*{Sqrt[7/2]*Cosh[(d-c)*t+c],
Sqrt[7]*Sinh[(d-c)*t+c]}+u*{Sqrt[5]*Cos[(b-a)*(1-t)
+a],Sqrt[5]*Sin[(b-a)*(1-t)+a]};
r={x,y,z}={R[[1]],Sqrt[(z^2+7-2*x^2)/3],R[[2]]};
NIntegrate[4*Dot[{z,y,x},Cross[D[r,u],D[r,t]]],
{u,0,1},{t,0,1}]
\end{Verbatim}
\end{tcolorbox}

El flujo sobre la superficie del sólido es la suma de las expresiones
\ref{ejercicio6h} y \ref{ejercicio6e}.
Como ${\rm div}\, {\rm F} =1$, el valor del flujo neto de ${\bf F}$ coincide con el volumen del sólido que se halló en la página \pageref{ejercicio6}, con la expresión \ref{ejercicio6}.

\item Hallar el flujo neto sobre el campo vectorial
${\bf F}(x,y,z) = \la 3x,y,3z \ra $ sobre la superficie frontera del sólido acotado por el paraboloide $x^2+y^2=3-z$ y el paraboloide hiperbólico $x^2-y^2=3z+6$. \index{Paraboloide!hiperbólico}

\begin{figure}[H]
\centering
\caption{Campo vectorial ${\bf F}(x,y,z)=\la 3x,y,3z\ra$}
\includegraphics[width=0.4\textwidth]{Img/T_10c.pdf} \hspace{5mm}
\includegraphics[width=0.4\textwidth]{Img/T_10d.pdf}
\fuentepropia
\end{figure}

Esta figura es el producto de ejecutar las siguientes instrucciones:

\begin{tcolorbox}[title=\bf Instrucción \verb"Mathematica"]\begin{Verbatim}[formatcom=\letraverbatim]
a=RegionPlot3D[x^2+y^2<3-z&&x^2-y^2<3*z+6,{x,-2,2},
{y,-3,3},{z,-5,3},PlotPoints->90,Mesh->False]
b=SliceVectorPlot3D[{3*x,y,3*z},{x^2+y^2==3-z,x^2-
y^2==3*z+6},{x,-Sqrt[15/4],Sqrt[15/4]},{y,-Sqrt[15/2],
Sqrt[15/2]},{z,-5,3},RegionFunction->Function[{x,y,z},
4*x^2+2*y^2<=15],VectorPoints->18,VectorStyle->Black,
VectorScale->{0.14,Scaled[0.25],Automatic},
PlotStyle->Opacity[0]];
Show[a,b]
\end{Verbatim}
\end{tcolorbox}

En la página \pageref{ejercicio4a} se trata esta región.
La expresión $4x^{2}+2y^{2}=15$ describe la curva de intersección en el plano $xy$. Esta región puede representarse mediante las expresiones
$x=\frac{\sqrt{15}}{2}u\cos v$, $y=\sqrt{\frac{15}{2}}u\sin v$, con $u\in \lbrack 0,1]$,
$v\in \lbrack 0,2\pi ]$.
Calcularemos el flujo neto sobre la superficie del sólido en dos partes: primero sobre la superficie del paraboloide y luego sobre el paraboloide hiperbólico.

\subsubsection{Flujo sobre el paraboloide}

Al reemplazar las expresiones anteriores en la ecuación $x^2+y^2=3-z$, se obtiene que
$z=\frac{15}{8}u^{2}\cos 2v-\frac{45}{8}u^{2}+3;$
por lo tanto, una parametrización de la superficie en mención es
\[{\bf r}\left( u,v\right) =\left\langle \frac{\sqrt{15}}{2}u\cos v,\sqrt{\frac{15}{2}}u\sin v,
\frac{15}{8}u^{2}\cos 2v-\frac{45}{8}u^{2}+3\right\rangle\!,\]
en donde $u\in \lbrack 0,1]$, $v\in \lbrack 0,2\pi ]$.
Con la función ${\bf r}$, al evaluarla en el campo vectorial se sigue que
\[{\bf F}\left({\bf r}\left( u,v\right) \right) =\left\langle \frac{3\sqrt{15}}{2}u\cos v,\sqrt{\frac{15}{2}}u\sin v,\frac{45}{8}u^{2}\cos 2v-\frac{135}{8}u^{2}+9\right\rangle\!.\]
% ${\bf F}\left({\bf r}\left( u,v\right) \right) =\left\langle \frac{3\sqrt{15}}{2}u\cos v,\sqrt{\frac{15}{2}}u\sin v,\frac{45}{8}u^{2}\cos 2v-\frac{135}{8}u^{2}+9\right\rangle ,$
Derivando parcialmente se obtienen los siguientes vectores;
\begin{eqnarray*}
{\bf r}_{u} &=&\left\langle \frac{1}{2}\sqrt{15}\cos v,\frac{1}{2}\sqrt{2}\sqrt{15}%
\sin v,\frac{15}{4}u\cos 2v-\frac{45}{4}u \right\rangle \\
{\bf r}_{v}&=& \left\langle -\frac{1}{2}\sqrt{15}u\sin v,\frac{1}{2}\sqrt{2}\sqrt{15}u\cos v,-\frac{15}{4}u^{2}\sin 2v \right\rangle \\
{\bf r}_{u}\times {\bf r}_{v}&=&\left\langle \frac{15}{4}\sqrt{2}\sqrt{15}u^{2}\cos v,\frac{15}{2}\sqrt{15}u^{2}\sin v,\frac{15}{4}\sqrt{2}u\right\rangle;
\end{eqnarray*}
entonces,
\[{\bf F}\left({\bf r}\left(u,v\right) \right) \cdot \left({\bf r}_{u}\times {\bf r}_{v}\right) =\frac{225\sqrt{2}}{32}u\left( u^{2}+\frac{24}{5}+5u^{2}\cos 2v\right)\!.\]
Por lo expuesto antes, el flujo sobre esta región está dado por
\begin{equation}\label{flujo5}
\int_{0}^{1}\int_{0}^{2\pi }\left( \frac{225\sqrt{2}}{32}u\left( u^{2}+ \frac{24}{5}+5u^{2}\cos 2v\right) \right) dvdu= \frac{2385}{64}\sqrt{2}\pi
\end{equation}
Este valor también se obtiene como producto de ejecutar las siguientes instrucciones:

\begin{tcolorbox}[title=\bf Instrucción \verb"Mathematica"]\begin{Verbatim}[formatcom=\letraverbatim]
r={x,y,z}={Sqrt[15]/2*u*Cos[v],Sqrt[15/2]*u*Sin[v],
3-x^2-y^2};
Integrate[Dot[Cross[D[r,u],D[r,v]],{3x,y,3z}],
{u,0,1},{v,0,2*Pi}]
\end{Verbatim}
\end{tcolorbox}

\subsubsection{Flujo sobre el paraboloide hiperbólico}

Con las mismas ecuaciones que describen el conjunto de puntos que satisfacen la condición $4x^{2}+2y^{2} \leq 15$,
la coordenada $z$ sobre esta superficie satisface que
\[z=\frac{x^2-y^2-6}{6}=\frac{15}{8}u^{2}\cos 2v-\frac{5}{8}u^{2}-2,\]
por lo tanto, utilizaremos la siguiente parametrización, para describir la porción del paraboloide hiperbólico:

{\small
\[{\bf r}\left( u,v\right) =\left\langle \frac{\sqrt{15}}{2}u\cos v,\sqrt{\frac{15}{2}}u\sin v,\frac{15}{8}u^{2}\cos 2v-\frac{5}{8}u^{2}-2\right\rangle\!,\]}

en donde $u\in \lbrack 0,1]$, $v\in \lbrack 0,2\pi ]$.

Con ella se consiguen los siguientes vectores:

{\small
\begin{align*}
{\bf r}_{u}&=\left\langle \frac{1}{2}\sqrt{15}\cos v,\frac{1}{2}\sqrt{2}\sqrt{15}\sin v,\frac{15}{4}u\cos 2v-\frac{5}{4}u\right\rangle \\
{\bf r}_{v}&= \left\langle -\frac{1}{2}\sqrt{15}u\sin v,\frac{1}{2}
\sqrt{2}\sqrt{15}u\cos v,-\frac{15}{4}u^{2}\sin 2v \right\rangle \\
{\bf r}_{v}\times {\bf r}_{u}&=\left\langle \frac{5}{4}\sqrt{2}\sqrt{15}u^{2}\cos v,-\frac{5}{2}\sqrt{3}\sqrt{5}u^{2}\sin v,-\frac{15}{4}\sqrt{2}u\right\rangle\!.
\end{align*}
}

Es de anotar que para conservar la orientación de la superficie se calculó ${\bf r}_{v}\times {\bf r}_{u}$. Al evaluar la parametrización en el campo vectorial, se sigue que

{\small
\[{\bf F}\left( {\bf r}\left( u,v\right) \right) =\left\langle \frac{3\sqrt{15}}{2}u\cos
v,\sqrt{\frac{15}{2}}u\sin v,\frac{45}{8}u^{2}\cos 2v-\frac{15}{8}u^{2}-6\right\rangle\!,\]}

entonces,

{\small
\[{\bf F}\left({\bf r}\left( u,v\right) \right) \cdot \left({\bf r}_{u}\times {\bf r}_{v}\right) =
\frac{75\sqrt{2}}{32}u\left( 5u^{2}+\frac{48}{5}+u^{2}\cos 2v\right)\!.\]}

Por lo señalado antes, el flujo sobre esta superficie está dado por
{\small
\begin{equation}\label{flujo6}
\int_{0}^{1}\int_{0}^{2\pi }\left( \frac{75\sqrt{2}}{32}u\left( 5u^{2} + \frac{48}{5}+u^{2}\cos 2v\right) \right) dvdu= \frac{1815}{64}\sqrt{2}\pi
\end{equation}
}

Con las siguientes instrucciones se hace el cálculo anterior:

\begin{tcolorbox}[title=\bf Instrucción \verb"Mathematica"]\begin{Verbatim}[formatcom=\letraverbatim]
r={x,y,z}={Sqrt[15]/2*u*Cos[v],
Sqrt[15/2]*u*Sin[v],(x^2-y^2-6)/6};
Integrate[Dot[Cross[D[r,v],D[r,u]],{3x,y,3z}],
{u,0,1},{v,0,2*Pi}]
\end{Verbatim}
\end{tcolorbox}

El flujo neto sobre toda la superficie del sólido es la suma de los resultados hallados en \ref{flujo5} y
\ref{flujo6}, esto es \[\frac{2385}{64}\sqrt{2}\pi +\frac{1815}{64}\sqrt{2}\pi = \frac{525}{8}\sqrt{2}\pi.\]
Dado que ${\rm div}{\bf F}=7$, evidentemente, el valor obtenido es $7$ veces el volumen del sólido que se obtuvo en la expresión \ref{ejercicio3} de la página \pageref{ejercicio3}

\item Hallar el flujo neto del campo vectorial ${\bf F}(x,y,z) = \la y,-x,z \ra $
sobre la superficie del sólido interior a la esfera con ecuacion $x^2+y^2+z^2=9 $ y al cilindro cardioidico
$r=1+\cos \theta$.

\begin{figure}[H]
\centering
\caption{Campo vectorial ${\bf F}(x,y,z)=\la y,-x,z\ra$, sobre $x^2+y^2+z^2=9 $ \& $r=1+\cos \theta$}
\label{cvFxyzsobrex2y2z2}
\includegraphics[width=0.4\textwidth]{Img/T_12c.pdf} \hspace{5mm}
\includegraphics[width=0.4\textwidth]{Img/T_12d.pdf}
\fuentepropia
\end{figure}

Las instrucciones que permiten obtener la figura~\ref{cvFxyzsobrex2y2z2} son las siguientes.

\begin{tcolorbox}[title=\bf Instrucción \verb"Mathematica"]\begin{Verbatim}[formatcom=\letraverbatim]
a=RegionPlot3D[(x^2+y^2-x)^2<(x^2+y^2)&&
x^2+y^2+z^2<9,{x,-1,2.2},{y,-2,2},{z,-3,3.2},
PlotPoints->90,Mesh->False,BoxRatios->{1,1,1}];
b=SliceVectorPlot3D[{y,-x,z},{(x^2+y^2-x)^2==x^2+y^2,
x^2+y^2+z^2==9},{x,-1,2.2},{y,-2,2},{z,-3,3.2},
RegionFunction->Function[{x,y,z},(x^2+y^2-x)^2<(x^2
+y^2)&&x^2+y^2+z^2<9],VectorPoints->17,
VectorStyle->Black,VectorScale->{0.12,Scaled[0.25],
Automatic},PlotStyle->Opacity[0]];
Show[a,b]
\end{Verbatim}
\end{tcolorbox}

A partir de la ecuación $r = 1 + \cos \theta$, se puede construir en el plano $xy$ la parametrización de la región que rellena esta cardioide. Usaremos las ecuaciones:
\[
x = u(1 + \cos \theta) \cos \theta, \quad
y = u(1 + \cos \theta) \sin \theta, \quad
u \in [0,1],\; \theta \in [0,2\pi].
\]

La región sobre la esfera se parametriza completamente hallando $z$, es decir:
\[
z = \sqrt{9 - x^2 - y^2} = \sqrt{9 - u^2(1 + \cos \theta)^2}.
\]

Por lo tanto, la siguiente función describe la porción sobre la esfera interior al cilindro cardioidico:
\[
{\bf r}(u,\theta) =
\left\langle
u(1 + \cos \theta) \cos \theta,\;
u(1 + \cos \theta) \sin \theta,\;
\sqrt{9 - u^2(1 + \cos \theta)^2}
\right\rangle\!,
\]
en donde $u \in [0,1],\; \theta \in [0,2\pi]$. Es sencillo ver que:
{\small
\[
{\bf F}({\bf r}(u,\theta)) =
\left\langle
u(1 + \cos \theta) \sin \theta,\;
- u(1 + \cos \theta) \cos \theta,\;
\sqrt{9 - u^2(1 + \cos \theta)^2}
\right\rangle\!.
\]
}

Con la parametrización elegida hallamos los siguientes vectores:
{\footnotesize
\begin{eqnarray*}
{\bf r}_{u} &=& \left \langle (\cos \theta +1) \cos \theta ,\left(\cos \theta +1\right) \sin \theta ,\frac{-4u\cos^2\left( \frac{\theta}{2} \right)}{\sqrt{9-u^2(1+\cos(\theta))^2}}\right\rangle \\
{\bf r}_{\theta}&=&\left\langle -u( 1+2\cos \theta)\sin \theta ,u(\cos \theta +\cos (2\theta)),\frac{u^{2}(\sin 2\theta+2\sin \theta)}{\sqrt{2}\sqrt{9-u^2(1+\cos(\theta))^2}} \right\rangle ,
\end{eqnarray*}
}
cuyo producto cruz es
{\small
\begin{multline*}
{\bf r}_{u}\times {\bf r}_{\theta}= \bigg \langle \frac{8u^2\cos^6(\frac{\theta}{2})\cos(\theta)}{\sqrt{9-u^2(1+\cos(\theta))^2}}, \frac{6u^2\cos^7(\frac{\theta}{2})\sin(\frac{\theta}{2})}{\sqrt{9-u^2(1+\cos(\theta))^2}},4u\cos^4\left(\frac{\theta}{2}\right )\bigg \rangle.
\end{multline*}
}
Además,
\[{\bf F}\left({\bf r}\left( u,\theta \right) \right) \cdot \left({\bf r}_{u}\times {\bf r}_{\theta}\right) =
4u \cos^4\left(\frac{\theta}{2}\right)\sqrt{9-u^2(1+\cos(\theta))^2}.\]
Por lo tanto, el flujo sobre la superficie de la esfera e interior al cilindro cardioidico, esta dado por
\begin{equation}\label{flujo7}
2\int_{0}^{1}\int_{0}^{2\pi}{\bf F}\left({\bf r}\left( u,\theta \right) \right) \cdot \left({\bf r}_{u}\times {\bf r}_{\theta}\right) d\theta du\approx 25.8179
\end{equation}
Se ha multiplicado por $2$ debido a que la región considerada es similar al hemisferio opuesto de la esfera.
El cálculo anterior se hizo en \verb"Mathematica" utilizando las siguientes instrucciones:

\begin{tcolorbox}[title=\bf Instrucción \verb"Mathematica"]\begin{Verbatim}[formatcom=\letraverbatim]
{x,y}={u*(1+Cos[t])*Cos[t],u*(1+Cos[t])*Sin[t]};
r={x,y,Sqrt[9-x^2-y^2]};
A=Dot[{y,-x,z},TrigExpand[Cross[D[r,u],D[r,t]]]];
2*NIntegrate[A,{u,0,1},{t,0,2*Pi}]
\end{Verbatim}
\end{tcolorbox}

Para calcular el flujo sobre el cilindro cardióidico, se debe parametrizar dicha región.
Al proyectar el cilindro sobre el plano $xy$, la curva resultante puede parametrizarse en la forma
$x=\left( 1+\cos \theta \right) \cos \theta $, $y=\left( 1+\cos \theta\right) \sin \theta $, $\theta \in [0,2\pi]$.
Sobre la esfera, la coordenada $z$ puede ser de dos formas
\[y = \pm \frac{\sqrt{15-4\cos(\theta)-\cos(2\theta)}}{\sqrt{2}}.\]
Con dos puntos de la forma;
{\small
\[P=\left(\left( 1+\cos \theta \right) \cos \theta, \left( 1+\cos \theta \right) \sin \theta,\frac{1}{\sqrt{2}}\sqrt{15-4\cos(\theta)-\cos(2\theta)} \right )\]
\[Q=\left(\left( 1+\cos \theta \right)\cos\theta,\left(1+\cos\theta\right)\sin\theta,-\frac{1}{\sqrt{2}}\sqrt{15-4\cos(\theta)-\cos(2\theta)} \right)\]
}

se puede hacer una combinación convexa que determine una función ${\bf r}$ que parametrice la región en mención. Esto es

{\footnotesize
\[{\bf r}(u,\theta ) =\bigg \langle \left( \cos \theta +1\right ) \cos \theta,\linebreak
(1+\cos\theta) \sin \theta, \frac{(1-2u)\sqrt{15-4\cos(\theta) -\cos(2\theta)}}{\sqrt{2}}\bigg\rangle,\]
}

entonces,
{\footnotesize
\[{\bf F}\left({\bf r}\left(u,\theta \right) \right)  \linebreak
=\big \langle \left( 1+\cos \theta \right) \sin \theta , -\left( \cos \theta +1\right) \cos \theta ,\frac{(1-2u)\sqrt{15-4\cos(\theta)-\cos(2\theta)}}{\sqrt{2}} \big \rangle.\]
}

Con la función ${\bf r}$ se obtienen los siguientes vectores:
{\footnotesize
\begin{eqnarray*}
{\bf r}_{u} &=& \left \langle 0,0,-\sqrt{2}\sqrt{15-4\cos(\theta)-\cos(2\theta)}\right\rangle \\
{\bf r}_{\theta}&=& \left\langle -\sin(\theta)-\sin(2\theta),\cos(\theta)+\cos(2\theta),\frac{(1-2u)(2\sin(\theta)+\sin(2\theta))}{\sqrt{2}\sqrt{15-4\cos(\theta)-\cos(2\theta)}} \right \rangle;
\end{eqnarray*}
}

además, estos vectores permiten calcular
\begin{multline*}
{\bf r}_{u}\times {\bf r}_{\theta } = \bigg\langle 2\sqrt{2}\cos^2(\theta/2)(-1+2\cos(\theta))\sqrt{15-4\cos(\theta)-\cos(2\theta)}, \\
\sqrt{2}(1+2\cos(\theta))\sqrt{15-4\cos(\theta)-\cos(2\theta)}\sin(\theta),0 \bigg\rangle .
\end{multline*}
La densidad de flujo está dada por
\[{\bf F}\left({\bf r}\left( u,\theta \right) \right) \cdot \left({\bf r}_{u}\times {\bf r}_{\theta}\right) =\frac{1}{\sqrt{2}}\sqrt{15 - 4 \cos(\theta) - \cos(2\theta)},\]
cuya integral proporciona el flujo de ${\bf F}$ sobre la superficie del cilindro cardióidico, esto es
\begin{equation}\label{flujo8}
\int_{0}^{1}\int_{0}^{2\pi} \frac{1}{\sqrt{2}}\sqrt{15 - 4 \cos(\theta) - \cos(2\theta)}\,
d\theta du=0.
\end{equation}
Estos cálculos se obtienen ejecutando las siguientes instrucciones:

\begin{tcolorbox}[title=\bf Instrucción \verb"Mathematica"]\begin{Verbatim}[formatcom=\letraverbatim]
{x,y,z}={(1+Cos[t])*Cos[t],(1+Cos[t])*Sin[t],
Sqrt[9-x^2-y^2]};
r=Simplify[(1-u)*{x,y,z}+u*{x,y,-z}];
A=Simplify[Dot[{y,-x,z},Cross[D[r,u],D[r,t]]]];
Integrate[A,{u,0,1},{t,0,2*Pi}]
\end{Verbatim}
\end{tcolorbox}

El flujo neto sobre la superficie del sólido considerado es la suma de las expresiones \ref{flujo7}
y \ref{flujo8}.
Dado que ${\rm div} {\bf F}=1$, este valor coincide con el volumen del sólido que se halló con la expresión \ref{ejercicio3}, página \pageref{ejercicio3}.

\item Hallar el flujo neto sobre el campo vectorial ${\bf F}(x,y,z) = \la x,1,y+z \ra $
sobre la superficie sólido acotado por el semicono $y=\sqrt{x^2+z^2}$ y la esfera $x^2+y^2+z^2=9$.

La siguiente figura ilustra esta situación.

\begin{figure}[H]
\centering
\caption{Campo vectorial ${\bf F}(x,y,z) = \la x,1,y+z \ra $ }
\includegraphics[width=0.4\textwidth]{Img/T_14c.pdf} \hspace{5mm}
\includegraphics[width=0.4\textwidth]{Img/T_14d.pdf}
\fuentepropia
\end{figure}

La figura se consigue ejecutando las siguientes instrucciones:

\begin{tcolorbox}[title=\bf Instrucción \verb"Mathematica"]\begin{Verbatim}[formatcom=\letraverbatim]
a=RegionPlot3D[y>Sqrt[x^2+z^2]&&x^2+y^2+z^2<9,
{x,-3,3},{y,0,3},{z,-3,3},PlotPoints->90,Mesh->False];
b=SliceVectorPlot3D[{x,1,y+z},{x^2+y^2+z^2==9,
y==Sqrt[x^2+z^2]},{x,-3,3},{y,0,3},{z,-3,3},
VectorPoints->18,VectorScale->{0.16,Scaled[0.25],
Automatic},VectorStyle->Black,RegionFunction->
Function[{x,y,z},x^2+z^2<=9/2],PlotStyle->Opacity[0]];
Show[a,b]
\end{Verbatim}
\end{tcolorbox}

Las superficies mencionadas se trataron en la página \pageref{ejercicio7a}.
Su intersección se caracteriza por la ecuación
$x^{2}+z^{2}=\frac{9}{2}$. Por lo tanto, la porción del cono se parametriza como
${\bf r}\left( u,v\right) =\left\langle u\cos v,u,u\sin v\right\rangle $
y la porción de la esfera se parametriza en la forma
${\bf r}\left(u,v\right) =\left\langle u\cos v,\sqrt{9-u^{2}},u\sin v\right\rangle $, en donde $u\in \lbrack 0,\frac{3}{\sqrt{2}}]$, $v\in \lbrack 0,2\pi ]$.
Calcularemos el flujo en cada región por separado.

\subsubsection{Flujo sobre la superficie del cono}

En este caso ${\bf F}\left( {\bf r}\left( u,v\right) \right) =\left\langle u\cos v,1,u+u\sin v\right\rangle $
y también
\begin{align*}
{\bf r}_{u} &= \left\langle \cos v, 1, \sin v \right\rangle \\
{\bf r}_{v} &=\left\langle -u\sin v,0,u\cos v\right\rangle \\
{\bf r}_{u}\times {\bf r}_{v}&= \left\langle u\cos v,-u,u\sin v\right\rangle\!,
\end{align*}
con lo cual,
${\bf F}\left({\bf r}\left(u,v\right)\right)\cdot \left({\bf r}_{u}\times {\bf r}_{v}\right) = u^{2}\sin v-u+u^{2}$.
Entonces, el flujo de ${\bf F}$ sobre la superficie del cono está dado por

\begin{equation}
\label{flujo9}
\int_{0}^{2\pi}\int_0^{\frac{3}{\sqrt{2}}}\left(u^{2}\sin v-u+u^{2}\right)dudv= \frac{9}{2}\pi \left( \sqrt{2}-1\right)
\end{equation}

\subsubsection{Flujo sobre la superficie de la esfera}

En este sentido,
${\bf F}\left( {\bf r}\left( u,v\right) \right) =\left\langle u\cos v,1,\sqrt{9-u^{2}}+u\sin v\right\rangle \!$, y con la parametrización escogida se obtienen los siguientes vectores:
\begin{align*}
{\bf r}_{u} &= \left\langle \cos v,-\frac{u}{\sqrt{9-u^{2}}},\sin v\right\rangle \\
{\bf r}_{v} &= \left\langle -u\sin v,0,u\cos v\right\rangle \\
{\bf r}_{v}\times {\bf r}_{u}&= \left\langle\frac{u^{2}\cos v}{\sqrt{9-u^{2}}},u,\frac{u^{2}\sin v}{\sqrt{9 -u^{2}}}\right\rangle\!.
\end{align*}
De lo anterior, se obtiene
${\bf F}\left( {\bf r}\left( u,v\right) \right) \cdot \left( {\bf r}_{u}\times {\bf r}_{v}\right) =\frac{u^{3}}{\sqrt{9-u^{2}}}+u+u^{2}\sin v$,
cuya integración con los límites establecidos proporciona el flujo de ${\bf F}$ sobre la porción de la esfera.
\begin{equation}\label{flujo10}
\int_{0}^{2\pi }\int_{0}^{\frac{3}{\sqrt{2}}}\left( \frac{u^{3}}{\sqrt{9-u^{2}}}+u+u^{2}\sin v\right) dudv= \frac{9}{2}\pi \left( 9-5\sqrt{2}\right)
\end{equation}

El flujo sobre toda la superficie es la suma de los valores obtenidos en \ref{flujo9}
y \ref{flujo10}. Como ${\rm div\,}{\bf F}=2$, es claro que este valor corresponde a dos veces el volumen del sólido el cual se obtuvo con la expresión \ref{ejercicio7}.

\item Hallar el flujo neto del campo vectorial ${\bf F}(x,y,z) = \la x,y,2z \ra $
sobre la superficie del sólido determinado por la expresión $\rho =1+\cos \phi$.

\begin{figure}[H]
\centering
\caption{Campo vectorial ${\bf F}(x,y,z)=\la x,y,2z \ra$}
\includegraphics[width=0.4\textwidth]{Img/T_13b.pdf} \hspace{5mm}
\includegraphics[width=0.4\textwidth]{Img/T_13c.pdf}
\fuentepropia
\end{figure}

La figura anterior se obtiene al ejecutar las siguientes instrucciones:

\begin{tcolorbox}[title=\bf Instrucción \verb"Mathematica"]\begin{Verbatim}[formatcom=\letraverbatim]
a=RegionPlot3D[(x^2+y^2+z^2-x)^2<x^2+y^2+z^2,
{x,-0.5,2.2},{y,-1.5,1.5},{z,-1.5,1.5},PlotPoints->90,
Mesh->False,BoxRatios->{1,1,1}];
b=SliceVectorPlot3D[{x,y,2*z},{(x^2+y^2+z^2-x)^2==
x^2+y^2+z^2},{x,-0.5,2.2},{y,-1.5,1.5},{z,-1.5,1.5},
VectorStyle->Black,VectorScale->{0.12,Scaled[0.25],
Automatic},VectorPoints->15,PlotStyle->Opacity[0]];
Show[a,b]
\end{Verbatim}
\end{tcolorbox}

Con la ecuación \(\rho = 1 + \cos \phi\) y las expresiones que definen las coordenadas esféricas, se propone la siguiente función:
\[
\mathbf{r}(\theta,\phi) =
\left\langle
(1 + \cos \phi)\sin \phi \cos \theta,\
(1 + \cos \phi)\sin \phi \sin \theta,\
(1 + \cos \phi)\cos \phi
\right\rangle
\]

la cual parametriza la superficie del sólido con $\phi \in[0,\pi], \theta\in[0,2\pi]$.
De esta función se obtiene
{\small
\begin{eqnarray*}
{\bf r}_{\theta} &=& \left\langle -\left(\cos\phi +1\right)\sin\theta\sin\phi,(\cos \phi +1)\cos\theta\sin\phi,0\right\rangle \\
{\bf r}_{\theta } &=& \big \langle -\left(\cos\phi +1\right)\sin\theta\sin\phi,(\cos \phi +1)\cos\theta\sin\phi, 0 \big\rangle \\
{\bf r}_{\phi } &=& \big\langle \left( \cos \phi +\cos (2\phi) \right)\cos \theta , \left( \cos \phi +\cos (2\phi) \right)\sin \theta,
\left( 2\cos \phi +1\right)\sin \phi\big \rangle, \\
\end{eqnarray*}
}

Además,
{\small
\[{\bf F}({\bf r}(\theta,\phi))=(1+\cos\phi) \left\langle\sin\phi\cos\theta,\sin\phi\sin\theta,2\cos\phi\right\rangle\!,\]

\begin{multline*}
{\bf r}_{\phi}\times {\bf r}_{\theta}=\bigg\langle\left(\cos\phi+1\right)\left(\sin2\phi+\sin\phi\right)\cos\theta\sin \phi,\\
\left( \cos \phi +1\right) \left( \sin 2\phi +\sin \phi \right) \sin \theta \sin \phi ,\left( \cos \phi +1\right)
^{2}\left( 2\cos \phi -1\right) \sin \phi \bigg\rangle .
\end{multline*}
}

Por lo tanto,

\[
\mathbf{F}(\mathbf{r}(\phi,\theta)) \cdot (\mathbf{r}_{u} \times \mathbf{r}_{\phi}) =
\frac{1}{2}(\cos \phi + 1)^{3} \left( \sin 3\phi - \sin 2\phi + 3\sin \phi \right)\!.
\]

Entonces, el flujo de \(\mathbf{F}\) sobre la superficie en mención es:
{\small
\[
\int_{0}^{\pi} \int_{0}^{2\pi}
\frac{1}{2}(\cos \phi + 1)^{3} \left( \sin 3\phi - \sin 2\phi + 3\sin \phi \right)
\, d\theta\, d\phi = \frac{32}{3}\pi.
\]
}

Como ${\rm div\,}{\bf F} =4$, es claro que el flujo corresponde a cuatro veces el volumen de la región, el cual se obtuvo en la fórmula \ref{ejercicio5}.

%\clearpage

\item Hallar el flujo neto del campo vectorial ${\bf F}(x,y,z) = \la x,y,z \ra $
sobre la superficie del sólido acotado por el paraboloide $4y=x^2+z^2$ y la esfera $x^2+y^2+z^2=12$.

Este sólido se consideró en la página \pageref{ejercicio8a}, la intersección de las superficies mencionadas y proyectada en el plano $xz$ se determina por la ecuación $x^2+z^2=8$.

La figura se obtiene mediante la ejecución de las siguientes instrucciones:

\begin{tcolorbox}[title=\bf Instrucción \verb"Mathematica"]\begin{Verbatim}[formatcom=\letraverbatim]
a=RegionPlot3D[4*y>x^2+z^2&&x^2+y^2+z^2<12,{x,-3,3},
{y,0,4},{z,-3,3},PlotPoints->90,Mesh->False];
b=SliceVectorPlot3D[{x,y,z},{x^2+y^2+z^2==12,
4*y==x^2+z^2},{x,-3,3},{y,0,3.7},{z,-3,3},
RegionFunction->Function[{x,y,z},x^2+z^2<=8],
VectorPoints->16,PlotStyle->Opacity[0],
VectorStyle->Black,VectorScale->{0.15,Scaled[0.25],
Automatic}];
Show[a,b]
\end{Verbatim}
\end{tcolorbox}

\begin{figure}[H]
\centering
\caption{Campo vectorial ${\bf F}(x,y,z)=\la x,y,z\ra$ sobre $4y=x^2+z^2$ \& $x^2+y^2+z^2=12$}
\includegraphics[width=0.4\textwidth]{Img/T_15c.pdf} \hspace{5mm}
\includegraphics[width=0.4\textwidth]{Img/T_15d.pdf}
\fuentepropia
\end{figure}

Con el propósito de parametrizar la superficie mencionada antes, es posible asignar
$x=\sqrt{8}u\cos v$, $z=\sqrt{8}u\sin v$, en donde $u\in [0,1]$ y $v \in [0,2\pi]$.
Es por esta razón que al sustituir en las ecuaciones respectivas utilizaremos las funciones
${\bf r}(u,v) = \langle\sqrt{8}u \cos v ,2u^2 ,\sqrt{8}u\sin v \rangle $, como parametrización del paraboloide y
${\bf r}(u,v) = \langle  \sqrt{8}u \cos v ,\sqrt{12-8u^2} ,\sqrt{8}u\sin v \rangle$, en el caso de la esfera.
Los cálculos se harán por separado.

\subsubsection{Flujo sobre la superficie de la esfera}

En este caso, con la parametrización escogida se obtiene
\begin{eqnarray*}
{\bf r}_u &=& \left \langle 2 \sqrt{2}\cos(v), -\frac{4u}{\sqrt{3-2u^2}},2\sqrt{2}\sin(v) \right \rangle \\
{\bf r}_v &=&\langle -2\sqrt{2}u\sin(v),0,2\sqrt{2}u\cos(v) \rangle \\
{\bf r}_v \times {\bf r}_u &=& \left \langle \frac{8\sqrt{2}u^2 \cos(v)}{\sqrt{3-2u^2}},8u,\frac{8\sqrt{2}u^2 \sin(v)}{\sqrt{3-2u^2}} \right\rangle\!.
\end{eqnarray*}

Por lo tanto, ${\bf F}({\bf r}(u,v)) = \langle \sqrt{8}u \cos v ,\sqrt{12-8u^2} ,\sqrt{8}u\sin v \rangle$ y
${\bf F}({\bf r}(u,v))\cdot({\bf r}_v \times {\bf r}_u)=\dfrac{48u}{\sqrt{3-2u^2}}$.
El flujo de ${\bf F}$ sobre la esfera está dado por
\begin{equation}\label{flujo12}
\int_0^1\int_0^{2\pi}\frac{48u}{\sqrt{3-2u^2}} dvdu= 48\pi (\sqrt{3}-1)
\end{equation}

Este resultado también se obtiene ejecutando  las siguientes instrucciones:

\begin{tcolorbox}[title=\bf Instrucción \verb"Mathematica"]\begin{Verbatim}[formatcom=\letraverbatim]
r={x,y,z}={Sqrt[8]*u*Cos[v],Sqrt[12-8*u^2],
Sqrt[8]*u*Sin[v]};
A=Simplify[Dot[{x,y,z},Cross[D[r,v],D[r,u]]]]
Integrate[A,{u,0,1},{v,0,2*Pi}]
\end{Verbatim}
\end{tcolorbox}

\subsubsection{Flujo sobre la superficie del paraboloide}

Con la función que parametriza el paraboloide, se obtiene

\begin{eqnarray*}
{\bf r}_u &=& \left \langle 2 \sqrt{2}\cos(v), 4u,2\sqrt{2}\sin(v) \right \rangle \\
{\bf r}_v &=&\langle -2\sqrt{2}u\sin(v),0,2\sqrt{2}u\cos(v) \rangle \\
{\bf r}_u \times {\bf r}_v &=& \left \langle 2\sqrt{2}u^2 \cos(v),-8u,2\sqrt{2}u^2 \sin(v) \right\rangle
\end{eqnarray*}

Luego, ${\bf F}({\bf r}(u,v)) = \langle \sqrt{8}u \cos v ,2u^2 ,\sqrt{8}u\sin v \rangle$ y
${\bf F}({\bf r}(u,v))\cdot({\bf r}_u \times {\bf r}_v)=16u^3 $.
El flujo de ${\bf F}$ sobre el paraboloide está determinado por
\begin{equation}\label{flujo11}
\int_0^1\int_0^{2\pi } 16 u^3 dvdu= 8\pi.
\end{equation}
El cálculo anterior se obtiene ejecutando las siguientes instrucciones:

\begin{tcolorbox}[title=\bf Instrucción \verb"Mathematica"]\begin{Verbatim}[formatcom=\letraverbatim]
r={x,y,z}={Sqrt[8]*u*Cos[v],2*u^2,Sqrt[8]*u*Sin[v]};
Integrate[Dot[{x,y,z},Cross[D[r,u],D[r,v]]],{u,0,1},
{v,0,2*Pi}]
\end{Verbatim}
\end{tcolorbox}

El flujo neto sobre la superficie del sólido es la suma de los valores obtenidos en \ref{flujo11}
y \ref{flujo12}. En este caso ${\rm div\,}{\bf F} =3$, el valor obtenido es tres veces el volumen obtenido en la expresión \ref{ejercicio8}.

\item Hallar el flujo neto del campo vectorial ${\bf F}(x,y,z)=\la x-2,z,-y \ra $
sobre la superficie del sólido acotado por el semicono
$x=\sqrt{y^2+z^2}$ y la esfera con ecuación $y^2+z^2+4x=12$.

\begin{figure}[H]
\centering
\caption{Campo vectorial ${\bf F}(x,y,z)=\la x-2,z,-y \ra$ sobre una superficie}
\includegraphics[width=0.4\textwidth]{Img/T_17d.pdf} \hspace{5mm}
\includegraphics[width=0.4\textwidth]{Img/T_17e.pdf}

\includegraphics[width=0.4\textwidth]{Img/T_17c.pdf}
\fuentepropia
\end{figure}

Las figuras se obtienen al ejecutar las siguientes instrucciones:

\begin{tcolorbox}[title=\bf Instrucción \verb"Mathematica"]\begin{Verbatim}[formatcom=\letraverbatim]
a=RegionPlot3D[x>Sqrt[y^2+z^2]&&y^2+z^2+4*x<12,
{x,0,4},{y,-3,3},{z,-3,3},PlotPoints->90,
Mesh->False];
b=SliceVectorPlot3D[{x-2,z,-y},{x==Sqrt[y^2+z^2],
y^2+z^2+4*x==12},{x,0,3.5},{y,-2,2},{z,-2,2},
VectorPoints->18,RegionFunction->Function[{x,y,z},
y^2+z^2<=4],VectorStyle->Black,VectorScale->{0.16,
Scaled[0.24],Automatic},PlotStyle->Opacity[0]];
Show[a,b]
\end{Verbatim}
\end{tcolorbox}

La intersección de estas dos superficies se consigue eliminando $y^{2}+z^{2}$ de las dos ecuaciones, lo que se reduce a $x^{2}+4x=12$ e implica que $x=2$.
La proyección en el plano $yz$ se determina por la ecuación $y^{2}+z^{2}=4$. Esta región con su interior se describe asignando $y=u\cos v$, $z=u\sin v$,
en donde $u\in \lbrack 0,2]$, $v\in \lbrack 0,2\pi ]$.
Sobre el cono se satis\-face que $x=\sqrt{y^{2}+z^{2}}=u$ y sobre el paraboloide se cumple que
$x=\frac{1}{4}\left( 12-y^{2}-z^{2}\right) =\frac{1}{4}\left(12-u^{2}\right) = 3-\frac{1}{4}u^{2}$,

Por lo descrito anteriormente, resulta apropiado escoger como parametrizaciones de la porción del paraboloide y del cono a las siguientes funciones, respectivamente.

\begin{eqnarray*}
{\bf r}\left( u,v\right) &=&\left\langle 3-\frac{1}{4}u^{2},\cos v,u\sin v\right\rangle\\
{\bf r}( u,v) &=&\left\langle u,\cos v,u\sin v\right\rangle\!.
\end{eqnarray*}
con $u\in [0,2]$, $v\in [0,2\pi ]$.
Calcularemos el flujo sobre cada una de las regiones componentes.
\subsubsection{Flujo sobre la porción del cono}

Es claro que \[{\bf F}({\bf r}(u,v))=\left\langle 1-\frac{1}{4}u^{2},u\sin v, -\cos v\right\rangle\] y con la parametrización escogida se sigue que
\begin{eqnarray*}
{\bf r}_{u}&=&\left\langle 1,0,\sin v\right\rangle \\
{\bf r}_{v}&=&\left\langle 0,-\sin v,u\cos v\right\rangle \\
{\bf r}_{u}\times {\bf r}_{v}&=& \left\langle \sin ^{2}v,-u\cos v,-\sin v\right\rangle\!.
\end{eqnarray*}
Entonces, ${\bf F}({\bf r}(u,v)) \cdot ({\bf r}_{u}\times {\bf r}_{v}) = u(2-u)$.
Con lo cual, el flujo está dado por
\begin{equation}\label{flujo15}
\int_0^2\int_0^{2\pi} u(2-u) du dv = \frac{8\pi}{3}
\end{equation}

\subsubsection{Flujo sobre la porción del paraboloide}

Con la parametrización escogida se sigue que
\[{\bf F}({\bf r}(u,v))=\left\langle u-2,u\sin v,-\cos v \right\rangle\] y de la función ${\bf r}$,
se obtienen los siguientes vectores.
\begin{eqnarray*}
{\bf r}_{u}&=&\left\langle -\frac{1}{2}u,0,\sin v\right\rangle \\
{\bf r}_{v}&=&\left\langle 0,-\sin v,u\cos v\right\rangle \\
{\bf r}_{u}\times {\bf r}_{v}&=&\left\langle \sin ^{2}v,\frac{1}{2}u^{2}\cos v,\frac{1}{2}u\sin v\right\rangle\!.
\end{eqnarray*}
Entonces, ${\bf F}({\bf r}(u,v)) \cdot ({\bf r}_{u}\times {\bf r}_{v}) = \frac{1}{4}u(u^2-4)$, por lo tanto, el flujo está dado por
\begin{equation}\label{flujo16}
\int_0^2\int_0^{2\pi}\frac{1}{4}\left (u-\frac{u^3}{4} \right ) du dv = 2\pi
\end{equation}
Sumando los resultados obtenidos en \ref{flujo15} y \ref{flujo16}, se obtiene el flujo neto sobre el sólido que se considera. Este valor coincide con el volumen hallado en la página \pageref{ejercicio10}, pues en este caso ${\rm div}\,{\bf F}=1$.

\item Hallar el flujo neto del campo vectorial ${\bf F}(x,y,z) = \left\langle yz,y-xz,xy\right\rangle $
sobre la superficie del sólido acotado por el paraboloide $x+3=y^2+z^2$
y el paraboloide hiperbólico $y^2-z^2=9x-3$.

La figura de esta región se produce con las siguientes instrucciones:

\begin{tcolorbox}[title=\bf Instrucción \verb"Mathematica"]\begin{Verbatim}[formatcom=\letraverbatim]
a=RegionPlot3D[x+3>y^2+z^2&&y^2-z^2>9*x-3,{x,-3,2},
{y,-2,2},{z,-2,2},PlotPoints->90,Mesh->False,
BoxRatios->{1,1,1}];
b=SliceVectorPlot3D[{y*z,y-x*z,x*y},{y^2-z^2==9*x-3,
x+3==y^2+z^2},{x,-3,1.2},{y,-2,2},{z,-2,2},
VectorPoints->20,VectorStyle->Black,
RegionFunction->Function[{x,y,z},4*y^2+5*z^2<=15],
VectorScale->{0.18,Scaled[0.3],Automatic},
PlotStyle->Opacity[0]];
Show[a,b]
\end{Verbatim}
\end{tcolorbox}

\begin{figure}[H]
\centering
\caption{${\bf F}(x,y,z)=\la yz,y-xz,xy\ra $ sobre $x+3=y^2+z^2$ \& $y^2-z^2=9x-3$}
\includegraphics[width=0.4\textwidth]{Img/T_16d.pdf} \hspace{5mm}
\includegraphics[width=0.4\textwidth]{Img/T_16e.pdf}

\includegraphics[width=0.4\textwidth]{Img/T_16f.pdf}
\fuentepropia
\end{figure}

El paraboloide y el paraboloide hiperbólico se intersecan en una curva cuya proyección en el plano $yz$, está dada por
$4y^{2}+5z^{2}=15$, esta ecuación representa una elipse.
Haciendo $y=\frac{\sqrt{15}}{2}u\cos v$, $z=\sqrt{3}u\sin v$ se des\-cribe esta región. También se cumplen las siguientes relaciones para el paraboloide.

\begin{eqnarray*}
x&=&\left( \frac{\sqrt{15}}{2}u\cos v\right)^{2}+\left(\sqrt{3}u\sin v\right)^{2}-3 =\frac{3}{8}u^{2}\cos 2v
+\frac{27}{8}u^{2}-3
\end{eqnarray*}
Asimismo el paraboloide hiperbólico
{\small
\begin{eqnarray*}
x&=&\frac{1}{9}\left( \left( \frac{\sqrt{15}}{2}u\cos v\right) ^{2}-\left( \sqrt{3}u\sin v\right) ^{2}+3\right)
= \frac{3}{8}u^{2}\cos 2v+\frac{1}{24}u^{2}+\frac{1}{3}.
\end{eqnarray*}
}

El flujo neto se hará por separado sobre las dos regiones que acotan el sólido.

\subsubsection{Flujo sobre la porción del paraboloide}
Por lo mencionado antes, resulta natural escoger la siguiente parametri\-zación para el paraboloide.

{\small
\[{\bf r}\left( u,v\right) = \langle \frac{3}{8}u^{2}\cos 2v+\frac{27}{8}u^{2}-3, \sqrt{\frac{15}{4}}u\cos v,\sqrt{3}u\sin v \rangle,\]}

con $u \in [0,1]$, $v \in [0,2\pi]$. De esta función se obtienen los siguientes vectores:

{\small
\begin{eqnarray*}
{\bf r}_{u}&=&\left\langle \frac{27}{4}u+\frac{3}{4}u\cos 2v,\frac{1}{2}\sqrt{15}\cos v,\sqrt{3}\sin v\right\rangle \\
{\bf r}_{v}&=&\left\langle -\frac{3}{4}u^{2}\sin 2v,-\frac{1}{2}\sqrt{15}u\sin v,\sqrt{3}u\cos v\right\rangle \\
{\bf r}_{u}\times {\bf r}_{v}&=&\left\langle -\frac{3}{2}\sqrt{5}u,\frac{15}{2}\sqrt{3}u^{2}\cos v, 3\sqrt{3}\sqrt{5}u^{2}\sin v\right\rangle
\end{eqnarray*}
}

Asimismo
\begin{multline*}
{\bf F}({\bf r}(u,v))=\bigg \langle \frac{3}{4}\sqrt{5}u^2\sin(2v),
\frac{1}{16}\sqrt{3}u(8\sqrt{5}\cos(v)+(48-51u^2)\sin(v) \\
3u^2\sin(3v)), \frac{3}{16}\sqrt{15}u\cos(v)(u^2(9+\cos(2v))-8)\bigg \rangle,
\end{multline*}
con lo cual
${\bf F}({\bf r}(u,v)) \cdot \left({\bf r}_{u}\times {\bf r}_{v}\right)=\frac{45}{4} u^3\cos(v)(\sqrt{5}\cos(v)-\sin(v))$.
Por lo tanto, el flujo sobre la porción del paraboloide que se considera es:
\begin{equation} \label{flujo13}
\int_{0}^{1}\int_{0}^{2\pi}\frac{45}{4} u^3\cos(v)(\sqrt{5}\cos(v)-\sin(v)) dvdu = \frac{45 \sqrt{5}}{16} \pi
\end{equation}
El cálculo anterior se puede hacer en \verb"Mathematica" ejecutando las siguien\-tes instrucciones:

\begin{tcolorbox}[title=\bf Instrucción \verb"Mathematica"]\begin{Verbatim}[formatcom=\letraverbatim]
{y,z,x}={Sqrt[15]/2*u*Cos[v],Sqrt[3]*u*Sin[v],
y^2+z^2-3};
F=TrigFactor[{y*z,y-x*z,x*y}]
Integrate[Dot[F,Cross[D[{x,y,z},v],D[{x,y,z},u]]],
{u,0,1},{v,0,2*Pi}]
\end{Verbatim}
\end{tcolorbox}

\subsubsection{Flujo sobre el paraboloide hiperbólico}

En este caso, la parametrización escogida es la siguiente:
\[
{\bf r}(u,v) =
\left\langle
\frac{3}{8}u^{2}\cos 2v + \frac{1}{24}u^{2} + \frac{1}{3},\;
\frac{\sqrt{15}}{2}u\cos v,\;
\sqrt{3}u\sin v
\right\rangle\!,
\]
en donde \( u \in [0,1],\; v \in [0,2\pi] \). También se puede ver que
\[
\mathbf{F}(\mathbf{r}(u,v)) =
\left\langle
\frac{3}{2}\sqrt{5}\,u,\
-\frac{5u^2\cos(v)}{2\sqrt{3}},\
\sqrt{\frac{5}{3}}\, u^2 \sin(v)
\right\rangle\!.
\]
Con la función ${\bf r}$ obtenemos los siguientes vectores:
\begin{eqnarray*}
{\bf r}_{u} &=&\left\langle \frac{1}{12}u+\frac{3}{4}u\cos 2v,\frac{1}{2}\sqrt{15}\cos v,\sqrt{3}\sin v\right\rangle\\
{\bf r}_{v}&=&\left\langle -\frac{3}{4}u^{2}\sin 2v,-\frac{1}{2}\sqrt{15}u\sin v,\sqrt{3}u\cos v\right\rangle \\
{\bf r}_{u}\times {\bf r}_{v}&=& \la \frac{3}{2}\sqrt{5}u,-\frac{5}{6}\sqrt{3} u^{2}\cos v,\frac{1}{3}\sqrt{15}u^{2}\sin v\ra,
\end{eqnarray*}
entonces,
{\small
\[{\bf F}({\bf r}(u,v)) \cdot \left({\bf r}_{u}\times {\bf r}_{v}\right)=\frac{5u^3}{24}\cos(v)(6\sqrt{5}\cos(v)-(62+u^2(1+9\cos(2v)))\sin(v)).\]
}
Por lo tanto, sobre la porción del paraboloide hiperbólico el flujo de ${\bf F}$ es
\begin{equation} \label{flujo14}
\int_0^1\int_0^{2\pi}{\bf F}({\bf r}(u,v)) \cdot \left({\bf r}_{u}\times{\bf r}_{v}\right)dvdu=-\frac{5 \sqrt{5}}{16}\pi
\end{equation}
Con la ayuda de las siguientes instrucciones, se puede obtener el cálculo anterior.

\begin{tcolorbox}[title=\bf Instrucción \verb"Mathematica"]\begin{Verbatim}[formatcom=\letraverbatim]
{y,z,x}={Sqrt[15]/2*u*Cos[v],Sqrt[3]*u*Sin[v],
(y^2-z^2+3)/9};
Integrate[Dot[{y*z,y-x*z,x*y},Cross[D[{x,y,z},u],
D[{x,y,z},v]]],{u,0,1},{v,0,2*Pi}]
\end{Verbatim}
\end{tcolorbox}

El flujo neto sobre la superficie del sólido es la suma de los valores obtenidos en \ref{flujo13} y
\ref{flujo14}, es decir $\frac{5 \sqrt{5}}{16} \pi$ . Como ${\rm div\,}{\bf F}=1$,
este valor coincide con el volumen del sólido hallado con la expresión \ref{ejercicio9}.
\end{parrafonumerado}

\begin{ejer}
Resuelva los siguientes problemas.
\begin{parrafonumerado}

\item Calcule el flujo del campo vectorial ${\bf F}(x,y,z)= \la x-2y,3x+y,z+1\ra $ sobre la superficie del hiperboloide de un manto $x^2 + y^2 - z^2 = 1$, con $-2 \leq z \leq 2$.

\begin{figure}[H]
\centering
\caption{${\bf F}(x,y,z)=\la x-2y,3x+y,z+1\ra$ sobre un hiperboloide}
\includegraphics[width=0.4\textwidth]{Img/T_35a.pdf} \hspace{5mm}
\includegraphics[width=0.4\textwidth]{Img/T_35b.pdf}
\end{figure}

\item Considere la superficie sobre el paraboloide  $z=x^2-2x+y^2$ interior al cilindro cardioidico $r=1+\cos t$.

\begin{figure}[H]
\centering
\caption{Superficie sobre un paraboloide}
\includegraphics[width=0.4\textwidth]{Img/T_36a.pdf} \hspace{5mm}
\includegraphics[width=0.4\textwidth]{Img/T_36b.pdf}
\end{figure}

Calcule el flujo del campo vectorial ${\bf F}(x,y,z)= \la x-2,y,2z\ra $ sobre la superficie del paraboloide.
$z=x^2-2x+y^2$, interior al cilindro cardioidico $r=1+\cos t$.

%

\begin{figure}[H]
\centering
\caption{Campo vectorial ${\bf F}(x,y,z)= \la x-2,y,2z\ra $ sobre un paraboloide}
\includegraphics[width=0.4\textwidth]{Img/T_36c.pdf} \hspace{5mm}
\includegraphics[width=0.4\textwidth]{Img/T_36d.pdf}
\end{figure}

Evalúe $\ds \int_S\rot{F}\cdot d{\bf S}$ y $\ds\int_C{\bf F}\cdot d{\bf r}$; en donde $S$
es la superficie en mención y $C$ es su frontera.

\begin{comment}
x=Cos[t]*(1+Cos[t]);
y=Sin[t]*(1+Cos[t]);
z=x^2-2*x+y^2;
a=ParametricPlot3D[{x,y,z},{t,0,2*Pi},PlotStyle->Red];

Show[SliceVectorPlot3D[{x-2,y,2z},{z==x^2-2*x+y^2},{x,-1/4,2},{y,-3/2,3/2},
{z,-1,1.1},VectorStyle->Black,RegionFunction->Function[{x,y,z},
(x^2+y^2-x)^2<=x^2+y^2],VectorScale->{0.1,Scaled[0.2],Automatic},VectorPoints->17],a]
\end{comment}

\item Calcule el flujo del campo vectorial ${\bf F}(x,y,z)= \la z-y,y,z^2 \ra $ sobre la superficie del paraboloide hiperbólico $9z=x^2-y^2$, interior al cilindro $x^2+y^2 \leq 16$ y exterior al cilindro cardioidico $r=1+\cos t$.

\begin{figure}[H]
\centering
\caption{Campo vectorial ${\bf F}(x,y,z)= \la z-y,y,z^2\ra $ sobre una silla de montar}
\includegraphics[width=0.4\textwidth]{Img/T_37a.pdf} \hspace{5mm}
\includegraphics[width=0.4\textwidth]{Img/T_37b.pdf}
\end{figure}

\begin{comment}
x=Cos[t]*(1+Cos[t]);
y=Sin[t]*(1+Cos[t]);
z=(x^2-y^2)/9;
a=ParametricPlot3D[{x,y,z},{t,0,2*Pi},
PlotStyle->{Thickness->0.005,Blend[{Blue,Black},1/3]}];
b=ParametricPlot3D[{4*Cos[t],4*Sin[t],16*Cos[2*t]/9},{t,0,
2*Pi},PlotStyle->{Thickness->0.005,
Blend[{Blue,Black},1/3]}];
datos=Table[{z-y,y,z^2},{x,-4,4},{y,-4,4},{z,-2,2}];
Show[ListSliceVectorPlot3D[datos,
ImplicitRegion[
9*z==x^2-y^2&&(x^2+y^2-x)^2>=x^2+y^2&&
x^2+y^2<=16,{x,y,z}],
DataRange->{{-4,4},{-4,4},{-2,2}},VectorPoints->15,
VectorStyle->{Thin,Black,Thick},
BoxRatios->{1,1,2/3}],a,b,PlotPoints->60]
\end{comment}

\end{parrafonumerado}
\end{ejer} 

\chapter{Anexos}
 \section{Funciones hiperbólicas} \input{anexo1} 
\section{Sumatorias} \input{anexo2} 
\section{Sólidos de revolución en coordenadas paramétricas} \input{anexo4} 
\section{Sistemas de coordenadas curvilíneas} \input{anexo3}
%\addcontentsline{toc}{chapter}{Índice alfabético}
\printindex

\printbibliography[heading=bibintoc]
\printbibliography


\chapter*{\lhseries Sobre los autores}
\addcontentsline{toc}{chapter}{Sobre los autores}

{\large\medseries\noindent Juan Nepomuceno Zambrano Caviedes}
\noindent
Licenciado en Matemáticas y Física y Especialista en Educación Matemática de la Universidad Surcolombiana,  Magister en matemáticas Aplicadas de la Universidad EAFIT. Profesor Asociado de la Universidad Distrital con más de treinta años de experiencia docente universitaria.
\faEnvelope~\url{jzambranoc@udistrital.edu.co}\\
%\textcolor{orcidlogocol}{\aiOrcid}~\url{https://orcid.org/0000-0001-8378-4215}
\\

{\large\medseries\noindent Efrén Ricardo Baquero Torres}
\noindent 
Matemático con Magister en Ciencias Matemáticas de La Universidad Nacional de Colombia. En más de dos décadas ha orientado cursos universitarios de Matemáticas en la Universidad Escuela Colombiana de Ingeniría Julio Garavito y en La Universidad Distrital Francisco José de Caldas, tiene un especial 
interés en el uso de nuevas tecnologías en su enseñanza.

\faEnvelope~\url{efren.baquero@escuelaing.edu.co}\\
%\textcolor{orcidlogocol}{\aiOrcid}~\url{https://orcid.org/0000-0001-8378-4215}
%%%%%%%%%%%%%%%%%%%%%%%%%%%
%\escolofon{el mes}{año}
%\escolofon{septiembre}{2024}
%\tableofcontents \mainmatter	
%\include{present2}  
\end{document}
